\documentclass[11pt,openany]{book}
\setcounter{secnumdepth}{3}
\usepackage[margin=1.1in]{geometry}
% Japanese edition: build with XeLaTeX (see the Makefile)
\usepackage{xeCJK}
\setCJKmainfont[ItalicFont=IPAexGothic,BoldFont=IPAexGothic]{IPAexMincho} % \emph and bold in Gothic: IPAex has no italic
\setCJKsansfont{IPAexGothic}
\setCJKmonofont{IPAexGothic}
\xeCJKsetup{CJKecglue={\hskip 0.15em plus 0.05em minus 0.05em}}
\usepackage{titlesec}
\newcommand{\chaptersuffix}{章}
\titleformat{\chapter}[display]{\normalfont\huge\bfseries}{\ifx\chaptersuffix\empty\appendixname\thechapter\else 第\thechapter\chaptersuffix\fi}{20pt}{\Huge}
\AtBeginDocument{\renewcommand{\contentsname}{目次}\renewcommand{\figurename}{図}\renewcommand{\tablename}{表}\renewcommand{\bibname}{参考文献}\renewcommand{\appendixname}{付録~}\renewcommand{\proofname}{証明}}
\usepackage{amsmath,amssymb,amsthm}
\usepackage{graphicx}
\usepackage{tikz}
\usetikzlibrary{arrows.meta,decorations.markings,decorations.pathmorphing,decorations.pathreplacing,calc}
\input{figures/icons}
\usepackage[unicode,colorlinks=true,linkcolor=blue!60!black,citecolor=blue!60!black]{hyperref}
\usepackage{microtype}
\setlength{\emergencystretch}{3em} % long names (licences, references) in narrow lines
\usepackage{empheq}
\usepackage{array}
\usepackage{booktabs}
\usepackage{multirow}
\usepackage{longtable}
\usepackage{circuitikz}
\definecolor{keyyellow}{rgb}{1.0,0.97,0.78}
% key equations: \begin{empheq}[box=\keybox]{equation} ... \end{empheq} -- light-yellow box, number stays outside
\newcommand{\keybox}[1]{{\setlength{\fboxsep}{7pt}\colorbox{keyyellow}{#1}}}
\usepackage[most]{tcolorbox}
% key statements (propositions etc.): the same light-yellow background, breakable across pages
\newtcolorbox{keyblock}{enhanced,breakable,colback=keyyellow,colframe=keyyellow,boxrule=0pt,arc=0pt,left=6pt,right=6pt,top=2pt,bottom=2pt,boxsep=0pt}
\renewcommand{\chaptermark}[1]{\markboth{\small 第\thechapter 章\quad #1}{}}
\renewcommand{\sectionmark}[1]{\markright{\small\thesection.\ #1}}
\renewcommand{\topfraction}{0.95}
\renewcommand{\textfraction}{0.05}
\renewcommand{\floatpagefraction}{0.75}

% part pages: title at the top third, and text may follow on the same page (used for the part introductions)
\makeatletter
\renewcommand\part{\if@openright\cleardoublepage\else\clearpage\fi\thispagestyle{plain}\if@twocolumn\onecolumn\@tempswatrue\else\@tempswafalse\fi\null\vspace{0.18\textheight}\secdef\@part\@spart}
\renewcommand\@endpart{\par\vspace{3em}\if@tempswa\twocolumn\fi}
\makeatother
\newtheorem{theorem}{定理}[chapter]
\newtheorem{proposition}{命題}[chapter]
\newtheorem{lemma}{補題}[chapter]
\theoremstyle{remark}
\newtheorem{remark}{注意}[chapter]
% exercises: \begin{exercise}[\lvlB] ... \end{exercise}, numbered "Exercise 4.3"; difficulty marks
\theoremstyle{definition}
\newtheorem{exercise}{問題}[chapter]
\newcommand{\lvlA}{$\star$}
\newcommand{\lvlB}{$\star\star$}
\newcommand{\lvlC}{$\star\star\star$}
% answer to an exercise in the Answers appendix: \answer{ex:NN:k}{text}
\newcommand{\answer}[2]{\par\noindent\textbf{\ref{#1}.}~#2\par\smallskip}
\newcommand{\dd}{\mathrm{d}}
% neuron type code: first four letters of the primary mediator
\newcommand{\nt}[1]{\textsf{\textbf{#1}}}
\newcommand{\mV}{\,\text{mV}}
\newcommand{\ms}{\,\text{ms}}
\newcommand{\mM}{\,\text{mM}}
\newcommand{\um}{\,\text{\textmu m}}
\newcommand{\ion}[2]{\mathrm{#1}^{#2}}
\newcommand{\Na}{\ion{Na}{+}}
\newcommand{\K}{\ion{K}{+}}
\newcommand{\Cl}{\ion{Cl}{-}}
\newcommand{\Ca}{\ion{Ca}{2+}}
\newcommand{\conc}[1]{[\mathrm{#1}]}

\title{\Huge 20ワットのコンピュータ\\[16pt] \Large 脳はどのように計算するか：エンジニアのための定量的アプローチ}
\hypersetup{pdftitle={20ワットのコンピュータ. 脳はどのように計算するか：エンジニアのための定量的アプローチ},pdfauthor={Konstantin Kladko}}
\author{\Large コンスタンチン・クラドコ（Konstantin Kladko）}
\date{}

\begin{document}
\frontmatter
\maketitle
\thispagestyle{empty}
\vspace*{\fill}
\noindent{\small \copyright~2026 Konstantin Kladko（コンスタンチン・クラドコ）。\\[4pt]
本書の本文・図・問題・解答は、クリエイティブ・コモンズ 表示—非営利—継承 4.0 国際ライセンス（CC BY-NC-SA 4.0）の下で公開されている：\url{https://creativecommons.org/licenses/by-nc-sa/4.0/deed.ja}。著者を表示し、派生物を同じライセンスで公開する限り、非営利目的で自由に複製し、授業に使い、翻訳し、改変することができる。\\[4pt]
モデルとビルドのプログラムは MIT ライセンスの下で公開されている。\\[4pt]
ソース、モデル、最新版：\url{https://github.com/kladkogex/20-watt-computer}。}
\clearpage

\tableofcontents
\chapter*{本書の読み方}
\addcontentsline{toc}{chapter}{本書の読み方}

本書は最初から順に読む必要はない。第1--4章は共通の基礎であり、ニューロンの種類、\nt{GLUT}と\nt{GABA}が何を計算するか、そしてそれらがどんな言語で話すかを扱う。その先で本書は枝分かれする。図~\ref{fig:roadmap}の地図は、どの章がどの章に依存しているかを示す。章$A$から章$B$への矢印は、$B$で$A$の概念と式を使うという意味である。地図には主要な依存関係だけを載せた。細かな後方参照は読むうえで妨げにならない。

\begin{figure}[h]
\centering
\begin{tikzpicture}[>=Stealth,scale=0.95,
  ch/.style={draw,thick,rounded corners=3pt,align=center,font=\scriptsize,text width=1.85cm,minimum height=0.62cm,inner sep=2pt},
  base/.style={ch,fill=blue!8}, bus/.style={ch,fill=orange!14}, sum/.style={ch,fill=gray!18},
  io/.style={ch,fill=green!10}, sort/.style={ch,fill=cyan!8}, lab/.style={ch,fill=violet!10},
  dep/.style={->,thick,gray!60!black}]
% foundation
\node[base] (c1) at (0,0) {1 ニューロン\\の種類};
\node[base] (c2) at (2.3,0) {2 \nt{GLUT}};
\node[base] (c3) at (4.6,0) {3 \nt{GLUT}と\nt{GABA}};
\node[base] (c4) at (6.9,0) {4 モールス\\符号};
% broadcast channels and the synthesis
\node[bus] (c5) at (4.6,-1.5) {5 \nt{SERO}};
\node[bus] (c6) at (7.3,-1.5) {6 \nt{DOPA}};
\node[bus] (c7) at (9.2,-0.9) {7 ドーパミン\\の諸理論};
\node[bus] (c8) at (9.2,-2.1) {8 \nt{NORA}};
\node[bus] (c9) at (11.5,-2.1) {9 \nt{ACET}};
\node[sum] (c10) at (13.8,-0.9) {10 マシン全体};
% inputs and outputs
\node[io] (c12) at (2.3,-4.1) {12 認識};
\node[io] (c11) at (4.6,-4.1) {11 入力};
\node[io] (c13) at (6.9,-4.1) {13 筋肉};
\node[io] (c14) at (9.2,-3.05) {14 痛み};
\node[io] (c15) at (9.2,-4.1) {15 運動の\\学習};
\node[io] (c16) at (9.2,-5.15) {16 体の\\化学};
% lab, sorting, sleep
\node[lab] (c17) at (11.5,-4.1) {17 デバッグ};
\node[lab] (c21) at (13.8,-4.1) {21 生きた\\マシン};
\node[lab] (c22) at (13.8,-5.15) {22 DishBrain};
\node[sort] (c18) at (11.5,-5.15) {18 回路素子\\としての型};
\node[sort] (c19) at (13.8,-6.2) {19 樹状突起\\の数};
\node[bus] (c20) at (11.5,-6.2) {20 睡眠};
% dependencies
\draw[dep] (c1) -- (c2); \draw[dep] (c2) -- (c3); \draw[dep] (c3) -- (c4);
\draw[dep] (c1) -- (c5); \draw[dep] (c4) -- (c6); \draw[dep] (c5) -- (c6);
\draw[dep] (c6) -- (c7); \draw[dep] (c6) -- (c8); \draw[dep] (c8) -- (c9);
\draw[dep] (c7) -- (c10); \draw[dep] (c9) -- (c10);
\draw[dep] (c4) -- (c11); \draw[dep] (c11) -- (c12); \draw[dep] (c11) -- (c13); \draw[dep] (c5) -- (c13);
\draw[dep] (c13) -- (c14); \draw[dep] (c13) -- (c15); \draw[dep] (c13) -- (c16);
\draw[dep] (c15) -- (c17); \draw[dep] (c9) -- (c17); \draw[dep] (c17) -- (c21); \draw[dep] (c21) -- (c22);
\draw[dep] (c16) -- (c18); \draw[dep] (c18) -- (c19); \draw[dep] (c16) -- (c20);
% legend
\begin{scope}[yshift=-7.25cm]
\foreach \s/\t [count=\i] in {base/基礎, bus/チャネルとモード, sum/まとめ, io/入力と出力, sort/ニューロンの分類, lab/実験室}{
  \pgfmathtruncatemacro{\col}{mod(\i-1,3)}\pgfmathtruncatemacro{\row}{(\i-1)/3}
  \node[\s,text width=0.3cm,minimum height=0.32cm] at ({\col*4.8+0.2},{-\row*0.55}) {};
  \node[anchor=west,font=\scriptsize] at ({\col*4.8+0.55},{-\row*0.55}) {\t};}
\end{scope}
\end{tikzpicture}
\caption{章の地図。矢印は、ある章からそれに依存する章へ向かう。これ以外の章間の参照は案内であって、依存関係ではない。}
\label{fig:roadmap}
\end{figure}

\section*{本書の読み進め方}

\begin{center}
\footnotesize
\setlength{\tabcolsep}{4pt}
\renewcommand{\arraystretch}{1.2}
\begin{tabular}{>{\raggedright}p{3.0cm}>{\raggedright}p{3.9cm}>{\raggedright\arraybackslash}p{6.4cm}}
\toprule
コース & 対象 & 章の順序 \\
\midrule
1クォーターの講義 & 教員、10週 & 第1週：第1--2章、2：第3章、3：第4章、4：第5--6章、5：第7--8章、6：第9章、7：第10章、8：第11--12章、9：第13章と第15章、10：第17章 \\
AIと機械学習 & ニューラルネットワークを知っていて、脳がどう学習するかを理解したい人 & 1--4、5（流し読み）、6、7、8、9、11（流し読み）、12、13（流し読み）、15 \\
エレクトロニクスとハードウェア & 回路設計者、ニューロモルフィックチップの開発者 & 1--4、5、11、13、16、18、19、17（第9章と第15章は流し読み） \\
ニューロインターフェースと生きた計算機 & ブレイン・コンピュータ・インターフェースや生きたニューロンを載せたチップを作る人 & 1、2、4、11、13、17（第9章と第15章は流し読み）、21、22（第12章は流し読み） \\
駆け足の概観 & 週末が空いている技術者 & 1、2、3、6、10。第6章から第4--5章への参照は文脈から理解できる \\
\bottomrule
\end{tabular}
\end{center}

「流し読み」とは、章の冒頭、黄色い枠内の命題、章末のまとめの表を読めば十分という意味である。

各章の終わりには問題がある。見積もり、導出、シミュレーション、設計の問題である。簡潔な解答は付録~\ref{sec:answers}にまとめた。記号はすべて付録~\ref{sec:notation}に、各章の主要な主張の根拠となる実験は付録~\ref{sec:supplement}にある。本書の数値は\texttt{models/}フォルダにある小さなプログラムで得たものであり、実行して書き換えることもできる。

\mainmatter

\chapter{ニューロンの種類}
\label{sec:neurontypes}

\section{ブラックボックスとしてのニューロン}

860億個のニューロンが入った袋を渡され~\cite{Azevedo2009}、それを山に仕分けてほしいと頼まれたとしよう。どう仕分けるだろうか。

見知らぬICが詰まった袋を渡されたエンジニアなら、パッケージの形で仕分けたりはしない。部品の入力はいくつか、出力はいくつか、出力に何を出すか、を問うはずである。ニューロンも同じように、回路図上のブロックとして描いてみよう（図~\ref{fig:blackbox}）。

\begin{figure}[t]
\centering
\begin{tikzpicture}[>=Stealth,x=1cm,y=1cm]
% the box
\node[draw,very thick,minimum width=2.6cm,minimum height=2.6cm,fill=blue!6] (N) at (0,0) {};
\node at (0,0.55) {\small ニューロン};
\node at (0,-0.15) {$\displaystyle u=\sum_i w_i x_i$};
\node at (0,-0.8) {\small $y=\Theta(u-\theta)$};
% inputs
\foreach \k/\y/\lab in {1/1.05/{x_1},2/0.55/{x_2},3/0.05/{x_3}}{
  \draw[->,thick,blue!60!black] (-4.2,\y) -- (-1.3,\y);
  \node[left] at (-4.2,\y) {$\lab$};
  \node[above,blue!60!black] at (-2.1,\y-0.06) {\scriptsize $w_{\k}$};}
\node at (-4.45,-0.4) {$\vdots$};
\node at (-2.1,-0.4) {$\vdots$};
\draw[->,thick,blue!60!black] (-4.2,-1.05) -- (-1.3,-1.05);
\node[left] at (-4.2,-1.05) {$x_n$};
\node[above,blue!60!black] at (-2.1,-1.11) {\scriptsize $w_{n}$};
\draw[decorate,decoration={brace,amplitude=5pt},gray!60!black] (-4.9,-1.2) -- (-4.9,1.2);
\node[left,align=right,gray!40!black] at (-5.1,0) {\small 入力\\[-2pt]\small $n\sim10^3$--$10^5$};
% single output wire, then fan-out
\draw[very thick,red!70!black] (1.3,0) -- (3.2,0);
\foreach \x in {1.6,2.2,2.34,2.9}{\draw[thick,red!70!black] (\x,0) -- (\x,0.4);}
\node[above right,font=\tiny,gray!30!black,inner sep=1pt] at (1.42,0.45) {カチッ\dots\ カチッカチッ\dots};
\node[below,red!70!black,align=center] at (2.25,-0.05) {\scriptsize 出力は$y$ひとつ};
\fill[red!70!black] (3.2,0) circle (0.06);
\foreach \y in {1.3,0.65,0,-0.65,-1.3}{
  \draw[->,thick,red!70!black] (3.2,0) -- (5.0,\y);}
\draw[decorate,decoration={brace,amplitude=5pt},gray!60!black] (5.3,1.45) -- (5.3,-1.45);
\node[right,align=left,gray!40!black] at (5.5,0) {\small $y$のコピーを\\[-2pt]\small $10^3$--$10^4$個の\\[-2pt]\small 別のニューロンへ};
\end{tikzpicture}
\caption{エンジニアの目で見たニューロン。入力$x_i$が多数あり、それぞれに重み$w_i$がつく。内部では重みつき和$u$を計算し、閾値$\theta$と比較する（$\Theta$は階段関数で、引数が正なら$1$、そうでなければ$0$）。出力$y$はひとつで、スパイクの列であり、分岐して数千の受け手に複製される。}
\label{fig:blackbox}
\end{figure}

ブロックの出力は数値ではなく、短い電圧パルスの列である。「カチッ」、休止、「カチッカチッ」、長い休止。どのスパイクも高さは同じなので、情報はすべてスパイクが\emph{いつ}起きるかにある。ブロックの内部には、第一次の粗い近似として、加算器とコンパレータがある。入力が重みつきで足し合わされ、和が閾値を超えるとブロックはスパイクを出す。次章では、これを連続時間で正しく動くモデルに置き換える。

出力はひとつだが、配線は枝分かれし、どの枝も\emph{まったく同じ}信号のコピーを運ぶ。エレクトロニクスでいうファンアウト（fan-out）である。大脳皮質のニューロンではファンアウトは数千のオーダーである。ICの論理ゲートは、ふつう数個のゲートを駆動するだけである。

では、出力配線は受け手に\emph{何を}届けるのか。スパイクは各枝の末端まで走り、そこで化学信号に変わる。枝は細胞間の狭い隙間に特定の物質、\emph{神経伝達物質}をひとつまみ放出し、それが受け手の入力電流を変える。これで仕分けの基準が決まる。\emph{ニューロンの出力がどの神経伝達物質を放出するか}である。

\section{出力ひとつに信号の型ひとつ}

この仕分けが意味をもつのは、各ニューロンの出力の型がひとつに限られる場合だけである。実際にそうなのか。実験によれば、ほぼそうである。

\begin{keyblock}
\begin{proposition}[ニューロン1個に出力の型1種類]
\label{prop:dale}
ほとんどすべてのニューロンは主要な神経伝達物質を1種類だけもち、出力のすべての枝でそれを放出する。したがって、ニューロンの型はその主要な神経伝達物質で決まる~\cite{Strata1999}。
\end{proposition}
\end{keyblock}

ニューロンの型は、\emph{主要な神経伝達物質の英語名の最初の4文字}で表すことにする。たとえば、主要な神経伝達物質がグルタミン酸（glutamate）であるニューロンは\nt{GLUT}ニューロンである。すべての型を表~\ref{tab:types}にまとめた。

\begin{table}[t]
\centering
\footnotesize
\setlength{\tabcolsep}{3pt}
\renewcommand{\arraystretch}{1.15}
\begin{tabular}{>{\raggedright}p{1.0cm}>{\raggedright}p{2.05cm}>{\raggedright}p{1.75cm}>{\raggedright}p{1.6cm}>{\centering}p{0.7cm}>{\raggedright}p{1.55cm}>{\raggedright}p{1.8cm}>{\raggedright\arraybackslash}p{3.2cm}}
\toprule
型 & 主要な神経伝達物質 & バス & 速さ & 符号 & 脳内のニューロン数 & ニューロン1個の出力数 & 計算での役割 \\
\midrule
\nt{GLUT} & グルタミン酸 & アドレス & 速い・遅い & $+$ & $\sim8\cdot10^{10}$ & $\sim10^4$ & 主要な正の重み \\
\nt{GABA} & GABA & アドレス & 速い・遅い & $-$ & $\sim3\cdot10^{9}$ & $10^3$--$10^4$ & 主要な負の重み \\
\nt{GLYC} & グリシン & アドレス & 速い & $-$ & $10^6$--$10^7$ & $10^2$--$10^3$ & 脊髄と脳幹での負の重み。グルタミン酸受容体の一種の第二の鍵 \\
\nt{ACET} & アセチルコリン & 両方 & 速い・遅い & $\pm$ & $\sim10^6$ & 筋肉：$10$--$10^3$、脳：$\sim10^5$ & 筋肉への出力。脳内では学習率（仮説） \\
\nt{DOPA} & ドーパミン & ブロード\newline キャスト & 遅い & $\pm$ & $\sim5\cdot10^{5}$ & $10^5$--$10^6$ & 報酬予測誤差$\delta$ \\
\nt{SERO} & セロトニン & ブロード\newline キャスト & 遅い（受容体の一種は速い） & $\pm$ & $\sim3\cdot10^{5}$ & $\sim10^5$ & 計画の時間的視野（仮説） \\
\nt{HIST} & ヒスタミン & ブロード\newline キャスト & 遅い & $+$ & $\sim6\cdot10^{4}$ & 推定なし & 全体への「覚醒せよ」信号 \\
\nt{NORA} & ノルアドレナリン & ブロード\newline キャスト & 遅い & $\pm$ & $\sim5\cdot10^{4}$ & $\sim10^5$ & ゲイン（仮説） \\
\nt{ADRE} & アドレナリン & ブロード\newline キャスト & 遅い & $\pm$ & $\sim10^4$ & 推定なし & ニューロンは少数。脳内での役割は不明 \\
\nt{ASPA} & アスパラギン酸 & アドレス & 速い & $+$ & 推定なし & 推定なし & 弱い正の重み。役割には異論がある \\
\bottomrule
\end{tabular}
\caption{ニューロンの型。脳内のニューロン数の多い順。\emph{バス}：アドレスバスでは、神経伝達物質は特定の受け手への狭い隙間に放出される。ブロードキャストでは、少数の発信元ニューロンから広い領域へ届く（\ref{sec:buses}節）。\emph{速さ}：速いはイオンチャネル型受容体経由でミリ秒、遅いは代謝型受容体経由で数百ミリ秒以上。\emph{符号}は通常の符号である。最終的に符号を決めるのは受け手であり（\ref{sec:receiver}節）、$\pm$は両方の符号がありうることを示す。\emph{脳内のニューロン数}は、ヒトの脳全体に含まれるその型のニューロンの数のオーダーである。ニューロンの総数は約$8.6\cdot10^{10}$である~\cite{Azevedo2009}。\nt{GLUT}ニューロンのほとんど、約$7\cdot10^{10}$個は小脳の小さな入力ニューロンである。\nt{GLUT}と\nt{GABA}の数は、この計数と大脳皮質における\nt{GABA}ニューロンの割合から求めた~\cite{Hendry1987}。少数派の型で直接数えられているのは\nt{HIST}~\cite{Airaksinen1991}と、\nt{DOPA}ニューロンの二つの群のうち主要なほう~\cite{Pakkenberg1991}だけなので、\nt{DOPA}の値は下限である。\nt{SERO}は二つの部分的な計数の和である~\cite{Baker1990,Baker1991}。\nt{NORA}、\nt{GLYC}、\nt{ACET}、\nt{ADRE}の値は直接の計数がない粗いオーダー推定である。\emph{ニューロン1個の出力数}は、1個のニューロンがもつシナプス終末、つまり出力配線の枝上にある神経伝達物質の放出点の数のオーダーである（図~\ref{fig:blackbox}）。ブロードキャスト型では終末を直接数えた例はない。\nt{DOPA}の値は、1本の出力配線の枝分かれが標的領域のどれだけを覆うかから導いたものであり~\cite{Matsuda2009}、それ以外の型の推定はさらに粗い。\nt{ACET}には二つの場合がある。筋肉を制御するニューロンと、脳内のブロードキャスト型ニューロンである。}
\label{tab:types}
\end{table}

% the pie is a table-type float with a figure caption, so it can never float ahead of Table~\ref{tab:types}
\begin{table}[tb]
\makeatletter\def\@captype{figure}\makeatother
\centering
\definecolor{vizglut}{HTML}{2A78D6}
\definecolor{vizgaba}{HTML}{E34948}
\definecolor{vizrest}{HTML}{8A8986}
\begin{tikzpicture}[>=Stealth]
% pie: GLUT 96.4 %, GABA 3.6 % (13.0 degrees), the rest 0.006 % (a hairline at 0 degrees)
\fill[vizglut] (0,0) -- (13:1.9) arc (13:360:1.9) -- cycle;
\fill[vizgaba] (0,0) -- (0:1.9) arc (0:13:1.9) -- cycle;
\draw[white,line width=1pt] (0,0) -- (13:1.9);
\draw[vizrest,line width=1.2pt] (0,0) -- (0:1.9);
\node[white,align=center,font=\small] at (190:0.95) {\nt{GLUT}\\$96.4\,\%$};
\draw[gray!60!black] (6.5:1.75) -- (2.05,2.1);
\node[anchor=west,font=\small] at (2.05,2.1) {\nt{GABA} $3.6\,\%$};
% zoom box for the remaining types
\draw[densely dashed,vizrest] (1.9,0) -- (3.1,1.65);
\draw[densely dashed,vizrest] (1.9,0) -- (3.1,-2.2);
\draw[vizrest,rounded corners=3pt] (3.1,-2.2) rectangle (10.1,1.65);
\node[anchor=west,font=\scriptsize] at (3.2,1.4) {その他すべての合計：$\approx0.006\,\%$、対数目盛};
\foreach \code/\lg/\y/\val/\lx in {GLYC/6.477/1.0/{$10^6$--$10^7$}/4.05,ACET/6.0/0.6/{$\sim10^6$}/3.0,DOPA/5.699/0.2/{$\sim5\cdot10^5$}/2.699,SERO/5.477/-0.2/{$\sim3\cdot10^5$}/2.477,HIST/4.778/-0.6/{$\sim6\cdot10^4$}/1.778,NORA/4.699/-1.0/{$\sim5\cdot10^4$}/1.699,ADRE/4.0/-1.4/{$\sim10^4$}/1.0}{
  \node[anchor=east,font=\scriptsize] at (4.35,\y) {\nt{\code}};
  \fill[vizrest] (4.45,\y-0.13) rectangle ({4.45+\lg-3},\y+0.13);
  \node[anchor=west,font=\scriptsize] at ({4.5+\lx},\y) {\val};}
\draw[vizrest,thick] (7.45,1.0) -- (8.45,1.0);
\draw[vizrest,thick] (7.45,0.92) -- (7.45,1.08);
\draw[vizrest,thick] (8.45,0.92) -- (8.45,1.08);
\draw[gray!60!black] (4.45,-1.72) -- (8.45,-1.72);
\foreach \e in {3,4,5,6,7}{
  \draw[gray!60!black] ({4.45+\e-3},-1.72) -- ({4.45+\e-3},-1.66);
  \node[below,font=\scriptsize] at ({4.45+\e-3},-1.72) {$10^{\e}$};}
\end{tikzpicture}
\caption{表~\ref{tab:types}の推定による脳内のニューロンの型の割合。円はほぼ全体が\nt{GLUT}で占められ、\nt{GABA}は細い扇形、それ以外の型はすべて合わせても境界上の髪の毛一本にすぎない。右側はその髪の毛一本をニューロン数の対数目盛で拡大したものである。\nt{GLYC}の棒は$10^6$--$10^7$の範囲の中央まで伸び、線分が範囲全体を示す。\nt{ASPA}は推定値がないので図にない。}
\label{fig:typepie}
\end{table}

これらの山の大きさがどれほど不揃いかは図~\ref{fig:typepie}でわかる。ニューロン1000個のうち$964$個が\nt{GLUT}、$36$個が\nt{GABA}であり、それ以外の型はすべて合わせても10万個あたり約6個にすぎない。

GABAはもともと4文字の略称である。主要な神経伝達物質になることがほとんどない物質、つまり神経ペプチド、気体、脂質、プリンには独自のコードを与えない。これらについては付録~\ref{sec:othermediators}で扱う。命題~\ref{prop:dale}の「ほとんど」という語は重要である。多くのニューロンは、主要な神経伝達物質と一緒に補助的な物質を放出する~\cite{Yao2023}。二つ目の速い神経伝達物質を、しかも逆の符号のものを放出するニューロンもある。\nt{DOPA}ニューロンの一部はGABAも放出する~\cite{Tritsch2012}。さらに、同じ出力の異なる枝から異なる神経伝達物質が出るニューロンもある~\cite{Zhang2015}。これらはすべて測定されているので、「ニューロン1個に神経伝達物質1種類」は例外のある規則であって、法則ではない。それでも、最初の区分としては検証に耐える。発現している遺伝子によってニューロンを5000以上の群に分けたマウス脳の細胞アトラスでも、ニューロンの大きなクラスはやはり、まず主要な神経伝達物質で分かれる~\cite{Yao2023}。

この規則からは、脳を人工ニューラルネットワークとただちに区別する帰結が導かれる。人工ネットワークでは、同じニューロンがある受け手には正の重みを、別の受け手には負の重みをもちうる。重み$w_{ij}$は行列の中の単なる数だからである。脳では、作用の符号はおもに神経伝達物質で決まり、ニューロンの神経伝達物質は1種類である。したがって、ニューロン$j$がある受け手を興奮させるなら、ほかの受け手もすべて興奮させる。ニューロン$j$から出る重み行列の\emph{列}全体が同じ符号をもつ。
\begin{empheq}[box=\keybox]{equation}
\operatorname{sign} w_{ij} \;=\; s_j \quad\text{すべての受け手 } i \text{ について}, \qquad s_j\in\{+1,-1\}.
\label{eq:dalesign}
\end{empheq}
ニューロンは興奮性（$s_j=+1$）と抑制性（$s_j=-1$）に分かれ、これは結合ではなくニューロン自体の性質である。興奮性ニューロンからの負の結合がネットワークに必要なら、仲介役を通すしかない。興奮性ニューロンが抑制性ニューロンを興奮させ、それが標的を抑制する。つまり、1段分の遅延をともなう負の重みである。

神経伝達物質は100種類以上知られているが、脳の仕事のほとんどは6種類ほどが担っており、それらはまったく異なる二つのクラスに分かれる。

\section{二つのバス：アドレスバスとブロードキャスト}
\label{sec:buses}

コンピュータネットワークには、メッセージを届ける方法が二つある。一つ目は\emph{アドレス指定}、つまりユニキャスト（unicast）である。パケットは1人の送り手から特定の受け手へ速く届き、ほかの誰にも見えない。二つ目は\emph{ブロードキャスト}（broadcast）である。送り手は一つのメッセージをネットワークの全ノードに一度に送る。ニューロンは両方の方法を使い、神経伝達物質はまさにこの基準で分かれる（図~\ref{fig:phone_radio}）。

\begin{figure}[t]
\centering
\begin{tikzpicture}[>=Stealth,scale=0.95]
% ---- unicast: point-to-point wiring
\node[above] at (2.0,3.3) {\small \textbf{アドレスバス}：グルタミン酸とGABA};
\foreach \i in {0,...,4}{
  \fill[blue!60!black] (0.2+0.9*\i,2.5) circle (0.1);
  \fill[blue!60!black] (0.2+0.9*\i,0.6) circle (0.1);}
\foreach \a/\b in {0/0,0/1,1/1,1/3,2/2,3/2,3/4,4/4,2/0}{
  \draw[->,blue!60!black] (0.2+0.9*\a,2.38) -- (0.2+0.9*\b,0.72);}
\fill[red!70!black] (1.1,1.55) circle (0.09);
\pic[scale=1.2] at (1.75,1.9) {letter};
\draw[->,red!70!black,thick] (1.1,1.55) -- (0.28,0.7);
\draw[->,red!70!black,thick] (1.1,1.55) -- (1.95,0.72);
\node[align=center,below] at (2.0,0.2) {\small 数十億のニューロン、\\[-2pt]\small それぞれに固有の宛先、\\[-2pt]\small 時間：ミリ秒};
% ---- broadcast: one small source, global targets
\begin{scope}[xshift=7.2cm]
\node[above] at (1.8,3.3) {\small \textbf{ブロードキャスト}：ドーパミンなど};
\foreach \i in {0,...,8}{\fill[blue!60!black] (-0.2+0.5*\i,2.75) circle (0.08);}
\fill[orange!85!black] (1.8,0.35) ellipse (0.35 and 0.18);
\pic[scale=1.5] at (2.7,0.75) {mast};
\foreach \x in {-0.2,0.3,0.8,1.3,1.8,2.3,2.8,3.3,3.8}{
  \draw[->,orange!85!black] (1.8,0.5) .. controls (1.8,1.3) and (\x,1.6) .. (\x,2.62);}
\node[align=center,below] at (1.8,0.1) {\small 数千から数十万のニューロン、\\[-2pt]\small 一つのメッセージを全員に、\\[-2pt]\small 時間：数百ミリ秒以上};
\end{scope}
\end{tikzpicture}
\caption{二つの伝達方式。左はアドレスバスで、封筒に宛先を書く郵便に似ている。赤は1個のニューロンとその受け手2個を示す。右はブロードキャストで、ラジオ局に似ている。発信元のニューロンは小さな群であり、同じメッセージを全員が受け取る。}
\label{fig:phone_radio}
\end{figure}

\emph{アドレスバス}はグルタミン酸とGABAである。圧倒的多数のニューロンがこれらを放出する。各出力は特定の細胞につながり、神経伝達物質は接点の狭い隙間の中だけで、しかも速く作用する。1ミリ秒で受け手のイオンチャネルが開き、数ミリ秒ですべてが終わる。これは\emph{データ}バスであり、脳が計算しているものがここを通る。

\emph{ブロードキャスト}のバスはドーパミン、セロトニン、ノルアドレナリン、そして一部はアセチルコリンである。これらを放出するのはごく小さなニューロン群だが、その一つひとつの出力は信じられないほど広く枝分かれする。神経伝達物質は狭い隙間ではなく細胞間の空間に放出されることが多く、そこで広がって周囲のすべてに作用する。作用は遅い。チャネルを直接開くのではなく、受け手の内部で化学反応の連鎖を起こす。ここを通るのはデータではなく、ネットワークの広い領域に共通の\emph{制御パラメータ}であり、それがネットワークの動作モードを変える。ゲイン、学習率、「今のはよかった」という信号である。そのため、このような神経伝達物質は\emph{修飾物質}（モジュレータ）と呼ばれる。こうしたチャネルはそれぞれ脳全体でただ一つの数だと長く考えられてきた。近年の測定はこの描像を精密にした。チャネルは1本の配線ではなく、並列な線を束ねた小さなハーネスである。同じ神経伝達物質の発信元ニューロンはサブシステムに分かれ、それぞれに固有の受け手と固有のメッセージがあり、その場でどれだけ放出するかは局所の信号によっても調整される~\cite{Ren2018,Poe2020}。こうした線の数は数本から数十本であって数十億ではないので、チャネルはやはりブロードキャストである。

この章から絵を一枚だけ持ち帰るなら、これである。脳はデータをアドレス指定で伝える巨大なネットワークであり、その上に制御信号を運ぶブロードキャストチャネルがいくつか張られている。プログラマなら見慣れた構成に気づくだろう。計算と、少数の制御変数の組である。それぞれの制御変数はネットワークの広い領域に共通である。

\section{\nt{GLUT}：正の重み}

グルタミン酸は脳の主要な興奮性神経伝達物質である。ニューロンの出力がグルタミン酸を放出すると、受け手ではナトリウムを内側に通すチャネルが開き、受け手の膜電位が上がって、受け手自身がスパイクを出す確率が高まる。図~\ref{fig:blackbox}の言葉でいえば、これは\emph{正の}重み$w_i>0$をもつ入力である。

グルタミン酸を出すのは誰か。長距離でデータを伝えるほとんどすべてのニューロンである。大脳皮質のニューロンの4分の3から5分の4、そして脳の異なる領域をつなぐニューロンの大部分がそうである。脳のネットワークは、何よりもまずグルタミン酸結合のネットワークである。

工学的な細部を一つ。放出の後、隙間の中のグルタミン酸濃度は数千倍に跳ね上がって約1ミリモーラーに達し、時定数約1ミリ秒で減衰する~\cite{Clements1992}。グルタミン酸は隙間から拡散して出ていき、専用のポンプタンパク質が細胞内へくみ戻す。なぜか。通信線路で各シンボルの後に信号をゼロに戻すのと同じ理由である。神経伝達物質を取り除かなければ、次のスパイクは「使用中」の線路に到着し、数ミリ秒間隔のスパイクは融合してしまう。

\section{\nt{GABA}：負の重み}

重みがすべて正のネットワークを想像してほしい。平均して1個のスパイクが次のスパイクを1個より多く生むなら、活動は指数関数的に増え、数ステップでネットワーク全体に広がる。電子回路の技術者なら、ループゲインが1を超える正帰還ループだとわかるだろう。回路は飽和する。そうならないためには負の重みが必要である。その主な供給源が$\gamma$-アミノ酪酸、略してGABAである。

GABAは受け手で塩化物イオンを通すチャネルを開く。塩化物イオンの平衡電位は静止電位付近かそれより少し低いところにあるので、開いた塩素チャネルは膜を静止電位付近に保ち、電圧を閾値まで上げさせない。これは\emph{負の}重み$w_i<0$をもつ入力である。GABAニューロンは大脳皮質のニューロンの5分の1から4分の1を占める~\cite{Hendry1987}。その出力は短く、すぐ隣のニューロンを抑制する。

脳における抑制は、単なる引き算よりはるかに多くのことをする。ネットワーク全体を動作点に保つ負帰還として働くのである。正常に動作している大脳皮質では、ニューロンに入る興奮性電流と抑制性電流がほぼ正確に釣り合っていると考えられている。この状態は理論で予言され~\cite{vanVreeswijk1996}、生きた皮質での電流測定でも見えており~\cite{Haider2006}、ニューロンは常に閾値付近にいる。不安定点の近くで強い負帰還によって安定化した回路は、小さな信号に速く敏感に反応する。昔からある工学の手法である。

\section{\nt{DOPA}：誤差信号}

最初のブロードキャストチャネルはドーパミンである。ヒトのドーパミンニューロンは約50万個、およそ17万個に1個である。これはドーパミンニューロンの二つの群のうち主要なほうで数えられた値なので~\cite{Pakkenberg1991}、下限である。その出力は、行動を選択する領域へと広がっていく。

このチャネルは何を伝えるのか。答えは、報酬を受け取る動物でドーパミンニューロンのスパイクを記録する実験から得られる~\cite{Schultz1997}。動物にときどき予告なしにジュースを1滴与えると、ドーパミンニューロンは短いスパイクのバーストで応答する。次に、ジュースの前にランプを点けるようにする。いつもジュースの1秒前である。動物がこの関係を学習すると、ニューロンは\emph{ランプ}にバーストで応答するようになり、予定どおりの時刻に来たジュースそのものにはまったく応答しなくなる。そしてランプの後にジュースを与えないと、ジュースが来るはずだった瞬間にニューロンは\emph{黙り込み}、発火率は通常より下がる。

三つの場合すべてを説明する規則はこうである（図~\ref{fig:rpe}）。ニューロンが伝えるのは、どれだけよいかではなく、\emph{予想よりどれだけよいか悪いか}である。

\begin{figure}[t]
\centering
\begin{tikzpicture}[>=Stealth,x=3.4cm,y=1cm]
\foreach \y/\lab in {2.8/{予告なしのジュース},1.4/{ランプの後のジュース},0/{ランプだがジュースなし}}{
  \draw[gray!70!black] (0,\y) -- (3,\y);
  \draw[densely dotted,gray!60!black] (0,\y+0.3) -- (3,\y+0.3);
  \node[left,align=right,font=\small] at (-0.03,\y+0.35) {\lab};}
% row 1: juice without a cue -> burst at the juice
\draw[very thick,orange!85!black] plot[domain=0:3,samples=120] (\x,{2.8+0.3+0.75*exp(-((\x-2.08)/0.07)^2)});
\pic[scale=0.8] at (2,2.58) {juice};
% row 2: cue then juice -> burst at the cue, nothing at the juice
\draw[very thick,orange!85!black] plot[domain=0:3,samples=120] (\x,{1.4+0.3+0.75*exp(-((\x-1.08)/0.07)^2)});
\pic[scale=0.85] at (1,1.15) {lamp}; \pic[scale=0.85] at (2,1.15) {juice};
% row 3: cue, juice omitted -> burst at the cue, dip when the juice was due
\draw[very thick,orange!85!black] plot[domain=0:3,samples=120] (\x,{0.3+0.75*exp(-((\x-1.08)/0.07)^2)-0.28*exp(-((\x-2.1)/0.12)^2)});
\pic[scale=0.85] at (1,-0.25) {lamp}; \pic[scale=0.85] at (2,-0.25) {nojuice};
\node[right,font=\scriptsize] at (2.36,0.1) {落ち込み};
\node[right,font=\scriptsize] at (2.13,3.75) {驚き！};
\node[left,font=\scriptsize] at (0.97,2.25) {ランプへのバースト};
\node[above,font=\scriptsize] at (2,1.75) {応答なし};
\draw[->,gray!70!black] (0,-0.75) -- (3.05,-0.75) node[right,black,font=\small] {$t$};
\draw[gray!60!black] (1,-0.8) -- (1,-0.7) node[below=2pt,font=\scriptsize] {ランプ、$1\,\text{s}$};
\draw[gray!60!black] (2,-0.8) -- (2,-0.7) node[below=2pt,font=\scriptsize] {ジュース、$2\,\text{s}$};
\end{tikzpicture}
\caption{三つの実験における\nt{DOPA}ニューロンの発火率（模式図、\cite{Schultz1997}による）。点線は一定の背景発火率。ランプは合図、しずくはジュース、斜線を引いたしずくは約束されたのに与えられなかったジュースである。応答は予測誤差~\eqref{eq:rpe}である。}
\label{fig:rpe}
\end{figure}

強化学習~\cite{SuttonBarto2018}を知っているプログラマなら、この量にすぐ気づくだろう。行動の選択を学習するエージェントは推定値$V(s)$を保持している。状態$s$にいるときに期待する報酬の量である。エージェントは各ステップの後、期待と実際に得たものを比べる。
\begin{empheq}[box=\keybox]{equation}
\delta_t \;=\; r_t \;+\; \gamma\,V(s_{t+1}) \;-\; V(s_t) ,
\label{eq:rpe}
\end{empheq}
そして推定値を修正する：$V(s_t)\leftarrow V(s_t)+\alpha\,\delta_t$。ここで$r_t$は得られた報酬、$\gamma<1$は割引率（将来の報酬が即時の報酬よりどれだけ安いか）、$\alpha$は学習率である。量$\delta_t$は\emph{時間差分誤差}と呼ばれる。

この量はドーパミンニューロンとまったく同じようにふるまう。予告なしのジュース：$V$はまだゼロで、$r>0$、$\delta>0$。学習済みのジュース：報酬はランプが予測しており、ジュースの直前の$V(s_t)$はすでに$r$に等しいので、$\delta=0$。ジュースが来ない：$V(s_t)=r$だが$r_t=0$なので、$\delta<0$。そしてバーストがランプに移るのは、まさにランプの瞬間に期待が跳ね上がるからである：$\gamma V(s_{t+1})-V(s_t)>0$。ここでのステップ$t,t+1,\dots$はアルゴリズム上の約束事にすぎない。生体の中に時間の刻みはなく、同じ量を連続時間で第\ref{sec:dofa}章で導く。

つまりドーパミンは、それをもとに学習すべきすべての相手に、スカラーの誤差信号$\delta$をブロードキャストで配信するものである。第一近似では、脳全体に1本の配線、各瞬間に一つの数である。この配線が実際にはどの程度ハーネスなのかは、第\ref{sec:dofaalt}章で検討する。

\section{その他のブロードキャストチャネル}

ほかの修飾物質については、その役割を同じ大域パラメータの言葉に翻訳する作業仮説があるだけである~\cite{Doya2002}。あくまで仮説として受け止めてほしい。

\emph{\nt{NORA}、ノルアドレナリン：ゲイン。}ノルアドレナリンニューロンはドーパミンニューロンよりさらに少なく、粗い推定で数万個だが、その出力は脳のほぼ全体に広がる。予想外のことや重要なことが起きると、これらは急激に活性化する。ノルアドレナリンは受け手のニューロンの伝達特性の傾きを変える。弱い入力はより弱く、強い入力はより強くなる。これは次のように測定される。受け手の背景スパイクが入力への応答より強く抑えられ、信号対雑音比が上がる~\cite{Foote1975NE}。単純なモデルではこれを一つの数で表す。ニューロンが入力電流$u$に発火率$f=F\bigl(g\cdot u\bigr)$で応答するなら、ノルアドレナリンは係数$g$を大きくする~\cite{AstonJones2005}（詳しくは第\ref{sec:nora}章）。$g$が大きいネットワークは「コントラスト高く」動作し、速く決定を下す。$g$が小さいと「ぼやけて」動作し、探索モードにある強化学習エージェントのように、より多くの選択肢を試す。

\emph{\nt{ACET}、アセチルコリン：学習率。}アセチルコリンは、ネットワークが現在の入力にどれだけ耳を傾け、自分が保存しているデータにどれだけ耳を傾けるかを制御する。アセチルコリンの濃度が高いと、感覚器からの入力が強まり、内部の「想起」結合が弱まる。同時に重みの変更が起こりやすくなる~\cite{Hasselmo2006}。大まかにいえば、これは「書き込み/読み出し」の切り替えスイッチであり、それとともに学習率$\alpha$でもある。

\emph{\nt{SERO}、セロトニン：計画の時間的視野。}セロトニンが何を伝えるかは最もよくわかっていない。仮説の一つは、これを式~\eqref{eq:rpe}の割引率$\gamma$と結びつける。セロトニンの濃度が高いことは、小さな報酬にすぐ飛びつかずに大きな報酬を待つ構えに対応する。セロトニンニューロンは一様にゆっくりと動作し、覚醒時には毎秒数スパイクで発火し、睡眠のある相ではほとんど黙り込む~\cite{JacobsAzmitia1992}。

6種類の神経伝達物質すべてを表~\ref{tab:transmitters}にまとめた。

\begin{table}[t]
\centering
\small
\begin{tabular}{>{\raggedright}p{2.4cm}>{\raggedright}p{3.0cm}>{\raggedright}p{2.0cm}>{\raggedright\arraybackslash}p{5.2cm}}
\toprule
ニューロンの型 & ニューロンの割合 & バス & 計算での役割 \\
\midrule
\nt{GLUT}（グルタミン酸） & 大多数 & アドレス & 正の重み、$w>0$ \\
\nt{GABA}（GABA） & 皮質で20--25\% & アドレス & 負の重み、$w<0$。安定化 \\
\nt{DOPA}（ドーパミン） & $\sim 6\cdot10^{-6}$ & ブロード\newline キャスト & 予測誤差$\delta$ \\
\nt{NORA}（ノルアドレナリン） & $\sim 6\cdot10^{-7}$ & ブロード\newline キャスト & ゲイン$g$（仮説） \\
\nt{ACET}（アセチルコリン） & 少数 & 両方 & 学習率$\alpha$。「書き込み/読み出し」（仮説） \\
\nt{SERO}（セロトニン） & $\sim 3\cdot10^{-6}$ & ブロード\newline キャスト & 割引$\gamma$（仮説） \\
\bottomrule
\end{tabular}
\caption{脳の主要な神経伝達物質を計算の言葉で表したもの。右の列は大幅な単純化である。グルタミン酸、GABA、ドーパミンについては信頼できるが、ほかは仮説である。}
\label{tab:transmitters}
\end{table}

\section{符号を決めるのは受け手}
\label{sec:receiver}

グルタミン酸は正の重み、GABAは負の重み、と言ったとき、実は少しごまかしていた。符号をもつ分子は一つもない。分子は単なる鍵である。それが捕らえられたとき何が起きるかを決めるのは\emph{錠}、つまり受け手の細胞にある受容体である。

最もよい例はドーパミンである。ドーパミンには二つのファミリーの受容体がある~\cite{Beaulieu2011}。一方では細胞の興奮を促し、他方では妨げる。行動を選択する領域では、一部のニューロンが第一の型の受容体を、別のニューロンが第二の型の受容体をもち、同じドーパミン信号が一方の群を強めると同時に他方の群を弱める。一つのブロードキャストパケットを、ネットワークのノードごとに異なって解釈するようなものである。

\begin{keyblock}
\begin{proposition}[メッセージの意味を決めるのは受け手]
\label{prop:receptor}
神経伝達物質の作用の符号と速さは、神経伝達物質そのものではなく、それが結合する受容体で決まる。受容体には二つの種類がある。\emph{イオンチャネル型}は、それ自体がイオンチャネルであり、1ミリ秒未満で開く。トランジスタのゲートのような、電流の直接制御である。\emph{代謝型}は細胞内で化学反応の連鎖を起こし、数百ミリ秒以上かけて作用する。むしろ設定レジスタへの書き込みに近く、細胞のその後の動作を変える。アドレスバスはおもに前者を通じて、ブロードキャストはおもに後者を通じて話す。
\end{proposition}
\end{keyblock}

それなら、なぜそもそも「グルタミン酸は興奮させる」と言えるのか。実際に出会うグルタミン酸受容体のほとんどすべてが興奮性であり、GABA受容体のほとんどすべてが抑制性だからである。これは自然法則ではなく、非常に信頼性の高い取り決めであり、規則~\eqref{eq:dalesign}はまさにその上に成り立っている。

\section{持ち帰るもの}

圧倒的多数のニューロンはアドレス指定のデータバスであり、重みの符号はニューロンごとに一つである。ごく少数派がブロードキャストチャネルであり、脳の広い領域にいくつかの共通の制御信号を送る。10万個のうちわずか2個ほどのニューロンに、残りのネットワーク全体がどう学習し、どのモードで動作するかがかかっている。エンジニアにとって、これは驚くことではない。どんなコンピュータでも制御レジスタはメモリセルよりも比べものにならないほど少ないが、それがなければマシンは動かない。

次章では、最も数の多い型である\nt{GLUT}ニューロンだけでできたネットワークが何を計算できるかを調べる。

\section{問題}

\begin{exercise}[\lvlA]
\label{ex:01:1}
\emph{山と配線。}表~\ref{tab:types}から（範囲は対数目盛での中央、\nt{ACET}の出力数は$10^5$）：(a)~ニューロン1個を人間1人とすると、\nt{GLUT}ニューロンは地球の人口の何倍になり、ブロードキャスト型全体はどれほどの規模の都市になるか。(b)~\nt{DOPA}、\nt{SERO}、\nt{NORA}、\nt{ACET}の終末の割合は、そのニューロンの割合の何倍か。
\end{exercise}

\begin{exercise}[\lvlA]
\label{ex:01:2}
\emph{ゼロへの復帰。}隙間のグルタミン酸は$c_0e^{-t/\tau_c}$で減衰する。$c_0=1\mM$、$\tau_c=1\ms$で、受け手は$10$~\textmu M以上の濃度を検知する。(a)~1回の放出の後、線路はどれだけの時間「使用中」になるか。(b)~ポンプが部分的に阻害され$\tau_c=10\ms$になった。線路が空くのが間に合う放出頻度の上限はいくらか。
\end{exercise}

\begin{exercise}[\lvlB]
\label{ex:01:3}
\emph{バーストの移動先。}図~\ref{fig:rpe}の実験で、ランプの後に状態$s_0\to\dots\to s_{K-1}$が刻み$\Delta$で続き、$T=K\Delta$、報酬$r$は最後の遷移で届き、ランプの時刻は予測できない。学習後の$V(s_k)$とランプへの応答$\delta$を求めよ。$\gamma=e^{-\Delta/\tau_\gamma}$とおき、$T=1$~s、$\Delta=0.1$~s、$\gamma=0.95$のときの$\tau_\gamma$を求めよ。
\end{exercise}

\begin{exercise}[\lvlB]
\label{ex:01:4}
\emph{シミュレーション：誤差による学習。}問題~\ref{ex:01:3}の連鎖を$K=10$、$\gamma=0.95$、$r=1$、規則$V(s_t)\leftarrow V(s_t)+\alpha\delta_t$でシミュレーションせよ。$\alpha=0.05$、$0.1$、$0.2$、$0.5$のとき、ランプへの応答がジュースへの応答を初めて上回る試行$n^*$はいくつか。$n^*$は$\alpha$にどう依存するか。
\end{exercise}

\begin{exercise}[\lvlB]
\label{ex:01:5}
\emph{設計：一つの数の配信。}数$\delta$を$10^{10}$個のニューロン全部に届けたい。中継役を含む各受け手は前の層から$10$個の終末を受け取る必要があり、1段ごとに$1\ms$かかる。ニューロン1個あたりの出力数が$10^4$（\nt{GLUT}）の場合と$10^{5.5}$（\nt{DOPA}）の場合で、最も経済的な中継ツリーのニューロン数と段数はいくつか。
\end{exercise}

\begin{exercise}[\lvlB]
\label{ex:01:6}
\emph{釣り合いはなぜ敏感か。}ネットワークは\nt{GLUT}が$80\,\%$、\nt{GABA}が$20\,\%$で、全対全に重み$J_E/N$と$-J_I/N$で結合し、$\tau\dot r=-r+Wr+h$である。$J_E=10$のとき、$W$のゼロでない唯一の固有値は$0.5$である。抑制性電流と興奮性電流の比を求め、ネットワークが雪崩に陥るには抑制を何パーセント弱めれば足りるかを求めよ。
\end{exercise}

\begin{exercise}[\lvlC]
\label{ex:01:7}
\emph{研究：\nt{SERO}は$\gamma$か。}問題~\ref{ex:01:3}より、ランプへの\nt{DOPA}ニューロンの応答は$re^{-T/\tau_\gamma}$である。遅延$T=1,\dots,5$~sをそれぞれ$n$回ずつ提示し、$\ln\delta$はばらつき$0.5$で測定される。$\tau_\gamma=2$~sと$2.4$~sを区別するには$n$はいくつ必要か（有意水準$0.05$、検出力$0.8$）。$\gamma$以外のどんな機構が$T$に対して同じ減衰を生むか。
\end{exercise}

\chapter{\nt{GLUT}ニューロンは何を計算するか}
\label{sec:glutcomputer}
\begin{chaptergoals}
\item 漏れ積分発火ニューロンがORと、時間のANDである同時検出器を作る仕組み；
\item なぜ\nt{GLUT}ニューロンだけでは単調関数しか計算できず、活動で活動を止められないのか；
\item オン・オフの配線ペアで\nt{GLUT}ニューロンが任意の論理関数を計算し、準備完了を示す仕組み；
\item 遅延線が時間差を場所に変える仕組みと、帰還がメモリとタイマーを作る仕組み。
\end{chaptergoals}

\section{ゲームのルール}

第\ref{sec:neurontypes}章で見たように、脳のニューロンの圧倒的多数は\nt{GLUT}である。これらは興奮させるだけで、重みは正だけである。このようなニューロンを好きなだけ用意して問おう。生体の中で起こりうることだけを許すとき、このネットワークは何を計算できるか。

生体には、すべての素子を一斉に切り替えるクロック発生器がない。各ニューロンは独立に連続時間で動作し、スパイクはまちまちの遅延で、そのつど届いては出ていく。\emph{入力}は光、音、圧力で発火する感覚ニューロンであり、\emph{出力}は筋肉を制御するニューロンである。その間にネットワークがあり、いつ計算を始めていつ終えるかを教えてくれる者はいない。電子回路の技術者にとって、これは素子の遅延が事前に正確にはわからない\emph{非同期}回路である。

ニューロンのモデルとしては、連続時間で正しく動作するもののうち最も単純なものを採る。ニューロンの膜電位$V$は静止時にゼロである。ニューロン$j$からスパイクが届くたびに$V$は$w_j\ge0$だけ跳ね上がり、その後は時定数$\tau$でひとりでにゼロへ戻っていく。
\begin{empheq}[box=\keybox]{equation}
\tau\,\frac{\dd V}{\dd t} \;=\; -V \;+\; \tau\sum_j w_j \sum_k \delta\bigl(t-t_j^{(k)}-d_j\bigr), \qquad w_j\ge 0 .
\label{eq:glutneuron}
\end{empheq}
ここで$t_j^{(k)}$はニューロン$j$のスパイクの時刻、$d_j$はそこからの経路の遅延、$\delta$はデルタ関数で、瞬間的な跳躍を表す。$V$が閾値$\theta$に達すると、ニューロンはスパイクを出し、$V$はゼロにリセットされ、約1ミリ秒の間ニューロンは再び発火できない。典型的な値は、$\tau$がニューロンによって1ミリ秒未満から20ミリ秒、遅延が1ミリ秒未満から数十ミリ秒である。これは\emph{漏れ積分発火ニューロン}、つまり閾値つきのRC回路である。意図的な単純化であることに注意しよう。皮質ニューロンでは入力の枝そのものが局所的なスパイクを出し~\cite{Smith2013}、1個のニューロンはむしろ小さな多層ネットワークのようにふるまう（第\ref{sec:synthesis}章）。この章の問いには、RC回路一つで足りる。

論理ゲートとの主な違いは漏れである。ニューロンはスパイクを永遠に覚えているのではなく、約$\tau$だけ覚えている。足し合わされるのは、この窓の中に届いたものだけである。

\section{ORとAND}

ゲートから始めよう。1個のスパイクで電圧が閾値を超える、つまり$w\ge\theta$なら、ニューロンはどの入力のどのスパイクにも発火する。これが\emph{OR}である。

ANDはもっと面白い。$w<\theta<2w$としよう。スパイク1個では足りず、2個必要である。だがこうなると、「両方届いたか」という問いは、\emph{どれだけの時間内に}かを決めない限り意味をもたない。最初のスパイクが時刻$0$に、2番目が時間$\Delta$の後に届いたとする。2番目が届いた時点で最初のスパイクの分は$we^{-\Delta/\tau}$残っており、$we^{-\Delta/\tau}+w\ge\theta$ならニューロンは発火する。ここから同時性の窓が得られる。
\begin{empheq}[box=\keybox]{equation}
\Delta \;\le\; \Delta_{\max} \;=\; \tau\,\ln\frac{w}{\theta-w} .
\label{eq:window}
\end{empheq}
これは論理的なANDではなく\emph{同時検出器}、つまり時間の中のANDである（図~\ref{fig:coincidence}）。$\theta=1.5w$なら窓は$\tau\ln2\approx0.7\tau$である。$\tau=10\ms$の皮質ニューロンでは約7ミリ秒になる。$\tau$が1ミリ秒未満の聴覚系のニューロンでは、窓は1ミリ秒未満に縮む。

\begin{figure}[t]
\centering
\begin{tikzpicture}[>=Stealth,x=0.9cm,y=1.6cm]
\foreach \dx/\lab/\c [count=\i] in {0.5/{同時}/red!70!black, 2.6/{同時でない}/blue!60!black}{
 \begin{scope}[xshift={(\i-1)*7.2cm}]
  \draw[->,gray!70!black] (0,0) -- (6.3,0) node[below left,black] {\small $t$};
  \draw[->,gray!70!black] (0,0) -- (0,1.75) node[left,black] {\small $V$};
  \draw[densely dashed,gray!60!black] (0,1.5) -- (6.0,1.5) node[right,black] {\small $\theta$};
  \draw[very thick,\c] (0,0) -- (1,0) -- (1,1)
      plot[domain=1:1+\dx,samples=30] (\x,{exp(-(\x-1)/1.5)});
  \pgfmathsetmacro{\v}{exp(-\dx/1.5)+1}
  \draw[very thick,\c] (1+\dx,{exp(-\dx/1.5)}) -- (1+\dx,{min(\v,1.5)});
  \ifdim\v pt<1.5pt
    \draw[very thick,\c] plot[domain=1+\dx:6,samples=40] (\x,{\v*exp(-(\x-1-\dx)/1.5)});
  \else
    \draw[very thick,\c] (1+\dx,1.5) -- (1+\dx,0) -- (6,0);
    \draw[->,thick,\c] (1+\dx,1.55) -- (1+\dx,1.72) node[above,font=\scriptsize] {スパイク};
  \fi
  \draw[->,thick] (1,-0.35) -- (1,-0.08); \draw[->,thick] (1+\dx,-0.35) -- (1+\dx,-0.08);
  \draw[decorate,decoration={brace,amplitude=3pt,mirror}] (1,-0.42) -- (1+\dx,-0.42) node[midway,below=3pt] {\scriptsize $\Delta$};
  \node[\c] at (3.5,-0.95) {\small \lab};
 \end{scope}}
\end{tikzpicture}
\caption{同時検出器、閾値$\theta=1.5w$。左では2番目のスパイクが$\Delta<\Delta_{\max}$の後に届き、和が閾値に達してニューロンが発火した。右では$\Delta>\Delta_{\max}$で、最初のスパイクの分はほとんど残っておらず、ニューロンは黙っている。}
\label{fig:coincidence}
\end{figure}

ANDを得るもう一つの方法は、正確な時刻ではなく持続時間を使うことである。光が点いている間、感覚ニューロンは高い頻度で一様にスパイクを出す。閾値にとってそのような入力1本では足りず2本なら足りるとすれば、ニューロンは両方が活動している間ずっと動作する。こうしてネットワークは\emph{レベル}で計算し、スパイクの到着の正確な時刻は問題にならない。以下の多くはこの方式である。

\section{できないこと}

次にNOTを作ってみよう。入力が黙っているときに動作するニューロンである。これはできない。入力スパイクがなければ式~\eqref{eq:glutneuron}には跳躍が一つもなく、電圧はゼロ、つまり閾値より下にとどまる。多数のニューロンからなるネットワークならどうか。やはりできない。しかもこれは、あらゆるネットワークについて一挙に証明できる。

発火率で考えるのが最も便利である。$r_i$をニューロン$i$の発火率、$I_i$を入力からの流入、$f_i$をその伝達特性、つまり発火率が流入の総和にどう依存するかとする。積分型ニューロンではこの特性は非減少である。流入が大きいほどスパイクは頻繁になる。平衡に達したネットワークの発火率は、次の方程式を満たす。
\begin{equation}
r_i \;=\; f_i\Bigl(\sum_j w_{ij}\,r_j + I_i\Bigr), \qquad w_{ij}\ge 0 .
\label{eq:ratenet}
\end{equation}

\begin{keyblock}
\begin{proposition}[単調性]
\label{prop:monotone}
ネットワーク~\eqref{eq:ratenet}が\nt{GLUT}ニューロンだけからなり、完全な静寂から出発するとする。入力を強めても、どのニューロンの定常発火率も減らない。このようなネットワークが計算するものはすべて入力の\emph{単調}関数である。
\end{proposition}
\end{keyblock}

証明。静寂$r^{(0)}=0$から始め、発火率を式~\eqref{eq:ratenet}の右辺に代入していく：$r^{(k+1)}=F\bigl(r^{(k)}\bigr)$。重みが非負で$f_i$が非減少なので、右辺$F$はどの引数についても減少しない。$r^{(1)}\ge0=r^{(0)}$だから、帰納法により$r^{(k+1)}\ge r^{(k)}$である。発火率は増えるだけで、式~\eqref{eq:ratenet}の最小解に到達する。入力$I$をより大きな$I'$に置き換えれば、各ステップで$r^{(k)}(I)\le r^{(k)}(I')$であり、極限でも同じである。実際のネットワークはステップではなく連続的に平衡に達するが、同じことが成り立つ。正の結合のもとでは、二つの試行の間に最初に成り立っていた不等式はずっと保たれる。

個々のスパイクについては、この主張は文字どおりには成り立たない。余分な入力スパイクがニューロンを少し早く発火させ、1ミリ秒の不応期のせいで次のスパイクが失われることがある。しかしこの効果は弱く当てにならないので、計算には使えない。発火率については単調性は厳密である。

帰結はインバータのない回路と同じである。NOTがない。排他的ORがない。2番目の入力を加えても出力を消すことはできない。そして最もやっかいなのは、\emph{活動で活動を止めることはできない}ことである。できるのは、活動を支えているものを取り除くことだけである。

\section{2本の配線：オンとオフ}

行き詰まりからの出口は生体自身が示している。多くの感覚器では、刺激は2組のニューロンで伝えられる。視覚では、ある細胞は光の増加に、別の細胞は光の減少に応答する~\cite{Schiller1992}。これらを\emph{オン}ニューロン、\emph{オフ}ニューロンと呼ぼう。ネットワークは1本の配線$x$の代わりにペア$(x,\bar x)$を受け取る。ネットワークが自分では計算できない「ないこと」は、センサ自身が知らせてくれるのである。

配線のペアがあれば、NOTはただで手に入る。配線を入れ替えればよい。ANDとORはそれぞれニューロン2個で作れる。
\begin{empheq}[box=\keybox]{equation}
\begin{aligned}
x\wedge y \;&\to\; \bigl(\;x\wedge y,\;\; \bar x\vee\bar y\;\bigr),\\
x\vee y \;&\to\; \bigl(\;x\vee y,\;\; \bar x\wedge\bar y\;\bigr).
\end{aligned}
\label{eq:dualrail}
\end{empheq}
右辺の演算はすべて単調なので、命題~\ref{prop:monotone}には反しない。

非同期回路にとって、2線式符号にはもう一つ、さらに重要な長所がある。配線のペアは三つの状態をとる。「はい」、「いいえ」、そして両方が黙っているときの\emph{「まだ答えがない」}である。受け手は計算が終わったことをクロックではなく信号そのものから知る。各ペアでちょうど1本の配線が動作し始めればよい。刺激と刺激の間はセンサが黙り、ネットワークは静寂に戻って次の回に備える。非同期エレクトロニクスでは、この手法を使って素子の遅延がどうであれ正しく動作する回路を作る~\cite{Sparso2001}。

\begin{keyblock}
\begin{proposition}[非同期計算]
\label{prop:dualrail}
各入力が「オン・オフ」の配線ペアで届くなら、\nt{GLUT}ニューロンのネットワークは入力の任意の論理関数を計算できる。答えも配線ペアで届き、その準備ができたことは各ペアで1本の配線が動作し始めたことでわかる。そのためにクロックは要らない。代償はニューロンがおよそ2倍になることである。
\end{proposition}
\end{keyblock}

例は排他的ORである。「二つの刺激のちょうど一方だけが変化した」。答えの正の配線は$(x\wedge\bar y)\vee(\bar x\wedge y)$、負の配線は$(x\wedge y)\vee(\bar x\wedge\bar y)$である。ニューロン6個で、どれも抑制を必要としない（図~\ref{fig:dualxor}）。

\begin{figure}[t]
\centering
\begin{tikzpicture}[>=Stealth,x=1cm,y=1cm,
  nrn/.style={circle,draw,very thick,fill=blue!6,minimum size=0.75cm,inner sep=0pt,font=\small},
  inp/.style={font=\small}]
\node[inp] (X)  at (0,3.3) {$x$};
\node[inp] (Xb) at (0,2.2) {$\bar x$};
\node[inp] (Y)  at (0,1.1) {$y$};
\node[inp] (Yb) at (0,0)   {$\bar y$};
\node[nrn] (A1) at (3.2,3.3) {$2$};
\node[nrn] (A2) at (3.2,2.2) {$2$};
\node[nrn] (A3) at (3.2,1.1) {$2$};
\node[nrn] (A4) at (3.2,0)   {$2$};
\node[nrn] (O)  at (7.2,2.75) {$1$};
\node[nrn] (Ob) at (7.2,0.55) {$1$};
\foreach \s/\t in {X/A1,Yb/A1,Xb/A2,Y/A2,X/A3,Y/A3,Xb/A4,Yb/A4}{\draw[->,thick,blue!60!black] (\s) -- (\t);}
\draw[->,thick,blue!60!black] (A1) -- node[pos=0.45,above,sloped,font=\scriptsize,black] {$x\wedge\bar y$} (O);
\draw[->,thick,blue!60!black] (A2) -- node[pos=0.45,below,sloped,font=\scriptsize,black] {$\bar x\wedge y$} (O);
\draw[->,thick,blue!60!black] (A3) -- node[pos=0.45,above,sloped,font=\scriptsize,black] {$x\wedge y$} (Ob);
\draw[->,thick,blue!60!black] (A4) -- node[pos=0.45,below,sloped,font=\scriptsize,black] {$\bar x\wedge\bar y$} (Ob);
\draw[->,very thick,red!70!black] (O) -- (8.6,2.75) node[right,black] {$x\oplus y$：「はい」};
\draw[->,very thick,red!70!black] (Ob) -- (8.6,0.55) node[right,black] {$\overline{x\oplus y}$：「いいえ」};
\node[below,font=\small,gray!40!black] at (3.2,-0.55) {AND：閾値$2$};
\node[below,font=\small,gray!40!black] at (7.2,-0.55) {OR：閾値$1$};
\node[below,font=\small,gray!40!black] at (0,-0.55) {配線ペア};
\end{tikzpicture}
\caption{2線式符号による排他的OR。\nt{GLUT}ニューロンだけで作る。丸の中の数は、入力1本の重みを単位とした閾値である。4個のANDニューロンが入力の4通りの組み合わせを認識し、2個のORニューロンがそれを答えの配線ペアにまとめる。}
\label{fig:dualxor}
\end{figure}

\section{クロックの代わりに遅延}

時間を使う計算には、配線の遅延と同時検出器があれば十分である。最もよい例は、脳が音の来た方向を知るしくみである。

横にある音源からの音は、遠い耳より近い耳に早く届く。その差は方向で決まり、音源が真正面ならゼロ、ヒトで真横なら$0.22\,\text{m}/343\,\text{m/s}\approx0.6\ms$である。脳は約10\emph{マイクロ秒}の差を区別する~\cite{KlumppEady1956,Grothe2010}。ニューロン自体の発火の精度は1ミリ秒未満の程度にすぎないのに、なぜか。

左耳からの配線を同時検出器の列に沿って左から右へ、右耳からの配線を右から左へ通す（図~\ref{fig:itd}）。信号は配線上を速さ$v$で進む。長さ$L$の列の左端から距離$x$にある検出器まで、左の信号は時間$x/v$、右の信号は$(L-x)/v$かかる。右耳に音が$\delta t$だけ遅れて届いたなら、二つの信号は次の点で同時に出会う。
\begin{empheq}[box=\keybox]{equation}
\frac{x}{v} \;=\; \frac{L-x}{v} + \delta t
\quad\Longrightarrow\quad
x \;=\; \frac{L}{2} + \frac{v\,\delta t}{2} .
\label{eq:itd}
\end{empheq}
発火するのは、二つの信号が窓~\eqref{eq:window}の中に届いた検出器だけである。発火した検出器の番号が、そのまま音源の方向である。時間の差が\emph{場所}に変わった。

\begin{figure}[t]
\centering
\begin{tikzpicture}[>=Stealth,
  nrn/.style={circle,draw,very thick,fill=blue!6,minimum size=0.7cm,inner sep=0pt}]
\foreach \k in {0,...,6}{\node[nrn] (D\k) at (1.4*\k,0) {\scriptsize $2$};}
\node[nrn,fill=red!15] at (D4) {\scriptsize $2$};
\draw[->,thick,blue!60!black] (-1.4,0.9) node[left,black] {\small 左耳} -- (9.2,0.9);
\draw[->,thick,orange!85!black] (9.8,-0.9) node[right,black] {\small 右耳} -- (-0.8,-0.9);
\foreach \k in {0,...,6}{
  \draw[->,blue!60!black] (1.4*\k,0.9) -- (D\k.north);
  \draw[->,orange!85!black] (1.4*\k,-0.9) -- (D\k.south);}
\draw[->,thick,red!70!black] (D4.north east) ++(0.1,0.05) -- ++(0.5,0.45) node[above right,font=\small,black] {発火};
\node[font=\small,gray!40!black] at (4.2,-1.5) {遅延が増える $\longrightarrow$ \hspace{2.2cm} $\longleftarrow$ 遅延が増える};
\pic[scale=1.4] at (-2.3,1.75) {speaker};
\node[font=\scriptsize,align=center] at (-2.3,2.35) {左の音源};
\end{tikzpicture}
\caption{音の方向の検出。両耳からの信号は互いに向かって進む。各丸は閾値$2$の同時検出器である。右耳に音が遅れて届くと、信号は式~\eqref{eq:itd}に従って中央より右で出会う。ネットワークの出力は発火したニューロンの番号である。ここでは音源が左にあり、左耳に音が先に届いた。}
\label{fig:itd}
\end{figure}

ここにはクロックも抑制もなく、2線式符号さえない。ネットワークは「はい」か「いいえ」で答えるのではなく、多数のニューロンのうち1個を指し示す。遅延線と同時検出器からなるこの回路は、鳥の聴覚脳幹で実際に働いているらしい~\cite{Carr1990}。マイクロ秒の精度は、個々のニューロンの精度からではなく、検出器が多数あってそれぞれが自分の遅延を見ていることから生まれる。哺乳類ではこのような遅延の列は見つかっていない。そこでは同じ課題を、\nt{GLYC}ニューロンからの時間的に精密な抑制で調整された同時検出器が解いているらしい~\cite{Brand2002,Grothe2010}。抑制が同時性の窓をどう狭めるかは、第\ref{sec:gaba}章で\nt{GABA}を例に検討する。グリシンも受け手の塩素チャネルを開いて同じように抑制する。

\section{メモリ}

コンピュータにはメモリが必要である。このようなネットワークでは、メモリとは自分自身を維持する活動である。互いに強く結合した\nt{GLUT}ニューロンの群をとり、それを一つの発火率$r$で表そう（図~\ref{fig:recurrentgroup}(a)）。平衡では$r=f(wr+I)$である。ここで$w$は群の自分自身への帰還の強さ、$I$は外部からの流入である。この方程式を、曲線$f(wr+I)$と直線$r$の交点として図で解く（図~\ref{fig:bistable}）。

\begin{figure}[t]
\centering
% (b): tau dr/dt = -r + f(w r + I), f as in fig:bistable, w=1, I=0.6; Euler dt=0.001 tau.
\begin{tikzpicture}[>=Stealth,
  nrn/.style={circle,draw,thick,fill=blue!6,minimum size=0.5cm,inner sep=0pt}]
\begin{scope}
\node[font=\small] at (-0.5,1.55) {(a)};
\foreach \k in {1,...,6}{\node[nrn] (n\k) at ({1.4+1.0*cos(60*\k)},{1.0*sin(60*\k)}) {};}
\foreach \a/\b in {1/2,2/3,3/4,4/5,5/6,6/1,1/4,2/5,3/6}{\draw[<->,blue!60!black] (n\a) -- (n\b);}
\foreach \k in {2,3,4}{\draw[->,thick,green!45!black] ($(n\k)+(-0.95,0)$) -- (n\k);}
\node[font=\small,green!45!black] at (-0.35,-0.45) {$I$};
\node[font=\Large] at (3.1,0) {$\approx$};
\node[nrn,minimum size=0.85cm,font=\small] (R) at (4.35,-0.1) {$r$};
\draw[->,thick,blue!60!black] (R) edge[loop above,looseness=6] node[above,font=\small] {$w$} (R);
\draw[->,thick,green!45!black] (3.55,-0.95) node[left,font=\small] {$I$} -- (R);
\end{scope}
\begin{scope}[xshift=6.3cm,y=2.6cm,x=1.05cm]
\node[font=\small] at (-0.6,1.25) {(b)};
\draw[->,gray!70!black] (0,0) -- (7.3,0) node[below left,font=\small,black] {$t/\tau$};
\draw[->,gray!70!black] (0,0) -- (0,1.12) node[left,font=\small,black] {$r$};
\foreach \t in {0,2,4,6}{\draw[gray!60!black] (\t,-0.02) -- (\t,0.02); \node[below,font=\scriptsize] at (\t,0) {\t};}
\draw[densely dotted,gray!60!black] (0,0.5) -- (7,0.5) node[right,font=\scriptsize,black,align=left] {書き込み\\閾値};
\fill[green!45!black,opacity=0.25] (1,-0.2) rectangle (2,-0.13);
\fill[gray!60!black,opacity=0.35] (1,-0.29) rectangle (1.5,-0.22);
\draw[very thick,blue!60!black] plot coordinates {(0.0,0.002) (0.1,0.002) (0.2,0.002) (0.3,0.002) (0.4,0.002) (0.5,0.002) (0.6,0.002) (0.7,0.002) (0.8,0.002) (0.9,0.002) (1.0,0.002) (1.1,0.083) (1.2,0.165) (1.3,0.242) (1.4,0.313) (1.5,0.378) (1.6,0.437) (1.7,0.491) (1.8,0.539) (1.9,0.583) (2.0,0.623) (2.1,0.643) (2.2,0.665) (2.3,0.688) (2.4,0.71) (2.5,0.732) (2.6,0.753) (2.7,0.773) (2.8,0.792) (2.9,0.81) (3.0,0.826) (3.2,0.855) (3.4,0.88) (3.6,0.9) (3.8,0.917) (4.0,0.931) (4.4,0.953) (4.8,0.967) (5.2,0.977) (5.6,0.984) (6.0,0.988) (6.5,0.992) (7.0,0.994)};
\draw[very thick,gray!60!black,densely dashed] plot coordinates {(0.0,0.002) (0.1,0.002) (0.2,0.002) (0.3,0.002) (0.4,0.002) (0.5,0.002) (0.6,0.002) (0.7,0.002) (0.8,0.002) (0.9,0.002) (1.0,0.002) (1.1,0.083) (1.2,0.165) (1.3,0.242) (1.4,0.313) (1.5,0.378) (1.6,0.358) (1.7,0.336) (1.8,0.313) (1.9,0.291) (2.0,0.269) (2.2,0.228) (2.4,0.191) (2.6,0.159) (2.8,0.133) (3.0,0.11) (3.2,0.091) (3.4,0.076) (3.6,0.063) (3.8,0.052) (4.0,0.043) (4.4,0.03) (4.8,0.021) (5.2,0.015) (5.6,0.011) (6.0,0.008) (6.5,0.006) (7.0,0.004)};
\node[font=\scriptsize,blue!60!black,anchor=west] at (3.3,0.8) {流入$1\tau$：書き込み成功};
\node[font=\scriptsize,gray!50!black,anchor=west] at (2.3,0.3) {流入$0.5\tau$：消えた};
\end{scope}
\end{tikzpicture}
\caption{自分自身と強く結合した群。(a)~互いに興奮させ合う多数の\nt{GLUT}ニューロンは、強さ$w$で自分自身を興奮させる発火率$r$の1個のニューロンのようにふるまう。外部からは流入$I$が来る。(b)~$w=1$、図~\ref{fig:bistable}と同じ曲線$f$のときの群の発火率の時間変化。流入$I=0.6$は、軸の下の帯で示した時間だけ与えられる。流入が$1\tau$続けば、群は不安定点$r=0.5$を越え、燃え上がり、流入が終わっても燃え続ける。$0.5\tau$だけなら、この点まで届かずに消えて元に戻る。ビットを書き込むには、流入はおよそ$0.7\tau$より長く続かなければならない。}
\label{fig:recurrentgroup}
\end{figure}

\begin{figure}[t]
\centering
\begin{tikzpicture}[>=Stealth,x=5cm,y=3.2cm]
\draw[->,gray!70!black] (0,0) -- (1.1,0) node[below left,black] {\small $r$};
\draw[->,gray!70!black] (0,0) -- (0,1.1);
\draw[thick,gray!60!black] (0,0) -- (1.05,1.05) node[right,black] {\small $r$};
\draw[very thick,blue!60!black] plot[domain=0:1.05,samples=80] (\x,{1/(1+exp(-(\x-0.5)/0.08))});
\draw[very thick,red!70!black,densely dashed] plot[domain=0:1.05,samples=80] (\x,{1/(1+exp(-(\x+0.3-0.5)/0.08))});
\fill[blue!60!black] (0.002,0.002) circle (0.018);
\draw[blue!60!black,fill=white,thick] (0.5,0.5) circle (0.018);
\fill[blue!60!black] (0.998,0.998) circle (0.018);
\node[below right,font=\small] at (0.02,0.0) {オフ};
\node[right,font=\small] at (0.52,0.46) {不安定};
\node[below,font=\small] at (0.99,0.96) {オン};
\node[anchor=east,font=\small,red!70!black] at (0.19,0.62) {$f(wr+I)$};
\node[anchor=west,font=\small,blue!60!black] at (0.45,0.1) {$f(wr)$};
\end{tikzpicture}
\caption{強い帰還をもつ\nt{GLUT}ニューロン群によるメモリ。実線は外部流入がない場合で、直線$r$との安定な交点が二つ（「オフ」と「オン」）、その間に不安定な交点が一つある。短い流入$I$（破線）があると、上の交点だけが残る。}
\label{fig:bistable}
\end{figure}

帰還が強ければ交点は三つある。下は静寂、上は群が全力で燃えている状態、中央は不安定である。二つの安定状態をもつ素子、つまりフリップフロップができた。これに1を書き込むのは簡単である。短い流入が曲線を持ち上げ、下の交点が消え、群は燃え上がる。そして流入が終わっても燃え続ける（図~\ref{fig:recurrentgroup}(b)）。流入が短すぎると群は不安定点まで届かず、消えて元に戻る。刺激が消えた後も数秒続くこのような自己維持活動は、実際に脳で観察されており、短期記憶の基礎と考えられている~\cite{Wang2001}。だが、秒単位の記憶の方法はこれだけではない。書き込んだ内容はシナプス自体の遅い変数にも保存できる。スパイクのバーストの後、第\ref{sec:code}章の「ツー」シナプスのような促通が約1秒続き、ネットワークは短い発火の合間にほとんど黙っていられる。フリップフロップではなく、充電されたコンデンサである~\cite{Mongillo2008}。記憶保持中のこのような「静かな」区間は観察されており~\cite{Stokes2015}、スパイクを出し続けない保存のほうがエネルギー的に安い。皮質がどの場合にどちらの方法を使うかは議論が続いている。

では、どう消去するのか。命題~\ref{prop:monotone}により、活動では消せない。何かを\emph{取り除く}しかない。一つ目の方法は支えを取り除くことである。群が別の群からの流入があるときだけ燃えているなら、記憶はその流入とともに消える。二つ目は疲労である。実際のシナプスでは長く働くと神経伝達物質の蓄えが枯渇し~\cite{Tsodyks1997}、ニューロンでは興奮性を下げる遅い電流が増えていく。帰還が弱まり、上の交点が消え、群は数秒後にひとりでに消える。信号でオンになり、時間だけでオフになるメモリができる。

\section{シーケンス}

クロックの代わりに、\emph{イベントで起動するタイマー}をもつことができる。\nt{GLUT}ニューロンの群を連鎖状につなぐ。1番目が2番目を興奮させ、2番目が3番目を、というように続く。外部信号が1番目の群を点火すると、活動の波が連鎖を走り抜ける。各段は前の段から1--2遅延分の後に、ただ一度だけ発火する。スパイク直後のニューロンは不応期にあり、そのとき波はもう先へ進んでいるからである。

このような連鎖はイベントからの時間を測る。$k$番目の段が発火したなら、起動から$k$遅延分が経過している。その出力を筋肉に送れば、ひとりでに展開する一連の動作が得られる。鳴禽では、歌っている間、脳のある部位の個々の\nt{GLUT}ニューロンが厳密に順番に、それぞれ歌の決まった瞬間に発火する。このような連鎖によく似ている~\cite{Hahnloser2002}。

クロックと違って、連鎖は起動されるまで黙っており、一度だけ動作し、つながっている者のためにだけ時間を測る。ネットワーク全体に共通のクロックは依然としてない。

\section{足りないもの}

抑制なしで\nt{GLUT}ニューロンだけから多くのものが作れる。だが三つのものが足りず、三つとも命題~\ref{prop:monotone}のせいである。

\emph{選択。}ネットワークは一つの方向、一つの行動を選ばなければならない。二つの刺激がほぼ等しいと、図~\ref{fig:itd}のネットワークでは両方の出力が発火し、弱いほうを抑えて強いほうを選ぶ手段がない。

\emph{リセット。}信号でメモリを消去することができない。

\emph{安定性。}結合がエンジニアの設計どおりにできるのではないネットワークでは、正帰還ループは避けられず、その中で活動が雪崩のように増えうる（第\ref{sec:neurontypes}章）。唯一のブレーキは、遅くて粗い同じ疲労である。

三つの問題はすべて一つの薬で治る。スパイクでスパイクを消すニューロンである。それが\nt{GABA}であり、必要なのは計算するためではなく、選択し、リセットし、計算を制御下に保つためである。

\section{問題}

\begin{exercise}[\lvlA]
\label{ex:02:1}
\emph{音のための検出器の列。}両耳からの信号は図~\ref{fig:itd}の配線を速さ$5$~m/sで進み、到着時刻の差の最大値は$0.64\ms$である。検出器は$\tau=0.5\ms$、$\theta=1.5w$のニューロン~\eqref{eq:glutneuron}で、隣どうしが$10$~\textmu sを区別できるように並んでいる。列には検出器が何個あり、一つの音にそのうち何個が発火するか。
\end{exercise}

\begin{exercise}[\lvlB]
\label{ex:02:2}
\emph{等しくない入力は窓をずらす。}$\tau=0.5\ms$、$\theta=1.5$の同時検出器~\eqref{eq:glutneuron}が、重み$1$と$0.8$の二つの入力から1個ずつスパイクを受け取る。同時性の窓とその中心を求めよ。一方の耳からの入力が他方より$20\,\%$弱いと、図~\ref{fig:itd}の列は何マイクロ秒誤るか。
\end{exercise}

\begin{exercise}[\lvlB]
\label{ex:02:3}
\emph{振動しないネットワーク。}非減少の$f_i$と$i\ne j$で$w_{ij}\ge0$をもつネットワーク$\tau\dot r_i=-r_i+f_i\bigl(\textstyle\sum_jw_{ij}r_j+I_i\bigr)$は単調で、安定的に振動しない。各サイクルの抑制性結合の数が偶数であるネットワークは、置換$r_i\to\pm r_i$でこれに帰着することを証明せよ。二つの群の相互抑制は振動しうるか。\nt{GLUT}--\nt{GABA}ループはどうか。
\end{exercise}

\begin{exercise}[\lvlB]
\label{ex:02:4}
\emph{設計：2線式加算器。}各ビットは配線ペア$(x,\bar x)$で伝えられる。和のペア$(S,\bar S)$と桁上げのペア$(C,\bar C)$を出す3ビットの全加算器を、重みと閾値を指定して\nt{GLUT}ニューロンで作れ。ニューロン8個に収めよ。図~\ref{fig:dualxor}のようにANDとORの素子で作ると何個必要か。
\end{exercise}

\begin{exercise}[\lvlB]
\label{ex:02:5}
\emph{シミュレーション：書き込みにかかる時間。}図~\ref{fig:recurrentgroup}の群：$\tau\dot r=-r+f(r+I)$、$f(u)=1/\bigl(1+e^{-(u-0.5)/0.08}\bigr)$で、書き込み前は下の状態にある。流入$I$を時間$T$だけ与える。$I=0.6$でビットを書き込む最小の$T$と、どんな$T$でもビットが書き込まれなくなる閾値$I_c$を数値的に求めよ。
\end{exercise}

\begin{exercise}[\lvlA]
\label{ex:02:6}
\emph{タイマーとしての連鎖。}連鎖の1段は\nt{GLUT}ニューロン$100$個で、遅延は$2\ms$、独立なばらつきは$0.2\ms$である。(a)~$1$~sを測るにはニューロンが何個必要で、相対誤差はいくらか。それは間隔にどう依存するか。(b)~実際のタイマーの誤差はどの間隔でも$5\,\%$である。どんなばらつきの源がそれを生むか。
\end{exercise}

\begin{exercise}[\lvlC]
\label{ex:02:7}
\emph{研究：フリップフロップかコンデンサか。}ニューロン$100$個がビットを$10$~s保持する。フリップフロップ：各ニューロンが毎秒$20$スパイク。コンデンサ：促通$u$は$\tau_F=1$~sで減衰し、$u\ge u_{\max}/2$ならビットが読め、リフレッシュは各ニューロンの3スパイクのバースト。コンデンサは何倍経済的か。群を$200\ms$黙らせるとビットはどうなるか。
\end{exercise}

\chapter{\nt{GLUT}と\nt{GABA}の協働}
\label{sec:gaba}

\section{ゲームのルール}

\nt{GLUT}ニューロンだけのネットワークは、選択することも、メモリをリセットすることも、活動の雪崩を抑えることもできない。すべて単調性のせいである（命題~\ref{prop:monotone}）。そこで、数で2番目に多いニューロンの型、\nt{GABA}を加えよう。ルールは前章と同じである。生体の中で起こりうることだけを使い、共通のクロックは使わない。

ニューロンのモデルは式~\eqref{eq:glutneuron}と同じだが、\nt{GABA}ニューロンからのスパイクは受け手の電圧を上げるのではなく下げる。重みは今や二つの符号をとり、符号は送り手が決める。規則~\eqref{eq:dalesign}である。

回路のパラメータをいくつか挙げる。\nt{GABA}ニューロンは皮質のニューロンの5分の1から4分の1を占める~\cite{Hendry1987}。その出力は短く、遠い領域ではなく隣のニューロンを抑制する。抑制は1ミリ秒後に届き、数ミリ秒続く。入力がなければ\nt{GABA}ニューロンは黙っている。誰かを抑制するには、自分がまず興奮を受け取らなければならない。ふつうは同じネットワークの\nt{GLUT}ニューロンからである。

そのため、\nt{GLUT}ニューロン$A$からニューロン$B$への負の結合は、仲介役を通して作るしかない。$A$が\nt{GABA}ニューロンを興奮させ、それが$B$を抑制する。符号の反転にはニューロン1個と遅延1段、約1ミリ秒がかかる。

抑制では、結合が\emph{誰から}来るかだけでなく、\emph{どこに}つくかも重要である。結合の着地点が、その結合の働きを大きく左右する~\cite{Markram2004}。\nt{GLUT}の出力はおもに受け手の入力の枝、つまり\emph{樹状突起}につき、データを運ぶ。そのような入力の一つひとつが和の一つの項である。ニューロンの細胞体への出力は和全体に作用する。多くの速い\nt{GABA}ニューロンはこうして受け手の応答を割り算し、受け手がいつ発火するかを決める。受け手のスパイクが生まれる場所への出力は最終決定であり、出力全体を一度に禁止する。そして他のニューロンの配線の終末への出力は、他の線には触れずに、1本の入力線の放出確率$p$を変える~\cite{Rudomin1999}。第\ref{sec:pain}章の痛みのゲートは、とくにこのしくみで働く。脳の外では、出力は筋線維（第\ref{sec:muscles}章）や臓器（第\ref{sec:chemoutput}章）につき、どこにもつかずに神経伝達物質を組織の体積中や血液中に放出するものもある（第\ref{sec:serocomputer}章と第\ref{sec:chemoutput}章）。

\section{NOTと禁止}

これでNOTは作れるか。ほぼ作れる。入力が黙っているときに動作するニューロンは、どこかから興奮を得なければならない。生きた脳にはそれがある。各ニューロンには、数千の他のニューロンからの背景活動が絶えず届いている。背景活動がニューロン$Y$を動作させておき、入力$x$が\nt{GABA}の仲介役を通じてそれを抑制するとしよう。これがNOTである。

しかし、より有用なのは\emph{禁止}、「$a$、ただし$b$のときは除く」である。
\begin{empheq}[box=\keybox]{equation}
y \;=\; a \wedge \bar b .
\label{eq:veto}
\end{empheq}
ニューロン$Y$は$a$から興奮を受け取り、$b$が興奮させる\nt{GABA}ニューロンから抑制を受け取る。ここでは背景活動は要らない。興奮は$a$自身が与える。禁止とORからすべてが組み立てられる。たとえば排他的OR：$x\oplus y=(x\wedge\bar y)\vee(\bar x\wedge y)$は、禁止ニューロン2個、\nt{GABA}の仲介役2個、ORニューロン1個でできる。

\begin{keyblock}
\begin{proposition}[\nt{GLUT}+\nt{GABA}の完全性]
\label{prop:gabacomplete}
\nt{GLUT}ニューロンと\nt{GABA}ニューロンからなるネットワークは、入力の任意の論理関数を計算でき、各ビットは1本の配線で伝えられる。第\ref{sec:glutcomputer}章の「オン・オフ」センサはもう必要ない。すべての入力が黙っているときに1を出すべき関数には、背景活動の源が必要である。
\end{proposition}
\end{keyblock}

証明：定数1があれば、禁止~\eqref{eq:veto}とORは関数的に完全である。NOT $x$は禁止「1、ただし$x$のときは除く」だからである。そして、すべての入力が黙っているときに0を出す関数は、定数1なしに禁止とORから組み立てられる。

\section{運動の方向}

時間の中の禁止は、第\ref{sec:glutcomputer}章の時間の中のANDと同様に、運動方向の検出器になる。

隣り合う二つのセンサ$A$と$B$が距離$\Delta x$だけ離れているとする。出力ニューロン$Y$は$B$から興奮を受け取り、$A$から\nt{GABA}の仲介役を通じて遅延$d$で抑制を受け取る（図~\ref{fig:direction}）。$A$から$B$へ速さ$v$で動く光点は時刻$0$に$A$を通過し、抑制は時刻$d$に$Y$に届く。光点は時刻$\Delta x/v$に$B$を通過し、そのとき興奮が届く。これらの時刻が窓~\eqref{eq:window}の精度で一致すれば、抑制が興奮を打ち消し、$Y$は黙る。これが起こるのは速さがほぼ次の値のときである。
\begin{empheq}[box=\keybox]{equation}
v^* \;=\; \frac{\Delta x}{d} .
\label{eq:vstar}
\end{empheq}
$B$から$A$へ動く場合は、$B$からの興奮が時刻$0$に届き、$A$からの抑制は時刻$\Delta x/v+d$になってやっと届くが、そのとき$Y$はもう発火している。

\begin{figure}[t]
\centering
\begin{tikzpicture}[>=Stealth,
  nrn/.style={circle,draw,very thick,fill=blue!6,minimum size=0.9cm,inner sep=0pt},
  gab/.style={nrn,fill=red!10},
  inh/.style={-{Bar[width=7pt]},thick,red!70!black}]
\node[nrn] (A) at (0,2) {$A$};
\node[nrn] (B) at (3,2) {$B$};
\node[gab] (G) at (0.9,0.6) {\small \nt{GABA}};
\node[nrn] (Y) at (3,-0.6) {$Y$};
\draw[->,thick] (A) -- (G);
\draw[inh] (G) -- node[below left=-2pt] {\scriptsize 遅延$d$} (Y);
\draw[->,thick] (B) -- (Y);
\draw[->,very thick,red!70!black] (Y) -- (4.6,-0.6) node[right,black] {\small 出力};
\draw[->,thick,orange!85!black] (-0.6,2.9) -- (3.6,2.9) node[right,black] {\small 沈黙};
\draw[<-,thick,orange!85!black] (-0.6,3.4) -- (3.6,3.4) node[right,black] {\small 応答};
\foreach \yy/\dd in {2.9/-1,3.4/1}{
  \foreach \k in {1,2,3}{\draw[orange!70!yellow,line width=0.8pt] (1.5+\dd*0.12+\dd*0.13*\k,\yy+0.1-0.1*\k+0.1) -- ++(\dd*0.25,0);}
  \fill[yellow!70!orange,draw=orange!70!black] (1.5,\yy) circle (0.14);}
\node[font=\scriptsize,align=center] at (1.5,3.95) {光点};
\end{tikzpicture}
\caption{運動方向の検出器。$B$は$Y$を興奮させ、$A$は\nt{GABA}の仲介役を通じて遅延$d$で$Y$を抑制する。$A$から$B$へ約$\Delta x/d$の速さで動くと、抑制が興奮を打ち消す。逆向きに動くと抑制は間に合わない。短い線のついた黄色の点は、センサの上を走る光である。}
\label{fig:direction}
\end{figure}

網膜では、運動方向の選択性はまさにこのように、つまり応答すべきでない運動が来る側からの非対称な抑制によって作られているらしい~\cite{Barlow1965}。

\section{引き算か割り算か}

これまで抑制は引き算をしてきた。実際には、そのしくみはもっと繊細である。

シナプスはチャネルを開く。つまりコンダクタンスをオンにする。興奮性シナプスの平衡電位は閾値よりはるか上、静止電位より約$70\mV$上にあり、開いたチャネルは電圧を上に引っ張る。抑制性の\nt{GABA}シナプスのチャネルは塩化物イオンを通し、その平衡電位は静止電位付近にあるので、開いたチャネルは電圧を下にではなく\emph{静止電位へ}引っ張る。電圧を静止電位から測れば、漏れコンダクタンス$g_L$、興奮性コンダクタンス$g_E$、抑制性コンダクタンス$g_I$をもつ膜の定常電圧は次のようになる。
\begin{empheq}[box=\keybox]{equation}
V_\infty \;=\; \frac{g_E\,E_E}{g_L+g_E+g_I} ,
\label{eq:shunt}
\end{empheq}
ここで$E_E\approx70\mV$は興奮性シナプスの平衡電位である。これは分圧器のオームの法則である。$g_E$は$E_E$へ、$g_L$と$g_I$はゼロへ引っ張る。

$g_I$は分子にマイナス符号つきで入るのではなく、分母に入る。抑制は興奮から\emph{引く}のではなく、興奮を\emph{割る}のである（図~\ref{fig:shunt}）。このような抑制はシャント抑制と呼ばれる。開いた塩素チャネルが膜を短絡するからである。ネットワークにとって、これはゲイン調整である。$g_I$が近隣の活動の総和に比例するなら、各ニューロンは自分の入力の絶対的な強さではなく、全体の活動に占めるその割合に応答する。全体の活動によるこのような割り算は脳の多くの部位で見られ、脳の標準的な演算の一つと考えられている~\cite{Carandini2012}。同じネットワークが弱い入力でも強い入力でも動作できる。

但し書きがある。式~\eqref{eq:shunt}はスパイクを出さないニューロンについての式である。発火しているニューロンでは電圧は常に閾値付近に保たれ、抑制性コンダクタンスを通る電流はほぼ一定なので、シャントはスパイクの\emph{発火率}にはおもに引き算として作用する~\cite{HoltKoch1997}。発火率が割り算されるのは、生きた皮質の揺らぎのある釣り合った状態のように、釣り合った興奮性と抑制性の背景コンダクタンスをニューロンに同時に加えた場合である~\cite{Chance2002}。つまり脳は割り算を、シャント1個からではなく、抑制と背景ノイズの組み合わせから得ている~\cite{Carandini2012}。

\begin{figure}[t]
\centering
\begin{tikzpicture}[>=Stealth,x=1.25cm,y=0.05cm]
\draw[->,gray!70!black] (0,0) -- (5.4,0) node[right,black] {\small $g_E/g_L$};
\draw[->,gray!70!black] (0,0) -- (0,68) node[above,black] {\small $V_\infty$, mV};
\foreach \v in {20,40,60}{\draw[gray!60!black] (-0.05,\v) -- (0.05,\v) node[left=3pt,font=\scriptsize] {\v};}
\foreach \x in {1,2,3,4,5}{\draw[gray!60!black] (\x,-1.5) -- (\x,1.5) node[below=5pt,font=\scriptsize] {\x};}
\draw[very thick,blue!60!black] plot[domain=0:5,samples=60] (\x,{70*\x/(1+\x)});
\draw[very thick,red!70!black] plot[domain=0:5,samples=60] (\x,{70*\x/(3+\x)});
\draw[thick,densely dashed,gray!50!black] plot[domain=0.56:5,samples=60] (\x,{70*\x/(1+\x)-25});
\node[right,font=\small,blue!60!black] at (5.02,58.3) {抑制なし};
\node[right,font=\small,red!70!black] at (5.02,45.5) {割り算：$g_I=2g_L$};
\node[right,font=\small,gray!40!black] at (5.02,32.5) {$25\mV$の引き算};
\end{tikzpicture}
\caption{シャント抑制は電圧から引くのではなく、電圧を割る。青の曲線は抑制なしの定常電圧~\eqref{eq:shunt}、赤は抑制性コンダクタンス$g_I=2g_L$のあるもの。赤の曲線は青の曲線を横に3倍引き延ばしたものであり、入力を$1+g_I/g_L=3$で割ったのと同じである。破線は一定の$25\mV$を引く抑制で、曲線全体が下がり、傾きは変わらない。}
\label{fig:shunt}
\end{figure}

もう一つの帰結もある。膜の時定数は$\tau_{\text{eff}}=C/(g_L+g_E+g_I)$であり、抑制はこれを短くする。同時性の窓~\eqref{eq:window}は$\tau$に比例するので、興奮と一緒に届いた抑制は\emph{窓を狭める}。入力がニューロンを直接興奮させると同時に、1ミリ秒遅れて\nt{GABA}の仲介役を通じて抑制する回路は、脳で最もありふれた回路の一つである。ニューロンは最初の1--2ミリ秒に届いたものにしか応答する暇がない~\cite{Pouille2001}。

\section{選択}

いよいよ\nt{GLUT}ネットワークに欠けていた最も重要なもの、複数の候補から一つを選ぶことである。入力$I_1$と$I_2$をもつ\nt{GLUT}ニューロンの群を二つとる。それぞれが自分の\nt{GABA}仲介役を通じて、強さ$\beta$で相手を抑制する（図~\ref{fig:wta}）。発火率$r_1,r_2$について次を得る。
\begin{equation}
\tau\,\frac{\dd r_1}{\dd t} = -r_1 + \bigl[I_1-\beta r_2\bigr]_+ ,\qquad
\tau\,\frac{\dd r_2}{\dd t} = -r_2 + \bigl[I_2-\beta r_1\bigr]_+ ,
\label{eq:wta}
\end{equation}
ここで$[u]_+=\max(u,0)$である。発火率は負にならない。

\begin{figure}[t]
\centering
\begin{tikzpicture}[>=Stealth,
  nrn/.style={circle,draw,very thick,fill=blue!6,minimum size=0.9cm,inner sep=0pt},
  gab/.style={nrn,fill=red!10,minimum size=0.75cm},
  inh/.style={-{Bar[width=7pt]},thick,red!70!black}]
\node[nrn] (E1) at (0,0) {$r_1$};
\node[nrn] (E2) at (4,0) {$r_2$};
\node[gab] (G1) at (1.4,1.1) {};
\node[gab] (G2) at (2.6,-1.1) {};
\draw[->,thick] (E1) -- (G1); \draw[inh] (G1) -- (E2);
\draw[->,thick] (E2) -- (G2); \draw[inh] (G2) -- (E1);
\draw[->,thick,blue!60!black] (-1.6,0) node[left,black] {$I_1$} -- (E1);
\draw[->,thick,blue!60!black] (5.6,0) node[right,black] {$I_2$} -- (E2);
\end{tikzpicture}
\caption{二者択一。\nt{GLUT}ニューロンの二つの群（青）が、\nt{GABA}の仲介役（赤）を通じて互いに抑制し合う。}
\label{fig:wta}
\end{figure}

両方のニューロンが動作している間、系~\eqref{eq:wta}は線形であり、そのふるまいは二つの固有値で決まる。両方の発火率が同じ向きにずれる場合は$-(1+\beta)/\tau$、逆向きにずれる場合は$-(1-\beta)/\tau$である。抑制が弱い$\beta<1$では両方とも負で、両方のニューロンが動作し、抑制は弱いほうを少し抑えるだけである。抑制が強い$\beta>1$では2番目の固有値が正になる。$r_1$と$r_2$の差はどんなものでも増大し、やがて一方のニューロンが完全に黙る。

\begin{keyblock}
\begin{proposition}[勝者総取り]
\label{prop:wta}
相互抑制が1より強い、$\beta>1$なら、両方の群が動作する状態は不安定である。ネットワークは一方の群が動作し他方が黙る状態に至る。発火率$r_1=I_1$の勝者は、$I_2<\beta I_1$である限り勝利を保つ。
\end{proposition}
\end{keyblock}

これをモデルで確かめた。$\beta=0.5$、入力$1.0$と$0.9$では、両方の群が発火率$0.73$と$0.53$で動作する。$\beta=2$で同じ入力なら、第2の群は完全に黙る。この選択には記憶がある。その後で入力を入れ替えても、$1.0<2\cdot0.9$なので、以前の勝者が勝者のままである。

100個の群でも同じである。それぞれが他のすべてを抑制し、一つだけが生き残る。

\section{メモリのリセット}

第\ref{sec:glutcomputer}章のメモリは疲労によってしか消えなかった。そこに、「リセット」信号で興奮し群を抑制する\nt{GABA}ニューロンを加えよう。抑制が図~\ref{fig:bistable}の曲線$f(wr+I)$を下げ、上の交点が消え、群は消える。そしてリセットが終わっても下の状態にとどまる。こうしてメモリは一つの信号で書き込まれ、別の信号で消去される。セット入力とリセット入力をもつフリップフロップと同じである。

さらに美しいのは、そのような群を二つ、図~\ref{fig:wta}のように相互抑制でつなぎ、それぞれに自己帰還を与えることである。一方の群の入力への信号がセルをその状態に切り替え、セルは最後の決定を覚えている。これは、たすき掛けに接続した二つのNOR素子によるラッチの直接の類似物である。

\section{安定性とリズム}

\nt{GLUT}ネットワークの三つ目の問題、雪崩が残っている。自己帰還$w_{EE}$をもつ\nt{GLUT}ニューロンの群と、第1の群が強さ$w_{IE}$で興奮させ、第1の群を強さ$w_{EI}$で抑制する\nt{GABA}ニューロンの群をとる。発火率$E$と$I$について次を得る。
\begin{equation}
\tau_E\frac{\dd E}{\dd t} = -E + w_{EE}E - w_{EI}I + h, \qquad
\tau_I\frac{\dd I}{\dd t} = -I + w_{IE}E ,
\label{eq:ei}
\end{equation}
ここで$h$は外部入力である~\cite{WilsonCowan1972}。抑制がなく$w_{EE}>1$なら、活動は際限なく増える。1個のスパイクが次のスパイクを1個より多く生むからである。抑制があれば、定常発火率
\begin{equation}
E^* \;=\; \frac{h}{1-w_{EE}+w_{EI}w_{IE}}
\label{eq:eistar}
\end{equation}
は、分母が正なら有限である。だがそれだけでは足りない。平衡はさらに引き寄せる力をもたなければならない。式~\eqref{eq:ei}の線形解析から二つの条件が得られる。

\begin{keyblock}
\begin{proposition}[安定性の条件]
\label{prop:ei}
平衡~\eqref{eq:eistar}は、抑制が次の条件を満たせば安定である。
\begin{enumerate}
\item[(i)] \emph{十分に強い}：$w_{EI}\,w_{IE} > w_{EE}-1$。
\item[(ii)] \emph{十分に速い}：$\tau_I < \tau_E/(w_{EE}-1)$。
\end{enumerate}
(i)が破れると、活動は雪崩のように増える。(i)が満たされ(ii)が破れると、抑制が遅れ、活動は振動し始める。
\end{proposition}
\end{keyblock}

条件~(i)は系~\eqref{eq:ei}の行列の行列式、条件~(ii)はそのトレースから来る。2番目の意味は、制御器を調整したことのある人なら誰にでもわかる。遅れた帰還はずれを打ち消すのではなく、揺り動かして大きくする。

\begin{figure}[t]
\centering
\begin{tikzpicture}[>=Stealth]
\draw[->,gray!70!black] (0,0) -- (10.9,0) node[below left,black] {\small $t$, ms};
\draw[->,gray!70!black] (0,0) -- (0,3.3) node[left,black] {\small $E$};
\foreach \t in {0,50,100,150}{\draw[gray!70!black] (\t*0.07,0.06) -- (\t*0.07,-0.06) node[below,font=\scriptsize] {\t};}
\input{figures/ei_traces}
\node[font=\small,gray!40!black,anchor=west] at (0.9,3.15) {抑制なし：雪崩};
\node[font=\small,blue!60!black,anchor=west] at (7.4,1.25) {速い抑制};
\node[font=\small,red!70!black,anchor=west] at (7.4,2.55) {遅い抑制};
\end{tikzpicture}
\caption{入力をオンにした後の、帰還$w_{EE}=1.5$、$\tau_E=10\ms$をもつ\nt{GLUT}群の発火率。抑制がないと（破線）活動は際限なく増える。強さ$w_{EI}w_{IE}=2$、$\tau_I=2\ms$の抑制があると（青の曲線）、活動はレベル$E^*=h/(1-1.5+2)=0.67h$に落ち着く。同じ強さでも遅い抑制$\tau_I=30\ms$では（赤の曲線）、活動は振動する。}
\label{fig:ei}
\end{figure}

皮質はまさにこの状態で動作していると考えられている。皮質自身の帰還は、抑制がなければ雪崩に陥るほど強く、それを動作点に保っているのは速い抑制だけである~\cite{Tsodyks1997isn}。

この状態には逆説的な特徴があり、それによって外から見分けることができる~\cite{Tsodyks1997isn}。式~\eqref{eq:ei}に、\nt{GABA}群に直接入る外部入力$h_I$を加えよう。定常発火率は次のようになる。
\begin{equation}
E^* \;=\; \frac{h-w_{EI}\,h_I}{D},\qquad
I^* \;=\; \frac{w_{IE}\,h+(1-w_{EE})\,h_I}{D},\qquad
D \;=\; 1-w_{EE}+w_{EI}w_{IE} .
\label{eq:isn}
\end{equation}
$w_{EE}>1$なら、2番目の式の$h_I$の係数は負である。\nt{GABA}ニューロンを後押しすると、その発火率は\emph{下がる}。原因はループにある。余分な抑制が\nt{GLUT}群を抑え、\nt{GLUT}群は\nt{GABA}ニューロンをあまり興奮させなくなり、\nt{GABA}ニューロンは得た分より多くを失う。図~\ref{fig:ei}の値（$w_{EE}=1.5$、$w_{EI}w_{IE}=2$、$D=1.5$）では、加えた入力1単位ごとに$I^*$が3分の1単位下がる。この特徴は皮質で実際に見つかった。視覚野のニューロンの応答が周囲の刺激で抑えられるとき、ニューロンが受け取る抑制性電流は増えるのではなく、興奮性電流と一緒に減る~\cite{Ozeki2009}。その後、同じ特徴は皮質のさまざまな領域と層で見つかっている~\cite{Sanzeni2020}。

条件~(ii)が破れたときの振動も脳で見られる。数ミリ秒の遅延をもつ「\nt{GLUT}が\nt{GABA}を興奮させ、\nt{GABA}が\nt{GLUT}を抑制する」ループは、周期$12$--$50\ms$程度（周波数$20$--$80$\,Hz）のリズムを生み、このような速いリズムは皮質でよく知られている~\cite{Cardin2009,Buzsaki2012}。これはコンピュータのクロックではない。リズムは局所的で、ネットワークが活動しているときにだけ現れる。だが数サイクルの間、リズムは時間を窓に区切り、ニューロンはその中で発火できる。

\section{まとめ}

\begin{table}[t]
\centering
\small
\begin{tabular}{>{\raggedright}p{3.3cm}>{\raggedright}p{4.4cm}>{\raggedright\arraybackslash}p{4.4cm}}
\toprule
& \nt{GLUT}のみ & \nt{GLUT} + \nt{GABA} \\
\midrule
NOT & 「オン・オフ」センサを通じてのみ & 禁止$a\wedge\bar b$。背景活動があればNOT \\
1ビットあたりの配線 & 2本 & 1本 \\
運動の方向 & なし & 遅延つきの禁止 \\
ゲイン調整 & なし & 割り算：シャントと背景ノイズの組み合わせ \\
同時性の窓 & $\sim\tau$ & 抑制によって狭まる \\
多数から一つを選ぶ & なし & 相互抑制、$\beta>1$ \\
メモリのリセット & 疲労によってのみ & \nt{GABA}を介した信号で \\
安定性 & 疲労によってのみ & 速い負帰還 \\
リズム & 不要 & 遅い抑制で生じる \\
\bottomrule
\end{tabular}
\caption{\nt{GABA}が\nt{GLUT}ニューロンのネットワークに加えるもの。}
\label{tab:gaba}
\end{table}

\nt{GABA}はネットワークに新しい\emph{計算}をほとんど加えない。これらはすべて2線式センサでも計算できた。加えるのは\emph{制御}である（表~\ref{tab:gaba}）。選択、リセット、ゲイン調整、そして安定性である。脳では興奮がデータを運び、抑制がどのデータを、いつ、どれだけの大きさで通すかを決める。皮質のニューロンの4個か5個に1個にとって、これは非常に合理的な分業である。

\section{問題}

\begin{exercise}[\lvlA]
\label{ex:03:1}
\emph{数値で見るシャント。}$E_E=70\mV$とする。式~\eqref{eq:shunt}から、$g_E=g_L$と$4g_L$のとき、抑制なしとシャント$g_I=2g_L$ありの$V_\infty$を求めよ。$g_E=g_L$でシャントと同じだけ下げる引き算なら、$g_E=4g_L$での$V_\infty$はいくらか。シャントは同時性の窓を何分の1に狭めるか。
\end{exercise}

\begin{exercise}[\lvlB]
\label{ex:03:2}
\emph{安定性条件の導出。}線形系~\eqref{eq:ei}の行列のトレースと行列式を求め、そこから命題~\ref{prop:ei}の二つの条件を導け。条件~(ii)の境界では、平衡は振動的に安定性を失う。図~\ref{fig:ei}の値でこの振動の周期を求めよ。
\end{exercise}

\begin{exercise}[\lvlB]
\label{ex:03:3}
\emph{遅延から生まれるリズム。}モデル~\eqref{eq:ei}の遅い抑制を、遅延$d$をもつ速い抑制に置き換える：$\tau_E\dot E=-E(t)-GE(t-d)+h$。$\tau_E=10\ms$のとき、振動を生む最小の遅延$d_c$と、$G=2$と$G=5$での振動の周期を求めよ。問題~\ref{ex:03:2}と違って、これは皮質の速いリズム$12$--$50\ms$に入るか。
\end{exercise}

\begin{exercise}[\lvlB]
\label{ex:03:4}
\emph{シミュレーション：記憶をともなう選択。}$\beta=2$、$\tau=1$で式~\eqref{eq:wta}を積分せよ。(a)~$I_1=1$で$I_2$を$0$から$3$までゆっくり上げてから下げる。どの$I_2$でネットワークは切り替わるか。(b)~入力の差$\varepsilon$を10分の1にすると、静寂からの決定時間はどれだけ増えるか。
\end{exercise}

\begin{exercise}[\lvlB]
\label{ex:03:5}
\emph{設計：速さの帯域の検出器。}図~\ref{fig:direction}の検出器は、$A$から$B$への走行時間が$5$--$20\ms$の運動に対して黙っていなければならない。遅延$d_k$の\nt{GABA}仲介役は、$A$のスパイクの後の区間$[d_k,\,d_k+5\ms]$で$Y$を禁止する。$B$からの興奮はすぐに届く。仲介役は何個必要で、遅延はそれぞれいくらか。
\end{exercise}

\begin{exercise}[\lvlB]
\label{ex:03:6}
\emph{禁止と背景活動。}「入力ごとに\nt{GABA}仲介役、$f=1$の入力の組ごとに禁止、全体のOR」という一般的な回路で、パリティ$x\oplus y\oplus z$にはニューロンが何個必要か。入力を直接受け取る\nt{GABA}と\nt{GLUT}の2個のニューロンでこれを作り、重みと閾値を示せ。
\end{exercise}

\begin{exercise}[\lvlC]
\label{ex:03:7}
\emph{設計：100個からの選択。}$100$個の\nt{GLUT}群のそれぞれが、自分以外のすべてを等しく抑制しなければならない。線形の\nt{GABA}仲介役は群から興奮を受けて群を抑制する。群どうしは互いに興奮させてよいが、自分自身は興奮させない。仲介役の最小数はいくつか。直接の興奮性結合がまったくない場合はどうか。
\end{exercise}

\begin{exercise}[\lvlC]
\label{ex:03:8}
\emph{研究：抑制が逆説的になるとき。}図~\ref{fig:ei}のモデル（$w_{EI}=2$、$w_{IE}=1$）で、\nt{GLUT}群の伝達関数を$f(u)=[u]_+^2$とし、\nt{GABA}群は線形のままとする。\nt{GABA}群への入力の追加がその群自身の発火率を下げるのは、入力$h$がどの値$h_c$を超えたときか。シミュレーションで確かめよ。
\end{exercise}

\chapter{モールス符号：\nt{GLUT}と\nt{GABA}のニューロンが話す言語について何がわかっているか}
\chaptermark{モールス符号}
\label{sec:code}

\section{記号が一つだけのアルファベット}

モールス符号には短点と長点（トンとツー）の二つの記号があり、その間に3種類の長さの休止がある。第\ref{sec:neurontypes}章で見たように、ニューロンのアルファベットはさらに貧しく、記号は一つしかない。スパイクは約1ミリ秒続き、1個のニューロンのスパイクは高さも形もすべて同じである。そのためメッセージのすべては休止に、つまり時刻の列$t_1,t_2,t_3,\dots$にある。長点のないモールス符号で、意味を運ぶのは短点のリズムだけである（図~\ref{fig:morse}）。

\begin{figure}[t]
\centering
\begin{tikzpicture}[>=Stealth,x=0.08cm,y=1cm]
% Morse letter: dot, dash, dot
\node[left,font=\small] at (-6,2.55) {モールス：};
\fill[black] (0,2.45) rectangle (6,2.65);
\fill[black] (14,2.45) rectangle (32,2.65);
\fill[black] (40,2.45) rectangle (46,2.65);
\node[right,font=\small,gray!40!black] at (52,2.55) {トン、ツー、トン：意味は長さと休止にある};
% stimulus onset
\node[left,font=\small] at (-6,0.4) {ニューロン：};
\draw[->,thick,green!45!black] (0,-0.55) -- (0,-0.05);
\node[below,font=\scriptsize,green!45!black] at (0,-0.55) {刺激};
\draw[gray!70!black] (0,0) -- (152,0) node[right,black] {\small $t$};
% spikes: single ones blue, the burst red
\foreach \s in {14,26,41,88,112,131}{\draw[very thick,blue!60!black] (\s,0) -- (\s,0.8);}
\foreach \s in {60,63,66,69}{\draw[very thick,red!70!black] (\s,0) -- (\s,0.8);}
% latency
\draw[<->,thick] (0,1.1) -- node[above,font=\scriptsize] {$t_1$：最初のスパイクの時刻} (14,1.1);
% burst
\draw[decorate,decoration={brace,amplitude=4pt},red!70!black] (58,0.95) -- (71,0.95) node[midway,above=4pt,font=\scriptsize] {バースト＝「ツー」};
% counting windows
\foreach \a/\b/\n in {0/50/3,50/100/5,100/150/2}{
  \draw[decorate,decoration={brace,amplitude=4pt,mirror},gray!50!black] (\a+1,-0.2) -- (\b-1,-0.2) node[midway,below=4pt,font=\small] {$n=\n$};}
\node[font=\scriptsize,gray!40!black] at (75,-1.05) {レート符号：各窓のスパイク数$n$};
\foreach \x in {0,50,100,150}{\draw[gray!60!black] (\x,-0.06) -- (\x,0.06);}
\end{tikzpicture}
\caption{モールス符号とニューロンの符号。同じスパイク列もいろいろに読める。$50\ms$の窓ごとにスパイクを数える（\ref{sec:rate}節）、遅延$t_1$~\eqref{eq:latency}を測る、あるいはバースト、つまり「ツー」~\eqref{eq:burst}に気づく。}
\label{fig:morse}
\end{figure}

そもそもどんな休止がありうるのか。下限は不応期で決まる。スパイクの後、ニューロンは約1ミリ秒の間は再び発火できないので、毎秒数百スパイクを超えることはない。上限は何もない。ニューロンは何分も黙っていられる。皮質の\nt{GLUT}ニューロンの発火率は毎秒1スパイク未満から数十スパイクの範囲にあり、ほとんどのニューロンはほとんどの時間を下限の近くで動作している。

第\ref{sec:glutcomputer}章と第\ref{sec:gaba}章では、スパイクで数を書き表すあれこれの方法を途中でその都度仮定してきた。ここで正面から問おう。\nt{GLUT}ニューロンと\nt{GABA}ニューロンは、互いにどんな言語で話しているのか。完全な答えはない。この言語は解読されていない。だが、わかっていることもいくつか確かにあり、そのほとんどは通信技術者のように考えれば導ける。通信路容量、雑音、受信機、伝送のコストである。

\section{どれだけ伝えられるか}

上限から始めよう。受け手がスパイクの時刻を精度$\Delta t$で区別するとする。$\Delta t$より小さなずれは受け手には区別できない。時間を不応期より短い長さ$\Delta t$のセルに刻み、各セルにはスパイクが1個あるかないかのどちらかだとする（図~\ref{fig:cells}）。平均発火率$r$では、スパイクがセルに入る確率は$p=r\Delta t$である。流れが最も多くの情報を運ぶのはセルが独立なときで、そのとき各セルはエントロピー$-p\log_2p-(1-p)\log_2(1-p)\approx p\log_2(e/p)$ビット（$p\ll1$）を与える。$\Delta t$で割ると、伝送速度の上限と、そのスパイク1個あたりの値が得られる~\cite{MacKay1952}。
\begin{empheq}[box=\keybox]{equation}
R_{\max} \;\approx\; r\,\log_2\frac{e}{r\,\Delta t}\ \ \text{bit/s},
\qquad
\frac{R_{\max}}{r} \;\approx\; \log_2\frac{e}{r\,\Delta t}\ \ \text{bit/スパイク}.
\label{eq:spikeentropy}
\end{empheq}
$r=10$スパイク/秒、$\Delta t=1\ms$では、スパイク1個あたり約8ビット、毎秒約80ビットである。

\begin{figure}[t]
\centering
\begin{tikzpicture}[>=Stealth,x=0.5cm,y=1cm]
% time cut into cells of width Delta t; a cell holds one impulse or none
\foreach \k in {0,...,23}{\draw[gray!55] (\k,0) rectangle (\k+1,0.7);}
\foreach \k in {2,7,8,15,21}{
  \fill[red!10] (\k,0) rectangle (\k+1,0.7);
  \draw[very thick,red!70!black] (\k+0.5,0.05) -- (\k+0.5,0.65);}
\foreach \k in {0,...,23}{
  \pgfmathtruncatemacro{\bit}{(\k==2)||(\k==7)||(\k==8)||(\k==15)||(\k==21)}
  \ifnum\bit=1 \def\bitcol{red!70!black}\else \def\bitcol{gray!60!black}\fi
  \node[font=\scriptsize,text=\bitcol] at (\k+0.5,-0.25) {\bit};}
\draw[<->,gray!60!black] (0,0.95) -- (1,0.95) node[midway,above,font=\scriptsize,black] {$\Delta t$};
\node[font=\small,anchor=east] at (-0.3,0.35) {スパイク};
\node[font=\small,anchor=east] at (-0.3,-0.25) {ビット};
\draw[decorate,decoration={brace,amplitude=5pt,mirror},gray!60!black] (0,-0.55) -- (24,-0.55)
  node[midway,below=6pt,font=\scriptsize,black,align=center] {数えるだけの受け手：「$n=5$」\\ セルが見える受け手：$24$個のうちどの$5$個か、つまり$\log_2\binom{24}{5}\approx15$ビット};
\end{tikzpicture}
\caption{容量~\eqref{eq:spikeentropy}はどこから来るか。時間を長さ$\Delta t$のセルに刻むと、各セルにはスパイクがあるか（$1$）ないか（$0$）のどちらかであり、スパイク列は2進の文字列になる。スパイクがセルに入る確率は$p=r\Delta t$である。セルが細かいほど文字列は長くなり、ビットも増える。スパイクを数えるだけのカウンタは、1が\emph{どこに}あるかに書かれたもののほとんどすべてを失う。}
\label{fig:cells}
\end{figure}

スパイク1個あたりのビット数は、時間精度に対数的にしか依存しない。精度$10\ms$でも、スパイク1個あたり約5ビットが残る。しかし受け手が1秒間のスパイクを数えるだけなら、$r=10$では1秒全体で2ビット未満しか得られない（\ref{sec:rate}節）。

\begin{keyblock}
\begin{proposition}[スパイク通信路の容量]
\label{prop:capacity}
平均発火率$r$のスパイク列を時間精度$\Delta t$で読むとき、運べる情報は~\eqref{eq:spikeentropy}を超えない。この容量のほとんどすべてはスパイクの\emph{時刻}にある。スパイクを数えるだけの受け手が同じ列から引き出すものは、時刻を1ミリ秒の精度で区別する受け手の数十分の1である。
\end{proposition}
\end{keyblock}

これは上限であって測定値ではない。刺激がわかっている感覚ニューロンでの測定では、スパイク1個あたり1ビットから数ビットであり、最もよい場合でその精度について計算した上限の半分ほどに達する~\cite{Rieke1997}。つまり、少なくとも神経系の入口では、スパイクの時刻は失われずに使われている。

\section{雑音のある通信路}

容量~\eqref{eq:spikeentropy}は雑音なしで計算した。雑音がどこで生じるかについては、三つの事実が確かにわかっている。

\emph{スパイク発生器そのものは精確である。}脳から組織片ごと取り出した皮質ニューロンに、時間的に速く変化する同じ電流を何度も与えると、スパイクは約1ミリ秒の精度で再現される~\cite{Mainen1995}。電流が一定なら、スパイクの時刻は試行ごとにばらけていく。精確な時刻を生むのはニューロン自体ではなく、入力の速い変化である。

\emph{シナプスは当てにならない。}\nt{GLUT}シナプスに達したスパイクは、ある確率$p$でしか神経伝達物質のひとまとまりを放出しない。海馬のシナプスでは半分以上のスパイクが受け手にまったく届かず、その原因は伝導の失敗ではなく、まさに放出のランダム性である~\cite{AllenStevens1994}。通信技術者にとって、これは消失通信路である。二つのニューロンの間にシナプスが複数あれば、ある程度は救われる。

\emph{生きた皮質では、スパイク列はランダムに見える。}スパイク間隔は、ランダムなポアソン過程とほぼ同じようにばらつく。同じ刺激を繰り返すと、窓内のスパイク数は分散が平均に近くなるように変動する~\cite{Softky1993}。

精確な素子がどうしてランダムな列を生むのか。ニューロンは数千の入力を受け取り、そのそれぞれが当てにならず、しかも第\ref{sec:gaba}章のように興奮と抑制が釣り合っている。膜電位は閾値付近で揺らぎ、スパイクの時刻はその揺らぎで決まる。これが雑音なのか、それとも私たちが読み方を知らないメッセージなのかは、大部分が未解決の問題である。「雑音」の一部はすでにメッセージだとわかった。視覚野では、ばらつきのかなりの部分が動物自身の運動に追従している~\cite{Stringer2019}。ここでの困難は原理的なものである。よく圧縮されたメッセージは、符号を知らない者にとってはランダムな列と区別がつかない。

\section{レート：遅いが確実}
\label{sec:rate}

最も単純な符号は\emph{レート符号}である。数は、ニューロンが窓$T$の間に出すスパイクの数に書かれている。消失も時刻の揺らぎも、この符号には怖くない。だが代償がある。ポアソン列なら、窓内のスパイク数$n$は平均$rT$、ばらつき$\sqrt{rT}$をもつ。
\begin{empheq}[box=\keybox]{equation}
\frac{\sigma_n}{\bar n} \;=\; \frac{1}{\sqrt{r\,T}} .
\label{eq:rateerror}
\end{empheq}
発火率を10パーセントの精度で知るには100個のスパイクが必要で、毎秒20スパイクなら5秒かかる。隣り合うレベルはばらつきの2倍ほど離れていれば区別できるので、発火率$r$までの範囲で時間$T$の間に区別できるレベルは約$\sqrt{rT}$個である。$r=10$、$T=1\,\text{s}$では3--4レベルで、2ビットに満たない。
脳はこんなに遅くは動かない。ヒトは一瞬だけ見せられた画像の中の動物を認識し、脳内の応答は約$150\ms$後に現れる~\cite{Thorpe1996}。粗く見積もって、この間に信号は約10段の逐次的な処理段を通り、1段あたり$10$--$20\ms$である。皮質のふつうの発火率では、各ニューロンは1段の間に0個か1個のスパイクしか出せず、その窓での発火率は定まらない。出口は二つある。多数のニューロンを使うか、時刻を見るかである。

\emph{多数のニューロン。}$N$個のニューロンが同じ数を運び、受け手がそれらのスパイクを足し合わせるとする。雑音が独立なら、式~\eqref{eq:rateerror}の$rT$は$NrT$に置き換わる。$r=20$の100個のニューロンは$15\ms$で30個のスパイクを出し、精度は約20パーセントになる。空間が時間の代わりをする。しかし隣り合う皮質ニューロンの雑音は独立ではない。隣どうしのペアのスパイク数の相関係数は、ふつう約$0.1$である~\cite{Zohary1994}。それが$c$なら、$N$個のニューロンの平均の分散は次のようになる。
\begin{empheq}[box=\keybox]{equation}
\sigma^2_N \;=\; \sigma^2_1\,\frac{1+(N-1)\,c}{N}
\;\xrightarrow[N\to\infty]{}\; c\,\sigma^2_1 .
\label{eq:correlated}
\end{empheq}

\begin{keyblock}
\begin{proposition}[平均化の限界]
\label{prop:pooling}
雑音が相関していると、群での平均化は誤差を$1/\sqrt{c}$分の1より小さくはできない。$c=0.1$ではこれは3分の1であり、この改善は数十個のニューロンですでに達成される。それ以上加えても無駄である。
\end{proposition}
\end{keyblock}

相関がすべて同じように有害なわけではない。精度を制限するのは、共通の雑音のうち\emph{信号そのものに似た}部分、つまり測定される量が変化したときと同じように群全体をずらす部分だけである~\cite{MorenoBote2014}。脳に実際にそのような雑音がどれだけあるかは、まだ調べられているところである。

多数のニューロンを使う方法はもう一つある。音の方向の検出器（図~\ref{fig:itd}）のように、\emph{どの}ニューロンが発火したかに数を書くことである。これは\emph{場所}による符号で、知られているうちで最も確実なものである。感覚ニューロンでは、配線の番号が皮膚のどの点に触れたか、音の高さはいくらかを告げ、これは何度も確かめられている。ニューロンの数の点では高価だが、速い。メッセージ1個に遅延1段で済む。

\section{時刻：速いが、どこから測るか}

二つ目の出口は、スパイクの時刻に数を書くことである。ニューロン~\eqref{eq:glutneuron}をとり、時刻$t=0$から、電圧をレベル$U$まで上げるような一定の入力を与える。電圧は$V=U\bigl(1-e^{-t/\tau}\bigr)$で増え、次の時刻に閾値$\theta$に達する。
\begin{empheq}[box=\keybox]{equation}
t_1 \;=\; \tau\,\ln\frac{U}{U-\theta}, \qquad U>\theta .
\label{eq:latency}
\end{empheq}
強い入力なら早いスパイク、弱い入力なら遅いスパイク、閾値以下ならスパイクなしである。$\tau=10\ms$では、入力$U=4\theta$で最初のスパイクは$2.9\ms$後、$U=2\theta$で$6.9\ms$後、$U=1.2\theta$で$17.9\ms$後に来る。たった1個のスパイクが数を運ぶ。

しかし$t_1$はどの時刻から測るのか。共通の時刻基準はなく、受け手は刺激がいつ始まったかを知らない。答えは三つある。

\emph{相対遅延。}二つのニューロンが同じ刺激を見るなら、その遅延の差$t_1^A-t_1^B$は入力だけで決まり、開始時刻にはよらない。時刻を比べるのは遅延つきの同時検出器である（図~\ref{fig:itd}）。$N$個のニューロンが1個ずつスパイクを出したなら、その到着の\emph{順序}そのものが最大$\log_2N!$ビットを運ぶ。10個のニューロンなら約22ビットであり、「発火したかどうか」という答えでは10ビット以下にしかならない。網膜では、出力ニューロンの最初のスパイクの相対遅延から画像を速く確実に復元できることが示されている~\cite{Gollisch2008}。

\emph{語頭としての一斉発火。}刺激の急な変化は、多数のニューロンをほぼ同時に発火させる。この一斉発火そのものが、モールス符号の語と語の間の長い休止のような開始の目印になり、受け手はそこから遅延を測る。

\emph{局所的なリズム。}第\ref{sec:gaba}章では、\nt{GLUT}--\nt{GABA}ループが周期$12$--$50\ms$の局所的なリズムを生んだ。これはクロックではないが、その1サイクルの中で、強い入力をもつニューロンは弱い入力をもつニューロンより早く発火でき、式~\eqref{eq:latency}は\emph{位相}による符号に変わる。数は、サイクルのどの部分でスパイクが来たかに書かれる。このような符号は観察されている。空間内の特定の場所を担当するニューロンは、動物が「自分の」場所を通過するにつれて、遅い局所リズムのサイクルのますます早い時点で発火する~\cite{OKeefe1993}。脳が位相をどれほど広く使っているかは、まだわかっていない。

\section{符号を選ぶのは受け手}

スパイク列には、発火率も、時刻も、順序もある。そのうちどれが読まれるかは受け手が決め、しかもただ一つのパラメータ、自分の時定数$\tau$で決める。

式~\eqref{eq:glutneuron}を時間で平均しよう。ニューロンに発火率$r_j$の独立な列が届くなら、平均電圧とその相対的な揺らぎは次のようになる。
\begin{empheq}[box=\keybox]{equation}
\bar V \;=\; \tau\sum_j w_j\,r_j ,
\qquad
\frac{\sigma_V}{\bar V} \;\approx\; \frac{1}{\sqrt{2\,r_{\text{in}}\,\tau}} ,
\label{eq:reader}
\end{empheq}
ここで2番目の式は重みが等しい場合のもので、$r_{\text{in}}$はすべての入力の発火率の合計である。$\tau$が大きいと、電圧はほとんど揺らがず、発火率の重みつき和に追従する。受け手は\emph{積分器}であり、発火率を読んで時刻を忘れる。$\tau$が小さいと、電圧はスパイクとスパイクの間に下がりきり、閾値に達するのは複数のスパイクが窓~\eqref{eq:window}の中に届いたときだけである。受け手は\emph{同時検出器}であり、時刻を読む（図~\ref{fig:reader}）。毎秒5スパイクの入力1000本、$\tau=10\ms$では揺らぎは約10パーセントで、ほぼ純粋な積分である。第\ref{sec:gaba}章のように抑制が$\tau$を$2\ms$に縮めると、揺らぎは20パーセント余りに増え、ニューロンは同時性に気づき始める。

\begin{figure}[t]
\centering
\begin{tikzpicture}[>=Stealth,x=0.115cm,y=1cm]
% common input: the same spike train for both receivers
\foreach \s in {6,21,22,38,52,55.5,59,62.5,66,69.5,73,76.5,80,83.5,87,90.5,94,97.5}{\draw[thick,gray!40!black] (\s,5.25) -- (\s,5.65);}
\draw[gray!70!black] (0,5.25) -- (105,5.25);
\node[left,font=\small] at (-1,5.45) {入力};
\node[above,font=\scriptsize,gray!40!black] at (13.5,5.65) {まばら};
\node[above,font=\scriptsize,gray!40!black] at (21.5,6.0) {ペア};
\node[above,font=\scriptsize,gray!40!black] at (75,5.65) {高頻度};
% axes and thresholds
\foreach \y/\th in {2.7/2.5,0/1.5}{
  \draw[gray!70!black] (0,\y) -- (106,\y) node[right,black] {\small $t$};
  \draw[densely dashed,gray!60!black] (0,\y+0.6*\th) -- (104,\y+0.6*\th) node[right,black,font=\small] {$\theta$};
}
\node[anchor=south west,font=\small] at (0,4.3) {$\tau=10\ms$：積分器};
\node[anchor=south east,font=\small] at (104,1.0) {$\tau=2\ms$：同時検出器};
% integrator: tau=10 ms, theta=2.5w
\begin{scope}[yshift=2.7cm]
\foreach \a/\b/\v in {0/6/0.000,6/21/1.000,21/22/1.223,22/38/2.107,38/52/1.425,52/55.5/1.351,55.5/59/1.952,59/62.5/2.376,62.5/66/0.000,66/69.5/1.000,69.5/73/1.705,73/76.5/2.201,76.5/80/0.000,80/83.5/1.000,83.5/87/1.705,87/90.5/2.201,90.5/94/0.000,94/97.5/1.000,97.5/104/1.705}{\draw[very thick,blue!60!black] plot[domain=\a:\b,samples=14] (\x,{0.6*\v*exp(-(\x-\a)/10)});}
\foreach \s/\p/\q in {6/0.000/1.000,21/0.223/1.223,22/1.107/2.107,38/0.425/1.425,52/0.351/1.351,55.5/0.952/1.952,59/1.376/2.376,62.5/1.674/2.500,62.5/2.500/0.000,66/0.000/1.000,69.5/0.705/1.705,73/1.201/2.201,76.5/1.551/2.500,76.5/2.500/0.000,80/0.000/1.000,83.5/0.705/1.705,87/1.201/2.201,90.5/1.551/2.500,90.5/2.500/0.000,94/0.000/1.000,97.5/0.705/1.705}{\draw[very thick,blue!60!black] (\s,0.6*\p) -- (\s,0.6*\q);}
\foreach \s in {62.5,76.5,90.5}{\draw[->,thick,red!70!black] (\s,1.58) -- (\s,1.95);}
\end{scope}
% coincidence: tau=2 ms, theta=1.5w
\begin{scope}[yshift=0cm]
\foreach \a/\b/\v in {0/6/0.000,6/21/1.000,21/22/1.001,22/38/0.000,38/52/1.000,52/55.5/1.001,55.5/59/1.174,59/62.5/1.204,62.5/66/1.209,66/69.5/1.210,69.5/73/1.210,73/76.5/1.210,76.5/80/1.210,80/83.5/1.210,83.5/87/1.210,87/90.5/1.210,90.5/94/1.210,94/97.5/1.210,97.5/104/1.210}{\draw[very thick,blue!60!black] plot[domain=\a:\b,samples=14] (\x,{0.6*\v*exp(-(\x-\a)/2)});}
\foreach \s/\p/\q in {6/0.000/1.000,21/0.001/1.001,22/0.607/1.500,22/1.500/0.000,38/0.000/1.000,52/0.001/1.001,55.5/0.174/1.174,59/0.204/1.204,62.5/0.209/1.209,66/0.210/1.210,69.5/0.210/1.210,73/0.210/1.210,76.5/0.210/1.210,80/0.210/1.210,83.5/0.210/1.210,87/0.210/1.210,90.5/0.210/1.210,94/0.210/1.210,97.5/0.210/1.210}{\draw[very thick,blue!60!black] (\s,0.6*\p) -- (\s,0.6*\q);}
\foreach \s in {22}{\draw[->,thick,red!70!black] (\s,0.98) -- (\s,1.35);}
\end{scope}
\foreach \x in {0,50,100}{\draw[gray!60!black] (\x,-0.06) -- (\x,0.06) node[below,font=\scriptsize] {\x\,ms};}
\end{tikzpicture}
\caption{同じ入力と二つの異なる受け手。各スパイクは電圧を$w$だけ上げ、その後はニューロンのモデル~\eqref{eq:glutneuron}のように電圧が流れ落ちる。赤い矢印は受け手のスパイクである。遅い受け手（$\tau=10\ms$、$\theta=2.5w$）はスパイクが\emph{多い}ところで発火し、ペアには気づかない。速い受け手（$\tau=2\ms$、$\theta=1.5w$）は窓~\eqref{eq:window}の中のペアにだけ発火し、高頻度でも同時でない列は見逃す。}
\label{fig:reader}
\end{figure}

測定によれば、活動している皮質では膜のコンダクタンスが静止時の数倍に増える~\cite{Destexhe2003}。つまりそこでは$\tau$が短く、動作中の皮質ニューロンは積分器よりも同時検出器に近い。

\begin{keyblock}
\begin{proposition}[符号を決めるのは受け手]
\label{prop:reader}
同じスパイク列が、遅い受け手にとっては発火率を、速い受け手にとっては時刻を、神経伝達物質が放出される体積にとっては数秒にわたって平滑化された濃度レベルを運ぶ（第\ref{sec:serocomputer}章）。どの符号が使われるかを決めるのは送り手ではなく受け手の時定数であり、それを抑制が調整する。
\end{proposition}
\end{keyblock}
\section{トンとツー：フィルタとしてのシナプス}

これまでシナプスは重みつきの配線にすぎなかった。実際には、シナプスの応答はそれ以前のスパイクに依存し、それがニューロンの言語に第二の記号を与える。

\emph{枯渇：シナプスは変化を伝える。}放出のたびに、放出準備のできた神経伝達物質の蓄えの割合$U$が使われ~\cite{Tsodyks1997}、蓄えは数百ミリ秒程度の時定数$\tau_{\text{rec}}$で回復する。発火率$r$では、スパイク1個への応答の振幅はほぼ$A(r)=A_0/(1+U r\tau_{\text{rec}})$に落ち着く。ここで$A_0$は休んだシナプスの応答である。受け手が単位時間に受け取る電流は次のようになる。
\begin{empheq}[box=\keybox]{equation}
r\,A(r) \;=\; \frac{A_0\,r}{1+U r\,\tau_{\text{rec}}}
\;\xrightarrow[r\to\infty]{}\; \frac{A_0}{U\tau_{\text{rec}}} .
\label{eq:depression}
\end{empheq}
$U=0.5$、$\tau_{\text{rec}}=0.8\,\text{s}$では、飽和は$1/(U\tau_{\text{rec}})=2.5$スパイク/秒からすでに始まる。発火率が5から20に増えても、定常電流は3分の1しか増えない。そのかわり、新しい発火率での最初のスパイクは蓄えがまだ以前のままのうちに届くので、電流は一瞬4倍に跳ね上がり、その後$\tau_{\text{rec}}/(1+Ur\tau_{\text{rec}})\approx90\ms$で減衰する（図~\ref{fig:depression}）。

このようなシナプスが伝えるのは発火率ではなくその\emph{変化}、本質的には相対変化$\Delta r/r$である~\cite{Abbott1997depression}。電子回路の技術者にとって、これは配線に直接組み込まれたハイパスフィルタである。

\begin{figure}[t]
\centering
\begin{tikzpicture}[>=Stealth,x=12.5cm,y=1cm]
% input rate
\draw[->,gray!70!black] (0,2.6) -- (0.9,2.6) node[right,black] {\small $t$};
\draw[->,gray!70!black] (0,2.6) -- (0,4.3) node[above,black] {\small $r$};
\draw[very thick,blue!60!black] (0,2.9) -- (0.2,2.9) -- (0.2,3.8) -- (0.8,3.8);
\node[left,font=\scriptsize] at (0,2.9) {$5$};
\node[left,font=\scriptsize] at (0,3.8) {$20$};
\node[right,font=\small,blue!60!black] at (0.4,4.1) {シナプスへの入力発火率};
% output current
\draw[->,gray!70!black] (0,0) -- (0.9,0) node[right,black] {\small $t$};
\draw[->,gray!70!black] (0,0) -- (0,2.45) node[left,black] {\small $rA$};
\draw[densely dashed,gray!60!black] (0,0.667) -- (0.8,0.667);
\draw[very thick,red!70!black] (0,0.5) -- (0.2,0.5) -- (0.2,2.0)
   plot[domain=0.2:0.8,samples=60] (\x,{0.667+(2.0-0.667)*exp(-(\x-0.2)/0.089)});
\node[left,font=\scriptsize] at (0,0.5) {$1.7$};
\node[left,font=\scriptsize] at (0,2.0) {$6.7$};
\node[right,font=\scriptsize] at (0.8,0.667) {$2.2$};
\node[right,font=\small,red!70!black] at (0.28,1.6) {受け手への電流：跳ね上がり、$\approx90\ms$で減衰};
\draw[gray!60!black] (0.2,-0.08) -- (0.2,0.08); \node[below,font=\scriptsize] at (0.2,0) {$0.2\,\text{s}$};
\draw[gray!60!black] (0.8,-0.08) -- (0.8,0.08); \node[below,font=\scriptsize] at (0.8,0) {$0.8\,\text{s}$};
\end{tikzpicture}
\caption{式~\eqref{eq:depression}による枯渇するシナプス、$U=0.5$、$\tau_{\text{rec}}=0.8\,\text{s}$。電流は毎秒$A_0$単位で表し、破線は定常電流である。定常電流は$1.7$から$2.2$に増えただけだが、跳躍の瞬間には$6.7$に達する。シナプスが受け手に伝えるのは、発火率がいくらかではなく、発火率が変わったということである。}
\label{fig:depression}
\end{figure}

\emph{促通：ツーとしてのバースト。}別のシナプスでは放出確率は小さいが、スパイクのたびに数十ミリ秒にわたって上がる。このようなシナプスでは、単発のスパイクはたいてい失われる。ところが\emph{バースト}、つまり数ミリ秒間隔の連続した数個のスパイクは、ほぼ確実に通過する。促通がなくても、バーストの$k$個のスパイクがすべて失われる確率は次のとおりである。
\begin{empheq}[box=\keybox]{equation}
P_{\text{loss}} \;=\; (1-p)^k ,
\label{eq:burst}
\end{empheq}
$p=0.2$では、単発のスパイクで$0.8$、4個のバーストで$0.41$である。促通はこの確率をさらに下げる。こうしてニューロンに第二の記号ができる。単発のスパイクがトン、バーストがツーである。枯渇するシナプスは休止の後の最初のスパイクを最もよく伝え、促通するシナプスはバーストを通して単発のトンを消す。バーストがより確実に伝わることは測定されている。脳が実際にそれを独立した記号として使っているというのは、もっともらしい仮説である~\cite{Lisman1997}。

\section{\nt{GABA}ニューロンはどう話すか}

皮質で最も数の多い\nt{GABA}ニューロンは速いタイプである~\cite{Tremblay2016}。そのスパイクは\nt{GLUT}ニューロンのものより短く、毎秒数百の頻度で一様に長時間出し続けることができ、ほとんど疲れない。そのシナプスはより確実である。各受け手との結合は多数の放出部位からなり、スパイクが失われることはほとんどない。その出力は受け手のスパイクが生まれる場所の近くにつくので、受け手が入力をどう足し合わせるかではなく、受け手が発火するかどうか、いつ発火するかを制御する。

速い\nt{GABA}ニューロン自身は近隣の多くの\nt{GLUT}ニューロンから興奮を受け取り、周囲の活動の総和を伝える。まさに第\ref{sec:gaba}章の割り算に必要なものである。さらに、速い\nt{GABA}ニューロンは互いに結合していて一緒に発火しやすい。前に述べた局所的なリズムを作っているのはこれらである~\cite{Cardin2009}。

モールス符号の言葉でいえば、\nt{GABA}ニューロンが受け持つのは文字ではなく\emph{休止}である。\nt{GLUT}スパイクの連続した流れを受け手が発火できる窓に区切り、同時性の窓を狭め、流れが大きくなりすぎると全体を抑える。

\begin{keyblock}
\begin{proposition}[役割の分担]
\label{prop:punctuation}
\nt{GLUT}ニューロンはメッセージの内容、つまり\emph{何を}と\emph{どれだけ}を運ぶ。速い\nt{GABA}ニューロンはその区切り、つまり\emph{いつ}話してよいか、\emph{どれだけの大きさで}話すかを運ぶ。もう一つのクラスは抑制するものを抑制することで、\emph{誰に対して}ゲートを開くかを決める。この分担の個々の部分は測定されている。一般規則としては作業仮説である。
\end{proposition}
\end{keyblock}

ほかに、出力が受け手の個々の入力につく、もっと遅い\nt{GABA}ニューロンもある。これは第\ref{sec:gaba}章の禁止~\eqref{eq:veto}として働く。ほかには触れずに、\nt{GLUT}スパイクの一つの流れを閉じる。

\nt{GABA}ニューロンを速いものと遅いものに分けるのは古い描像であり、不完全である。現在、皮質では少なくとも三つの大きなクラスが区別されており~\cite{Tremblay2016}、3番目のクラスは意外な構造をもつ。それは\nt{GLUT}ニューロンではなく別の\nt{GABA}ニューロン、とりわけ入力を閉じるものを抑制する~\cite{Pfeffer2013}。抑制の抑制は二重否定である。このようなニューロンは、受け手に興奮性のスパイクを一つも送らずにゲートを\emph{開く}。これをオンにするのは皮質の他の領域からの信号とブロードキャストチャネルなので、これはネットワークの一区画全体のイネーブル入力として働く。それが活動している間、禁止は解除され、流れが通る~\cite{Pi2013}。付録~\ref{sec:othermediators}では、この手法を脱抑制と呼んでいる。
\section{経済的な言語}

ニューロンの言語について確かにわかっている最後のことは、それが高価だということである。スパイクは、それが引き起こすシナプスの仕事と合わせて、約$10^9$個のATP分子を要し（げっ歯類の皮質での推定~\cite{Attwell2001}）、分子1個は約$10^{-19}$ジュールを与える。したがってスパイク1個のコストは$E_{\text{sp}}\sim10^{-10}$ジュールのオーダーである。脳全体は約$20$ワットを消費し、その半分、$P\sim10$ワットが信号の伝達に使われるとすると、ニューロン1個の平均発火率は次のようになる。
\begin{empheq}[box=\keybox]{equation}
\bar r \;\approx\; \frac{P}{N\,E_{\text{sp}}}
\;\sim\; \frac{10\ \text{W}}{8.6\cdot10^{10}\cdot10^{-10}\ \text{J}}
\;\approx\; 1\ \text{スパイク/秒}.
\label{eq:energy}
\end{empheq}
皮質についてのより詳しい推定では、さらに小さな値になる~\cite{Lennie2003}。脳の符号は\emph{スパース}でなければならない。速く動作できるのはニューロンのごく一部だけである。

式~\eqref{eq:spikeentropy}からわかるように、これはそれほど悪いことではない。精度$1\ms$では、毎秒100スパイクで動作するニューロンのスパイクは5ビット以下しか運ばないが、毎秒1スパイクで動作するニューロンのスパイクは最大11ビットを運ぶ。スパイクに代価を払うなら、まれに話すほうが得である。皮質は実際に精度をエネルギーと引き換えにする。食物が不足すると、ニューロンはスパイクの発火率を保ったままシナプス電流への支出を減らし、その応答は精度が下がる~\cite{Padamsey2022}。

モールス符号では、頻繁に使う文字は短い。英語の文章で最も頻繁な文字Eは短点1個である。ニューロンの符号は、わかっている限り、同じ規則に別の形で従っている。予想どおりのことにはスパイクをあまり使わない~\cite{Barlow1961}。一定の刺激のもとでニューロンは発火率を徐々に下げ、第\ref{sec:glutcomputer}章のセンサはレベルではなく変化に応答し、第\ref{sec:neurontypes}章の\nt{DOPA}ニューロンは報酬予測誤差だけを伝える。

\begin{keyblock}
\begin{proposition}[経済性]
\label{prop:economy}
エネルギーはニューロンの平均発火率を毎秒1スパイク程度に制限する。そのため\nt{GLUT}ニューロンの言語はスパースであり、スパイクをニュース、つまり変化したことや予測されなかったことに使う。エネルギーの限界は測定されている。脳の符号が1ジュールあたりのビット数を最大にするよう調整されているというのは、十分な根拠のある仮説である。
\end{proposition}
\end{keyblock}

\section{まとめ}

\begin{table}[t]
\centering
\footnotesize
\setlength{\tabcolsep}{4pt}
\renewcommand{\arraystretch}{1.15}
\begin{tabular}{>{\raggedright}p{2.6cm}>{\raggedright}p{2.3cm}>{\raggedright}p{3.0cm}>{\raggedright}p{2.0cm}>{\raggedright\arraybackslash}p{3.6cm}}
\toprule
書き方 & 運ぶもの & 読む者 & メッセージ1個の時間 & わかっていること \\
\midrule
発火したニューロンの番号（場所による符号） & どの値か & あらゆる受け手、同時検出器 & 遅延1段 & 感覚ニューロンと音の方向の検出で確立 \\
1個のニューロンの発火率 & 量 & $\tau$の大きい積分器、体積（第\ref{sec:serocomputer}章） & 数百ms--数秒 & 確立。$10$--$20\ms$の処理段には遅すぎる \\
群の発火率 & 量 & 数千の入力をもつニューロン & $10$--$50\ms$ & 確立。改善は相関~\eqref{eq:correlated}で制限される \\
最初のスパイクの時刻、順序 & 量、順序 & 遅延つきの同時検出器 & 遅延1段 & 神経系の入口で確立。皮質では議論中 \\
局所リズムの中の位相 & 量、位置 & 共通のリズムをもつニューロン & サイクル$12$--$50\ms$以上 & 一部の領域で観察。一般的な役割は仮説 \\
バースト（「ツー」） & 「重要」：確実な配送 & 促通するシナプス & $\sim10\ms$ & 確実性は測定済み。独立の記号というのは仮説 \\
発火率の変化 & ニュース & 枯渇するシナプス & $\sim0.1\,\text{s}$ & 確立 \\
速い\nt{GABA}ニューロンのスパイク & いつ、どれだけの大きさで & 近隣のすべての\nt{GLUT}ニューロン & $1$--$30\ms$ & リズムと割り算は測定済み。規則としての「句読点」は仮説 \\
\bottomrule
\end{tabular}
\caption{\nt{GLUT}ニューロンと\nt{GABA}ニューロンがスパイクでメッセージを書き表す方法。右の列は、測定されたことと推測されていることを区別している。}
\label{tab:code}
\end{table}

表~\ref{tab:code}は私たちが知っていることをまとめたものであり、知らないこともそこから見えてくる。皮質が群の発火率だけでなく、スパイクの精確な時刻をどれだけ使っているかはわからない。ニューロンに「単語」、つまり多数のニューロンのスパイクの安定した組み合わせで、全体として読まれるものがあるかどうかもわからない。そして、言語が脳の部位によって同じかどうかもわからない。モールス符号は一晩で覚えられる。その表がわかっているからである。ニューロンの言語の表はこれから作らなければならず、それを作るのはおそらく通信技術者だろう。

\section{問題}

\begin{exercise}[\lvlA]
\label{ex:04:1}
\emph{数値で見る容量。}$\Delta t=1\ms$とする。(a)~式~\eqref{eq:energy}のスパイクのコストのもとで、発火率毎秒$1$スパイクのニューロンは1ジュールあたり何ビットを与えるか。(b)~$r=10$で、1秒間のスパイクを数えるだけの受け手が引き出す情報は、上限~\eqref{eq:spikeentropy}の何分の1か。
\end{exercise}

\begin{exercise}[\lvlB]
\label{ex:04:2}
\emph{最適な発火率。}ニューロンは電力$P_0+rE_{\text{sp}}$を消費する。ここで$P_0$は$20$~Wの残り半分をすべてのニューロンで割ったもの、$E_{\text{sp}}=10^{-10}$~Jである。$\Delta t=1\ms$の式~\eqref{eq:spikeentropy}による1ジュールあたりのビット数$R_{\max}(r)/(P_0+rE_{\text{sp}})$が最大になる発火率$r$はいくらか。
\end{exercise}

\begin{exercise}[\lvlB]
\label{ex:04:3}
\emph{どの相関が有害か。}ニューロン$N=1000$個、分散$\sigma^2=1$、ペアの相関$c=0.1$。等しい傾き$f'=\mathbf 1$の場合と、$\mathbf 1^{\mathsf T}f'=0$となる$\pm1$の傾きの場合について、フィッシャー情報量$\mathcal I=f'^{\mathsf T}\Sigma^{-1}f'$（$\Sigma$は共分散行列）を計算せよ。両者は何倍異なるか。
\end{exercise}

\begin{exercise}[\lvlB]
\label{ex:04:4}
\emph{シミュレーション：同時検出器。}ニューロン~\eqref{eq:glutneuron}が毎秒$5$スパイクのポアソン入力を$1000$本受け取り、$w=1$である。閾値$\theta$は、自由な電圧がそれを毎秒1回横切るように選ぶ。$\tau=10$と$2\ms$のとき、ニューロンを確率$1/2$で発火させる同時入力の数$m$をシミュレーションで求めよ。
\end{exercise}

\begin{exercise}[\lvlB]
\label{ex:04:5}
\emph{設計：相対変化の検出器。}シナプス~\eqref{eq:depression}を設計せよ。毎秒$40$スパイクでの定常電流は毎秒$10$スパイクのときより$10\,\%$以上高くなってはならず、$40$への跳躍の後、電流は時定数$50\ms$以下で新しいレベルに落ち着かなければならない。最小の$\tau_{\text{rec}}$はいくらで、そのときの$U$はいくらか。
\end{exercise}

\begin{exercise}[\lvlA]
\label{ex:04:6}
\emph{最初のスパイクの時刻とバースト。}(a)~式~\eqref{eq:latency}から、最初のスパイクがちょうど時定数1個分の後に来る入力を求めよ。それは何に依存するか。(b)~放出確率$p=0.2$のとき、すべてが失われる確率を$5\,\%$未満にするには、バーストに何個のスパイクが必要か。
\end{exercise}

\begin{exercise}[\lvlC]
\label{ex:04:7}
\emph{研究：配置の中に何ビットあるか。}ニューロンを独立なセル~\eqref{eq:spikeentropy}とし、$r=10$、$\Delta t=1\ms$とする。(a)~$20$セルの「単語」のエントロピーを、スパイク数のエントロピーと残りに分解せよ。(b)~$N=30$と$10^4$個の記録から、単語の頻度によるエントロピーの推定をシミュレーションせよ。推定値はどれだけ過小になるか。
\end{exercise}

\chapter{\nt{SERO}ニューロンは何を計算するか}
\label{sec:serocomputer}

\section{最初のブロードキャストチャネル}

これまで本書のニューロンはすべてアドレスバスで話していた。\nt{GLUT}と\nt{GABA}は決まった受け手にスパイクを送り、そのようなネットワークが何を計算し、どんな言語で話すかを私たちは調べてきた。ここからは\ref{sec:buses}節の第二のバス、ブロードキャストバスに移る。その最初のチャネルが\nt{SERO}である。これまでと同じく、生きた生体で実際に起こることだけを使う。

この章では、以後すべてのブロードキャストチャネルが使う三つの部品が登場する。一定の背景入力があれば止められるまで自分で動き続けるニューロン、実は組織の体積である配線、そして個々のスパイクの代わりに使うゆっくりした濃度である。第\ref{sec:dofa}--\ref{sec:acet}章の\nt{DOPA}、\nt{NORA}、\nt{ACET}も同じ部品から組み立てられる。

エンジニアの目で見ると、\nt{SERO}ニューロンには重要な違いが四つある。

\emph{第一に、符号は受け手が選ぶ。}セロトニンには両方の符号の受容体があり、規則~\eqref{eq:dalesign}はここでは成り立たない。符号を決めるのは送り手ではなく、受け手の受容体である（命題~\ref{prop:receptor}）~\cite{BarnesSharp1999}。

\emph{第二に、一定の背景があれば\nt{SERO}ニューロンは自分で動く。}覚醒時には毎秒数回、一定のリズムでスパイクを出し、スパイクごとの後押しを必要としない。つまりペースメーカーである。ただし一定の背景入力は欠かせない。この入力がない脳スライスでは、こうした細胞の大半は沈黙している。ノルアドレナリンを$\alpha_1$受容体経由で与えると一定のリズムが戻り、その受容体を遮断すると消える~\cite{VandermaelenAghajanian1983}。生体ではこの背景は\nt{NORA}から来る。回路として見れば電源を要する発振器であり、電源がある限り、抑制で止められるまで自分で動く。

\emph{第三に、遅い。}セロトニン受容体はほぼすべて代謝型である。自分でチャネルを開くのではなく、細胞内で一連の反応を起動する。応答は遅く、主要な受容体を通る抑制電流はおよそ1秒のうちに立ち上がって減衰する~\cite{CourtneyFord2016}。これは\nt{GLUT}の速いシナプスの応答より数十倍長い。

\emph{第四に、神経伝達物質を配線ではなく体積に放出する。}\nt{SERO}ニューロンの出力が届く領域では、セロトニンは狭い間隙ではなく細胞外空間に出て、周囲に広がることが多い~\cite{Bunin1998}。受け手が感じるのは濃度であり、分子がどの送り手から来たかは知りようがない。

このような時間スケールでは個々のスパイクはもはや意味をもたないので、ネットワークは発火率$r_j$と濃度で記述する。ニューロン$j$が伝達物質を放出する体積内のセロトニン濃度$c$は、放出によって増え、伝達物質が回収されることで減る。
\begin{empheq}[box=\keybox]{equation}
\tau_c\,\frac{\dd c}{\dd t} \;=\; -c \;+\; \kappa\sum_j r_j ,
\qquad
r_i \;=\; f_i\bigl(b_i + s_i\,g\,c\bigr), \quad s_i=\pm1 .
\label{eq:seroneuron}
\end{empheq}
これはスパイク列のもう一つの読み方であり、命題~\ref{prop:reader}を補う。体積には個々のスパイクもその順序も見えず、時間$\tau_c$で平均した発火率だけが見える。第一式は膜と同じRC回路である。$\tau_c$は体積内での伝達物質の寿命で、直接には測定されておらず、ここでは1秒程度とおく。$\kappa$は1スパイクあたりの伝達物質の量である。第二式は受け手を記述する。$b_i$は受け手自身の入力で、ペースメーカーでは正である（\nt{NORA}からの一定の背景入力もここに含まれる）。$g$は受容体の感度、符号$s_i$は受け手が選び、$f_i$は非減少の伝達特性である。

\section{1ニューロンでNOT}

抑制性受容体（$s=-1$）をもつペースメーカーを考える。周囲にセロトニンがなければ、自身の入力$b$で動く。セロトニンが現れると入力は$gc$だけ減り、$gc>b$ならニューロンは沈黙する。これがNOTである。入力がないときに出力がある。使ったのはニューロン1個である（図~\ref{fig:serogates}）。命題~\ref{prop:monotone}はここには当てはまらない。あれは正の重みを前提としていた。

同じペースメーカーに二つの異なる体積からセロトニンが作用すれば、少なくとも一方にセロトニンがあるとき沈黙する。これはNORであり、NORからは任意の論理関数を組み立てられる。\nt{GLUT}ネットワークの2線式符号はここでは要らない。信号がないことは、止められるまで動き続けるニューロンが計算する。

\begin{figure}[t]
\centering
\begin{tikzpicture}[>=Stealth,
  nrn/.style={circle,draw,very thick,fill=blue!6,minimum size=0.95cm,inner sep=0pt},
  pace/.style={nrn,double,double distance=1.2pt},
  inh/.style={-{Bar[width=7pt]},thick,red!70!black}]
\node[pace] (N1) at (0,0) {};
\draw[inh] (-1.6,0) node[left,black] {$c$} -- (N1);
\draw[->,very thick,red!70!black] (N1) -- (1.3,0) node[right,black] {$\bar c$};
\node[below=0.55cm] at (0,0) {\small NOT};
\node[pace] (N2) at (5,0) {};
\draw[inh] (3.4,0.45) node[left,black] {$c_1$} -- (N2);
\draw[inh] (3.4,-0.45) node[left,black] {$c_2$} -- (N2);
\draw[->,very thick,red!70!black] (N2) -- (6.3,0) node[right,black] {$\overline{c_1\vee c_2}$};
\node[below=0.55cm] at (5,0) {\small NOR};
\end{tikzpicture}
\caption{\nt{SERO}ニューロンによる論理。二重丸はペースメーカーで、一定の背景入力があれば抑制されるまで自分で動く。横棒付きの線は抑制性受容体（$s=-1$）。$c$、$c_1$、$c_2$は受け手の周囲の体積におけるセロトニン濃度。左がNOT、右がNOR。}
\label{fig:serogates}
\end{figure}

\begin{keyblock}
\begin{proposition}[\nt{SERO}論理の完全性]
\label{prop:serocomplete}
ネットワークに抑制性受容体をもつペースメーカーがあり、異なる体積への放出が混ざらないなら、ネットワーク~\eqref{eq:seroneuron}は任意の論理関数を計算でき、各ビットは1本の配線で伝わる。
\end{proposition}
\end{keyblock}

論理の面では\nt{SERO}ネットワークは\nt{GLUT}より強い。しかし「異なる体積への放出が混ざらない」という条件は脳では満たされない（\ref{sec:volume}節）。

\section{負帰還}

セロトニンニューロンは、自分自身の上に抑制性受容体をもつ~\cite{Sprouse1987}。放出したセロトニンが自分に戻ってきて、自分にブレーキをかける。エンジニアにはおなじみの回路、負帰還増幅器である（図~\ref{fig:serofeedback}）。

何が測定済みで、何がモデルかを区別しておく。\nt{SERO}ニューロンが集まる縫線核そのものの中では、セロトニンは共通の体積ではなく、シナプスを通じて一対一で近隣を抑制する。トランスポーターがすばやく回収し、間隙の外に溜まらせないからである。1回の放出は発火に短い休止を生むだけで、平均発火率は変わらない~\cite{CourtneyFord2016}。したがって以下の、共通の体積を通じて閉じたループはモデルである。このような帰還が平均発火率のレベルで働いたら何をするかを計算する。

\begin{figure}[t]
\centering
\begin{tikzpicture}[>=Stealth,
  nrn/.style={circle,draw,very thick,fill=blue!6,minimum size=0.95cm,inner sep=0pt},
  pace/.style={nrn,double,double distance=1.2pt},
  blk/.style={draw,very thick,rounded corners=2pt,fill=orange!10,minimum height=0.9cm,minimum width=2.2cm,align=center,font=\small},
  inh/.style={-{Bar[width=7pt]},thick,red!70!black}]
\node[pace] (N) at (0,0) {};
\node[blk] (V) at (3.6,0) {体積\\$\tau_c$};
\draw[->,very thick] (N) -- node[above] {\small $r$} (V);
\draw[->,very thick] (V) -- (6.0,0) node[right] {\small $c$：すべての受け手へ};
\draw[inh] (V.south) -- ++(0,-0.9) -| node[pos=0.25,below] {\small $-g\,c$} (N.south);
\draw[->,thick,blue!60!black] (-1.7,0) node[left,black] {\small $b$} -- (N);
\end{tikzpicture}
\caption{自身の伝達物質による帰還をもつ\nt{SERO}ニューロン。濃度$c$はすべての受け手に届き、ニューロン自身も抑制する。モデルではこれはレベル調節器である。脳でループがこの役割を果たすというのは仮説である。}
\label{fig:serofeedback}
\end{figure}

このループが何をするかを計算しよう。伝達特性を線形とし、$u>0$で$f(u)=au$とする。平衡では$c=\kappa r$かつ$r=a(b-gc)$である。一方を他方に代入すると
\begin{empheq}[box=\keybox]{equation}
r \;=\; \frac{a\,b}{1+G}, \qquad G \;=\; a\,g\,\kappa .
\label{eq:feedback}
\end{empheq}
$G$はループゲインである。これはどんな負帰還増幅器にも共通の式であり、同じ二つの利点をもたらす。

\emph{ばらつきへの強さ。}ニューロンの興奮性$a$が割合$\varepsilon$だけ変わったとする。帰還がなければ発火率も同じ割合で変わる。帰還があり$G\gg1$なら、発火率$r\approx b/(g\kappa)$は$a$にほとんど依存せず、その変化は$1+G$分の1に縮む。$G=9$ならばらつきは10分の1に抑えられる。

\emph{速さ。}$r=a(b-gc)$を式~\eqref{eq:seroneuron}の第一式に代入すると$\tau_c\,\dd c/\dd t=-(1+G)\,c+\kappa ab$となる。濃度は時定数$\tau_c/(1+G)$で自分のレベルに近づく。帰還がない場合の$1+G$倍速い。

これが\nt{SERO}系が担いうる最初の計算である。サーモスタットが温度を保つように、脳全体のセロトニンのレベルを設定値に保つ。これは仮説である。実験では、自己抑制は今のところ平均発火率を変えない短い休止としてしか見えていない。

\section{スパイクからレベルへ}

第二の計算は式~\eqref{eq:seroneuron}の第一式に隠れている。ニューロンは個々のスパイクを出すが、受け手が感じるのは、それらを時間$\tau_c$で平滑化した濃度である。源の発火率が変わると、濃度は遅れてそれに追従する。
\begin{equation}
c(t) \;=\; \frac{\kappa}{\tau_c}\int_{-\infty}^{t} e^{-(t-t')/\tau_c}\,\sum_j r_j(t')\,\dd t' .
\label{eq:lowpass}
\end{equation}
これは直近の数$\tau_c$にわたる源の活動の移動平均、すなわちローパスフィルタである（図~\ref{fig:lowpass}）。最も簡単なD/A変換器も、パルスをRC回路に通してその頻度をレベルに変える。\nt{SERO}系は同じことを数十万個のニューロンで一度に行う。

この量は同時に記憶でもある。源が沈黙した後も、濃度はさらに$\tau_c$程度の時間保たれる。これはビットではなく、なめらかに消えていく痕跡である。

\begin{figure}[t]
\centering
\begin{tikzpicture}[>=Stealth,x=1.45cm,y=1cm]
\draw[gray!70!black] (0,3.2) -- (8.1,3.2);
\node[left,font=\small] at (-0.05,3.4) {スパイク};
\node[above,font=\scriptsize,gray!40!black] at (1,3.65) {毎秒$2$回};
\node[above,font=\scriptsize,gray!40!black] at (3.5,3.65) {毎秒$8$回};
\node[above,font=\scriptsize,gray!40!black] at (6.5,3.65) {毎秒$2$回};
\draw[->,gray!70!black] (0,0) -- (8.3,0) node[right,black] {\small $t$};
\draw[->,gray!70!black] (0,0) -- (0,2.75) node[left,black] {\small $c$};
\foreach \s in {0.136,0.529,0.760,1.223,1.714,1.748,1.755,2.663,2.701,2.734,3.414,3.493,3.719,3.800,3.928,3.948,4.074,4.327,4.420,4.589,4.728,4.736,4.914,5.025,5.205,5.220,6.224,6.544,7.178}{\draw[thick,blue!60!black] (\s,3.2) -- (\s,3.6);}
\foreach \a/\b/\v in {0.000/0.136/0.000,0.136/0.529/1.000,0.529/0.760/1.675,0.760/1.223/2.330,1.223/1.714/2.466,1.714/1.748/2.509,1.748/1.755/3.425,1.755/2.663/4.402,2.663/2.701/2.775,2.701/2.734/3.672,2.734/3.414/4.553,3.414/3.493/3.306,3.493/3.719/4.055,3.719/3.800/4.235,3.800/3.928/4.905,3.928/3.948/5.316,3.948/4.074/6.211,4.074/4.327/6.476,4.327/4.420/6.028,4.420/4.589/6.493,4.589/4.728/6.483,4.728/4.736/6.642,4.736/4.914/7.589,4.914/5.025/7.351,5.025/5.205/7.579,5.205/5.220/7.331,5.220/6.224/8.221,6.224/6.544/4.012,6.544/7.178/3.914,7.178/8.000/3.076}{\draw[very thick,orange!85!black] plot[domain=\a:\b,samples=10] (\x,{0.25*\v*exp(-(\x-\a))});}
\foreach \s/\p/\q in {0.136/0.000/1.000,0.529/0.675/1.675,0.760/1.330/2.330,1.223/1.466/2.466,1.714/1.509/2.509,1.748/2.425/3.425,1.755/3.402/4.402,2.663/1.775/2.775,2.701/2.672/3.672,2.734/3.553/4.553,3.414/2.306/3.306,3.493/3.055/4.055,3.719/3.235/4.235,3.800/3.905/4.905,3.928/4.316/5.316,3.948/5.211/6.211,4.074/5.476/6.476,4.327/5.028/6.028,4.420/5.493/6.493,4.589/5.483/6.483,4.728/5.642/6.642,4.736/6.589/7.589,4.914/6.351/7.351,5.025/6.579/7.579,5.205/6.331/7.331,5.220/7.221/8.221,6.224/3.012/4.012,6.544/2.914/3.914,7.178/2.076/3.076}{\draw[very thick,orange!85!black] (\s,0.25*\p) -- (\s,0.25*\q);}
\draw[thick,densely dashed,black] plot[domain=0:2,samples=30] (\x,{0.25*2*(1-exp(-\x))})
  plot[domain=2:5,samples=40] (\x,{0.25*(8-6.271*exp(-(\x-2)))})
  plot[domain=5:8,samples=40] (\x,{0.25*(2+5.688*exp(-(\x-5)))});
\foreach \x in {0,2,5,8}{\draw[gray!60!black] (\x,-0.05) -- (\x,0.05) node[below=3pt,font=\scriptsize] {\x\,s};}
\end{tikzpicture}
\caption{スパイクからレベルへ。上は\nt{SERO}ニューロンのスパイク。2秒間はまばら、3秒間は密、その後ふたたびまばらになる。下は$\tau_c=1\,\text{s}$のときのセロトニン濃度~\eqref{eq:lowpass}。破線はスパイクが完全に等間隔だった場合の同じ量。}
\label{fig:lowpass}
\end{figure}

\section{配線とは体積である}
\label{sec:volume}

命題~\ref{prop:serocomplete}の条件に戻ろう。近くで別々の送り手が放出したセロトニンは、一つの濃度に足し合わされる。

\emph{ここでの配線は線維ではなく、組織の体積である。}二つの信号が別々のままでいられるのは、伝達物質が広がりきる距離よりも離れた領域に放出される場合だけである。拡散係数$D$の分子は、時間$\tau$のうちに
\begin{empheq}[box=\keybox]{equation}
\ell \;\sim\; \sqrt{D\,\tau}\, .
\label{eq:diffusionlength}
\end{empheq}
だけ移動する。細胞間隙の屈曲は小分子の拡散をおよそ$2.6$倍遅くする~\cite{Sykova2008}。セロトニン自体の組織内拡散係数は測定されていない。分子の大きさから、数$\times10^{-6}$~cm$^2$/sと見積もる。信号の寿命が$\tau\sim1$~s（これも推定）なら、$\ell\sim\sqrt{3\cdot10^{-6}}$~cm $\approx20$~\textmu mとなる。このネットワークでのアドレスは\emph{場所}である。受け手は、自分がどこにいるかによってしか、信号が誰から来たかを知ることができない。

実際の\nt{SERO}ニューロンは、こうした個別のアドレスがほとんど残らないように作られている。それぞれの出力は脳の広大な範囲にばらまかれ、1個の細胞の枝が互いに遠く離れた多くの領域に届く~\cite{GagnonParent2014}。1細胞あたりの放出部位は推定でおよそ$10^5$（表~\ref{tab:types}）で、直接数えられたわけではない。異なる\nt{SERO}ニューロンの「配線」は重なり合い、混ざる。それでも完全に1本の配線に融合するわけではない。\nt{SERO}ニューロンはいくつかのサブシステムに分かれ、それぞれ異なる領域に出力を送り、同じ出来事に異なる応答をする~\cite{Ren2018}。しかしそうしたサブシステムの数は数個であって、数十万ではない。

\begin{keyblock}
\begin{proposition}[独立な配線の数]
\label{prop:volumewires}
容積伝達のネットワークでは、独立な信号の数を制限するのはニューロンの数ではなく、ニューロンが伝達物質を放出する、大きさ$\ell\sim\sqrt{D\tau}$の重ならない領域の数である。各ニューロンの出力が大きな体積にばらまかれていれば、ネットワーク全体は少数の共通変数しか伝えない。
\end{proposition}
\end{keyblock}

したがって、命題~\ref{prop:serocomplete}の論理回路は、脳では組み立てる材料がない。一方、調節器~\eqref{eq:feedback}とフィルタ~\eqref{eq:lowpass}には個別の配線は要らず、共通の量が一つあれば足りる。フィルタは投射先領域での体積への放出の測定から直接導かれる~\cite{Bunin1998}。レベル調節器は仮説である。

\section{計画の時間幅}
\label{sec:horizon}

共通でゆっくりした量が一つあると、何の役に立つだろうか。よい候補は、ネットワーク全体で共通であるべきで、数秒から数分より速くは変わらないパラメータである。そのようなパラメータは第\ref{sec:neurontypes}章ですでに登場した。誤差の式~\eqref{eq:rpe}の係数$\gamma$で、将来の報酬が即時の報酬よりどれだけ安いかを表す。連続時間ではその役割を\emph{時間幅}$\tau_\gamma$が担う。時間$D$待たなければならない報酬$R$は、いま$R\,e^{-D/\tau_\gamma}$の価値をもつ。

時間幅は待つ価値があるかどうかを決める。小さな報酬$r_0$をすぐ取るか、大きな報酬$R$を待つかを選べるとする。待つほうが得なのは
\begin{empheq}[box=\keybox]{equation}
D \;<\; \tau_\gamma\,\ln\frac{R}{r_0} .
\label{eq:patience}
\end{empheq}
の間である。$\tau_\gamma=2\,\text{s}$で遅延報酬が3倍なら、$2.2\,\text{s}$まで待つ価値がある。時間幅が2倍なら$4.4\,\text{s}$までである。忍耐は時間幅に比例する。

式~\eqref{eq:patience}は指数割引に基づいている。秒単位の遅延で測定された割引は、双曲線$R/(1+D/\tau_\gamma)$に近い。サルでは、選択における遅延報酬の価値も、その予告信号に対する\nt{DOPA}ニューロンの応答も、このように下がる~\cite{KobayashiSchultz2008}。双曲線では条件~\eqref{eq:patience}は$D<\tau_\gamma\,(R/r_0-1)$に置き換わる。報酬の比への依存は対数ではなく線形になり、同じ数値で待機の限界はほぼ2倍に伸びる。遅延がランダムなら同じ指数関数から双曲線が得られ（問題~\ref{ex:05:7}）、忍耐が時間スケールに比例することは変わらない。

仮説によれば、この時間幅を設定するのが\nt{SERO}である~\cite{Doya2002}。その役割に必要なものはそろっている。ネットワーク全体で一つの共通レベルがあり、それが数秒で変わる。直接の実験もある。セロトニンニューロンを人工的に活性化すると、動物は遅延報酬をあきらめるまでより長く待ち~\cite{Miyazaki2014}、報酬がまれになった場所でもより長く報酬を得ようとし続ける~\cite{Lottem2018}。セロトニン濃度がネットワーク内部でどのように$\tau_\gamma$に変わるのかは分かっていない。次章では、\nt{DOPA}が送り出す誤差に時間幅が入る。

\section{まとめ}

\begin{table}[t]
\centering
\small
\begin{tabular}{>{\raggedright}p{3.6cm}>{\raggedright}p{4.3cm}>{\raggedright\arraybackslash}p{4.3cm}}
\toprule
& \nt{GLUT} & \nt{SERO} \\
\midrule
反応時間 & ミリ秒 & 1秒未満から1秒 \\
作用の符号 & $+$のみ & $+$と$-$、受け手が選ぶ \\
入力がないとき & 沈黙 & 一定の背景入力があれば自分で動く \\
アドレス指定 & 一対一 & 体積、アドレスは場所 \\
NOT & 「オフ」検出器経由のみ & ニューロン1個 \\
記憶 & 自己維持する活動、疲労で消える & 濃度、時間$\tau_c$で消える \\
主な計算 & 同時性、遅延、論理、系列 & 平滑化、レベル調節と時間幅$\tau_\gamma$（仮説） \\
\bottomrule
\end{tabular}
\caption{\nt{GLUT}ニューロンだけのネットワークと\nt{SERO}ニューロンだけのネットワークが計算するもの。}
\label{tab:glutvssero}
\end{table}

論理素子としては（表~\ref{tab:glutvssero}）\nt{SERO}ニューロンは\nt{GLUT}より豊かである。しかし通信系としては一斉同報であり、遅く、アドレスもほとんどない。だからプロセッサになろうとはせず、第\ref{sec:neurontypes}章で述べた少数の共通量を平滑化して配る。仮説によれば、計画の時間幅もその一つである。ペースメーカー、体積、ゆっくりした濃度には、この後さらに3回、\nt{DOPA}、\nt{NORA}、\nt{ACET}で出会う。

\section{問題}

\begin{exercise}[\lvlA]
\label{ex:05:1}
\emph{配線とは体積である。}セロトニンについて$D=3\cdot10^{-6}$~cm$^2$/sとする（\ref{sec:volume}節）。\eqref{eq:diffusionlength}から、寿命$1\,\text{s}$の信号の$\ell$と、伝達物質が$5$~cm広がるのにかかる時間を求めよ。では\nt{SERO}のブロードキャスト性を支えているのは、拡散か、それとも放出部位$\sim10^5$の配線の分岐か。
\end{exercise}

\begin{exercise}[\lvlA]
\label{ex:05:2}
\emph{レベル調節器。}\eqref{eq:seroneuron}で$f(u)=au$のとき、相対感度$S=(a/r)\,\partial r/\partial a$が$G\gg1$に限らず厳密に$1/(1+G)$に等しいことを示せ。$G=9$で興奮性$a$が$20\,\%$増えた。線形化せずに求めると、$r$は何パーセント変わるか。
\end{exercise}

\begin{exercise}[\lvlB]
\label{ex:05:3}
\emph{ループ内の遅い受容体。}\nt{SERO}受容体は遅れて応答する。$r(t)=a\bigl(b-gc(t-T)\bigr)$、体積は\eqref{eq:seroneuron}に従う。$G>1$のとき、ループが安定なのは$T<T^*=\tau_c\arccos(-1/G)/\sqrt{G^2-1}$の場合に限ることを示せ。$\tau_c=1\,\text{s}$、$G=2$と$G=9$について$T^*$を求めよ。遅延が数百ミリ秒のとき、どれだけのばらつき抑制が可能か。
\end{exercise}

\begin{exercise}[\lvlB]
\label{ex:05:4}
\emph{レベルのショットノイズ。}各スパイクは\eqref{eq:lowpass}に従い、$c$に$\tau_c=1\,\text{s}$の指数関数を加える。$3\cdot10^5$個の\nt{SERO}ニューロンすべてが毎秒$2$回のポアソンスパイクを一つの体積に放出するとき、$c$の変動係数を求めよ。毎秒$4$スパイクのニューロン1個で$\mathrm{CV}=10\,\%$を得るには、$\tau_c$はいくらであればよいか。
\end{exercise}

\begin{exercise}[\lvlB]
\label{ex:05:5}
\emph{シミュレーション：帰還対ノイズ。}発火率$r_i=a_i[b-gc]_+$（$a_i$は$[2.8,\,5.2]$で一様）の$N=50$個のポアソン型ペースメーカーと、$\kappa=1/N$、$\tau_c=1\,\text{s}$の体積~\eqref{eq:seroneuron}をシミュレートせよ。帰還（$b=1$、$g=2.25$）は、$b=0.1$、$g=0$の場合に比べてレベル$c$のCVを何分の1にするか。ニューロンの発火率はそろうか。
\end{exercise}

\begin{exercise}[\lvlB]
\label{ex:05:6}
\emph{設計：\nt{SERO}論理によるXOR。}ゲートは、$b-g\sum_kc_k>0$なら$r_0$を、そうでなければ$0$を自分の体積$\tau_c\,\dd c/\dd t=-c+\kappa r$に出すペースメーカーで、NORを計算する。このゲート5個でXORを組み、$\tau_c=1\,\text{s}$、$G'=g\kappa r_0/b=2$での整定時間を求めよ。
\end{exercise}

\begin{exercise}[\lvlB]
\label{ex:05:7}
\emph{遅延が未知のときの忍耐。}(a)~動物は、待ち時間が$6\,\text{s}$以内なら3倍の報酬を待つことに応じる。これはどんな時間幅$\tau_\gamma$を意味するか。(b)~遅延$D$が平均$\bar D$の指数分布に従うとする。待つほうが得なのは$\bar D<\tau_\gamma(R/r_0-1)$のときであることを示し、$\tau_\gamma=2\,\text{s}$、$R=3r_0$で固定遅延の場合と比べよ。
\end{exercise}

\begin{exercise}[\lvlC]
\label{ex:05:8}
\emph{濃度はどのように時間幅を設定するか。}\nt{DOPA}ニューロンは$V$を、直接には重み$k_1$で、$\tau_d=0.1\,\text{s}$で平滑化する\nt{GABA}仲介ニューロン経由では重み$-k_2$で受け取る。ゆっくりした$V$では和は$(k_1-k_2)V+k_2\tau_d\dot V$となる。$\tau_\gamma=2\,\text{s}$の\eqref{eq:ctd}に合うように$k_1$、$k_2$を選べ。\nt{SERO}は$k_1$を$m$倍する。時間幅を2倍にするには$m$をどれだけ変えればよいか。
\end{exercise}

\chapter{学習するニューラルネット：\nt{DOPA}を加える}
\chaptermark{学習するネットワーク}
\label{sec:dofa}

\section{ゲームのルール}

これまでの章の回路では、重みはすべてエンジニアが設定した。生体にエンジニアはいない。重みは動作しながら自分で決まり、しかもネットワークが何か役に立つことをするように決まらなければならない。これを学習という。

これまでのルールに一つ加える。シナプスは自分の重みを自分で変え、知ることができるのは\emph{自分のすぐ近く}で起きていることだけである。送り手の活動$x$、受け手の活動$y$、そしてそこまで届く物質である。シナプスに個別の修正値を送ってくれる者はいない。

これで人工ニューラルネットワークの主要な手法、誤差逆伝播法は除外される。そこでは出力の誤差が、順方向の伝播の後の別フェーズで、同じ重みを通ってネットワークを逆向きに進み、各重みがそれぞれの修正値を受け取る。そのためには順方向の配線を正確になぞる逆向きの配線と、共通のクロックが必要である。どちらも脳では見つかっていない。少なくとも人工ネットワークにおける形では~\cite{Lillicrap2020,WhittingtonBogacz2019}。

重みの変化は、ゆっくりした連続的な過程として書く。
\begin{equation}
\frac{\dd w}{\dd t} \;=\; \eta\,\Phi(\text{シナプスが知っていること}),
\label{eq:plasticity}
\end{equation}
ここで$\eta$は学習率で、1スパイクの間に重みがほとんど変わらないほど小さい。この章全体が、関数$\Phi$がどんなものでありうるかについての話である。

\section{重みはどこに蓄えられるか}
\label{sec:weightstore}

「重みを変える」シナプスとは何だろうか。結合には二つの側がある。送り手の配線の終末（\emph{シナプス前}側）と、受容体をもつ受け手の膜の一部（\emph{シナプス後}側）で、その間に狭い間隙がある。\nt{GLUT}結合では、受け手の膜はふつう\emph{スパイン}、つまり受け手の入力枝にある小さな突起の上にある。結合の重みは三つの因子に分解すると便利である~\cite{delCastillo1954}。
\begin{empheq}[box=\keybox]{equation}
w \;=\; N_s\,p\,A_1 .
\label{eq:quantal}
\end{empheq}
ここで$N_s$は二つのニューロン間の放出部位の数、$p$はスパイクが1か所で伝達物質の1パケットを放出する確率（第\ref{sec:code}章の$p$と同じ）、$A_1$は1パケットに対する受け手の応答、つまりそれを捕らえる受容体の数である。因子$p$は送り手の側に、$A_1$は受け手の側にあり、$N_s$には両側が必要である。新しい放出部位が生え、古いものは消え、これは生涯続く。

\emph{決めるのは受け手である。}典型的な\nt{GLUT}シナプスでは、グルタミン酸受容体の一つが、二つの条件がそろったときだけ電流を通す。送り手からグルタミン酸が来ていること、そして受け手の膜がすでに興奮していることである~\cite{Nowak1984}。これは1個の分子に組み込まれたヘッブ則であり、第\ref{sec:glutcomputer}章の同時検出器がシナプスの入口そのものに置かれたものである。作動した受容体は受け手にカルシウムを流し込み、カルシウムが重みの変化を起動する。入力と出力が同時に見える場所は受け手だけなので、決定は受け手で下される。

\emph{決定を実行するのはどちらの側でもよい。}最もよく調べられている長期増強は受け手の側で起こる。膜に新しい受容体が組み込まれ~\cite{Malinow2002}、スパインが大きくなる~\cite{Matsuzaki2004}。つまり$A_1$が増える。しかし受け手は$p$を変えることもできる。間隙を逆向きに送り手へ向かう物質を放出し（付録~\ref{sec:othermediators}の逆行チャネル）、送り手は伝達物質の放出を減らす~\cite{WilsonNicoll2001}。そして第\ref{sec:code}章の短期的な重みの変化、枯渇と促通は、送り手が自分の最近の活動に応じて自ら行う$p$の変化である。

エンジニアから見ると、重みレジスタは二つのチップに分かれている。学習則を実行するのは受け手だが、結果は自分の側にも、逆行チャネルを通じて送り手の側にも書き込める。したがって本章の規則でいう「局所的」とは、結合の片側ではなく結合全体を指す。

\section{教師なしで学ぶ}

シナプスにできる最も単純なことは、送り手と受け手が一緒に活動したときに強まることである。
\begin{equation}
\frac{\dd w}{\dd t} \;=\; \eta\,x\,y .
\label{eq:hebb}
\end{equation}
これがヘッブ則である~\cite{Hebb1949}。局所的で、クロックも要らない。しかし第\ref{sec:glutcomputer}章の\nt{GLUT}ネットワークと同じ問題を抱えている。正帰還である。強い重みは$y$を大きくし、大きな$y$は重みをさらに強め、重みは際限なく増える。治療法は重みに対する負帰還である。重みを$y^2$に比例して減らしもする。入力$\mathbf x$と線形出力$y=\mathbf w\cdot\mathbf x$をもつニューロンでは
\begin{empheq}[box=\keybox]{equation}
\frac{\dd\mathbf w}{\dd t} \;=\; \eta\,y\,\bigl(\mathbf x - y\,\mathbf w\bigr) .
\label{eq:oja}
\end{empheq}
この規則も局所的である。シナプスは自分の入力$x_i$、出力$y$、重み$w_i$を知っている。

\begin{keyblock}
\begin{proposition}[ヘッブ則が見つけるもの]
\label{prop:oja}
規則~\eqref{eq:oja}は重みベクトルを、入力が最も大きく変動する方向に沿った単位長のベクトル、すなわち入力の相関行列$C=\langle\mathbf x\mathbf x^{\mathsf T}\rangle$の主固有ベクトルに導く~\cite{Oja1982}。ニューロンは、入力を圧縮できる最も表現力の高い一つの量を出すことを学ぶ。
\end{proposition}
\end{keyblock}

証明。\eqref{eq:oja}を入力について平均すると$\dd\mathbf w/\dd t=\eta\bigl(C\mathbf w-(\mathbf w^{\mathsf T}C\mathbf w)\,\mathbf w\bigr)$となる。不動点は$C$の単位長の固有ベクトルである。$\mathbf w$が固有値$\lambda_k$の固有ベクトル$\mathbf e_k$に近いとき、別の固有ベクトル$\mathbf e_j$方向の小さな摂動は$e^{\eta(\lambda_j-\lambda_k)t}$のように成長する。安定なのは、すべての$\lambda_j<\lambda_k$となる点、つまり主固有ベクトルだけである。

スパイクのタイミングを見るヘッブ則の変種もある。送り手のスパイクが受け手のスパイクより$s$ミリ秒\emph{前}に来れば重みは$A_+e^{-s/\tau_+}$だけ増え、$s$ミリ秒\emph{後}なら$A_-e^{-s/\tau_-}$だけ減る。$\tau_\pm$は20ミリ秒程度である。この規則は\nt{GLUT}ニューロンのシナプスで測定されている~\cite{BiPoo1998}。応答の\emph{原因でありえた}結合を強め、遅れて来た結合を弱める（図~\ref{fig:stdp}）。同じ興奮の系列を何度もネットワークに流せば、各段から次の段への結合が自然に育つ。第\ref{sec:glutcomputer}章の連鎖が、エンジニアなしで組み上がる。

\begin{figure}[t]
\centering
\begin{tikzpicture}[>=Stealth,x=0.07cm,y=1.3cm]
\draw[->,gray!70!black] (-66,0) -- (68,0) node[right,black] {\small $s$};
\draw[->,gray!70!black] (0,-1.2) -- (0,1.35) node[above,black] {\small 重みの変化};
\draw[very thick,blue!60!black] plot[domain=0.5:60,samples=50] (\x,{exp(-\x/20)});
\draw[very thick,red!70!black] plot[domain=-60:-0.5,samples=50] (\x,{-0.6*exp(\x/20)});
\draw[blue!60!black,thick] (0.5,0) -- (0.5,1);
\draw[red!70!black,thick] (-0.5,0) -- (-0.5,-0.6);
\foreach \x in {-40,-20,20,40}{\draw[gray!60!black] (\x,-0.04) -- (\x,0.04);}
\node[below,font=\scriptsize] at (20,-0.04) {$20\ms$};
\node[above,font=\scriptsize] at (-20,0.04) {$-20\ms$};
\node[align=left,font=\small,blue!60!black,anchor=west] at (22,0.95) {送り手が先：\\ 結合が強まる};
\node[align=right,font=\small,red!70!black,anchor=east] at (-12,-0.95) {送り手が後：\\ 結合が弱まる};
\end{tikzpicture}
\caption{タイミングによるヘッブ則。横軸は$s=t_{\text{post}}-t_{\text{pre}}$、つまり受け手のスパイクが送り手のスパイクからどれだけ遅れたか。$\tau_\pm=20\ms$、$A_-=0.6A_+$。幅数十ミリ秒の窓は、同時性の窓~\eqref{eq:window}と同じ桁である。}
\label{fig:stdp}
\end{figure}

しかしこの節の規則はどれも、\emph{頻繁な}ことを教えるのであって、\emph{役に立つ}ことを教えるのではない。$x$と$y$しか知らないシナプスには、行動がジュースをもたらしたのか痛みをもたらしたのかが分からない。結果から学ぶには、シナプスに第三の信号が必要である。

\section{第三の因子}

第三の信号を運ぶのは\nt{DOPA}である。ドーパミンニューロンは一つの数、報酬予測誤差$\delta$~\eqref{eq:rpe}をブロードキャストする（第\ref{sec:neurontypes}章）。重みの変化を、送り手の活動、受け手の活動、$\delta$の三つの因子の積に比例させよう。

難しいのは、報酬がネットワークが行動を選んだ時点ではなく、数百ミリ秒から数秒後に来ることである。$\delta$が届く頃には積$xy$はとうにゼロになっていて、シナプスには自分が関与したことの記憶が必要になる。最も簡単な記憶は、おなじみのRC回路である。
\begin{empheq}[box=\keybox]{equation}
\tau_e\,\frac{\dd e}{\dd t} \;=\; -e \;+\; x\,y ,
\qquad
\frac{\dd w}{\dd t} \;=\; \eta\,\delta(t)\,e(t) .
\label{eq:threefactor}
\end{empheq}
量$e$を\emph{適格度トレース}と呼ぶ~\cite{Fremaux2016}。同時の活動はシナプスに印を残し、それは時間$\tau_e$で消えていく。その間にドーパミンが来れば重みは$\delta$の符号で変わり、来なければ印は跡形もなく消える（図~\ref{fig:trace}）。\nt{GLUT}シナプスで測定されている。同時の活動の後1〜2秒以内に届いたドーパミンはそれを結合の強化に変えるが、それより遅く届いたものは変えない~\cite{Yagishita2014}。第三の信号がドーパミンではなく受け手自身の強い興奮である場合にも、秒単位の可塑性の窓が見つかっている~\cite{Bittner2017}。

\begin{figure}[t]
\centering
\begin{tikzpicture}[>=Stealth,x=7cm,y=1cm]
\fill[orange!12] (0.7,-0.2) rectangle (0.8,5.2);
% rows: x*y, delta, traces, weights
\foreach \y/\lab in {4.4/{$x\,y$},3.1/{$\delta$},1.4/{トレース $e$},0/{重み $w$}}{
  \draw[gray!70!black] (0,\y) -- (1.42,\y);
  \node[left,font=\small] at (-0.02,\y+0.3) {\lab};}
\draw[very thick,blue!60!black] (0,4.4) -- (0,4.9) -- (0.2,4.9) -- (0.2,4.4);
\node[right,font=\scriptsize,blue!60!black] at (0.21,4.8) {送り手と受け手が同時に活動};
\draw[very thick,orange!85!black] (0,3.1) -- (0.7,3.1) -- (0.7,3.7) -- (0.8,3.7) -- (0.8,3.1) -- (1.4,3.1);
\node[right,font=\scriptsize,orange!85!black] at (0.81,3.6) {ドーパミンが到着};
% traces, each in units of its own maximum
\draw[very thick,blue!60!black] plot[domain=0:0.2,samples=20] (\x,{1.4+1.1*(1-exp(-\x))/(1-exp(-0.2))})
  plot[domain=0.2:1.4,samples=60] (\x,{1.4+1.1*exp(-(\x-0.2))});
\draw[thick,densely dashed,gray!50!black] plot[domain=0:0.2,samples=30] (\x,{1.4+1.1*(1-exp(-\x/0.05))/(1-exp(-4))})
  plot[domain=0.2:1.4,samples=80] (\x,{1.4+1.1*exp(-(\x-0.2)/0.05)});
\node[right,font=\scriptsize,blue!60!black] at (1.0,2.15) {$\tau_e=1\,\text{s}$};
\node[right,font=\scriptsize,gray!40!black] at (0.3,1.65) {$\tau_e=50\ms$};
% weights
\draw[very thick,blue!60!black] (0,0.25) -- (0.7,0.25) -- (0.8,0.8) -- (1.4,0.8) node[right,font=\scriptsize] {$\tau_e=1\,\text{s}$：重みが増加};
\draw[thick,densely dashed,gray!50!black] (0,0.2) -- (1.4,0.2) node[right,font=\scriptsize,gray!40!black] {$\tau_e=50\ms$：変化なし};
% time axis marks
\foreach \x/\l in {0/0,0.2/0.2,0.7/0.7,1.4/1.4}{\draw[gray!60!black] (\x,-0.05) -- (\x,0.05) node[below=3pt,font=\scriptsize] {\l\,s};}
\draw[<->,thick,gray!50!black] (0.2,5.35) -- node[above,font=\scriptsize,black] {報酬の遅延 $0.5\,\text{s}$} (0.7,5.35);
\end{tikzpicture}
\caption{規則~\eqref{eq:threefactor}による適格度トレース。同時の活動は$0.2\,\text{s}$続き、その終了の0.5秒後にドーパミンが来る。トレースはそれぞれの最大値に対する割合で示す。遅いトレース（$\tau_e=1\,\text{s}$）はこの時点でまだ生きているが、速いトレース（$\tau_e=50\ms$）はとうに消えている。}
\label{fig:trace}
\end{figure}

\begin{keyblock}
\begin{proposition}[ブロードキャスト信号による学習]
\label{prop:threefactor}
規則~\eqref{eq:threefactor}は、直近の$\tau_e$の間にネットワークの動作に関与したシナプスだけを、共通の信号$\delta$が指定する方向に変える。ブロードキャスト信号と局所的なトレースが組み合わさって、宛先をもった修正になる。行動と結果の間の遅延は、$\tau_e$を超えない限り許される。
\end{proposition}
\end{keyblock}

\section{なぜうまくいくのか：探索としてのノイズ}
\label{sec:dofagradient}

規則~\eqref{eq:threefactor}はなぜ改善をもたらすのか。答えは第\ref{sec:code}章のノイズにある。ニューロンの出力を$y=f(u)+\xi$とする。ここで$u=\sum_jw_jx_j$、$\xi$は平均ゼロ、分散$\sigma^2$のランダムなゆらぎである。報酬$R$は$y$に依存する。トレースとして単なる$x_jy$ではなくそのランダムな部分$x_j\xi$をとり、第三の因子として誤差$\delta=R-\bar R$をとる。$\bar R$は期待報酬である。ノイズが小さければ$R\approx R\bigl(f(u)\bigr)+R'\,\xi$なので
\begin{empheq}[box=\keybox]{equation}
\bigl\langle \delta\,\xi\,x_j \bigr\rangle \;=\; \sigma^2\,R'\,x_j \;=\; \frac{\sigma^2}{f'(u)}\,\frac{\partial\langle R\rangle}{\partial w_j} .
\label{eq:perturbation}
\end{empheq}
平均的なステップは期待報酬の勾配に沿って進む~\cite{Williams1992}。誰も微分を計算していない。ネットワークがランダムにずれ、ドーパミンがそれでよくなったかを伝え、関与したシナプスが得をする方向に動いた（図~\ref{fig:perturbation}）。第\ref{sec:code}章でメッセージの伝達を妨げていたノイズが、ここでは探索になる。

\begin{figure}[t]
\centering
\begin{tikzpicture}[>=Stealth,x=2cm,y=3cm]
\draw[->,gray!70!black] (0,0) -- (4.2,0) node[below left,font=\small,black] {出力 $y$};
\draw[->,gray!70!black] (0,0) -- (0,1.2) node[left,font=\small,black] {報酬 $R$};
% reward hill
\draw[very thick,blue!60!black] plot[domain=0:4,samples=60] (\x,{max(0,1-((\x-3)/2.2)^2)});
\node[font=\scriptsize,blue!60!black,anchor=south] at (3,1.02) {最良の出力};
% noise around f(u)
\draw[thick,gray!60!black] plot[domain=0.6:2.4,samples=50] (\x,{0.28*exp(-((\x-1.5)/0.3)^2/2)});
\draw[densely dashed,gray!60!black] (1.5,0) -- (1.5,0.535);
\node[below,font=\small] at (1.5,0) {$f(u)$};
\node[font=\scriptsize,gray!50!black] at (1.5,-0.2) {ゆらぎ $\xi$};
% expected reward
\draw[densely dotted,gray!60!black] (0.5,0.535) node[left,font=\scriptsize,black] {$\bar R$} -- (2.1,0.535);
% two random deviations
\fill[green!45!black] (2.1,0.833) circle (0.035);
\draw[very thick,green!45!black] (2.1,0.535) -- (2.1,0.833);
\node[font=\scriptsize,green!45!black,anchor=west,align=left] at (2.18,0.6) {$\xi>0$：よい、\\ $\delta>0$};
\fill[red!70!black] (0.9,0.089) circle (0.035);
\draw[very thick,red!70!black] (0.9,0.535) -- (0.9,0.089);
\node[font=\scriptsize,red!70!black,anchor=east,align=right] at (0.83,0.27) {$\xi<0$：悪い、\\ $\delta<0$};
% resulting step
\draw[->,very thick,orange!85!black] (1.55,0.4) -- (1.95,0.4) node[right,font=\scriptsize,black] {ステップ};
\end{tikzpicture}
\caption{探索としてのノイズ~\eqref{eq:perturbation}。ニューロンの出力は$f(u)$のまわりでゆらぐ。右へのランダムなずれ（緑）は期待より多くを得て$\delta>0$、左へのずれ（赤）は少なく$\delta<0$。どちらの場合も積$\delta\,\xi$は正で、重みは報酬が増える右へ動く。平均するとステップは傾き$R'$に比例する。山の頂上では両側へのずれが同じように悪く、ステップは打ち消し合う。}
\label{fig:perturbation}
\end{figure}

\begin{keyblock}
\begin{proposition}[逆向きの配線なしの勾配法]
\label{prop:gradient}
ネットワークに活動のランダムなずれがあり、ブロードキャスト信号が報酬予測誤差に等しく、各状況で平均ゼロであるなら、規則~\eqref{eq:threefactor}は平均として重みを期待報酬の勾配に沿って変える。そのために逆向きの配線も独立した学習フェーズも要らない。
\end{proposition}
\end{keyblock}

これで、ドーパミンが報酬そのものではなく報酬予測の\emph{誤差}を伝える理由が分かる。第三の因子が報酬そのもの$R\ge0$だったとしても\eqref{eq:perturbation}の平均は変わらないが、二つの問題が生じる。第一に、有用な部分に項$\bar R\,\xi x_j$が加わる。平均はゼロだが、ノイズは大きい。第二に、こちらのほうが深刻だが、本物のシナプスはトレース$x_jy$のランダムな部分とそうでない部分を分けられない。入力が同じなら$x_jf(u)$は試行ごとに同じ数である。その状況で$\delta$の平均がゼロなら、トレースのこの部分は何も変えない。しかし$\delta\ge0$なら、この部分はネットワークがしたことすべてを、正しいことも誤ったことも、繰り返し強化してしまう。したがって$\bar R$は、生涯の平均ではなく\emph{その状況での}期待値でなければならない。それを計算するのは、後で述べるクリティックの仕事である。

これをモデルで確かめた。\nt{GLUT}ニューロンの二つの群が、図~\ref{fig:wta}のように相互抑制によって二つの行動の一方を選ぶ。「開始」信号から各群への重みは最初は等しく、どちらが勝つかはノイズが決める。行動$A$は確率$0.8$で、行動$B$は確率$0.2$でジュースをもたらし、ジュースは選択の0.5秒後に来る。$\tau_e=1\,\text{s}$、$\delta=R-\bar R$では、500試行で$A$を選ぶ割合は最初の100試行の$0.65$--$0.8$から最後の100試行の$0.96$--$1.0$に上がり、重みは$0.85$--$1.35$の範囲にとどまった。報酬の遅延より短い$\tau_e=50\ms$では、ネットワークは何も学ばなかった。$A$を選ぶ割合は半分前後のままだった。期待値を引かない$\delta=R$でもネットワークは$A$を選ぶようになったが、勝者の重みは同じ500試行で1000倍になり、なお増え続けた。そして同じ$\delta=R$で学習率を10倍にすると、3回の実行のうち1回で、ネットワークは最初のランダムな選択、しかも悪いほうを永久に固定してしまった。

\section{配線1本の代償}

信号$\delta$は全員に共通だが、それが評価するノイズは、結果に影響するすべての$N$個のニューロンのゆらぎの和である。各シナプスは$\delta$の中に、自分の有用な部分と、$N-1$個の隣人のノイズを見る。自分の分を取り出すには多くの試行で平均しなければならず、学習時間は$N$に比例して伸びる~\cite{Werfel2005}。誤差逆伝播法はこのような代償を払わない。

\begin{keyblock}
\begin{proposition}[ブロードキャストの代償]
\label{prop:broadcastcost}
共通の信号一つによる学習の時間は、結果にノイズが影響するニューロンの数に比例して伸びる。このような学習は、少数の選択肢から選ぶ小さな回路には向くが、数十億の重みの調整には向かない。
\end{proposition}
\end{keyblock}

脳はどうやら実際にこのように仕事を分けている。\nt{DOPA}ニューロンの出力は主に行動を選ぶ領域に向かい（第\ref{sec:neurontypes}章）、重みの圧倒的多数がある皮質の巨大な\nt{GLUT}ネットワークは、主として局所的な規則で、外部の教師なしに学習する。かつてはそれをヘッブ則~\eqref{eq:oja}に帰着させていた。現在では、皮質には自分の目標、すなわち次の入力を\emph{予測する}ことがあるというデータが増えている。その場合、誤差は予測と実際の入力の差としてその場で計算される~\cite{Keller2018}。教師はネットワーク自体に組み込まれている。第三の信号が受け手自身の強い興奮である秒単位の可塑性の窓~\cite{Bittner2017}は、シナプスに局所的な誤差信号が実際にありうることを示している。それが皮質の一般的な規則であるかどうかは仮説である。

\section{誤差はどこから来るか：連続時間のクリティック}

残る問題は、\nt{DOPA}ニューロン自身がどのように$\delta$を計算するかである。強化学習のプログラムでは、予測誤差はステップ$t,t+1,\dots$で計算される。生体にステップはないので、連続時間で書く~\cite{Doya2000}。$r(t)$を報酬の流れ、$V(t)$を将来の全報酬の期待値とする。遠い報酬ほど価値は低い。
\begin{equation}
V(t) \;=\; \int_t^{\infty} e^{-(s-t)/\tau_\gamma}\,r(s)\,\dd s .
\end{equation}
微分すると$\dd V/\dd t=V/\tau_\gamma-r$となる。この等式の残差が予測誤差である。
\begin{empheq}[box=\keybox]{equation}
\delta(t) \;=\; r(t) \;+\; \frac{\dd V}{\dd t} \;-\; \frac{V}{\tau_\gamma} .
\label{eq:ctd}
\end{empheq}
定数$\tau_\gamma$は計画の時間幅で、係数$\gamma$の連続版である。仮説によればこれを設定するのは\nt{SERO}であり（\ref{sec:horizon}節）、そのとき\eqref{eq:patience}はどれだけ待つ価値があるかの規則になる。三つの場合を確かめよう（図~\ref{fig:ctdcases}）。ジュースを約束するランプが点いた。$V$が跳ね上がり、$\dd V/\dd t$が大きく、バーストが出る。約束のジュースが来た。$r>0$だが、同じ瞬間に期待は尽きて$\dd V/\dd t\approx-r$となり、バーストは出ない。ジュースが来なかった。期待が報酬なしに消え、$\dd V/\dd t<0$で、落ち込みが出る。

\begin{figure}[t]
\centering
\begin{tikzpicture}[>=Stealth,x=1.15cm,y=1cm]
\foreach \col/\ttl/\lampon/\juice in {0/{予期しないジュース}/0/1, 1/{約束のジュース}/1/1, 2/{ジュースが来ない}/1/0}{
 \begin{scope}[xshift=\col*4.6cm]
  \node[font=\small] at (1.5,3.55) {\ttl};
  \foreach \yy in {2.4,1.2,0}{\draw[gray!60!black] (0,\yy) -- (3.1,\yy);}
  % reward r
  \ifnum\juice=1
    \fill[green!45!black] (1.95,2.4) rectangle (2.05,2.85);
    \pic[scale=0.55] at (2.0,3.1) {juice};
  \else
    \pic[scale=0.55] at (2.0,3.1) {nojuice};
  \fi
  % expectation V and error delta
  \ifnum\lampon=1
    \pic[scale=0.55] at (1.0,3.1) {lamp};
    \draw[very thick,blue!60!black] (0,1.2) -- (1,1.2) -- (1,1.45)
      plot[domain=1:2,samples=20] (\x,{1.2+0.25*exp((\x-1)/1.5)}) -- (2,{1.2+0.25*exp(1/1.5)}) -- (2,1.2) -- (3.1,1.2);
    \draw[very thick,red!70!black] (0,0) -- (0.95,0) -- (0.95,0.5) -- (1.05,0.5) -- (1.05,0) -- (3.1,0);
    \ifnum\juice=0
      \draw[very thick,red!70!black] (1.95,0) -- (1.95,-0.45) -- (2.05,-0.45) -- (2.05,0);
      \node[font=\scriptsize,red!70!black,anchor=west] at (2.1,-0.35) {落ち込み};
    \else
      \node[font=\scriptsize,gray!50!black,anchor=north,align=center] at (2.0,-0.08) {$r$と$V$の低下が\\ 打ち消し合う};
    \fi
  \else
    \draw[very thick,blue!60!black] (0,1.2) -- (3.1,1.2);
    \draw[very thick,red!70!black] (0,0) -- (1.95,0) -- (1.95,0.5) -- (2.05,0.5) -- (2.05,0) -- (3.1,0);
  \fi
  \ifnum\col=0
    \node[font=\small,anchor=east] at (-0.1,2.55) {$r$};
    \node[font=\small,anchor=east] at (-0.1,1.35) {$V$};
    \node[font=\small,anchor=east] at (-0.1,0.1) {$\delta$};
  \fi
 \end{scope}}
\end{tikzpicture}
\caption{誤差~\eqref{eq:ctd}の三つの場合。上から報酬$r$、期待$V$、誤差$\delta$。左では期待がなく、ジュースはまるごと予期しないものである。中央ではランプが$V$を跳ね上げる。これは$\dd V/\dd t$の跳躍であり、$\delta$のバーストはランプの時点に来る。その後$V$は時間幅の要求どおりに増え、$\delta=0$となる。ジュースの瞬間、報酬$r$と$V$の低下が打ち消し合う。右では$V$が報酬なしに下がり、落ち込みになる。これは図~\ref{fig:rpe}と同じバースト、移行、落ち込みだが、今度はそれらがどこから来るかが見える。}
\label{fig:ctdcases}
\end{figure}

手持ちの回路で\eqref{eq:ctd}を作るには何が要るか。三つある。

\emph{推定値$V$そのもの。}これを出すのは\nt{GLUT}ニューロンの群、つまりクリティックである。その重みは同じ$\delta$を使う同じ規則~\eqref{eq:threefactor}で学習する。$\delta>0$なら期待が低すぎたのであり、直前に活動していた結合が強まる。

\emph{微分$\dd V/\dd t$。}これはレベルではなく変化に応答する回路なら何でも計算できる。第\ref{sec:gaba}章の方向検出器のような、直接の興奮と\nt{GABA}仲介ニューロンを通る遅れた抑制の組み合わせでもよいし、第\ref{sec:code}章の枯渇するシナプス~\eqref{eq:depression}でもよい。

\emph{1本の配線で符号を運ぶ。}誤差は正にも負にもなるが、発火率は正にしかならない。解決策は第\ref{sec:serocomputer}章の\nt{SERO}ニューロンと同じで、\nt{DOPA}ニューロンはペースメーカーである。毎秒数回の一定のスパイクは$\delta=0$、バーストは$\delta>0$、休止は$\delta<0$を意味する。負の誤差の余地は小さい。発火率はゼロより下げられないからである。実際の\nt{DOPA}ニューロンでも、落ち込みはバーストより浅い~\cite{Bayer2005}。

ドーパミンが$\delta$のようにふるまうことは測定されている。脳が$\dd V/\dd t$を正確にどう計算するかは仮説である。

\section{アクターとクリティック}

すべてを組み合わせよう（図~\ref{fig:actorcritic}）。状況ニューロンは、\nt{GABA}を通じて勝者を選ぶ行動群（命題~\ref{prop:wta}）、つまり\emph{アクター}を興奮させ、同時に$V$を出すクリティックも興奮させる~\cite{SuttonBarto2018}。\nt{DOPA}ニューロンは報酬$r$と$V$を受け取り、\eqref{eq:ctd}を計算し、$\delta$をアクターとクリティックのすべてのシナプスにブロードキャストする。

\begin{figure}[t]
\centering
\begin{tikzpicture}[>=Stealth,
  nrn/.style={circle,draw,very thick,fill=blue!6,minimum size=0.9cm,inner sep=0pt,font=\small},
  gab/.style={circle,draw,very thick,fill=red!10,minimum size=0.6cm,inner sep=0pt},
  dofa/.style={circle,draw,very thick,double,double distance=1.2pt,fill=orange!15,minimum size=1.35cm,inner sep=0pt,font=\scriptsize},
  inh/.style={-{Bar[width=7pt]},thick,red!70!black},
  bc/.style={->,thick,densely dashed,orange!85!black}]
\node[nrn] (S1) at (0,2.4) {$x_1$};
\node[nrn] (S2) at (0,1.2) {$x_2$};
\node[nrn] (S3) at (0,0) {$x_3$};
\node[left=0.15cm,align=right,font=\small] at (-0.5,1.2) {状況};
\node[nrn] (A) at (4,2.6) {$A$};
\node[nrn] (B) at (4,1.0) {$B$};
\node[gab] (G) at (5.2,1.8) {};
\draw[->,thick] (A) -- (G); \draw[->,thick] (B) -- (G);
\draw[inh] (G) to[bend left=35] (A);
\draw[inh] (G) to[bend right=35] (B);
\node[above,font=\small] at (4.6,3.05) {アクター：行動の選択};
\node[nrn] (V) at (4,-1.1) {$V$};
\node[below,font=\small] at (4,-1.6) {クリティック};
\foreach \s in {S1,S2,S3}{\foreach \t in {A,B,V}{\draw[->,gray!60!black] (\s) -- (\t);}}
\node[dofa] (D) at (8.0,-1.1) {\nt{DOPA}};
\draw[->,thick] (V) -- node[above,font=\scriptsize] {$V,\ \dd V/\dd t$} (D);
\draw[->,thick,green!45!black] (10.0,-1.1) node[right,black,font=\small] {報酬 $r$} -- (D);
\pic[scale=1.1] at (10.6,-0.5) {juice};
\draw[bc] (D.north) -- (8.0,4.0) -- node[above,font=\small,black] {$\delta$：アクターとクリティックの全シナプスへ} (2.2,4.0) -- (2.2,3.0);
\draw[bc] (D.south) -- (8.0,-2.4) -- (2.2,-2.4) -- (2.2,-0.6);
\end{tikzpicture}
\caption{行動の選択を学ぶネットワーク。灰色の矢印は重みが変わる\nt{GLUT}シナプス。赤い丸は\nt{GABA}ニューロン。二重丸は\eqref{eq:ctd}に従って$\delta$を計算する\nt{DOPA}ペースメーカー。破線の矢印は$\delta$のブロードキャスト。}
\label{fig:actorcritic}
\end{figure}

伝達物質の作用の符号は受け手が選び（命題~\ref{prop:receptor}）、行動選択の領域では、ニューロンの一部はあるファミリーのドーパミン受容体をもち、一部は別のファミリーをもつ。スライスでの実験では、ドーパミンと同時の活動がそろうと、前者ではシナプスが強まり、後者では弱まる~\cite{Shen2008}。モデルでいえば、後者では\eqref{eq:threefactor}の$\delta$の符号が反転している。前者は\emph{すべきこと}を、後者は\emph{すべきでないこと}を学ぶ。

\section{まとめ}

\begin{table}[t]
\centering
\footnotesize
\setlength{\tabcolsep}{4pt}
\renewcommand{\arraystretch}{1.15}
\begin{tabular}{>{\raggedright}p{2.6cm}>{\raggedright}p{2.7cm}>{\raggedright}p{3.1cm}>{\raggedright\arraybackslash}p{5.0cm}}
\toprule
規則 & シナプスが知ること & ネットワークが学ぶこと & 分かっていること \\
\midrule
発火率によるヘッブ則~\eqref{eq:oja} & $x$、$y$、自分の重み & 入力の圧縮：主方向 & 同時活動による強化は測定済み。重みを制限する機構は研究中 \\
タイミングによるヘッブ則 & $\sim20\ms$以内の$x$と$y$のスパイクの順序 & 因果的な結合、系列 & \nt{GLUT}シナプスで測定済み \\
三要素則~\eqref{eq:threefactor} & $x$、$y$、トレース$e$、共通信号$\delta$ & 報酬をもたらす行動 & 活動後数秒以内のドーパミンの影響は測定済み。選択学習の主要な機構であることは十分な根拠のある仮説 \\
クリティック~\eqref{eq:ctd} & 同上 & 報酬の期待$V$ & ドーパミンと$\delta$の類似は測定済み。$\dd V/\dd t$を計算する回路は仮説 \\
誤差逆伝播法 & 重みごとの個別の修正値 & 何でも学べるが、逆向きの配線とクロックが必要 & 脳ではこの形では見つかっていない。近似が探されている \\
\bottomrule
\end{tabular}
\caption{\nt{GLUT}、\nt{GABA}、\nt{DOPA}ニューロンからなるネットワークの学習則。}
\label{tab:learning}
\end{table}

学習則を表~\ref{tab:learning}にまとめる。\nt{GLUT}はデータを運び、\nt{GABA}は流れを制御し、\nt{DOPA}はネットワークのどの変化を残すかを決める。

\section{問題}

\begin{exercise}[\lvlA]
\label{ex:06:1}
\emph{トレースの寿命。}同時の活動$xy=1$が$e=0$から始まって$T_a=0.2\,\text{s}$続き、その終了の$d=0.5\,\text{s}$後にドーパミンが来る。(a)~この時点のトレース~\eqref{eq:threefactor}は、$\tau_e=1\,\text{s}$のとき$\tau_e=50\ms$のときの何倍か。(b)~トレースが最大になるのはどの$\tau_e$か。
\end{exercise}

\begin{exercise}[\lvlA]
\label{ex:06:2}
\emph{ヘッブ則のドリフト。}送り手と受け手は毎秒$10$スパイクの独立なポアソン列を出し、スパイクの各ペアは$\tau_\pm=20\ms$、$A_-=0.6A_+$の図~\ref{fig:stdp}の窓に従って重みを変える。まったくランダムな活動1時間で、重みは$A_+$の何倍増えるか。ドリフトが消えるのはどの$\tau_-$か。
\end{exercise}

\begin{exercise}[\lvlB]
\label{ex:06:3}
\emph{共分散ではなく相関。}命題~\ref{prop:oja}の規則は$C=\langle\mathbf x\mathbf x^{\mathsf T}\rangle$の主固有ベクトルを見つける。発火率は正である。$\mathbf x=\mathbf m+\mathbf n$、$\mathbf m=(\mu,0)$、ノイズの共分散は$\operatorname{diag}(s_1^2,s_2^2)$とする。どんな条件でニューロンは$\mathbf e_1$を学ぶか。$\mu=1$、$s_1=0.1$、$s_2=0.5$では何を学ぶか。
\end{exercise}

\begin{exercise}[\lvlA]
\label{ex:06:4}
\emph{バーストの中の時間幅。}ランプが予測できない時刻に点き、大きさ$R$のジュースが$L=1\,\text{s}$後に来る。\eqref{eq:ctd}で学習済みの期待について、ランプとジュースの間で$\delta\equiv0$であることを示し、$\tau_\gamma=2\,\text{s}$と$\tau_\gamma=4\,\text{s}$のときのランプへのバーストの面積を求めよ。
\end{exercise}

\begin{exercise}[\lvlB]
\label{ex:06:5}
\emph{ブロードキャストの代償はどこから来るか。}$N$個のニューロンが独立にゆらぎ$\xi_i\sim\mathcal N(0,\sigma^2)$、誤差は$\delta=a\sum_i\xi_i$、シナプス$j$は勾配を$\delta\,\xi_jx_j$で推定する。1試行のSN比を求めよ。ニューロン1個が1試行$5\,\text{s}$で$100$試行で学習するとき、$N=10^4$と$N=10^{10}$では学習にどれだけかかるか。
\end{exercise}

\begin{exercise}[\lvlB]
\label{ex:06:6}
\emph{設計：$\delta$を計算する回路。}\nt{DOPA}ニューロンは$r$と、重み$k_1$の$V$と、$\tau_d$で平滑化した$V$のコピーを重み$-k_2$で受け取る。(a)~ゆっくりした$V$で出力が\eqref{eq:ctd}になるように$k_1$、$k_2$を選べ。(b)~$V=V_0e^{t/\tau_\gamma}$では真の$\delta=0$である。$\tau_\gamma=2\,\text{s}$のとき、回路の偽の誤差が$V/\tau_\gamma$の$5\,\%$以下になるのはどの$\tau_d$か。
\end{exercise}

\begin{exercise}[\lvlB]
\label{ex:06:7}
\emph{シミュレーション：報酬遅延の限界。}\texttt{models/\allowbreak dofa\_learning.py}のコピーの\texttt{run()}に報酬遅延$d$を加えよ。$d=0.5\,\text{s}$で$\tau_e=0.05$、$0.2$、$1$、$4\,\text{s}$のとき、および$\tau_e=1\,\text{s}$、$d=3\,\text{s}$のとき、$500$試行の最後の100試行で$A$を選ぶ割合を求めよ。報酬まで生き残るトレースの割合（問題~\ref{ex:06:1}）がどれだけになると学習は消えるか。
\end{exercise}

\begin{exercise}[\lvlC]
\label{ex:06:8}
\emph{ブロードキャストの代償は回避できるか。}$N$個のニューロンで$y_i=w_i+\xi_i$、$\xi_i\sim\mathcal N(0,\sigma^2)$、報酬$R=-\sum_i(y_i-w_i^*)^2$、規則$\Delta w_i=\eta\,(R-\langle R\rangle)\,\xi_i$とする。最良の$\eta$で誤差$\sum_i(w_i-w_i^*)^2$が$1/e$になるのに何試行かかるか。$N$個のうちランダムな$m$個だけがゆらぐ場合はどうか。スパースな探索は役に立つか。
\end{exercise}

\chapter{ドーパミンに基づく学習の現代的な代替理論}
\chaptermark{\nt{DOPA}の他の理論}
\label{sec:dofaalt}
\begin{chaptergoals}
\item 誤差への応答が非対称な\nt{DOPA}ニューロンが報酬の分布全体を記録する仕組み；
\item 一部の\nt{DOPA}ニューロンが誤差の符号ではなく重要性や危険を伝えること；
\item 1本の配線が学習させる速い誤差と、行動の速さを決めるゆっくりしたレベルを運ぶ仕組み；
\item なぜ\nt{DOPA}の信号は一つの数ではなく束なのか、誤差そのものに挑む理論はどれか。
\end{chaptergoals}

\section{配線を実際に流れているもの}

第\ref{sec:dofa}章では、\nt{DOPA}系は1本の配線であり、そこを一つの数、報酬予測誤差$\delta$~\eqref{eq:ctd}が流れていた。この描像は多くを説明するが、ドーパミンの測定が精密になるほど、そこに収まらないものが見つかる。不快なものに報酬と同じように応答するニューロンがある。何も予期しないことが起きていないのに、動物が目標に向かって走る間ドーパミン濃度が上がる。\nt{DOPA}ニューロンごとにふるまいが違う。

エンジニアにとって、これらの発見はすべて一つの問いに帰着する。\emph{ブロードキャストバス上の信号の形式は何か。}一つの数か、ベクトルか。配線上の信号は一つか、周波数で分けられた複数か。$\delta$のように未来について語るのか、過去について語るのか。以下に現代の六つの答えを示す。それぞれについて、測定されたことと推測されていることを分ける。

\section{一つの数ではなく分布}

誤差~\eqref{eq:ctd}はクリティックに報酬の\emph{平均の}期待を教える。しかし確実な半滴のジュースと、確率2分の1の1滴とは平均が同じである。両者を区別するには報酬の分布全体が必要である。

最も簡単な課題を考える。予告信号の後にランダムな大きさ$r$の報酬が来る。\nt{DOPA}ニューロンが多数あり、$i$番目のニューロンはそれぞれ自分の推定値$V_i$を担当していて、報酬の瞬間の誤差は$\delta_i=r-V_i$であるとする。さらに、各ニューロンは正の誤差を重み$\alpha_i^+$で、負の誤差を重み$\alpha_i^-$で扱うとする。報酬のたびに推定値は次だけ動く。
\begin{equation}
\Delta V_i \;=\; \eta\,\bigl(\alpha_i^+\,[\delta_i]_+ \;-\; \alpha_i^-\,[-\delta_i]_+\bigr).
\label{eq:asymmetric}
\end{equation}
平均して正の修正と負の修正がつり合うと、推定値は変わらなくなる。
\begin{empheq}[box=\keybox]{equation}
\alpha_i^+\,\bigl\langle[r-V_i]_+\bigr\rangle \;=\; \alpha_i^-\,\bigl\langle[V_i-r]_+\bigr\rangle ,
\qquad
\kappa_i \;=\; \frac{\alpha_i^+}{\alpha_i^++\alpha_i^-} .
\label{eq:expectile}
\end{empheq}
$\kappa_i=1/2$ならこれは通常の平均である。$\kappa_i>1/2$ならニューロンは「楽観主義者」で、推定値は分布の上端寄りにずれる。$\kappa_i<1/2$なら「悲観主義者」である。

例：ジュースが確率$1/2$で来て、1滴の大きさは$1$とする。すると\eqref{eq:expectile}から$\kappa_i\cdot\tfrac12(1-V_i)=(1-\kappa_i)\cdot\tfrac12V_i$、つまり$V_i=\kappa_i$となる（図~\ref{fig:expectiles}）。$\kappa_i$の異なるニューロンの集まりは、統計学で分位点の集まりが分布を記録するのと同じように、報酬の分布を記録する。

\begin{figure}[t]
\centering
\begin{tikzpicture}[>=Stealth,x=7cm,y=3cm]
\draw[->,gray!70!black] (-0.08,0) -- (1.15,0) node[right,black] {\small 報酬};
\draw[->,gray!70!black] (-0.08,0) -- (-0.08,0.75) node[left,black] {\small 確率};
\fill[blue!25] (-0.03,0) rectangle (0.03,0.5);
\fill[blue!25] (0.97,0) rectangle (1.03,0.5);
\node[below,font=\small] at (0,0) {$0$};
\node[below,font=\small] at (1,0) {$1$};
\node[left,font=\scriptsize] at (-0.08,0.5) {$1/2$};
\foreach \k/\c/\lab/\h in {0.2/blue!60!black/{悲観主義者、$\kappa=0.2$}/0.62, 0.5/gray!50!black/{平均、$\kappa=0.5$}/0.42, 0.8/red!70!black/{楽観主義者、$\kappa=0.8$}/0.22}{
  \draw[very thick,\c] (\k,0) -- (\k,\h);
  \node[above,font=\scriptsize,\c] at (\k,\h) {\lab};}
\end{tikzpicture}
\caption{\nt{DOPA}ニューロンの集まりに記録された報酬の分布。報酬は確率$1/2$で$0$または$1$（青い棒）。割合$\kappa$のニューロンは推定値$V=\kappa$に至る~\eqref{eq:expectile}。}
\label{fig:expectiles}
\end{figure}

測定されていること。\nt{DOPA}ニューロンごとに「反転点」、つまり応答がバーストから落ち込みに切り替わる報酬の大きさが異なり、正と負の誤差への応答の非対称性が大きいニューロンほど、\eqref{eq:expectile}が要求するとおり大きな報酬で反転する~\cite{Dabney2020}。ニューロン群の応答から分布の形を復元できる。これがばらつきの副作用ではなく主な機能であることは仮説である。

\begin{keyblock}
\begin{proposition}[非対称性から分布へ]
\label{prop:expectile}
\nt{DOPA}ニューロンが正の誤差を扱う割合$\kappa_i$で異なれば、その推定値は報酬分布の異なる点~\eqref{eq:expectile}に収束する。ニューロン群は平均ではなく分布を伝え、したがってリスクも伝える。
\end{proposition}
\end{keyblock}

\section{誤差だけでなく重要性も}

誤差の式~\eqref{eq:ctd}によれば、不快な出来事（打撃、苦い味）は落ち込みを起こすはずである。期待より悪い結果だからである。\nt{DOPA}ニューロンの一部はそのとおりにふるまうが、別の一部は不快なものに報酬と同じく\emph{バースト}で応答する~\cite{Matsumoto2009}。これらのニューロンが伝えるのは誤差の符号ではなく、何か\emph{重要な}こと、つまり予期しない、強い、注意を要することが起きたということである~\cite{BrombergMartin2010}。

エンジニアにとっては、これを二つの信号と考えると分かりやすい。第一は符号付きの誤差$\delta$で、重みをどちらへ変えるかを示す。第二はその絶対値$|\delta|$、あるいは出来事の新奇性で、重みを\emph{どれだけ強く}変えるか、つまり学習率を示す。予期しない出来事の後は速く学ぶべきである。意味の上では、第\ref{sec:neurontypes}章のノルアドレナリンによるゲインに近い。

第三の可能性もある。別の軸のための別の配線である。行動選択の領域の一つに向かう\nt{DOPA}ニューロンは、報酬ではなく刺激の新奇性と強さに応答し、動物が危険を避けることを学ぶのに必要である~\cite{Menegas2018}。この配線を流れるのは「よいか悪いか」ではなく「危険か否か」である。

測定されていること。\nt{DOPA}ニューロンの群ごとに不快なものと新しいものへの応答が異なり、出力ごとに異なることを教える。脳が重要性の信号を正確にどう使うのか、学習率としてか、注意の信号としてか、危険の軸に沿った別の誤差としてかは議論されている。

\section{教えるだけでなく急き立てる}

ドーパミンには即時の作用もある。多いほど、動物は進んで、速く行動する。これを学習とどう両立させるか。

エンジニアの答えは\emph{周波数による分離}である。数百ミリ秒続く速いバーストと落ち込みは$\delta$を運び、学習させる。数分かけて変わるゆっくりしたレベルは別の量、単位時間あたりの平均報酬$\bar R$を運ぶ~\cite{Niv2007}。時間$\tau_a$で行う行動が労力$C/\tau_a$を要するとする。速いほど高くつく。一方、遅れた1秒ごとに$\bar R$のコストがかかる。その1秒で得られたはずの報酬である。総コスト$C/\tau_a+\bar R\,\tau_a$は次で最小になる。
\begin{empheq}[box=\keybox]{equation}
\tau_a^{*} \;=\; \sqrt{\frac{C}{\bar R}} .
\label{eq:vigor}
\end{empheq}
環境が豊かなほど時間は高価になり、速く行動するほうが得になる。$\bar R$が2倍になれば行動時間は$\sqrt2\approx1.4$分の1になる。$\bar R$は第\ref{sec:dofa}章で差し引いた期待報酬そのものであることに注意しよう。

投射先でのドーパミン放出は、発火率が変わらなくても変わりうる。出力そのものにある局所的な信号が放出を調整するからである~\cite{Mohebi2019}。しかしゆっくりした成分はスパイク自体にもある。次節の実験では、報酬に近づくにつれてのなめらかな増加が\nt{DOPA}ニューロンの発火率にも見える~\cite{Kim2020}。つまりゆっくりしたレベルは、発火率と放出の局所的な調整という二つの源からなり、どちらが主かはおそらく出力と課題による。

\begin{keyblock}
\begin{proposition}[1本の配線に二つの信号]
\label{prop:multiplex}
\nt{DOPA}系は、時間スケールで分けられた少なくとも二つの量を伝える。学習させる速い誤差$\delta$と、時間のコスト~\eqref{eq:vigor}を、したがって行動の速さを決めるゆっくりしたレベルである。速い成分と遅い成分が異なるふるまいをすることは測定されている。遅い成分がまさに$\bar R$に等しいことは仮説である。
\end{proposition}
\end{keyblock}

\section{接近に伴う増加}

動物が通路を遠くの報酬へ走るとき、行動選択の領域のドーパミン濃度は、目標が近づく数秒の間なめらかに上昇する~\cite{Hamid2016}。第\ref{sec:dofa}章によれば、こうなるはずはない。すべては計画どおりに進み、$\delta$はゼロのはずである。

第一の説明。この上昇は$\delta$ではなく期待$V$そのもの、つまり前節の「目標は近い、がんばる価値がある」という信号である~\cite{Hamid2016}。第二の説明は$\delta$を保持する~\cite{Kim2020}。期待が正しければ、時間幅$\tau_\gamma$のもとで報酬$R$への道のりにおいて期待は$V=Re^{-(T-t)/\tau_\gamma}$のように増え、誤差の式~\eqref{eq:ctd}により$\dd V/\dd t-V/\tau_\gamma=0$となる。しかし遠くでは期待が低すぎ、近づくにつれて推定が精密になるとしよう。すると期待はもっと急に、たとえば$\tau_v<\tau_\gamma$として$V=Re^{-(T-t)/\tau_v}$のように増え、
\begin{empheq}[box=\keybox]{equation}
\delta(t) \;=\; \frac{\dd V}{\dd t}-\frac{V}{\tau_\gamma} \;=\; V\Bigl(\frac{1}{\tau_v}-\frac{1}{\tau_\gamma}\Bigr) \;>\; 0 .
\label{eq:ramp}
\end{empheq}
誤差は正で、$V$とともに増える。これがあのなめらかな上昇である（図~\ref{fig:ramp}）。$\tau_\gamma=4\,\text{s}$、$\tau_v=1\,\text{s}$なら、毎秒$\delta=0.75\,V$となる。この説は仮想通路で、動物を瞬時に目標の近くへ移すことで検証された。ドーパミンはなめらかなレベルではなく、微分にふさわしくバーストで応答した~\cite{Kim2020}。

\begin{figure}[t]
\centering
\begin{tikzpicture}[>=Stealth,x=1.6cm,y=2.6cm]
\draw[->,gray!70!black] (-4.2,0) -- (0.5,0) node[below left,black] {\small $t$};
\draw[->,gray!70!black] (-4.2,0) -- (-4.2,1.2);
\foreach \x/\l in {-4/{$T-4$},-2/{$T-2$},0/{$T$}}{\draw[gray!60!black] (\x,-0.03) -- (\x,0.03); \node[below,font=\scriptsize] at (\x,0) {\l\,s};}
\draw[thick,densely dashed,gray!60!black] plot[domain=-4:0,samples=40] (\x,{exp(\x/4)});
\draw[very thick,blue!60!black] plot[domain=-4:0,samples=60] (\x,{exp(\x)});
\draw[very thick,red!70!black] plot[domain=-4:0,samples=60] (\x,{0.75*exp(\x)});
\node[anchor=south east,font=\small,gray!50!black] at (-1.3,0.77) {$\tau_v=\tau_\gamma$での$V$：$\delta=0$};
\node[right,font=\small,blue!60!black] at (0.05,1.0) {$\tau_v=1\,\text{s}$での$V$};
\node[right,font=\small,red!70!black] at (0.05,0.72) {\eqref{eq:ramp}による$\delta$};
\end{tikzpicture}
\caption{予測誤差としての、報酬への接近に伴うドーパミンの増加。$\tau_\gamma=4\,\text{s}$。破線は時間幅の要求どおりに増える期待（$\delta=0$）、青はもっと急に増える期待、赤は誤差~\eqref{eq:ramp}。}
\label{fig:ramp}
\end{figure}

測定されていること。なめらかな増加は、ドーパミン濃度にも\nt{DOPA}ニューロンの発火率にもあり~\cite{Kim2020}、状況の瞬間的な跳躍はドーパミンの跳躍を起こす。二つの説明のどちらが正しいのか、あるいは出力ごとに両方が正しいのかは議論されている。

\section{後ろを振り返る}

ここまでの理論はすべて\emph{前を}見ている。現代の理論の一つは方向を逆にする~\cite{Jeong2022}。脳は予告信号から報酬を予測することを学ぶのではなく、報酬を得てから\emph{振り返る}ことを学ぶとする。どの出来事が他より頻繁にそれに先行したか。この結びつきの尺度は二つの確率の差である。
\begin{empheq}[box=\keybox]{equation}
M \;=\; P(\text{報酬の前の予告信号}) \;-\; P(\text{ランダムな窓の予告信号}) .
\label{eq:retro}
\end{empheq}
予告信号が報酬なしにもよく現れるなら第二項が大きく、結びつきは弱い。この理論では、ドーパミンが伝えるのは予測誤差ではなく、その出来事がどれだけ報酬のありそうな\emph{原因}であるかである。

振り返りの機構は、第\ref{sec:dofa}章の三要素則と同じものを要求する。報酬の瞬間にその前に何があったかを読み出せるよう、各出来事は消えていくトレースを残さなければならない。違いはシナプスではなく、\nt{DOPA}が何を伝えるかにある。

理論~\eqref{eq:retro}と誤差~\eqref{eq:ctd}が異なる予測をする実験では、ドーパミン放出がいくつかの場合に前者に従った~\cite{Jeong2022}。しかし別の実験では、学習中のドーパミン応答は報酬の時点から予告信号の時点へと徐々に移った。誤差~\eqref{eq:ctd}による学習が要求するとおりである~\cite{Amo2022}。どちらの理論が脳を記述するかはまだ分からない。

\section{報酬だけではない誤差}

最後の拡張は、誤差が\emph{何についての}ものかに関する。期待していたジュースを、同じ価値で味の違う別のジュースに替えると、価値は変わらず、誤差~\eqref{eq:ctd}はゼロである。ところが\nt{DOPA}ニューロンはこの置き換えに応答する~\cite{Takahashi2017}。また人工的に起こしたドーパミンのバーストは、報酬なしに、二つの中立な刺激の結びつきを動物に学ばせることができる~\cite{Sharpe2017}。どうやら\nt{DOPA}は報酬だけでなく、\emph{何が}起こるかの予測誤差も伝えている。

形式的には、これは\eqref{eq:ctd}の単純な一般化である。報酬の流れ$r$一つの代わりに、起きていることの特徴$\varphi_k$（味、場所、音）の組をとり、それぞれに固有の期待$M_k$と固有の誤差をもたせる~\cite{Gardner2018}。
\begin{empheq}[box=\keybox]{equation}
\delta_k(t) \;=\; \varphi_k(t) \;+\; \frac{\dd M_k}{\dd t} \;-\; \frac{M_k}{\tau_\gamma} ,
\qquad k=1,\dots,K .
\label{eq:vectortd}
\end{empheq}
報酬は特徴の一つにすぎず、誤差のベクトルは何がよいかだけでなく、何が何に続くか、つまり世界のモデルを教える。

\nt{DOPA}ニューロンの不均一性はこれと整合する。一つの課題の中で、ニューロンごとに異なる量を運ぶ。報酬に関係するもの、運動に関係するもの、注意の向きに関係するものがある~\cite{Engelhard2019}。

\begin{keyblock}
\begin{proposition}[配線ではなく配線の束]
\label{prop:bundle}
\nt{DOPA}系の信号は一つの数ではない。ニューロンにわたって（分布~\eqref{eq:expectile}、異なる特徴~\eqref{eq:vectortd}、危険の別軸）、また時間にわたって（速い誤差とゆっくりしたレベル~\eqref{eq:vigor}）分解されている。スカラー誤差~\eqref{eq:ctd}はこの信号の第一近似であり、閾値付き加算器がニューロンの第一近似であるのと同じである。
\end{proposition}
\end{keyblock}

\section{まとめ}

\begin{table}[t]
\centering
\footnotesize
\setlength{\tabcolsep}{4pt}
\renewcommand{\arraystretch}{1.15}
\begin{tabular}{>{\raggedright}p{2.5cm}>{\raggedright}p{3.2cm}>{\raggedright}p{3.6cm}>{\raggedright\arraybackslash}p{4.2cm}}
\toprule
理論 & 配線を流れるもの & 主な測定 & 分かっていること \\
\midrule
予測誤差（第\ref{sec:dofa}章） & スカラー$\delta$~\eqref{eq:ctd} & バースト、予告信号への移行、落ち込み & 測定済み。第一近似 \\
分布 & 非対称性の異なる$\delta_i$の組 & ニューロンごとに異なる反転点 & 実験との一致あり。ばらつきの役割は仮説 \\
重要性 & $|\delta|$、新奇性。危険の軸 & 一部のニューロンでの不快なものへのバースト & 測定済み。使われ方は議論中 \\
時間のコスト & ゆっくりしたレベル$\bar R$ & ドーパミンは行動を速める。遅い成分は出力での放出にも発火率にもある & 速い成分と遅い成分の分離は測定済み。$\bar R$は仮説 \\
接近に伴う増加 & $V$、または推定の精密化に伴う$\delta$~\eqref{eq:ramp} & なめらかな増加、通路での移動時の跳躍 & 増加は測定済み。説明は議論中 \\
後ろを振り返る & 因果的結びつきの尺度~\eqref{eq:retro} & 二つの理論の予測が分かれる実験 & 結果が互いに矛盾 \\
誤差のベクトル & 特徴ごとの$\delta_k$~\eqref{eq:vectortd} & 味の変化への応答。報酬なしの学習 & 測定済み。ベクトル形式は仮説 \\
\bottomrule
\end{tabular}
\caption{\nt{DOPA}系が伝えるもの：第\ref{sec:dofa}章の標準的な描像とその現代的な拡張。}
\label{tab:dofaalt}
\end{table}

表~\ref{tab:dofaalt}の理論は互いに排他的ではない。ほぼすべてが予測誤差を基礎として残し、その形式を拡張している。本格的に競合するのは振り返りだけで、この論争は実験が決着をつける。エンジニアにとっての結論は単純である。配線が実は宛先の異なる多数の配線であれば、命題~\ref{prop:broadcastcost}のブロードキャストの代償は下がる。

\section{問題}

\begin{exercise}[\lvlA]
\label{ex:07:1}
\emph{1滴の分布の点。}大きさ$1$のジュースが確率$p=0.2$で来て、そうでなければ報酬は$0$である。\eqref{eq:expectile}から$\kappa_i=0.2$、$0.5$、$0.8$の$V_i$を求めよ。この報酬の中央値はいくらか。点~\eqref{eq:expectile}は分位点とどう違うか。
\end{exercise}

\begin{exercise}[\lvlB]
\label{ex:07:2}
\emph{2個のニューロンから分布を。}報酬は$0$か$x$で、$x$の確率は$p$である。規則~\eqref{eq:asymmetric}に従う2個のニューロンが$V(1/2)=0.25$、$V(0.8)=0.5$を報告した。$p$、$x$、報酬の標準偏差を復元せよ。$\kappa=0.2$のニューロンは何を報告するか。
\end{exercise}

\begin{exercise}[\lvlB]
\label{ex:07:3}
\emph{シミュレーション：非対称性から分布へ。}$\kappa_i=0.1,\dots,0.9$の9個の\nt{DOPA}ニューロンを、規則~\eqref{eq:asymmetric}、$\alpha_i^+=2\kappa_i$、$\alpha_i^-=2(1-\kappa_i)$、$[0,1]$上一様な報酬でシミュレートせよ。$V_i$のばらつきは$\eta$にどう依存するか。この分布を、等確率の2滴$0.25$と$0.75$から区別するにはどんな$\eta$が必要か。
\end{exercise}

\begin{exercise}[\lvlA]
\label{ex:07:4}
\emph{時間のコスト。}(a)~\eqref{eq:vigor}を導き、最適点での総コストを求めよ。(b)~脳が$\bar R$を$k$倍の誤差で推定した。$k=2$と$k=4$での相対的な払いすぎを求めよ。ドーパミンのゆっくりしたレベルは$\bar R$をどれだけ正確に伝えなければならないか。
\end{exercise}

\begin{exercise}[\lvlB]
\label{ex:07:5}
\emph{移動時のバースト。}期待は$\tau_v=1\,\text{s}$、$\tau_\gamma=4\,\text{s}$で\eqref{eq:ramp}に従って増え、動物は地点$T-2\,\text{s}$から$T-1\,\text{s}$へ瞬時に移される。$\delta$のバーストの面積を$R$に対する割合で求め、移動前のランプの何秒分にあたるかを求めよ。ドーパミンが$V$そのものを伝えるなら、移動は何を示すか。
\end{exercise}

\begin{exercise}[\lvlB]
\label{ex:07:6}
\emph{設計：配線上の二つの信号。}\nt{DOPA}の配線を、毎秒$s=1/60$スパイク/sの割合で増えるレベルと、$A=3$スパイクの出来事が流れる。トレース$\tau_e=1\,\text{s}$のシナプスは、発火率から$\tau_f$で平滑化したそのコピーを差し引く。出来事の損失が$10\,\%$以下で、かつトレースの時間内でのレベルの寄与が出来事の寄与の$10\,\%$未満になる$\tau_f$はどれか。
\end{exercise}

\begin{exercise}[\lvlB]
\label{ex:07:7}
\emph{実験の設計：振り返り対誤差。}窓の中で予告信号は確率$0.1$で現れ、その後の報酬は確率$0.8$である。$\Delta P=P(\text{報酬}\mid\text{予告})-P(\text{報酬}\mid\text{なし})$と\eqref{eq:retro}の$M$を、(a)~予告信号なしの報酬を窓あたり$0.08$加えた場合、(b)~報酬なしの予告信号の頻度を2倍にした場合で比べよ。
\end{exercise}

\begin{exercise}[\lvlC]
\label{ex:07:8}
\emph{配線の束とブロードキャストの代償。}問題~\ref{ex:06:8}のモデルで、$N$個のニューロンが自分の配線をもつ$K$群に分けられ、各配線には他群の誤差が割合$\rho$だけ漏れ込む。学習定数が$N/K+2+\rho^2N(1-1/K)$に等しいことを示せ。$K\to\infty$、$N=10^{10}$、$\rho=0.01$で何試行になるか。
\end{exercise}

\chapter{いつ探索し、いつ決めるか：\nt{NORA}を加える}
\chaptermark{\nt{NORA}を加える}
\label{sec:nora}
\begin{chaptergoals}
\item 一つのゲインのつまみがニューロンをアナログ増幅器からコンパレータに変える仕組み；
\item なぜゲインを上げて役立つのは増幅器の後で生じるノイズに対してだけなのか；
\item 重みを変えずにゲインが選択ネットワークを「思案中」から「決定済み」にする仕組み；
\item なぜゲインは逆温度として働き、報酬はゲインに対して逆U字になるのか。
\end{chaptergoals}

\section{足りないもの}

第\ref{sec:dofa}章でネットワークはノイズのおかげで学習した。活動のランダムなずれが探索であり、\nt{DOPA}はそれでよくなったかを伝えた。しかしそこでは一つのパラメータが未設定のまま残った。\emph{どれだけ}探索するかである。ノイズが強すぎると、ネットワークはあちこちに振れ、すでに知っていることを使わない。弱すぎると、毎回同じものを選び、近くのどこかがよくなったことに気づかない。強化学習ではこれを活用と探索の選択と呼び、どのアルゴリズムにもそのためのつまみがある。第\ref{sec:neurontypes}章で、脳ではこのつまみを回すのは\nt{NORA}だと仮定した。

ヒトのノルアドレナリンニューロンは数万個で（表~\ref{tab:types}）、ほぼすべてが一つの小さな集団に集まっている。その出力は脳のほぼ全体に広がるが、集団の部分ごとに、ある程度異なる領域を担当していることが分かっている~\cite{Poe2020}。これは\nt{DOPA}よりさらに細いブロードキャストチャネルである。動作モードは二つある~\cite{AstonJones2005}。\emph{持続}モードでは、ニューロンは毎秒1回未満から数回の発火率で規則的にスパイクを出し、この発火率は覚醒度とともにゆっくり変わる。\emph{相動}モードでは、現在の課題にとって重要な出来事に、ほぼ意思決定の瞬間に短いバーストで応答する。便利だが不正確なテストポイントが瞳孔径である。瞳孔径はこれらのニューロンの活動とともに変わるが、他の領域の活動~\cite{Joshi2016}や皮質のコリン作動性出力~\cite{Reimer2016}とともにも変わる。瞳孔が示すのは全体的な興奮レベルであり、\nt{NORA}だけではない。

\section{ゲイン}

測定されているのは次のことである。ノルアドレナリンはニューロンの背景活動を刺激への応答より強く抑え、信号対背景比が上がる~\cite{Foote1975NE}。個々のニューロンの特性の傾きが文字どおり係数倍されることは、実験からは導かれない。第\ref{sec:neurontypes}章と同じく、この効果を再現する簡単なモデルをとる。ノルアドレナリンは受け手の伝達特性の傾きを変える~\cite{ServanSchreiber1990}。すると弱い閾値下の入力（背景）はさらに小さな出力を、強い入力はさらに大きな出力を生む。ニューロンは発火率で応答するとする。
\begin{empheq}[box=\keybox]{equation}
r \;=\; F\bigl(g\,(u-\theta)\bigr), \qquad F(z)=\frac{r_{\max}}{1+e^{-z}} ,
\label{eq:gain}
\end{empheq}
ここで$u$は入力電流、$\theta$は閾値、$g$はゲインで、モデルではこれを\nt{NORA}が制御する。$g$が小さいと、発火率は広い範囲で入力になめらかに従う。ニューロンは\emph{アナログ増幅器}である。$g$が大きいと、閾値より上のほぼどんな入力も最大発火率を、下なら0を生む。ニューロンは\emph{コンパレータ}、ほぼディジタル素子である（図~\ref{fig:gaincurves}）。一つのつまみが、電子回路なら別々の回路が担う二つのモードの間で同じニューロンを切り替える。

\begin{figure}[t]
\centering
\begin{tikzpicture}[>=Stealth,x=1.1cm,y=2.4cm]
\draw[->,gray!70!black] (-3.3,0) -- (3.4,0) node[below left,black] {\small $u$};
\draw[->,gray!70!black] (-3.3,0) -- (-3.3,1.2) node[left,black] {\small $r$};
\draw[densely dashed,gray!60!black] (0,0) -- (0,1.1);
\node[below,font=\small] at (0,0) {$\theta$};
\draw[densely dotted,gray!60!black] (-3.3,1) -- (3.2,1) node[right,black,font=\small] {$r_{\max}$};
\draw[very thick,blue!60!black] plot[domain=-3.2:3.2,samples=60] (\x,{1/(1+exp(-0.8*\x))});
\draw[very thick,orange!85!black] plot[domain=-3.2:3.2,samples=60] (\x,{1/(1+exp(-2*\x))});
\draw[very thick,red!70!black] plot[domain=-3.2:3.2,samples=120] (\x,{1/(1+exp(min(-8*\x,9)))});
\node[blue!60!black,font=\small,anchor=west] at (0.5,0.12) {小さい$g$：増幅器};
\node[red!70!black,font=\small,anchor=east] at (-0.3,0.85) {大きい$g$：コンパレータ};
\end{tikzpicture}
\caption{ゲイン$g$の三つの値での伝達特性~\eqref{eq:gain}。}
\label{fig:gaincurves}
\end{figure}

大きなゲインはノイズに対して役立つか。いつもではない。二つの入力が$\Delta u$だけ異なり、どちらが大きいかを言わなければならないとする。ノイズ$\sigma_{\text{pre}}$が増幅器の\emph{前で}入力に加わるなら、信号もノイズも増幅され、その比は変わらない。ノイズ$\sigma_{\text{post}}$が増幅器の\emph{後で}加わるなら（第\ref{sec:code}章のように、信頼性の低いシナプスやネットワーク自体のランダムなスパイクから）、SN比
\begin{empheq}[box=\keybox]{equation}
d' \;=\; \frac{g\,\Delta u}{\sqrt{g^2\sigma_{\text{pre}}^2+\sigma_{\text{post}}^2}}
\label{eq:gainsnr}
\end{empheq}
は、$g\sigma_{\text{pre}}$が$\sigma_{\text{post}}$に並ぶまで$g$とともに増える~\cite{ServanSchreiber1990}。

\begin{keyblock}
\begin{proposition}[ゲインが役立つとき]
\label{prop:gainnoise}
ゲインを上げることで\nt{NORA}が入力の識別を改善できるのは、増幅器の\emph{後で}、ネットワーク自体の中で生じるノイズに対してだけである。入力と一緒に来たノイズはゲインでは除けない。
\end{proposition}
\end{keyblock}

\section{二者択一でのゲイン}

第\ref{sec:gaba}章の選択に戻ろう。入力$I_1$、$I_2$をもつ\nt{GLUT}ニューロンの二つの群が、強さ$\beta$で互いに抑制し合う。今度は各群にゲイン$g$がある。選択の方程式~\eqref{eq:wta}の括弧内の式が$g$倍される。両方の群が動いている間、定常発火率は$r_1=g(I_1-\beta r_2)$、$r_2=g(I_2-\beta r_1)$から求まる。これらを足し引きすると、和$r_1+r_2=g(I_1+I_2)/(1+g\beta)$と差$r_1-r_2=g(I_1-I_2)/(1-g\beta)$が得られる。その比は、ネットワークが一方の入力をどれだけはっきり好むかを表す。
\begin{empheq}[box=\keybox]{equation}
\frac{r_1-r_2}{r_1+r_2} \;=\; \varepsilon\,\frac{1+g\beta}{1-g\beta}, \qquad \varepsilon=\frac{I_1-I_2}{I_1+I_2}, \quad g\beta<1 .
\label{eq:wtagain}
\end{empheq}
入力の小さな差$\varepsilon$は$(1+g\beta)/(1-g\beta)$倍に増幅され、この係数は$g\beta$が1に近づくと際限なく大きくなる。$g\beta>1$では、命題~\ref{prop:wta}と同じく、両方の群が動く状態は不安定である。一方が勝ち、他方は沈黙する（図~\ref{fig:pitchfork}）。

\begin{figure}[t]
\centering
\begin{tikzpicture}[>=Stealth,x=3.2cm,y=1.5cm]
\draw[->,gray!70!black] (0,0) -- (2.15,0) node[below left,black] {\small $g\beta$};
\draw[->,gray!70!black] (0,-1.25) -- (0,1.35) node[above,black] {\small $(r_1-r_2)/(r_1+r_2)$};
\draw[gray!60!black] (1,-0.04) -- (1,0.04); \node[below right,font=\small] at (1,0) {$1$};
\node[left,font=\scriptsize] at (0,1) {$1$}; \node[left,font=\scriptsize] at (0,-1) {$-1$};
\draw[very thick,blue!60!black] plot[domain=0:0.905,samples=60] (\x,{0.05*(1+\x)/(1-\x)}) -- (1,1) -- (2.05,1);
\draw[very thick,blue!60!black] (1.105,-1) -- (2.05,-1);
\draw[thick,densely dashed,gray!60!black] (1,0) -- (2.05,0);
\node[font=\small,align=center] at (0.45,-0.62) {思案中：\\両方の群が\\動いている};
\node[font=\small,align=center] at (1.55,0.45) {決定済み：\\一方だけが動く};
\node[font=\small,blue!60!black,anchor=south west] at (1.05,-0.97) {弱い入力が先に勝つと、そのまま保たれる};
\end{tikzpicture}
\caption{ゲインを変えたときの二者択一。入力の差は$\varepsilon=0.05$。$g\beta>1$では両方の決定が安定で（実線。弱い入力の勝利は$g\beta>(1+\varepsilon)/(1-\varepsilon)\approx1.1$で）、ネットワークはすでに下した決定を保つ。対称な状態（破線）は不安定である。}
\label{fig:pitchfork}
\end{figure}

\begin{keyblock}
\begin{proposition}[ゲインが選択のモードを切り替える]
\label{prop:gainwta}
相互抑制$\beta$をもつ選択ネットワークで、ゲイン$g$はネットワークを境界$g\beta=1$の向こうへ移す。その下ではネットワークは「思案中」である。両方の群が動き、入力の差は\eqref{eq:wtagain}に従って増幅される。その上ではネットワークは「決定済み」である。一方の群だけが動き、決定は保たれる。$g$を変えることで、\nt{NORA}は重みに一つも触れずにネットワークをこの二つのモードの間で切り替えられる。
\end{proposition}
\end{keyblock}

\section{温度としてのゲイン}

次に、ネットワーク自体の中、増幅器の後で生じるノイズを選択に加えよう。どちらの群が勝つかは、量$g(u_A-u_B)+\eta$の符号で決まる。ここで$u_A-u_B$は入力の差、つまり第\ref{sec:dofa}章で学習した重みの差であり、$\eta$はスケール$\sigma$のノイズである。ノイズがロジスティック分布に従うなら（ガウス分布でも曲線はほぼ同じ）、$A$を選ぶ確率は
\begin{empheq}[box=\keybox]{equation}
P(A) \;=\; \frac{1}{1+e^{-g\,(u_A-u_B)/\sigma}} .
\label{eq:softmax}
\end{empheq}
プログラマーならここに温度$T=\sigma/g$のソフトマックス選択を見てとるだろう。ゲインは逆温度である。$g$が小さいと、明らかによい選択肢でも悪いほうよりわずかに多く選ばれるだけである。ネットワークは多く探索する。$g$が大きいと、既知の最良の選択肢がほぼ常に選ばれる。ネットワークは知っていることを活用する。

\eqref{eq:wtagain}と\eqref{eq:softmax}の$g$が何かを明確にしておく。これは選択の瞬間における、現在の課題の入力の差に対する\emph{実効}ゲインである。\nt{NORA}の二つのモードはこれに逆向きに作用する。決定の瞬間の相動性バーストはこれを上げ、ネットワークは選択を固める。逆に高い持続性活動は探索と一緒に現れる。ヒトでは当て推量の選択の前にベースラインの瞳孔が広く~\cite{JepmaNieuwenhuis2011}、ラットでは前帯状皮質への\nt{NORA}入力が選択をランダムなモードに移す~\cite{Tervo2014stochastic}。したがって高い持続性\nt{NORA}には\emph{小さな}実効$g$が、バーストには大きな実効$g$が対応する。

\section{どれだけ探索するか}

どのゲインが最もよいか。モデルで確かめた。ネットワークは\eqref{eq:softmax}に従って二つの行動の一方を選ぶ。どちらもある確率で報酬をもたらし、ネットワークは第\ref{sec:dofa}章のように、選んだ行動の報酬からその確率を推定する。ときどき世界が変わり、二つの確率が引き直される。変わるのは選んだ行動の場合もあれば、ネットワークが長く試していない行動の場合もある。後者についてネットワークが知る方法は探索しかない。図~\ref{fig:invertedU}は、一定の$g$をいろいろ変えたとき、ネットワークが最良の可能な報酬のどれだけの割合を得るかを示す。

\begin{figure}[t]
\centering
\begin{tikzpicture}[>=Stealth,x=1.2cm,y=1cm]
\draw[->,gray!70!black] (0,0) -- (7.6,0) node[below left,black] {\small $g$};
\draw[->,gray!70!black] (0,0) -- (0,4.4) node[above,black,font=\small] {最良の報酬に対する割合};
\foreach \x/\l in {0/0.5,1/1,2/2,3/4,4/8,5/16,6/32,7/64}{\draw[gray!60!black] (\x,-0.06) -- (\x,0.06); \node[below,font=\scriptsize] at (\x,0) {\l};}
\foreach \v/\y in {0.8/0.8,0.9/2.4,1.0/4.0}{\draw[gray!60!black] (-0.06,\y) -- (0.06,\y); \node[left,font=\scriptsize] at (0,\y) {\v};}
\draw[very thick,blue!60!black] plot[mark=*,mark size=1.3pt] coordinates {(0,0.368) (1,0.880) (2,1.728) (3,2.784) (4,3.456) (5,3.616) (6,3.488) (7,3.264)};
\draw[very thick,red!70!black] plot[mark=*,mark size=1.3pt] coordinates {(0,0.368) (1,0.672) (2,1.184) (3,1.840) (4,2.176) (5,1.920) (6,1.456) (7,1.104)};
\node[blue!60!black,font=\small,anchor=south] at (5,3.7) {穏やかな世界};
\node[red!70!black,font=\small,anchor=north] at (5.1,1.25) {速く変わる世界};
\draw[->,thick,blue!60!black] (5,3.25) -- (5,3.55);
\draw[->,thick,red!70!black] (4,1.8) -- (4,2.1);
\node[font=\small\itshape,gray!35!black,anchor=west] at (0.15,2.3) {やみくもに探す};
\node[font=\small\itshape,gray!35!black] at (6.45,2.55) {変化に気づかない};
\end{tikzpicture}
\caption{一定のゲイン$g$での、最良の可能な報酬に対する割合（$g$の目盛りは対数）。青の曲線は世界が平均1000試行に1回変わる場合、赤は20試行に1回の場合。矢印は最良の$g$を示す。速く変わる世界ではそれは半分になる。4000試行の実行60回の平均。}
\label{fig:invertedU}
\end{figure}

$g=0.5$ではネットワークは知識をほとんど使わず、可能な報酬のおよそ4分の3を得る。$g$が非常に大きいと変化を見逃す。最良の値は、世界が1000試行に1回変わるとき$g=16$（割合$0.976$）、20試行に1回変わるとき$g=8$（$0.886$）である。

\begin{keyblock}
\begin{proposition}[逆U字]
\label{prop:explore}
ネットワークが集める報酬は、ゲインに対して逆U字形に依存する。$g$が小さすぎると試行を探索に費やし、大きすぎると変化に気づかない。最良の$g$は、世界が速く変わるほど小さい。
\end{proposition}
\end{keyblock}

逆U字は実験でも観察される。動物が課題を最もよく解くのは、\nt{NORA}ニューロンの持続性活動が中程度で、相動性応答が明瞭なときであり、持続性活動が高いと気が散り、課題の間を切り替える~\cite{AstonJones2005}。これは前節のモードの違いと整合する。決定の瞬間の相動性バーストは、選ぶべきときにちょうど大きな実効$g$を与える。バーストのない高い持続性活動は、無関係な入力も含めて何でも増幅するので、現在の課題にとっての実効$g$は下がる。これが探索である。

\section{誰がつまみを回すのか}

最良のゲインが世界の変化の速さに依存するなら、それを調整できるとよい。しかし\nt{NORA}ニューロンは世界が変わったことをどうして知るのか。自然な手がかりは\emph{驚き}である。予測誤差がふだんより大きくなった。二種類の不確実性を区別することが大切である。報酬が単にランダムなら、誤差は常に大きく、それは予期されている。一方、誤差がふだんの大きさに比べて急に増えたなら、何かが変わったのである。ある仮説によれば、\nt{NORA}が伝えるのはまさにこの予期しない不確実性であり、予期された不確実性は\nt{ACET}が伝える~\cite{YuDayan2005}。

この信号で何をするか。ゲインを下げて探索を増やしてもよい。古くなった推定を速く忘れるよう学習率を上げてもよい。実験では、急な変化の後に人はより速く学び、そのとき瞳孔が広がる~\cite{Nassar2012}。瞳孔は\nt{NORA}だけでなく\nt{ACET}の指標でもあることを思い出そう~\cite{Reimer2016}。両方を同時にしてもよい。さらに別の仮説によれば、\nt{NORA}の相動性バーストは\emph{割り込み}として働く。プロセッサの共通の割り込み線のように、ネットワークが現在の状態をリセットし、新しい状況に合わせて組み直す合図である~\cite{BouretSara2005}。

見つからなかったことも正直に述べておく。モデルで二つの簡単な調整法を試した。平滑化した誤差が長期平均を超えたら$g$を下げる方法と、学習率を上げる方法である。長く静止したり頻繁に変わったりする世界で、どちらの方法も最良の設定では最良の報酬の$0.926$--$0.927$を得た。これは最良の一定の$g$（$g=8$で$0.925$）や、一定だが十分大きな学習率（$0.928$）とほとんど同じであり、設定が悪いとそれより少なかった。図~\ref{fig:invertedU}の曲線は頂上付近でなだらかで、このような世界では調整することがほとんどない。調整が割に合うのは、モードごとに大きく異なる設定が要る場合のはずである。\nt{NORA}がどんな制御則を実装しているのかは、まだ確立していない。

\section{まとめ}

\begin{table}[t]
\centering
\footnotesize
\setlength{\tabcolsep}{4pt}
\renewcommand{\arraystretch}{1.15}
\begin{tabular}{>{\raggedright}p{3.0cm}>{\raggedright}p{4.4cm}>{\raggedright\arraybackslash}p{6.0cm}}
\toprule
\nt{NORA}がすること & 方法 & 分かっていること \\
\midrule
ニューロンの応答の傾きを変える & \eqref{eq:gain}のゲイン$g$ & SN比の向上は測定済み。乗算モデルは単純化 \\
選択を切り替える & ネットワークを$g\beta=1$の向こうへ移す & モデル~\eqref{eq:wtagain}から導かれる。閾値の直接測定はない \\
どれだけ探索するかを決める & $g$は逆温度~\eqref{eq:softmax}。持続性は小さな実効$g$（探索）、バーストは大きな実効$g$（決定） & 逆U字と二つの動作モードは測定済み。探索との関係は十分な根拠のある仮説 \\
変化を伝える & 予期しない不確実性 & 変化の後に瞳孔と学習率が増すことは測定済み。それが\nt{NORA}の信号であることは仮説 \\
割り込み、リセットする & 予期しない出来事への相動性バースト & 重要な出来事へのバーストは測定済み。「割り込み」は仮説 \\
\bottomrule
\end{tabular}
\caption{\nt{GLUT}、\nt{GABA}、\nt{DOPA}のネットワークに\nt{NORA}が加えるもの。}
\label{tab:nora}
\end{table}

\nt{DOPA}はネットワークに重みを\emph{どこへ}変えるかを伝える。\nt{NORA}は重みを一つも変えないが、ネットワークがすでに知っていることを\emph{どう}使うかを決める（表~\ref{tab:nora}）。一つのつまみ、ゲインが、ニューロンを増幅器からコンパレータへ、選択ネットワークを「思案中」から「決定済み」へ、選択を探索から活用へと変える。\nt{NORA}が世界の変化の速さをどう知るのかは未解決の問題である。

\section{問題}

\begin{exercise}[\lvlA]
\label{ex:08:1}
\emph{脳全体に一つのつまみ。}(a)~表~\ref{tab:types}から、\nt{NORA}ニューロン1個あたり脳のニューロンが何個あるかを見積もれ。(b)~増幅器の前のノイズが$\sigma_{\text{pre}}=0.5$、後のノイズが$\sigma_{\text{post}}=4$である。\eqref{eq:gainsnr}の$d'$が上限の$90\,\%$に達するのはどの$g$か。
\end{exercise}

\begin{exercise}[\lvlB]
\label{ex:08:2}
\emph{ゲイン付きの選択。}$\tau\,\dd r_1/\dd t=-r_1+g[I_1-\beta r_2]_+$、$\tau\,\dd r_2/\dd t=-r_2+g[I_2-\beta r_1]_+$、$\tau=20\ms$。(a)~$g\beta=0.5$と$0.9$で、$r_1-r_2$が減衰する時間はいくらか。(b)~$\varepsilon=0.05$のとき、それより下で両方の群が動き、それより上で弱い入力の勝利が安定になる$g\beta$を求めよ。
\end{exercise}

\begin{exercise}[\lvlA]
\label{ex:08:3}
\emph{ノイズからソフトマックスへ。}$K$個の群はそれぞれ自分のガンベルノイズ$P(\eta_k<z)=\exp(-e^{-z/\sigma})$を受け、$gu_k+\eta_k$が最大のものが勝つ。$P(k)$を求めよ。$u_A-u_B=0.1$、$\sigma=1$の二つの選択肢のうちよいほうが$90\,\%$の確率で選ばれるのはどの$g$か。
\end{exercise}

\begin{exercise}[\lvlB]
\label{ex:08:4}
\emph{最良のゲインの見積もり。}図~\ref{fig:invertedU}のモデルで、変化後の報酬確率は$[0,1]$上で一様である。ネットワークは悪いほうの行動を$e^{g\Delta}$試行に1回試す。これを$1/h$に等しいとおくと$g^*\approx\ln(1/h)/\bar\Delta$となる。$\bar\Delta=\langle|p_A-p_B|\rangle$と、$h=0.001$、$0.01$、$0.05$での$g^*$を求めよ。
\end{exercise}

\begin{exercise}[\lvlB]
\label{ex:08:5}
\emph{シミュレーション：ゲインと変化の速さ。}\texttt{models/\allowbreak nora\_explore.py}のコピー（ファイルを現在のフォルダに書き出す）で、$h$が$0.001$から$0.2$までの$g^*(h)$を求めよ。問題~\ref{ex:08:4}の見積もりはどこで正しいか。変化が速いとき、$g^*$が減らなくなるのはなぜか。
\end{exercise}

\begin{exercise}[\lvlB]
\label{ex:08:6}
\emph{設計：選択ネットワークへの割り込み。}問題~\ref{ex:08:2}のネットワークは$\beta=2$、$\tau=20\ms$、$g=1$で、入力が今は$I_1=0.9$、$I_2=1.0$なのに$r_1=0.9$、$r_2=0$を保っている。\nt{NORA}は時間$T_p$の間$g$を$g_{\text{lo}}$に下げる。$g_{\text{lo}}=0$のときの最小の$T_p$を解析的に、$g_{\text{lo}}=0.25$のときをモデルで求めよ。
\end{exercise}

\begin{exercise}[\lvlC]
\label{ex:08:7}
\emph{調整が割に合うとき。}オラクルは穏やかな時期か変化の多い時期かを知っていて、それぞれに自分の$g$を設定する。(a)~本章の世界（各時期$1000$試行、$h=0.001$と$0.05$）で、最良の一定$g$に対するその利得をモデルで求めよ。(b)~利得がより大きい世界を構成し、驚きによる調整がそのどれだけを取り戻すかを求めよ。
\end{exercise}

\chapter{いつ書き、いつ読むか：\nt{ACET}を加える}
\chaptermark{\nt{ACET}を加える}
\label{sec:acet}

\section{足りないもの}

メモリチップには、アドレス線とデータ線のほかにもう1本の線、書き込み許可がある。これがオフの間、メモリは読み出されるだけである。オンになると、データ線上にあるものが書き込まれる。この線がなければメモリは働かない。読み出すたびに同時に何かが書き込まれるなら、読んだもの、誤って読んだものまで、すぐに書き戻されてしまう。

脳ではこの問題はチップよりも深刻である。第\ref{sec:glutcomputer}章では、記憶とは自身の結合によってオンに保たれる\nt{GLUT}ニューロンの群だった（図~\ref{fig:bistable}）。第\ref{sec:dofa}章では、これらの結合が動作中にそのままヘッブ則で学習することを見た。つまり同じシナプスが古いものを保持し、新しいものも受け入れる。そして独立した書き込みフェーズはなく、クロックもない。ネットワークは、見ているものを記憶すべき瞬間と、知っていることを思い出すべき瞬間をどう区別するのか。第\ref{sec:neurontypes}章で、この線を担うのは\nt{ACET}だと仮定した。この章では、そのような線がどう働くべきか、なぜそれなしでは済まないのかを確かめる。

\section{1本の配線の三つの作用}

脳の内部で働くアセチルコリンニューロンは少なく、いくつかの小さな群にまとまっているが、その出力は皮質全体に広がる。これは\nt{DOPA}や\nt{NORA}と同じブロードキャストチャネルである。皮質のアセチルコリンレベルは注意深い覚醒時に高く、徐波睡眠で低い~\cite{Hasselmo1999}。ドーパミンと同じく、信号の意味は受け手の受容体が決める（命題~\ref{prop:receptor}）。アセチルコリンには、チャネルを開く速い受容体と、細胞内の反応を起動する遅い受容体の両方がある。ここで重要な作用は三つある。

\emph{第一に、内部結合を弱める。}嗅皮質のスライスでの実験で、アセチルコリンは同じ領域のニューロン間の結合による伝達を強く抑え、外から入ってくる入力にはほとんど触れないことが示された~\cite{HasselmoBower1992}。

\emph{第二に、入力を強める。}アセチルコリンはニューロンの興奮性を高め、いくつかの入力経路の伝達を促進する。感覚器官からの信号がネットワーク自身の活動より大きくなる~\cite{Hasselmo2006}。

\emph{第三に、重みの変化を容易にする。}アセチルコリンのレベルが高いと、結合は変わりやすくなる~\cite{Hasselmo2006}。

エンジニアにとって、三つの作用はすべて一つの操作、「読む」モードから「書く」モードへの切り替えである。ネットワークは自分自身を聞くのをやめ、入力を聞き始め、聞いたことを記憶する。

\section{補完するネットワーク}

互いに結合した$N$個の\nt{GLUT}ニューロンをとる。その結合に、いくつかのパターン（同時に活動する$pN$個のニューロンの組）をヘッブ則で保存する~\cite{Hebb1949}。このネットワークにパターンの半分を与えると、結合が残りの半分を興奮させる。ネットワークは、図~\ref{fig:bistable}の群が自身を維持したのと同じように、一部からパターンを\emph{補完する}。これは連想記憶であり、内容の一部がアドレスになる。第\ref{sec:gaba}章と同じく、\nt{GABA}は活動するニューロンの総数を$pN$程度に保ち、二つのパターンが同時に点灯しないようにする。

アセチルコリンのレベルを$0$から$1$の$a$で表し、その三つの作用をそのままモデルに書き込む。
\begin{empheq}[box=\keybox]{equation}
\tau\frac{\dd r_i}{\dd t} = -r_i + F\Bigl(\tfrac{1+a}{2}\,x_i + (1-a)\sum_j W_{ij}\,r_j - g\Bigr),
\qquad
\frac{\dd W_{ij}}{\dd t} = \eta\,a\,(r_i-p)(r_j-p) .
\label{eq:acetnet}
\end{empheq}
ここで$r_i$はニューロン$i$の発火率、$x_i$は感覚器官からのその入力、$F$は閾値型の伝達特性、$g$は\nt{GABA}による全体的な抑制で、活動ニューロンが$pN$より多いと増える。係数$\frac{1+a}{2}$は入力の強化、$1-a$は内部結合の弱化、$\eta a$は学習率である。第二式はまばらなパターンに便利な形のヘッブ則~\eqref{eq:hebb}である~\cite{Tsodyks1988}。両方のニューロンがふだんより活動すると重みは増え、一方だけが活動すると減る。重みの負の部分は、いつものように\nt{GABA}仲介ニューロンが担う。クロックはない。ニューロンも重みも連続的に変わる。

\begin{figure}[t]
\centering
\begin{tikzpicture}[>=Stealth,
  nrn/.style={circle,draw,very thick,fill=blue!6,minimum size=0.8cm,inner sep=0pt,font=\small},
  acet/.style={circle,draw,very thick,double,double distance=1.2pt,fill=orange!15,minimum size=1.35cm,inner sep=0pt,font=\scriptsize},
  bc/.style={->,thick,densely dashed,orange!85!black}]
\foreach \i/\y in {1/2.4,2/1.2,3/0}{
  \node[nrn] (N\i) at (3,\y) {$r_\i$};
  \draw[->,thick,blue!60!black] (0,\y) node[left,black] {\small $x_\i$} -- (N\i);}
\node[left,align=right,font=\small] at (-0.5,1.2) {感覚器官\\から};
\draw[->,thick,gray!60!black] (N1) to[bend left=30] (N2);
\draw[->,thick,gray!60!black] (N2) to[bend left=30] (N1);
\draw[->,thick,gray!60!black] (N2) to[bend left=30] (N3);
\draw[->,thick,gray!60!black] (N3) to[bend left=30] (N2);
\draw[->,thick,gray!60!black] (N1.east) .. controls (4.6,2.0) and (4.6,0.4) .. (N3.east);
\node[right,font=\small,gray!40!black,align=left] at (4.7,1.2) {内部\\結合 $W$};
\node[acet] (A) at (1.2,-1.9) {\nt{ACET}};
\draw[bc] (A) -- (1.6,-0.15);
\node[font=\small,orange!60!black,anchor=east] at (0.95,-0.9) {$+$ 入力};
\draw[bc] (A) -- (4.3,0.6);
\node[font=\small,orange!60!black,anchor=west] at (4.3,0.1) {$-$ 内部結合、$+$ 学習率};
\node[right,font=\small,align=left] at (2.2,-1.9) {高レベル：「書く」\\低レベル：「読む」};
\end{tikzpicture}
\caption{書き込み許可線としての\nt{ACET}。ニューロン$r_i$は感覚器官からの入力$x_i$を受け（青い矢印）、パターンを保存している内部結合$W$（灰色）を通じて互いに興奮させる。アセチルコリン（破線の矢印）は、\eqref{eq:acetnet}のように入力を強め、内部結合を弱め、重みの変化を容易にする。}
\label{fig:acetnet}
\end{figure}

\section{書き込みと読み出しは互いに邪魔をする}

書き込みと読み出しが本当に衝突する課題で、モデル~\eqref{eq:acetnet}を確かめよう。$200$個のニューロンのネットワークが、活動ニューロン$20$個のパターンを五つすでに保存している。ここで六つ目を記憶しなければならないが、それは古いパターンの一つと半分重なっている。10個のニューロンが共通である。これはいつも起きることである。新しい部屋はなじみの部屋に似ていて、新しい顔は古い顔に似ている。

\emph{低い$a$での書き込み。}まずアセチルコリンの最初の二つの作用を第三の作用から切り離す。学習率はフルのままにし、変えるのは入力の強化と内部結合の弱化だけにする。$a=0$では、入力が共通の10個のニューロンを点灯させ、それらが内部結合を通じて\emph{古い}パターンの残り半分を点灯させる。\nt{GABA}は両方が同時に点灯することを許さず、自身の結合に支えられた古いパターンが、入力だけに支えられた新しいパターンを押しのける。ネットワークは目の前にあるものではなく覚えているものを見て、ヘッブ則はその見たものを律儀に書き込む。

結局、新しいパターンはまったく記憶されなかった。その特有の半分を与えても、ネットワークは何も思い出さない。そして古いパターンは壊れている。特有の半分で呼び出すと、平均して10個中6個の無関係なニューロンを引き連れてくる。$a=0.1$では新しいパターンは記憶されるようになるが、古いものと融合し、破損はさらにひどくなる。二つのそれぞれが、自分の半分で呼び出されると約8個の無関係なニューロンを引き連れる。$a=0.2$からは書き込みはきれいになり、想起時の余分なニューロンは1個未満になる（10回の実行の平均）。

\emph{本来の学習率での書き込み。}今度はアセチルコリンが、\eqref{eq:acetnet}のように学習率も決めるとする。1回の提示で新しいパターンが丸ごと記憶されるのは$a\ge0.9$の場合だけである。$a=0.75$では10分の8、$a\le0.5$ではまったく記憶されない。

\emph{読み出し。}五つのパターンがきれいに書き込まれたネットワークにパターンの半分を与え、それから取り去る。$a\le0.2$ではネットワークはパターンを丸ごと補完し、自力で保つ。$a=0.3$ではほぼ丸ごと。$a\ge0.4$では内部結合が弱められすぎていて、手がかりが消えると活動も消える（図~\ref{fig:acetsim}）。

\begin{figure}[t]
\centering
\begin{tikzpicture}[>=Stealth,x=9cm,y=3.2cm]
\draw[->,gray!70!black] (0,0) -- (1.1,0) node[right,black] {\small $a$};
\draw[->,gray!70!black] (0,0) -- (0,1.2);
\foreach \x in {0,0.2,0.4,0.6,0.8,1.0}{\draw[gray!60!black] (\x,-0.02) -- (\x,0.02); \node[below,font=\scriptsize] at (\x,0) {$\x$};}
\foreach \y in {0,0.5,1}{\draw[gray!60!black] (-0.01,\y) -- (0.01,\y); \node[left,font=\scriptsize] at (0,\y) {$\y$};}
\node[rotate=90,font=\small] at (-0.1,0.6) {復元されたパターンの割合};
\draw[very thick,blue!60!black] plot[mark=*,mark size=1.6pt] coordinates {(0,1.00) (0.1,1.00) (0.2,1.00) (0.3,0.96) (0.4,0.04) (0.5,0.00) (0.75,0.00) (1.0,0.00)};
\draw[very thick,red!70!black] plot[mark=*,mark size=1.6pt] coordinates {(0.1,0.00) (0.2,0.00) (0.3,0.00) (0.5,0.00) (0.75,0.80) (0.9,1.00) (1.0,1.00)};
\node[font=\small,blue!60!black,anchor=west,align=left] at (0.03,0.5) {読み出し：\\半分から\\パターンを};
\node[font=\small,red!70!black,anchor=east,align=right] at (1.0,0.35) {書き込み：\\新パターンを\\1回で};
\end{tikzpicture}
\caption{モデル~\eqref{eq:acetnet}：ニューロン$200$個、活動ニューロン$20$個のパターン、10回の実行の平均。青い点は読み出し。手がかりを取り去った後、半分からパターンのどれだけの割合が復元されたか。赤い点は書き込み。アセチルコリンが学習率も決める場合に、古いパターンと半分似た新しいパターンのどれだけの割合が1回の提示の後に思い出されるか。読み出しは$a\le0.3$で、書き込みは$a\ge0.9$で働く。}
\label{fig:acetsim}
\end{figure}

\begin{keyblock}
\begin{proposition}[書き込みと読み出しには別々のモードが要る]
\label{prop:writeread}
同じ結合が古いものを保持し新しいものを学習し、アセチルコリンが学習率も決めるモデル~\eqref{eq:acetnet}では、書き込みと読み出しがどちらも同じようにうまくいくレベル$a$は存在しない。書き込みには内部結合を弱めなければならない。さもないとネットワークは入力ではなく思い出したものを書き込む。読み出しには内部結合を強めなければならない。さもないと一部からパターンを補完できない。スイッチが必要であり、1本のブロードキャスト配線がその役を果たしうる。
\end{proposition}
\end{keyblock}

アセチルコリンが学習率を変えないなら、両方の仕事に使える狭い帯$0.2\le a\le0.3$がある。しかしそのときネットワークは読み出しのたびにも学習し、読んだものをその誤りごと書き直してしまう。学習中に内部結合を抑えるという考え自体は、スライスでの測定に基づくモデルからとったものである~\cite{HasselmoAndersonBower1992}。上の数値は私たちの簡略化したモデルについてのものであり、脳でモードの境界が正確にどこにあるかは分かっていない。

\section{目をどこまで信じるか}

「書く／読む」のスイッチは、より一般的な問い、入力をどれだけ信じ、記憶をどれだけ信じるかの極端な場合である。この問いにはエンジニアの厳密な答えがあり、それは1本の配線がなぜモードと学習率を同時に制御するのかを明らかにする。

ネットワークが、ゆっくりさまよう量$s(t)$を追跡しているとする。その分散は1秒あたり$q$だけ増える。感覚器官は$y(t)=s(t)+$（強度$R$のノイズ）を報告する。連続時間での最良の推定$\hat s$はカルマン--ビューシーフィルタで与えられる~\cite{KalmanBucy1961}。
\begin{empheq}[box=\keybox]{equation}
\frac{\dd\hat s}{\dd t} \;=\; \frac{P}{R}\,\bigl(y-\hat s\bigr),
\qquad
\frac{\dd P}{\dd t} \;=\; q-\frac{P^2}{R} ,
\label{eq:kalman}
\end{empheq}
ここで$P$は推定誤差の分散、つまりネットワーク自身が知っていることにどれだけ自信がないかである。第一式はニューロンモデル~\eqref{eq:glutneuron}の膜と同じRC回路だが、時定数$R/P$が変わる。ネットワークに自信がないと$P$は大きく、時定数は小さく、推定はすばやく入力に従う。これが「書く」である。自信があると$P$は小さく、推定はほとんど入力を聞かずに記憶にしがみつく。これが「読む」である。

定常状態では$P=\sqrt{qR}$で、記憶の時定数は$\sqrt{R/q}$である。世界が100秒で1単位ずれ、$q=0.01$、感覚器官のノイズが$R=1$なら、記憶は約10秒保つべきである。

ここで重要なのは、一つの数$P/R$が二つのことを同時に決めることである。記憶に対する入力の重みと、保持している推定が変わる速さである。これはまさに\eqref{eq:acetnet}でアセチルコリンがすることである。仮説~\cite{YuDayan2005}によれば、アセチルコリンはまさに$P$に似た量、\emph{予期された}不確実性、つまりネットワークがあらかじめ知っている不確実性をネットワークに伝える。このチャネルが常に遅いわけではない。\nt{ACET}ニューロンの一部は報酬と罰に20ミリ秒ほどで応答し、強化が予期しないものであるほど強く応答する~\cite{Hangya2015}。

\begin{keyblock}
\begin{proposition}[入力の重みと学習率に一つのつまみ]
\label{prop:kalman}
最適フィルタ~\eqref{eq:kalman}では、入力を記憶と混ぜる重みと学習率は同じ量$P/R$である。したがってそれらを一つの信号で制御するのは合理的である。世界の性質が変わらなければ、この信号は一定レベル$\sqrt{q/R}$に落ち着き、調整が必要なのは世界の性質が変わったときだけである。
\end{proposition}
\end{keyblock}

命題の第二文は第\ref{sec:nora}章の結果を説明する。そこでは驚きによる学習率の調整は最良の一定値に勝てなかった。変化の性質が一定の世界では、最良の学習率は一定値そのものである。利得が生じるのは、世界が突然別物になったときである。そのとき推定は急に入力から大きく外れるが、$P$はそれを知らない。正しい行動は$P$を急に上げ、それによって書き込みモードに移ることである。第\ref{sec:nora}章で論じた仮説によれば、そのような\emph{予期しない}変化を伝えるのは\nt{NORA}である~\cite{YuDayan2005}。分業は自然に決まる。\nt{NORA}は世界が変わったことに気づき、\nt{ACET}は新しい状況が学習されるまでネットワークを書き込みモードに保つ。

\section{読み出しモードとしての睡眠}

アセチルコリンの高レベルが書き込みモードなら、低レベルのときネットワークは何をするのか。徐波睡眠ではアセチルコリンのレベルが下がり、内部結合は抑制から解放され、感覚器官からの入力はほとんどない。読み出しモードのネットワークは、手がかりなしに、自分に書き込まれているものを再生する。活動の記録によれば、日中一緒に働いたニューロンは徐波睡眠でふたたび一緒に発火し~\cite{WilsonMcNaughton1994}、しかも同じ順序で発火する~\cite{Skaggs1996}。仮説~\cite{Hasselmo1999}によれば、これらの再生は日中にすばやく学習したものを長期記憶に書き写す（再生の短いバーストを人工的に延ばすと記憶がよくなる~\cite{FernandezRuiz2019}）。日中は入力からの書き込み、夜間は内部からの書き直しである。エンジニアにとっては、日中に高速バッファに書き込み、夜間にそれを低速の不揮発性メモリにフラッシュするのに似ている。

\section{まとめ}

\begin{table}[t]
\centering
\footnotesize
\setlength{\tabcolsep}{4pt}
\renewcommand{\arraystretch}{1.15}
\begin{tabular}{>{\raggedright}p{3.2cm}>{\raggedright}p{4.2cm}>{\raggedright\arraybackslash}p{6.0cm}}
\toprule
\nt{ACET}がすること & 方法 & 分かっていること \\
\midrule
内部結合を弱める & \eqref{eq:acetnet}の係数$1-a$ & スライスで測定済み \\
入力を強める & 係数$\frac{1+a}{2}$ & 興奮性と入力経路の伝達の増大は測定済み \\
学習を容易にする & 学習率$\eta a$ & 測定済み \\
「書く／読む」を切り替える & 命題~\ref{prop:writeread} & モデルから導かれる。モードの直接測定はない \\
予期された不確実性を伝える & \eqref{eq:kalman}の$P$ & 仮説 \\
睡眠中の再生のためにネットワークを解放する & 徐波睡眠での低い$a$ & 睡眠中の低レベルと再生は測定済み。長期記憶への書き直しは仮説 \\
\bottomrule
\end{tabular}
\caption{\nt{GLUT}、\nt{GABA}、\nt{DOPA}、\nt{NORA}のネットワークに\nt{ACET}が加えるもの。}
\label{tab:acet}
\end{table}

\nt{DOPA}はネットワークに重みをどこへ変えるかを伝え、\nt{NORA}は知っていることをどれだけ断固として使うかを伝え、\nt{ACET}は世界を聞くか自分を聞くかを伝える（表~\ref{tab:acet}）。これは表~\ref{tab:types}の最後のブロードキャスト制御線、書き込み許可である。これがなければ、パターンを読み出すのと同じ結合にパターンを保存する記憶は、読み出すたびに自分を壊してしまう。

\section{問題}

\begin{exercise}[\lvlA]
\label{ex:09:1}
\emph{どれだけ覚えておくか。}$q=0.01$、$R=1$のフィルタ~\eqref{eq:kalman}は約$10$~s覚えている。(a)~センサノイズの振幅が2倍になった場合、(b)~世界のさまよいが100倍速くなった場合について、$P_\infty$と記憶$\sqrt{R/q}$を求めよ。(c)~ネットワークが1分覚えているには、ノイズは何倍にならなければならないか。
\end{exercise}

\begin{exercise}[\lvlB]
\label{ex:09:2}
\emph{変化の後の書き込みモード。}世界が変わり、$q=0.01$、$R=1$のフィルタ~\eqref{eq:kalman}の不確実性が$P_0\to\infty$に跳ね上がった。(a)~$P$はどんな時定数で$P_\infty$に戻るか。(b)~何秒後に学習率が定常値を$10\,\%$以内でしか上回らなくなるか。
\end{exercise}

\begin{exercise}[\lvlB]
\label{ex:09:3}
\emph{一定の学習率。}ネットワークは$\dd\hat s/\dd t=(y-\hat s)/T$で学習する。世界はさまよい、$\dd s/\dd t=w$、$y=s+n$で、$w$と$n$は強度$q$と$R$の白色ノイズである。最良の$T$とそのときの誤差の分散を求めよ。$T$が$2$倍、$10$倍ずれていると、分散は何倍になるか。
\end{exercise}

\begin{exercise}[\lvlB]
\label{ex:09:4}
\emph{書き込みの境界はどこか。}\texttt{models/\allowbreak acet\_memory.py}では閾値$\theta=0.35$、パターン内の重み$w=0.041$（理想的には$0.045$）である。新しいパターンは古いパターンと$m=10$個のニューロンを共有する。\eqref{eq:acetnet}でどの$a$なら、古いパターンの残りのニューロンがこの$m$個から点灯しないか（閾値は鋭いとする）。
\end{exercise}

\begin{exercise}[\lvlB]
\label{ex:09:5}
\emph{シミュレーション：記憶容量。}\texttt{models/\allowbreak acet\_memory.py}の関数\texttt{read\_sweep}を変更せよ。$a=1$で$M$個のパターン（$M$は$5$から$100$）を書き込み、次に$a=0$で最初の10個を半分から読み出す。どの$M$で誤りが現れるか。限界$M_{\max}$は$N$と$p$にどう依存するか。
\end{exercise}

\begin{exercise}[\lvlB]
\label{ex:09:6}
\emph{設計：漏れのない書き込み。}学習率$\eta a$では、ネットワークは読み出し時にも学習する。$a\le0.3$でゼロ、$a\ge0.9$で$\eta a$に劣らない単調な$\eta(a)\le\eta$を提案せよ。検証：無関係なニューロンを3個含む手がかりで$a=0.2$で$20$回読み出した後、そのうち何個がパターンに入ったか。
\end{exercise}

\begin{exercise}[\lvlC]
\label{ex:09:7}
\emph{夜間にバッファをどう書き直すか。}(a)~睡眠中は$a=0.05$で、再生は$0.1\,\text{s}$続く（日中は$a=1$で$0.2\,\text{s}$）。学習率$\eta a$と問題~\ref{ex:09:6}の$\eta(a)$のそれぞれで、日中の1回の提示と同じ重み変化を蓄積するには何回の再生が要るか。(b)~ストレージはバッファより$100$倍遅く学習し、一晩にパターンを$0.1\,\text{s}$ずつ$20$回聞く。何晩必要か。
\end{exercise}

\chapter{マシン全体}
\chaptermark{マシン全体}
\label{sec:synthesis}
\begin{chaptergoals}
\item \nt{GLUT}、\nt{GABA}と四つのブロードキャストチャネルが一つの三層マシンになる仕組み；
\item すべてのつまみ$\delta$、$\tau_\gamma$、$g$、$a$を一組の式にまとめる方法；
\item 世界が変わるときにつまみがどう協調しうるか、そしてどの部分が仮説か；
\item このマシンが普通のコンピュータとどこが違うか、そして何がまだ分かっていないか。
\end{chaptergoals}

\section{何を組み立ててきたか}

ここまで表~\ref{tab:types}のニューロン型を一つずつ取り上げ、それが計算機に何を付け加えるかを問うてきた。ルールは毎回同じだった。生きた生物に実際にあるものだけを使い、共通のクロックは使わない。いよいよ部品を一つのマシンに組み上げ、それらがどう協調して働くかを見よう。

第\ref{sec:neurontypes}章で描いた全体像は裏づけられ、より精密になった。\nt{GLUT}ニューロンの巨大なアドレス網がデータを運ぶ。\nt{GABA}ニューロンがそのデータの流れを制御する。そして網の上には四つのブロードキャストチャネルがあり、それぞれが自分のつまみを持つ。\nt{DOPA}はどの重み変化を残すかを告げ、\nt{SERO}は全体のレベルを保ち、仮説によれば計画の時間幅も決める。\nt{NORA}は探索か決定かを選び、\nt{ACET}は書き込みか読み出しかを選ぶ（図~\ref{fig:machine}）。

\begin{figure}[t]
\centering
\begin{tikzpicture}[>=Stealth,
  blk/.style={draw,very thick,rounded corners=3pt,align=center,font=\small},
  reg/.style={draw,very thick,double,double distance=1.2pt,circle,minimum size=1.25cm,inner sep=0pt,font=\scriptsize,fill=orange!15},
  bc/.style={->,thick,densely dashed,orange!85!black}]
% sensors, bus, actions
\node[blk,fill=green!8,text width=1.7cm] (S) at (0.9,0) {センサ\\\scriptsize オン／オフ};
\draw[very thick,rounded corners=4pt,fill=blue!5] (2.6,-1.35) rectangle (9.0,1.35);
\node[font=\small] at (5.8,0.95) {\nt{GLUT}バス：データ};
\node[font=\scriptsize,align=center] at (4.3,0.25) {同時性、遅延、\\論理（第\ref{sec:glutcomputer}章）};
\node[font=\scriptsize,align=center] at (7.4,0.25) {記憶、符号\\（第\ref{sec:code}, \ref{sec:acet}章）};
\draw[thick,red!70!black,fill=red!8,rounded corners=2pt] (2.8,-1.15) rectangle (8.8,-0.55);
\node[font=\scriptsize,red!60!black] at (5.8,-0.85) {\nt{GABA}（第\ref{sec:gaba}章）：禁止、選択、除算、リズム};
\node[blk,fill=blue!5,text width=2.0cm] (A) at (10.9,0) {行動の\\選択};
\draw[->,very thick] (S) -- (2.6,0);
\draw[->,very thick] (9.0,0) -- (A);
\draw[->,very thick] (A) -- (12.6,0) node[right,font=\small] {筋肉};
\pic[scale=1.1] at (13.2,-0.5) {spring};
\pic[scale=1.15] at (0.4,0.85) {eyepic};
\pic[scale=1.05] at (1.35,0.85) {speaker};
% registers
\node[reg] (D) at (3.4,3.2) {\nt{DOPA}};
\node[reg] (C) at (5.6,3.2) {\nt{SERO}};
\node[reg] (N) at (7.8,3.2) {\nt{NORA}};
\node[reg] (H) at (10.0,3.2) {\nt{ACET}};
\node[font=\scriptsize,align=center,above=1pt] at (D.north) {誤差 $\delta$};
\node[font=\scriptsize,align=center,above=1pt] at (C.north) {レベル、$\tau_\gamma$};
\node[font=\scriptsize,align=center,above=1pt] at (N.north) {ゲイン $g$};
\node[font=\scriptsize,align=center,above=1pt] at (H.north) {書き込み $a$};
\draw[bc] (D.south) -- (3.4,1.35);
\draw[bc] (C.south) -- (5.6,1.35);
\draw[bc] (N.south) -- (7.8,1.35);
\draw[bc] (H.south) -- (10.0,2.1) -- (8.8,1.35);
\draw[bc] (H.south east) to[bend left=20] (A.north east);
\draw[->,thick] (2.85,1.35) -- (2.85,2.75) -- (D.west) node[pos=0.25,left,font=\scriptsize] {$V$};
\draw[->,thick,green!45!black] (0.4,3.2) node[left,black,font=\small] {報酬 $r$} -- (2.2,3.2) -- (D.west);
\pic[scale=0.9] at (-0.45,3.7) {juice};
\draw[->,thick,densely dotted,gray!50!black] ([yshift=0.45cm]D.north) to[bend left=18] node[above,font=\scriptsize,black] {意外さ} ([yshift=0.45cm]N.north);
\end{tikzpicture}
\caption{マシンの全体像。\nt{GLUT}バスはデータをセンサから行動の選択へ運び、\nt{GABA}がその内部の流れを制御する。二重丸はブロードキャストチャネル、破線の矢印はその信号を網全体へ配ることを表し、チャネルごとに自分のつまみがある。\nt{DOPA}は報酬 $r$ と、バス内部の評価器（クリティック）からの期待値 $V$ を受け取る（第\ref{sec:dofa}章）。点線の矢印は、\nt{NORA}が異常に大きな誤差から意外さを知るという仮説を表す（第\ref{sec:nora}章）。}
\label{fig:machine}
\end{figure}

\section{すべてのつまみを一つの式に}

つまみは一つの模式的な式にまとめられる。これは新しいモデルではなく、第\ref{sec:glutcomputer}--\ref{sec:acet}章のモデルの要約であり、どの記号もそれぞれの章から取ってある。
\begin{empheq}[box=\keybox]{equation}
\begin{aligned}
\tau\,\frac{\dd r_i}{\dd t} \;&=\; -r_i + F\Bigl(g\,\Bigl[\tfrac{1+a}{2}\,x_i + (1-a)\sum_j W_{ij}\,r_j - \sum_k M_{ik}\,q_k - \theta\Bigr]\Bigr),\\
\frac{\dd W_{ij}}{\dd t} \;&=\; \eta\,a\,\delta(t)\,e_{ij}(t),
\qquad \tau_e\,\frac{\dd e_{ij}}{\dd t} = -e_{ij} + r_i\,r_j,
\qquad \delta = r + \frac{\dd V}{\dd t} - \frac{V}{\tau_\gamma} .
\end{aligned}
\label{eq:machine}
\end{empheq}
ここで $r_i$ は\nt{GLUT}ニューロンの発火率、$x_i$ はセンサからの入力、$W_{ij}\ge0$ はそれらの間の重みである（第\ref{sec:glutcomputer}章）。$q_k$ は\nt{GABA}ニューロンの発火率、$M_{ik}\ge0$ はその抑制の強さ（第\ref{sec:gaba}章）。$e_{ij}$ は適格度トレース、$\delta$ は\nt{DOPA}の誤差（第\ref{sec:dofa}章）。$\tau_\gamma$ は仮説上\nt{SERO}が決める時間幅、$g$ は\nt{NORA}のゲイン（第\ref{sec:nora}章）、$a$ は\nt{ACET}のレベル（第\ref{sec:acet}章）である。

\eqref{eq:machine}に何が\emph{ない}かを見てほしい。クロックはない。すべては微分方程式で書かれ、各ニューロンは自分で変化する。独立したメモリもない。記憶するのは計算に使われるのと同じ $W_{ij}$ と、同じ活動 $r_i$ である。圧倒的多数のシナプスには個別の補正もない。その学習は局所的な適格度トレースと一つの共通の数との積である。出力ごとに専用の誤差配線を持つのは、運動を学習する回路だけである（第\ref{sec:motorlearning}章）。そして制御信号はわずか四種類、$\delta$, $\tau_\gamma$, $g$, $a$ で、これが八百億のニューロンを受け持つ。実際にはどれも一つの数ではなく、並列な線の小さな束である（命題~\ref{prop:bundle}）。だが束の中の線は数本にすぎない。

\begin{keyblock}
\begin{proposition}[アーキテクチャ]
\label{prop:architecture}
現在の理解の範囲で、脳というマシンは三つの層で記述できる。データは\nt{GLUT}のアドレスバスを流れる。データの流れを方向づけ、制限し、正規化するのはアドレス型の\nt{GABA}ニューロンである。動作モードを決めるのは数本のブロードキャストチャネルで、それぞれが網の広い領域に共通な並列線の小さな束である。最初の二層は確立している。ブロードキャストチャネル間の役割分担は、大部分が仮説である。
\end{proposition}
\end{keyblock}

\section{すべての型を一つの表に}

表~\ref{tab:synthesis}は本書のすべてのニューロン型をまとめる。自分の章を持たなかった四つの型も一行ずつ占める。そのうち二つはマシンの出力を扱う章で役割を得た。\nt{GLYC}は脊髄のループで（第\ref{sec:muscles}章）、\nt{ADRE}は化学的出力で（第\ref{sec:chemoutput}章）。残る二つの計算上の役割は、単純であるか未知である。ニューロンの型を決めない物質は付録~\ref{sec:othermediators}にまとめた。

\begin{table}[t]
\centering
\footnotesize
\setlength{\tabcolsep}{3pt}
\renewcommand{\arraystretch}{1.15}
\begin{tabular}{>{\raggedright}p{1.0cm}>{\raggedright}p{1.0cm}>{\raggedright}p{3.8cm}>{\raggedright}p{1.4cm}>{\raggedright}p{2.6cm}>{\raggedright\arraybackslash}p{3.5cm}}
\toprule
型 & 章 & 何を計算するか & 時間 & バス & 分かっていること \\
\midrule
\nt{GLUT} & \ref{sec:glutcomputer}, \ref{sec:code}, \ref{sec:sensors}, \ref{sec:motorlearning}, \ref{sec:dendrites} & データ：同時性、遅延、論理、記憶、系列 & ms & アドレス & 確立 \\
\nt{GABA} & \ref{sec:gaba}, \ref{sec:code}, \ref{sec:pain}, \ref{sec:motorlearning} & 流れの制御：禁止、選択、除算、安定性、リズム；痛みのゲート & ms & アドレス & 確立；発火率の除算は背景ノイズとの組み合わせで \\
\nt{SERO} & \ref{sec:serocomputer}, \ref{sec:pain}, \ref{sec:sleep} & レベル調整器、平滑化；時間幅 $\tau_\gamma$；痛みの音量；睡眠のモード & 秒 & 容積 & 自己抑制は測定済み；時間幅は仮説 \\
\nt{DOPA} & \ref{sec:dofa}, \ref{sec:dofaalt} & 予測誤差 $\delta$；遅いレベルは時間の値段 & $0.1$\,s と分 & ブロードキャスト & $\delta$ は測定済み；拡張は第\ref{sec:dofaalt}章 \\
\nt{NORA} & \ref{sec:nora}, \ref{sec:pain}, \ref{sec:chemoutput}, \ref{sec:sleep} & ゲイン $g$：探索か決定か；意外さ；臓器への「アクセル」 & 秒 & ブロードキャスト & SN比の向上は測定済み；探索での役割は仮説 \\
\nt{ACET} & \ref{sec:acet}, \ref{sec:muscles}, \ref{sec:chemoutput}, \ref{sec:sleep} & 書き込みか読み出しか；学習速度；筋肉への出力；臓器への「ブレーキ」 & 秒 & 両方 & 三つの効果は測定済み；モード切り替えは仮説 \\
\nt{GLYC} & \ref{sec:muscles}, \ref{sec:pain} & 脊髄と脳幹での負の重み：拮抗筋の禁止 & ms & アドレス & 確立 \\
\nt{HIST} & --- & 全体への「起きていよ」信号 & 遅い & ブロードキャスト & 計算上の役割は未研究 \\
\nt{ADRE} & \ref{sec:chemoutput} & 血液を通じた全身への「アクセル」；脳内では不明 & 遅い & ブロードキャスト & 脳の外では確立；脳内の役割は不明 \\
\nt{ASPA} & --- & 弱い正の重み & ms & アドレス & 役割には異論あり \\
\bottomrule
\end{tabular}
\caption{本書のすべてのニューロン型：それぞれが何を計算し、それがどこまで分かっているか。}
\label{tab:synthesis}
\end{table}

\section{つまみはどう協調するか}

これまで各章は一つのつまみだけを回してきた。今度は一つの出来事をマシン全体で追ってみよう（図~\ref{fig:timeline}）。動物はずっと前に、行動 $A$ がジュースをもたらすと学習している。ある時点で世界が変わる。$A$ はもうジュースをもたらさず、代わりに、動物がほとんど試さない行動 $B$ がジュースをもたらすようになる。

\emph{誤差。} $A$ の後にジュースが来ず、\nt{DOPA}ニューロンは発火の落ち込みで応える。誤差の式~\eqref{eq:ctd}で $\delta<0$ である。いま $A$ を選んだばかりのシナプスは適格度トレースを持っており、弱まる。

\emph{意外さ。} 一度の落ち込みはよくある偶然である。だが落ち込みが繰り返され、誤差がふだんより大きくなる。第\ref{sec:nora}章の仮説では、これこそが\nt{NORA}の信号、「世界が変わった」である。ゲイン $g$ が下がって選択が柔らかくなり、ノイズによって $B$ が試される頻度が上がる。

\emph{書き込み。} 第\ref{sec:acet}章の仮説では、変化の後に\nt{ACET}が上がり、網を書き込みモードに切り替える。古い習慣を保持する内部結合は弱められ、センサからの入力は強められ、重みは容易に変わる。

\emph{新しい習慣。} $B$ を試すと、誰も予期していなかったジュースが来る。発火のバースト、$\delta>0$ であり、$B$ を選んだシナプスが強まる。こうした試行が何回か続くと、期待値 $V$ は新しい世界に追いつき、誤差は再び小さくなる。

\emph{復帰。} 意外さが去り、\nt{NORA}は高いゲインを戻し、選択は再び硬くなる。\nt{ACET}は下がり、網は読み出しモードで学習結果を使う。この間ずっと、仮説によれば\nt{SERO}が時間幅 $\tau_\gamma$、つまりどれほど先のジュースまで待つ価値があるかを決めている。

\begin{figure}[t]
\centering
\begin{tikzpicture}[>=Stealth,x=1.1cm,y=0.95cm]
\foreach \y/\lab in {4/{$\delta$（\nt{DOPA}）},3/{$g$（\nt{NORA}）},2/{$a$（\nt{ACET}）},1/{$B$ の選択}}{
  \draw[gray!60!black] (0,\y-0.35) -- (10,\y-0.35);
  \node[left,font=\scriptsize] at (0,\y) {\lab};}
\draw[densely dashed,gray!60!black] (2,0.4) -- (2,4.55) node[above,font=\scriptsize,black] {世界が変わった};
\draw[densely dashed,gray!60!black] (5,0.4) -- (5,4.55) node[above,font=\scriptsize,black] {$B$ で初めてのジュース};
\pic[scale=0.85] at (2,5.25) {nojuice};
\pic[scale=0.85] at (5,5.25) {juice};
% delta: dips after the change, burst at the first juice for B, then back to zero
\draw[thick,red!70!black] (0,3.65) -- (2.3,3.65) -- (2.3,3.3) -- (2.5,3.3) -- (2.5,3.65) -- (3.2,3.65) -- (3.2,3.3) -- (3.4,3.3) -- (3.4,3.65) -- (4.1,3.65) -- (4.1,3.35) -- (4.3,3.35) -- (4.3,3.65) -- (5.0,3.65) -- (5.0,4.35) -- (5.2,4.35) -- (5.2,3.65) -- (5.8,3.65) -- (5.8,4.1) -- (6.0,4.1) -- (6.0,3.65) -- (6.6,3.65) -- (6.6,3.85) -- (6.8,3.85) -- (6.8,3.65) -- (10,3.65);
% gain: high, drops after surprise, recovers
\draw[thick,blue!60!black] (0,3.2) -- (3.0,3.2) -- (3.4,2.75) -- (5.8,2.75) -- (7.2,3.2) -- (10,3.2);
% ACh: low, rises after the change, falls after learning
\draw[thick,green!45!black] (0,1.75) -- (2.8,1.75) -- (3.3,2.25) -- (6.8,2.25) -- (7.8,1.75) -- (10,1.75);
% choice of B
\draw[thick,black] (0,0.7) -- (3.0,0.7) plot[domain=3:10,samples=40] (\x,{0.7+0.6/(1+exp(-(\x-5.3)*1.8))});
\draw[->,gray!70!black] (0,0.2) -- (10.2,0.2) node[below left,font=\scriptsize,black] {時間};
\end{tikzpicture}
\caption{一つの出来事をマシン全体で追う。これは第\ref{sec:dofa}--\ref{sec:acet}章のモデルと仮説にもとづく定性的な模式図であり、計算結果ではない。変化の後、\nt{DOPA}は落ち込みで応える。落ち込みが繰り返されると、仮説では意外さの信号が上がり、\nt{NORA}がゲインを下げ、\nt{ACET}が書き込みを入れる。$B$ で初めてのジュースがバーストを生み、選択は $B$ に移り、つまみは元に戻る。}
\label{fig:timeline}
\end{figure}

どのつまみも、他のつまみが何をしているかを知らないことに注意しよう。それぞれは自分の手がかり（誤差、その異常な大きさ、不確かさ）に反応し、動作の順序はひとりでに決まる。各ブロックが入力のそろった時点で動作する非同期回路と同じである。我々の知る限り、このような協調動作の全体はまだ検証されていない。個々のつまみは測定されているが、それらの協調は仮説である。

\section{このマシンはコンピュータとどこが違うか}

\begin{table}[t]
\centering
\footnotesize
\setlength{\tabcolsep}{4pt}
\renewcommand{\arraystretch}{1.15}
\begin{tabular}{>{\raggedright}p{2.2cm}>{\raggedright}p{4.2cm}>{\raggedright\arraybackslash}p{6.6cm}}
\toprule
& 通常のコンピュータ & 脳 \\
\midrule
時間 & 共通のクロック、約1ギガヘルツ & 共通のクロックはない。各ニューロンは自分で変化し、窓は出来事と局所的なリズムが決める（第\ref{sec:glutcomputer}, \ref{sec:gaba}, \ref{sec:code}章） \\
素子 & 信頼できる2値ゲート & アナログのしきい素子。シナプスはスパイクの大部分を失う（第\ref{sec:code}章） \\
結合の符号 & 任意 & 送り手のニューロンが決める（第\ref{sec:neurontypes}章） \\
記憶 & プロセッサとは別 & 計算と同じ結合の中と、自己維持する活動の中（第\ref{sec:glutcomputer}, \ref{sec:acet}章） \\
学習 & 誤差逆伝播 & 局所的な規則、一つのブロードキャストされる数 $\delta$（第\ref{sec:dofa}章）と、運動回路の出力への誤差配線（第\ref{sec:motorlearning}章） \\
制御 & レジスタと命令バス & 四つのブロードキャストチャネル、それぞれ数本の線の束（第\ref{sec:serocomputer}, \ref{sec:dofa}--\ref{sec:acet}章） \\
エネルギー & プロセッサ1個に数十〜数百ワット & 脳全体で約 $20$ ワット、ニューロンあたり平均で毎秒約1スパイク（第\ref{sec:code}章） \\
\bottomrule
\end{tabular}
\caption{脳と通常のコンピュータ。}
\label{tab:computer}
\end{table}

表~\ref{tab:computer}は主な違いをまとめる。どれも自然が直しそびれた欠点ではなく、一つの設計の一部である。共通のクロックがなければ、「オン／オフ」センサと同時検出器が必要になる。信頼できない素子には、多数のニューロンによる冗長性が必要になる。計算と一体の記憶には、書き込みが読み出しを壊さないように\nt{ACET}が必要になる。逆向きの配線のない学習には、\nt{DOPA}とノイズが必要になる。そして毎秒1スパイクというエネルギーの上限が、言語全体を倹約的にし、スパイクをニュースのために使わせる。

\section{分かっていないこと}

このまとめを締めくくるのに最も誠実なのは、分かっていないことの一覧である。どの項目もエンジニアにとっての課題である。

\emph{符号。} スパイクチャネルの容量は分かっているし、受け手がメッセージのどの部分を読むかを選ぶことも分かっている（第\ref{sec:code}章）。だが皮質がスパイクの正確なタイミングをどこまで使っているのか、ニューロンに「単語」があるのかは分かっていない。

\emph{多層を通した学習。} ブロードキャスト信号による学習は遅く、結果に影響するニューロンが一つ増えるたびにさらに遅くなる（命題~\ref{prop:broadcastcost}）。では皮質は多層回路の数十億の重みをどう調整するのか。これは分かっていない。層の間で誤差をスパイクのバーストが運ぶモデルはある~\cite{Payeur2021}。答えの一部は、ニューロン自身の樹状突起の枝の中での計算にあるのかもしれない。そこでは一つのニューロンが小さな多層ネットワークのようにふるまう。皮質ニューロン一つのモデルの応答を再現するのに、人工ネットワークは5〜8層を要する~\cite{Beniaguev2021}。ヒトのニューロンの枝は排他的論理和さえ計算できる~\cite{Gidon2020}。もう一つの候補は自己教師あり学習である。皮質が自分自身の入力を予測するように学習するなら、誤差は各層でその場に生じ、共通の信号はいらない（第\ref{sec:dofa}章）~\cite{Keller2018}。

\emph{ブロードキャストチャネルが実際に何を運ぶか。} \nt{DOPA}には標準的な描像とその六つの拡張がある（第\ref{sec:dofaalt}章）。\nt{NORA}と\nt{ACET}にはもっともらしい仮説がある（第\ref{sec:nora}, \ref{sec:acet}章）。\nt{SERO}については大部分が推測である。

\emph{時間的な協調。} 関数をどう計算するか、出来事からの時間をどう計るか、局所的なリズムで流れをどう窓に切り分けるかは見てきた。だが共通のクロックを持たない巨大な網が遠く離れた領域の働きをどう協調させるのかは、部分的にしか分かっていない。

\emph{エネルギー。} すべて合わせて20ワット。符号がなぜ疎でなければならないかは分かっている（命題~\ref{prop:economy}）。だが同じことを同じ電力でこなすマシンを作れる者は、まだいない~\cite{MehonicKenyon2022}。

脳はエンジニアが組み立てたものではない計算機だが、その法則はどのエンジニアにも見覚えがある。RC回路、閾値、フィードバック、フィルタ、バス、ブロードキャストレジスタ。その回路図を我々はまだ一部しか読めていない。続きを読むのは、回路図を読める人たちである。

\section{本書のこの先}

この章はマシンの計算の中核をまとめた。残りの章はその外へ出る。第\ref{sec:sensors}章と第\ref{sec:recognition}章は入力についてである。センサが光、音、分子をどうスパイクに変え、マシンがその中にどう物体を見分けるか。第\ref{sec:muscles}--\ref{sec:chemoutput}章は出力とそれを支えるものについてである。筋肉、警報信号としての痛み、専用の誤差配線による運動学習、そして体への遅い化学的出力。第\ref{sec:debugging}章は、ブレイン・コンピュータ・インターフェースを含め、マシンを外からどうデバッグするかを扱う。第\ref{sec:eetypes}章と第\ref{sec:dendrites}章は、今度は電気的な機能と入力の数でニューロンをもう一度分類する。第\ref{sec:sleep}章はマシンが世界から切り離されたときに何をするかを、第\ref{sec:livingmachine}章と第\ref{sec:dishbrain}章はチップ上の生きたニューロンで組み立てたマシンを扱う。

\section{問題}

この章の問題は、異なる章の結果を結びつける。どの問題も少なくとも二つの章にもとづく。

\begin{exercise}[\lvlA]
\label{ex:10:1}
\emph{バスとレジスタ。} 四つのブロードキャストチャネルはそれぞれ約5本の線の束である（命題~\ref{prop:bundle}）。各線は毎秒1回変化し、4つのレベルを区別できる。\nt{GLUT}バス（$8.6\cdot10^{10}$ 個のニューロン、発火率は\eqref{eq:energy}、精度は\eqref{eq:spikeentropy}により $1\ms$）が1秒で送る量を、制御系が送るには何年かかるか。
\end{exercise}

\begin{exercise}[\lvlB]
\label{ex:10:2}
\emph{一つのループに二つのつまみ。} ループ~\eqref{eq:ei}で、\eqref{eq:machine}のつまみが興奮側にかかるとする：$\tau_E\dot E=-E+g[(1-a)w_{EE}E-w_{EI}I+h]$。\nt{GABA}はつまみの制御を受けない。図~\ref{fig:ei}の数値で、\nt{NORA}のバーストが $g$ を $6$ に上げる。ループが安定となる最小の $a$ はいくつか。\nt{NORA}と\nt{ACET}のつまみは独立に回せるか。
\end{exercise}

\begin{exercise}[\lvlB]
\label{ex:10:3}
\emph{ブロードキャストの代償はどこから来るか。} $N$ 個のニューロンの出力 $y_k=f(u_k)+\xi_k$ に分散 $\sigma^2$ の独立なガウスノイズがあり、$R\approx R_0+\sum_kR'_k\xi_k$（$R'_k$ はすべて等しい）、\nt{DOPA}は $\delta=R-R_0$ をブロードキャストする。1回の試行での補正 $\delta\,\xi_jx_j$ のSN比を求めよ。試行数は $N$ とともにどう増えるか。
\end{exercise}

\begin{exercise}[\lvlB]
\label{ex:10:4}
\emph{音と閃光を結びつける。} 音は $5\ms$ 後にバスに届き、閃光は $30\ms$ に初発スパイクの遅延~\eqref{eq:latency}（$\tau=10\ms$、$U$ は $1.2\theta$ から $4\theta$）を加えた時間後に届く。各入力の背景は毎秒 $5$ スパイクである。すべてのコントラストに対して、音の線の遅延と最小の同時性窓を選べ。背景は毎秒何回の誤った結びつけを生むか。
\end{exercise}

\begin{exercise}[\lvlB]
\label{ex:10:5}
\emph{シミュレーション：変化に気づくのは誰か。} クリティックは $y=s+\text{ノイズ}$（$R=1$）を見る。確率 $1/200$ で $s$ は分散 $J^2=4$ の新しい値へ跳ぶ。$q=0$ で、$E\leftarrow0.9E+\delta^2/(P+R)$ が $20$ を超えたら $P$ を $J^2$ に引き上げるカルマンフィルタをシミュレーションせよ。その誤差は最良の一定の $\alpha$ の場合よりどれだけ小さいか。
\end{exercise}

\begin{exercise}[\lvlA]
\label{ex:10:6}
\emph{習慣を変えるのに何回の試行が要るか。} 網は $Q_A=0.8$, $Q_B=0.2$ を学習し、\eqref{eq:softmax}に従い $\sigma=1$, $g=8$ で選ぶ。いま $A$ は確率 $0.2$ でジュースを出す。$Q_A\leftarrow Q_A+\alpha(r-Q_A)$ で、$Q_B$ は変わらない。$\alpha=0.1$ と $0.5$ のとき、$A$ を何回選べば $A$ を選ぶ確率が $0.6$ まで下がるか。\nt{NORA}が $g$ を $2$ に下げたらどうか。
\end{exercise}

\begin{exercise}[\lvlC]
\label{ex:10:7}
\emph{配線の代わりに束。} 出力 $y_k=w_k+\xi_k$、報酬 $R=-\sum_ky_k^2$ とする。束の1本の線が $G$ 個のニューロンからなるグループにその出力の誤差を配り、$w_k\leftarrow w_k+\eta\,\delta_l\xi_k$ とする。最良の $\eta$ で $\sum w_k^2=\varepsilon\sum w_k^2(0)$ に達するには $\approx(G+2)\ln(1/\varepsilon)$ 回の試行が要ることを示せ。$\varepsilon=0.01$ で、$10^6$ 個のニューロンの回路が5本の線で、3秒に1回の試行で学習するとどれだけかかるか。
\end{exercise}

\chapter{マシンの入力：目、耳、その他のセンサはどう接続されているか}
\chaptermark{マシンの入力}
\label{sec:sensors}

\section{変換器としてのセンサ}

図~\ref{fig:machine}では、データは左側の「センサ」ブロックからマシンに入ってきた。今度はそのブロックを開けてみよう。エンジニアから見れば、どの感覚器官も\emph{変換器}である。アナログの入力段があり、最後にA/D変換器がある。ただし出力のデジタル値は数ではなく、スパイクである。

感覚器官の構成はどれも同じである（図~\ref{fig:transducer}）。光、圧力、振動、分子といった物理量が、感覚細胞の膜にある受容体タンパク質に作用する。このタンパク質はゲートとして働き、自らイオンチャネルを開くか、短い反応の連鎖を介して開く。\emph{受容器電流}が流れ、それとともに膜にアナログ電圧が生じる。電圧はゲイン調整を経てスパイク符号器に入る。符号器はおなじみの積分型ニューロン~\eqref{eq:glutneuron}で、レベルを発火率とスパイクの時刻に変える。出力はほぼ例外なく同じである。目、耳、嗅覚、皮膚の感覚ニューロンはグルタミン酸を放出するので、入力はそのまま\nt{GLUT}バスに接続される。

\begin{figure}[t]
\centering
\begin{tikzpicture}[>=Stealth,
  blk/.style={draw,very thick,rounded corners=2pt,align=center,font=\footnotesize,minimum height=1.0cm,text width=1.9cm,fill=blue!5}]
\node[align=center,font=\small] (P) at (0.1,0) {光、音、\\圧力、\\分子};
\node[blk] (R) at (3.0,0) {受容体\\ゲート};
\node[blk] (G) at (6.5,0) {ゲイン\\調整};
\node[blk] (E) at (10.0,0) {スパイク\\符号器};
\node[align=center,font=\small] (B) at (13.4,0) {\nt{GLUT}\\バス};
\draw[->,very thick] (P) -- (R);
\foreach \ic/\xx in {sun/-1.1,speaker/-0.3,press/0.5,molecule/1.3}{\pic[scale=0.85] at (\xx,1.05) {\ic};}
\draw[->,very thick] (R) -- node[above,font=\scriptsize] {電流} (G);
\draw[->,very thick] (G) -- node[above,font=\scriptsize] {レベル} (E);
\draw[->,very thick,red!70!black] (E) -- node[above,font=\scriptsize,black] {スパイク} (B);
\draw[decorate,decoration={brace,amplitude=5pt,mirror},gray!60!black] (1.8,-0.8) -- (7.7,-0.8) node[midway,below=5pt,font=\scriptsize,black] {アナログ部};
\draw[decorate,decoration={brace,amplitude=5pt,mirror},gray!60!black] (8.8,-0.8) -- (11.2,-0.8) node[midway,below=5pt,font=\scriptsize,black] {スパイク};
\end{tikzpicture}
\caption{変換の連鎖としての感覚器官。受容体タンパク質がチャネルを開いて電流を生む。ゲイン調整がそれを状況に合わせる。符号器~\eqref{eq:glutneuron}がレベルをスパイクに変え、スパイクは\nt{GLUT}バスへ出ていく。目と耳では最初の段はスパイクをまったく出さない。遠くへ送る必要が生じるまで、信号はアナログのままである。}
\label{fig:transducer}
\end{figure}

ある細部は電子回路の技術者にはなじみ深い。目と耳では、いちばん最初の細胞はスパイクをまったく出さない。基板上のアナログ段のように、なめらかな電圧を隣の細胞に渡すだけである。スパイクが現れるのは、信号を遠くへ運ぶ必要がある場所だけである。アナログ信号は数ミリの距離で減衰し、ノイズを拾う。一方スパイクは、第\ref{sec:neurontypes}章で見たように、同じ高さのまま配線の終端まで届く。デジタル化は長い線路が始まる場所で行われる。

\section{物理の限界で}

脳のセンサは感度の物理的限界で働いている。暗所の光センサは\emph{光子1個}に応答する。光の粒子1個の吸収が約1ピコアンペアの電流を生み、それは自身のノイズを背景にしても検出できる~\cite{Baylor1979}。音のセンサは、自分の感覚毛束の1ナノメートル未満の変位を検出する~\cite{Hudspeth1989}。

光が少ないとき、精度を制限するのはセンサではなく光そのものである。光子はポアソン流としてランダムに到来する。蓄積時間 $T$ の間にセンサに平均 $N=\Phi T$ 個の光子が入るなら、その数は $\sqrt N$ だけゆらぐ。増分 $\Delta N$ を検出するにはこのゆらぎの数倍でなければならず、識別できる明るさの変化の割合は次のようになる：
\begin{empheq}[box=\keybox]{equation}
\frac{\Delta I}{I} \;\approx\; \frac{k}{\sqrt{\Phi\,T}} .
\label{eq:photonnoise}
\end{empheq}
これはスパイク数についての式~\eqref{eq:rateerror}と同じである。光子のショットノイズとスパイクのショットノイズは同じ法則に従う。暗所で目が精度を上げる方法は一つしかない。光子をより長く蓄積することであり、実際に蓄積時間は十分の数秒まで伸びる。だから薄暗がりではものの見え方が遅くなる。

\section{ダイナミックレンジ：ゲイン調整}

センサにとって最も難しい工学的課題はレンジである。照度は星明かりの夜から日の当たる雪原まで約10桁変わり、音の強さは聴覚のしきい値から痛みの限界まで12桁変わる。一方、スパイクの発火率はゼロから毎秒数百までで、3桁にも満たない。このような入力をこのような出力に線形に収めることはできない。

解決法はどの受信機とも同じ、\emph{自動ゲイン調整}である。目の感覚細胞の応答は次の式でよく記述される：
\begin{empheq}[box=\keybox]{equation}
r \;=\; r_{\max}\,\frac{I}{I+\sigma} ,
\label{eq:nakarushton}
\end{empheq}
ここで $\sigma$ は応答が半分に達する明るさである~\cite{NakaRushton1966}。\eqref{eq:nakarushton}自体は2〜3桁しかカバーしない。だが $\sigma$ は一定ではない。明るい背景で長く働くと、平均の明るさを追って大きくなる。応答曲線は明るさの対数軸に沿ってずれ、そのつど急峻な部分を現在の入力の位置に置く（図~\ref{fig:adaptation}）。これは第\ref{sec:gaba}章のシャント抑制と同じ、全体の活動による除算である~\cite{Carandini2012}。ただしここでの分母は、直近数秒の平均の明るさである。

\begin{figure}[t]
\centering
\begin{tikzpicture}[>=Stealth,x=1.25cm,y=2.8cm]
\draw[->,gray!70!black] (-2,0) -- (6.4,0) node[below left,font=\small,black] {$\log_{10} I$};
\draw[->,gray!70!black] (-2,0) -- (-2,1.15) node[left,font=\small,black] {$r/r_{\max}$};
\foreach \u in {-2,0,2,4,6}{\draw[gray!60!black] (\u,-0.02) -- (\u,0.02); \node[below,font=\scriptsize] at (\u,0) {$\u$};}
\draw[densely dashed,gray!60!black] (-2,0.5) -- (6.2,0.5);
\foreach \s/\c/\lab in {0/blue!60!black/{暗い背景},2/green!45!black/{中程度},4/red!70!black/{明るい}}{
  \pgfmathsetmacro{\lo}{max(-2,\s-3)}
  \draw[very thick,\c] (-2,0) -- (\lo,0) plot[domain=\lo:6.2,samples=80] (\x,{1/(1+10^(\s-\x))});
  \node[font=\scriptsize,\c,anchor=west] at (\s+0.15,0.42) {\lab};}
\foreach \s/\ic in {0/moon,2/cloud,4/sun}{\pic at (\s+0.6,0.64) {\ic};}
\end{tikzpicture}
\caption{\eqref{eq:nakarushton}による光センサのゲイン調整。各曲線は明るさの2〜3桁をカバーする。より明るい背景に移ると $\sigma$ が大きくなり、曲線は対数軸に沿ってずれる。ずれていく曲線全体でレンジ全体をカバーする。}
\label{fig:adaptation}
\end{figure}

ゲイン調整には、第\ref{sec:code}章でおなじみの代償がある。センサが伝えるのは明るさではなく、\emph{背景に対する}明るさである。一定の入力は次第に応答を起こさなくなり、変化のたびに応答が起きる。マシンの入力は最初から、レベルではなく変化について語る。ここにも命題~\ref{prop:economy}と同じスパイクの節約がある~\cite{Barlow1961}。第\ref{sec:glutcomputer}章のオン型とオフ型のセンサもここから来る~\cite{Schiller1992}。一方は明るくなったことを、他方は暗くなったことを伝える。

エンジニアはこの手法を\emph{イベントカメラ}で再現した。このカメラはクロックに合わせてフレームを撮るのではない。各画素が、共通のクロックなしに独立して、明るさの対数が所定の割合だけ変化したときに「明るく」または「暗く」のイベントを出す~\cite{Lichtsteiner2008}。これはまさに第\ref{sec:glutcomputer}章のオン／オフの2線式符号をシリコンで実現したものであり、通常のイメージセンサには得られないレンジと速さをカメラに与える~\cite{Gallego2022}。

\section{目：1本のケーブルへの圧縮}

目にはおよそ $10^8$ 個の光センサがあるが~\cite{Curcio1990}、画像を脳へ運ぶ出力ニューロンは約100万個しかない。信号は目を出る前に約100分の1に圧縮される。圧縮の主な手法はすでにおなじみである。各出力ニューロンは、小さな斑点の明るさとその周囲の明るさとの差に応答する。これは第\ref{sec:gaba}章と同じ、抑制による近傍の減算である。画像の一様な部分はほとんどコストがかからず、エッジと変化が伝えられる。出力ニューロンは種類ごとに特化している。明るくなったことを伝えるもの、暗くなったことを伝えるもの、特定の方向の運動を伝えるものがある（図~\ref{fig:direction}）。マウスではこうした出力チャネルが30種類以上数えられている~\cite{Baden2016}。

できあがったケーブルの伝送容量はどれほどか。モルモットでは出力ニューロンの記録から、$\sim10^5$ 個のニューロンで毎秒約 $875$ キロビットが測定された。ヒトの100万個の出力ニューロンに換算すると、毎秒 $10^7$ ビットのオーダーになる~\cite{Koch2006}。昔の10メガビットのネットワークケーブルと同じである。ニューロンあたりの平均発火率が毎秒数スパイクなら、1スパイクあたり数ビットとなり、第\ref{sec:code}章で得た値と一致する。

\section{耳：フィルタバンク}

耳が解く課題は別である。音を周波数に分解しなければならない。エンジニアならバンドパスフィルタのバンクを置くだろう。耳はまさにそれを機械的に行う。音の振動は弾性のある隔壁の上を伝わり、隔壁の性質は長さに沿ってなめらかに変わるので、各部分がそれぞれの周波数で共振する。その部分にあるセンサは、自分の帯域にだけ応答する（図~\ref{fig:filterbank}）。周波数は場所に沿ってほぼ対数的に並ぶ。どのオクターブにも同程度の割合のセンサが割り当てられる。応答したセンサの番号が周波数であり、第\ref{sec:code}章の場所符号である。

\begin{figure}[t]
\centering
\begin{tikzpicture}[>=Stealth,x=1.35cm,y=2.2cm]
\draw[->,gray!70!black] (-0.3,0) -- (7.6,0) node[above left,font=\small,black] {周波数、Hz};
\draw[->,gray!70!black] (-0.3,0) -- (-0.3,1.15) node[left,font=\small,black] {応答};
\foreach \k/\f in {0/125,1/250,2/500,3/1000,4/2000,5/4000,6/8000}{
  \draw[gray!60!black] (\k,-0.02) -- (\k,0.02); \node[below,font=\scriptsize] at (\k,0) {\f};
  \draw[thick,blue!60!black] plot[domain={max(\k-1.2,-0.3)}:{\k+0.7},samples=60] (\x,{1/(1+((\x-\k)/(\x<\k?0.45:0.2))^4)});}
% a piano keyboard: seven white keys per octave, like the octaves on the axis
\foreach \i in {0,...,48}{\draw[gray!50!black,fill=white,line width=0.3pt] ({\i/7},-0.44) rectangle ({(\i+1)/7},-0.26);}
\foreach \o in {0,...,6}{\foreach \j in {0,1,3,4,5}{\fill[black] ({\o+(\j+1)/7-0.035},-0.35) rectangle ({\o+(\j+1)/7+0.035},-0.26);}}
\end{tikzpicture}
\caption{バンドパスフィルタのバンクとしての耳（模式図）。各曲線は、ある部分のセンサがさまざまな周波数の音にどう応答するかを示す。中心周波数は、対数目盛上の鍵盤のように1オクターブおきに並ぶ。実際のフィルタは高周波側の斜面が急で、低周波側がなだらかである。下は鍵盤で、耳と同じく、どのオクターブにも同じ幅が割り当てられている。}
\label{fig:filterbank}
\end{figure}

耳の共振器は受動的ではない。一部のセンサは自ら音に合わせて隔壁を揺らし、弱い振動を振幅で数千倍（$66$--$76$\,dB）増幅する。音が大きくなると増幅は下がる~\cite{Ruggero1997,Robles2001}。エンジニアから見れば、これは自励発振の境界ぎりぎりに置かれた正帰還増幅器である。第\ref{sec:gaba}章のネットワークと同じく、境界ぎわで生きている。境界の近くでは、系は小さな信号にとりわけ敏感になる。音量の圧縮も同じ増幅器が行う。小さな音は大きく、大きな音は小さく増幅される。

耳にはもう一つの符号がある。低い周波数では、ニューロンのスパイクは音と位相をそろえて出る。ニューロンは毎回の波で発火するわけではないが、発火するときは波の決まった部分で発火する。モルモットでは位相への同期は約 $600$\,Hz で弱まり始め、約 $3.5$\,kHz まではまだ認められる~\cite{PalmerRussell1986}。こうしてスパイクの時刻には音の正確な時間が保たれ、そこから両耳への音の到達時間差が得られる。第\ref{sec:glutcomputer}章のネットワークはこの差から方向を求める~\cite{Grothe2010}。高い周波数ではスパイクが波に追いつけず、場所符号だけが残る。一つの耳が第\ref{sec:code}章の二つの符号を両方使っている。ニューロンが追いつけるところでは時間を、追いつけないところでは場所を。

\section{その他のセンサ}

\begin{table}[t]
\centering
\footnotesize
\setlength{\tabcolsep}{3pt}
\renewcommand{\arraystretch}{1.15}
\begin{tabular}{>{\raggedright}p{1.9cm}>{\raggedright}p{2.6cm}>{\raggedright}p{4.0cm}>{\raggedright\arraybackslash}p{4.6cm}}
\toprule
感覚 & 何を測るか & 入力の仕組み & 本書の手法 \\
\midrule
視覚 & 光 & $\sim10^8$ 個のセンサ、$\sim100$ 分の1に圧縮、$\sim10^7$ bit/s & ゲイン調整、近傍の減算、2本の配線、運動の方向 \\
聴覚 & 空気の圧力 & 機械的フィルタバンク、能動的増幅器 & 場所符号と時間符号、遅延 \\
触覚 & 圧力、振動 & 速く順応するセンサと遅く順応するセンサ & 「ハイパスフィルタとローパスフィルタ」の対 \\
温度 & 熱 & 温センサと冷センサが別々 & 2線式符号 \\
嗅覚 & 分子 & 受容体は $\sim400$ 種類、各センサは1種類 & 組み合わせ符号：匂いは応答した種類の集合 \\
味覚 & 食物中の物質 & 五つの味質：甘味、苦味、塩味、酸味、うま味 & ラベル付き回線；一部の細胞はグルタミン酸ではなくATPやセロトニンを放出 \\
身体の位置 & 筋肉の伸張 & 伸張の長さと速さのセンサ & 運動制御器へのフィードバック \\
\bottomrule
\end{tabular}
\caption{センサはどう接続されているか。3列目の仕組みはすべて測定済みである。右の列はそれらを前の章の手法と結びつける。}
\label{tab:sensors}
\end{table}

残りの感覚は表~\ref{tab:sensors}にまとめた。ほとんどどの行にも、前の章の手法が見てとれる。触覚は一対のフィルタを使う。ある圧センサは接触に応答してすぐに沈黙し、別の圧センサは圧力がある間ずっと応答を保つ。温度は温と冷で別々の2本の配線で伝えられる。最も変わった入力は嗅覚である。ヒトの嗅覚受容体は約400種類あり、各嗅覚ニューロンはそのうち1種類だけを持つ~\cite{Mombaerts2004}。匂いは1種類ではなく一組の種類を活性化し、メッセージは\emph{どの}種類が応答したかである。エンジニアから見れば、これは長さ約400の2値ベクトルであり、とりうる値の数は天文学的である。

\begin{keyblock}
\begin{proposition}[マシンの入力]
\label{prop:sensors}
どの感覚も変換器であり、感度の物理的限界で働き、巨大なダイナミックレンジを自動ゲイン調整で圧縮し、\nt{GLUT}バスを通じて何よりも\emph{変化}を伝える。そのためにセンサはマシン自身と同じ手法を使う。2線式符号、場所符号、時間符号、そして背景による除算である。センサの仕組みは測定されている。その符号が1スパイクあたりの情報を最大にするよう調整されているというのは、十分な根拠のある仮説である。
\end{proposition}
\end{keyblock}

入力も、他のすべてと同じく非同期に働く。各センサは伝えることがあるときにスパイクを出し、感覚ごとに遅延も違う。音は数ミリ秒後に、光は数十ミリ秒後に届く。共通のクロックがないのに、マシンがそれらをどう一つの世界像にまとめるのか。これは第\ref{sec:synthesis}章の時間的協調の課題の一つである。

\section{問題}

\begin{exercise}[\lvlA]
\label{ex:11:1}
\emph{光子をどれだけ蓄積するか。} 薄暗がりでセンサに毎秒 $100$ 個の光子が入る。増分はゆらぎ~\eqref{eq:photonnoise}の $3$ 倍あれば検出できる。(a)~1個のセンサが明るさの $10\,\%$ の変化を検出するのにどれだけの時間が要るか。(b)~$0.2\,\text{s}$ で間に合わせるには何個のセンサを足し合わせればよいか。空間分解能は何分の1に落ちるか。
\end{exercise}

\begin{exercise}[\lvlA]
\label{ex:11:2}
\emph{1スパイクに何ビットあるか。} (a)~目は $10^6$ 個の出力ニューロンで、毎秒 $3$--$5$ スパイクの発火率で $\sim10^7$ bit/s を送る。$\Delta t=1\ms$ での容量~\eqref{eq:spikeentropy}のどれだけを使っているか。(b)~ある匂いが $\approx400$ 種類の受容体のうち $20$ 種類を活性化する。この組は何ビットを運ぶか。二つのランダムな匂いに共通する種類は平均で何個か。
\end{exercise}

\begin{exercise}[\lvlB]
\label{ex:11:3}
\emph{1本の曲線~\eqref{eq:nakarushton}で何ができるか。} (a)~応答が $r_{\max}$ の $10\,\%$ から $90\,\%$ まで上がるのは、明るさのどの範囲か。それは何桁か。(b)~明るさの10桁をカバーするには、ずれていく曲線の位置がいくつ要るか。音の大きさの12桁ではどうか。
\end{exercise}

\begin{exercise}[\lvlB]
\label{ex:11:4}
\emph{耳の時間符号の限界。} 両耳への音の到達時間差は最大 $0.22\,\text{m}/343\,\text{m/s}$ である。(a)~スパイクが純音と位相をそろえて出るとき、どの周波数まで方向を一意に決められるか。(b)~スパイクの時刻は $0.5\ms$ ゆらぐが、$20$~\textmu s を区別する必要がある。毎秒 $100$ スパイクのニューロンを $50\ms$ の間に何個平均すればよいか。
\end{exercise}

\begin{exercise}[\lvlB]
\label{ex:11:5}
\emph{目のケーブルに載せるイベントカメラ。} $10^6$ 画素、出力 $10^7$ bit/s のカメラが、画素の $\ln I$ が $\ln1.2$ だけ変化したときにイベント（アドレスと符号）を送る。長さ $500$ 画素、明るさの段差 $4$ 倍のエッジは、最大どれだけの速さで動けるか。同じ出力で、1画素 $8$ ビットの通常のカメラなら毎秒何フレームになるか。
\end{exercise}

\begin{exercise}[\lvlB]
\label{ex:11:6}
\emph{シミュレーション：ゲイン調整。} センサ~\eqref{eq:nakarushton}が $\sigma$ を $\tau=1\,\text{s}$ で、対数で $\tau\,\dd\ln\sigma/\dd t=\ln I-\ln\sigma$、または線形に $\tau\,\dd\sigma/\dd t=I-\sigma$ と調整する。背景は $1\to100\to1$ と跳ぶ。それぞれの法則で、上への跳びと下への跳びの後、$+10\,\%$ のプローブへの応答が順応後の値の半分まで回復するのに何秒かかるか。
\end{exercise}

\begin{exercise}[\lvlC]
\label{ex:11:7}
\emph{曲線はどんな世界に合わせてあるか。} 出力のノイズが小さく一定のとき、明るさについての情報が最大になるのは $r(I)=r_{\max}\Phi_I(I)$ のときである。ここで $\Phi_I$ は明るさの分布関数である。(a)~曲線~\eqref{eq:nakarushton}が最適となるのはどんな分布か。その $\log_{10}I$ の広がりはどれほどか。(b)~出力のノイズがポアソン型のとき、最適な曲線はどれか。
\end{exercise}

\chapter{画像と動画の認識}
\chaptermark{認識}
\label{sec:recognition}

\section{課題：画素が違っても同じもの}

第\ref{sec:sensors}章では、画像をマシンの入力まで運んだ。目の約100万個の出力ニューロンと、主にエッジと変化を伝える圧縮された流れである。だがセンサと行動の間には、これまで触れずにきた段階がある。絵の中のネコは画素の集まりではない。半フレームずらし、回転させ、遠ざけ、照明を半分消せば、ほとんどすべての画素が変わる。それでも答えは「ネコ」のままでなければならない。エンジニアはこれを\emph{不変性}と呼ぶ。分類器の答えは、平行移動、拡大縮小、回転、照明によって変わってはならない。

難しさは簡単な実験でよく分かる。$24$ 点からなる1次元の「網膜」と、どこにでも置ける5点ずつの二つの図形を考える。1か所で記憶したテンプレートと入力を比べる分類器の正答率は $49$--$52\%$ で、コイン投げと変わらない。ずれがあると、画素ごとの比較は完全に破綻する。別の構成が必要である。

\section{段のパイプライン}

脳がこれをどれほど速く行うかは、第\ref{sec:code}章ですでに知っている。一瞬だけ示された絵の中の動物はおよそ $150\ms$ で認識され、信号は十数段を1段あたり $10$--$20\ms$ で通過する~\cite{Thorpe1996}。各段でニューロンが出せるのはスパイク0個か1個である。つまり速い認識とは、長い熟考も繰り返しもない、パイプラインの1回の通過である。問題は各段が何をするかである。

視覚の最初の段での測定~\cite{HubelWiesel1962}が示した答えは、交互に現れる二つの演算からなる（図~\ref{fig:conveyor}）。

\emph{選択性。} $S$ 段のニューロンは入力の小さな領域を見て、そこに特定の部品の組み合わせ（このエッジと、その隣のあのエッジ）があるときだけ発火する。これは部品についての AND であり、どの部品一つが与える入力よりも閾値が高い、第\ref{sec:glutcomputer}章のしきいニューロンである。

\emph{不変性。} $C$ 段のニューロンは、同じ組み合わせを異なる場所で探す多数の $S$ ニューロンの出力を集め、それがどこか一か所でも見つかれば発火する。これは位置についての OR、より正確には最大値の選択である。

1次元モデルでは、二つの演算は次のように書ける：
\begin{empheq}[box=\keybox]{equation}
s_{f,p} \;=\; \Bigl[\sum_i t_{f,i}\,x_{p+i} - \theta\Bigr]_+ ,
\qquad
c_f \;=\; \max_p\, s_{f,p} ,
\label{eq:scstage}
\end{empheq}
ここで $x$ は入力、$t_{f,i}$ は特徴 $f$ のテンプレート、$p$ は位置、$\theta$ は閾値である。第1式は応答を選択的にし、第2式は位置に依存しないようにする。多数の入力の最大値は、第\ref{sec:gaba}章の道具で得られる。相互抑制は最も強い入力を残し（命題~\ref{prop:wta}）、全体の活動による除算は明るさの違いによる応答の差をならす~\cite{Carandini2012}。

\begin{figure}[t]
\centering
\begin{tikzpicture}[>=Stealth,
  px/.style={draw,minimum size=0.36cm,inner sep=0pt},
  on/.style={px,fill=blue!55!black},
  snode/.style={circle,draw,very thick,minimum size=0.62cm,inner sep=0pt,font=\scriptsize,fill=blue!6},
  cnode/.style={circle,draw,very thick,minimum size=0.8cm,inner sep=0pt,font=\scriptsize,fill=red!10}]
% two inputs: the same shape at two positions
\foreach \row/\sh/\lab in {0/1/{フレーム1},1/7/{フレーム2}}{
  \begin{scope}[yshift=-\row*1.0cm]
  \node[left,font=\scriptsize] at (-0.25,0) {\lab};
  \foreach \i in {0,...,13}{\node[px] at (\i*0.4,0) {};}
  \foreach \d in {0,1,3,4}{\node[on] at ({(\sh+\d)*0.4},0) {};}
  \end{scope}}
% S stage: detectors of the shape at several positions
\foreach \k in {0,...,5}{
  \node[snode] (S\k) at ({0.8+0.8*\k},1.7) {$S$};}
\node[font=\scriptsize,left] at (-0.25,1.7) {$S$ 段：部品の AND};
\draw[->,thick,blue!60!black] (1.2,0.25) -- (S1.south);
\draw[->,thick,orange!85!black] (3.6,0.25) -- (S4.south);
\node[font=\scriptsize,blue!60!black,anchor=east] at (1.25,0.85) {フレーム1};
\node[font=\scriptsize,orange!85!black,anchor=west] at (3.9,0.85) {フレーム2};
% C stage
\node[cnode] (C) at (2.8,3.25) {$C$};
\node[font=\scriptsize,left] at (-0.25,3.25) {$C$ 段：位置の OR};
\foreach \k in {0,...,5}{\draw[->,gray!60!black] (S\k.north) -- (C);}
\draw[->,very thick,red!70!black] (C.north) -- ++(0,0.6) node[above,font=\scriptsize,black] {「図形あり」：どちらのフレームでも};
% next stages
\draw[decorate,decoration={brace,amplitude=5pt},gray!60!black] (6.3,1.4) -- (6.3,-1.3);
\node[right,font=\scriptsize,align=left] at (6.5,0.05) {同じ物体、\\違う画素};
\draw[->,thick,densely dashed,gray!60!black] (3.6,3.25) -- (5.9,3.25) node[right,font=\scriptsize,black,align=left] {次の $S$, $C$ の対：\\より大きく、より複雑、\\より不変};
\end{tikzpicture}
\caption{パイプラインの一対の段。$S$ ニューロンは同じ部品の組み合わせを異なる場所で探し、$C$ ニューロンはそれがどこかで見つかれば応答する~\eqref{eq:scstage}。フレーム1とフレーム2の図形は異なる $S$ ニューロンを興奮させるが、同じ $C$ を興奮させる。次の段の対は、同じことをより大きく複雑な組み合わせについて繰り返す。}
\label{fig:conveyor}
\end{figure}

同じ1次元モデルで、一対の段~\eqref{eq:scstage}は図形をどの位置でも誤りなく認識する。入力の各点を確率 $0.1$ でランダムに反転させると正答率は $86\%$、確率 $0.2$ では $70\%$ である。コイン投げは相変わらず半分である。このような対をいくつか積み重ねてみよう。最初の対はエッジを、次の対はエッジの組み合わせを、さらに上は物体の部分を、最上段は物体全体を探す。対を重ねるごとに組み合わせは大きくなり、答えは物体がどこにどう置かれているかにますます依存しなくなる~\cite{Riesenhuber1999}。

\begin{keyblock}
\begin{proposition}[選択性と不変性のパイプライン]
\label{prop:conveyor}
速い認識は十数段の1回の通過である。各段の対で、部品の AND~\eqref{eq:scstage}が応答を選択的にし、位置の OR がずれに依存しないようにする。この二つの演算を交互に繰り返すと、物体を区別しつつ、それがどこにどう見えるかにほとんど依存しない答えが得られる。最初の段の仕組みと通過の速さは測定されている。連鎖全体がこの二つの演算に帰着するというのは単純化である。
\end{proposition}
\end{keyblock}

この仕組みに見覚えがあるなら、それで正しい。まさにこの交互の繰り返し、つまりすべての場所で特徴を探し、次に近くの場所どうしをまとめることを、エンジニアは人工の\emph{畳み込み}ネットワークに移した~\cite{Fukushima1980}。最近は逆向きの流れもある。物体認識を学習した畳み込みネットワークが、視覚の上位段のニューロン応答について知られている最良のモデルであることが分かった~\cite{Yamins2014}。最上段の結果は簡単に読み出せる。最終段の約100個のニューロンの発火率の合計から、線形の判定器が提示後0.1秒で物体を認識する~\cite{Hung2005}。これは第\ref{sec:code}章の集団のレート符号である。

\section{いない教師：動画から学ぶ}

畳み込みネットワークは、ラベルの付いた数百万枚の画像で学習する。脳にラベルを与えてくれる者はほとんどいない。子どもは何時間も物を見るが、「カップ」という言葉を聞くのはたまにである。では $C$ 段は、どの $S$ ニューロンをまとめるべきかをどうやって知るのか。「ここにあるこの図形」と「あそこにあるこの図形」をまとめるには、それが同じ図形だとすでに知っていなければならないはずである。

ヒントは動画そのものにある。世界は連続的に変わる。視野の中の物体は、ずれ、回転し、照明が変わる間も同じ物体であり続け、視線が別の物体に移るのはときどきにすぎない。\emph{ゆっくり}変わるものはおそらく物体に属し、速く変わるものはその位置と見え方に属する~\cite{WiskottSejnowski2002}。したがって $C$ ニューロンには、第\ref{sec:dofa}章の適格度トレース付きヘッブ則で十分である。同時に来たものだけでなく、続けて来たものを結びつければよい~\cite{Foldiak1991}：
\begin{empheq}[box=\keybox]{equation}
\tau_{\text{tr}}\,\frac{\dd\bar y_k}{\dd t} \;=\; -\bar y_k + y_k ,
\qquad
\frac{\dd\mathbf w_k}{\dd t} \;=\; \eta\,\bar y_k\,\bigl(\mathbf x - \mathbf w_k\bigr) .
\label{eq:traceinvariance}
\end{empheq}
ここで $y_k$ は抑制による競合に勝った $C$ ニューロン $k$ の活動、$\bar y_k$ はそのトレースで、規則~\eqref{eq:threefactor}と同じRC回路である。$\mathbf x$ は $S$ ニューロンの出力である。物体の最初の見え方で勝ったニューロンは、2番目の見え方が来たときにもまだそれを覚えていて、それも自分に引き寄せる。

\begin{figure}[t]
\centering
\begin{tikzpicture}[>=Stealth,x=1.0cm,y=1.0cm]
\foreach \i/\lab in {0/1,1/2,2/3,3/4}{
  \node[draw,thick,fill=blue!8,minimum width=0.9cm,minimum height=0.55cm,font=\scriptsize] at (\i+0.5,2.2) {$A$@\lab};}
\foreach \i/\lab in {0/4,1/3,2/2,3/1}{
  \node[draw,thick,fill=orange!12,minimum width=0.9cm,minimum height=0.55cm,font=\scriptsize] at (\i+6.5,2.2) {$B$@\lab};}
\node[font=\scriptsize,gray!40!black] at (5.0,2.2) {休止};
\draw[->,gray!70!black] (0,0) -- (10.4,0) node[below left,font=\scriptsize,black] {時間};
% traces
\draw[thick,blue!60!black] (0,0.05) .. controls (0.4,1.2) and (0.8,1.3) .. (1.2,1.35) -- (4,1.4) .. controls (4.6,0.5) and (5.2,0.15) .. (6,0.08) -- (10,0.05);
\draw[thick,orange!85!black] (0,0.05) -- (6,0.05) .. controls (6.4,1.2) and (6.8,1.3) .. (7.2,1.35) -- (10,1.4);
\node[font=\scriptsize,blue!60!black,anchor=west] at (1.4,1.65) {トレース $\bar y_1$：$A$ の通過中ずっと保たれる};
\node[font=\scriptsize,orange!85!black,anchor=west] at (7.2,1.65) {トレース $\bar y_2$};
\draw[decorate,decoration={brace,amplitude=4pt},gray!60!black] (4.0,2.6) -- (0,2.6);
\draw[decorate,decoration={brace,amplitude=4pt,mirror},gray!60!black] (0,-0.35) -- (4,-0.35) node[midway,below=4pt,font=\scriptsize,black] {$A$ のすべての見え方がニューロン1に結びつく};
\end{tikzpicture}
\caption{動画が不変性をどう教えるか（模式図）。物体 $A$ が四つの位置を通過し（$A$@1 は位置1の $A$）、休止をはさんで物体 $B$ が続く。$A$ の最初の見え方で勝ったニューロンのトレース~\eqref{eq:traceinvariance}は通過の間ずっと保たれ、$A$ のすべての見え方が一つのニューロンに結びつく。休止がトレースを消し、$B$ は別のニューロンに割り当てられる。}
\label{fig:tracevideo}
\end{figure}

これをモデルで確かめた（図~\ref{fig:tracevideo}）。二つの $C$ ニューロンが $16$ 個の $S$ ニューロン（二つの図形×八つの位置）からの入力を奪い合う。図形はすべての位置を1位置あたり $50\ms$ で通過し、次に $0.3\,\text{s}$ の休止、そして次のランダムな図形が来る。ラベルは一切ない。トレースが短い $\tau_{\text{tr}}=5\ms$ のとき、ニューロンは入力を\emph{位置}で分け合い、答えが図形と一致する割合はコイン投げと変わらない。10回の試行の平均で $52\%$ である。トレースが一つの位置の提示時間と同程度、$\tau_{\text{tr}}\approx20\ms$ 以上になると、各ニューロンは自分の図形の\emph{すべての}位置を集める（$20$, $50$, $200\ms$ では10回の試行すべてで）。不変性が動画だけから学習されたのである。物体間の休止は重要である。休止がないと、トレースが次の物体に流れ込み、二つを混同させる。

同じことは実験でも見られる。眼があるところから別のところへ跳ぶ間に物体をこっそり差し替えると、上位段のニューロンは1〜2時間のうちに、差し替えた物体に元の物体と同じように応答し始める。続けて来たものを結びつけるのである~\cite{LiDiCarlo2008}。脳の中の不変性は、確かに時間の連続性から学習される。視覚の学習全体がこの規則でどこまで説明できるかは、仮説である。

\section{動画とは運動である}

動画は学習のためのフレームの流れであるだけでなく、それ自体が課題でもある。何が、どちらへ、どれだけ速く動いているか。最初の一歩はおなじみの、第\ref{sec:gaba}章の方向検出器である（図~\ref{fig:direction}）。そこでは遅れて来る抑制が一方向の運動を打ち消す。このような検出器が視野全体に並び、その先は形の特徴と同じ扱いを受ける。同じ方向の検出器を異なる場所から多数集める段は、大きな物体の運動に、それがどこにあっても応答する。これは同じ AND と OR の対~\eqref{eq:scstage}であり、特徴が「エッジ」ではなく「時間 $d$ の間にずれたエッジ」であるだけである。

動画はまた予測可能でもある。次のフレームは前のフレームとほとんど同じである。\emph{予測符号化}の仮説では、上位の段は下位の段に予測を送り、上へ送られるのは主に誤差、つまり予測が捉えきれなかった部分である~\cite{RaoBallard1999}。これは命題~\ref{prop:economy}と同じスパイクの節約である。すでに分かっていることではなく、ニュースに支払う。個々の観察はこの仮説と整合するが、パイプラインの一般的な仕組みとしては証明されていない。

\section{脳と畳み込みネットワーク}

類似はどこまで及ぶのか。一つの物体の速い認識については、確かに深い~\cite{Yamins2014}。だが脳には、単純な畳み込みネットワークにないものもある。

\emph{フィードバック。} 脳のパイプラインは順方向だけではない。どの段にも後ろ向きと横向きの結合がある。重なりや悪い照明のある難しい画像では、上位段の応答はより遅れて形成され、その遅い応答は順方向のネットワークよりフィードバックのあるネットワークのほうがよく記述する~\cite{Kar2019}。1回の通過は易しい場合を解き、難しい場合には何周かが必要になる。

\emph{頑健性。} 人工ネットワークは、目には見えない画素の改ざんでだませる~\cite{Szegedy2014}。人に見えない変化が「パンダ」を「テナガザル」に変える~\cite{Goodfellow2015}。ヒトの視覚はこのような改ざんにはるかに強いが、無敵ではない。多くのネットワークを同時にだます画像は、短時間提示ではヒトの選択もずらす~\cite{Elsayed2018}。頑健性がなぜこれほど違うのかは未解決の問題である。候補はノイズ、全体の活動による除算、フィードバックである。

\emph{教師。} ネットワークにはラベル付きの例が数百万必要だが、脳に要るのは主にラベルなしの動画~\eqref{eq:traceinvariance}と、何が重要かについての\nt{DOPA}からのわずかなヒントである。

\emph{エネルギー。} 脳全体で約 $20$ ワット（命題~\ref{prop:economy}）、認識はその一部にすぎない。学習済みの畳み込みネットワークはGPU上で数百ワットを使う。

\section{錯視：マシンはどこでなぜ間違えるか}
\label{sec:illusions}

他人の回路を理解したいエンジニアは、ふつうでない信号を入れて、どこで間違えるかを見る。視覚については、そうした信号はずっと前に考案されている。錯視である。錯視は故障ではない。どの錯視も、ふだんの状況では正しく働くマシンの特定の機構を指し示し、特別に選んだ画像でその機構が正体を現す。

\emph{妥当な期待。} 入力はノイズを含むので、マシンはそれを世界でふつう起こることで補う。遅い運動は速い運動より多い。入力が信頼できないとき、たとえば画像のコントラストが低いとき、速さの推定は遅いほうへずれ、淡い物体は明るい物体より遅く動いて見える~\cite{Weiss2002}。明るさの錯視の多くも同じように説明される。我々が見るのは画素の明るさではなく、その明るさが過去に最もよく対応していた表面である~\cite{Purves2011}。ある人には白と金に、別の人には青と黒に見えるドレスの写真も同じ機構である。画像はあいまいで、人は無意識に想定する照明の点で分かれる~\cite{LaferSousa2015,Wallisch2017}。個人差は測定されている。ここで脳が確率論の規則に従って入力と期待を組み合わせているというのは、十分な根拠のある仮説である。

\emph{ゲイン調整。} 滝を1分間見つめてから、動かない岩を見てみよう。岩が上へ流れていく。「下」と「上」のチャネルは\eqref{eq:dualrail}のような配線の対である。ゲイン調整~\eqref{eq:nakarushton}が疲れたチャネルを弱め、対のつり合いがずれたのである。周囲が中心の見え方を変える傾きやコントラストの錯視は、第\ref{sec:gaba}章と同じく、近傍の活動による除算で説明される~\cite{Schwartz2009}。慎重さを教える例もある。格子の白い帯の交差点に見える暗い斑点は、長く近傍の単純な抑制で説明されてきたが、格子を変形した実験によって、その説明が成り立たないことが示された~\cite{SchillerCarvey2005}。

\emph{遅延の補償。} 目から上位段までの経路には数十ミリ秒かかり、動く物体はその間に移動してしまう。その横で静止した点が光ると、両者は一致していたのに、物体は閃光より前にあるように見える~\cite{Nijhawan1994}。一つの説明では、マシンは遅延を補償するために運動を前に延長する。第\ref{sec:muscles}章と同じく、フィードバックが遅れるなら予測しなければならない。この錯視には説明がいくつかあり、論争は決着していない。

\emph{タイミング符号の誤り。} 色の急な移り変わりのある輪からなる静止画が、回転しているように見える。コントラストの高い部分は淡い部分より速く処理され、遅延の差~\eqref{eq:latency}を方向検出器が運動として読む~\cite{Conway2005}。錯視を起動するのは、目の小さな不随意の跳躍と瞬きである。跳躍のたびに入力が更新され、検出器は再び「運動」を見る~\cite{OteroMillan2012}。これは第\ref{sec:code}章のタイミング符号の読み違いである。

\emph{一つの中の二つの絵。} ある面をこちらに向けたり別の面を向けたりする立方体の枠や、左右の目に違う絵を見せた場合、知覚は数秒ごとに切り替わる。最良のモデルは、二つの解釈の相互抑制、遅い疲労、ノイズである~\cite{MorenoBote2007}。つまり第\ref{sec:muscles}章で歩行のリズムを生んだのと同じ仕組み~\eqref{eq:halfcentre}である。切り替わりのタイミングはノイズと疲労が決める。モデル~\cite{MorenoBote2007}では主役はノイズで、疲労は補助にすぎない。

では人工ネットワークはどうか。動画の次のフレームを予測するよう学習したネットワークは、同じ静止した輪に回転を見て、対照画像には見ない~\cite{Watanabe2018}。画像のノイズ除去と鮮鋭化を学習したネットワークは、ヒトと同じ色と明るさの錯視にかかる~\cite{GomezVilla2020}。つまり錯視の一部は、神経組織の特性ではなく、マシンが学習する課題そのものから生じる。どの錯視が脳だけに固有なのかは、どんな視覚モデルにとってもよい試験になる。

\begin{keyblock}
\begin{proposition}[機構の痕跡としての錯視]
\label{prop:illusions}
錯視は、ふだんの世界で役に立つ機構がふつうでない入力を受けたときに生じる。期待が信頼できない入力に勝ち、ゲイン調整がチャネルの対をずらし、遅延の補償が物体を前へ押し出し、遅延の差が運動として読まれ、ノイズと疲労をともなう相互抑制が切り替わる。どの錯視も、自分の機構を指し示す試験信号である。錯視そのものは測定されている。その説明は、十分な根拠のあるものから論争中のものまである。
\end{proposition}
\end{keyblock}

\section{まとめ}

\begin{table}[t]
\centering
\footnotesize
\setlength{\tabcolsep}{4pt}
\renewcommand{\arraystretch}{1.15}
\begin{tabular}{>{\raggedright}p{3.0cm}>{\raggedright}p{4.6cm}>{\raggedright\arraybackslash}p{5.2cm}}
\toprule
マシンは何をするか & どうやって & 分かっていること \\
\midrule
$150\ms$ で認識する & 十数段の1回の通過 & 速さは測定済み \\
応答を選択的にする & 部品の AND、$S$ 段~\eqref{eq:scstage} & 最初の段で測定済み \\
応答を不変にする & 位置の OR、$C$ 段；抑制と除算 & 最初の段で測定済み；上位段については単純化 \\
答えを出す & 最終段の約100個のニューロンの発火率、線形の読み出し & 測定済み \\
ラベルなしで学ぶ & 連続した動画上のトレース付きヘッブ則~\eqref{eq:traceinvariance} & 物体の差し替えが応答を変える（測定済み）；主要な規則としては仮説 \\
運動を見る & 方向検出器、次に同じ AND／OR の対 & 測定済み \\
予測できるものを節約する & 上へは予測誤差が送られる & 仮説 \\
難しい場合を解く & フィードバック、何周か & 遅い応答とフィードバックの役割は測定済み \\
予測どおりに間違える & 錯視：期待、ゲイン調整、遅延、ノイズと疲労をともなう相互抑制 & 錯視は測定済み；説明は根拠のあるものから論争中のものまで \\
\bottomrule
\end{tabular}
\caption{マシンは画像と動画をどう認識するか。}
\label{tab:recognition}
\end{table}

認識は、すでに手元にある部品で組み立てられている（表~\ref{tab:recognition}）。閾値が部品の AND を、相互抑制が位置の OR を、除算が明るさへの非依存を、トレース付きヘッブ則がいない教師の代わりを与える。エンジニアは視覚から選択性と不変性の交互の繰り返しを借り、それをもとに畳み込みネットワークを作った。逆方向には、脳はまだ肝心なもの、つまりラベルなしの動画からほとんどエネルギーを使わずに学ぶ能力を明け渡していない。

\section{問題}

\begin{exercise}[\lvlA]
\label{ex:12:1}
\emph{1段での発火率。} パイプラインの1段は $15\ms$ 続き、ニューロンは毎秒 $20$ スパイクで発火し、ノイズは $c=0.1$~\eqref{eq:correlated}で相関している。100個のニューロンの集団の発火率の相対誤差と、集団をいくら大きくしたときのその極限はいくらか。1段で区別できる発火率のレベルはいくつか。
\end{exercise}

\begin{exercise}[\lvlA]
\label{ex:12:2}
\emph{1回の認識のエネルギー。} $150\ms$ の1回の通過で、$10^7$ 個のニューロンからなる10段がそれぞれ高々1スパイクを出し、1スパイクのコストは $10^{-10}$~J である。(a)~この $150\ms$ の間の脳全体のエネルギーのうち、認識はどれだけの割合を占めるか。(b)~画像を $10\ms$ で認識する $300$~W のアクセラレータは何倍多く使うか。
\end{exercise}

\begin{exercise}[\lvlB]
\label{ex:12:3}
\emph{$S$ 段の閾値。} $24$ 点の「網膜」、図形 $A=11011$ と $B=10101$、テンプレート~\eqref{eq:scstage}は図形の1の位置で $+1$、0の位置で $-1$ とする。(a)~$S$ ニューロンが1点の反転を許容しつつ、他方の図形には沈黙する閾値の範囲はどこか。(b)~$100\times100$ の画像上で $5\times5$ の特徴10個に要る $S$ ニューロンの数はいくつか。
\end{exercise}

\begin{exercise}[\lvlB]
\label{ex:12:4}
\emph{二つの絵の一方はどれだけ続くか。} 切り替えの仕組み~\eqref{eq:halfcentre}で $\tau\ll\tau_a$ とし、グループ1が勝っている間は $r_2=0$, $r_1=I-a_1$ とする。(a)~$I=1$, $\beta=2$, $b=2$, $\tau_a=0.4\,\text{s}$ で勝利の持続時間を求めよ。(b)~絵が3秒ごとに切り替わるのはどの $\tau_a$ か。
\end{exercise}

\begin{exercise}[\lvlB]
\label{ex:12:5}
\emph{シミュレーション：不変性の代償。} \texttt{models/\allowbreak recognition\_models.py} の関数 \texttt{two\_stage}（$4000$ 試行）を使い、$N=24$ で段の対 $S$, $C$ が4回に1回間違える点の反転確率 $p$ を求めよ。$p=0.1$ で、$N=8$ から $N=192$ まで正答率はどう変わるか。
\end{exercise}

\begin{exercise}[\lvlB]
\label{ex:12:6}
\emph{シミュレーション：トレースの窓。} \texttt{models/\allowbreak recognition\_models.py} の関数 \texttt{trace\_learning} を、休止を $0.3\,\text{s}$ として10回実行し、$5\ms$ から $2\,\text{s}$ までのどの $\tau_{\text{tr}}$ ですべての実行が誤りなく不変性を学習するかを求めよ。この窓の上限はどこから来るか。
\end{exercise}

\begin{exercise}[\lvlB]
\label{ex:12:7}
\emph{どこででも運動を。} 間隔 $0.1^\circ$ の点の列の隣り合う点に、第\ref{sec:gaba}章の方向検出器を置く。遅延 $d$ の $A$ からの抑制は、$B$ からの興奮とのずれが $5\ms$ 以内なら興奮を打ち消す。その上に $C$ 段がある。毎秒 $5$ から $20^\circ$ の逆向きの運動で $C$ を沈黙させるには、どの遅延が必要か。
\end{exercise}

\begin{exercise}[\lvlC]
\label{ex:12:8}
\emph{画素の改ざん。} きれいな図形について \texttt{two\_stage} モデル（$N=24$）の答えを変えるには、何点反転すればよいか。同じくずれに不変な分類器「入力の1が $3.5$ 個より多ければ $A$」ではどうか。$p=0.1$ でのその正答率は、段の対の $0.86$ と比べてどうか。
\end{exercise}

\chapter{マシンの出力：脳はどう筋肉を制御するか}
\chaptermark{マシンの出力}
\label{sec:muscles}

\section{速い出力}

マシンが世界の中で素早く行うことは、すべて筋肉で行われる。言葉、視線、一歩、微笑み、どれも筋肉の収縮である。もう一つの遅い出力である体の化学は、第\ref{sec:chemoutput}章で扱う。図~\ref{fig:machine}で出力は「筋肉」という矢印だった。今度はその矢印の先に何があるかを見よう。

連鎖の最後の環は\nt{ACET}ニューロンであり、表~\ref{tab:types}で「筋肉への出力」と書かれていたものである。各筋線維はちょうど1個のこうしたニューロンから入力を受け、1個のニューロンは数百から2000本ほどの線維からなるグループを制御する~\cite{Feinstein1955mu}。細かな調整を要する筋肉ではこの単位は小さく、脚の大きな筋肉では大きい。ニューロンとその線維を合わせて\emph{運動単位}と呼ぶ。これは脳が個別にオンにできる筋肉の最小の部分である。

筋肉上のニューロンのシナプスは、第\ref{sec:code}章の皮質のシナプスとはまったく違う。あちらではスパイクの大部分が失われた。こちらではスパイクのたびに、線維を発火させるのに必要な量の数倍のアセチルコリンが放出される~\cite{WoodSlater2001}。これは\emph{安全率}を持つシナプスであり、ニューロンのスパイク1個がちょうど1回の線維の収縮になる。マシンの内部では信頼できない結合を許容し、誤りを冗長性で直すことができる。出力では誤りは動きそのものを損なうので、信頼性はハードウェアに直接組み込まれている。

\section{力：発火率と単位の数}

1個のスパイクは単収縮（トウィッチ）と呼ばれる1回の収縮を生む。力は数十ミリ秒で立ち上がり、また下がる。そのよい近似は次の式である~\cite{Fuglevand1993}：
\begin{equation}
h(t) \;=\; F_1\,\frac{t}{T_c}\,e^{\,1-t/T_c},
\label{eq:twitch}
\end{equation}
ここで $T_c$ は立ち上がり時間で $50\ms$ 程度、$F_1$ は単収縮1回の力である。スパイクがより頻繁に来ると、単収縮は足し合わされる：
\begin{empheq}[box=\keybox]{equation}
F(t) \;=\; \sum_k h\bigl(t-t_k\bigr) .
\label{eq:forcesum}
\end{empheq}
これは第\ref{sec:serocomputer}章でスパイクを濃度に変えたのと同じローパスフィルタであり、違うのは出力が濃度ではなく力であることだけである。筋肉はスパイクから力へのD/A変換器である（図~\ref{fig:tetanus}）。我々のモデルでは、毎秒5スパイクのとき単収縮はほとんど重ならず、力は $0.2F_1$ から $1.1F_1$ まで跳ねる。毎秒15スパイクではすでに $1.8F_1$ と $2.2F_1$ の間に保たれ、毎秒40スパイクではほぼ平らな約 $5.4F_1$ になる。実際の筋肉は高頻度のスパイクで飽和するので、線形の和~\eqref{eq:forcesum}が成り立つのは毎秒数十スパイクまでである。

\begin{figure}[t]
\centering
\begin{tikzpicture}[>=Stealth,x=20cm,y=0.55cm]
\draw[->,gray!70!black] (0,0) -- (0.66,0) node[right,font=\small,black] {$t$, s};
\draw[->,gray!70!black] (0,0) -- (0,6.4) node[left,font=\small,black] {$F/F_1$};
\foreach \t in {0,0.2,0.4,0.6}{\draw[gray!60!black] (\t,-0.1) -- (\t,0.1); \node[below,font=\scriptsize] at (\t,0) {\t};}
\foreach \f in {1,3,5}{\draw[gray!60!black] (-0.004,\f) -- (0.004,\f); \node[left,font=\scriptsize] at (0,\f) {\f};}
\draw[thick,blue!60!black] plot file {figures/muscle_twitch_5.dat};
\draw[thick,green!45!black] plot file {figures/muscle_twitch_15.dat};
\draw[thick,red!70!black] plot file {figures/muscle_twitch_40.dat};
\node[font=\scriptsize,blue!60!black,anchor=west] at (0.605,0.9) {5 スパイク/s};
\node[font=\scriptsize,green!45!black,anchor=west] at (0.605,2.1) {15 スパイク/s};
\node[font=\scriptsize,red!70!black,anchor=west] at (0.605,5.4) {40 スパイク/s};
\end{tikzpicture}
\caption{D/A変換器としての筋肉：1個の運動単位の規則的なスパイクに対する力~\eqref{eq:forcesum}、$T_c=50\ms$。まばらなスパイクは個々の単収縮を生み、頻繁なスパイクは融合して平らな力になる。図~\ref{fig:lowpass}で頻繁なスパイクが融合して平らな濃度になったのと同じである。}
\label{fig:tetanus}
\end{figure}

脳には力のつまみが二つある。各単位のスパイクの発火率と、オンになっている単位の数である。そして2番目のつまみは実に巧妙にできている。筋肉の中の単位はそれぞれ異なり、最も弱いものは最も強いものの100分の1である。より大きな力が必要になると、単位は\emph{大きさの順に}、小さいものから大きいものへとオンになる~\cite{Henneman1957}。この順序を計算する者はいない。小さなニューロンは同じ共通入力でより興奮しやすいだけであり、オンになる順序はハードウェア自体が決めている。

これで何が得られるか。順番に並べた各単位が一つ前より $q$ 倍強いとし、$q$ は1より少し大きいとする。オンになっている全単位の力は等比数列の和であり、そのほぼすべてが最後の数単位による。すると次にオンになる単位が加える力は
\begin{empheq}[box=\keybox]{equation}
\frac{\Delta F}{F} \;\approx\; q-1 \;=\; \text{const} ,
\label{eq:sizeprinciple}
\end{empheq}
つまり力の刻みは力そのものに比例する。弱い力は細かく、強い力は粗く刻まれる。これは\emph{浮動小数点数}である。指数はオンになっている単位の数で、仮数はそのスパイクの発火率で決まる。第\ref{sec:sensors}章のセンサと同じ対数目盛である。マシンの入力と出力は、世界を相対的な割合で測る。

浮動小数点には裏面がある。力のばらつきも力そのものに比例して大きくなる。強い動きは弱い動きより必然的に不正確になる。有力な仮説の一つによれば、まさにこの信号依存のノイズが、我々の動きがなめらかな理由を説明する。なめらかな指令が終点でのばらつきを最小にするのである~\cite{HarrisWolpert1998}。

\section{引けるが押せない：関節あたり2本の配線}

筋肉は引くことしかできない。\nt{GLUT}ニューロンが負の信号を出せないのと同じく、押すことはできない。解決法は第\ref{sec:glutcomputer}章でおなじみの\emph{2本の配線}である。各関節は、逆向きに引く一対の筋肉が受け持つ（図~\ref{fig:antagonists}）。力を $F_1$, $F_2$、モーメントアームを $\ell$ とすると、関節のトルクと剛性は
\begin{empheq}[box=\keybox]{equation}
M \;=\; \ell\,(F_1-F_2), \qquad K \;\approx\; \kappa_m\,(F_1+F_2),
\label{eq:antagonists}
\end{empheq}
となる。ここで $\kappa_m$ は筋肉の力とその剛性とを結ぶ係数である。緊張した筋肉は弛緩した筋肉よりばねとして硬い。力の差がトルクを、和が剛性を決める~\cite{Hogan1984}。2本の配線で、脳は二つの独立な量を制御する。関節をどちらへ回すかと、どれだけ硬く保つかである。こぼさずにカップを持つとき、我々は両方の筋肉を同時に緊張させる。トルクは同じで剛性が高くなり、不意の衝撃で手がぶれにくくなる。

\begin{figure}[t]
\centering
\begin{tikzpicture}[>=Stealth,
  nrn/.style={circle,draw,very thick,minimum size=0.95cm,inner sep=0pt,font=\scriptsize},
  inh/.style={-{Bar[width=7pt]},thick,red!70!black}]
% the joint: a bar on a pivot, two springs pulling down on either side
\fill[gray!30] (4.6,-0.25) -- (5.4,-0.25) -- (5,0.15) -- cycle;
\draw[line width=3pt,gray!50!black] (2.4,0.2) -- (7.6,0.2);
\fill[black] (5,0.2) circle (0.08);
\draw[decorate,decoration={coil,segment length=5pt,amplitude=4pt},very thick,blue!60!black] (3.0,0.2) -- (3.0,-1.6);
\draw[decorate,decoration={coil,segment length=5pt,amplitude=4pt},very thick,orange!85!black] (7.0,0.2) -- (7.0,-1.6);
\draw[very thick] (2.5,-1.6) -- (3.5,-1.6); \draw[very thick] (6.5,-1.6) -- (7.5,-1.6);
\node[font=\small,blue!60!black] at (2.35,-0.75) {$F_1$};
\node[font=\small,orange!85!black] at (7.65,-0.75) {$F_2$};
\draw[->,thick] (5.9,0.75) arc (60:120:1.8) node[above,font=\scriptsize,pos=0.5] {$M=\ell(F_1-F_2)$};
% motor neurons and reflex circuit
\node[nrn,fill=green!12] (M1) at (1.6,-3.0) {\nt{ACET}};
\node[nrn,fill=green!12] (M2) at (8.4,-3.0) {\nt{ACET}};
\node[nrn,fill=red!10] (G) at (5.0,-3.0) {\nt{GLYC}};
\node[nrn,fill=blue!6] (S) at (1.6,-5.0) {\nt{GLUT}};
\draw[->,very thick] (M1) -- (3.0,-1.7);
\draw[->,very thick] (M2) -- (7.0,-1.7);
\draw[->,thick] (S) -- (M1);
\draw[->,thick] (S) -- (G);
\draw[inh] (G) -- (M2);
\draw[->,thick,densely dashed,gray!60!black] (3.15,-1.2) to[bend left=25] node[pos=0.35,right,font=\scriptsize,black] {伸張} (S.north east);
\draw[->,thick,blue!60!black] (-0.3,-3.0) node[left,font=\scriptsize,black,align=right] {脳からの\\指令} -- (M1);
\draw[->,thick,blue!60!black] (10.3,-3.0) node[right,font=\scriptsize,black,align=left] {脳からの\\指令} -- (M2);
\end{tikzpicture}
\caption{一つの関節に二つの筋肉と、最も速いフィードバックループ。\nt{ACET}ニューロンは脳からの指令を受け、関節を逆向きに引く。力の差が関節を回し、和が関節を硬くする。左の筋肉の伸張センサ（\nt{GLUT}）はその筋肉自身の\nt{ACET}ニューロンを興奮させ、\nt{GLYC}の介在ニューロンを通じて拮抗筋のニューロンを抑制する。第\ref{sec:gaba}章の禁止である。}
\label{fig:antagonists}
\end{figure}

\section{遅延のあるフィードバック}

筋肉はやみくもに働くわけではない。どの筋肉にも表~\ref{tab:sensors}で触れた伸張センサがあり、最も短いループは脊髄の中で直接閉じている（図~\ref{fig:antagonists}）。筋肉が不意に伸ばされると、伸張センサがその筋肉自身の\nt{ACET}ニューロンを興奮させ、筋肉は収縮して関節を元に戻す。同時に、抑制性の介在ニューロンを通じて拮抗筋がオフになり、邪魔をしないようにする。この介在ニューロンとして働くのが\nt{GLYC}である。表~\ref{tab:types}のうち自分の章を持たなかった型であり、その居場所はまさにここ、脊髄の速いループにある。

どのループにも遅延 $d$ がある。脊髄のループは腕で約 $25$--$30\ms$、皮質を通るループはその1.5〜3倍、目を通るループは $100$--$200\ms$ である。そして遅延のあるフィードバックは、命題~\ref{prop:ei}で見たように系を揺らす。$d$ 秒前の情報にもとづいて偏差 $x$ を速さ $G$ で修正する最も単純な制御器、$\dd x/\dd t=-G\,x(t-d)$ の安定条件は次のとおりである~\cite{Hayes1950}：
\begin{empheq}[box=\keybox]{equation}
G\,d \;<\; \frac{\pi}{2} ,
\label{eq:delaystable}
\end{empheq}
そして境界では系は周波数 $1/(4d)$ で振動する。これをモデルで確かめた。$d=30\ms$ のとき、$G=40$（毎秒）では偏差は減衰し、境界 $52$ をわずかに超える $G=55$（毎秒）では振動が増大する。

ここから二つの帰結が出る。第一に、ループが長いほど修正は遅くなければならない。遅延150ミリ秒の目を通して手を直すなら、約0.1秒より速くは直せない。そうしないと手が震える。第二に、脊髄ループの安定境界 $1/(4\cdot30\ms)\approx8$ 回/秒は、健康な人なら誰にでもある細かな震え、毎秒8〜12回に近い。伸張ループはこの震えの有力な原因の一つである~\cite{McAuleyMarsden2000}。

フィードバックが遅れるなら、脳はどうやって速く正確に動くのか。広く受け入れられた仮説によれば、\emph{予測する}のである。体の内部モデルを持ち、センサを待たずに指令の結果がどうなるかを計算する~\cite{WolpertGhahramani2000}。センサが必要なのは、一つ一つの動きではなくモデルを修正するためである。エンジニアならこれを、状態オブザーバ付きのフィードフォワード制御と呼ぶだろう。

\section{運動のリズム発生器}

歩行、呼吸、咀嚼はリズミカルな運動である。そのためにクロックは必要か。必要ない。脊髄と脳幹には、リズミカルな入力が何も来なくても自らリズムを生み出す小さなネットワークがある~\cite{MarderBucher2001}。最も単純なそのようなネットワークは、すでに手元にある部品で組み立てられる。図~\ref{fig:wta}のように\nt{GABA}または\nt{GLYC}の介在ニューロンを介して互いに抑制し合い、第\ref{sec:glutcomputer}章の記憶のようにゆっくり疲労する、二つの\nt{GLUT}ニューロン群である：
\begin{empheq}[box=\keybox]{equation}
\tau\,\frac{\dd r_i}{\dd t} = -r_i + \bigl[I - \beta\,r_j - a_i\bigr]_+ ,
\qquad
\tau_a\,\frac{\dd a_i}{\dd t} = -a_i + b\,r_i ,
\label{eq:halfcentre}
\end{empheq}
ここで $a_i$ は群 $i$ の疲労、$j$ はもう一方の群である。$\beta>1$ の相互抑制が勝者を選ぶ（命題~\ref{prop:wta}）。勝者は働き、疲れる。疲労が入力を食いつぶすと、すでに休んでいた敗者が前に出て、今度は自分が疲れ始める。二つの群は交代で働き、それぞれが対の一方の筋肉を制御する（図~\ref{fig:halfcentre}）。

\begin{figure}[t]
\centering
\begin{tikzpicture}[>=Stealth,x=3.2cm,y=2.4cm]
\draw[->,gray!70!black] (0,0) -- (4.2,0) node[right,font=\small,black] {$t$, s};
\draw[->,gray!70!black] (0,0) -- (0,1.2) node[left,font=\small,black] {$r$};
\foreach \t in {0,1,2,3,4}{\draw[gray!60!black] (\t,-0.03) -- (\t,0.03); \node[below,font=\scriptsize] at (\t,0) {\t};}
\draw[thick,blue!60!black] plot file {figures/halfcentre1.dat};
\draw[thick,orange!85!black] plot file {figures/halfcentre2.dat};
\node[font=\scriptsize,blue!60!black,anchor=west] at (1.6,1.28) {群1：「前へ」};
\node[font=\scriptsize,orange!85!black,anchor=west] at (2.7,1.28) {群2：「後ろへ」};
\end{tikzpicture}
\caption{モデル~\eqref{eq:halfcentre}によるリズム発生器：$I=1$, $\beta=2$, $b=2$, $\tau=20\ms$, $\tau_a=0.4\,\text{s}$。相互抑制と疲労を持つ二つの群は、入力が一定の信号であるにもかかわらず、周期 $0.84\,\text{s}$ で交代に働く。}
\label{fig:halfcentre}
\end{figure}

我々のモデルでは、一定の入力のもとで二つの群は周期 $0.84\,\text{s}$ で交代する。これは疲労の時定数 $\tau_a$ のほぼ2倍である。上からの一定の指令「歩け」が、その場でリズムに変わる。これも第\ref{sec:gaba}章の\nt{GLUT}--\nt{GABA}ループのリズムと同じく局所的なリズムである。ネットワークが働いているときだけ生じ、自分が制御する筋肉にだけテンポを与える。

\begin{keyblock}
\begin{proposition}[マシンの出力]
\label{prop:muscles}
脳は制御器の階層として筋肉を制御する。最上位で行動が選ばれ（第\ref{sec:dofa}章）、その下で局所的な発生器が一定の指令をリズムに変え、脊髄の速いループが関節を保持する。力は浮動小数点で指定される。指数は大きさの順にオンになる単位の数、仮数はスパイクの発火率である。各関節は引く筋肉の対で制御され、その力の差がトルクを、和が剛性を決める。これらの仕組みは測定されている。脳が体の内部モデルに頼っているというのは、十分な根拠のある仮説である。
\end{proposition}
\end{keyblock}

\section{まとめ}

\begin{table}[t]
\centering
\footnotesize
\setlength{\tabcolsep}{4pt}
\renewcommand{\arraystretch}{1.15}
\begin{tabular}{>{\raggedright}p{3.0cm}>{\raggedright}p{4.6cm}>{\raggedright\arraybackslash}p{5.2cm}}
\toprule
出力は何をするか & どうやって & 分かっていること \\
\midrule
スパイクを力に変える & 単収縮の和~\eqref{eq:forcesum}、ローパスフィルタ & 測定済み；線形の和は単純化 \\
力を加減する & 大きさの順の単位の動員、浮動小数点~\eqref{eq:sizeprinciple} & 動員の順序は測定済み \\
引くだけで押す & 筋肉の対：トルクは差、剛性は和~\eqref{eq:antagonists} & 測定済み \\
関節を保持する & 伸張ループ、\nt{GLYC}による拮抗筋の禁止 & 測定済み \\
震えない & $Gd<\pi/2$~\eqref{eq:delaystable} & 制御理論から導かれる；震えへの寄与は有力な原因 \\
速く動く & 内部モデルによる予測 & 十分な根拠のある仮説 \\
リズミカルに歩き、呼吸する & リズム発生器~\eqref{eq:halfcentre} & 発生器は測定済み；最も単純な回路はモデル \\
\bottomrule
\end{tabular}
\caption{マシンは筋肉をどう制御するか。}
\label{tab:muscles}
\end{table}

マシンの出力は、他のすべてと同じ手法でできている（表~\ref{tab:muscles}）。符号の代わりに2本の配線、DACの代わりにローパスフィルタ、線形目盛の代わりに対数目盛、クロック発生器の代わりに相互抑制と疲労。入力（第\ref{sec:sensors}章）と出力は世界を相対的な割合で測り、その間では本書が扱うマシンが働いている。

\section{問題}

\begin{exercise}[\lvlA]
\label{ex:13:1}
\emph{数字で見る浮動小数点。} 筋肉に $N=100$ 個の運動単位があり、力は $f_n=f_1q^{\,n-1}$、最も強いものは最も弱いものの $100$ 倍である。(a)~$q$、相対刻み~\eqref{eq:sizeprinciple}、ダイナミックレンジをデシベルで求めよ。(b)~同じ総力を持つ $100$ 個の等しい単位からなる「線形DAC」では、力 $10f_1$ での刻みはいくらか。
\end{exercise}

\begin{exercise}[\lvlA]
\label{ex:13:2}
\emph{安全率の値段。} 筋線維上のシナプスは1スパイクごとに平均 $m$ のポアソン数の量子を放出し、線維は $n_{\text{th}}=20$ 量子以上で発火する。安全率は $S=m/n_{\text{th}}$ である。毎秒 $20$ スパイクで働く単位は、$S=2$ と $S=4$ で1日に何回失敗するか。
\end{exercise}

\begin{exercise}[\lvlB]
\label{ex:13:3}
\emph{フィルタとしての筋肉。} $T_c=50\ms$ の単収縮~\eqref{eq:twitch}は線形系のインパルス応答である。(a)~その伝達関数 $H(s)$ と遮断周波数を求めよ。(b)~規則的なスパイクのどの頻度で、力の脈動の振れ幅が平均の力の $10\%$ になるか。
\end{exercise}

\begin{exercise}[\lvlB]
\label{ex:13:4}
\emph{遅延より速くは直せない。} 制御器 $\dd x/\dd t=-G\,x(t-d)$。(a)~振動なしで偏差が最も速く減衰するのは $Gd=1/e$ のときで、その速さは $1/d$ であることを示せ。(b)~脊髄ループ $d=30\ms$ と目を通るループ $d=150\ms$ について、この $G$ と減衰の時定数を求めよ。
\end{exercise}

\begin{exercise}[\lvlB]
\label{ex:13:5}
\emph{リズム発生器の周期。} $\tau\ll\tau_a$ のとき、モデル~\eqref{eq:halfcentre}の周期 $T$ は方程式 $(\beta-1)(1-E^{2+b})/(\beta-E)=b\,(1-E^{1+b})/(1+b)$ で決まる。ここで $E=e^{-T/(2\tau_a)}$。(a)~$\beta=b=2$, $\tau_a=0.4$~s のときの $T$ を求めよ。(b)~$T$ が入力 $I$ によらず、リズムは $b>\beta-1$ のときだけ可能であることを示せ。
\end{exercise}

\begin{exercise}[\lvlB]
\label{ex:13:6}
\emph{発生器のシミュレーション。} \texttt{models/\allowbreak muscle\_models.py} から関数 \texttt{halfcentre} をインポートし（ファイル自体は実行しない）、最初の2周期を捨てて、$b=0.8$, $1.05$, $1.2$, $1.5$, $2$, $4$ および $I=0.5$, $1$, $2$ で周期を測れ。「もっと速く歩け」という指令は $I$ を通じてテンポを変えられるか。
\end{exercise}

\begin{exercise}[\lvlB]
\label{ex:13:7}
\emph{剛性対反射。} 関節が伸張ループに囲まれている：$\dd\theta/\dd t=-a\theta(t)-G\theta(t-d)$、$a=K/B$、$d=30\ms$。これを問題~\ref{ex:13:4}に帰着させ、偏差が時定数 $15\ms$ で振動なしに減衰する最小の $a$ と、そのとき必要な $G$ を求めよ。筋肉の同時収縮を強めたら $G$ をどう変えるべきか。
\end{exercise}

\begin{exercise}[\lvlC]
\label{ex:13:8}
\emph{テンポのつまみ。} 問題~\ref{ex:13:5}と~\ref{ex:13:6}によれば、発生器~\eqref{eq:halfcentre}の入力 $I$ は振れ幅だけを変え、テンポは変えない。ある $I_{\min}$ 未満ではリズムがなく、それを超えると $I$ とともに周期が減り、$I$ を3倍にすれば少なくとも半分になるような、本書の部品を使った最小限のモデル変更を提案せよ。$\tau\ll\tau_a$ で $I_{\min}$ を求め、シミュレーションで確かめよ。
\end{exercise}

\chapter{痛み：警報信号}
\chaptermark{痛み}
\label{sec:pain}

\section{測定ではなく警報}

第\ref{sec:sensors}章のセンサはどれも世界を測っていた。光がどれだけあるか、音の高さはどうか、圧力はどれほどか。痛みのしくみは違う。痛みが伝えるのは「そこに何があるか」ではなく、「危ない、すべてを放り出せ」である。エンジニアから見れば、これは測定チャネルではなく\emph{警報信号}である。マシンがいま何をしていても、それを突き抜けて届かなければならない最優先の線路である。

警報は三つの性質で測定チャネルと異なり、痛みはその三つをすべて備えている。第一に、順応によって弱めにくい。痛みのセンサは触覚のセンサよりはるかに順応しにくく、軽い損傷の後にはかえって敏感になる~\cite{LaMotte1982hyper}。強い加熱の後の応答の低下が測定されているのは、その一つの型だけである~\cite{Treede1995heat}。光センサは一定の背景のもとで沈黙する（図~\ref{fig:adaptation}）が、1分で黙る警報は役に立たない。第二に、しきい値が高い。痛みのセンサは作用が危険に近いときにしか発火しない。例えば加熱に対しては、センサの型によってしきい値は $41^\circ$ から $46$--$48^\circ$ である~\cite{Treede1995heat, Weidner1999cfib}。第三に、そしてこれがこの章の主題だが、警報の音量を決めるのはセンサだけでなく、マシン自身でもある。同じ傷でも状況によって痛み方が違う。これがすでにおなじみのどんな部品で組み立てられているかを見よう。

痛みのセンサは、第\ref{sec:sensors}章の他の感覚ニューロンと同じく\nt{GLUT}ニューロンである。その一部はグルタミン酸とともに、付録~\ref{sec:othermediators}のサブスタンス~P を放出し、伝達をゆっくりと強める。

\section{2本の配線：速い痛みと遅い痛み}

指をぶつけると、痛みを2回感じる。最初は短く鋭い痛み、1秒ほど後に鈍く長い痛みである。これは2本の別々の配線である。速い配線は絶縁されていて、スパイクを速さ $v$ 毎秒 $5$ から $30$ メートルで伝える。遅い配線は細くむき出しで、毎秒 $0.5$--$2$ メートルである。足からの信号は約1メートル進むので、到着時間
\begin{empheq}[box=\keybox]{equation}
t \;=\; \frac{L}{v}
\label{eq:paintime}
\end{empheq}
は、速い配線では数十ミリ秒、遅い配線では約1秒になる。2本の配線は二つのメッセージを運ぶ。速い配線は\emph{どこで}\emph{いつ}を伝える。そのスパイクは第\ref{sec:code}章の時間符号のように、より早く、より正確に届く。遅い配線は\emph{どれほど悪いか}を伝え、その長く鈍い痛みは損傷が治まるまでマシンを警報モードに保つ。

\section{脊髄のゲート}

痛みの信号が最初に着くのは脊髄であり、そこで信号は\emph{ゲート}を通る~\cite{MelzackWall1965}。このゲートを構成する具体的なニューロンが見つかったのは比較的最近である~\cite{Duan2014}。痛みを脳へ送る出力\nt{GLUT}ニューロンには、痛みの配線と触覚の配線が集まる。そしてその隣に、触覚によって興奮する抑制性の介在ニューロン、\nt{GABA}または\nt{GLYC}がある（図~\ref{fig:gate}）。触覚はゲートを閉じる。元のゲート理論~\cite{MelzackWall1965}では、逆に痛みがこの介在ニューロンをオフにし、強い痛みはゲートを開く。これは理論の仮定であり、見つかったゲートのニューロンで直接示されてはいない。

\begin{figure}[t]
\centering
\begin{tikzpicture}[>=Stealth,
  nrn/.style={circle,draw,very thick,minimum size=1.0cm,inner sep=0pt,font=\scriptsize},
  inh/.style={-{Bar[width=7pt]},thick,red!70!black},
  bc/.style={->,thick,densely dashed,orange!85!black}]
\node[nrn,fill=blue!6] (Z) at (6,0) {\nt{GLUT}};
\node[nrn,fill=red!10] (Q) at (3.6,-1.6) {\nt{GABA}};
\draw[->,very thick,red!70!black] (0,0.6) node[left,font=\small,black,align=right] {痛み $\nu$} -- (5.45,0.15);
\draw[->,very thick,blue!60!black] (0,-2.6) node[left,font=\small,black,align=right] {触覚 $\mu$} -- (2.9,-2.0);
\draw[->,thick,blue!60!black] (1.8,-2.35) -- (5.5,-0.35) node[pos=0.8,below right=-2pt,font=\scriptsize] {$w_{z\mu}$};
\draw[inh] (1.6,0.5) -- (3.35,-1.1) node[pos=0.5,left=1pt,font=\scriptsize] {$w_{q\nu}$};
\draw[inh] (Q) -- (Z) node[pos=0.5,above left=-1pt,font=\scriptsize] {$\beta$};
\draw[->,very thick] (Z) -- (8.2,0) node[right,font=\small,align=left] {脳へ：$z$};
\draw[bc] (6,2.1) node[above,font=\scriptsize,black,align=center] {上から：\nt{SERO}, \nt{NORA}, オピオイド} -- (Z.north) node[pos=0.45,right,font=\scriptsize,black] {ゲイン $g$};
\end{tikzpicture}
\caption{モデル~\eqref{eq:gate}による脊髄の痛みのゲート。出力\nt{GLUT}ニューロンは痛み $\nu$ と、触覚 $\mu$ からの弱い興奮を受ける。抑制性の介在ニューロン（\nt{GABA}または\nt{GLYC}）は触覚で興奮し、ゲート理論の仮定では痛みでオフになる。上からはブロードキャストチャネルを通じてゲイン $g$ が来る。}
\label{fig:gate}
\end{figure}

これを第\ref{sec:gaba}章と同じく発火率で書こう。介在ニューロンの発火率 $q$ と出力の発火率 $z$ は
\begin{empheq}[box=\keybox]{equation}
q \;=\; \bigl[w_{q\mu}\,\mu + q_0 - w_{q\nu}\,\nu\bigr]_+ ,
\qquad
z \;=\; \bigl[\,g\,(\nu + w_{z\mu}\,\mu) - \beta\,q\,\bigr]_+ ,
\label{eq:gate}
\end{empheq}
である。ここで $\nu$ は痛みの入力、$\mu$ は触覚の入力、$w$ は結合の重み（第1添字が受け手、第2添字が送り手）、$\beta$ は抑制の強さ、$q_0$ は介在ニューロン自身の背景活動、$g$ は上から決められるゲイン（後述）である。これは第\ref{sec:gaba}章の禁止~\eqref{eq:veto}、「痛み、ただし触覚があるときは除く」であり、ゲート理論から仮定として取り入れた一つの修正がある。強い痛みは自ら禁止ニューロンをオフにする。

$w_{q\mu}=1$, $q_0=0.2$, $w_{q\nu}=0.5$, $w_{z\mu}=0.3$, $\beta=1.5$ でモデルを計算した（図~\ref{fig:gatecurves}）。触覚がなければ、痛みは $\nu\approx0.18$ から通り始める。触覚があるとしきい値は $0.86$ に上がり、$\nu=1$ で出力は $1$ から $0.25$ へと4分の1に落ちる。だが強い痛み $\nu=2$ では、触覚はもう何も変えない。ゲートは開け放たれている。実生活でも同じである。打ち身をさするのは効くが、重いけがでは効かない。痛みのない触覚だけではゲートを通らない。触れただけで警報は上がらない。

\begin{figure}[t]
\centering
\begin{tikzpicture}[>=Stealth,x=5.2cm,y=1.35cm]
\draw[->,gray!70!black] (0,0) -- (2.12,0) node[right,font=\small,black] {痛み $\nu$};
\draw[->,gray!70!black] (0,0) -- (0,4.3) node[left,font=\small,black] {$z$};
\foreach \t in {0,0.5,1,1.5,2}{\draw[gray!60!black] (\t,-0.05) -- (\t,0.05); \node[below,font=\scriptsize] at (\t,0) {\t};}
\foreach \f in {1,2,3,4}{\draw[gray!60!black] (-0.01,\f) -- (0.01,\f); \node[left,font=\scriptsize] at (0,\f) {\f};}
\draw[very thick,red!70!black] plot file {figures/pain_gate_notouch.dat};
\draw[very thick,blue!60!black] plot file {figures/pain_gate_touch.dat};
\draw[very thick,orange!85!black,densely dashed] plot file {figures/pain_gate_sens.dat};
\draw[very thick,red!70!black] (0.08,3.9) -- (0.2,3.9) node[right,font=\scriptsize,black] {触覚なし};
\draw[very thick,blue!60!black] (0.08,3.45) -- (0.2,3.45) node[right,font=\scriptsize,black] {触覚あり};
\draw[very thick,orange!85!black,densely dashed] (0.08,3.0) -- (0.2,3.0) node[right,font=\scriptsize,black] {感作の後、$g=2$};
\end{tikzpicture}
\caption{痛みの強さに対するゲート~\eqref{eq:gate}の出力。触覚はしきい値を上げ、弱い痛みと中程度の痛みを弱めるが、強い痛みは弱めない。感作（破線）はゲインを2倍にする。同じ痛みが2倍の音量で伝えられ、しきい値は下がる。}
\label{fig:gatecurves}
\end{figure}

\section{上からの音量つまみ}

ゲートは下からだけ制御されるのではない。脳幹からは下行路が来て、\eqref{eq:gate}のゲイン $g$ を変える。\nt{SERO}, \nt{NORA}、そして付録~\ref{sec:othermediators}のオピオイドペプチドである~\cite{Fields2004}。これらは第\ref{sec:serocomputer}--\ref{sec:acet}章と同じブロードキャストチャネルであり、ここでのつまみは警報の音量である。つまみはどちらの向きにも回せる。同じ下行回路が痛みを弱めることも強めることもできる。

マシンはなぜ警報を弱める必要があるのか。警報は割り込みであり、割り込みがいつも適切とは限らないからである。捕食者から逃げる動物は、傷ついた足のせいで止まってはならない。まず走り、傷をなめるのはその後である。楽になるという期待も痛みを弱め、その効果の一部はオピオイド受容体を遮断すると消える~\cite{Levine1978}。脳で計算された期待が、オピオイドのチャネルを下ってゲートに届くのである。この描像そのものは測定されている。脳が音量をどう決めているのかは仮説である。

\section{感作：学習する警報}

損傷の後、ゲートの出力ニューロンはより敏感になる。同じ入力がより大きな出力を生み、しきい値は下がる~\cite{Woolf1983}。モデル~\eqref{eq:gate}ではこれはゲイン $g$ の上昇であり、その最も単純な記述は、痛みそのものが駆動する遅いRC回路である：
\begin{empheq}[box=\keybox]{equation}
\tau_s\,\frac{\dd g}{\dd t} \;=\; -(g-1) + k_s\,z ,
\label{eq:sensitization}
\end{empheq}
ここで $\tau_s$ は感度が正常に戻るまでの時間である。感作は損傷刺激が止んだ後も続くので~\cite{Woolf1983}、$\tau_s$ は痛みそのものの持続時間より長い。$k_s$ は駆動の強さである。組織が損傷している間はゲートの出力が $g$ を1より上に保ち、傷の近くで起こることはすべてより大きな音量で伝えられる（図~\ref{fig:gatecurves}の破線：$g=2$ でしきい値は $0.11$ に下がり、出力は2倍になる）。これは警報の合理的な設定である。損傷が治るまでは、用心しすぎるほうがよい。

エンジニアから見れば、これは遅い漏れを持つ正帰還である。痛みがゲインを上げ、ゲインが痛みを上げる。ループが弱いうちは、傷が治れば $g$ は正常に戻る。ループが強いと、警報は損傷がなくなった後もオンの状態で固まりうる。図~\ref{fig:bistable}の強いフィードバックを持つ群と同じである。モデル~\eqref{eq:sensitization}は単純化である。中枢性感作が存在することは測定されている。

\section{教師として、割り込みとしての痛み}

マシンの中で痛みには、警報のほかに二つの仕事がある。

\emph{教師。} 痛みは負の報酬であり、誤差の式~\eqref{eq:ctd}には $r<0$ として入る。第\ref{sec:dofa}章で述べたことはすべてここでも成り立つ。痛みの前触れは前もって誤差を起こすようになり、網は痛みにつながるものを避けることを学ぶ。ヒトでの実験は、痛みからの学習中の脳活動がまさに時間差分誤差に従うことを示し~\cite{Seymour2004}、\nt{DOPA}ニューロンの一部は不快な出来事にも応答する（第\ref{sec:dofaalt}章）。

\emph{割り込み。} 痛みは注意をいまの課題から引きはがす。これは副作用ではなく、痛みの目的である~\cite{EcclestonCrombez1999}。第\ref{sec:nora}章では、同じ「割り込み」の役割を\nt{NORA}の急なバーストについて論じた。そして最も速い応答は脳にさえ届かない。熱いものから手を引っ込める反射は、図~\ref{fig:antagonists}と同じ筋肉の対と同じ拮抗筋の禁止によって、脊髄の中で直接閉じている。

\begin{keyblock}
\begin{proposition}[警報信号]
\label{prop:pain}
痛みは測定ではなく、優先度の高い警報信号である。測定チャネルよりはるかに順応しにくく、2本の配線（速いほうは「どこ」、遅いほうは「どれほど悪いか」）を通り、触覚が閉じる（そしてゲート理論の仮定では強い痛みが開く）ゲートを通過し、その音量は上からブロードキャストチャネルが決める。痛みによるゲートの開放を除き、これらの仕組みは測定されている。単純なモデル~\eqref{eq:gate}と~\eqref{eq:sensitization}は単純化である。
\end{proposition}
\end{keyblock}

\section{まとめ}

\begin{table}[t]
\centering
\footnotesize
\setlength{\tabcolsep}{4pt}
\renewcommand{\arraystretch}{1.15}
\begin{tabular}{>{\raggedright}p{3.0cm}>{\raggedright}p{4.9cm}>{\raggedright\arraybackslash}p{4.9cm}}
\toprule
痛みは何をするか & どうやって & 分かっていること \\
\midrule
警報を上げる & 高いしきい値、触覚より弱い順応 & しきい値は測定済み；順応の比較は推定 \\
「どこ」と「どれほど悪いか」を伝える & 速さの違う2本の配線~\eqref{eq:paintime} & 測定済み \\
ゲートを通る & \nt{GABA}／\nt{GLYC}による禁止~\eqref{eq:gate} & ゲートは測定済み；痛みによる開放は理論の仮定；モデルは単純化 \\
音量を変える & 下行性の\nt{SERO}, \nt{NORA}、オピオイド & 測定済み；音量を選ぶ規則は仮説 \\
損傷の後に学習する & ゲインの上昇~\eqref{eq:sensitization} & 感作は測定済み；モデルは単純化 \\
回避を教える & \eqref{eq:ctd}の負の報酬 & 時間差分誤差との類似は測定済み \\
作業に割り込む & 注意の奪取；引っ込め反射 & 測定済み \\
\bottomrule
\end{tabular}
\caption{警報信号としての痛み。}
\label{tab:pain}
\end{table}

痛みはすでに見てきた部品で組み立てられている（表~\ref{tab:pain}）。\nt{GLUT}センサ、2本の配線、抑制性介在ニューロンによる禁止、ブロードキャストの音量つまみ、遅いRC回路、予測誤差、割り込み。新しいのは一つだけである。これはマシンの入力のうち、音量をマシン自身が決める唯一のものである。持ち主が弱めることも設定し直すこともできる警報なのである。

\section{問題}

\begin{exercise}[\lvlA]
\label{ex:14:1}
\emph{2本の配線。} 信号は足から来る、$L=1$~m。(a)~一斉射撃が、速さが範囲~\eqref{eq:paintime}を埋める同じ型の配線の束を通る。脊髄に着くまでに、速い型と遅い型でそれぞれどれだけ引き伸ばされるか。(b)~$15$~m/s と $1$~m/s の配線の痛みの波は、差が $0.1$~s 以上なら区別できる。どれだけの距離から融合するか。
\end{exercise}

\begin{exercise}[\lvlB]
\label{ex:14:2}
\emph{触覚が痛むとき。} 本文のパラメータのゲート~\eqref{eq:gate}で、感作によってゲイン $g$ が上がる。痛みのない触覚が、どんな強さ $\mu$ でもゲートを通らないのは $g$ がいくつまでか。触覚 $\mu=1$ が通り始めるのはどの $g$ か。
\end{exercise}

\begin{exercise}[\lvlB]
\label{ex:14:3}
\emph{感作の安定性。} 入力は一定で、ゲートの出力は $g$ について線形：$z=gA-C$、$A=\nu+w_{z\mu}\mu$、$C=\beta q\ge0$。(a)~モデル~\eqref{eq:sensitization}の平衡点 $g^*$ とその安定条件を求めよ。(b)~$k_s=0.5$ で、触覚なしと触覚 $\mu=1$ ありのそれぞれについて、それを超えるとモデルが暴走する痛みの強さを求めよ。
\end{exercise}

\begin{exercise}[\lvlA]
\label{ex:14:4}
\emph{痛みの前触れ。} 時刻 $0$ の信号が時刻 $T=2$~s の痛みを予告する。\eqref{eq:ctd}で $r(t)=-P\,\delta(t-T)$、$\tau_\gamma=2$~s。(a)~信号の時点での $\delta$ の振れの面積を求めよ。(b)~時刻 $t_e=1$~s の行動が痛みを取り消す。この時点での $\delta$ の振れはいくらか。それは網に何を教えるか。
\end{exercise}

\begin{exercise}[\lvlB]
\label{ex:14:5}
\emph{ラッチする警報。} \texttt{models/\allowbreak pain\_gate.py} のゲートにダイナミクス~\eqref{eq:sensitization}と飽和 $z\le4$ を加えよ。触覚 $\mu=1$ は常にあり、損傷 $\nu=2$ が時間 $T$ 続いた後 $\nu=0$ になる。$k_s=4$, $4.6$, $5$, $6$, $8$ について、警報がオンのまま残る最小の $T$（$\tau_s$ 単位）をシミュレーションで求めよ。
\end{exercise}

\begin{exercise}[\lvlB]
\label{ex:14:6}
\emph{ゲートの設計。} $\beta=1.5$, $w_{z\mu}=0.3$, $g=1$ で、痛みのしきい値が触覚なしで $0.2$、触覚 $\mu=1$ ありで $1.0$ となり、$\nu_\times=2.5$ で触覚が痛みを弱めなくなるように、\eqref{eq:gate}の $q_0$, $w_{q\nu}$, $w_{q\mu}$ を選べ。痛みのない触覚がゲートを通らないのはゲイン $g$ がいくつまでか。
\end{exercise}

\begin{exercise}[\lvlC]
\label{ex:14:7}
\emph{音量つまみ。} 作業は単位時間あたり価値 $V$ をもたらし、損傷 $\nu$ を抱えての作業はコスト $c\nu$ がかかるので、$\nu>V/c$ なら中断するのが得である。痛みは $z>\theta=0.5$ で割り込む。(a)~$V/c=0.25$, $0.5$, $2$ で、触覚なしのゲート~\eqref{eq:gate}にこのしきい値を与えるゲイン $g^*$ はいくらか。(b)~マシンは本書のどの信号から $V$ と $c$ を推定できるか。
\end{exercise}

\chapter{マシンは動き方をどう学ぶか}
\chaptermark{マシンは動き方をどう学ぶか}
\label{sec:motorlearning}

\section{三人の教師}

第\ref{sec:muscles}章では、マシンが筋肉をどう動かすかを見て、一つの仮説で締めくくった。速く正確な動きには、指令の結果を予測する体の内部モデルが必要である~\cite{WolpertGhahramani2000}。このモデルはどこから来るのか。学習しなければならない。それも、これまでとは違うやり方で。

本書にはすでに二人の教師が登場した。一人目は\emph{教師なし}である。ヘッブ則（第\ref{sec:dofa}章）は、よく一緒に起こることを強め、入力を圧縮する。二人目は\emph{報酬}である。\nt{DOPA}は全員に一つの数、「期待より良かったか悪かったか」を配る。運動には三人目の教師、\emph{誤差}が必要である。手がカップから左へ2センチ外れた。これは「悪かった」ではなく、各出力にどちらへどれだけ間違えたかを告げるベクトルである。エンジニアならこれを教師あり学習と呼ぶだろう。

二人目と三人目の教師の違いは、全員に1本の配線か、出力ごとに専用の配線かの違いである。命題~\ref{prop:broadcastcost}では、一つの共通の数による学習は、ノイズが結果に影響するニューロンが一つ増えるごとに遅くなった。一方、各出力 $i$ が自分の誤差 $e_i$ を受け取るなら、重みは最も単純な規則で変えられる：
\begin{empheq}[box=\keybox]{equation}
\frac{\dd w_{ij}}{\dd t} \;=\; -\eta\,e_i(t)\,x_j(t) .
\label{eq:lms}
\end{empheq}
これは適応フィルタの技術で古くから知られた最小二乗則（LMS則）である~\cite{WidrowHoff1960}。重みは自分の入力と自分の出力の誤差との積に比例して変わり、平均すると二乗誤差の勾配をたどる。第\ref{sec:dofa}章のような探索のためのノイズは要らない。補正の方向は誤差の配線を通じて直接届く。そして他の出力の誤差はシナプスに来ないので、速さは出力の数とともに落ちない。

\begin{keyblock}
\begin{proposition}[三人の教師]
\label{prop:teachers}
ヘッブ則は教師なしで、よく起こることを教える。\nt{DOPA}の信号は網全体に一つの数で、得になることを教え、ニューロンが増えるごとに遅くなる。各出力への誤差の配線は規則~\eqref{eq:lms}によって正確さを教え、その速さは出力の数によらない。三つ目の方法の代償は、出力ごとの専用配線と、正解を知っている誰かである。
\end{proposition}
\end{keyblock}

\section{誤差の配線}

運動について正解を知りうるのは誰か。センサ自身である。手が目標から左へずれたなら、それを目と第\ref{sec:sensors}章の伸張センサが伝える。必要なのは、この誤差を正しい出力へ届ける回路である。

脳には、まさにこのために作られたように見える領域がある~\cite{Marr1969,Albus1971}。その回路は異様なほど規則正しく、ほとんど結晶のようで、その中に規則~\eqref{eq:lms}を容易に見てとれる（図~\ref{fig:errorwire}）。

\emph{入力の拡張層。} 指令と体の状態についての信号は、小さな\nt{GLUT}ニューロンの巨大な層に届く。その数は約700億で、脳の残り全部を合わせたより多い~\cite{Azevedo2009}。各ニューロンはいくつかの入力を混ぜ合わせ、信号は非常に長いベクトルに展開される。エンジニアにはなじみの手法である。入力の次元を上げると、新しい空間ではほとんどどんな望みの関係も単純な重み付き和で得られる~\cite{Cover1965}。

\emph{出力ニューロン。} 重み付き和を計算するのは約3000万個の大きな出力ニューロンで~\cite{Andersen1992cereb}、それぞれが拡張層から10万ほどの入力を持つ。そしてこれらは\nt{GABA}ニューロンである。学習の結果は興奮ではなく抑制として先へ送られる。第\ref{sec:gaba}章の禁止と同じである。

\emph{誤差の配線。} 各出力ニューロンは、もう一つ特別な入力を受ける。ちょうど1本の\nt{GLUT}の配線で、めったに発火しない（毎秒1回ほど）が、発火するたびに出力ニューロンに強力なバーストを起こす~\cite{Ito2001}。これが $e_i$ である。バーストの確率は運動誤差の方向に依存する~\cite{Herzfeld2018}。誤差のバーストと入力の活動が一致すると、その入力は弱まる。まさに\eqref{eq:lms}の項 $-\eta\,e_i x_j$ である。

\begin{figure}[t]
\centering
\begin{tikzpicture}[>=Stealth,
  gran/.style={circle,draw,thick,fill=blue!6,minimum size=0.32cm,inner sep=0pt},
  outn/.style={circle,draw,very thick,fill=red!10,minimum size=1.1cm,inner sep=0pt,font=\scriptsize},
  inh/.style={-{Bar[width=7pt]},very thick,red!70!black}]
% inputs
\foreach \y/\lab in {1.6/{指令},0.4/{体の状態}}{
  \draw[->,thick,blue!60!black] (-0.6,\y) node[left,font=\scriptsize,black] {\lab} -- (0.35,\y);}
% expansion layer
\foreach \i in {0,...,11}{\node[gran] (g\i) at ({0.8+0.28*mod(\i,2)},{-0.3+0.23*\i}) {};}
\foreach \i in {0,...,11}{\pgfmathsetmacro{\yy}{mod(\i,3)>0 ? 1.6 : 0.4}\draw[gray!60!black] (0.35,\yy) -- (g\i);}
\node[font=\scriptsize,align=center] at (0.95,-1.05) {拡張層\\$\sim7\cdot10^{10}$ \nt{GLUT}};
% parallel wires to the output neuron
\foreach \i in {0,...,11}{\draw[->,gray!70!black] (g\i.east) -- (4.1,{1.0+0.07*(\i-5.5)});}
\node[outn] (P) at (4.6,1.0) {\nt{GABA}};
\node[font=\scriptsize,align=center,above=2pt] at (P.north) {出力ニューロン\\$\sim10^5$ 個の入力 $x_j$、重み $w_{ij}$};
\draw[inh] (P) -- (6.8,1.0) node[right,font=\scriptsize,align=left] {運動への\\補正};
% error wire
\draw[->,very thick,red!70!black,densely dashed] (4.6,-1.4) node[below,font=\scriptsize,black,align=center] {1本の誤差の配線 $e_i$\\\nt{GLUT}、$\sim1$ 回/s} -- (P.south);
\end{tikzpicture}
\caption{脳が実装していると思われる教師あり学習の回路。指令と体の状態は\nt{GLUT}ニューロンの巨大な層によって長いベクトル $x_j$ に展開され、出力\nt{GABA}ニューロンがそれを重み $w_{ij}$ で足し合わせる。ただ1本の誤差の配線が $e_i$ を運ぶ。誤差と活動中の入力が一致すると、規則~\eqref{eq:lms}のとおり、その入力の重みが弱まる。}
\label{fig:errorwire}
\end{figure}

一つの難しさは第\ref{sec:dofa}章でおなじみである。誤差は、それを招いた入力より後に届く。外れたことが見えるのは指令の数十〜数百ミリ秒後である。したがってシナプスは自分が活動していたことを覚えていなければならない。規則~\eqref{eq:threefactor}と同じく適格度トレースが必要である。そして実験によれば、このトレースの長さは課題ごとのフィードバックの遅延に合わせてある。誤差が100ミリ秒後に届くところでは、誤差の100ミリ秒前に活動していた入力が最も強く弱められる~\cite{Suvrathan2016}。

\begin{keyblock}
\begin{proposition}[誤差の配線]
\label{prop:errorwire}
図~\ref{fig:errorwire}の回路では、各出力ニューロンはちょうど1本の誤差の配線を受け、誤差と入力が一致するとその入力が弱まる。これは、巨大な次元に拡張された入力に適用した最小二乗則~\eqref{eq:lms}である。回路の構成と一致による弱化は測定されている。回路全体が教師あり学習として働くというのは、有力な仮説である。
\end{proposition}
\end{keyblock}

\section{適応フィルタとしての内部モデル}

このような回路は何を学ぶのか。最も自然な答えは予測である。入力として指令 $u(t)$ が、拡張層のように時定数の異なるフィルタの組を通って来るとする：$x_j(t)=(g_j*u)(t)$。出力はそれらの重み付き和で、センサが伝えるであろう値の予測である：
\begin{empheq}[box=\keybox]{equation}
\hat y(t) \;=\; \sum_j w_j\,(g_j*u)(t), \qquad e(t) \;=\; \hat y(t) - y(t) ,
\label{eq:adaptivefilter}
\end{empheq}
そして重みは\eqref{eq:lms}によって学習される。これは\emph{適応フィルタ}である。未知の系を再現するように、自分の伝達関数を自ら調整する回路である~\cite{Fujita1982}。ここでの未知の系は、質量、弾性、遅延を持つ自分自身の体である。学習したフィルタこそが第\ref{sec:muscles}章の内部モデルであり、センサを待たずに、センサが何を伝えるかを告げる。

このモデルは二重に役に立つ。第一に、その予測を遅れたフィードバックの代わりに使えば、安定条件~\eqref{eq:delaystable}に手を縛られなくなる。モデルにもとづいて素早く修正できる。第二に、予測と実際の差は\emph{予期しないこと}の信号である。マシンが自分でしたことはすべてモデルが予測して差し引き、残るのは世界がしたことである。ある仮説によれば、自分で自分をくすぐれないのはそのためである~\cite{Blakemore1998}。

\section{世界が変わるとき：速い過程と遅い過程}

この回路を、簡単に行える実験で試してみよう。人が目標に手を伸ばしているとき、世界が不意に変わる。例えば手に横向きの力がかかる。最初の数回の動きは外れ、やがて外れは小さくなる。力を取り除くと、手はまず\emph{反対の}方向に外れる。モデルが力を学習し、それを補償し続けているのである。これは、単に「うまくできるようになった」のではなく、モデルが学習されたことの直接の証拠である。

実験はさらに微妙なことも示す。二つの過程が同時に学習している。速く学び速く忘れる速い過程と、ゆっくり学び長く覚えている遅い過程である~\cite{Smith2006}。補正は動きのたびに、出来事ごとに更新される：
\begin{empheq}[box=\keybox]{equation}
e = p - (x_f + x_s), \qquad x_f \leftarrow A_f\,x_f + B_f\,e, \qquad x_s \leftarrow A_s\,x_s + B_s\,e ,
\label{eq:tworate}
\end{empheq}
ここで $p$ は外乱、$x_f$ と $x_s$ は二つの過程が学習した補正、$A$ は補正が次の動きまでどれだけ保たれるか、$B$ は誤差からどれだけ強く学ぶかである。速い過程では $B$ が大きく、$A$ は1よりかなり小さい。遅い過程ではその逆である。

\begin{figure}[t]
\centering
\begin{tikzpicture}[>=Stealth,x=0.034cm,y=1.9cm]
\draw[->,gray!70!black] (0,0) -- (355,0) node[right,font=\small,black] {動き};
\draw[->,gray!70!black] (0,-1.15) -- (0,1.25);
\foreach \v in {-1,1}{\draw[gray!60!black] (-3,\v) -- (3,\v); \node[left,font=\scriptsize] at (0,\v) {$\v$};}
\foreach \n in {100,200,300}{\draw[gray!60!black] (\n,-0.04) -- (\n,0.04); \node[below,font=\scriptsize] at (\n,0) {\n};}
\draw[thin,gray!70!black] plot file {figures/tworate_p.dat};
\draw[thick,orange!85!black] plot file {figures/tworate_fast.dat};
\draw[thick,blue!60!black] plot file {figures/tworate_slow.dat};
\draw[very thick,black] plot file {figures/tworate_net.dat};
\node[font=\scriptsize,gray!50!black] at (120,1.12) {外乱 $p$};
\node[font=\scriptsize,black,anchor=west] at (238,0.52) {和：復帰};
\node[font=\scriptsize,blue!60!black,anchor=west] at (150,0.39) {遅い $x_s$};
\node[font=\scriptsize,orange!85!black,anchor=west] at (60,0.08) {速い $x_f$};
\draw[densely dashed,gray!60!black] (226,-1.15) -- (226,1.25);
\node[font=\scriptsize,align=center,anchor=south west] at (228,-1.1) {誤差は\\もう来ない};
\end{tikzpicture}
\caption{モデル~\eqref{eq:tworate}による二つの学習過程、$A_f=0.92$, $B_f=0.1$, $A_s=0.996$, $B_s=0.02$。外乱 $p=1$ で200回の動き、次に和がゼロに戻るまで逆向きの $p=-1$ で6回。破線の後は誤差が届かなくなり、和はひとりでに古い補正へ戻る。速い過程は逆向きの外乱を忘れ、遅い過程は順向きの外乱を覚えている。}
\label{fig:tworate}
\end{figure}

モデル~\eqref{eq:tworate}は、それなしでは理解しにくい実験を説明する（図~\ref{fig:tworate}）。我々の計算では、外乱のもとでの200回の動きで補正の和は $0.84$ に達し、その大部分 $0.63$ を遅い過程が担う。次に逆向きの外乱での6回の動きが和をゼロに戻す。速い過程はマイナスの $-0.53$ に行き、遅い過程はまだプラスの $0.46$ にある。ここで誤差を伝えるのをやめると、速い過程は数十回の動きで自分の補正を忘れ、遅い過程ははるかに長く保つので、和は40回の動きで\emph{ひとりでに} $0.37$ まで上がる。訓練なしに古い技能が戻ってくる。この自発的回復はヒトで測定されている。一つの過程のモデルではこれは出せない。誤差がなければ単にゼロへ減衰するだけである。

なぜ過程が二つなのか。もっともらしい工学的な答えを、第\ref{sec:acet}章の最適フィルタが与える。体はさまざまな理由でさまざまな速さで変わる。筋肉は数分で疲れ、数か月で成長し衰える。外乱に速いものと遅いものがあるなら、最良の推定は時定数の異なる二つのフィルタからなり、それぞれが自分の外乱を捉える~\cite{Kording2007}。速い過程は一時的な補正、「腕が疲れた」であり、遅い過程は「腕が変わった」である。これはまだ仮説だが、1日で学んだ技能が翌日までに一部失われる理由と、2回目にはより速く学べる理由を説明する。

\section{まとめ}

\begin{table}[t]
\centering
\footnotesize
\setlength{\tabcolsep}{4pt}
\renewcommand{\arraystretch}{1.15}
\begin{tabular}{>{\raggedright}p{3.0cm}>{\raggedright}p{4.6cm}>{\raggedright\arraybackslash}p{5.2cm}}
\toprule
回路は何をするか & どうやって & 分かっていること \\
\midrule
教師ありで学ぶ & 各出力への誤差の配線、規則~\eqref{eq:lms} & 回路の構成と一致による弱化は測定済み；教師あり学習は有力な仮説 \\
入力を拡張する & 700億個の\nt{GLUT}ニューロン & ニューロン数は測定済み；拡張の役割は理論 \\
誤差の遅延を考慮する & 遅延に合わせた適格度トレース & 測定済み \\
内部モデルを作る & 適応フィルタ~\eqref{eq:adaptivefilter} & 十分な根拠のある仮説 \\
二つのテンポで学ぶ & 速い過程と遅い過程~\eqref{eq:tworate} & 後効果と自発的回復は測定済み；二つの過程はモデル \\
\bottomrule
\end{tabular}
\caption{マシンは動き方をどう学ぶか。}
\label{tab:motorlearning}
\end{table}

これでマシンには三人の教師がそろった（表~\ref{tab:motorlearning}）。ヘッブはよく起こることを圧縮し、\nt{DOPA}は一つの数で得になるものを選ぶことを教え、誤差の配線は正確さを教える。しかも速く教える。各出力に自分の誤差が届くからである。脳はおそらく三つすべてを使い、それぞれを最も安上がりな場所で使っている。ブロードキャスト信号は少数の決定に、専用の配線は、正解をセンサ自身が伝える体の精密な調整に。

\section{問題}

\begin{exercise}[\lvlA]
\label{ex:15:1}
\emph{誤差の配線を持つ回路の容量。} 回路に入力 $10^5$ 個の出力ニューロンが $3\cdot10^7$ 個あり、それぞれ専用の誤差の配線（$1$~スパイク/s）がある。$n$ 個の入力を持つニューロンは $\approx2n$ 個のベクトルの任意のラベル付けを学習できる~\cite{Cover1965}。(a)~回路はいくつの連想を保持するか。(b)~$\Delta t=10\ms$ で\eqref{eq:spikeentropy}により、ニューロンの容量を満たす最短時間を見積もれ。
\end{exercise}

\begin{exercise}[\lvlB]
\label{ex:15:2}
\emph{最小二乗則の速さ。} 規則~\eqref{eq:lms}のモードは速さ $\eta\lambda_k(C)$ で減衰する。白色ノイズ上の5個のフィルタ~\eqref{eq:adaptivefilter}では $C_{jk}=2\sqrt{\tau_j\tau_k}/(\tau_j+\tau_k)$ である。$\tau_j$ が $2$ 倍ずつ（$10$--$160\ms$）増える場合と $4$ 倍ずつ増える場合に、最も速い速さと最も遅い速さの比を求めよ。
\end{exercise}

\begin{exercise}[\lvlB]
\label{ex:15:3}
\emph{二つの過程の解析。} 図~\ref{fig:tworate}のパラメータでモデル~\eqref{eq:tworate}を $\mathbf x_{n+1}=M\mathbf x_n+\mathbf b\,p$ の形に書け。(a)~$M$ の固有値から、動きの回数で測った時定数を求めよ。(b)~一定の $p=1$ での定常な $x_f$, $x_s$ とその和を求めよ。
\end{exercise}

\begin{exercise}[\lvlB]
\label{ex:15:4}
\emph{シミュレーション：2回目はより速い。} \texttt{models/\allowbreak motor\_learning.py} のパラメータでモデル~\eqref{eq:tworate}を使い、$p=1$ で $x_f+x_s\ge0.8$ になるまでの動きの回数を求めよ。(a)~未学習のマシン。(b)~図~\ref{fig:tworate}のように、$p=1$ で $200$ 回動いた後、和がゼロになるまで逆向きの外乱を加えた後。一つの過程のモデルは、このように速くなりうるか。
\end{exercise}

\begin{exercise}[\lvlB]
\label{ex:15:5}
\emph{内部モデルを持つ制御器。} 対象 $\dd x/\dd t=ku$ は遅延 $d$ で観測される。制御器 $u=-G\hat x$ は $\hat x(t)=x(t-d)+\hat k\int_{t-d}^{t}u\,\dd s$ と予測する。(a)~$\hat k=k=1$, $d=150\ms$ で、時定数 $20\ms$ を与える $G$ を求め、限界~\eqref{eq:delaystable}と比べよ。(b)~$\hat k=0.4k$ でこの系は安定か。
\end{exercise}

\begin{exercise}[\lvlB]
\label{ex:15:6}
\emph{適応フィルタが筋肉を学ぶ。} 対象は面積1の単収縮~\eqref{eq:twitch}、$T_c=50\ms$。フィルタ~\eqref{eq:adaptivefilter}は $\tau_j=12.5$, $25$, $50$, $100$, $200\ms$ のRC回路。入力は $10\ms$ ごとに新しい正規乱数。$\eta=50$ と $200$ で、規則~\eqref{eq:lms}が誤差を $0.5\%$ まで下げるのに何秒かかるか。$300$~s 後でも重みが最良値から遠いのはなぜか。
\end{exercise}

\begin{exercise}[\lvlC]
\label{ex:15:7}
\emph{最適フィルタとしての二つの過程。} 外乱は過程 $p^{f,s}_{n+1}=A_{f,s}\,p^{f,s}_n+w^{f,s}_n$（ノイズの分散 $q_f$, $q_s$）の和で、誤差は分散 $R$ のノイズとともに観測される。カルマンフィルタは\eqref{eq:tworate}を与える。$A_f=0.92$, $A_s=0.996$ から $B_f=0.1$, $B_s=0.02$ が出るのは、$q_f/R$, $q_s/R$ がいくつのときか。$q_s$ を4倍にすると $B_f$, $B_s$ はどう変わるか。
\end{exercise}

\chapter{第二の出力：脳は体の化学をどう制御するか}
\chaptermark{化学出力}
\label{sec:chemoutput}

\section{速い出力と遅い出力}

第~\ref{sec:muscles}~章では、マシンの速い出力である筋肉を見た。だがマシンには遅い第二の出力もある。脳は体温、心拍数、水分と糖の量を必要な範囲に保ち、体を走る準備や眠る準備に切り替える。そのために関節を動かすのではなく、化学信号を放出する。経路は二つある。一つは内臓、すなわち心臓、血管、腸、腺へ直接向かう神経である。もう一つは血液である。一部のニューロンや腺は、物質を隣の細胞とのすき間ではなく、血流に直接放出する。

血液は、本書に登場したバスの中で最も広域のブロードキャストバスである。第~\ref{sec:serocomputer}--\ref{sec:acet}~章のブロードキャストチャネルは脳の広い領域に信号を配った。血液は体のすべての細胞に配る。血液はおよそ1分で全身を一周するので、メッセージは数十秒で全体に届き、物質が分解されるまで数分から数時間とどまる。この配信にはアドレスがまったくない。命題~\ref{prop:receptor}と同じく、メッセージの意味は受け手が決める。応答するのは、その物質の受容体をもつ細胞だけである。

\section{アクセルとブレーキ：臓器への2本の配線}

最良の例は心臓である。心臓には、第~\ref{sec:serocomputer}~章の\nt{SERO}ニューロンのように独自のペースメーカーがある。ペースメーカーは自分でリズムを刻み、神経をすべて切っても心臓は毎分約100回拍動する。脳はリズムを生み出すのではなく、調整するだけである。しかも符号が逆の2本の配線で調整する（図~\ref{fig:gasbrake}）。一方を通るのは\nt{NORA}で、これは\emph{アクセル}としてリズムを速める。もう一方を通るのは\nt{ACET}で、これは\emph{ブレーキ}として遅くする。大まかに書けば
\begin{empheq}[box=\keybox]{equation}
f \;\approx\; f_0 \;+\; a_S\,S \;-\; a_P\,P ,
\label{eq:gasbrake}
\end{empheq}
である。ここで$f_0\approx100$拍/分はペースメーカー固有のリズム、$S$と$P$はアクセルとブレーキの活動である。安静時はブレーキが踏まれており、心臓はふつう毎分60--70回程度で拍動する。負荷がかかるとブレーキを離し、アクセルを踏む。

\begin{figure}[t]
\centering
\begin{tikzpicture}[>=Stealth,
  blk/.style={draw,very thick,rounded corners=2pt,align=center,font=\small},
  pace/.style={circle,draw,very thick,double,double distance=1.2pt,minimum size=1.8cm,inner sep=0pt,font=\scriptsize,fill=red!8,align=center},
  inh/.style={-{Bar[width=7pt]},thick,red!70!black}]
\node[blk,fill=blue!5,text width=1.6cm] (B) at (0,0) {脳の\\指令};
\node[pace] (H) at (6.2,0) {ペース\\メーカー};
\draw[->,very thick,orange!85!black] (B.north east) to[bend left=20] node[above,font=\scriptsize,black,align=center] {\nt{NORA}：アクセル、秒単位} (H.north west);
\draw[inh,very thick,blue!60!black] (B.south east) to[bend right=20] node[below,font=\scriptsize,black,align=center] {\nt{ACET}：ブレーキ、1秒未満} (H.south west);
\draw[->,very thick] (H) -- (8.4,0) node[right,font=\small] {心拍数 $f$};
\node[blk,fill=orange!10,text width=1.6cm] (A) at (0,-2.6) {\nt{ADRE}\\血中へ};
\draw[->,thick,orange!85!black] (B) -- (A);
\foreach \y/\lab in {-2.0/{心臓},-2.6/{血管},-3.2/{肝臓、筋肉}}{
  \draw[->,thick,densely dashed,orange!85!black] (A.east) -- (4.3,\y) node[right,font=\scriptsize,black] {\lab};}
\end{tikzpicture}
\caption{心臓のペースメーカーへ向かう符号が逆の2本の配線。アクセル\nt{NORA}は遅く、ブレーキ\nt{ACET}は速い。下段は同じアクセルをブロードキャストバスで送る経路である。\nt{ADRE}細胞がアドレナリンを血中に放出し、適合する受容体をもつすべての臓器が信号を受け取る（破線の矢印）。}
\label{fig:gasbrake}
\end{figure}

これは第~\ref{sec:glutcomputer}~章の2線式符号であり、第~\ref{sec:muscles}~章の筋肉の対でもある。ただし今回の配線が制御するのは関節ではなく、ペースメーカーのテンポである。そして2本の配線は速さが違う。どちらもローパスフィルタとして働くが、ブレーキはアクセルより数倍速く変化を通す~\cite{Berger1989}。\nt{ACET}の経路は短い。受容体に結合したタンパク質が、細胞内のセカンドメッセンジャーを介さずにカリウムチャネルを直接開く~\cite{Logothetis1987gbg}。一方\nt{NORA}は細胞内の反応の連鎖を通じて作用する。エンジニアならここに、速さの異なる二つのアクチュエータを組み合わせる手法を見て取るだろう。速い方が精密な即時調整を、遅い方が大きく長い変化を受け持つ。

ここで\nt{ADRE}の役割も見えてくる。表~\ref{tab:types}では、脳内での役割は不明と書いた。脳の外では明確である。アクセル経路の細胞の一部は、臓器へ配線を送らず、アドレナリンを直接血中に放出する。これは同じアクセルをブロードキャストバスで送るものである。数十秒で心臓、血管、肝臓、筋肉に同時に届き、数分間とどまる。アドレス付きのアクセル\nt{NORA}は一つの臓器を調整し、ブロードキャストのアクセル\nt{ADRE}は全身を走るモードに切り替える。

\section{フィードバック付きカスケード}

多くのホルモンは一つの腺からではなく、カスケードで放出される。脳の底部にある小さなニューロン群が最初の信号を血中に放出する。それが中継の腺に第二の信号を放出させ、第二の信号が最終段の腺に、体に作用するホルモンを放出させる。そして最終段のホルモンは前段に戻ってそれを抑える。こうして、負帰還で囲まれた3段の増幅器ができる。これはレベル調節器であり、式~\eqref{eq:feedback}の\nt{SERO}系と同じ種類のものである。

各段を時定数$\tau$、利得$g$のRC回路で、フィードバックを係数$k$で表す。
\begin{equation}
\tau\,\frac{\dd x_1}{\dd t} = -x_1 + g\bigl[u - k\,x_3\bigr]_+ ,\qquad
\tau\,\frac{\dd x_2}{\dd t} = -x_2 + g\,x_1 ,\qquad
\tau\,\frac{\dd x_3}{\dd t} = -x_3 + g\,x_2 ,
\label{eq:cascade}
\end{equation}
ここで$u$は脳の指令、$x_3$は最終段のホルモンのレベルである。ループ利得は$K=g^3k$である。平衡では$x_3=g^3u/(1+K)$となり、どのフィードバック調節器でもそうであるように、大きなループ利得はレベルを各段の利得のばらつきに鈍感にする。しかし各段は位相をずらす。角周波数$\omega=\sqrt3/\tau$では、同一の3段が合わせて半回転の位相遅れと8分の1への減衰を生む。ここから制御理論の古典的な結果が得られる。
\begin{empheq}[box=\keybox]{equation}
K \;<\; 8 ,
\label{eq:cascadestable}
\end{empheq}
これを超えると調節器は周期約$2\pi\tau/\sqrt3\approx3.6\,\tau$で発振し始める。

\begin{keyblock}
\begin{proposition}[精度と安定性のトレードオフ]
\label{prop:cascade}
比例フィードバックをもつ3段カスケードでは、レベルの誤差は$1/(1+K)$であり、安定なのは$K<8$のときに限られる。したがってこの調節器は、およそ10パーセントより高い精度でレベルを保てない。利得を上げようとすれば、調節器は発振器に変わる。
\end{proposition}
\end{keyblock}

これをモデル~\eqref{eq:cascade}で確かめた（図~\ref{fig:cascade}）。$K=4$では、レベルはスイッチオン後に少し振動してから落ち着く。$K=7$でも落ち着くが、ゆっくりである。$K=12$では振動は減衰せず、際限なく成長することもない。段は負の信号を出せないので、カスケードは平均レベルの$0.83$倍から$1.24$倍の間を周期$3.7$--$3.9\,\tau$で規則正しく脈動する状態に落ち着く。

\begin{figure}[t]
\centering
\begin{tikzpicture}[>=Stealth,x=0.3cm,y=1.2cm]
\draw[->,gray!70!black] (0,0) -- (41.5,0) node[right,font=\small,black] {$t/\tau$};
\draw[->,gray!70!black] (0,0) -- (0,2.8) node[left,font=\small,black] {$x_3/x_3^*$};
\foreach \t in {0,10,20,30,40}{\draw[gray!60!black] (\t,-0.05) -- (\t,0.05); \node[below,font=\scriptsize] at (\t,0) {\t};}
\foreach \f in {1,2}{\draw[gray!60!black] (-0.3,\f) -- (0.3,\f); \node[left,font=\scriptsize] at (0,\f) {\f};}
\draw[densely dashed,gray!60!black] (0,1) -- (40,1);
\draw[thick,blue!60!black] plot file {figures/cascade_K4.dat};
\draw[thick,red!70!black] plot file {figures/cascade_K12.dat};
\node[font=\scriptsize,blue!60!black,anchor=west] at (16,0.55) {$K=4$：落ち着く};
\node[font=\scriptsize,red!70!black,anchor=west] at (16,1.75) {$K=12$：脈動する};
\end{tikzpicture}
\caption{式~\eqref{eq:cascade}によるフィードバック付き3段カスケード。最終段のホルモンのレベルを平衡値$x_3^*$に対する比で示す。ループ利得が8未満ならレベルは落ち着き、8を超えるとカスケードは自ら規則正しい脈動を生み出す。}
\label{fig:cascade}
\end{figure}

実際のカスケードのホルモンは、確かにパルス状に放出される。たとえばストレスホルモンのレベルは約1時間の周期で振動する。ある仮説によれば、この脈動は遅れを伴う段間のフィードバックそのものが生み出しており、別のリズム発生器は必要ない~\cite{Walker2010}。これは命題~\ref{prop:ei}の\nt{GLUT}--\nt{GABA}ループや、条件~\eqref{eq:delaystable}の筋伸張ループと同じ振動の原因、すなわち遅れる負帰還である。

\section{ホルモンのパルス符号}

パルスは副作用であるだけでなく、符号でもありうる。生殖を制御するニューロンは、自らの信号を約1時間に1回、短い分量ずつ血中に放出する。受け手の腺に同じ信号を分量ずつではなく連続して与えると、腺はパルスのときのように全力で働くのではなく、沈黙する。再び1時間に1回の分量で与えると、働きは回復する~\cite{Belchetz1978}。つまり受け手が読んでいるのはレベルではなく、分量の\emph{頻度}である。

エンジニアにとってここに謎はない。物質に長くさらされると疲れて応答しなくなる受容体は、第~\ref{sec:code}~章の枯渇するシナプス~\eqref{eq:depression}と同じくハイパスフィルタである。連続信号は抑え込み、変化は通す。このような受け手には、パルスでしか情報を伝えられない。情報はレベルから、分量の頻度と形へ移る。さらに、分量の頻度が違えば同じ腺の異なる細胞への作用も違うので、1本の化学的な線で複数の量を伝えられる。第~\ref{sec:dofaalt}~章で、1本の\nt{DOPA}配線に速い成分と遅い成分が流れていたのと同じである。

\section{概日リズム：モードを切り替える遅い発振器}

体で最も遅いリズムは概日リズムである。これは細胞内の遅れを伴う負帰還が生み出す。遺伝子がタンパク質を作り、そのタンパク質が数時間遅れて自らの産生を抑える~\cite{TakahashiJS2017}。またしても遅れる負帰還であり、本書で4度目である。そしてまたリズムが生じる。今回の周期は約1日である。主発振器は脳の底部にある数万個のニューロンからなる小さな群で、そのリズムが体の他の細胞を同期させる。

ヒトにおけるこの発振器の固有周期はちょうど1日ではなく、平均$24.18$時間である~\cite{Czeisler1999}。毎日、眼のセンサ（第~\ref{sec:sensors}~章）から届く光がこれを合わせ直す。電子技術者にとってこれは\emph{位相同期ループ}（PLL）である。固有周波数$\omega_0$の発振器が周波数$\omega$の外部信号に同期し、位相差$\varphi$は次の方程式に従う。
\begin{empheq}[box=\keybox]{equation}
\frac{\dd\varphi}{\dd t} \;=\; (\omega-\omega_0) \;-\; K_\varphi\,\sin\varphi .
\label{eq:pll}
\end{empheq}
$|\omega-\omega_0|<K_\varphi$なら同期は保たれる。周期$24.18$時間の発振器は毎日約11分ずれを補正する必要がある。朝の光はふつうこの補正に十分である。極夜や光のない洞窟では同期がなくなり、リズムは固有周期で進むようになる。

概日発振器はコンピュータのクロック発振器ではない。計算のステップを数えることも、ニューロンのスパイクを同期させることもしない。これが切り替えるのは\emph{モード}である。睡眠と覚醒、ホルモンのレベル、体温、摂食の準備。これは第~\ref{sec:serocomputer}--\ref{sec:acet}~章のレジスタと同じ種類の遅いブロードキャストレジスタであり、ただその値が太陽に合わせたスケジュールで変わる。

\begin{keyblock}
\begin{proposition}[化学出力]
\label{prop:chemoutput}
マシンの遅い出力は、速い出力と同じ部品で作られている。臓器は符号が逆で速さの異なる2本の配線、すなわちアクセル\nt{NORA}とブレーキ\nt{ACET}を受け取り、全身はアドレナリン\nt{ADRE}を含む血中のブロードキャスト信号を一斉に受け取る。ホルモンのレベルは負帰還付きカスケードが保ち、その精度は安定性によって制限される。一部のホルモンは分量の頻度で符号化される。概日リズムは、光で位相同期する遅い発振器である。経路の構造、分量投与の実験、周期は測定済みである。カスケードの脈動をフィードバックそのもので説明するのは仮説である。
\end{proposition}
\end{keyblock}

\section{まとめ}

\begin{table}[t]
\centering
\footnotesize
\setlength{\tabcolsep}{4pt}
\renewcommand{\arraystretch}{1.15}
\begin{tabular}{>{\raggedright}p{3.1cm}>{\raggedright}p{4.8cm}>{\raggedright\arraybackslash}p{4.9cm}}
\toprule
化学出力の役割 & 方法 & わかっていること \\
\midrule
臓器を調整する & アクセル\nt{NORA}とブレーキ\nt{ACET}~\eqref{eq:gasbrake} & 測定済み。ブレーキはアクセルより速い \\
全身を走るモードにする & アドレナリン\nt{ADRE}を血中へ & 測定済み \\
ホルモンのレベルを保つ & フィードバック付きカスケード、$K<8$~\eqref{eq:cascadestable} & カスケードとフィードバックは測定済み。ループ自体による脈動は仮説 \\
分量の頻度で伝える & ハイパスフィルタとしての疲れる受容体 & 分量への依存は測定済み \\
1日の中でモードを切り替える & 光で位相同期する発振器~\eqref{eq:pll} & 周期と光による同期は測定済み \\
\bottomrule
\end{tabular}
\caption{マシンは体の化学をどう制御するか。}
\label{tab:chemoutput}
\end{table}

マシンには二つの出力がある（表~\ref{tab:chemoutput}）。速い出力は筋肉で、ミリ秒単位、アドレス付きで、関節ごとに届く。遅い出力は体の化学で、秒、分、日の単位で働き、2本の配線で臓器へ、血液で全身へ一斉に届く。どちらの出力の部品も同じである。符号が逆の2本の配線、ペースメーカー、フィードバックとその遅れ。そしてそれは、マシンの内部を組み立てている部品と同じである。

\section{問題}

\begin{exercise}[\lvlA]
\label{ex:16:1}
\emph{フィルタとしての血液。}ホルモンは半減期$t_{1/2}=10$~minで血中から除去され、$T=60$~minごとに短い分量で放出される。濃度の最大値は最小値の何倍か。分量の基本波は何分の1に減衰するか。最大値が最小値のわずか2倍になるのは$t_{1/2}$がいくつのときか。
\end{exercise}

\begin{exercise}[\lvlA]
\label{ex:16:2}
\emph{概日同期。}発振器~\eqref{eq:pll}の固有周期は$24.18$~h、光による位相シフトは1日最大1時間、すなわち$K_\varphi=2\pi/24$~rad/dayである。定常位相差$\varphi^*$を分単位で求め、緩和時間を求めよ。外部信号が$\pm6$~h跳んだあと、ずれが1時間未満になるまで何日かかるか。
\end{exercise}

\begin{exercise}[\lvlB]
\label{ex:16:3}
\emph{精度と安定性。}線形化したカスケード~\eqref{eq:cascade}を同一の$n$段に延ばす。臨界利得$K_c(n)$と、$n=3$および$n=6$で達成できる最小誤差$1/(1+K_c)$を求めよ。$n\to\infty$でこの誤差は何に近づくか。
\end{exercise}

\begin{exercise}[\lvlB]
\label{ex:16:4}
\emph{アクセルとブレーキ。}式~\eqref{eq:gasbrake}で、アクセルとブレーキの寄与は$\tau_S=2$~s、$\tau_P=0.4$~sのRC回路の出力であり、指令は最大$80$拍/分と$60$拍/分とする。$f_0=100$、安静時はブレーキ$35$、アクセル$0$で$f=65$である。$f=120$に達するのにかかる時間はいくらか。ブレーキを離さずアクセルだけで動かすとどうなるか。
\end{exercise}

\begin{exercise}[\lvlB]
\label{ex:16:5}
\emph{シミュレーション：遅れのあるカスケード。}\texttt{models/\allowbreak chem\_models.py}の関数\texttt{cascade}（コピーして使い、ファイル自体は実行しない）のフィードバックに遅れ$d$を加え、第1段が$x_3(t-d)$を見るようにせよ。$d/\tau=0$、$0.5$、$1$、$2$で$K_c$と境界での周期を求め、$0.9K_c$と$1.1K_c$でシミュレーションにより確かめよ。
\end{exercise}

\begin{exercise}[\lvlB]
\label{ex:16:6}
\emph{分量を読む受け手。}受容体は疲れる：$y=[c-a]_+$、$\tau_a\,\dd a/\dd t=c-a$、$\tau_a=30$~min。ホルモンは$6$~minずつの分量で周期$T$で届き、平均投与量$\bar c$は同じとする。$T=15$、$60$、$120$~minおよび$T\to\infty$での平均応答を求めよ。濃度が一定値$\bar c$のときはどうか。
\end{exercise}

\begin{exercise}[\lvlC]
\label{ex:16:7}
\emph{発振器か、フィードバックか。}カスケード~\eqref{eq:cascade}の遅れるフィードバックが生む脈動と、独立したペースメーカーによる脈動とを区別する実験を提案せよ。$K$を1.5倍しか変えられず、周期の測定精度が$5\%$のとき、周期のずれから二つの仮説を区別できるか。
\end{exercise}

\chapter{マシンのデバッグ}
\chaptermark{マシンのデバッグ}
\label{sec:debugging}

\section{回路図のないマシンをどうデバッグするか}

回路図なしで動作中の基板を渡されたエンジニアは、四つの手法でデバッグする。\emph{プローブ}を当てて配線に何が流れているかを見る。\emph{テスト信号を入れて}何が変わったかを見る。回路の一部を\emph{切り離して}何が動かなくなったかを見る。そして\emph{制御レジスタを読んで}、基板がどのモードで動いているかを知る。本書全体が脳の回路図を読む試みである。本章では、エンジニアが四つの手法のうちどれをすでに使えるのか、そして前章までの内容から何が期待できるのかを見る。

今日このようなデバッグの最も目立つ例は、ブレイン・コンピュータ・インターフェースである。ニューロンのスパイクを読み取って外部の機械への指令に変える装置、あるいは逆に外から脳へ信号を入れる装置である。本章の大部分はこれを扱い、とくにNeuralink社の取り組みを取り上げる。

\section{プローブ：電極には何が聞こえるか}

脳のプローブは、組織に刺し込んだ細い電極である。電極には、半径およそ100マイクロメートル以内にあるニューロン、ふつうは1個から数個のスパイクが聞こえる。電極上のスパイクは、数十から数百マイクロボルト、幅約1ミリ秒のパルスであり、その波形から隣り合うニューロンを区別できる。最もよく聞こえるのは大きな\nt{GLUT}ニューロンである。皮質ではこれが多数派であり（表~\ref{tab:types}）、電流も強い。

すぐに主要な困難が見えてくる。今日の優れたプローブは数百から数千の電極をもつが、脳のニューロンは$8.6\cdot10^{10}$個ある。盗聴しているのはマシンのおよそ1億分の1にすぎない。これほどわずかな標本から、なぜ何かがわかるのか。答えは第\ref{sec:code}章にある。隣り合うニューロンのノイズは相関しており、群での平均化には限界がある（命題~\ref{prop:pooling}）。100個のニューロンは1000個とほぼ同じだけの情報を運ぶ。さらに、運動皮質が解く課題は低次元である。腕の関節は多くなく（第\ref{sec:muscles}章）、腕を制御する数百個のニューロンの活動は、わずか10--20個の共通変数の空間に収まる~\cite{Gallego2017}。100ニューロンのプローブは、これらの変数のほぼすべてを聞き取れる。もし脳が各ニューロンに独立したメッセージを格納していたら、このような盗聴は役に立たなかっただろう。

\section{データストリーム：プローブ上での圧縮}

プローブが出すデータは膨大である。スパイクの波形を見るには、各電極を毎秒約$2$万回、約10ビットの精度でデジタル化する必要がある。1000個の電極では
\begin{empheq}[box=\keybox]{equation}
\begin{aligned}
R_{\text{raw}} \;&=\; N\,f_s\,b \;\approx\; 10^3\cdot2\cdot10^4\cdot10 \;\approx\; 2\cdot10^8\ \text{bit/s},\\
R_{\text{spk}} \;&\approx\; N\,r\,\log_2\frac{e}{r\,\Delta t} \;\approx\; 8\cdot10^4\ \text{bit/s}.
\end{aligned}
\label{eq:probedata}
\end{empheq}
となる。第2の見積もりは、$r=10$スパイク/s、精度$\Delta t=1\ms$でのスパイク列の容量~\eqref{eq:spikeentropy}である。スパイク列に実際に含まれるものすべてが、2500分の1の量に収まる。埋め込み型装置にとってこれは決定的な違いである。頭の中で許される電力では、組織を加熱せずに生のデータを皮膚越しに無線で送ることはできない。したがって埋め込み型装置は、電極のそばのチップ上でスパイクを検出し、イベント、すなわちチャネル番号とスパイクの時刻だけを外に送る必要がある。Neuralinkのシステムの最初の報告では、まだそこまで到達していなかった。チップは信号を増幅してデジタル化し、生のデータはすべてケーブルで外へ送られ、スパイクは外部のプログラムが検出していた。チップ上の圧縮と無線接続は計画として挙げられていた~\cite{Musk2019}。

これはおなじみの手法である。眼は長いケーブルに送る前に信号を100分の1に圧縮する（第\ref{sec:sensors}章）。イベントカメラはフレームではなく変化を送る。配線に収めるために、プローブは脳自身と同じことをせざるをえない。スパイクで、しかも新しいことだけを伝えるのである。

\section{Neuralinkは何をしているか}

Neuralinkの装置は、これらの制約に合わせて組み立てたプローブである~\cite{Musk2019}。硬い針の配列の代わりに、髪の毛より細い柔軟なスレッドを多数使い、各スレッドに沿って電極を並べる。最初の報告では$96$本のスレッドに最大$3072$個の電極、1本あたり$32$個であった。同社によれば、ヒトに埋め込んだ装置は$64$本のスレッドで約1000チャネルである。柔軟なスレッドは組織の損傷が少なく、頭蓋内で脳が常にわずかに動くことにもよく耐える。スレッドはロボットが血管を避けながら挿入する。前述のとおり、最初の報告では生のデータ~\eqref{eq:probedata}はケーブルで外へ送られていた。同社によれば、ヒト用の装置は別の構成である。筐体は完全に皮膚の下にあり、スパイクは内部で検出され、データは無線で送られ、電力はワイヤレスで供給される。

最初の装置の目的は読み取りである。スレッドを腕の運動を計画する皮質領域に置き、デコーダがスパイクを画面上のカーソルの動きに変える。腕が麻痺した人が、動きを想像してコンピュータを操作する。発想自体は新しくない。100電極の硬い配列を使うインターフェースは、以前からカーソル操作~\cite{Hochberg2006}、手書きの動きを想像して文字を書くこと~\cite{Willett2021}、さらには話そうとする試みを毎分約60語の速さで画面上のテキストに変換すること~\cite{Willett2023}を可能にしてきた。Neuralinkの新しさはエンジニアリングにある。チャネル数、ワイヤレス化、密閉筐体、ロボットによる挿入、すなわち研究室の外で何年も装着できるプローブを作る試みである。我々の知る限り、ヒトでの試験結果は主に同社自身が公表しており、査読付き論文ではない。同社は特に、挿入後に一部のスレッドがずれ、動作するチャネルの数が減ったと報告している。デバッグではよくある不具合である。基板に対して動くプローブは、もう別の配線を聞いている。

同社の第二の方向は書き込みである。視力を失った人の視覚皮質に信号を入れる視覚用の装置である。これについては後の書き込みの節で述べる。

\section{読み取り：受け手としてのデコーダ}

インターフェースのデコーダは、命題~\ref{prop:reader}の意味での受け手である。メッセージのどの部分を読むかは自分で決める。実際には、ほとんどすべてのインターフェースが\emph{群の発火率}を読む。各チャネルのスパイクを数十ミリ秒の窓で数え、意図した動きが得られるように選んだ重みで足し合わせる。これは第\ref{sec:rate}節の群のレート符号であり、偶然選ばれたのではない。窓あたり数スパイクでは1個のニューロンの発火率は決まらないが、100個の和ならもう使える。

どれだけの情報を取り出せるか。「想像上の手」による書字は毎分約90文字であった~\cite{Willett2021}。つまり毎秒1.5文字であり、1文字5ビットとしても毎秒8ビット未満である。Neuralinkによれば、カーソル操作は毎秒約10ビットを与える。一方、上限~\eqref{eq:spikeentropy}は\emph{1個}のニューロンで毎秒数十ビット、100個なら数千ビットである。インターフェースは、盗聴したニューロンが運べるはずの情報のごく一部しか取り出していない。ボトルネックは配線ではない。群が運ぶ独立変数の数（命題~\ref{prop:pooling}）と、第\ref{sec:code}章で見たように完全には解読されていない符号をデコーダがどれだけ理解しているかである。

第\ref{sec:motorlearning}章でおなじみの第二の特徴もある。デコーダと脳は互いから学ぶ。脳はカーソルがどこへ動いたかを見て誤差を受け取り、指令を調整する。これは誤差配線による運動学習と同じである。一方エンジニアは、スレッドの下のニューロンが変わり、ネットワークの結合が学習し直されるので、ときどきデコーダの重みを選び直す。こうして互いに合わせ合う二人の学習者ができる。この組が暴走しないようにするには、学習速度を分けておくとよい。一方は速く、他方は遅く学ぶようにする。

\begin{keyblock}
\begin{proposition}[1000ニューロンのプローブはなぜ働くか]
\label{prop:probe}
ニューロンのごく一部を盗聴するだけのインターフェースが動きを制御できるのは、群の活動が低次元で、そのノイズが相関しているからである。数百個のニューロンが共通変数のほぼすべてを運ぶ。同じ理由で、インターフェースが取り出すのは毎秒数ビットにすぎない。これはスパイク列の容量ではなく、課題の独立変数の数で決まる。インターフェースの性能は測定済みである。符号が理解されたときデコーダがどれだけ良くなるかは未解決の問題である。
\end{proposition}
\end{keyblock}

\section{書き込み：読み取りより難しい}

マシンに信号を入れるのは、読むより難しい。電流を流す電極は1個のニューロンではなく、周囲のすべてのニューロンと配線を、種類も符号も区別せずに興奮させる。\emph{どの}ニューロンが\emph{いつ}発火したかが重要な符号（第\ref{sec:code}章）を、これで書き込むことはできない。大まかに\emph{どこで}、\emph{どれだけ強く}を指定できるだけである。

視覚にはこれで部分的に足りる。視覚皮質の電極に電流を流すと、人は視野の特定の位置に光る点を見る。最大16個の点を同時に点灯させると、視力を失った人がいくつかの文字を認識し、物体の境界を区別できた~\cite{Fernandez2021}。これは場所による符号であり、書き込むことができる。しかし第\ref{sec:sensors}章を思い出してほしい。眼が脳に送るのはピクセルではなく、圧縮された記述である。エッジ、変化、明るくなることと暗くなることが、別々の配線で送られる。皮質に「光の点」を書き込むのは、入力にまったく別の符号を待っているマシンにピクセルを書き込むことである。だから人工視覚は、電極の数から期待されるより今のところ粗い。電極の要らないテスト信号もある。錯視は眼を通して普通でない入力を与え、マシンを通常のモードから外す。そして誤りの一つ一つが、それぞれの機構を指し示す（第\ref{sec:illusions}節）。

\section{プローブに見えないもの}

電極が聞くのはアドレスバス、つまり\nt{GLUT}ニューロンのスパイクと、それより聞き取りにくい小さな\nt{GABA}ニューロンのスパイクである。しかし第\ref{sec:serocomputer}--\ref{sec:acet}章で見たように、マシン全体の動作モードはブロードキャストレジスタ、すなわち\nt{SERO}、\nt{DOPA}、\nt{NORA}、\nt{ACET}の濃度が決める。電極にはこれが聞こえない。発信源のニューロンはごくわずかであり、その信号はプローブのそばのスパイクではなく、体積中の濃度だからである。データバスは見えても制御レジスタが見えないデバッガは、いつまでも驚き続けるだろう。同じ入力が異なる応答を生む。どこかで$g$、$a$、$\tau_\gamma$が変わったからである。濃度は組織内の化学センサで直接測れるが~\cite{Bunin1998}、インターフェースではそうしたセンサはまだ使われていない。マシンの第二の遅い出力、すなわち血中のホルモン（第\ref{sec:chemoutput}章）も、プローブの視野のまったく外にある。これは電極ではなく採血で測る。

プローブには重みも見えない。そして記憶と学習したことのほぼすべては、まさに重みに格納されている。重みは応答の変わり方から推測するしかない。また、プローブには人が感じるような痛みも見えない。第\ref{sec:pain}章で見たように、警報信号の強さはマシン自身、すなわち脊髄のゲートと上位の調節器が決めるので、センサからどれほど痛いかは読み取れない。

\section{まとめ：デバッグと未解決の問題}

\begin{table}[t]
\centering
\footnotesize
\setlength{\tabcolsep}{4pt}
\renewcommand{\arraystretch}{1.15}
\begin{tabular}{>{\raggedright}p{2.2cm}>{\raggedright}p{3.6cm}>{\raggedright}p{3.4cm}>{\raggedright\arraybackslash}p{4.0cm}}
\toprule
デバッグ手法 & インターフェースがすること & 本書による説明 & わかっていること \\
\midrule
プローブ & 数百--数千個のニューロンを聞く & 低次元性と相関したノイズ（第\ref{sec:code}章） & 測定済み \\
プローブ上での圧縮 & 生の信号の代わりにイベントを送る & スパイク列の容量~\eqref{eq:spikeentropy} & 工学的に解決済み \\
読み取り & 群の発火率$\to$動き、文字、発話 & デコーダは受け手、群のレート符号 & 動作する。毎秒数ビット \\
共同学習 & 脳とデコーダが互いに合わせ合う & 誤差配線（第\ref{sec:motorlearning}章） & 観察されている。調整の規則は研究課題 \\
書き込み & 電流が視野に点を灯す & 場所の符号は書き込めるが、時間の符号は書けない（第\ref{sec:sensors}章） & 点は測定済み。実用的な視覚はまだ \\
レジスタの読み取り & ほとんど行われていない & ブロードキャストチャネル（第\ref{sec:serocomputer}--\ref{sec:acet}章） & 化学センサは実験段階 \\
\bottomrule
\end{tabular}
\caption{マシンのデバッグ：ブレイン・コンピュータ・インターフェースにできることと、本書の各章がそれについて語ること。}
\label{tab:debugging}
\end{table}

表~\ref{tab:debugging}が本章をまとめる。ブレイン・コンピュータ・インターフェースは、すでにアドレスバスにプローブを当て、毎秒数ビットを読み取れる。それがそもそも働くのは、第\ref{sec:code}章の低次元性と相関したノイズのおかげである。マシンへの書き込みは粗くしかできない。マシンの入力符号は圧縮され、個々の配線にアドレス付けされているからである。そして制御レジスタと重み、すなわちマシンを学習可能で調整可能にしているものは、プローブにはまだほとんど見えない。

第\ref{sec:synthesis}章の未解決の問題は、どれもデバッガへの課題でもある。どの符号を読むべきかは、プローブがより高密度になり、デコーダが発火率だけでなくスパイクのタイミングを使えるようになったときに決着する。多層ネットワークがどう学習するかは、複数の層に同時にプローブを置き、学習中に応答がどう変わるかを見れば確かめられる。ブロードキャストチャネルが何を伝えているかは、電気的なプローブに化学的なプローブが加わったときに見えてくる。本書は、マシンの回路図をその振る舞いと、どのエンジニアも知っている法則から読み取る試みである。いま作られているデバッガが、この回路図をプローブで確かめることを可能にするだろう。

\section{問題}

\begin{exercise}[\lvlA]
\label{ex:17:1}
\emph{無線チャネルの予算。}皮膚越しの伝送は1ビットあたり$1$~nJかかり、送信機に使える電力は$10$~mW以下である。$N=1000$電極の生データ~\eqref{eq:probedata}とスパイク列を送るのに必要な電力はそれぞれいくらか。生データなら1ビットあたり何ジュールに抑える必要があるか。
\end{exercise}

\begin{exercise}[\lvlA]
\label{ex:17:2}
\emph{100ニューロンに何ビットあるか。}デコーダは$N$個のニューロンのスパイクを$50\ms$の窓で足し合わせる。$r=20$~スパイク/s、ノイズの対相関$c=0.1$、信号は群の発火率を（二乗平均で）$30\%$変える。$N=10$、$100$、$1000$、$N\to\infty$での窓あたりのレート$\tfrac12\log_2(1+\text{SNR})$を見積もれ。$c=0$、$N=1000$ではどうか。
\end{exercise}

\begin{exercise}[\lvlB]
\label{ex:17:3}
\emph{キーボード型インターフェースの速度。}インターフェースは毎秒1.5回、$N=32$文字から1文字を選ぶ。意図した文字を選ぶ確率は$P$、誤りは等確率とする。$P=0.99$、$0.94$、$0.8$で毎秒何ビット伝えるか。$P=0.94$で$10$~bit/sを得るには毎秒何文字必要か。
\end{exercise}

\begin{exercise}[\lvlB]
\label{ex:17:4}
\emph{シミュレーション：群のデコーダ。}$N$個のニューロンで$\lambda_i=[b+m\,\mathbf v\cdot\mathbf p_i]_+$、$b=20$、$m=15$~スパイク/s、窓$50\ms$、$\mathbf v$の各成分の標準偏差$0.5$とする。全ニューロンが共通ノイズ入りの$\mathbf v+\boldsymbol\xi$、$\sigma_\xi=0.2$を見る。群ベクトルによる不偏デコーダをシミュレートせよ。二つのノイズが等しくなる$N$はいくつか。$N\to\infty$で窓あたり成分あたり何ビット得られるか。
\end{exercise}

\begin{exercise}[\lvlB]
\label{ex:17:5}
\emph{二人の学習者。}脳は指令$m=b\,v^*$を、デコーダは$v=d\,m$を出す。$\langle v^{*2}\rangle=1$。両者は試行ごとに$\tfrac12(v-v^*)^2$について学習率$\eta$で勾配を1ステップ下る。$b=1.2$、$d=1$から始め、$\eta=0.8$、$1.1$、$1.5$、$2$でシミュレートせよ。それぞれ単独なら$\eta<2$で安定である。組が安定なのは$\eta$がいくつのときか。
\end{exercise}

\begin{exercise}[\lvlB]
\label{ex:17:6}
\emph{設計：プローブのパイプライン。}$1024$チャネル、毎秒$2$万サンプル、ニューロンは$10$~スパイク/s、無線は$1$~nJ/bitとする。サンプルの白色ノイズに対してどの閾値なら、偽スパイクがチャネルあたり$0.1$~Hz以下になるか。$20\ms$ごとのスパイク数を$2$ビットで送るとき、電力はいくらで、$10$ビットの生サンプルの何分の1か。
\end{exercise}

\begin{exercise}[\lvlC]
\label{ex:17:7}
\emph{インターフェースの速度限界。}帯域$W=5$~Hzの$D=10$変数について、$\text{SNR}\approx0.9$（信号に似たノイズ）と$90$（$1000$ニューロンの独立ノイズ）の場合に$R\le DW\log_2(1+\text{SNR})$を見積もれ。測定値は約$10$~bit/sである。このギャップの二つの説明、信号に似たノイズと符号の無知とを区別する実験は何か。
\end{exercise}

\chapter{電気的機能によるニューロンの分類}
\chaptermark{素子としてのニューロン}
\label{sec:eetypes}
\begin{chaptergoals}
\item ニューロンがどんな素子として働くか：積分器、フィルタ、変換器、発振器、ラッチ；
\item 1個の細胞がバンドパス共振器や、\nt{SERO}がオンにするラッチになる仕組み；
\item 漏れのあるニューロンからフィードバックで作る漏れのない積分器と、精密な調整が要る理由；
\item なぜこれらは型でなくモードであり、イオンチャネルとブロードキャストチャネルが選ぶのか。
\end{chaptergoals}

第\ref{sec:neurontypes}章では、出力が放出する神経伝達物質によってニューロンを分類した。電子技術者なら別の分け方をするだろう。ニューロンが自分の信号に何をするか、つまりどんな素子として働くかによる分類である。本書の主要な素子は第\ref{sec:glutcomputer}章から知っている。漏れ積分ニューロン~\eqref{eq:glutneuron}、すなわちコンパレータ付きのRC回路である。しかし脳には他の素子もある。表~\ref{tab:eetypes}にそれらを四つのグループにまとめた。ほとんどはすでに本書に登場している。

\begin{center}
\footnotesize
\setlength{\tabcolsep}{3pt}
\renewcommand{\arraystretch}{1.15}
\begin{longtable}{>{\raggedright}p{3.1cm}>{\raggedright}p{3.3cm}>{\raggedright}p{4.7cm}>{\raggedright\arraybackslash}p{2.4cm}}
\caption{素子としてのニューロン：名前、電子回路での対応物、素子の働き、本書での登場箇所。}\label{tab:eetypes}\\
\toprule
ニューロン & 素子 & 働き & 本書での箇所 \\
\midrule
\endfirsthead
\toprule
ニューロン & 素子 & 働き & 本書での箇所 \\
\midrule
\endhead
\bottomrule
\endfoot
\multicolumn{4}{l}{\emph{フィルタと積分器}} \\
漏れ積分ニューロン & コンパレータ付きRC積分器 & 窓$\tau$内の入力を足し合わせ、閾値で発火する & 第\ref{sec:glutcomputer}章、\eqref{eq:glutneuron} \\
同時検出器 & 時間窓付きAND & 入力がほぼ同時に届いたときだけ発火する。$\tau$がごく小さい & 第\ref{sec:glutcomputer}、\ref{sec:code}章、\eqref{eq:window} \\
相動性ニューロン & ハイパスフィルタ、エッジ検出器 & 入力の変化に応答し、その後沈黙する & 第\ref{sec:sensors}章 \\
順応するニューロン & 自動ゲイン制御付きハイパスフィルタ & 一定の入力では発火率が下がる & 第\ref{sec:sensors}章、\eqref{eq:nakarushton} \\
共振器 & 共振回路、バンドパスフィルタ & 自分の周波数の入力に最もよく応答する & 第\ref{sec:resonator}節 \\
漏れのない積分器 & オペアンプ積分器 & 速度を位置に変えて保持する & 第\ref{sec:perfectint}節 \\
\midrule
\multicolumn{4}{l}{\emph{変換器}} \\
持続性ニューロン & 電圧--周波数変換器 & 発火率が入力に線形に追従し、ほとんど疲れない & 第\ref{sec:code}章 \\
段階的電位ニューロン（スパイクなし） & アナログ増幅器、バッファ & 滑らかな電圧を短い距離だけ伝える & 第\ref{sec:sensors}章 \\
整流ニューロン & 半波整流器 & 発火率は負にならない、$[u]_+$ & 第\ref{sec:glutcomputer}、\ref{sec:gaba}章 \\
オン--オフの対 & 整流器の対、2線式符号 & 一方は「増加」に、他方は「減少」に応答する & 第\ref{sec:glutcomputer}章、\eqref{eq:dualrail} \\
シャント抑制を受けるニューロン & アナログ除算器、ゲイン制御 & 入力から引くのではなく、入力を割る & 第\ref{sec:gaba}章、\eqref{eq:shunt} \\
\midrule
\multicolumn{4}{l}{\emph{発振器と時間}} \\
ペースメーカー & 弛張発振器、マルチバイブレータ & 一定の頻度で自らスパイクを出す & 第\ref{sec:serocomputer}、\ref{sec:dofa}章 \\
入力付きペースメーカー & 電圧制御発振器 & 固有のリズムを入力が速めたり遅めたりする & 第\ref{sec:chemoutput}章 \\
バースト発火ニューロン & ゲート付きバースト発生器 & 高頻度スパイクのバースト、つまり「ツー」を出す & 第\ref{sec:code}章、\eqref{eq:burst} \\
位相固定ニューロン & 位相検出器、位相同期ループ & 入力波の特定の位相で発火する & 第\ref{sec:sensors}、\ref{sec:chemoutput}章 \\
遅延のある軸索 & 遅延線 & 信号を決まった時間だけ遅らせる & 第\ref{sec:glutcomputer}章、図~\ref{fig:itd} \\
\midrule
\multicolumn{4}{l}{\emph{記憶と切り替え}} \\
双安定ニューロン & シュミットトリガ、ラッチ & 短い入力で長く続く状態に移る & 第\ref{sec:bistable}節 \\
樹状突起の枝をもつニューロン & マルチプレクサ、小さな多層ネットワーク & 許可入力があるときだけ枝が入力を通す & 第\ref{sec:synthesis}章 \\
\end{longtable}
\end{center}

表全体より重要な但し書きが一つある。これらは別々の細胞ではなく、別々の\emph{モード}である。同じニューロンが、入力が遅ければ積分器になり、抑制が$\tau$を短くすれば同時検出器になる（第\ref{sec:code}章）。覚醒中はペースメーカーで、睡眠中は黙った受信器になる。どの素子になるかは、膜のイオンチャネルと、それを調整するブロードキャストチャネルが決める。

\section{本書にまだ出ていない四つの素子}

\subsection{共振器}
\label{sec:resonator}

一部のニューロンには、電圧のあらゆる変化に逆らう遅い電流がある。電圧が上がれば電流はそれを引き下げ、下がれば引き上げる。ただし遅れを伴う。電子技術者から見ると、この電流はインダクタンスのように振る舞い、膜容量と合わせて共振回路を作る。ニューロンは特定の周波数、ふつう数Hzから十数Hzの入力に最も強く応答し、それより遅い入力にも速い入力にも弱く応答する~\cite{HutcheonYarom2000}。これは1個の細胞の中のバンドパスフィルタであり、ノイズの多い入力から自分の周波数のリズム、たとえば第\ref{sec:gaba}章の局所リズムを選び出す。

\subsection{漏れのない積分器}
\label{sec:perfectint}

漏れ積分器が入力を覚えているのは時間$\tau$、つまり数十ミリ秒だけである。しかし眼を静止させておくには、自分の位置を数秒間覚えている必要がある。「回転せよ」という指令は短時間の速度として来るが、眼は新しい位置に保たなければならない。ネットワークはこれを、第\ref{sec:glutcomputer}章の記憶と同じフィードバックで解く。時定数$\tau$のニューロン群が重み$w$で自らを興奮させるなら、$\tau\,\dd r/\dd t=-r+w\,r+I$であり、群の時定数は
\begin{empheq}[box=\keybox]{equation}
\tau_{\text{eff}} \;=\; \frac{\tau}{1-w} .
\label{eq:perfectint}
\end{empheq}
となる。$\tau=0.1\,\text{s}$、$w=0.995$なら$20\,\text{s}$になる。個々には10分の1秒で忘れるニューロンから作った、ほぼ理想的な積分器である~\cite{Seung1996}。代償は精密な調整である。$w$が1をわずかに超えれば積分器は飽和し、わずかに下回ればすぐに忘れる。そのため、このような積分器には第\ref{sec:motorlearning}章で述べたような学習による絶え間ない調整が必要である。

\subsection{双安定ニューロン}
\label{sec:bistable}

第\ref{sec:glutcomputer}章と第\ref{sec:gaba}章では、フリップフロップをニューロン群で組み立てた。しかし1個の細胞の中のフリップフロップもある。一部のニューロンには、いったんオンになると自ら興奮を維持する電流がある。短い入力でニューロンは長く続く発火状態、すなわち\emph{プラトー}に移り、抑制で静止状態に戻る。これはシュミットトリガである。ニューロンは二つの安定状態とその間のヒステリシスをもつ。たとえば筋肉を制御する\nt{ACET}ニューロンがこう振る舞い、この性質はブロードキャストチャネルでオンになる。セロトニンがニューロンに届くとプラトーが現れる~\cite{Hounsgaard1988}。同じ細胞が、\nt{SERO}があればラッチ、なければ普通の積分器である。

\subsection{積分器型と共振器型}

理論は興奮性細胞を二つのクラスに分ける~\cite{Izhikevich2007}。\emph{積分器型}は、ゼロから始まる電圧制御発振器に似ている。閾値をわずかに超えたニューロンはいくらでもゆっくり発火でき、発火率は入力とともに滑らかに上がる。このようなニューロンは届いたものをすべて足し合わせ、入力の大きさを発火率でよく伝える。\emph{共振器型}は、最低周波数のある発振器に似ている。ゆっくり発火することはできず、閾値の入力ですぐに有限の頻度でスパイクを出し、自分の周波数の入力を好む。前者は第\ref{sec:code}章のレート符号の自然な素子であり、後者はタイミング符号の自然な素子である。

\begin{keyblock}
\begin{proposition}[1個のニューロンは多くの素子になる]
\label{prop:eetypes}
電気的機能から見ると、ニューロンは漏れのある積分器と漏れのない積分器、同時検出器、ハイパスフィルタとバンドパスフィルタ、電圧--周波数変換器、整流器、除算器、発振器、遅延線、フリップフロップとして働く。ある細胞がどの素子になるかは、そのイオンチャネルと、それを調整するブロードキャストチャネルが決める。モード自体は測定済みである。脳がこうした切り替えを計算にどこまで使っているかは、主に仮説である。
\end{proposition}
\end{keyblock}

エンジニアにとって、これはプログラマブルロジックアレイに似ている。ハードウェアは一つで、素子の働きは構成で決まる。ただしここでは構成を書き込み時にではなく動作中に変える。変えるのは、第\ref{sec:serocomputer}--\ref{sec:acet}章で扱った遅いブロードキャストチャネルである。

\section{問題}

\begin{exercise}[\lvlA]
\label{ex:18:1}
\emph{漏れのない積分器の許容誤差。}$\tau=0.1$~sのニューロンが~\eqref{eq:perfectint}に従って眼の位置を$T=1$~sのあいだ、どちら向きにも$5\%$以内の誤差で保つ。(a)~$w$の許容範囲を求めよ。(b)~$w$は$K$個のシナプスの重みの和で、各々に独立な$10\%$の誤差がある。和のばらつきが許容範囲に収まるのは$K$がいくつのときか。
\end{exercise}

\begin{exercise}[\lvlB]
\label{ex:18:2}
\emph{回路としての共振器。}第\ref{sec:resonator}節の遅い電流を変数$W$で表す：$C\,\dd V/\dd t=-g_LV-g_wW+I$、$\tau_w\,\dd W/\dd t=V-W$。これが並列枝$R_w=1/g_w$、$L_w=\tau_w/g_w$をもつ$RC$回路であることを示し、$C/g_L=10\ms$、$\tau_w=100\ms$、$\alpha=g_w/g_L=4$で共振周波数と、そこでの$|Z|$の$|Z(0)|$に対する利得を求めよ。
\end{exercise}

\begin{exercise}[\lvlB]
\label{ex:18:3}
\emph{積分器はどう発火し始めるか。}(a)~ニューロン$\tau\,\dd V/\dd t=-V+\mu$、$\tau=10\ms$、閾値$\theta$、ゼロへのリセット。発火率が$1$~Hzになるのは$\mu/\theta-1$がいくつのときか。(b)~モデル$\tau\,\dd v/\dd t=v^2+\mu$（スパイクは$v$の$+\infty$への発散、リセットは$-\infty$へ）は$\mu=0.01$でどの発火率を与えるか。
\end{exercise}

\begin{exercise}[\lvlB]
\label{ex:18:4}
\emph{シミュレーション：積分器型と共振器型。}問題~\ref{ex:18:2}のモデルで$g_L=1$、$\tau_m=10\ms$、$\tau_w=100\ms$、閾値$1$、リセット$V\to0$（$W$は変えない）とし、入力$\mu=1.5\sin2\pi ft$、$f=1$--$40$~Hzを与える。$\alpha=0$と$\alpha=4$で、ニューロンがスパイクを出す周波数帯はどこか。条件$1.5\,|Z(f)|>1$と比べよ。
\end{exercise}

\begin{exercise}[\lvlB]
\label{ex:18:5}
\emph{1個の細胞のラッチ。}プラトー電流をもつニューロン：$\tau\,\dd r/\dd t=-r+F(wr+I)$、$F(x)=\min(1,\max(0,x))$、$\tau=10\ms$、$w=1.5$、背景入力$I=-0.2$。(a)~パルス$I=0.3$がニューロンを$r=0$からプラトーに移すのに必要な最短の長さはいくらか。(b)~長いパルス$I=-0.2-B$がニューロンを静止状態に戻すのに必要な最小の振幅$B$はいくらか。
\end{exercise}

\begin{exercise}[\lvlB]
\label{ex:18:6}
\emph{積分器の自己調整。}注視中は入力$I=0$で、滑りセンサがずれの速度$e=\dd r/\dd t$を与え、$\dd w/\dd t=-\eta\,e\,r$とする。$w=0.95$、$\tau=0.1$~s、$\langle r^2\rangle=1$のとき、$60$~sの注視で$|1-w|$が問題~\ref{ex:18:1}の許容範囲に入る$\eta$を求めよ。眼の回転中にも学習すると何が壊れるか。
\end{exercise}

\begin{exercise}[\lvlC]
\label{ex:18:7}
\emph{なぜ素子を組み替えるのか。}ブロードキャストチャネルが問題~\ref{ex:18:2}のモデルの$\alpha$を$0$と$4$の間で切り替える。電流の白色ノイズ中の弱い正弦波について、$11$~Hzと$1$~Hzで共振器型は積分器型よりSN比で何倍有利か。モードの切り替えそのものが有利になる課題を考えよ。
\end{exercise}

\chapter{樹状突起の数によるニューロンの分類}
\chaptermark{樹状突起の数}
\label{sec:dendrites}

\section{回路パラメータとしての入力数}

第~\ref{sec:neurontypes}~章ではニューロンを出力が放出するものによって、第~\ref{sec:eetypes}~章ではどんな素子として働くかによって分類した。エンジニアがすぐに気づく第三の特徴もある。\emph{ニューロンの入力はいくつあるか}である。入力を集めるのは樹状突起、すなわち他のニューロンの軸索が接続する枝である（第~\ref{sec:dofa}~章、第~\ref{sec:weightstore}~節）。樹状突起がまったくないニューロンもあれば、1本か2本のものも、10万個の入力を集める木をもつものもある。電子技術者の言葉ではこれはファンインであり、ほとんどの場合、課題に見合っている。

樹状突起の数と大きさがなぜそれほど重要かは、簡単な見積もりからわかる。膜は比容量$c_m\approx1$~\textmu F/cm$^2$のコンデンサであり、ニューロンと樹状突起を合わせた面積$A$が大きいほど、電圧を$\Delta V$上げるのに運ぶべき電荷が多い。漏れを無視すれば、入力電流$I$がこれを行う時間は
\begin{empheq}[box=\keybox]{equation}
t \;\approx\; \frac{c_m\,A\,\Delta V}{I} .
\label{eq:dendriteload}
\end{empheq}
である。短い樹状突起が数本の小さなニューロンは、面積が約$2\cdot10^{-5}$~cm$^2$、容量が約$20$~pFである。数nAの電流を流す1個の大きなシナプスは、これを0.1ミリ秒未満で$20\mV$持ち上げる。大きな木をもつニューロンの面積は約15倍で、容量は約$300$~pFである。約$0.1$~nAの普通のシナプスなら充電に$60\ms$かかる。これは漏れの時定数よりはるかに長いので、実際には決して充電できない。このようなニューロンには数十個の同時入力が必要である。ここでの数値は桁の見積もりだが、結論はそれに左右されない。\emph{入力が少なければ速く正確、多ければ遅く総和で働く}。

\begin{figure}[t]
\centering
\begin{tikzpicture}[>=Stealth,
  soma/.style={circle,draw,very thick,fill=blue!6,minimum size=0.55cm,inner sep=0pt},
  ax/.style={->,very thick,red!70!black},
  den/.style={very thick,blue!60!black}]
% 0 dendrites: sensor on a line
\begin{scope}[xshift=0cm]
\draw[den,decorate,decoration={snake,amplitude=1pt,segment length=4pt}] (-0.9,0) -- (0,0);
\node[font=\scriptsize] at (-0.9,0.3) {センサ};
\draw[ax] (0,0) -- (1.0,0);
\draw[thick] (0.1,0) -- (0.1,0.45);
\node[soma,minimum size=0.4cm] at (0.1,0.65) {};
\node[font=\scriptsize,align=center] at (0.05,-0.75) {$0$\\センサ付き線路};
\end{scope}
% 1 dendrite
\begin{scope}[xshift=3.3cm]
\draw[den] (-0.9,0) -- (-0.28,0);
\node[soma] at (0,0) {};
\draw[ax] (0.28,0) -- (0.9,0);
\node[font=\scriptsize,align=center] at (0,-0.75) {$1$\\バッファ、中継器};
\end{scope}
% 2 dendrites: one per ear
\begin{scope}[xshift=6.2cm]
\draw[den] (-0.28,0) -- (-0.9,0); \draw[den] (0.28,0) -- (0.9,0);
\node[soma] at (0,0) {};
\draw[ax] (0,-0.28) -- (0,-0.6);
\node[font=\scriptsize] at (-0.9,0.25) {左}; \node[font=\scriptsize] at (0.9,0.25) {右};
\node[font=\scriptsize,align=center] at (0,-1.0) {$2$\\時間のAND};
\end{scope}
% ~4 short dendrites: mixer
\begin{scope}[xshift=9.0cm]
\foreach \a in {45,135,225,315}{\draw[den] (\a:0.28) -- (\a:0.7); \fill[blue!60!black] (\a:0.72) circle (0.06);}
\node[soma] at (0,0) {};
\draw[ax] (0.28,0) -- (1.0,0);
\node[font=\scriptsize,align=center] at (0,-1.0) {$\sim4$\\ミキサ};
\end{scope}
% many: summing tree
\begin{scope}[xshift=12.0cm]
\foreach \a in {60,90,120}{\draw[den] (0,0.28) -- ++(\a:0.55) coordinate (b); \foreach \d in {-20,20}{\draw[den] (b) -- ++(\a+\d:0.35);}}
\foreach \a in {-120,-60}{\draw[den] (0,-0.28) -- ++(\a:0.45);}
\node[soma] at (0,0) {};
\draw[ax] (0.28,0) -- (1.0,0);
\node[font=\scriptsize,align=center] at (0,-1.0) {$10^4$--$10^5$\\加算器};
\end{scope}
\end{tikzpicture}
\caption{入力数による五つのクラス。青は樹状突起（入力）、赤は軸索（出力）。入力が少ないほど、ニューロンは個々の信号を速く正確に伝える。多いほど、足し合わせと分類が得意になる。細胞の形の図ではなく模式図である。}
\label{fig:dendriteclasses}
\end{figure}

\section{樹状突起ゼロ：線路の端のセンサ}

触覚、痛み、筋伸張の感覚ニューロン（第~\ref{sec:sensors}、\ref{sec:pain}、\ref{sec:muscles}~章）には、細胞体に樹状突起が1本もない。軸索は細胞体から1本の小枝として出て、すぐに二つに分かれる。一方の枝は皮膚や筋肉へ向かい、その末端自体がセンサになる。もう一方は脊髄へ向かう。スパイクはセンサから細胞体を経ずにまっすぐ脊髄へ走る。エンジニアの目には、遠端にセンサがある通信線路であり、細胞体は電源ユニットのように脇にぶら下がっている。線路に給電するが、伝送には加わらない。利点は速さと信頼性である。経路上に、信号を足し合わせて送るかどうかを改めて判断する場所が一つもない。

光と音のセンサはさらに極端な例である。これらは入力を集めるニューロンではなく変換器そのものであり、入力はシナプスではなく受容体タンパク質である（図~\ref{fig:transducer}）。

\section{樹状突起1本：バッファ}

樹状突起1本と軸索1本をもつニューロンは、1本の入力線を受け取って先へ渡す。嗅覚ニューロンがそうである。唯一の樹状突起が鼻の表面へ伸び、ただ1種類の受容体をもつ~\cite{Mombaerts2004}。光センサと眼の出力ニューロンの間にある網膜の中継細胞も、耳のセンサと脳をつなぐニューロンもそうである。エンジニアならバッファ、あるいは中継器と呼ぶだろう。計算はせず、1個の信号を届ける。ときには形を変える。たとえば第~\ref{sec:sensors}~章のように、センサの滑らかな電圧をスパイクに変える。

\section{樹状突起2本：時間のAND}

音の方向を決める同時検出器（図~\ref{fig:itd}）は、哺乳類ではしばしば反対方向を向いたちょうど2本の樹状突起をもつ。一方には左耳からの入力が、他方には右耳からの入力が届く~\cite{Grothe2010}。なぜ入力を分けるのか。式~\eqref{eq:shunt}を思い出そう。膜の一か所で多くの興奮性チャネルが開いていると、そこの電圧は平衡電位に近づき、入力が一つ増えるごとに加わる分は小さくなる。したがって、片耳からの多くの入力を1本の樹状突起に集めるとそれは飽和し、それだけではニューロンを閾値に届かせられない。一方、各樹状突起からの1個ずつの入力はほぼ線形に足し合わされる。2本の樹状突起はニューロンをより厳密なANDに変える。信号が両側から同時に来たときだけ発火せよ、である。モデルでは、このように入力を分けることが実際に同時検出を改善することが示されている~\cite{AgmonSnir1998}。

\section{短い樹状突起が数本：ミキサと中継器}

\emph{ミキサ。}第~\ref{sec:motorlearning}~章の運動学習の回路では、約700億個の小さな\nt{GLUT}ニューロンが入力を非常に長いベクトルに拡張する。それぞれの樹状突起は短いものが約4本だけで、各々の先端は1個の入力の終末をつかむ「爪」になっている。つまり各ミキサは多くの入力のうち約4個しか見ず、その4個に対する小さなAND素子として働く。$M$本の入力線から選べる4個組は$\binom{M}{4}$通りあり、$M=100$ならほぼ400万である。なぜそれほど必要なのか。$n$入力の閾値素子が二つのクラスに正しく分けられるランダムなパターンは、たかだか
\begin{empheq}[box=\keybox]{equation}
P_{\max} \;\approx\; 2n
\label{eq:cover}
\end{empheq}
個である~\cite{Cover1965}。同じ回路の出力ニューロンはミキサから約$10^5$個の入力を受けるので、\eqref{eq:cover}によれば上限は約$2\cdot10^5$個の異なる対応である。これは理想的な素子の限界である。興奮性シナプスのように重みがすべて正で、ノイズへの余裕も必要なら、同じニューロンの見積もりは約$4\cdot10^4$個の対応になる~\cite{Brunel2004}。入力の少ないミキサが必要なのは、まさに大きな加算器に多数の異なる、ほぼ独立な入力を与えるためである~\cite{Marr1969,Albus1971}。

\emph{中継器。}耳から方向検出器への経路には、短い樹状突起を少数もつニューロンがある。これらは巨大なシナプスをわずか数個、あるものは実質的に1個だけ受け取る。\eqref{eq:dendriteload}によれば、これはまさに必要なものである。小さな容量と大きな電流により、ニューロンは入力スパイクごとに1ミリ秒未満で発火し、時間精度をほとんど失わない。そうした中継器の一つは\nt{GLUT}を受け取り\nt{GLYC}を出す。数十マイクロ秒の精度をもつインバータであり、方向検出器に速い抑制を供給する~\cite{BorstSoria2012}。

\section{数千の入力：加算器と分類器}

尺度の反対の端には、1万個の入力をもつ皮質の\nt{GLUT}ニューロンと、10万個の入力をもつ運動学習回路の出力\nt{GABA}ニューロンがある。\eqref{eq:dendriteload}によれば、こうしたニューロンは遅く、個々の入力には応答できない。足し合わせるのである。これは第~\ref{sec:glutcomputer}~章の積分ニューロンであり、分類器を組み立てる加算器である。大きな木は新しいものも加える。その枝はそれぞれ閾値をもつ別々のサブ回路として働くので、1個のニューロンが小さな多層ネットワークのように振る舞う~\cite{Beniaguev2021,Gidon2020}。

\section{出力でもある樹状突起}

軸索がまったくなく、樹状突起が入力と出力の両方を兼ねるニューロンがある。信号を受け取り、自ら神経伝達物質を放出する。最もよく調べられた例は、運動の方向の検出にかかわる網膜の\nt{GABA}ニューロンである（第~\ref{sec:gaba}~章、図~\ref{fig:direction}）。このニューロンでは、各枝が単独で細胞体から自分の先端へ向かう動きに応答し、その先端から抑制を出す~\cite{Euler2002}。1個のニューロンが、枝ごとに一つずつ、複数の独立した方向検出器になっている。エンジニアから見れば、これは1個の素子ではなく、1個のパッケージに複数のチャネルを収めたICである。

\section{まとめ}

\begin{table}[t]
\centering
\footnotesize
\setlength{\tabcolsep}{3pt}
\renewcommand{\arraystretch}{1.15}
\begin{tabular}{>{\raggedright}p{1.9cm}>{\raggedright}p{3.4cm}>{\raggedright}p{2.6cm}>{\raggedright}p{1.8cm}>{\raggedright\arraybackslash}p{3.8cm}}
\toprule
樹状突起 & 例 & 素子 & 入力数 & わかっていること \\
\midrule
$0$ & 触覚、痛み、伸張のセンサ & 端にセンサのある線路 & --- & 確立 \\
$1$ & 嗅覚ニューロン、網膜の中継細胞、耳のニューロン & バッファ、中継器 & $1$--少数 & 確立 \\
$2$ & 音の方向検出器 & 時間のAND & 2束 & 構造は確立。入力を分ける利点はモデルで示された \\
$\sim4$ & 運動学習回路のミキサ & 小さなAND素子 & $\sim4$ & 確立。入力拡張の役割は根拠の強い仮説 \\
少数、短い & 耳からの経路の中継器 & 中継器、インバータ & $1$--数個の巨大シナプス & 確立 \\
数千 & 皮質の\nt{GLUT}ニューロン、運動学習の出力ニューロン & 加算器、分類器 & $10^4$--$10^5$ & 確立。サブ回路としての枝は個別の例で測定 \\
出力を兼ねる樹状突起 & 網膜の方向\nt{GABA}ニューロン & 多チャネルIC & 多数 & 測定済み \\
\bottomrule
\end{tabular}
\caption{樹状突起の数によるニューロンの分類。}
\label{tab:dendrites}
\end{table}

\begin{keyblock}
\begin{proposition}[入力数は課題に従う]
\label{prop:dendrites}
個々の信号の速さと精度が必要な場所、すなわちセンサ、中継器、同時検出器では、ニューロンの樹状突起は少なく、入力も少なく、シナプスは大きい。小さな容量~\eqref{eq:dendriteload}はすばやく充電される。足し合わせと分類が必要な場所では、ニューロンは数千の入力をもつ大きな木をもつ。クラスの構造は測定済みである。計算上の説明は根拠が強いが、多くのクラスについては仮説のままである。
\end{proposition}
\end{keyblock}

表~\ref{tab:dendrites}は、これまでの二つの分類、すなわち神経伝達物質による分類（表~\ref{tab:types}）と素子による分類（表~\ref{tab:eetypes}）を補う。三つとも同じ素子を別々の側から見ている。何を出すか、何を計算するか、どれだけ聞くかである。

\section{問題}

\begin{exercise}[\lvlA]
\label{ex:19:1}
\emph{大きなニューロンに入力はいくつ必要か。}$C=300$~pF、$\tau_m=20\ms$のニューロンが、それぞれ$0.1$~nAの電流源であるシナプスを受け、閾値は静止電位より$20\mV$高い。(a)~そもそも閾値に達するには、$10\ms$で達するには、$5\ms$で達するには、同時に何個のシナプスが働く必要があるか。(b)~$C=20$~pFのニューロンは$4$~nAのシナプス1個でどれだけの時間で閾値に達するか。
\end{exercise}

\begin{exercise}[\lvlA]
\label{ex:19:2}
\emph{なぜ4なのか。}ミキサは$M=100$本から選んだ$k$本の線のANDで、パターンでは半数の線が活動し、ミキサは$7\cdot10^{10}$個ある。(a)~$k=2$、$4$、$40$でパターンに応答するミキサは何個か。(b)~$10^5$入力の出力ニューロンについて、ミキサは$M=100$本の直接の線に比べ対応の数~\eqref{eq:cover}を何倍にするか。
\end{exercise}

\begin{exercise}[\lvlB]
\label{ex:19:3}
\emph{$2n$はどこから来るか。}原点を通る超平面は、$n$次元空間の一般の位置にある$P$個の点を$C(P,n)=2\sum_{k=0}^{n-1}\binom{P-1}{k}$通りに二つのクラスに分ける。(a)~$P=2n$で$2^P$通りの分け方のちょうど半分が分離可能であることを示せ。(b)~$n=100$、$P=180$と$220$で分離可能な割合はいくらか。
\end{exercise}

\begin{exercise}[\lvlB]
\label{ex:19:4}
\emph{厳密なANDとしての2本の樹状突起。}樹状突起は漏れ$g_L$をもつ区画で、各入力はそこに電位$E_E$のコンダクタンス$g_0=0.5\,g_L$を開く。細胞体は$\kappa(V_1+V_2)$を見て、閾値は$\theta=1.1\,\kappa E_E$である。なぜ1本の樹状突起にいくつ入力があってもニューロンは発火しないのか。各樹状突起にいくつの入力が必要か。
\end{exercise}

\begin{exercise}[\lvlB]
\label{ex:19:5}
\emph{中継器の設計。}スパイクのジッタは$\sigma_t=\sigma_V/\dot V$で、$\dot V$は閾値付近での電圧の傾きである。ノイズ$\sigma_V=1\mV$で$\sigma_t\le20$~\textmu sが必要とし、漏れは無視する。$C=300$~pFのニューロンには$0.1$~nAの同期したシナプスが何個必要か。$C=20$~pFで$4$~nAのシナプス1個のニューロンのジッタはいくらか。
\end{exercise}

\begin{exercise}[\lvlB]
\label{ex:19:6}
\emph{シミュレーション：長い樹状突起の充電。}長さ$L=\ell/\lambda$、両端が閉じた受動ケーブル$\tau_m\partial_tV=\lambda^2\partial_x^2V-V+i/g$をシミュレートせよ（$40$区画、オイラー法）。遠端に短い電流パルスを入れる。$L=0.2$、$1$、$2$のとき、近端はいつ、同じ電荷を樹状突起全体に広げたときの電圧の何割まで上がるか。
\end{exercise}

\begin{exercise}[\lvlC]
\label{ex:19:7}
\emph{研究：最適な入力数。}容量$C=C_0+nc_s$、同時に働く入力は電流$I_s$のものが$m$個、ニューロンは$P\approx2n$個の対応を記憶する~\eqref{eq:cover}。$m$一定の場合と割合一定$m=\varphi n$の場合に、応答時間は$P$にどう依存するか。コスト基準を提案し、中継器と分類器それぞれの最適な$n$を求めよ。
\end{exercise}

\chapter{睡眠中にニューロンは何をしているか}
\chaptermark{睡眠}
\label{sec:sleep}

\section{睡眠は電源オフではなく、別のモードである}

電源の切れたコンピュータを見たエンジニアは、何もしていないと言うだろう。眠っている脳ではそうはいかない。睡眠中もニューロンは黙っていない。深い睡眠では皮質はスパイクを出し、レム睡眠では日中とほぼ同じだけ出す。変わるのはニューロンが\emph{働くかどうか}ではなく、\emph{どう}働くかである。睡眠はマシンのモードの集まりであり、それを切り替えるのは日中と同じブロードキャストチャネルである。

各モードについて「レジスタの状態」を書き出せる（表~\ref{tab:sleepmodes}）。日中は\nt{ACET}、\nt{NORA}、\nt{SERO}が働き、センサはつながり、筋肉は指令に従う。深いノンレム睡眠では三つのチャネルがすべて静まり、感覚器官からの入力は弱まる~\cite{Hasselmo1999}。皮質でのアセチルコリン放出は下がり、\nt{NORA}ニューロンと\nt{SERO}ニューロンは日中より低い頻度でスパイクを出す~\cite{Marrosu1995,AstonJonesBloom1981,McGintyHarper1976}。レム睡眠では奇妙な様子になる。\nt{ACET}は再び日中のレベルまで上がるが~\cite{Marrosu1995}、\nt{NORA}と\nt{SERO}はほぼ完全に沈黙する~\cite{AstonJonesBloom1981,McGintyHarper1976,JacobsAzmitia1992}。レム睡眠中は体の筋肉が切り離される。筋肉を制御する\nt{ACET}ニューロンが\nt{GABA}と\nt{GLYC}によって強く抑制されるからである~\cite{BrooksPeever2012}。第\ref{sec:gaba}章と同じ禁止を、出力全体に一度にかけたものである。

\begin{table}[t]
\centering
\footnotesize
\setlength{\tabcolsep}{4pt}
\renewcommand{\arraystretch}{1.15}
\begin{tabular}{>{\raggedright}p{3.1cm}>{\raggedright}p{2.9cm}>{\raggedright}p{3.2cm}>{\raggedright\arraybackslash}p{3.4cm}}
\toprule
& 覚醒 & ノンレム睡眠 & レム睡眠 \\
\midrule
\nt{ACET}（書き込み、第\ref{sec:acet}章） & 高い & 低い & 高い \\
\nt{NORA}（ゲイン、第\ref{sec:nora}章） & 中程度、バースト & 低い & ほぼ沈黙 \\
\nt{SERO}（第\ref{sec:serocomputer}章） & 中程度 & 低い & ほぼ沈黙 \\
センサからの入力 & 接続 & 弱まる & 弱まる \\
筋肉への出力 & 接続 & 弱まる & 抑制で遮断 \\
皮質の活動 & 一様 & 遅い振動：活動--静止 & 日中のように一様 \\
\bottomrule
\end{tabular}
\caption{三つのモードにおけるマシンのレジスタの状態。すべての行は測定済みである。\nt{ACET}、\nt{NORA}、\nt{SERO}のレベルは睡眠段階ごとの直接記録による~\cite{Marrosu1995,AstonJonesBloom1981,McGintyHarper1976}。}
\label{tab:sleepmodes}
\end{table}

\section{いつ眠るか：ヒステリシス付きリレー}

マシンがいつ睡眠に切り替わるかを決めるものは何か。二つの部分からなる単純な回路がよく働く~\cite{Borbely1982}。第一の部分は\emph{睡眠圧}$S$である。起きている間にたまり、眠っている間に解消される。これは日中に充電され、夜に放電されるコンデンサである。
\begin{empheq}[box=\keybox]{equation}
\frac{\dd S}{\dd t} \;=\; \frac{1-S}{\tau_r}\ \ \text{（覚醒）},\qquad
\frac{\dd S}{\dd t} \;=\; -\frac{S}{\tau_d}\ \ \text{（睡眠）} .
\label{eq:sleeppressure}
\end{empheq}
第二の部分は二つの閾値である。$S$が上の閾値$H(t)$に達するとマシンは眠りに落ち、$S$が下の閾値$L(t)$まで下がると目覚める。どちらの閾値も、第\ref{sec:chemoutput}章の概日リズムに合わせて滑らかに揺れる。こうしてヒステリシス付きリレー、すなわち第\ref{sec:bistable}節のシュミットトリガができる。コンデンサと合わせると、概日リズム~\eqref{eq:pll}に位相同期する弛張発振器になる。

\begin{figure}[t]
\centering
\begin{tikzpicture}[>=Stealth,x=0.16cm,y=4.2cm]
\foreach \a/\b in {0/4.3,21.2/29.3,45.4/53.4,69.5/72}{\fill[blue!8] (\a,0) rectangle (\b,1);}
\draw[->,gray!70!black] (0,0) -- (75,0) node[right,font=\small,black] {$t$, h};
\draw[->,gray!70!black] (0,0) -- (0,1.08) node[left,font=\small,black] {$S$};
\foreach \t in {0,24,48,72}{\draw[gray!60!black] (\t,-0.02) -- (\t,0.02); \node[below,font=\scriptsize] at (\t,0) {\t};}
\draw[thick,densely dashed,red!70!black] plot file {figures/sleep_H.dat};
\draw[thick,densely dashed,green!45!black] plot file {figures/sleep_L.dat};
\draw[very thick,blue!60!black] plot file {figures/sleep_S.dat};
\node[font=\scriptsize,red!70!black,anchor=west] at (73,0.72) {$H$：入眠};
\node[font=\scriptsize,green!45!black,anchor=west] at (73,0.24) {$L$：覚醒};
\node[font=\scriptsize,blue!60!black] at (25.25,0.93) {睡眠};
\node[font=\scriptsize,blue!60!black] at (49.4,0.93) {睡眠};
\end{tikzpicture}
\caption{$\tau_r=18.2$~h、$\tau_d=4.2$~hでの睡眠圧$S$~\eqref{eq:sleeppressure}。日中$S$は上の閾値$H$まで充電され、夜（網掛け）には下の閾値$L$まで放電される。二つの閾値は概日リズムとともに揺れる。マシンはおのずと約$16$時間の覚醒と$8$時間の睡眠のサイクルに落ち着く。}
\label{fig:twoprocess}
\end{figure}

我々のモデルでは、通常の値$\tau_r=18.2$~h、$\tau_d=4.2$~hで、マシンはおのずと$16.1$時間の覚醒と$8.0$時間の睡眠に落ち着く（図~\ref{fig:twoprocess}）。モデルは奇妙に見えることも説明する。40時間眠らずにいると睡眠圧は$0.90$に達するが、モデルでの回復睡眠は$8.8$時間にすぎず、ふだんより1日長く眠るわけではない。放電は指数関数的なので、高い電荷はすばやく抜ける。「借り」は睡眠の長さではなく深さで返される。

物理的に$S$の役割を果たすのは何か。有力な候補は、付録~\ref{sec:othermediators}の「負荷カウンタ」であるアデノシンである。その濃度は覚醒中に上がり、睡眠中に下がる~\cite{PorkkaHeiskanen1997,Dunwiddie2001}。睡眠圧がまさにアデノシンであり、それ以外の何物でもないというのは仮説である。

\section{ノンレム睡眠：弛張発振器としてのネットワーク全体}

深い睡眠では、皮質のニューロンは群全体で交互に働く。1秒に1回より少ない頻度で、数百ミリ秒のあいだ一斉にオンになり（\emph{活動状態}）、一斉に沈黙する（\emph{静止状態}）。この振動の周波数は$1$~Hz未満で、麻酔下では$0.3$--$0.4$~Hzが最も多い~\cite{Steriade1993}。この遅い振動は、すでに手元にある部品で組み立てられる。自己への強いフィードバックをもつ\nt{GLUT}ニューロン群、すなわち二つの安定状態をもつ第\ref{sec:glutcomputer}章の記憶を取り、第\ref{sec:muscles}章のリズム発生器のような遅い疲労を加える。
\begin{empheq}[box=\keybox]{equation}
\tau\,\frac{\dd r}{\dd t} = -r + F\bigl(w\,r + I - a\bigr),
\qquad
\tau_a\,\frac{\dd a}{\dd t} = -a + b\,r .
\label{eq:updown}
\end{empheq}
群はオンになり、働き、疲れる。疲労$a$が上の状態を食いつぶし、群は静止に落ちる。静止中に疲労が回復すると下の状態が消え、群は再びオンになる。睡眠圧と同じ弛張発振器ができるが、およそ8万倍速い。

\begin{figure}[t]
\centering
\begin{tikzpicture}[>=Stealth,x=2.9cm,y=1.0cm]
\begin{scope}[yshift=1.9cm]
\draw[->,gray!70!black] (0,0) -- (4.1,0);
\draw[->,gray!70!black] (0,0) -- (0,1.25) node[left,font=\small,black] {$r$};
\draw[thick,blue!60!black] plot file {figures/updown_sleep.dat};
\node[font=\scriptsize,anchor=west] at (4.15,0.5) {睡眠：$b=5$};
\end{scope}
\draw[->,gray!70!black] (0,0) -- (4.1,0) node[right,font=\small,black] {$t$, s};
\draw[->,gray!70!black] (0,0) -- (0,1.25) node[left,font=\small,black] {$r$};
\draw[thick,red!70!black] plot file {figures/updown_wake.dat};
\node[font=\scriptsize,anchor=west] at (4.15,0.5) {日中：$b=1$};
\foreach \t in {0,1,2,3,4}{\draw[gray!60!black] (\t,-0.05) -- (\t,0.05); \node[below,font=\scriptsize] at (\t,0) {\t};}
\end{tikzpicture}
\caption{二つのモードにおける同じ群~\eqref{eq:updown}：$w=5$、$I=2$、$\theta=2.5$、$\tau=10\ms$、$\tau_a=0.3$~s。疲労が強いと（$b=5$、ノンレム睡眠のように）、群は周期$1.1$~sで活動と静止の間を振動する（モデルの値。皮質では周期は1秒より長く、麻酔下では約$3$~s）。疲労を弱めると（$b=1$。アセチルコリンの作用がこれにあたる）、同じ群が一様に働く。}
\label{fig:updown}
\end{figure}

では、ネットワークを振動から一様な日中の動作に移すものは何か。アセチルコリンは、ニューロンで疲労を生む遅い電流を抑える~\cite{McCormick1992}。我々のモデルでは、疲労が強い$b=5$のとき群は周期$1.1$~sで振動し（測定された振動より速いが同じ桁である）、整数個の周期で平均すると時間の約$35\%$働く。疲労が弱い$b=1$では、同じ群が一様に働く（図~\ref{fig:updown}）。一つのレジスタ、つまり\nt{ACET}のレベルが、発振器を増幅器に切り替える。睡眠中は遅い振動の上に、毎秒$7$--$15$回のより速いリズムのバーストも乗る。これを生むのは皮質と中間の中継核との間の\nt{GLUT}--\nt{GABA}ループであり~\cite{SteriadeMcCormickSejnowski1993}、第\ref{sec:gaba}章のリズムの親戚である。

\section{リプレイ：1日の早送り}

第\ref{sec:acet}章で、ノンレム睡眠中には日中に一緒に働いたニューロンが再び一緒に、同じ順序で発火するのを見た。ここにエンジニアリング上の細部を一つ加えよう。リプレイは\emph{早送り}で行われる。日中に数秒かかった系列が、睡眠中には1秒の数分の1で再生される~\cite{Nadasdy1999}。海馬ではおよそ20倍速く~\cite{LeeWilson2002}、皮質では6--7倍速い~\cite{Euston2007}。しかも再生はでたらめな時刻に起こるのではない。ある領域でのリプレイの短いバーストは、活動と静止の振動や、皮質のリズムの速いバーストにタイミングを合わせる。この協調を人為的に強めると記憶が良くなる~\cite{Maingret2016}。これはもう偶然の一致ではなく、因果関係である。

なぜマシンは1日を早送りするのか。エンジニアは\emph{破滅的忘却}と呼ばれる困難を知っている。新しいことを次々に順番に学ぶネットワークは、古いことを壊してしまう。救いになるのは\emph{インターリーブ}である。新しいことを古いことと混ぜて、少しずつ、何度も学ぶ。二つの学習系の仮説~\cite{McClelland1995}によれば、速い記憶が1日をまるごと記録し、睡眠中にそれを少しずつ、混ぜながら、何度も遅い皮質に再生する。これは第\ref{sec:acet}章の「速いバッファ--遅いストレージ」の組と同じであり、第\ref{sec:motorlearning}章の速い学習者と遅い学習者の組とも同じである。リプレイとその協調は測定済みである。その目的が遅い記憶のインターリーブ学習だというのは、根拠の強い仮説である。

\section{重みの再正規化}

日中は、第\ref{sec:dofa}章のヘッブ則が主に結合を強める。強めるだけなら重みは増え続け、それとともにエネルギー消費が増え、感度が落ちる。すべての入力が強いニューロンは、入力を区別できなくなる。\emph{シナプス恒常性}仮説~\cite{TononiCirelli2014}によれば、睡眠が必要な理由の一つは重みを動作レベルに戻すことである。すべての結合が同じ比率で弱められるので、結合の\emph{比}、つまり学習した内容は保たれる。エンジニアにとってこれは規則~\eqref{eq:oja}のような正規化だが、動作中にではなく別の時間に行われる。

直接の測定はこれと矛盾しない。マウスでは、皮質シナプスの接触面積は睡眠後に覚醒後より約$18\%$小さく、減少は大きさに比例し、最も大きく安定した接触はほとんど変わらない~\cite{deVivo2017}。重みのスケーリングは測定済みである。それが睡眠の主な役割だというのは仮説である。

\section{レム睡眠：入力と出力から切り離されたマシン}

レム睡眠では皮質はほぼ日中と同じように働くが、感覚器官からの入力は弱まり、筋肉への出力は抑制で遮断されている。エンジニアにはおなじみのモード、\emph{ベンチ試験}である。回路は全力で動くが、入力は内部の信号発生器につながれ、出力は何も壊さないようにアクチュエータから外されている。ある仮説によれば、夢はまさにこうした、日中に組み立てた世界モデルの内部での試運転である。そのときネットワークが正確に何をしているのか、学習しているのかどうかはわかっていない。ここで測定されているのは表~\ref{tab:sleepmodes}に書いたことだけである。どのチャネルがオンで、どれが沈黙し、出力が遮断されていること。

\section{保守と局所睡眠}

睡眠中は細胞間のすき間が広がり、脳から代謝産物がより速く洗い流されると長く考えられてきた~\cite{Xie2013}。別の方法で洗浄を測定した最近の実験は逆の結果を得た。睡眠中は洗浄が遅いというのである~\cite{Miao2024}。この問題はいま未解決であり、我々も未解決のままにしておく。

より確かな事実がもう一つある。睡眠はマシン全体の性質ではなく、局所ネットワークの性質である。動物を長く眠らせずにおくと、動物が起きて動いている間にも、皮質の個々の部位がノンレム睡眠のように短時間静止に落ち込み、その瞬間に動物は誤りを犯しやすくなる~\cite{Vyazovskiy2011}。各ネットワークは自分で疲れ、自分で切り替わる。全体のモードが生じるのは、ブロードキャストチャネルがすべてのネットワークを一斉に切り替えるときである。

\begin{keyblock}
\begin{proposition}[モードの集まりとしての睡眠]
\label{prop:sleep}
睡眠中もニューロンはオフにならない。ブロードキャストチャネルがマシンを別のモードに移す。いつ眠るかは、概日リズムに同期したヒステリシス付きリレー~\eqref{eq:sleeppressure}が決める。ノンレム睡眠ではネットワークが弛張発振器~\eqref{eq:updown}になり、1日を早送りで再生する。レム睡眠では入力と出力から切り離されて働く。モード、リプレイ、重みのスケーリングは測定済みである。それらの目的、すなわちインターリーブ学習と再正規化は、根拠の強い仮説である。
\end{proposition}
\end{keyblock}

\section{まとめ}

\begin{table}[t]
\centering
\footnotesize
\setlength{\tabcolsep}{4pt}
\renewcommand{\arraystretch}{1.15}
\begin{tabular}{>{\raggedright}p{3.2cm}>{\raggedright}p{4.4cm}>{\raggedright\arraybackslash}p{5.2cm}}
\toprule
起きること & 方法 & わかっていること \\
\midrule
モードの切り替え & \nt{ACET}、\nt{NORA}、\nt{SERO}のレジスタ（表~\ref{tab:sleepmodes}） & 測定済み \\
いつ眠るか & コンデンサ$S$とヒステリシス付きリレー~\eqref{eq:sleeppressure} & モデルは実験をよく記述する。$S$としてのアデノシンは仮説 \\
遅い振動 & フィードバックと疲労をもつ群~\eqref{eq:updown} & 振動は測定済み。モデルは最も単純な回路 \\
1日の早送り & $6$--$20$倍速のリプレイ、振動と協調 & 測定済み。協調を強めると記憶が良くなる \\
リプレイの目的 & 遅い記憶のインターリーブ学習 & 根拠の強い仮説 \\
重みの再正規化 & 睡眠中の結合のスケーリング & 接触の縮小は測定済み。睡眠の主な役割とするのは仮説 \\
レム睡眠 & 入力と出力を切り離した動作 & モードは測定済み。目的は不明 \\
洗浄 & 細胞間のすき間の拡大 & 矛盾する結果 \\
局所睡眠 & 個々の部位が静止に落ちる & 測定済み \\
\bottomrule
\end{tabular}
\caption{睡眠中にニューロンは何をしているか。}
\label{tab:sleep}
\end{table}

表~\ref{tab:sleep}が本章をまとめる。睡眠はマシンの作業の休止ではなく、動作中の保守であることがわかった。同じブロードキャストチャネルによるモードの切り替え、歩行リズムと同じ部品でできた発振器、遅い記憶のための1日の早送り、そして重みの再正規化である。日中にマシンは書き込み、夜には書いたものを書き直して整理する。

\section{問題}

\begin{exercise}[\lvlA]
\label{ex:20:1}
\emph{概日リズムのないリレー。}閾値が一定で$H=0.67$、$L=0.17$（図~\ref{fig:twoprocess}のモデルの閾値の平均値）とする。$\tau_r=18.2$~h、$\tau_d=4.2$~hのとき、\eqref{eq:sleeppressure}から覚醒と睡眠の長さ、発振器の周期を導け。それでも揺れる閾値のモデルが$16.1+8.0=24$~hに落ち着くのはなぜか。この効果は第\ref{sec:chemoutput}章のどの素子に関係するか。
\end{exercise}

\begin{exercise}[\lvlA]
\label{ex:20:2}
\emph{睡眠の借り。}睡眠圧$S_0$で始まった睡眠は$t_s=\tau_d\ln(S_0/L)$続く。$\tau_d=4.2$~h、$L$は目覚めるときの下の閾値である。(a)~$40$~h眠らずにいたあと$S_0=0.90$で、睡眠は$8.8$~h続いた。$L$はいくらだったか。通常の$L\approx0.092$なら睡眠はどれだけ続くか。(b)~一晩でシナプス接触の面積は$18\%$減る。日中に重みはどれだけ増えなければならないか。
\end{exercise}

\begin{exercise}[\lvlB]
\label{ex:20:3}
\emph{閾値の設計。}$\tau_r=18.2$~h、$\tau_d=4.2$~hのリレー~\eqref{eq:sleeppressure}がちょうど$16$~hの覚醒と$8$~hの睡眠を与えるように、一定の閾値$H$と$L$を選べ。ある夜$4$時間で起こされたとき、このマシンは閾値が揺れるマシンより何が劣るか。
\end{exercise}

\begin{exercise}[\lvlB]
\label{ex:20:4}
\emph{遅い振動の周期。}モデル~\eqref{eq:updown}で$F(x)=1/\bigl(1+e^{-(x-\theta)/k}\bigr)$、$w=5$、$I=2$、$\theta=2.5$、$k=0.25$、$b=5$、$\tau\ll\tau_a=0.3$~sとする。(a)~活動と静止が消える曲線$r=F(wr+I-a)$の折り返し点$a_{\text{up}}$と$a_{\text{down}}$を求めよ。(b)~活動中$r\approx1$、静止中$r\approx0$として、周期と活動の割合を求めよ。
\end{exercise}

\begin{exercise}[\lvlB]
\label{ex:20:5}
\emph{振動はいつ止まるか。}問題~\ref{ex:20:4}のモデルで、不動点は曲線$a=wr+I-F^{-1}(r)$と直線$a=br$の交点にある。交点が折り返し点の間の中間の枝にある限り振動が続く。それが成り立つのは$b$がどの範囲のときか。\nt{ACET}は$b$を変えることで、どのように発振器を増幅器に切り替えるか。
\end{exercise}

\begin{exercise}[\lvlB]
\label{ex:20:6}
\emph{シミュレーション：借りか位相か。}\texttt{sleep\_models.py}で、$48$~h後の最初の目覚めから、マシンを$D=24$、$32$、$40$、$48$、$64$~h起こしておく（\texttt{stay\_awake}、\texttt{days=8}）。各$D$で回復睡眠はどれだけ続くか。その長さを決めているものは何か。
\end{exercise}

\begin{exercise}[\lvlB]
\label{ex:20:7}
\emph{シミュレーション：忘却とインターリーブ。}線形ニューロン$y=\mathbf w\cdot\mathbf x$、$n=40$が規則$\mathbf w\leftarrow\mathbf w-0.5\,(y-y^*)\,\mathbf x$で学習する。課題AとBはそれぞれ$10$個のパターン、$x_j\sim\mathcal N(0,1/n)$、$y^*=\pm1$である。Aを$200$エポック、続いてBを$200$エポック学習したあと、Aでの二乗平均誤差はいくらか。AとBを混ぜて$200$エポック学習したあとではどうか。
\end{exercise}

\begin{exercise}[\lvlC]
\label{ex:20:8}
\emph{研究：インターリーブか再正規化か。}閾値をもつ問題~\ref{ex:20:7}のニューロンに、夜間の二つの処理を考える。(i)~すべての重みに$0.82$を掛ける。(ii)~古いパターンを新しいパターンと混ぜて繰り返す。それぞれは重みの比とニューロンの応答に何をするか。睡眠後のシナプスの大きさから、二つの仮説をどう区別できるか。
\end{exercise}

\chapter{生きたニューロンでできたマシン}
\chaptermark{生きたニューロンでできたマシン}
\label{sec:livingmachine}
\begin{chaptergoals}
\item なぜブロードキャストチャネルのない培養は一斉発火を起こし、その間隔は何で決まるのか；
\item 生きたネットワークがフィードバックで応答を変える仕組みと、まだ証明されていないこと；
\item オルガノイドが、線形読み出しを外側のシリコンで学習させるリザバーになる仕組み；
\item 培養がこれまでどう教えられてきたか、なぜ教師がまだシャーレの外にいるのか。
\end{chaptergoals}

\section{回路図の見えるテストベンチ}

第\ref{sec:debugging}章では、回路図のないマシンをデバッグした。プローブが拾えるのは800億個のニューロンのうち1000個であり、残りはすべて推測するしかない。この立場に置かれたエンジニアなら、単純なことをするだろう。回路の小さな一片をブレッドボードに組み、部品を一つ残らず見えるようにするのである。ニューロンでも同じことができる。

大脳皮質のニューロンをばらばらにし、電極アレイを載せたチップの上で育てる。電極の数は数十から数万である。数週間のうちに細胞は突起を伸ばし、数千から数十万個のニューロンからなるネットワークにつながる。大部分は\nt{GLUT}で、少数の\nt{GABA}が混じる。その割合は標本によって異なり、海馬の培養ではニューロンの約 $6\%$ である\cite{Benson1994}。各電極は、第\ref{sec:debugging}章のプローブのように聴くこともでき、テスト信号のように電流を流すこともできる。テストベンチのもう一つの種類が\emph{オルガノイド}である。ヒトの幹細胞から育てた、大きさ1ミリほどの三次元の神経組織の塊である\cite{Lancaster2013}。こうしたテストベンチの上で、生きたネットワークは何か月も動き続ける。オルガノイドの実験をネットワーク越しに遠隔で100日以上続けて行える施設もある\cite{Jordan2024}。この分野は\emph{生物学的計算}、あるいは\emph{オルガノイド知能}と呼ばれる\cite{Smirnova2023}。

テストベンチには脳との大きな違いが二つある。一つは小ささで、ニューロンの数は脳の10万分の1である。もう一つは、ブロードキャストチャネルがないことである。大脳皮質から作った培養には\nt{DOPA}、\nt{SERO}、\nt{NORA}、\nt{ACET}のニューロンがない。データバスとフロー制御はあるが制御レジスタのないマシンである。これに何ができるかを見ていこう。

\section{ネットワークが自分でしていること}

放っておかれた培養は、黙ってもいないし、一様にノイズを出してもいない。ときどきネットワークのほぼ全体が一度に発火し、つまり\emph{一斉発火}し、また静まる。一斉発火の間隔は、培養によって1秒から数分である\cite{Wagenaar2006}。エンジニアには見慣れた光景である。強い正帰還をもち負荷のない増幅器は発振器になる。

仕組みはすでに知っている。\nt{GLUT}ニューロン同士の強い相互結合は、第\ref{sec:gaba}章で命題\ref{prop:ei}の条件が破れたときのように雪崩を起こす。雪崩はシナプスの伝達物質の蓄えを急速に使い果たし\eqref{eq:depression}、ネットワークは静かになる。蓄えは時定数 $\tau_{\text{rec}}$ で回復し、新しい雪崩に足りる割合 $x^*$ まで戻ると、次の一斉発火が起こる。一斉発火の間隔は
\begin{empheq}[box=\keybox]{equation}
T_{\text{burst}} \;\approx\; \tau_{\text{rec}}\,\ln\frac{1}{1-x^*} .
\label{eq:burstinterval}
\end{empheq}
第\ref{sec:code}章の $\tau_{\text{rec}}=0.8\,\text{s}$ と $x^*=0.9$ ではおよそ2秒、$x^*=0.99$ ではおよそ4秒になる。これは本書の $\tau_{\text{rec}}$ を使ったモデルの推定であり、測定された間隔の短い側に入る。微小培養での直接測定では、一斉発火のあと伝達物質の蓄えが回復するのに数十秒かかる\cite{CohenSegal2011}。つまりそこでは $\tau_{\text{rec}}$ がもっと大きい。その $\tau_{\text{rec}}$ を使えば、式\eqref{eq:burstinterval}は数十秒から数分を与え、遅い培養と合う。これは弛張型の発振器であり、第\ref{sec:muscles}章のリズム発生器の親戚である。あちらでは群の疲労が発振器を再充電したが、ここでは伝達物質の蓄えがその役をする。

一斉発火は、ネットワークがすることのないときにすることである。生きた脳では、似た光景が深い睡眠（第\ref{sec:sleep}章）と、本物の入力が入ってくる前の発達中の脳に見られる。この類似は示唆的だが、仕組みが同じだというのは仮説である。

\section{ドーパミンの教師なしでネットワークを教える}

培養には\nt{DOPA}がないので、「良いか悪いか」の信号は自分たちで、電極から与えるしかない。実験によれば、テストベンチ上のネットワークは確かに応答を変える。そこで最も興味深いのは、\emph{何で}教えるかである。

\emph{黙る教師。}1本の電極から1〜3秒に1回ネットワークを刺激し、別の電極に刺激後 $50\pm10\ms$ で目的の応答が現れるまで続ける。応答が現れたら、刺激を5分間止める。このサイクルを数回繰り返すと、刺激に対してすぐに目的の応答が出るようになる\cite{Shahaf2001}。ここでの報酬は静けさである。刺激はネットワークを揺さぶり、刺激の停止はネットワークを正しい応答を出した状態のまま残す。これは第\ref{sec:dofa}章のランダムな摂動による勾配降下（命題\ref{prop:gradient}）であり、揺さぶられていること自体が誤差の役をする。正しくなるまで揺さぶるのである。

\emph{テニスをするネットワーク。}電極が密に並んだテストベンチの実験で、約100万個のニューロンが、ラケット1本のコンピュータテニスゲームにつながれた\cite{Kagan2022}。ボールの位置は「感覚」電極に電流として与えられた。ボールがどこにあるかは場所で、どれだけ遠いかは頻度で符号化された。二つの「運動」領域の活動がラケットを上下に動かした。学習の規則は変わっていた。ラケットがボールを打ち返すと、ネットワークは短い\emph{予測可能な}刺激を受けた。ミスすると、数秒間の\emph{予測不能な}ランダム電流を受けた。20分のプレイでラリーは長くなった。セッション内の有意な伸びは完全なフィードバックのときだけだった。刺激の代わりに無音にした場合も、セッションの終わりにはフィードバックなしより上手にプレイしたが、効果は小さかった。そして数日にわたるセッション間の安定した学習はなかった。この実験の詳しい分析は第\ref{sec:dishbrain}章にある。

著者らは、ネットワークは自分の入力が予測可能になるように行動する、という仮説でこの結果を説明した\cite{Friston2010}。第\ref{sec:code}章のスパイクの節約や第\ref{sec:recognition}章のフレームの予測と同じ考えである。入力が予想どおりなら、マシンの消費は少なくて済む。しかし実験そのものも、とりわけ培養の「知性」についての著者らの言葉も、厳しい批判を招いた。効果は小さく、もっと単純にも説明できるというのである\cite{Balci2023}。公正な評価はこうである。テストベンチ上の培養がフィードバックによって振る舞いを変えることは測定済みである。それがまさに自分の入力を予測するように学習しているというのは仮説である。

\emph{体を動かすネットワーク。}同じように、培養を単純な仮想の「体」につなぎ、訓練刺激の時空間パターンを選ぶことで、体を目標へ動かすよう教えた実験もある\cite{Bakkum2008}。これらの実験のどれにもドーパミンはない。教師の役をするのは、エンジニアが選ぶ電流のパターンである。教師がプログラムやドーパミンそのものになると何が変わるかは、\ref{sec:dishdopamine}--\ref{sec:cartpole}節で見る。

\begin{figure}[t]
\centering
\begin{tikzpicture}[>=Stealth,
  blk/.style={draw,very thick,rounded corners=2pt,align=center,font=\footnotesize,minimum height=0.8cm,fill=blue!5}]
% (a) closed loop
\node[font=\small] at (1.5,4.3) {(a) 閉ループ};
\node[blk,text width=2.6cm] (G) at (1.5,3.3) {ゲーム：\\ボールとラケット};
\draw[very thick,fill=orange!10,rounded corners=3pt] (0.5,0.2) rectangle (2.5,1.6);
\foreach \x in {0.7,1.0,...,2.35}{\foreach \y in {0.4,0.7,1.0,1.3}{\fill[gray!60] (\x,\y) circle (0.04);}}
\draw[->,thick,blue!60!black] (G.west) -| (-0.5,0.9) -- (0.5,0.9);
\node[font=\scriptsize,blue!60!black,anchor=east] at (-0.6,2.1) {ボール $\to$ 電流};
\draw[->,thick,red!70!black] (2.5,0.9) -- (3.5,0.9) |- (G.east);
\node[font=\scriptsize,red!70!black,anchor=west] at (3.6,2.1) {スパイク $\to$ ラケット};
\node[font=\scriptsize] at (1.5,-0.2) {電極アレイ上のニューロン};
\node[font=\scriptsize,align=center] at (1.5,-0.85) {ヒット：予測可能な刺激\\ミス：ランダム電流};
% (b) reservoir
\begin{scope}[xshift=7.9cm]
\node[font=\small] at (1.6,4.3) {(b) リザバー};
\node[blk,text width=1.1cm] (X) at (-0.4,1.0) {入力\\$u(t)$};
\draw[very thick,fill=orange!10] (1.6,1.0) circle (0.8);
\foreach \a/\r in {20/0.35,80/0.5,150/0.3,210/0.55,280/0.4,330/0.2,110/0.15}{\fill[red!60!black] ({1.6+\r*cos(\a)},{1.0+\r*sin(\a)}) circle (0.05);}
\node[font=\scriptsize] at (1.6,-0.2) {オルガノイド：教師なし};
\node[blk,text width=1.8cm,fill=green!8] (W) at (4.0,1.0) {読み出し\\$W_{\text{out}}$};
\draw[->,thick] (X) -- (0.8,1.0);
\draw[->,thick] (2.4,1.0) -- node[above,font=\scriptsize] {$\mathbf r(t)$} (W);
\draw[->,thick] (W) -- (4.0,2.3) node[above,font=\scriptsize] {$y(t)$};
\node[font=\scriptsize,green!40!black] at (4.0,0.3) {教師あり学習};
\end{scope}
\end{tikzpicture}
\caption{生きたネットワークに計算させる二つの方法。(a) 閉ループ：ボールの位置は電流として与えられ、運動領域のスパイクがラケットを動かし、打った後の予測可能な刺激またはランダムな刺激が教師となる。(b) リザバー\eqref{eq:reservoir}：オルガノイドには学習信号を与えない。オルガノイドは入力を、記憶をもつ多数の非線形応答に分解する。誤差から学習するのは外側の線形読み出し一つだけである。ただしオルガノイド自身の結合もこのとき変化する。教師なしで。}
\label{fig:livingmachine}
\end{figure}

\section{生きたリザバー}

生きたネットワークの使い方はもう一つある。まったく教えないことである。工学ではこれを\emph{リザバー計算}と呼ぶ\cite{Maass2002,JaegerHaas2004}。記憶をもつ出来合いの非線形システムを持ってくる。入力への応答が豊かで、直近の過去に依存しさえすれば何でもよい。入力 $u(t)$ がそれを揺り動かし、応答ベクトル $\mathbf r(t)$ が得られる。学習するのは線形の読み出し
\begin{empheq}[box=\keybox]{equation}
y(t) \;=\; W_{\text{out}}\,\mathbf r(t) ,
\label{eq:reservoir}
\end{empheq}
だけであり、その重みは第\ref{sec:motorlearning}章の最小二乗則\eqref{eq:lms}で決める。リザバーは、第\ref{sec:motorlearning}章の小さな\nt{GLUT}ニューロンがしていたことをする。入力を長いベクトルに分解し、目的のクラスが平面で分けられるようにするのである（第\ref{sec:dendrites}章）。

電極アレイ上のオルガノイドは、こうしたリザバーの役に向いていた。電流パターンとして与えた音声は、線形読み出しを一つ学習させたあと、話者を約 $78\%$ の精度で区別できた\cite{Cai2023}。$78\%$ という精度は著者ら自身が報告し、報道が繰り返した数字である。有用な結果だが、その中のどこに「知恵」があるかを理解しておくことが大切である。オルガノイドは非線形性と減衰する記憶を与え、誤差による学習は外側のシリコンで行われる（図\ref{fig:livingmachine}）。ただしオルガノイドは不変のままではない。著者らによれば、訓練の間にオルガノイド自身の機能的結合が変化し、彼らはこれをオルガノイド内の教師なし学習と呼んでいる\cite{Cai2023}。精度のどれだけがこの変化によるもので、どれだけが読み出しによるものかは切り分けられていない。

\section{テストベンチがマシンについて語ること}

\emph{エネルギー。}推定\eqref{eq:energy}によれば、スパイク1回のコストは約 $10^{-10}$~J である。100万個のニューロンからなる培養が1秒に1スパイクの頻度で発火すると、信号伝達に使う電力は
\begin{empheq}[box=\keybox]{equation}
P \;\approx\; 10^{6}\times1\,\text{s}^{-1}\times10^{-10}\,\text{J} \;=\; 10^{-4}\,\text{W},
\label{eq:dishpower}
\end{empheq}
すなわち0.1ミリワットである。これらのスパイクを刺激し、記録し、処理する電子回路は、桁違いに多くを使う。第\ref{sec:debugging}章と同じく、ボトルネックは計算ではなく、計算との通信である。

\emph{足りない部品。}テストベンチは、制御レジスタなしで、\nt{GLUT}のバスと\nt{GABA}のフロー制御にどれだけのことができるかを示す。自己組織化し、一斉発火を生み、フィードバックによって応答を変え、よいリザバーになる。それらなしにできないのは、何を成功とみなすかを自分で決めることである。この信号を、テストベンチではエンジニアが与え、脳では、第\ref{sec:dofa}--\ref{sec:acet}章で見たように、ブロードキャストチャネルが与える。

\emph{倫理。}オルガノイドはヒトの細胞から育てられ、複雑になるほど、そうした組織が何かを感じうるかという問いは重くなる。工学的な見方は部分的な答えを示す。第\ref{sec:pain}章の痛みは一つのセンサの信号ではなく、ゲート、音量調節器、ブロードキャストチャネルを備えた警報システム全体であり、それらはテストベンチにはない。しかしこれは推論であって証明ではない。この分野自身も、倫理的評価を実験と並行して、最初から行うことを提案している\cite{Smirnova2023}。

\begin{keyblock}
\begin{proposition}[テストベンチ上の生きたネットワーク]
\label{prop:livingmachine}
電極アレイ上の\nt{GLUT}ニューロンと\nt{GABA}ニューロンのネットワークは、自ら一斉発火\eqref{eq:burstinterval}を生み、フィードバックによって応答を変え、外で学習させた読み出しのためのリザバー\eqref{eq:reservoir}になりうる。これらはすべて測定済みである。そうしたネットワークが何らかの一般的な規則、たとえば自分の入力を予測するという規則で学習するというのは仮説であり、その「知性」をめぐる大げさな言葉を大多数の専門家は受け入れていない。
\end{proposition}
\end{keyblock}

\section{光のフラッシュで放出するドーパミン}
\label{sec:dishdopamine}

培養にドーパミンを教師として与えた直接の実験として私たちが知っている唯一のものは、研究者が遠隔で接続するオルガノイドのプラットフォームで行われた\cite{Jordan2024}。オルガノイドを浸す液に、光感受性の「ケージ」に閉じ込めたドーパミンを加えた。ケージに入った分子は受容体に作用しない。ケージドドーパミンの濃度は $1$~mM である。ネットワークに報酬を与えたいときは紫外線を当てる。ケージが壊れ、ドーパミンが空間に出ていく。

ループは次のようにできている。同じ刺激を何度も与え、それが以前より多くのスパイクを引き起こしたら、フラッシュを点けてドーパミンを放出する。$240$ 回の繰り返しで、誘発スパイクの数は全体として増えた。著者ら自身、この実験一つだけでは増加がドーパミンによるとは結論できないと書いている\cite{Jordan2024}。

エンジニアにとって、これはシャーレの中に三要素則\eqref{eq:threefactor}を組む最初の試みである。送り手と受け手の活動が適格度トレースを残し、共通の信号が変化を定着させるかどうかを決める。しかしここでの信号は正にしかならない。報酬があるかないかであり、負の誤差 $\delta<0$ をこの装置は出せない。そのうえドーパミンは広がり、ゆっくりとしか除去されない。\eqref{eq:lowpass}の濃度と同じである。そのため頻繁なフラッシュは単に全体のレベルを上げるだけかもしれない。学習をこれと区別するには二つの対照が必要である。ネットワークの応答と無関係なランダムな時刻に同じ数のフラッシュを与える対照と、ケージドドーパミンなしで同じフラッシュを与える対照である。さらに休止後の確認も必要である。ドーパミンが去ったあとも変化が残っているかどうかである。

\section{ドーパミンがシリコンを教える}
\label{sec:biohybrid}

逆向きの実験では、生きた細胞が電子回路を教える\cite{Keene2020}。ドーパミンを分泌する細胞を、有機ニューロモルフィック素子の上の微小流路で育てる。この素子はトランジスタで、そのチャネルのコンダクタンスがシナプスの重みの役をする。細胞から出たドーパミンは素子の電極で酸化され、それがコンダクタンスを変える。この装置はシナプス間隙から伝達物質を除去する仕組みも再現しているので、重みは長く変えておくことも、元に戻すこともできる。

回路としては、これは化学入力の漏れ積分器である。重みはドーパミンの流入を蓄積し、「汲み出し」が忘れる速さを決める。\eqref{eq:lowpass}と同じ RC 回路である。ここで肝心なのは接続の向きである。第\ref{sec:debugging}章のプローブにブロードキャストレジスタが見えないのは、それが化学的だからである。この素子はまさにその化学レジスタを読み取り、電気的な重みに変える。ブレイン・コンピュータ・インターフェースにとって、これはデータバスだけでなく制御レジスタも聴き取れるセンサーへの道である。

\section{自前のドーパミンニューロンをもつオルガノイド}
\label{sec:midbrainorg}

ヒトの幹細胞から、\nt{DOPA} ニューロンのある脳領域に似たオルガノイドを育てることができる。そこには働くドーパミンニューロンがあり、ドーパミンを分泌する\cite{Jo2016}。こうしたオルガノイドは主に病気の研究のために育てられている。そのドーパミンを別のネットワークの教師として使った実験は、私たちには見つからなかった。

生きたニューロンでできたマシンにとって、これは足りない部品である。第\ref{sec:livingmachine}章のテストベンチは、制御レジスタなしのデータバスとフロー制御だった。ここにはブロードキャスト信号の源がすでにある。足りないのは批評家、つまり予測誤差\eqref{eq:ctd}を計算してこれらのニューロンを制御するネットワークである。大脳皮質のオルガノイドをこうしたオルガノイドとつなぎ、誤差に応じてドーパミンが放出されるようにすることは未解決の課題であり、今のところ実験ではなく仮説である。

\section{ドーパミンなしで棒を立てるオルガノイド}
\label{sec:cartpole}

最近の実験で最も完全なものは、ドーパミンをまったく使わない\cite{Robbins2026}。マウス大脳皮質のオルガノイドを、「台車の上の棒」（倒立振子）の課題と閉ループでつないだ。オルガノイドの活動が台車を動かし、棒の位置は刺激として戻される。試行の合間に、オルガノイドは短い高頻度の学習刺激を受ける。それをどれにするかは強化学習プログラムが選ぶ。

結果は測定されている。
\begin{itemize}
\item 大多数のオルガノイドで、プログラムが選んだ学習信号は、ランダムに選んだ信号や信号なしより良い成績をもたらした。
\item $45$ 分の休止後、改善は残らなかった。
\item AMPA と NMDA の両方のグルタミン酸受容体を遮断すると、改善はなくなった。
\end{itemize}

最後の項目は、学習したものがどこにあるかを語る。\nt{GLUT} シナプスの中であり、\ref{sec:weightstore}節で入力と出力の同時検出器として働く受容体を通じてである。2番目の項目は何が足りないかを示す。速い変化はあるが定着がない。第\ref{sec:motorlearning}章の、遅い学習者を欠いた速い学習者と同じである。そして教師の役割は変わった。電流パターンを選ぶエンジニアはプログラムに置き換わったが、教師は依然としてシャーレの外にいる。

電極アレイ上の培養を学習させたもっと古い実験の中にも、ドーパミンを使ったものは一つも見つからなかった。学習信号はどれも電気刺激だった。

\section{誰が培養を教えるのか}

\begin{table}[t]
\centering
\footnotesize
\setlength{\tabcolsep}{4pt}
\renewcommand{\arraystretch}{1.15}
\begin{tabular}{>{\raggedright}p{2.7cm}>{\raggedright}p{2.5cm}>{\raggedright}p{3.2cm}>{\raggedright\arraybackslash}p{4.4cm}}
\toprule
実験 & 教えるのは誰か & 学習信号 & わかっていること \\
\midrule
目的の応答で刺激を止める & エンジニア & 刺激の終了 & 応答が定着する（測定済み） \\
DishBrain（第\ref{sec:dishbrain}章） & エンジニア & 予測可能な電流またはランダム電流 & セッション内のラリーの有意な伸びは完全なフィードバックのときだけ、無音では効果が小さい（測定済み）。セッション間では不安定。原因は議論中 \\
台車の上の棒 & 強化学習プログラム & プログラムが選ぶ高頻度刺激 & ランダムより良いが、休止後に残らない（測定済み） \\
フラッシュで放出するドーパミン & 「スパイクが増えたら報酬」という規則 & 光で解放されるドーパミン & スパイクの増加は観察。ドーパミンの役割は未証明 \\
バイオハイブリッドシナプス & 生きた細胞 & 細胞からのドーパミン & 素子の重みの長期変化と復帰（測定済み） \\
\nt{DOPA} ニューロンをもつオルガノイド & --- & --- & ドーパミン源はある。教師としては未使用 \\
\bottomrule
\end{tabular}
\caption{チップ上の生きたニューロンを誰が何で教えるか。}
\label{tab:dishteachers}
\end{table}

培養自体が学習するどの実験でも、教師は今のところ外にいる（表\ref{tab:dishteachers}）。エンジニアか、プログラムか、エンジニアが定めた規則である。ドーパミンがシャーレに入ったのは一度だけで、その役割はまだ証明されていない。第\ref{sec:dofa}章のように批評家、誤差、適格度トレースを備えた本物のブロードキャストチャネルをシャーレの中に組むことが次の一歩であり、まだ誰も踏み出していない。


\section{まとめ}

\begin{table}[t]
\centering
\footnotesize
\setlength{\tabcolsep}{4pt}
\renewcommand{\arraystretch}{1.15}
\begin{tabular}{>{\raggedright}p{2.8cm}>{\raggedright}p{4.2cm}>{\raggedright\arraybackslash}p{5.8cm}}
\toprule
実験 & 生きたネットワークがすること & わかっていること \\
\midrule
入力のないネットワーク & 間隔1秒から数分の一斉発火\eqref{eq:burstinterval} & 測定済み。シナプスの枯渇による機構は広く受け入れられたモデル \\
黙る教師 & 刺激を止めると目的の応答が定着する & 測定済み \\
テニスゲーム & ラリーが長くなる。セッション内の有意な伸びは完全なフィードバックのときだけ & 振る舞いの変化は測定済み。セッション間の学習は不安定。効果の大きさと入力予測による説明には異論がある \\
仮想の体 & 刺激パターンを選べば目標へ動く & 測定済み \\
リザバー & 読み出し\eqref{eq:reservoir}のための非線形性と記憶 & 測定済み。誤差による学習は外側のシリコンで。オルガノイドの結合も変化する \\
エネルギー & 100万ニューロンあたり約 $10^{-4}$~W~\eqref{eq:dishpower} & \eqref{eq:energy}による推定 \\
\bottomrule
\end{tabular}
\caption{生きたニューロンでできたマシンが示したこと。}
\label{tab:livingmachine}
\end{table}

生きたニューロンでできたマシン（表\ref{tab:livingmachine}）は、制御レジスタなしのデータバスとフロー制御に何ができるかを見せるテストベンチである。できることは多いが、何が良いかを自分で決めることはできない。テストベンチではエンジニアが自分の電流パターンでこの役を担い、脳では、本書の中盤で扱ったわずかなブロードキャスト配線が担う。

\section{問題}

\begin{exercise}[\lvlA]
\label{ex:21:1}
\emph{一斉発火の間隔。}一斉発火は蓄えをゼロまで使い果たし、その後 $\tau_{\text{rec}}\,\dd x/\dd t=1-x$ に従って回復し、$x=x^*$ で次の一斉発火が始まる。$\tau_{\text{rec}}=0.8$~s。(a)~$x^*=0.9$ と $0.99$ のときの間隔を求めよ。(b)~閾値 $x^*$ が $0.99$ の周りで $\pm0.005$ ランダムに変わる。間隔はどの範囲に入るか。
\end{exercise}

\begin{exercise}[\lvlA]
\label{ex:21:2}
\emph{テストベンチのエネルギー。}(a)~推定\eqref{eq:dishpower}は脳全体（$8.6\cdot10^{10}$ ニューロン）でどれだけの電力を与えるか。(b)~テストベンチは $1024$ 個の電極を $20$~kHz、1サンプル $10$~bit で記録する。1 bit あたり $1$~nJ のとき、このデータの伝送は培養自体の電力の何倍か。
\end{exercise}

\begin{exercise}[\lvlB]
\label{ex:21:3}
\emph{設計：一斉発火を抑える。}多数の電極からのまばらな刺激で、各シナプスが頻度 $r_b$ で伝達物質を放出する。蓄えは $\tau_{\text{rec}}\,\dd x/\dd t=1-x-Ur_b\tau_{\text{rec}}x$、$U=0.5$、$\tau_{\text{rec}}=0.8$~s。蓄えが $x^*=0.9$、$0.99$ に届かなくなる最小の $r_b$ はいくらか。そのときシナプスはどれだけ弱まるか。
\end{exercise}

\begin{exercise}[\lvlB]
\label{ex:21:4}
\emph{リザバーの記憶はニューロン数を超えない。}出力 $\mathbf r(t)$ が $N$ 個のリザバーに、平均ゼロ・分散1の白色入力 $u(t)$ を与える。$m_k$ を読み出し $\mathbf w_k\cdot\mathbf r(t)$ と $u(t-k)$ の相関の2乗の最大値とする。記憶容量 $\mathrm{MC}=\sum_{k\ge0}m_k\le N$ を証明せよ。
\end{exercise}

\begin{exercise}[\lvlB]
\label{ex:21:5}
\emph{シミュレーション：シリコンのリザバー。}リザバー $\mathbf r(t+1)=\tanh(W\mathbf r(t)+\mathbf v\,u(t)+\mathbf c)$、$N=30$、$W$ はスペクトル半径 $0.9$ のランダム行列、$v_i\sim U(-0.5,0.5)$、$c_i\sim U(-0.2,0.2)$、$u\sim U(-1,1)$、読み出しはリッジ回帰。問題\ref{ex:21:4}の記憶容量を $\tanh$ ありとなしで求めよ。どちらのリザバーがより多く覚えているか。
\end{exercise}

\begin{exercise}[\lvlB]
\label{ex:21:6}
\emph{生きたリザバーの読み出し。}オルガノイドは $1024$ チャネルと、2クラスの訓練記録 $P=240$ 個を与える。(a)~問題\ref{ex:19:3}の式で、ランダムなラベル付けが分離可能になる確率を $1\%$ 以下に抑えるには、特徴量 $n$ をいくつまで残せるか。(b)~$100$ 個の検証記録で精度 $78\%$ は、当て推量より何シグマ高いか。
\end{exercise}

\begin{exercise}[\lvlC]
\label{ex:21:7}
\emph{研究：ブロードキャストチャネルのない教師。}第\ref{sec:dofa}章の三要素則を $8$--$1024$ 個の電極でテストベンチ上に実現する刺激プロトコルと、読み出しを凍結して重みの学習とネットワーク状態のずれを区別する対照実験を提案せよ。どんな結果が学習の仮説を反証するか。
\end{exercise}

\chapter{DishBrain}
\chaptermark{DishBrain}
\label{sec:dishbrain}

\section{シャーレの中のネットワークがテニスをする}

DishBrain とは、電極付きのチップ上で育てた生きたニューロンの培養が、ラケット1本の簡略化したテニスをするテストベンチに著者らがつけた名前である\cite{Kagan2022}。生きたニューロンでできたマシンの実験の中で最も議論を呼んだものであり（こうしたマシンの概観は第\ref{sec:livingmachine}章）、本書にとってはとりわけ扱いやすい。ここでは小さなネットワーク全体に手が届く。入力はすべて実験者が決め、出力はすべて記録される。これまでの章の問い、つまり信号はどう符号化されているか、誰が教えるか、プローブに何が見えるかを、目も筋肉も \nt{DOPA} ニューロンももたないネットワークに問うことができる。エンジニアが他人のテストベンチを調べるように、この実験を分解しよう。回路、入力、出力、教師、結果、論争の順である。

\section{テストベンチ}

ニューロンはマウス胎児の大脳皮質から採る（チップあたり約 $800\,000$ 個）か、ヒトの幹細胞から育てる（チップあたり約100万個）。数週間チップ上で育つうちに、ニューロンは互いに突起を伸ばし合い、自らスパイクとバーストを出し始める。その下の $8$~mm$^2$ の区画には $26\,000$ 本の白金電極が並ぶ。同時に記録できるのはそのうち最大 $1024$ 本で、刺激に使うのは実際には8本である。

培養の周りにはループが閉じている（図\ref{fig:dishloop}）。コンピュータがゲームをシミュレートする。ボールが飛び、壁で跳ね返り、ラケットが上下に動く。ボールの位置は8本の「感覚」電極の電流に変換される。二つの「運動」領域のスパイクを $10\ms$ の窓ごとに数え、ラケットを動かす。ヒットやミスのたびに、培養はもう一つの刺激、つまりフィードバックを受ける。$10\ms$ の窓はテストベンチのコンピュータのクロックであり、培養のクロックではない。ネットワーク自体は、本書のすべてのネットワークと同じく非同期で動いている。

\begin{figure}[t]
\centering
\begin{tikzpicture}[>=Stealth,
  blk/.style={draw,very thick,rounded corners=3pt,align=center,font=\small,minimum height=1.2cm},
  lab/.style={font=\scriptsize,align=center}]
\node[blk,fill=blue!6,text width=3.4cm] (C) at (0,0) {チップ上の培養\\\scriptsize $\sim10^6$ ニューロン\\\scriptsize 電極 $26\,000$ 本};
\node[blk,fill=orange!10,text width=3.0cm] (G) at (7.2,0) {コンピュータ：\\テニスゲーム};
\draw[->,very thick,blue!60!black] (G.north west) to[bend right=25] node[lab,above] {ボール：\emph{場所}は8本のどの電極か、\\\emph{頻度}は 4--40 Hz} (C.north east);
\draw[->,very thick,red!70!black] (C.south east) to[bend right=25] node[lab,below] {ラケット：「上」と「下」の領域の\\$10\ms$ ごとのスパイク数} (G.south west);
\draw[->,thick,densely dashed,green!45!black] (G.east) -- ++(0.9,0) |- ++(0,-2.6) -| node[lab,pos=0.25,below] {ヒット：予測可能な刺激、ミス：予測不能な刺激} (C.south);
\end{tikzpicture}
\caption{DishBrain のループ。青い矢印は入力で、ボールの位置は場所による符号とレート符号で書かれる（第\ref{sec:code}章）。赤は出力で、二つの領域のカウントの差がラケットを動かす。緑の破線は各ラリー後のフィードバックで、テストベンチ唯一の「教師」である。}
\label{fig:dishloop}
\end{figure}

\section{入力と出力}

\emph{入力}は第\ref{sec:code}章の二つの符号で同時に書かれる。8本のうちどの電極を刺激するかがボールの高さを表す。これは第\ref{sec:sensors}章のセンサーと同じ、場所による符号である。刺激パルスの頻度がボールの遠さを表す。ボールが遠い壁にあるときは毎秒 $4$ 回、ラケット側の壁に近づくにつれて線形に増え、毎秒 $40$ 回になる。これはレート符号で、ボールのパルスの振幅は $75$~mV である。培養はどちらの符号もあらかじめ「知って」はいない。符号は実験者が決めたもので、ネットワークはそれを読むことを学ばなければならない。

\emph{出力}は第\ref{sec:debugging}章のインターフェースと同じように読む。一方の運動領域のスパイクはラケットを上へ、もう一方は下へ押す。最も単純に書けば、各窓で
\begin{empheq}[box=\keybox]{equation}
\Delta p \;=\; c\,\bigl(n_\uparrow - n_\downarrow\bigr),
\label{eq:dishreadout}
\end{empheq}
ここで $n_\uparrow$、$n_\downarrow$ は二つの領域の $10\ms$ あたりのスパイク数、$c$ はテストベンチの係数である。これは第\ref{sec:glutcomputer}章の2線式の符号である。動きの向きを担うのは信号の符号ではなく、2本の配線のどちらが大きいかである。

この出力のノイズを見よう。ある領域が合計で毎秒 $R$ 個のスパイクを出すなら、窓 $T=10\ms$ の中には平均 $RT$ 個あり、\eqref{eq:rateerror}により相対的なばらつきは $1/\sqrt{RT}$ である。$R=1000$ スパイク/秒なら各窓で約 $30$ パーセント、$R=100$ なら100パーセントを超える。ラケットは避けがたく震え、ネットワークが学習できる方法はただ一つ、カウントの差の\emph{平均}を正しい向きにずらすことである。これは第\ref{sec:debugging}章のあらゆるインターフェースを制約するのと同じ法則であり、今回はループの両側に効いている。

\section{ドーパミンのない教師}

このテストベンチで最も興味深いのは教師である。大脳皮質ニューロンの培養には \nt{DOPA} ニューロンがなく、テストベンチは「予想より良い」という信号をまったく送らない。フィードバックは別物である。ネットワークがボールを打ち返すと、8本の感覚電極すべてに同時に短い\emph{予測可能な}刺激が与えられる。$75$~mV のパルスを毎秒 $100$ 回、$0.1\,\text{s}$ の間である。ミスすると、著者らが\emph{予測不能}と呼ぶ刺激が4秒続く。$150$~mV のパルスを毎秒 $5$ 回、ランダムな時刻に、ランダムに選んだ電極へ与える。その後、刺激のない休止が4秒あり、ボールが再スタートする\cite{Kagan2022}。

著者らは、ネットワークは自分の入力が予測可能になるように行動する、という仮説から出発した\cite{Friston2010}。エンジニアにとって意味は単純である。4秒間のランダムなノイズから逃れる方法は、培養にはただ一つ、ボールを打ち返すことしかない。同じ考えは、第\ref{sec:code}章のスパイクの節約と第\ref{sec:recognition}章のフレームの予測にも出てきた。

しかし、ここで本当に予測可能性が必要なのだろうか。もっと粗い仕組みもありうる。ミスの後の予測不能な刺激はネットワーク内の最近の変化を単に\emph{かき混ぜ}、ヒットの後の穏やかな刺激はそれに触れない、としよう。ネットワークの設定を一つのベクトル $\boldsymbol\psi$ で表し、設定 $\boldsymbol\psi$ でネットワークがミスする確率を $P_{\text{miss}}(\boldsymbol\psi)$ とする。揺さぶりが対称である、つまり $\boldsymbol\psi$ から $\boldsymbol\psi'$ へ移る確率がその逆と等しいなら、定常状態では各設定から出ていく流れと入ってくる流れが釣り合い、ネットワークが設定 $\boldsymbol\psi$ で過ごす時間の割合は
\begin{empheq}[box=\keybox]{equation}
\rho(\boldsymbol\psi) \;\propto\; \frac{1}{P_{\text{miss}}(\boldsymbol\psi)} .
\label{eq:dishdensity}
\end{empheq}
ミスが10分の1の設定は、10倍長く保たれる。重みをどちらに変えるべきかをネットワークに教える者はいない。良い設定がただ壊れにくいだけである。

これをモデルで確かめた。培養を二つの数に還元する。ボールが高さ $y$ に来たときラケットは点 $a\,y+b$ に来て、外れが $0.1$ 未満なら打ち返す。ミスの後は両方の数がランダムに揺さぶられ、ヒットの後はそのまま残る。$2000$ ラリーの間に、打ち返した割合は最初の200回の $0.21$ から最後の500回の $0.38$ に伸びた（40個の「培養」、図\ref{fig:dishmodel}）。\emph{毎回の}ラリーの後に揺さぶると、フィードバックに情報がなく、割合は約 $0.12$ にとどまる。揺さぶりがまったくなければ $0.19$ のままである。テストベンチの実験はこれと合っている。セッション内の有意な伸びは完全なフィードバックのときだけだった。刺激の代わりに無音にした場合も、セッションの終わりにはフィードバックなしより上手にプレイしたが、効果は小さかった\cite{Kagan2022}。ただし、この条件でもミスの後には長い休止が、ヒットの後には短い休止が続くので、結果による違いは完全には消えていないことに注意したい。

\begin{figure}[t]
\centering
\begin{tikzpicture}[x=1cm,y=5cm]
\draw[->,gray!70!black] (0,0) -- (10.6,0);
\draw[->,gray!70!black] (0,0) -- (0,0.5) node[above,font=\small,black] {打ち返した割合};
\foreach \v in {0.1,0.2,0.3,0.4}{\draw[gray!60!black] (-0.06,\v) -- (0.06,\v); \node[left,font=\scriptsize] at (0,\v) {\v};}
\foreach \x/\f/\l/\lab in {1.8/0.21/0.38/{ミスの後に\\ノイズ},5.2/0.13/0.11/{毎ラリーの後に\\ノイズ},8.6/0.19/0.19/{ノイズなし}}{
  \fill[blue!35] (\x-0.8,0) rectangle (\x-0.05,\f);
  \fill[red!60!black] (\x+0.05,0) rectangle (\x+0.8,\l);
  \node[below,font=\scriptsize,align=center] at (\x,-0.01) {\lab};
  \node[above,font=\scriptsize] at (\x-0.425,\f) {\f};
  \node[above,font=\scriptsize] at (\x+0.425,\l) {\l};}
\fill[blue!35] (7.4,0.47) rectangle (7.7,0.495); \node[right,font=\scriptsize] at (7.75,0.482) {最初の200回};
\fill[red!60!black] (7.4,0.41) rectangle (7.7,0.435); \node[right,font=\scriptsize] at (7.75,0.422) {最後の500回};
\end{tikzpicture}
\caption{「ミスでかき混ぜる」モデル（\texttt{dishbrain\_model.py}）：$2000$ ラリーの最初と最後で打ち返した割合。40個のモデル培養の平均。学習が進むのは、揺さぶりがミスの後に来てヒットの後には来ないときだけである。セッション内の有意な伸びが完全なフィードバックのときだけだった実験と同じである。}
\label{fig:dishmodel}
\end{figure}

\begin{keyblock}
\begin{proposition}[かき混ぜる教師]
\label{prop:dishscramble}
失敗がネットワークを揺さぶり、成功がそのままにしておくなら、ネットワークはひとりでに、誤りの少ない設定へと漂っていく。ある設定で過ごす時間は失敗の確率に反比例する\eqref{eq:dishdensity}。このような学習には、報酬信号も、その符号も、予測可能性そのものも要らない。成功の後と失敗の後でフィードバックが異なれば十分である。培養の振る舞いの変化は測定済みである。シャーレの中で働いている仕組みが入力の予測なのか、かき混ぜなのかは確定していない。
\end{proposition}
\end{keyblock}

第\ref{sec:dofa}章と比べてほしい。そこでもノイズは探索の役をしたが、ドーパミンには符号があった。誤差 $\delta$ が重みをどちらにずらすかを告げ、一歩は勾配\eqref{eq:perturbation}に沿って進んだ。ここには符号がない。あるのは「残す」か「揺さぶる」かだけである。この探索はより粗く遅いが、ネットワークへの要求は少ない。\nt{DOPA} ニューロンも、批評家も、数秒にわたる適格度トレースも要らない。

\section{測定されたことと論争点}

各ゲームは $20$ 分続き、最初の $5$ 分と最後の $15$ 分を比較した。完全なフィードバックを受けた培養では、ラリーが長くなり、即座のミスが減った。細胞なし、ゲームなし、コンピュータの「ラケット」を使った対照では、そうしたことは起こらなかった。ヒト細胞の培養は、平均するとマウスの培養より長く打ち返し続けた。改善は統計的に確かである（マウス9個とヒト11個の培養、各群100回を超えるゲーム）が、小さい。ネットワークは良いプレーヤーにはならず、ランダムより少し良くなっただけである。確かなのはセッション内の改善であり、数日にわたるセッション間の安定した学習は著者らには得られなかった\cite{Kagan2022}。同じグループの次の研究では、プレイ中の培養の活動が\emph{臨界}状態、つまり減衰と雪崩の境界、第\ref{sec:gaba}章で見たあの縁の上の生き方に近づき、近づくほどプレイが良くなることが見つかった\cite{Habibollahi2023}。

主な論争を呼んだのは数字ではなく言葉である。論文の題は、シャーレの中のニューロンが「知性」（sentience）を示すと主張した。研究者のグループが反論した。閉ループでの振る舞いの変化はまだ知性でも感覚でもなく、結論はもっと控えめに述べるべきだ、と\cite{Balci2023}。工学的な観点からは、未解決の問いが二つあり、どちらも仕組みに関するものである。一つ目：ネットワークは学習しているのか、つまり重みが長期的に変わっているのか、それともフィードバックが現在の状態をずらしているだけなのか。二つ目：本当に予測可能性が必要なのか、それとも命題\ref{prop:dishscramble}のように、成功と失敗への応答の違いがあれば十分なのか。すべての電極に手が届くテストベンチなら、両方に答えられる。おそらくそれがこのテストベンチの最大の価値である。

\section{テストベンチの仕様：再現のために}
\label{sec:dishspec}

以下に、このテストベンチを組み直すのに必要なものをすべてまとめる。装置、ループ、入力の符号化、出力の読み出し、フィードバック、実験プロトコルである。各パラメータには出典を記す。「論文」は、値を論文\cite{Kagan2022}の本文から採ったことを示す。「選択」は、論文に記載がなく、再現するには自分で決めなければならないことを示す。既定値を提案し、それをどう確かめるかを述べる。

\subsection*{装置と培養}

\begin{center}
\footnotesize
\setlength{\tabcolsep}{4pt}
\renewcommand{\arraystretch}{1.15}
\begin{tabular}{>{\raggedright}p{4.0cm}>{\raggedright}p{6.9cm}>{\raggedright\arraybackslash}p{1.6cm}}
\toprule
パラメータ & 値 & 出典 \\
\midrule
チップ & MaxOne アレイ：$8$~mm$^2$ の区画に白金電極 $26\,000$ 本 & 論文 \\
記録 & 同時に最大 $1024$ チャネル、毎秒 $20\,000$ サンプル & 論文 \\
刺激 & 技術的には最大 $32$ 電極、実験では独立な $8$ 電極 & 論文 \\
細胞 & マウス胎児の大脳皮質。幹細胞由来のヒトニューロン（2通りの培養法）。対照として HEK293T 細胞（ニューロンではない） & 論文 \\
播種 & チップあたり約 $10^6$ 個 & 論文 \\
実験時の培養日数 & マウス：$\sim10$ 日から。ヒト：培養法により $14$--$24$ 日または $\sim80$ 日 & 論文 \\
\bottomrule
\end{tabular}
\end{center}

\subsection*{ループと時間}

ループは $10\ms$（$200$ サンプル）の窓ごとに動く。窓ごとに運動領域の電極のスパイクを数え、ゲームを更新し、次の刺激を出す。培養内のスパイクから応答刺激までの遅れは約 $5\ms$ である。プレイ中、各電極の信号は遮断周波数 $100$~Hz の2次ベッセル高域通過フィルタを通る。信号の絶対値は $1$~Hz の1次ベッセル低域通過フィルタで平滑化され、スパイクの閾値はこの平滑化した絶対値に比例する。背景の標準偏差の6倍という閾値を論文が挙げているのは、個別の活動測定についてであり、ゲームについてではない。刺激が培養の応答として数えられないよう、$75$~mV を超える大きなパルスが同時に $15$ 個より多く来たときは、すべての信号を消去する。以上はすべて論文による。

\subsection*{ボールの符号化}

\begin{center}
\footnotesize
\setlength{\tabcolsep}{4pt}
\renewcommand{\arraystretch}{1.15}
\begin{tabular}{>{\raggedright}p{3.4cm}>{\raggedright}p{7.5cm}>{\raggedright\arraybackslash}p{1.6cm}}
\toprule
パラメータ & 値 & 出典 \\
\midrule
感覚領域 & $8$ 本の電極。隣り合う電極が隣り合うボール位置を表すように配置 & 論文 \\
場所による符号 & 電極番号 $k=1,\dots,8$ がラケット中心に対するボールの位置を表す & 論文 \\
どの座標か & ボールの垂直位置。再現ではラケットに対する位置とし、$y-p\in[-\tfrac12,\tfrac12]$ で $k=\lceil 8(y-p+\tfrac12)\rceil$ & 論文。式は選択 \\
頻度による符号 & $4$ から $40$~Hz の刺激頻度がラケットまでの距離を運ぶ & 論文 \\
目盛りの向き & ボールが反対側の壁にあるとき $4$~Hz、ラケット側の壁に近づくにつれ線形に $40$~Hz まで：$f=4+36\,(1-x/L)$。$x$ はラケット側の壁までの距離、$L$ はコートの幅 & 論文 \\
パルス波形 & 短い矩形の二相性パルス。正の相が先、負の相が後 & 論文 \\
振幅 & $75$~mV。抑制された細胞を無理に発火させないように選ばれた & 論文 \\
各相の長さ & 記載なし。刺激システムの標準設定を使い、実験中は変えない & 選択 \\
\bottomrule
\end{tabular}
\end{center}

\subsection*{ラケットの読み出し}

論文では運動領域は二つあり、一方のスパイクはラケットを上へ、他方は下へ動かす。細胞密度と活動の違いを補正するため、領域の寄与には重みがつけられる\cite{Kagan2022}。正確な式は示されていない。再現には規則\eqref{eq:dishreadout}、すなわち $10\ms$ の窓ごとに $\Delta p=c\,(n_\uparrow-n_\downarrow)$ を使い、安静時の平均活動で領域をそろえる重みをつける（選択）。$n_\uparrow$ と $n_\downarrow$ を、ゲーム前に測ったその領域の窓あたり平均カウントで割るのである。係数 $c$ は、平均カウント1個分の差がラケットを1窓あたりコートの高さの $1\%$ だけ動かすように選ぶ（選択）。すると一方の領域が一定に優勢なら、ラケットは $1$~s でコートを横断する。

チップ上の領域の配置は論文本文に座標で与えられておらず、模式図があるだけである。再現には重ならない電極群が三つあれば足りる。中央の感覚電極8本と、その両側の記録電極群で、一方が「上」、他方が「下」である（選択）。

\subsection*{ゲーム}

コートの大きさ、ラケットの長さ、ボールの速さは論文に記載がない。ボールが一定の速さで飛び、壁で跳ね返るとあるだけである。再現には（選択）：コート $1\times1$、右の壁に長さ $0.2$ のラケット（モデル~\texttt{models/dishbrain\_model.py} と同じく $|y-p|<0.1$ でヒット）、ボールは $1$~s でコートを横切る。ラリーで一度も打ち返さずにミスしたものを「エース」、3回を超えて連続で打ち返したラリーを「長いラリー」と数える（論文）。

\subsection*{フィードバック}

\begin{center}
\footnotesize
\setlength{\tabcolsep}{4pt}
\renewcommand{\arraystretch}{1.15}
\begin{tabular}{>{\raggedright}p{2.6cm}>{\raggedright}p{8.3cm}>{\raggedright\arraybackslash}p{1.6cm}}
\toprule
イベント & 与えるもの & 出典 \\
\midrule
ヒット & 予測可能な刺激：感覚電極 $8$ 本すべてに同時に、$75$~mV、$100$~Hz、$100\ms$。その間、通常のボールの刺激は中断する & 論文 \\
ミス & 予測不能な刺激：$150$~mV、$5$~Hz、$4$~s、ランダムな時刻に、$8$ 本からランダムに選んだ電極へ。その後 $4$~s 刺激なし、ボールはランダムな方向にスタート & 論文 \\
ランダムさの法則 & 記載なし。再現では、パルスの時刻を平均頻度 $5$~Hz のポアソン過程とする & 選択 \\
\bottomrule
\end{tabular}
\end{center}

\subsection*{プロトコルと対照}

1セッションは $20$~min で、最初の $5$~min と最後の $15$~min を比べる（論文）。結果の指標は三つ。平均ラリー長、エースの割合、長いラリーの割合である。実験条件（論文）：
\begin{itemize}
\item \emph{刺激}：上に述べた完全なループ。
\item \emph{無音}：ヒットやミスの刺激の代わりに同じ時間まったく刺激せず、その後ボールはランダムな方向にスタートする。
\item \emph{フィードバックなし}：ミスの後、ボールは単に跳ね返って飛び続け、ボール位置の刺激は続く。
\item \emph{安静}：ゲームは進み、ラケットは活動に従って動くが、刺激はまったくない。
\item \emph{対照}：培地だけのチップ。HEK293T 細胞。「コンピュータの中の培養」、つまり細胞の代わりのシミュレーション。
\end{itemize}
主系列の標本数：マウス培養 $9$ 個（$101$ セッション）、ヒト培養 $11$ 個（$138$）、培地のみのチップ $6$ 枚（$80$）、安静の培養 $20$ 個（$42$）、シミュレーション $3$ 個（$38$）、計 $399$ セッション。フィードバックの系列では、1日 $3$ セッションを $3$ 日で $486$ セッション、条件は「刺激」「無音」「フィードバックなし」（$n=27$、$17$、$15$）。最初の $5$ 分から最後の $15$ 分への平均ラリー長の伸びは、生きた培養でだけ得られた。フィードバックの系列で時間とともに有意に伸びたのは「刺激」条件だけだった。セッションの終わりには「無音」条件も「フィードバックなし」を上回ったが、効果は小さく、3日間でセッション間の安定した学習はなかった\cite{Kagan2022}。

\subsection*{テストベンチの1サイクル}

$10\ms$ ごとに、テストベンチのコンピュータは次のことを行う。
\begin{enumerate}
\item 直前の窓における二つの運動領域のスパイク数 $n_\uparrow$、$n_\downarrow$ を読み、安静時のその領域の平均カウントで正規化する。
\item ラケットを $\Delta p=c\,(n_\uparrow-n_\downarrow)$ だけ動かし、コートの端で制限する。
\item ボールを1ステップ進め、上、下、左の壁で跳ね返す。
\item ボールが右の壁に達したら、$|y-p|<0.1$ ならヒットで、ラリーのカウントを1増やし、ヒットの刺激（$100\ms$）を与える。そうでなければミスで、ラリーを記録し、ミスの刺激（$4$~s）を与え、$4$~s 休止してボールを再スタートする。
\item そうでなければ、ボールの垂直位置から電極 $k$ を、ラケット側の壁までの距離から頻度 $f$ を選び、次の窓まで電極 $k$ に $75$~mV の二相性パルスを頻度 $f$ で与える。
\end{enumerate}
これだけあれば、公表されたものと互換なテストベンチを組み、本章の中心的な問いを確かめられる。培養を教えるのは予測可能性なのか、それともヒットとミスへの応答の単なる違いなのか。そのためには論文の三つの条件に、論文にない四つ目を加えるとよい。ミスの後に、予測不能な刺激と同じ長さとエネルギーの\emph{予測可能な}刺激を与える条件である。それでも培養が学習するなら、効いているのは予測可能性ではなく違いである。


\section{まとめ}

\begin{table}[t]
\centering
\footnotesize
\setlength{\tabcolsep}{4pt}
\renewcommand{\arraystretch}{1.15}
\begin{tabular}{>{\raggedright}p{2.1cm}>{\raggedright}p{4.2cm}>{\raggedright}p{1.9cm}>{\raggedright\arraybackslash}p{4.6cm}}
\toprule
テストベンチの部位 & 仕組み & 本書の章 & わかっていること \\
\midrule
入力 & 場所（8本のどの電極か）と頻度（4--40 Hz） & \ref{sec:code}, \ref{sec:sensors} & 実験者が決めたもの \\
出力 & 二つの領域の $10\ms$ ごとのカウントの差\eqref{eq:dishreadout} & \ref{sec:glutcomputer}, \ref{sec:debugging} & 実験者が決めたもの。ノイズは\eqref{eq:rateerror}に従う \\
教師 & ヒットの後に予測可能な刺激、ミスの後に予測不能な刺激。ドーパミンはない & \ref{sec:dofa} & セッション内の有意な伸びは完全なフィードバックのときだけ。無音では効果が小さい。セッション間では不安定 \\
仕組み & 入力の予測か、ミスでのかき混ぜ\eqref{eq:dishdensity}か & \ref{sec:dofa}, \ref{sec:recognition} & 未確定。どちらも結果を説明する \\
ネットワークの状態 & プレイ中に臨界状態へ近づく & \ref{sec:gaba} & プレイの質との関係は測定済み \\
「知性」 & 著者らの主張 & --- & 示されていない。反論がある \\
\bottomrule
\end{tabular}
\caption{エンジニアの目で見た DishBrain。}
\label{tab:dishbrain}
\end{table}

DishBrain は、生きたニューロンでできたマシンの最も単純な形である。入力は符号で与えられ、出力はカウントで読まれ、教師は成功と失敗への応答の違いである（表\ref{tab:dishbrain}）。目も筋肉もドーパミンもないネットワークがほんの少しプレイを覚えた。それだけで、このネットワークは本書のアイデアを試すテストベンチになる。仕組みが予測であれ、かき混ぜであれ、第三の何かであれ、答えは第\ref{sec:debugging}章が勧めるやり方で見つかるだろう。ようやく全体が見えるマシンの上で、プローブとテスト信号と部品の切り離しを使って。

\section{問題}

\begin{exercise}[\lvlA]
\label{ex:22:1}
\emph{カウントをどれだけずらせばよいか。}二つの領域は合わせて $R_\uparrow+R_\downarrow=2000$ スパイク/秒を出し、流れはポアソン過程である。(a)~$R=1000$ と $R=100$ のとき、一つの領域の $10\ms$ あたりのカウントの相対的なばらつきはいくらか。(b)~ボールが飛ぶ $1$~s の間、変位\eqref{eq:dishreadout}は足し合わされる。平均の変位がばらつきの2倍になる最小の $R_\uparrow-R_\downarrow$ はいくらか。
\end{exercise}

\begin{exercise}[\lvlA]
\label{ex:22:2}
\emph{学習を見るには何ラリー必要か。}ラリーは独立で、打ち返した割合は二項分布に従う。打ち返した割合の伸びが差の標準偏差の2倍を超えるには、二つの期間それぞれに何ラリー必要か。(a)~$0.20$ から $0.25$ へ。(b)~モデルのように $0.21$ から $0.38$ へ。
\end{exercise}

\begin{exercise}[\lvlB]
\label{ex:22:3}
\emph{\eqref{eq:dishdensity}の導出。}設定 $\boldsymbol\psi$ は有限集合を動く。ミス（確率 $P(\boldsymbol\psi)>0$）のときネットワークは対称な確率 $K(\boldsymbol\psi,\boldsymbol\psi')$ で $\boldsymbol\psi'$ に移り、ヒットのときはとどまる。(a)~$\rho\propto1/P$ が定常分布であることを証明せよ。(b)~$P=0.9$ と $P=0.1$ の二つの設定のとき、ミスの割合はいくらか。
\end{exercise}

\begin{exercise}[\lvlB]
\label{ex:22:4}
\emph{モデルの吸収領域。}ラケットは $p=\min\bigl(1,\max(0,ay+b)\bigr)$ に、ボールは高さ $y\sim U(0,1)$ に来て、$|p-y|<h=0.1$ ならヒット。(a)~$|b|<h$ かつ $|a-1+b|<h$ ならネットワークがすべてのボールを打ち返すことを証明し、この領域の面積を求めよ。(b)~$a\sim U(-1,1)$、$b\sim U(0,1)$ のときの平均の打ち返し割合を数値的に求めよ。
\end{exercise}

\begin{exercise}[\lvlB]
\label{ex:22:5}
\emph{シミュレーション：設定の周りの柵。}\texttt{dishbrain\_model.py} のコピーで、$200$ 個の培養にそれぞれ $20\,000$ ラリーを与えよ。最後の $500$ 回で打ち返した割合はいくらか。揺さぶりのたびに設定を長方形 $a\in[-1,2]$、$b\in[-1,1]$ に切り詰めたらどうなるか。違いを問題\ref{ex:22:3}で説明せよ。
\end{exercise}

\begin{exercise}[\lvlB]
\label{ex:22:6}
\emph{入力符号の設計。}高さの誤差が $h=0.1$ 未満ならボールは打ち返される。ネットワークは入力を $T=0.25$~s で読み、頻度 $r$ までの範囲で約 $\sqrt{rT}$ 段階を区別でき、刺激は毎秒 $40$ パルス以下である。(a)~8本の電極での場所による符号で足りるか。(b)~1本の電極で高さを頻度として伝えられるか。
\end{exercise}

\begin{exercise}[\lvlC]
\label{ex:22:7}
\emph{研究：予測か、かき混ぜか。}モデルを $d$ 個の調整可能な数（$p=\sum_{i<d}a_iy^i$）に一般化せよ。打ち返し割合 $0.5$ に達するまでのラリー数は、かき混ぜと誤差の符号を使う規則\eqref{eq:perturbation}とで、$d$ とともにどう増えるか。テストベンチでのどんな実験が二つの仮説を区別するか。
\end{exercise}

\appendix
\renewcommand{\chaptersuffix}{}
\chapter{ニューロンの型を決めない物質}
\label{sec:othermediators}

第\ref{sec:neurontypes}章では、ニューロンを主な神経伝達物質で分類し、10の型を得た（表\ref{tab:types}）。ここで話を複雑にしておこう。ニューロンが放出するのは主な伝達物質だけではなく、脳では主な伝達物質以外の物質も働いている。それらはネットワークが信号を伝え、蓄える仕方を変える。

アドレスバスとブロードキャストバス（\ref{sec:buses}節）のほかに、3本目のバスがある。\emph{逆向きチャネル}である。いくつかの物質は送り手ではなく\emph{受け手}が放出し、間隙を逆向きに送り手の出力側へ戻り、送り手が次に放出する伝達物質の量を変える。通信ネットワークでいえば、送り手が自分の送信を調整するのに使う受信確認に似ている。ニューロンが結合の重みを変えるときに使うのがまさにこのチャネルである。送り手の出力側はこれを通じて受け手の活動、つまり第\ref{sec:dofa}章の学習則の2番目の因子を知る。

既知の神経伝達物質の完全なリストは長い。100を超える物質があり、その大部分は短いタンパク質の鎖、神経ペプチドである。しかしどれもいくつかのクラスに収まり、クラス内の物質は似た振る舞いをする。主な伝達物質にならない物質を、エンジニアが問うであろう三つの問い、つまり信号がどのバスを通るか、速く効くか、ふつうの符号は何か、とともに表\ref{tab:othermediators}にまとめた。これらはニューロンの型を決めない。主な伝達物質と一緒に放出されるか、ニューロン以外から放出されるか、逆向きに進むかのいずれかだからである。

\begin{table}[t]
\centering
\footnotesize
\setlength{\tabcolsep}{4pt}
\renewcommand{\arraystretch}{1.12}
\begin{tabular}{>{\raggedright}p{2.0cm}>{\raggedright}p{2.6cm}>{\raggedright}p{2.55cm}>{\raggedright}p{1.3cm}>{\centering}p{0.8cm}>{\raggedright\arraybackslash}p{4.2cm}}
\toprule
クラス & 物質 & バス & 速さ & 符号 & 計算での役割 \\
\midrule
アミノ酸 & D-セリン & 局所 & 遅い & AND & グルタミン酸受容体の第二の鍵：「AND」条件 \\
\midrule
モノアミン & 微量アミン & ブロードキャスト & 遅い & $\pm$ & モノアミンチャネルの微調整 \\
\midrule
\multirow{2}{2.0cm}{プリン}
 & ATP & アドレス & 速・遅 & $+$ & 共伝達物質。ニューロンとグリアの間の信号 \\
 & アデノシン & 容積 & 遅い & $-$ & 活動とともに蓄積する：負荷カウンタ\cite{Dunwiddie2001} \\
\midrule
\multirow{8}{2.0cm}{神経ペプチド（100種以上）}
 & エンケファリン、エンドルフィン、ダイノルフィン & 容積 & 遅い & $-$ & 伝達の抑制 \\
 & サブスタンス P & 容積 & 遅い & $+$ & 遅い増強 \\
 & ニューロペプチド Y & 容積 & 遅い & $-$ & 遅い抑制 \\
 & ソマトスタチン & 容積 & 遅い & $-$ & 遅い抑制。GABA と共放出 \\
 & 血管作動性腸管ペプチド（VIP） & 容積 & 遅い & $+$ & 脱抑制：抑制性ニューロンの抑制 \\
 & コレシストキニン & 容積 & 遅い & $\pm$ & 一部の抑制性ニューロンで GABA と共放出 \\
 & オレキシン & ブロードキャスト & 遅い & $+$ & 覚醒モードの安定化 \\
 & オキシトシン、バソプレシン & ブロードキャスト & 遅い & $\pm$ & 行動の大域的な設定 \\
\midrule
\multirow{2}{2.0cm}{気体}
 & 一酸化窒素 NO & 逆向き、容積 & 遅い & $\pm$ & 膜を通り抜ける。重みが変わるときの逆向き信号\cite{Garthwaite2008} \\
 & CO、H$_2$S & 容積 & 遅い & $\pm$ & 役割はほとんど研究されていない \\
\midrule
脂質 & 内因性カンナビノイド（アナンダミド、2-AG） & 逆向き & 遅い & $-$ & 受け手が送り手の放出を減らす：重みのフィードバック\cite{WilsonNicoll2001} \\
\bottomrule
\end{tabular}
\caption{ニューロンの型を決めない物質。\emph{バス}、\emph{速さ}、\emph{符号}は表\ref{tab:types}と同じ。加えて、容積バスは伝達物質が放出場所の周りの細胞間隙に広がること、逆向きは受け手から送り手へ戻ること、局所は放出場所の近くで作用することを表す。「AND」は許可信号である。それだけでは電流を起こさないが、それがないと別の入力が働かない。論理ゲート「AND」の2番目の入力のようなものである。神経ペプチドはほぼ常に主な伝達物質と一緒に、主に高いスパイク頻度のときに放出される\cite{vandenPol2012}。}
\label{tab:othermediators}
\end{table}

\chapter{記号}
\label{sec:notation}

ニューロンの型は、主な神経伝達物質の名前の最初の4文字で表す。\nt{GLUT}：グルタミン酸、\nt{GABA}：$\gamma$-アミノ酪酸、\nt{GLYC}：グリシン、\nt{ACET}：アセチルコリン、\nt{DOPA}：ドーパミン、\nt{SERO}：セロトニン、\nt{HIST}：ヒスタミン、\nt{NORA}：ノルアドレナリン、\nt{ADRE}：アドレナリン、\nt{ASPA}：アスパラギン酸（表\ref{tab:types}）。

文字の数は量の数より少ないので、同じ文字が章によって別のものを表すことがある。下の表では、意味ごとに1行を割り当てた。右の列は、その記号を導入した式である。

\begin{center}
\footnotesize
\setlength{\tabcolsep}{4pt}
\renewcommand{\arraystretch}{1.12}
\begin{longtable}{>{\raggedright}p{2.0cm}>{\raggedright}p{9.6cm}>{\raggedright\arraybackslash}p{2.4cm}}
\toprule
記号 & 意味 & 導入箇所 \\
\midrule
\endfirsthead
\toprule
記号 & 意味 & 導入箇所 \\
\midrule
\endhead
\bottomrule
\endfoot
$a$ & 線形伝達特性の傾き、$f(u)=au$ & \eqref{eq:feedback} \\
$a$ & \nt{ACET} のレベル：$0$ は読み出し、$1$ は書き込み & \eqref{eq:acetnet} \\
$a_i$, $a$ & 群の疲労：リズム発生器で。徐波睡眠のシーソーで & \eqref{eq:halfcentre}, \eqref{eq:updown} \\
$a_S$, $a_P$ & 心拍リズムのアクセルとブレーキに対する感度 & \eqref{eq:gasbrake} \\
$A(r)$, $A_0$ & 頻度 $r$ でのスパイクに対するシナプスの応答。休んだシナプスでの応答 & \eqref{eq:depression} \\
$A_1$ & 伝達物質1包に対する受け手の応答 & \eqref{eq:quantal} \\
$A$ & 樹状突起を含むニューロンの膜面積 & \eqref{eq:dendriteload} \\
$A_f$, $A_s$ & 次の運動まで保持される補正の割合（速い過程と遅い過程） & \eqref{eq:tworate} \\
$b_i$ & ペースメーカー自身の流入 & \eqref{eq:seroneuron} \\
$b$ & リズム発生器や睡眠のシーソーで、働きが群を疲れさせる速さ & \eqref{eq:halfcentre}, \eqref{eq:updown} \\
$b$ & デジタル化の1サンプルあたりのビット数 & \eqref{eq:probedata} \\
$B_f$, $B_s$ & 補正が誤差から学習する強さ（速い過程と遅い過程） & \eqref{eq:tworate} \\
$c$ & 空間内の伝達物質の濃度 & \eqref{eq:seroneuron} \\
$c$ & 隣り合うニューロンのノイズの相関係数 & \eqref{eq:correlated} \\
$c$ & カウントの差がラケットを動かす係数 & \eqref{eq:dishreadout} \\
$c_f$ & $C$ ニューロンの応答：すべての位置にわたる $s_{f,p}$ の最大値 & \eqref{eq:scstage} \\
$c_m$ & 膜の比容量 & \eqref{eq:dendriteload} \\
$C$ & 努力のコスト：時間 $\tau_a$ での行動のコストは $C/\tau_a$ & \eqref{eq:vigor} \\
$d'$ & 二つの入力の弁別度：差をノイズで割ったもの & \eqref{eq:gainsnr} \\
$d$, $d_j$ & 信号経路上の遅延 & \eqref{eq:glutneuron}, \eqref{eq:vstar} \\
$d$ & 筋肉とのフィードバックループの遅延 & \eqref{eq:delaystable} \\
$D$ & 拡散係数 & \eqref{eq:diffusionlength} \\
$D$ & \nt{GLUT}--\nt{GABA} ネットワークの頻度の式の分母 & \eqref{eq:isn} \\
$D$ & 遅延報酬の待ち時間 & \eqref{eq:patience} \\
$e$, $e_{ij}$ & シナプスの適格度トレース & \eqref{eq:threefactor} \\
$e$, $e_i$ & 誤差配線を通って届く、出力や運動の誤差 & \eqref{eq:lms}, \eqref{eq:adaptivefilter}, \eqref{eq:tworate} \\
$E_{\text{spk}}$ & シナプスの働きを含むスパイク1回のエネルギー & \eqref{eq:energy} \\
$E$ & \nt{GLUT} 群の頻度 & \eqref{eq:ei} \\
$E_E$ & 興奮性シナプスの平衡電位 & \eqref{eq:shunt} \\
$f$, $F$ & 伝達特性：入力の関数としての頻度 & \eqref{eq:ratenet}, \eqref{eq:gain} \\
$f$, $f_0$ & 心拍数。ペースメーカー自身のリズム & \eqref{eq:gasbrake} \\
$f_s$ & 電極1本のサンプリング周波数 & \eqref{eq:probedata} \\
$F(t)$, $F_1$ & 筋力。単収縮1回の力 & \eqref{eq:forcesum} \\
$F_1$, $F_2$ & 関節を逆向きに引く二つの筋肉の力 & \eqref{eq:antagonists} \\
$g$ & 受容体の濃度に対する感度 & \eqref{eq:seroneuron} \\
$g_L$, $g_E$, $g_I$ & 漏れ、興奮、抑制のコンダクタンス & \eqref{eq:shunt} \\
$g$ & \nt{NORA} が設定するゲイン & \eqref{eq:gain}, \eqref{eq:machine} \\
$g$ & 記憶ネットワークでの \nt{GABA} による全体の抑制 & \eqref{eq:acetnet} \\
$g$ & 上から設定される痛みのゲートのゲイン & \eqref{eq:gate} \\
$g$ & ホルモンカスケード1段のゲイン & \eqref{eq:cascade} \\
$g_j$ & 適応フィルタの入力側にある、固有の時定数をもつフィルタ $j$ & \eqref{eq:adaptivefilter} \\
$G$ & フィードバックループのゲイン & \eqref{eq:feedback} \\
$G$ & 遅延のあるループでの修正の速さ & \eqref{eq:delaystable} \\
$h$, $h_I$ & \nt{GLUT} 群と \nt{GABA} 群への外部入力 & \eqref{eq:ei}, \eqref{eq:isn} \\
$h(t)$ & 筋肉の単収縮1回 & \eqref{eq:twitch} \\
$H(t)$ & 睡眠圧の上側の閾値：ここに達するとマシンは眠る & \eqref{eq:sleeppressure} \\
$I$, $I_i$ & ニューロンや群への外部からの流入 & \eqref{eq:ratenet}, \eqref{eq:wta}, \eqref{eq:updown} \\
$I$ & \nt{GABA} 群の頻度 & \eqref{eq:ei} \\
$I$ & 刺激の明るさや強度 & \eqref{eq:photonnoise}, \eqref{eq:nakarushton} \\
$I$ & 膜を充電する入力電流 & \eqref{eq:dendriteload} \\
$k$ & バースト内のスパイク数 & \eqref{eq:burst} \\
$k$ & 増分がノイズの何倍でなければならないか & \eqref{eq:photonnoise} \\
$k$ & ホルモンカスケードのフィードバック係数 & \eqref{eq:cascade} \\
$k_s$ & 痛みが感作を押し上げる強さ & \eqref{eq:sensitization} \\
$K$ & 結果の特徴の数 & \eqref{eq:vectortd} \\
$K$ & 関節の剛性 & \eqref{eq:antagonists} \\
$K$ & ホルモンカスケードのループゲイン、$K=g^3k$ & \eqref{eq:cascade}, \eqref{eq:cascadestable} \\
$K_\varphi$ & 概日発振器が外部信号に同調する強さ & \eqref{eq:pll} \\
$L$ & 遅延線の長さ & \eqref{eq:itd} \\
$L$ & 痛みのセンサーからの配線の長さ & \eqref{eq:paintime} \\
$L(t)$ & 睡眠圧の下側の閾値：ここに達するとマシンは目覚める & \eqref{eq:sleeppressure} \\
$M$ & 予兆と報酬の因果関係の尺度 & \eqref{eq:retro} \\
$M_k$ & 特徴 $k$ の期待値 & \eqref{eq:vectortd} \\
$M_{ik}$ & \nt{GABA} ニューロン $k$ からの抑制の強さ & \eqref{eq:machine} \\
$M$ & 関節のトルク & \eqref{eq:antagonists} \\
$M$ & ミキサーの入力線の数 & \eqref{eq:cover} \\
$n$, $\bar n$ & 窓内のスパイク数。その平均 & \eqref{eq:rateerror} \\
$n$ & 閾値素子の入力数 & \eqref{eq:cover} \\
$n_\uparrow$, $n_\downarrow$ & 「上」と「下」の領域の窓あたりのスパイク数 & \eqref{eq:dishreadout} \\
$N$ & ニューロン数 & \eqref{eq:correlated}, \eqref{eq:energy} \\
$N$ & プローブの電極数 & \eqref{eq:probedata} \\
$N_s$ & 2個のニューロンの間の放出部位の数 & \eqref{eq:quantal} \\
$p$ & スパイク1回あたりの伝達物質放出確率 & \eqref{eq:burst} \\
$p$ & 記憶ネットワークで活動しているニューロンの割合 & \eqref{eq:acetnet} \\
$p$ & 補償を学習する外乱 & \eqref{eq:tworate} \\
$P$ & 信号に使われる電力 & \eqref{eq:energy}, \eqref{eq:dishpower} \\
$P$ & 保持している推定値の不確かさ & \eqref{eq:kalman} \\
$P$ & 「ブレーキ」の活動：心臓への \nt{ACET} 配線 & \eqref{eq:gasbrake} \\
$P_{\max}$ & 閾値素子が二つのクラスに分けられるランダムなパターンの数 & \eqref{eq:cover} \\
$P_{\text{miss}}$ & ミスの確率 & \eqref{eq:dishdensity} \\
$q$ & 世界が変わる速さ & \eqref{eq:kalman} \\
$q_k$ & \nt{GABA} ニューロン $k$ の頻度 & \eqref{eq:machine} \\
$q_0$ & 抑制性介在ニューロンの背景活動 & \eqref{eq:gate} \\
$q$ & 痛みのゲートにおける抑制性介在ニューロンの頻度 & \eqref{eq:gate} \\
$q$ & 次の運動単位が前の運動単位の何倍強いか & \eqref{eq:sizeprinciple} \\
$r$, $r_i$ & ニューロンのスパイク頻度 & \eqref{eq:ratenet} \\
$r$, $r_t$ & 報酬（報酬の流れ） & \eqref{eq:rpe}, \eqref{eq:ctd} \\
$\mathbf r(t)$ & リザバーの応答ベクトル & \eqref{eq:reservoir} \\
$R$, $\bar R$ & 得られた報酬。その期待値、または単位時間あたりの平均 & \eqref{eq:perturbation}, \eqref{eq:vigor} \\
$R$ & フィルタ入力のノイズの分散 & \eqref{eq:kalman} \\
$R$, $r_0$ & 遅延報酬と即時報酬 & \eqref{eq:patience} \\
$R_{\max}$ & 最大伝送速度、bit/s & \eqref{eq:spikeentropy} \\
$r_{\max}$ & 最大スパイク頻度 & \eqref{eq:gain}, \eqref{eq:nakarushton} \\
$R_{\text{raw}}$, $R_{\text{spike}}$ & プローブのデータ量：生データとスパイクのみ、bit/s & \eqref{eq:probedata} \\
$s_j$ & ニューロン $j$ の出力の符号 & \eqref{eq:dalesign} \\
$s_i$ & 受け手の受容体の符号 & \eqref{eq:seroneuron} \\
$s$, $\hat s$ & 世界の状態。その推定値 & \eqref{eq:rpe}, \eqref{eq:kalman} \\
$s_{f,p}$ & 位置 $p$ における特徴 $f$ への $S$ ニューロンの応答 & \eqref{eq:scstage} \\
$S$ & 「アクセル」の活動：心臓への \nt{NORA} 配線 & \eqref{eq:gasbrake} \\
$S$ & 睡眠圧 & \eqref{eq:sleeppressure} \\
$t_1$ & 最初のスパイクの時刻 & \eqref{eq:latency} \\
$t_k$ & $k$ 番目のスパイクの時刻 & \eqref{eq:forcesum} \\
$t_{f,i}$ & 特徴 $f$ のテンプレート & \eqref{eq:scstage} \\
$T$ & スパイクを数える窓。光子の蓄積時間 & \eqref{eq:rateerror}, \eqref{eq:photonnoise} \\
$T_c$ & 筋肉の単収縮の立ち上がり時間 & \eqref{eq:twitch} \\
$T_{\text{burst}}$ & 培養の一斉発火の間隔 & \eqref{eq:burstinterval} \\
$u$ & ニューロンへの入力の総和 & \eqref{eq:gain} \\
$u$, $u(t)$ & 脳の指令：適応フィルタの入力として。ホルモンカスケードへの指令として & \eqref{eq:adaptivefilter}, \eqref{eq:cascade} \\
$u(t)$ & リザバーの入力 & \eqref{eq:reservoir} \\
$U$ & 一定の入力のレベル & \eqref{eq:latency} \\
$U$ & スパイク1回で放出される伝達物質の蓄えの割合 & \eqref{eq:depression} \\
$v$ & 配線中の信号の速さ。光点の動く速さ & \eqref{eq:itd}, \eqref{eq:vstar} \\
$V$ & 膜電圧 & \eqref{eq:glutneuron}, \eqref{eq:shunt} \\
$V$ & 将来の報酬の期待値 & \eqref{eq:rpe}, \eqref{eq:ctd} \\
$w$, $w_{ij}$, $W_{ij}$ & 結合の重み & \eqref{eq:glutneuron}, \eqref{eq:ratenet}, \eqref{eq:acetnet}, \eqref{eq:lms}, \eqref{eq:adaptivefilter} \\
$\mathbf w_k$ & $C$ ニューロン $k$ の入力の重み & \eqref{eq:traceinvariance} \\
$w_{q\mu}$, $w_{q\nu}$, $w_{z\mu}$ & 痛みのゲートの結合の重み & \eqref{eq:gate} \\
$W_{\text{out}}$ & リザバーの線形読み出しの重み & \eqref{eq:reservoir} \\
$x$, $y$ & 論理入力と論理出力 & \eqref{eq:dualrail}, \eqref{eq:veto} \\
$x$, $y$ & 送り手と受け手の活動 & \eqref{eq:hebb} \\
$x$ & 遅延線上の検出器の位置 & \eqref{eq:itd} \\
$x_i$ & 感覚器からの入力 & \eqref{eq:acetnet} \\
$x_p$ & 位置 $p$ の入力 & \eqref{eq:scstage} \\
$\mathbf x$ & $S$ ニューロンの出力、$C$ ニューロンの入力 & \eqref{eq:traceinvariance} \\
$x_j$ & 適応フィルタの入力 $j$：フィルタ $g_j$ を通った指令 & \eqref{eq:lms}, \eqref{eq:adaptivefilter} \\
$x_f$, $x_s$ & 速い過程と遅い過程の補正 & \eqref{eq:tworate} \\
$x_1$, $x_2$, $x_3$ & ホルモンカスケード3段のレベル。$x_3$ は最終ホルモン & \eqref{eq:cascade} \\
$x^*$ & 新しい雪崩が始まる、伝達物質の蓄えの回復割合 & \eqref{eq:burstinterval} \\
$y$ & フィルタのノイズを含む入力 & \eqref{eq:kalman} \\
$y$, $\hat y$ & センサーの信号。適応フィルタによるその予測 & \eqref{eq:adaptivefilter} \\
$y_k$, $\bar y_k$ & $C$ ニューロン $k$ の活動。そのトレース & \eqref{eq:traceinvariance} \\
$y(t)$ & リザバーの読み出しの出力 & \eqref{eq:reservoir} \\
$z$ & 痛みのゲートの出力の頻度 & \eqref{eq:gate} \\
\midrule
$\alpha_i^\pm$ & 正の誤差と負の誤差の重み & \eqref{eq:asymmetric} \\
$\beta$ & 相互抑制の強さ & \eqref{eq:wta} \\
$\beta$ & 痛みのゲートの出力に対する抑制の強さ & \eqref{eq:gate} \\
$\gamma$ & 割引率 & \eqref{eq:rpe} \\
$\delta(\cdot)$ & デルタ関数：瞬間的な跳び & \eqref{eq:glutneuron} \\
$\delta$, $\delta_t$ & 報酬予測誤差 & \eqref{eq:rpe}, \eqref{eq:ctd} \\
$\delta t$ & 両耳への音の到達時間差 & \eqref{eq:itd} \\
$\Delta$, $\Delta_{\max}$ & スパイクの間隔。同時性の窓 & \eqref{eq:window} \\
$\Delta t$ & 受け手が時間を区別できる精度 & \eqref{eq:spikeentropy} \\
$\Delta F$ & 次の運動単位による力の増分 & \eqref{eq:sizeprinciple} \\
$\Delta I$ & 区別できる明るさの増分 & \eqref{eq:photonnoise} \\
$\Delta p$ & 窓あたりのラケットの移動 & \eqref{eq:dishreadout} \\
$\Delta u$ & 二つの群の入力の差 & \eqref{eq:gainsnr} \\
$\Delta V$ & 膜電圧をどれだけ上げる必要があるか & \eqref{eq:dendriteload} \\
$\Delta x$ & 隣り合うセンサーの間隔 & \eqref{eq:vstar} \\
$\varepsilon$ & 入力の相対的な差 & \eqref{eq:wtagain} \\
$\eta$ & 学習率 & \eqref{eq:plasticity}, \eqref{eq:lms}, \eqref{eq:traceinvariance} \\
$\eta$ & 増幅器の後のネットワーク内のノイズ & \eqref{eq:softmax} \\
$\mu$ & 痛みのゲートへの触覚入力 & \eqref{eq:gate} \\
$\nu$ & 痛みの入力 & \eqref{eq:gate} \\
$\theta$ & 発火の閾値 & \eqref{eq:window}, \eqref{eq:scstage} \\
$\kappa$ & スパイク1回あたりの伝達物質の量 & \eqref{eq:seroneuron} \\
$\kappa_i$ & 正の誤差と負の誤差の重みの比 & \eqref{eq:expectile} \\
$\kappa_m$ & 筋力と筋肉の剛性の関係 & \eqref{eq:antagonists} \\
$\ell$ & 伝達物質が広がる距離 & \eqref{eq:diffusionlength} \\
$\ell$ & 筋肉が関節を引くモーメントアーム & \eqref{eq:antagonists} \\
$\xi$ & ニューロン出力のランダムな揺らぎ & \eqref{eq:perturbation} \\
$\rho(\boldsymbol\psi)$ & 設定 $\boldsymbol\psi$ で過ごす時間の割合 & \eqref{eq:dishdensity} \\
$\sigma$ & ばらつき、ノイズの尺度 & \eqref{eq:rateerror}, \eqref{eq:softmax} \\
$\sigma$ & 応答が最大の半分になる明るさ & \eqref{eq:nakarushton} \\
$\sigma_{\text{pre}}$, $\sigma_{\text{post}}$ & 増幅器の前と後のノイズ & \eqref{eq:gainsnr} \\
$\tau$ & 膜の時定数 & \eqref{eq:glutneuron} \\
$\tau$ & ホルモンカスケード1段の時定数 & \eqref{eq:cascade} \\
$\tau_{\text{eff}}$ & 自己興奮する群の時定数 & \eqref{eq:perfectint} \\
$\tau_c$ & 空間内の伝達物質の寿命 & \eqref{eq:seroneuron} \\
$\tau_E$, $\tau_I$ & \nt{GLUT} 群と \nt{GABA} 群の時定数 & \eqref{eq:ei} \\
$\tau_e$ & 適格度トレースの寿命 & \eqref{eq:threefactor} \\
$\tau_{\text{tr}}$ & $C$ ニューロンの活動のトレースの寿命 & \eqref{eq:traceinvariance} \\
$\tau_s$ & 損傷後に感度が戻るまでの時間 & \eqref{eq:sensitization} \\
$\tau_r$, $\tau_d$ & 覚醒中に睡眠圧がたまる時間と、睡眠中に解消される時間 & \eqref{eq:sleeppressure} \\
$\tau_{\text{rec}}$ & 伝達物質の蓄えの回復時間 & \eqref{eq:depression} \\
$\tau_\gamma$ & 計画の視野 & \eqref{eq:ctd} \\
$\tau_v$ & 期待値が更新される速さ & \eqref{eq:ramp} \\
$\tau_a$ & リズム発生器や睡眠のシーソーにおける疲労の時間 & \eqref{eq:halfcentre}, \eqref{eq:updown} \\
$\tau_a^*$ & 行動を実行する最適な時間 & \eqref{eq:vigor} \\
$\Phi$ & 学習則の右辺 & \eqref{eq:plasticity} \\
$\Phi$ & 光子の流束 & \eqref{eq:photonnoise} \\
$\varphi_k$ & 得られた結果の特徴 $k$ & \eqref{eq:vectortd} \\
$\varphi$ & 概日発振器と外部信号の位相差 & \eqref{eq:pll} \\
$\boldsymbol\psi$ & DishBrain モデルにおけるネットワークの設定 & \eqref{eq:dishdensity} \\
$\omega$, $\omega_0$ & 外部信号（光）の周波数。概日発振器の固有周波数 & \eqref{eq:pll} \\
\end{longtable}
\end{center}

\chapter{解答}
\label{sec:answers}

各章末の問題の略解を示す。完全な解答は教員用の別冊にある。

\answer{ex:01:1}{(a)~$8\cdot10^{10}$個の\nt{GLUT}は地球の人口の$10$倍。ブロードキャスト型は$\approx1.9\cdot10^6$個で、ウィーンほどの都市。(b)~終末は$\approx8\cdot10^{14}$個、うちブロードキャスト型は$\approx3\cdot10^{11}$個で割合は$\approx4\cdot10^{-4}$。これらのニューロンの割合（$2.2\cdot10^{-5}$）の$\approx16$倍。}
\answer{ex:01:2}{(a)~$\tau_c\ln(c_0/c_*)\approx4.6\ms$。(b)~$T\ge\tau_c\ln101$なら線路は空く。放出は毎秒$\approx220$回ではなく$\approx22$回まで。}
\answer{ex:01:3}{$V(s_k)=\gamma^{K-1-k}r$。$\delta_{\text{ランプ}}=\gamma^Kr=re^{-T/\tau_\gamma}\approx0.60\,r$（$\Delta$によらない）、$\tau_\gamma\approx1.95\,\text{s}$。}
\answer{ex:01:4}{$n^*=93$、$47$、$25$、$12$。小さな$\alpha$では$n^*\approx5/\alpha$。}
\answer{ex:01:5}{\nt{GLUT}：$1\to10\to10^4\to10^7$で、ニューロン$\approx10^7$個、$4$段、$4\ms$。\nt{DOPA}：$1\to10\to3.2\cdot10^5$で、ニューロン$\approx3\cdot10^5$個、$3$段で、$30$分の1。\nt{DOPA}ニューロンの数と同じオーダーである。}
\answer{ex:01:6}{$\lambda=f_EJ_E-f_IJ_I$より$J_I=37.5$。電流の比は$f_IJ_I/f_EJ_E=0.94$。抑制をわずか$6.7\,\%$弱めるだけで雪崩（$\lambda\ge1$）になる。}
\answer{ex:01:7}{各状態で遅延ごとに$n\ge57$回の提示。同じ減衰は、たとえば$T$とともに増える内部の時間推定の不正確さからも生じうる。}

\answer{ex:02:1}{長さ$3.2$~mm、間隔$25\um$の列で、検出器は$\approx129$個。発火するのは幅$v\,\tau\ln2=1.7$~mmの帯、つまり$\approx69$個で、列の半分以上。方向は帯の中心。}
\answer{ex:02:2}{窓は$-\tau\ln\frac{w_2}{\theta-w_1}\le\delta\le\tau\ln\frac{w_1}{\theta-w_2}$、すなわち$[-0.235,\,0.178]\ms$。中心は$c=\frac\tau2\ln\frac{w_1(\theta-w_1)}{w_2(\theta-w_2)}=-28$~\textmu s。方向は強い耳の側へ$28$~\textmu sずれる。分解能のほぼ3刻み分である。}
\answer{ex:02:3}{$s_iw_{ij}s_j\ge0$となる$s_i$が必要で、それはサイクルが偶数のとき、かつそのときに限り存在する。相互抑制は帰着する（持続的な振動はない）。\nt{GLUT}--\nt{GABA}ループは帰着しない（振動しうる）。}
\answer{ex:02:4}{6個のニューロン$g_k=[x+y+z\ge k]$と$\bar g_k=[\bar x+\bar y+\bar z\ge4-k]$、$k=1,2,3$。$C=g_2$、$\bar C=\bar g_2$、$S=[g_1+\bar g_2+g_3\ge2]$、$\bar S=[\bar g_1+g_2+\bar g_3\ge2]$で、計8個。ANDとORからなら$12$個。}
\answer{ex:02:5}{$T_{\min}(0.6)=0.718\,\tau$。$I_c=0.225$で、閾値付近では$T_{\min}\propto(I-I_c)^{-1/2}$。}
\answer{ex:02:6}{(a)~$500$段、ニューロン$5\cdot10^4$個。ばらつきは$4.5\ms$（$0.45\,\%$）、相対誤差は$\propto T^{-1/2}$。(b)~すべての段に共通な、ばらつき$5\,\%$のランダムな遅延係数。}
\answer{ex:02:7}{$\tau_F\ln2=0.69\,\text{s}$ごとにリフレッシュ：ニューロン1個あたり毎秒$4.3$スパイク対$20$で、$4.6$倍経済的（スパイク1個あたり$10^{-10}$~Jとして$4\cdot10^{-7}$対$2\cdot10^{-6}$~J）。フリップフロップはビットを失う。コンデンサは、リフレッシュから$0.49\,\text{s}$以内に黙らせ始めた場合にビットを保つ（確率$0.71$）。}

\answer{ex:03:1}{$g_E=g_L$で$35\to17.5\mV$、$g_E=4g_L$で$56\to40\mV$。引き算なら$38.5\mV$になる。シャントは$g_E$を$3$で割る。窓は半分に狭まる（$\tau_{\text{eff}}$が$2$分の1になる）。}
\answer{ex:03:2}{$\operatorname{tr}=\frac{w_{EE}-1}{\tau_E}-\frac1{\tau_I}$、$\det=\frac{1-w_{EE}+w_{EI}w_{IE}}{\tau_E\tau_I}$。境界では$\omega=\sqrt{\det}$、$\tau_I=20\ms$、周期$73\ms$。}
\answer{ex:03:3}{$d_c=\tau_E\arccos(-1/G)/\sqrt{G^2-1}$、周期$2\pi\tau_E/\sqrt{G^2-1}\in(2d_c,4d_c)$。$G=2$：$d_c=12.1\ms$、周期$36\ms$。$G=5$：$d_c=3.6\ms$、周期$12.8\ms$で、ちょうど速いリズムの帯域である。}
\answer{ex:03:4}{(a)~上げるときは$I_2=\beta I_1=2$、下げるときは$I_2=I_1/\beta=0.5$で、幅$1.5$のヒステリシス。(b)~$\varepsilon$が1桁小さくなるごとに$\tau\ln10/(\beta-1)=2.3\tau$ずつ増える。}
\answer{ex:03:5}{3個、$d_k=5$、$10$、$15\ms$。禁止は$15\ms$を覆う必要があり、1個あたり$5\ms$。}
\answer{ex:03:6}{$n+M+1=3+4+1=8$個。2個なら：\nt{GABA} $=[x+y+z\ge2]$、出力 $=[x+y+z-2\,\nt{GABA}\ge1]$。}
\answer{ex:03:7}{$\binom84=70<100\le\binom94=126$より仲介役は$9$個（シュペルナーの定理）。直接の結合がなければ$100$個：$\beta(J-I)$の階数が$n$だから。}
\answer{ex:03:8}{$h_c=\frac1{2w_{EE}}+\frac{w_{EI}w_{IE}-w_{EE}}{4w_{EE}^2}=\frac7{18}\approx0.39$。シミュレーションでは$h=0.38$と$0.40$の間で符号が変わる。}

\answer{ex:04:1}{(a)~$\approx1.1\cdot10^{11}$~bit/J。(b)~$\log_2\sqrt{10}\approx1.7$ビット対$81$ビットで、$\approx50$分の1。}
\answer{ex:04:2}{$r^*=(P_0/E_{\text{spk}})\ln\bigl(1/(r^*\Delta t)\bigr)$、$P_0/E_{\text{spk}}=1.16\,\text{s}^{-1}$より、$r^*\approx6$スパイク/s、$7.4\cdot10^{10}$~bit/J。}
\answer{ex:04:3}{$\mathcal I=\frac{N}{1+(N-1)c}=9.9$（極限は$10$）対$\mathcal I=\frac{N}{1-c}=1111$で、$\approx110$倍。相関が有害なのは、信号に似ているときだけである。}
\answer{ex:04:4}{$\theta\approx68.7$と$19.9$（$w$を単位として）、$m=19$と$m=10$。速い受け手は、平均入力が5分の1なのに、同時入力がほぼ半分で足りる。}
\answer{ex:04:5}{$U\tau_{\text{rec}}\ge0.725\,\text{s}$かつ$\tau_{\text{rec}}(1-2U)\le0.05\,\text{s}$で、これには$U\ge0.483$が必要。最小の$\tau_{\text{rec}}=0.725\,\text{s}$は$U=1$のとき（$U=0.5$なら$1.45\,\text{s}$）。}
\answer{ex:04:6}{(a)~$\theta\,e/(e-1)\approx1.58\,\theta$で、$\tau$と$\theta$によらない。(b)~スパイク$14$個。}
\answer{ex:04:7}{(a)~単語あたり$1.62$ビット：スパイク数が$0.77$、その配置が$0.84$。(b)~$1.13$ビットと$1.59$ビット。$30$個の単語からの推定はほぼ3分の1を失う。}

\answer{ex:05:1}{$\ell\approx17\um$、$\approx100$日。脳全体への配信を担うのは配線の分岐であり、拡散は各放出部位を$\sim20\um$ぼかすだけである。}
\answer{ex:05:2}{$r=ab/(1+G)$、$S=1/(1+G)$は厳密。$+20\,\%$ではなく$+1.7\,\%$。}
\answer{ex:05:3}{$T^*$は$G=2$で$1.21\,\text{s}$、$G=9$で$0.19\,\text{s}$。現実的には$G\sim2$--$4$で、抑制は$3$--$5$分の1。}
\answer{ex:05:4}{総レートを$\lambda$として$\mathrm{CV}=1/\sqrt{2\lambda\tau_c}\approx9\cdot10^{-4}$。ニューロン1個なら$\tau_c=12.5\,\text{s}$が必要。}
\answer{ex:05:5}{どちらの場合も$c\approx0.40$。$\mathrm{CV}=0.050$対$0.158$で、$\sqrt{10}$分の1。個々のニューロンのレートは$a_i$に比例したままで、調節されるのは和だけである。}
\answer{ex:05:6}{$N_1=\overline{A\vee B}$、$N_2=\overline{A\vee N_1}$、$N_3=\overline{B\vee N_1}$、$N_4=\overline{N_2\vee N_3}$、XOR $=N_5=\overline{N_4}$。各切り替えに$\tau_c\ln2$かかり、ゲート$4$段の経路でXOR 1回に$2.8\,\text{s}$。}
\answer{ex:05:7}{(a)~$\tau_\gamma=6/\ln3\approx5.5\,\text{s}$。(b)~$D<2.2\,\text{s}$に対して$\bar D<4\,\text{s}$。遅延が予測できないと忍耐強くなる。}
\answer{ex:05:8}{$k_2=1/\tau_d=10\,\text{s}^{-1}$、$k_1=1/\tau_d-1/\tau_\gamma=9.5\,\text{s}^{-1}$。$\tau_\gamma=1/(k_2-mk_1)$なので、$+2.6\,\%$で2倍になる（$+5.3\,\%$で時間幅は無限大）。}

\answer{ex:06:1}{$e=(1-e^{-T_a/\tau_e})e^{-d/\tau_e}$。(a)~$0.110$対$4.5\cdot10^{-5}$で、$\approx2.5\cdot10^3$倍。(b)~$\tau_e^*=T_a/\ln(1+T_a/d)=0.59\,\text{s}$。}
\answer{ex:06:2}{速度は毎秒$\nu_x\nu_y(A_+\tau_+-A_-\tau_-)=0.8A_+$、1時間で$\approx2900A_+$。$\tau_-=\tau_+/0.6=33\ms$。}
\answer{ex:06:3}{$\mu^2+s_1^2>s_2^2$のとき$\mathbf e_1$。ここでは$1.01>0.25$で、$\mathbf e_2$方向の散らばりが$25$倍大きいにもかかわらず、ニューロンは$\mathbf e_1$を学ぶ。規則が見つけるのは共分散行列ではなく相関行列の主固有ベクトルである。}
\answer{ex:06:4}{ランプとジュースの間で$V=Re^{-(t_c+L-t)/\tau_\gamma}$なので$\dot V=V/\tau_\gamma$。バーストの面積は$Re^{-L/\tau_\gamma}$で、$0.61R$と$0.78R$。}
\answer{ex:06:5}{SN比は$1/(N+1)$、試行数は$\propto N+1$。$5\cdot10^5$試行$\approx1$か月、$5\cdot10^{11}$試行$\approx8\cdot10^4$年。}
\answer{ex:06:6}{(a)~$k_2=1/\tau_d$、$k_1=1/\tau_d-1/\tau_\gamma$。(b)~偽の誤差は$-(V/\tau_\gamma)\,\varepsilon/(1+\varepsilon)$、$\varepsilon=\tau_d/\tau_\gamma$で、$\tau_d\le0.105\,\text{s}$。}
\answer{ex:06:7}{$0.46$、$0.90$、$0.99$、$0.95$、$d=3\,\text{s}$で$0.70$。点はトレースの割合$F=(1-e^{-T_{\text{cue}}/\tau_e})e^{-d/\tau_e}$の一本の曲線に乗る。$F\gtrsim0.1$で学習し、$F<10^{-2}$でランダムな選択になる。}
\answer{ex:06:8}{$\eta^*=1/\bigl(2\sigma^2(N+2)\bigr)$で、$N+2$試行。$N$個中$m$個がゆらぐ場合は$N(m+2)/m\ge N+2$試行で、スパース性は役に立たない。}

\answer{ex:07:1}{$V=\kappa p/\bigl(\kappa p+(1-\kappa)(1-p)\bigr)$：$0.059$、$0.20$、$0.50$。中央値は$0$だが、点~\eqref{eq:expectile}は$\kappa$とともになめらかに変わる。}
\answer{ex:07:2}{$p=1/3$、$x=0.75$、$\sigma_r=0.35$、$V(0.2)=0.083$。}
\answer{ex:07:3}{$V=\sqrt\kappa/(\sqrt\kappa+\sqrt{1-\kappa})$で、$0.25$から$0.75$まで。ばらつきは$\propto\sqrt\eta$。2滴では$V=0.25+0.5\kappa$で、両端のニューロンで分布の差は$0.05$なので、$\eta\lesssim0.01$が必要。}
\answer{ex:07:4}{(a)~$J^*=2\sqrt{C\bar R}$。(b)~払いすぎは$\bigl(k^{1/2}+k^{-1/2}\bigr)/2-1$で、$k=2$で$6\,\%$、$k=4$で$25\,\%$。「2倍以内」の精度で足りる。}
\answer{ex:07:5}{面積$R(e^{-1}-e^{-2})\approx0.23R$のバーストで、ランプ$2.3\,\text{s}$分にあたる。ドーパミンが$V$を伝えるならバーストはなく、レベルが階段状に$e$倍になるだけである。}
\answer{ex:07:6}{出来事は重みを$\propto A\tau_f/(\tau_e+\tau_f)$だけ動かし、レベルは誤差$s\tau_f$を生む。$9\,\text{s}\le\tau_f\le17\,\text{s}$。}
\answer{ex:07:7}{初期状態で$\Delta P=0.80$、$M=0.90$。(a)~$\Delta P=0.74$（$-8\,\%$）、$M=0.43$（$-52\,\%$）。(b)~$\Delta P=0.40$（$-50\,\%$）、$M=0.80$（$-11\,\%$）。二重解離である。}
\answer{ex:07:8}{$N_{\text{eff}}\to2+\rho^2N\approx10^6$試行。配線1本の場合の$10^4$分の1だが、$1\,\%$の漏れでもなお100万試行かかる。}

\answer{ex:08:1}{(a)~$\approx1.7\cdot10^6$。(b)~$g=\frac{0.9}{\sqrt{0.19}}\,\sigma_{\text{後}}/\sigma_{\text{前}}\approx16.5$。答えは雑音の比だけで決まる。}
\answer{ex:08:2}{(a)~固有値は$-(1\pm g\beta)/\tau$。差は$\tau/(1-g\beta)$で減衰する。$g\beta=0.5$で$40\ms$（入力の差は3倍に増幅）、$0.9$で$200\ms$（$19$倍）。(b)~$0.905$と$1.105$。二つの境界の間では強い入力だけが勝つ。}
\answer{ex:08:3}{$P(k)=e^{gu_k/\sigma}/\sum_le^{gu_l/\sigma}$、つまり\eqref{eq:softmax}。$g=10\ln9\approx22$。}
\answer{ex:08:4}{$\bar\Delta=1/3$、$g^*\approx3\ln(1/h)$：$21$、$14$、$9$。モデルでは$16$--$23$、$11$、$8$。}
\answer{ex:08:5}{$g^*$は$h=0.001$の$16$--$23$から$h=0.2$の$6$--$8$まで。見積もり$3\ln(1/h)$は$h\le0.05$で1.5倍以内で正しい。$h\gtrsim\eta/10$では規則$Q\leftarrow Q+\eta(R-Q)$が探索で得たものを取り込む時間がなく、$g^*$は約$8$で止まる。}
\answer{ex:08:6}{$g_{\text{lo}}=0$で$T_p=\tau\ln9=44\ms$（モデルも$44\ms$）、$g_{\text{lo}}=0.25$で$70\ms$。ゲインをほぼ0にする$2$--$3\tau$のバーストで足りる。}
\answer{ex:08:7}{(a)~最良の一定値$0.925$（$g=8$）、オラクル$0.929$（$16/8$）。得られるものはほとんどない。(b)~たとえば報酬がまれな穏やかな時期（$p\in[0,0.3]$）と$h=0.1$の変化の多い時期。$0.850$対$0.872$。調整は$\eta=0.2$で約4分の1、$\eta=0.05$でほぼすべてを取り戻す。}

\answer{ex:09:1}{(a)~$P_\infty=0.2$、記憶$20\,\text{s}$。(b)~$P_\infty=1$、記憶$1\,\text{s}$。(c)~記憶は雑音の振幅に比例するので、6倍。}
\answer{ex:09:2}{$P(t)=P_\infty\coth(t\sqrt{q/R})$。(a)~$\frac12\sqrt{R/q}=5\,\text{s}$で、記憶の半分。(b)~$t=\frac12\ln21\,\sqrt{R/q}\approx15\,\text{s}$。高い\nt{ACET}はこれだけ保たれなければならない。}
\answer{ex:09:3}{分散は$qT/2+R/(2T)$。$T^*=\sqrt{R/q}$で分散は$\sqrt{qR}=P_\infty$、フィルタ~\eqref{eq:kalman}と同じ。増加は$(\kappa+1/\kappa)/2$倍で、$1.25$と$5.05$。}
\answer{ex:09:4}{$a>a_w=1-\theta/(mw)$：$w=0.041$で$0.15$、$w=0.045$で$0.22$。}
\answer{ex:09:5}{誤りは$M\approx40$--$60$で現れ、$M=80$で余分なニューロン$1.9$個、$M=100$でパターンの割合$0.86$。他のパターンからの干渉の分散は$\approx(M-1)p(1-p)/N$なので、$M_{\max}\approx80$、$M_{\max}\propto N/p$。}
\answer{ex:09:6}{たとえば$\eta(a)=\eta\min\bigl(1,\max(0,(a-0.3)/0.6)\bigr)$。$\eta a$では無関係なニューロン3個がすべてパターンに入った。$\eta(a)$では1個も入らず、$a\ge0.9$での書き込みは同じ（割合$1.00$）。}
\answer{ex:09:7}{(a)~$40$回。問題~\ref{ex:09:6}の$\eta(a)$では決して蓄積しない。(b)~$200$回の再生で、$10$晩。}

\answer{ex:10:1}{バスは $\approx10^{12}$ bit/s（1スパイクあたり $11.4$ ビット）、制御系は $\approx40$ bit/s。$2.5\cdot10^{10}\,\text{s}$、つまり約800年。}
\answer{ex:10:2}{$a>1/3$（$a=0$ では $\approx67$~Hz の増大振動）。独立には回せない。安定性は積 $g(1-a)$ で決まる。}
\answer{ex:10:3}{平均 $\sigma^2R'x_j$、分散 $(N+1)\,x_j^2\sigma^4R'^2$。比は $1/(N+1)$ で、$n\approx z^2(N+1)\propto N$ 回の試行が要る。}
\answer{ex:10:4}{$L_v\in[32.9,\,47.9]\ms$：遅延 $35.4\ms$、窓 $\pm7.5\ms$。誤った結びつけは毎秒 $2\Delta_{\max}r_ar_v\approx0.38$ 回。}
\answer{ex:10:5}{最良の $\alpha=0.2$ での $0.180$ に対して $0.098$。$45\,\%$ 小さい。}
\answer{ex:10:6}{$n=\ln\bigl(\ln1.5/(0.6g)\bigr)/\ln(1-\alpha)$：$24$ 回と $4$ 回。$g=2$ では $11$ 回と $2$ 回。}
\answer{ex:10:7}{$S'=S\bigl(1-4\eta\sigma^2+4\eta^2\sigma^4(G+2)\bigr)$、最良は $\eta\sigma^2=1/(2(G+2))$。$\approx10^6$ 回の試行で約1か月。}

\answer{ex:11:1}{(a)~$T=k^2/\bigl(\Phi(\Delta I/I)^2\bigr)=9\,\text{s}$。(b)~センサ $45$ 個、斑点は $\approx7\times7$、分解能は $\approx7$ 倍悪くなる。}
\answer{ex:11:2}{(a)~1スパイクあたり $2$--$3.3$ ビット、上限（$9.1$--$9.8$ ビット）の $22$--$34\,\%$。(b)~$\log_2\binom{400}{20}\approx111$ ビット。共通の種類は平均 $1$ 個。}
\answer{ex:11:3}{(a)~$\sigma/9\le I\le9\sigma$、$1.9$ 桁。(b)~$6$ と $7$。}
\answer{ex:11:4}{(a)~$f<343/0.44\approx780$~Hz。(b)~スパイク $625$ 個、ニューロン $125$ 個。}
\answer{ex:11:5}{1イベントあたり $21$ ビット、毎秒 $4.8\cdot10^5$ イベント、エッジの1画素あたり $7$ イベント：毎秒 $136$ 画素。通常のカメラは毎秒 $1.25$ フレーム。}
\answer{ex:11:6}{対数では上下とも $\approx1.0\,\text{s}$（理論値 $0.96\tau$）。線形では上へ $0.2\,\text{s}$、下へ $3.0\,\text{s}$（理論値 $0.18\tau$ と $3.0\tau$）。}
\answer{ex:11:7}{(a)~$\ln I$ が中央値 $\ln\sigma$、尺度 $1$ のロジスティック分布に従うとき。広がりは自然対数で $\pi/\sqrt3$、$0.79$ 桁。(b)~$r=r_{\max}\Phi_I(I)^2$。}

\answer{ex:12:1}{$0.60$（相関がなければ $0.18$）、極限は $\sqrt{c/(rT)}\approx0.58$。100個のニューロンで区別できるのは「ある／ない」だけ。}
\answer{ex:12:2}{(a)~$10^8$ スパイク、$\approx10\,\text{mJ}$：脳全体の $3$~J の $0.3\,\%$。(b)~$3$~J、$\approx300$ 倍。}
\answer{ex:12:3}{(a)~$\theta_A\in[2,3)$, $\theta_B\in[1,2)$。他方の図形は高々 $2$ と $1$ しか与えない。(b)~$10\cdot96^2\approx9.2\cdot10^4$。}
\answer{ex:12:4}{(a)~$T_d=0.85\,\tau_a=0.34\,\text{s}$。(b)~$T_d\propto\tau_a$ で $I$ によらない：$\tau_a\approx3.5\,\text{s}$（シミュレーションでは $3.1\,\text{s}$）。}
\answer{ex:12:5}{$p\approx0.17$。正答率はゆっくり下がる：$N=8$, $12$, $24$, $48$, $96$, $192$ で $0.90$, $0.88$, $0.86$, $0.84$, $0.82$, $0.79$。}
\answer{ex:12:6}{$\tau_{\text{tr}}=20$, $50$, $200\ms$ で10回すべて誤りなし（$30\ms$ では10回中1回失敗）。$5$--$10\ms$ では約 $0.5$、$0.5$--$2\,\text{s}$ では $0.86$ まで下がる。上限：トレースは休止の間に消えなければならない、$e^{-0.3\,\text{s}/\tau_{\text{tr}}}\ll1$。}
\answer{ex:12:7}{遅延一つでは足りない。窓 $\pm5\ms$ は $5$ から $20\ms$ までの $\Delta x/v$ をカバーしない。$5<d_1\le10\ms$ と $15\le d_2\le d_1+10\ms$ の二つの抑制の枝が要る。}
\answer{ex:12:8}{\texttt{two\_stage}：$\Delta=2$。1点の反転で同点、2点で答えが変わる。1の個数を数える分類器：1点の反転。$p=0.1$ で $0.57$。}

\answer{ex:13:1}{(a)~$q=100^{1/99}\approx1.048$、刻み $\approx4.8\%$。レンジ $\approx2.2\cdot10^3$（$67$~dB）。(b)~刻みは $22f_1$、つまり $10f_1$ で $220\%$。}
\answer{ex:13:2}{$P_{\text{fail}}\approx1.8\cdot10^{-4}$ と $2.8\cdot10^{-16}$：$S=2$ で1日約 $300$ 回、$S=4$ で1日 $5\cdot10^{-10}$ 回。}
\answer{ex:13:3}{(a)~$H(s)=eF_1T_c/(1+sT_c)^2$、遮断周波数 $1/(2\pi T_c)\approx3.2$~Hz の2次フィルタ。(b)~$r\approx22$~スパイク/s（基本波だけなら $\approx20$）。}
\answer{ex:13:4}{(a)~$Gd=1/e$ で根 $s=-1/d$ は重根となり、ほかのどの $G$ でも最も右の根はそれより左に来ない。(b)~$d=30\ms$：$G\approx12$~s$^{-1}$、$\tau=30\ms$。$d=150\ms$：$G\approx2.5$~s$^{-1}$、$\tau=150\ms$。時定数は遅延に等しい。}
\answer{ex:13:5}{(a)~$E\approx0.427$、$T\approx1.70\,\tau_a=0.68$~s（$\tau=20\ms$ のモデルでは $0.84$~s）。(b)~方程式は $I$ を含まない。リズムは $b>\beta-1$ で存在する。}
\answer{ex:13:6}{$b=0.8$ ではリズムなし。$b=1.05$, $1.2$, $1.5$, $2$, $4$ で $T\approx2.73$, $1.74$, $1.19$, $0.84$, $0.45$~s。どの $I$ でも $T=0.840$~s：入力は振れ幅だけを変え、テンポは変えない。}
\answer{ex:13:7}{置換 $s=u-a$ により、最大の減衰速度は $G=e^{-1-ad}/d$ のとき $a+1/d$。$a\ge33.3$~s$^{-1}$、$G=e^{-2}/d\approx4.5$~s$^{-1}$。同時収縮が強い（$a$ が大きい）ほど $G$ を小さくとる必要がある。}
\answer{ex:13:8}{自由回答の問題。例えば疲労の駆動にしきい値を入れた $b[r_i-r_0]_+$ は、$b=4$, $r_0=0.2$ で $I_{\min}=\beta b r_0/(b-\beta+1)\approx0.53$ を与え、周期は $I=0.6$ で $1.8$~s、$I=3$ で $0.52$~s となる。}

\answer{ex:14:1}{(a)~$0.17$~s と $1.5$~s。(b)~$11$~cm 未満の距離から。}
\answer{ex:14:2}{$g\le\beta w_{q\mu}/w_{z\mu}=5$ である限り、痛みのない触覚はどんな $\mu$ でも通らない。触覚 $\mu=1$ は $g>6$ で通る。強い感作はふつうの接触を痛みに変える。}
\answer{ex:14:3}{(a)~$g^*=(1-k_sC)/(1-k_sA)$、$k_sA<1$ で安定。(b)~触覚なしでは $\nu\ge2$、触覚ありではすでに $\nu\ge1.7$ で。}
\answer{ex:14:4}{$V(t)=-Pe^{-(T-t)/\tau_\gamma}$。(a)~$-Pe^{-T/\tau_\gamma}\approx-0.37P$。(b)~$+Pe^{-(T-t_e)/\tau_\gamma}\approx+0.61P$。これは $t_e$ とともに $+P$ まで増え、痛みを避けることを教える。}
\answer{ex:14:5}{$k_s=4$ ではラッチしない（第2の安定状態は $k_s\ge4.58$ で現れる）。$k_s=4.6$, $5$, $6$, $8$ で $T_{\min}\approx4.35$, $1.40$, $0.65$, $0.32\,\tau_s$。}
\answer{ex:14:6}{$w_{q\nu}=4/9\approx0.444$, $q_0=2/9\approx0.222$, $w_{q\mu}=49/45\approx1.089$。痛みのない触覚は $g=49/9\approx5.4$ まで安全。}
\answer{ex:14:7}{(a)~しきい値は $\nu_c\ge q_0/w_{q\nu}$ のとき $\nu_c(g)=\theta/g$、そうでなければ $(\theta+\beta q_0)/(g+\beta w_{q\nu})$。$g^*\approx2.45$, $1.0$, $0.25$。(b)~未解決の問題。}

\answer{ex:15:1}{(a)~ニューロンあたり $2\cdot10^5$、回路全体で $6\cdot10^{12}$。(b)~配線1本あたり $\approx8$~bit/s、少なくとも $2.5\cdot10^4$~s $\approx7$ 時間。}
\answer{ex:15:2}{$2$ 倍刻み：隣どうしの相関 $0.943$、$\lambda$ は $4.18$ から $7.3\cdot10^{-4}$ まで、比は $\approx5.7\cdot10^3$。$4$ 倍刻み：相関 $0.8$、比は $\approx94$。}
\answer{ex:15:3}{(a)~$\lambda=0.988$ と $0.808$：$82$ 回と $4.7$ 回の動き。(b)~$x_s^*=0.690$, $x_f^*=0.172$、和 $0.862$。}
\answer{ex:15:4}{(a)~$116$ 回。(b)~$19$ 回。一つの過程のモデルでは、和がゼロなら節約は生じない。}
\answer{ex:15:5}{(a)~$G=50$~s$^{-1}$。モデルなしの限界 $10.5$~s$^{-1}$ と比べよ。(b)~安定ではない。$G=50$~s$^{-1}$ で臨界遅延は $d_c\approx103\ms$。$d=150\ms$ では $\hat k>0.445k$ が必要（どんな $G$, $d$ でも $\hat k>k/2$）。}
\answer{ex:15:6}{誤差は $6$--$30$~s で $0.5\%$ 未満に下がり、床は $\approx(6$--$9)\cdot10^{-4}$（最良の重みでは $7.5\cdot10^{-4}$）。$\eta=50$ の重みは、$C$ の条件の悪い方向（問題~\ref{ex:15:2}）に沿って最良値からまだ最大 $0.2$ ずれている。}
\answer{ex:15:7}{$\mathbf B=A\mathbf K$。$q_f/R\approx0.035$, $q_s/R\approx8.6\cdot10^{-4}$。$4q_s$ では $B_f\approx0.088$, $B_s\approx0.047$：遅い過程は2倍以上速く学び、速い過程はより遅くなる。}

\answer{ex:16:1}{血液は$\tau=t_{1/2}/\ln2=14.4$~minのRC回路である。最大と最小の比は$2^{T/t_{1/2}}=64$、基本波は$0.55$に減衰する。比が$2$になるのは$t_{1/2}=T=60$~minのとき。}
\answer{ex:16:2}{$\varphi^*=\arcsin\bigl((\omega-\omega_0)/K_\varphi\bigr)\approx41$~min（$\omega-\omega_0\approx0.047$~rad/day）、緩和時間は$1/(K_\varphi\cos\varphi^*)\approx3.9$~day。$+6$~hと$-6$~hのずれのあとは、それぞれ$7.7$~dayと$7.9$~day。}
\answer{ex:16:3}{$K_c(n)=\sec^n(\pi/n)$：$8$と$2.37$、誤差は$11\%$と$30\%$。$n\to\infty$で$K_c\to1$、誤差$\to50\%$。}
\answer{ex:16:4}{両極端の指令（アクセル$80$、ブレーキ$0$）で$0.76$~s後に$f=120$となる。「アクセルだけ」では$2.33$~sで、3倍長い。}
\answer{ex:16:5}{$3\arctan(\omega\tau)+\omega d=\pi$、$K_c=(1+\omega^2\tau^2)^{3/2}$：$K_c=8$、$3.54$、$2.50$、$1.76$、周期は$3.63$、$5.46$、$6.86$、$9.27\,\tau$。$0.9K_c$では振動が減衰し、$1.1K_c$では持続振動に落ち着く。}
\answer{ex:16:6}{$0.60\bar c$、$0.87\bar c$、$0.90\bar c$、極限は$0.91\bar c$。濃度が一定なら応答はゼロである。このような受け手は分量しか読まない。}
\answer{ex:16:7}{自由課題。フィードバックならリズムは$K>8$でしか現れず、周期は段の$\tau$に比例し、$K$にはほとんど依存しない（$K=8.5$で$3.64\,\tau$、$K=40$で$4.23\,\tau$）。$K$を1.5倍変えても周期は$2$--$4\%$しか動かず、誤差より小さい。決め手はループを開くこと（第1段の入力に最終段のホルモンを一定量与えると、この種のリズムは消える）か、周期の異なる位相で短時間ホルモンを追加したときの応答である。}

\answer{ex:17:1}{$0.2$~Wに対して$81$~\textmu W。生データなら$50$~pJ/bit以下が必要。}
\answer{ex:17:2}{$c=0.1$では$5.6$、$8.7$、$9.2$~bit/s、極限$9.3$~bit/s。$c=0$、$N=1000$では$65$~bit/s。}
\answer{ex:17:3}{$B=\log_2N+P\log_2P+(1-P)\log_2\frac{1-P}{N-1}=4.87$、$4.37$、$3.29$~bit/文字、すなわち$7.3$、$6.6$、$4.9$~bit/s。$10$~bit/sには毎秒$2.3$文字が必要。}
\answer{ex:17:4}{成分あたりの誤差分散は$2b/(Nm^2T)+\sigma_\xi^2$。両者の寄与が等しいのは$N\approx90$。極限は窓あたり成分あたり$\tfrac12\log_2(1+0.25/0.04)\approx1.4$~bit。}
\answer{ex:17:5}{組は$\eta_bd^2+\eta_db^2<2$、つまりおよそ$\eta<1$で収束する。$0.8$では収束、$1.1$では周期$2$の持続振動、$1.5$では非周期的、$2$では発散。学習者の速度は分けるべきである。}
\answer{ex:17:6}{閾値は$\approx4.4\sigma$。$102$~kbit/s、$0.10$~mWで、生データの$205$~mWの$2000$分の1。飽和する（$3$スパイクを超える）のは$5.7\cdot10^{-5}$のサンプルだけ。}
\answer{ex:17:7}{自由課題。$\text{SNR}\approx0.9$で約$46$~bit/s、独立ノイズで約$330$~bit/s。信号に似たノイズが原因なら、取り出せる情報はチャネル数を増やすと飽和する。符号の無知が原因なら、より良いデコーダ（たとえばスパイクのタイミングを使うもの）で増える。}

\answer{ex:18:1}{(a)~$0.9949\le w\le1.0049$、すなわち$|1-w|\lesssim0.005$。(b)~$K\ge400$個のシナプス。}
\answer{ex:18:2}{$(\omega_R\tau_w)^2=\sqrt{\alpha(\alpha+2+2\varrho)}/\varrho-1$、$\varrho=\tau_m/\tau_w$。$f_R\approx11.1$~Hz、$|Z(\omega_R)|/|Z(0)|=4.6$。}
\answer{ex:18:3}{(a)~$f=1/\bigl(\tau\ln\frac{\mu}{\mu-\theta}\bigr)$。$1$~Hzには$\mu/\theta-1\approx3.7\cdot10^{-44}$が必要。(b)~$T=\pi\tau/\sqrt\mu$、$f\approx3.2$~Hz。$f\propto\sqrt\mu$である。}
\answer{ex:18:4}{積分器型は$f\lesssim17.7$~Hzで発火し（ローパスフィルタ）、共振器型は$5.5$--$22$~Hzの帯域でだけ発火する（バンドパスフィルタ）。}
\answer{ex:18:5}{(a)~$T_{\min}=\frac{\tau}{w-1}\ln\frac{0.5}{0.3}\approx10.2\ms$。(b)~$B>w-1-0.2=0.3$。}
\answer{ex:18:6}{調整の時定数は$\tau/(\eta\langle r^2\rangle)$、$\eta=\tau\ln10/60\,\text{s}\approx3.8\cdot10^{-3}$。$I\neq0$では重みがずれるので、$e$から指令のコピー$I/\tau$を引く必要がある。}
\answer{ex:18:7}{自由課題。共振器型の有利さは$11$~Hzでわずか$1.35$倍で、$1$~Hzでは$17$分の1と不利になる。組み替えの主な利点は、閾値による帯域の選択にある（問題~\ref{ex:18:4}）。}

\answer{ex:19:1}{(a)~シナプス1個で$6.7\mV$（$R=66.7$~M$\Omega$）なので、そもそも閾値に達するには$4$個以上が必要（3個では漸近的にしか届かない）。$10\ms$では$8$個、$5\ms$では$14$個。(b)~$0.1\ms\ll\tau_m$、漏れによる補正は$\sim0.25\%$。}
\answer{ex:19:2}{(a)~$1.75\cdot10^{10}$（ミキサの4分の1が何にでも応答する）、$4.4\cdot10^9$、$0.06$（$k=40$ではほとんど応答しない）。(b)~$10^3$倍：$2\cdot10^5$対$200$。}
\answer{ex:19:3}{(a)~二項係数の対称性により$C(2n,n)=2^{2n-1}$。(b)~$0.93$と$0.09$。$P$についての遷移幅は$\sim\sqrt{2n}$、相対幅は$\sim1/\sqrt n$。}
\answer{ex:19:4}{1本の樹状突起からでは、入力がいくつあっても$V_{\text{soma}}<\kappa E_E<\theta$。$g_0=0.5g_L$では各樹状突起に$n\ge3$個が必要。}
\answer{ex:19:5}{$I\ge C\sigma_V/\sigma_t=15$~nA、すなわち$150$個の同期シナプスが必要で、非現実的である。小さなニューロンなら$1$~nAで足り、$4$~nAのシナプスではジッタは$5$~\textmu s。}
\answer{ex:19:6}{近端での最大：$0.03\,\tau_m$で$0.97$（$L=0.2$）、$0.32\,\tau_m$で$0.66$（$L=1$）、$0.79\,\tau_m$で$0.33$（$L=2$）。全容量による見積もり~\eqref{eq:dendriteload}は$L\ll1$でしか正しくない。}
\answer{ex:19:7}{自由課題。$m$を固定すると、応答時間は記憶する$P$に比例して増える。割合$m=\varphi n$を固定すると$n$に依存しなくなる。大きなニューロンの遅さは、大きさそのものではなく、多くの入力を足し合わせることの帰結である。}

\answer{ex:20:1}{$t_{\text{w}}=\tau_r\ln\frac{1-L}{1-H}=16.8$~h、$t_{\text{s}}=\tau_d\ln\frac HL=5.8$~h、周期$22.5$~h。発振器を$24$~hに引き込むのは閾値の概日的な揺れ、すなわち位相同期~\eqref{eq:pll}である。}
\answer{ex:20:2}{(a)~$L=0.111$。$L=0.092$なら睡眠は$9.6$~h続く。(b)~$18\%$ではなく$22\%$（$1/0.82=1.22$）。}
\answer{ex:20:3}{$H=\frac{p-1}{p-1/q}=0.623$、$L=H/q=0.093$、ただし$p=e^{16/\tau_r}$、$q=e^{8/\tau_d}$。$4$~hで起こされると位相は永久にずれる（入眠が$7.2$~h早まる）。閾値が揺れるなら位相は2--3日で戻る。}
\answer{ex:20:4}{(a)~$r_\pm=\frac12\bigl(1\pm\sqrt{1-4k/w}\bigr)=0.947,\ 0.053$。$a_{\text{up}}=3.51$、$a_{\text{down}}=0.49$。(b)~活動$\approx0.33$~s、静止$\approx0.59$~s、周期$\approx0.93$~s、活動の割合$0.36$。}
\answer{ex:20:5}{$a_{\text{up}}/r_+<b<a_{\text{down}}/r_-$：$3.71<b<9.20$。\nt{ACET}は$b$を$3.71$未満に下げる。交点は上の枝に移り、活動状態（$r\approx1$）が安定になる。}
\answer{ex:20:6}{$D=24$、$32$、$40$、$48$、$64$で$4.9$、$6.8$、$8.8$、$5.5$、$8.9$~h。長さを決めるのは借りではなく、入眠の概日位相である。}
\answer{ex:20:7}{Bのあと、Aでの誤差は平均$0.51$（$20$セットで$0.19$から$1.08$。応答ゼロなら$1$）。混ぜて学習すると、どちらの誤差も$10^{-4}$未満。}
\answer{ex:20:8}{自由課題。一様なスケーリングは重みの比を保つが、閾値ニューロンの応答を保つのは閾値も同じ比率でスケーリングされる場合に限る。インターリーブしたリプレイは比を変える。大きな接触が変わらないことは、純粋に一様なスケーリングと矛盾する。}

\answer{ex:21:1}{(a)~$T=\tau_{\text{rec}}\ln\frac{1}{1-x^*}$：$1.84$~s と $3.68$~s。(b)~$3.36$ から $4.24$~s。感度は $\dd T/\dd x^*=\tau_{\text{rec}}/(1-x^*)=80$~s で、$x^*\to1$ では一斉発火は不規則になる。}
\answer{ex:21:2}{(a)~$8.6$~W、\eqref{eq:energy}と同じ。(b)~$205$~Mbit/s、$0.2$~W。培養（$10^{-4}$~W）の $2\cdot10^3$ 倍。}
\answer{ex:21:3}{$x_{\text{ss}}=1/(1+Ur_b\tau_{\text{rec}})<x^*$ となるのは、$x^*=0.9$ では $r_b>(1/x^*-1)/(U\tau_{\text{rec}})=0.28$~Hz、$x^*=0.99$ では $0.025$~Hz のとき。シナプスの強さは $x_{\text{ss}}$ まで、つまり $10\%$ と $1\%$ 下がる。}
\answer{ex:21:4}{$u(t-k)$ は正規直交で、$m_k=\|P_Vu(t-k)\|^2$。ここで $P_V$ は出力が張る $N$ 次元部分空間への射影である。ベッセルの不等式から $\sum_k m_k\le N$。}
\answer{ex:21:5}{$\tanh$ ありで $\mathrm{MC}\approx9$--$11$、線形リザバーで $15$--$18$（シード3個）。どちらも $\le30$。非線形性の代償は記憶である。}
\answer{ex:21:6}{(a)~$n\le102$。(b)~$5.6$ シグマ。}
\answer{ex:21:7}{未解決の課題。鍵となる対照：読み出しを凍結し、刺激なしの休止後に改善が残るかを確かめる。残らなければ、変わったのは重みではなく状態である。}

\answer{ex:22:1}{(a)~$32\%$ と $100\%$。(b)~$R_\uparrow-R_\downarrow\ge2\sqrt{2000\,\text{Hz}/1\,\text{s}}\approx89$~Hz。合計頻度のわずか $4.5\%$ である。}
\answer{ex:22:2}{(a)~$n\ge4(p_1q_1+p_2q_2)/\Delta^2\approx560$。(b)~$\approx56$。}
\answer{ex:22:3}{(a)~詳細釣り合い：$\rho PK$ は対称である。(b)~ミスの割合は調和平均 $1/\langle1/P\rangle=0.18$ で、フィードバックなしの $0.5$ に比べて小さい。}
\answer{ex:22:4}{(a)~面積 $4h^2=0.04$ の平行四辺形（$p$ の制限を考慮すると数値的に $0.045$）。(b)~$\approx0.18$。}
\answer{ex:22:5}{$0.49$：一部の培養は、ほぼ必ずミスする領域に入り込み、そこをさまよう。制限をつけると $1.00$：有界な集合では、密度 $\propto1/P$ は吸収領域に集まる。}
\answer{ex:22:6}{(a)~$\ge5$ 段階が必要で、$8$ 本の電極で足りる。(b)~不可能：$r_{\max}\ge25/T=100$~Hz が必要になる。}
\answer{ex:22:7}{未解決の課題。ランダム探索の時間は $d$ とともにおよそ体積比 $(L/h)^d$ のように増え、誤差の符号を使う探索では $d$ に比例して増える。仮説を区別するのは入れ替え実験である。ミスの後に予測可能だが強い刺激、ヒットの後に予測不能だが弱い刺激を与える。各群に数十個の培養が必要である。}


\chapter{実験の補遺}
\label{sec:supplement}

ここでは章ごとに、主な主張とその根拠をまとめる。根拠とは、それを測定した実験、何を測定したか、どこで発表されたかである。右の列の状態は、主張がどれだけ確かかを示す。\emph{測定済み}は直接の実験があること、\emph{理論}は本文や古典的な研究で導かれた数学的帰結、\emph{モデル}は本書のモデルによる計算結果（スクリプトは \texttt{models/} フォルダにある）、\emph{推定}は直接の測定のない桁の見積もり、\emph{仮説}はもっともらしいがまだ実験で証明されていない説明である。実験がない場合は、そう明記した。

\section{第\ref{sec:neurontypes}章　ニューロンの種類}

この章の数値は、ヒトの脳で細胞の直接計数があるところではそれに基づく。「ニューロン1個に神経伝達物質1種類」という規則は、細胞アトラスと、例外も見つけた実験に基づく。\nt{DOPA}の役割はスパイクの記録に基づき、その他のブロードキャストチャネルの役割は仮説に基づく。

{\footnotesize
\setlength{\tabcolsep}{3pt}\renewcommand{\arraystretch}{1.15}
\begin{longtable}{>{\raggedright}p{0.5cm}>{\raggedright}p{4.0cm}>{\raggedright}p{5.6cm}>{\raggedright}p{2.3cm}>{\raggedright\arraybackslash}p{1.7cm}}
\toprule
No. & 主張 & 実験と測定されたもの & 出典 & 状態 \\
\midrule\endfirsthead
\toprule
No. & 主張 & 実験と測定されたもの & 出典 & 状態 \\
\midrule\endhead
1 & ヒトの脳のニューロンは約$8.6\cdot10^{10}$個で、\nt{GLUT}ニューロンのほとんどは小脳の小さなニューロンである（表~\ref{tab:types}） & 等方性分画法：組織を溶かして核の均一な懸濁液にし、ニューロンマーカーNeuNをもつ核を数える。成人男性のニューロンは$(86.1\pm8.1)\cdot10^9$個。そのうち大脳皮質にあるのは$19\,\%$にすぎず、大部分は小脳にある & \cite{Azevedo2009} & 測定済み \\
2 & ほとんどすべてのニューロンは主要な神経伝達物質を1種類もつ（命題~\ref{prop:dale}）が、例外もある & マウス脳のトランスクリプトーム・アトラス：約$4\cdot10^6$個の細胞、$5322$クラスタ。合成と小胞充填の遺伝子から各クラスタに神経伝達物質を割り当てた。例外は直接測定されている。スライス標本の\nt{DOPA}ニューロンは同じ小胞トランスポーターを通じてGABAも放出し、線条体のニューロンを抑制する。\nt{DOPA}軸索の一部では、グルタミン酸とドーパミンが同じ配線の異なる終末から出る（電子顕微鏡と光遺伝学） & \cite{Strata1999,Yao2023,Tritsch2012,Zhang2015} & 測定済み \\
3 & 重み行列の列全体が同じ符号をもつ~\eqref{eq:dalesign} & 命題~\ref{prop:dale}と、グルタミン酸受容体はほぼすべて興奮性でGABA受容体は抑制性であることから章内で導出。「1個のニューロンのすべての受け手が同じ符号の重みを受け取る」ことを一般規則として直接検証した例はない。反例は\nt{DOPA}ニューロンによるGABAの共放出（2行目）と、異なる符号の受容体をもつ受け手（11行目） & 章内の導出 & 理論 \\
4 & 大脳皮質のニューロンの4分の3から5分の4が\nt{GLUT}、5分の1から4分の1が\nt{GABA} & サル大脳皮質の$10$領域で、皮質の全層を貫く幅$50$\,\textmu mの柱状領域をGABAで免疫染色。$8$領域ではGABAニューロンが全ニューロンの約$25\,\%$、一次視覚野では$20\,\%$未満 & \cite{Hendry1987} & 測定済み \\
5 & \nt{GLUT}ニューロンの出力は約$10^4$個 & 剖検脳の立体解析学：若年男性の新皮質にシナプス$1.64\cdot10^{14}$個（電子顕微鏡、ダイセクター法）、ニューロン約$2.3\cdot10^{10}$個。比をとるとニューロン1個あたり$\sim7\cdot10^3$シナプス。平均すれば入力数は出力数に等しい & \cite{Tang2001synapses,Pakkenberg1997} & 測定済み \\
6 & \nt{DOPA}ニューロンの出力は$10^5$--$10^6$個の終末に枝分かれする。これは標的の被覆範囲からの推定であって計数ではない & ラットの個々の\nt{DOPA}ニューロンの膜をウイルスで標識し、軸索を完全に再構成：1個のニューロンが線条体の体積の$0.45$から$5.7\,\%$（平均$2.7\,\%$）を覆う。測定したのは被覆範囲で、終末数ではない。終末数は被覆範囲からオーダーとして導いたもので、終末の直接計数はない。\nt{SERO}と\nt{NORA}の終末数は粗い推定 & \cite{Matsuda2009} & 測定済み \\
7 & ブロードキャスト型ニューロンは少ない：\nt{DOPA} $\sim5\cdot10^5$（二つの群のうち主要なほう、下限）、\nt{SERO} $\sim3\cdot10^5$（二つの部分計数の和）、\nt{HIST} $\sim6\cdot10^4$（表~\ref{tab:types}、図~\ref{fig:typepie}） & ヒトでの立体解析学的計数：黒質の色素含有ニューロン$550\,000$個（腹側被蓋野を除く）。背側縫線核のニューロン$235\,000$個（核の全ニューロン）と、橋被蓋のセロトニンニューロンさらに$125\,000$個。視床下部のヒスタミンニューロン約$64\,000$個。\nt{NORA}、\nt{GLYC}、\nt{ACET}、\nt{ADRE}には直接の計数がなく、表の値は粗いオーダー推定である & \cite{Pakkenberg1991,Baker1990,Baker1991,Airaksinen1991} & 測定済み \\
8 & 隙間のグルタミン酸は約1ミリモーラーまで跳ね上がり、$\sim1\ms$で減衰する & シナプス伝達中に、速く解離する競合的阻害薬がNMDA受容体から追い出される程度から求めた（海馬ニューロンの培養）：ピーク$1.1$\,mM、減衰時定数$1.2\ms$ & \cite{Clements1992} & 測定済み \\
9 & 動作中の皮質では興奮性電流と抑制性電流がほぼ釣り合い、ニューロンは閾値付近にいる & 理論：強い結合をもつネットワークは自ら釣り合いの状態に入る。実験：フェレット前頭前皮質のin vivo記録で、「アップ」状態の間、コンダクタンスが$\sim21$\,nS変化してもシナプス電流の反転電位は数百ミリ秒にわたり約$-37$\,mVに保たれる。興奮と抑制が一緒に増減する & \cite{vanVreeswijk1996,Haider2006} & 測定済み \\
10 & \nt{DOPA}ニューロンは報酬予測誤差~\eqref{eq:rpe}を伝える（図~\ref{fig:rpe}） & サルの\nt{DOPA}ニューロンのスパイク：予告なしのジュースにバースト。学習後は合図にバーストし、予測されたジュースには応答しない。約束されたジュースが与えられないと発火率が落ち込む。定量的な実験では、発火率は現在の報酬と過去の報酬の加重平均との差に比例するが、それは正の誤差に限られる。式~\eqref{eq:rpe}そのものは時間差分アルゴリズムである & \cite{Schultz1997,Bayer2005,SuttonBarto2018} & 測定済み \\
11 & 符号と速さを決めるのは受容体：イオンチャネル型は1ミリ秒未満、代謝型ははるかに遅い（命題~\ref{prop:receptor}） & 海馬ニューロンの膜パッチに短いグルタミン酸パルスを与えると、イオンチャネル型受容体の電流は$0.2$--$0.6\ms$で立ち上がり（$20$から$80\,\%$）、$2$--$3\ms$で減衰する。錐体ニューロンの遅い抑制性電位は、GABAの代謝型受容体の選択的阻害薬で消える。ドーパミン受容体には逆の作用をもつ二つのファミリーがある。線条体では、D1受容体をもつニューロンはシナプスの増強にドーパミンを必要とし、D2受容体をもつニューロンは減弱にドーパミンを必要とする & \cite{Colquhoun1992,DutarNicoll1988,Beaulieu2011,Shen2008} & 測定済み \\
12 & ブロードキャストチャネルは1本の配線ではなく、少数の線の束である（\ref{sec:buses}節） & マウス背側縫線核のセロトニンニューロンをウイルスと遺伝学的手法で標識：扁桃体に出力するニューロンと前頭皮質に出力するニューロンは核内の異なる場所にあり、互いに補い合う領域に枝分かれし、不快な刺激に逆向きに応答する。総説：青斑核（\nt{NORA}）のニューロンの亜群は異なる過程を符号化する & \cite{Ren2018,Poe2020} & 測定済み \\
13 & \nt{NORA}がゲイン$g$を決める & 適応的ゲインの仮説。カテコールアミンが伝達特性の傾きを上げるネットワークモデル。ネットワーク全体で「ノルアドレナリンとともに$g$が上がる」ことの直接測定はない。測定されているのは、サルの瞳孔径が青斑核ニューロンのスパイクに追従すること & \cite{AstonJones2005,ServanSchreiber1990,Joshi2016} & 仮説 \\
14 & \nt{ACET}が「書き込み/読み出し」を切り替え、学習率を決める & 仮説。根拠：ラット嗅皮質のスライスで、アセチルコリン作動薬（$100$\,\textmu M）は内部の「想起」シナプスを$60\,\%$以上弱めるが、入力シナプスは$15\,\%$未満しか弱めない & \cite{Hasselmo2006,HasselmoBower1992} & 仮説 \\
15 & \nt{SERO}が割引率$\gamma$を決める。セロトニンニューロンは一様に毎秒数スパイクで発火し、睡眠のある相ではほぼ沈黙する & 「セロトニン$=\gamma$」仮説。最も強い論拠：マウス背側縫線核のセロトニンニューロンを光遺伝学的に活性化すると、遅延報酬を待つのを途中でやめる回数が減る。発火率と急速眼球運動睡眠中の沈黙は測定されている（記録の総説） & \cite{Doya2002,Miyazaki2014,JacobsAzmitia1992} & 仮説 \\
\bottomrule
\end{longtable}}

\section{第\ref{sec:glutcomputer}章　\nt{GLUT}ニューロンは何を計算するか}

この章の論理的な主張は、非負の重みからなる回路についての定理であり、章の中で導いたものである。実験が裏づけるのは、ニューロンのモデルと、これらの定理に従って組んだ回路（同時検出器、遅延線、自己維持する群、連鎖）が脳に実際に存在することである。

{\footnotesize
\setlength{\tabcolsep}{3pt}\renewcommand{\arraystretch}{1.15}
\begin{longtable}{>{\raggedright}p{0.5cm}>{\raggedright}p{4.0cm}>{\raggedright}p{5.6cm}>{\raggedright}p{2.3cm}>{\raggedright\arraybackslash}p{1.7cm}}
\toprule
No. & 主張 & 実験と測定されたもの & 出典 & 状態 \\
\midrule\endfirsthead
\toprule
No. & 主張 & 実験と測定されたもの & 出典 & 状態 \\
\midrule\endhead
1 & ニューロンは閾値とリセットをもつRC回路である~\eqref{eq:glutneuron} & ラット皮質の錐体ニューロン（スライス）の細胞体に、生きた脳のシナプス入力の流れに似た雑音電流を注入した。平均発火率の電流の平均とばらつきへの依存性は、遅い発火率適応を加えれば漏れ積分発火ニューロンで記述できる。$\tau$と遅延の範囲は実験を引用せずに章で示した & \cite{Rauch2003} & 測定済み \\
2 & モデル~\eqref{eq:glutneuron}は単純化であり、入力の枝そのものが局所スパイクを出す & マウス視覚野の錐体ニューロンの樹状突起からin vivoで直接記録：視覚刺激がNMDA受容体に依存する局所的な樹状突起スパイクを起こし、それを抑えるとニューロンの方位選択性が広がる & \cite{Smith2013} & 測定済み \\
3 & 同時性の窓$\Delta_{\max}=\tau\ln\frac{w}{\theta-w}$~\eqref{eq:window}。$\theta=1.5w$では$0.7\tau$ & モデル~\eqref{eq:glutneuron}からの帰結。狭い加算窓の直接測定は速い抑制をもつニューロンについてある（第\ref{sec:gaba}章、その表の3行目） & 章内の導出 & 理論 \\
4 & \nt{GLUT}ニューロンだけのネットワークは入力の単調関数しか計算しない（命題~\ref{prop:monotone}） & 章内の証明：非減少の右辺による静寂からの反復。これはモデルについての主張であり、抑制のない生きたネットワークで検証した実験はない & 章内の導出 & 理論 \\
5 & 感覚器は配線ペアを与える。刺激の増加に応答する細胞と減少に応答する細胞がある & 網膜：双極細胞の段階ですでに錐体の信号はオンチャネルとオフチャネルに分かれる。サルでオン双極細胞を薬理学的に遮断：両チャネルは皮質まで別々に進み、遮断後は光の増加の検出が悪化するが、減少の検出は悪化しない & \cite{Schiller1992} & 測定済み \\
6 & 2線式符号があれば、\nt{GLUT}ニューロンのネットワークは任意の論理関数を共通のクロックに頼らず非同期に計算でき、代償はニューロン数が2倍になること（命題~\ref{prop:dualrail}、~\eqref{eq:dualrail}、図~\ref{fig:dualxor}） & 2線式符号と信号自体による完了検出は、遅延に依存しない非同期回路の標準的な技法である。ニューロン6個の排他的OR回路は章の中で組んだ。脳がまさに2線式符号で論理関数を計算していることを示す実験はない & \cite{Sparso2001} & 理論 \\
7 & ヒトで両耳への音の到着時刻の差は最大$\approx0.6\ms$であり、脳は約10マイクロ秒を区別する & 二肢選択実験の聴取者：時間差の弁別閾（正答率75\,\%）は、帯域$150$--$1700$\,Hzの雑音で$9$\,\textmu s、$1$\,kHzの純音で$11$\,\textmu s、クリック音で$28$\,\textmu s。上限$0.6\ms$は章の計算、$0.22\,\text{m}/343\,\text{m/s}$ & \cite{KlumppEady1956} & 測定済み \\
8 & 鳥では遅延線と同時検出器によって時間差が場所に変わる~\eqref{eq:itd}（図~\ref{fig:itd}） & メンフクロウ：両耳からの軸索は同時検出器の核に反対側から入る。軸索に沿ったスパイクの到着時刻は深さとともに規則的に変わり、核を横切る遅延の幅（$700$\,\textmu mで$160$\,\textmu s）は両耳間遅延の範囲と一致する & \cite{Carr1990} & 測定済み \\
9 & 哺乳類では遅延の列は見つかっていない。同じ課題を、\nt{GLYC}ニューロンからの時間的に精密な抑制をもつ同時検出器が解く & スナネズミの内側上オリーブでのin vivo記録：応答は方向の地図をなさない。マイクロピペットでグリシンとその阻害薬を投与すると、時間的に精密な\emph{グリシン}抑制が遅延の動作範囲を決めていることがわかる。ここでの抑制は\nt{GABA}ではなく\nt{GLYC}による & \cite{Brand2002,Grothe2010} & 測定済み \\
10 & 強い帰還をもつ群はフリップフロップであり、自己維持活動は刺激の後も数秒続く（図~\ref{fig:bistable}） & 記憶した点への遅延眼球運動課題（遅延$1$--$6$\,s）中のサル前頭前皮質：課題に関係する$170$個のニューロンのうち$87$個が遅延中に発火率を変え、そのうち$79\,\%$では記憶した点の方向に依存する。この活動を二つの安定状態をもつ帰還が維持しているというのは理論である & \cite{Funahashi1989,Wang2001} & 測定済み \\
11 & 流入が$\approx0.7\tau$より長く続けばビットが書き込まれる（図~\ref{fig:recurrentgroup}(b)） & 方程式$\tau\,\dd r/\dd t=-r+f(wr+I)$を$w=1$、$I=0.6$でオイラー法により計算。\texttt{models/}にスクリプトはない。実験はない & 章内の計算 & モデル \\
12 & 記憶はほとんどスパイクなしにシナプスの促通に保存できる & 理論：終末の残留カルシウムが書き込んだ内容を約1秒保持する。根拠：フェレットの内側前頭前皮質（複数ニューロンの同時記録）には、後効果の長い促通性シナプスで結ばれた錐体ニューロンのサブネットワークがある。記憶保持中の「静かな」区間の観察の総説。皮質がいつどちらの方法を使うかは未解決の問題 & \cite{Mongillo2008,WangMarkram2006,Stokes2015} & 仮説 \\
13 & 記憶は疲労によって消去される。神経伝達物質の蓄えが枯渇する & 皮質錐体ニューロンのペア記録：高頻度のスパイクでシナプスの応答が低下し、低下の速さは放出確率で決まる。これこそが自己維持する群を消すことは直接には示されていない & \cite{Tsodyks1997} & 測定済み \\
14 & 群の連鎖はイベントで起動するタイマーであり、鳴禽ではニューロンが厳密に順番に発火する & 睡眠中と歌唱中のキンカチョウの高次発声中枢で、同定されたニューロンを記録：運動核に出力を送る各ニューロンは、歌の一つの正確な瞬間に短いバーストを1回出し、ニューロンは順に発火する & \cite{Hahnloser2002} & 測定済み \\
\bottomrule
\end{longtable}}

\section{第\ref{sec:gaba}章　\nt{GLUT}と\nt{GABA}の協働}

この章の論理的・力学的な主張は、線形解析とオームの法則から章の中で導いたものである。実験が示すのは、それらが依拠する回路、つまり遅延つきの禁止、興奮に続く直接の抑制、抑制による安定化、\nt{GLUT}--\nt{GABA}ループのリズムが、脳で実際に働いていることである。

{\footnotesize
\setlength{\tabcolsep}{3pt}\renewcommand{\arraystretch}{1.15}
\begin{longtable}{>{\raggedright}p{0.5cm}>{\raggedright}p{4.0cm}>{\raggedright}p{5.6cm}>{\raggedright}p{2.3cm}>{\raggedright\arraybackslash}p{1.7cm}}
\toprule
No. & 主張 & 実験と測定されたもの & 出典 & 状態 \\
\midrule\endfirsthead
\toprule
No. & 主張 & 実験と測定されたもの & 出典 & 状態 \\
\midrule\endhead
1 & \nt{GABA}ニューロンは皮質のニューロンの5分の1から4分の1 & サル皮質の$10$領域でGABAを免疫染色：$8$領域でニューロンの約$25\,\%$、一次視覚野で$20\,\%$未満。皮質の抑制性ニューロンの多様性についての総説 & \cite{Hendry1987,Markram2004} & 測定済み \\
2 & 抑制が何をするかは着地点で大きく決まる。樹状突起、細胞体、スパイクの発生点、他の配線の終末 & 総説：皮質の\nt{GABA}ニューロンのクラスは受け手上の標的で区別される。錐体ニューロンの細胞体と樹状突起からの同時記録：直接の抑制は樹状突起より細胞体ではるかに強い。脊髄で他の配線の終末への抑制は、1本の入力線の神経伝達物質放出を減らす（測定の総説） & \cite{Markram2004,Pouille2001,Rudomin1999} & 測定済み \\
3 & 仲介役を通じた符号の反転には遅延1段、約1ミリ秒がかかる。興奮に続いて届く抑制は同時性の窓を狭める & ラット海馬スライスの錐体ニューロン：仲介役1個を介した抑制は直接の興奮に短い遅延で続き、二つの興奮性入力が閾値まで加算される窓は$2\ms$未満。この抑制が弱い樹状突起では窓が広い & \cite{Pouille2001} & 測定済み \\
4 & 禁止$a\wedge\bar b$~\eqref{eq:veto}はORと定数1とともに関数的に完全である（命題~\ref{prop:gabacomplete}） & 章内の証明。古典的な結果：抑制性入力をもつ閾値素子のネットワークは任意の論理式を実現する。モデルについての主張であり、実験はない & \cite{McCullochPitts1943} & 理論 \\
5 & 遅延つきの禁止は、最適速度$v^*=\Delta x/d$~\eqref{eq:vstar}の運動方向検出器になる & ウサギの網膜：方向選択性は、「ヌル」方向の運動で興奮を打ち消す抑制によって作られる。選択性をもつ細胞の興奮性入力と抑制性入力を別々に記録すると、スターバーストアマクリン細胞（\nt{GABA}）が非対称に、つまり「ヌル」側に向いた突起だけで抑制していることがわかった。$v^*$の式は章の導出 & \cite{Barlow1965,Fried2002} & 測定済み \\
6 & シャント抑制は電圧から引くのではなく割る~\eqref{eq:shunt}（図~\ref{fig:shunt}） & 塩化物イオンの平衡電位が静止電位付近にあるときの、三つのコンダクタンスからなる分圧器のオームの法則。スパイクを出さないニューロンについて & 章内の導出 & 理論 \\
7 & スパイクの発火率に対してシャントは引き算として作用し、割り算は抑制と釣り合った背景ノイズの組み合わせから得られる & モデル：発火しているニューロンではシャントを通る電流はほぼ一定で、発火率への効果は減算的。実験：ラット体性感覚野の錐体ニューロン（スライス）にダイナミッククランプで背景の興奮性・抑制性コンダクタンスを注入した。両者を釣り合わせて同時に増やすと、発火率をあまりずらさずに応答の傾きが割り算される。総説：全体の活動による割り算は脳の多くの部位で見られる & \cite{HoltKoch1997,Chance2002,Carandini2012} & 測定済み \\
8 & $\beta>1$では相互抑制がただ一つの勝者を選ぶ（命題~\ref{prop:wta}、~\eqref{eq:wta}） & 固有値$-(1\pm\beta)/\tau$は章の導出。$\beta=0.5$での発火率$0.73$と$0.53$は固定点$r_{1,2}=(I_{1,2}-\beta I_{2,1})/(1-\beta^2)$。\texttt{models/}にスクリプトはない。章では直接の実験は挙げていない & 章内の導出 & 理論 \\
9 & \nt{GABA}を介した信号によるメモリのリセット。相互抑制する二つの群はラッチである & 図~\ref{fig:bistable}と命題~\ref{prop:wta}から従う。実験はない & 章内の導出 & 理論 \\
10 & \nt{GLUT}--\nt{GABA}ループの平衡は、抑制が十分に強く十分に速ければ安定である（命題~\ref{prop:ei}） & 二集団モデル~\eqref{eq:ei}。条件(i)と(ii)はその行列の行列式とトレースであり、章で導いた & \cite{WilsonCowan1972} & 理論 \\
11 & $w_{EE}=1.5$、$w_{EI}w_{IE}=2$のとき、速い抑制（$\tau_I=2\ms$）は$E^*=0.67h$を保ち、遅い抑制（$30\ms$）は振動を起こす（図~\ref{fig:ei}） & 式~\eqref{eq:ei}の数値解、スクリプト\texttt{figures/make\_ei.py} & 計算 & モデル \\
12 & 皮質は抑制安定化の状態で動作する。その特徴は、\nt{GABA}ニューロンを後押しするとその発火率が下がること~\eqref{eq:isn} & 理論が逆説的な応答を予言した。実験：ネコ一次視覚野の細胞内記録で、周囲の刺激で応答が抑えられるとき、抑制性コンダクタンスは増えずに興奮性コンダクタンスと一緒に減る。マウスの視覚野、体性感覚野、運動野でPVニューロンを光遺伝学的に興奮させると、刺激がなくても、また軽い麻酔下でも、抑制性ニューロンの発火率が下がる。図~\ref{fig:ei}の値での係数$-1/3$は章の導出 & \cite{Tsodyks1997isn,Ozeki2009,Sanzeni2020} & 測定済み \\
13 & 「\nt{GLUT}が\nt{GABA}を興奮させ、\nt{GABA}が\nt{GLUT}を抑制する」ループは周期$12$--$50\ms$の速いリズムを生む & マウス体性感覚野の速い\nt{GABA}ニューロンを$8$--$200$\,Hzで光遺伝学的に刺激すると、ガンマリズム（$20$--$80$\,Hz、周期$12$--$50\ms$）が選択的に強まる。\nt{GLUT}ニューロンの刺激はより遅いリズムだけを強める & \cite{Cardin2009,Buzsaki2012} & 測定済み \\
\bottomrule
\end{longtable}}

\section{第\ref{sec:code}章　モールス符号}

この章の上限の見積もりは情報理論と計算によるものである。測定されているのは、通信路のどこでノイズが生じるか、脳がどれだけ速く動作するか、スパイクがどれだけのコストをもつかである。皮質が実際にどの符号を使っているかという問いは未解決のままである。

{\footnotesize
\setlength{\tabcolsep}{3pt}\renewcommand{\arraystretch}{1.15}
\begin{longtable}{>{\raggedright}p{0.5cm}>{\raggedright}p{4.0cm}>{\raggedright}p{5.6cm}>{\raggedright}p{2.3cm}>{\raggedright\arraybackslash}p{1.7cm}}
\toprule
No. & 主張 & 実験と測定されたもの & 出典 & 状態 \\
\midrule\endfirsthead
\toprule
No. & 主張 & 実験と測定されたもの & 出典 & 状態 \\
\midrule\endhead
1 & 皮質の\nt{GLUT}ニューロンの発火率は毎秒1スパイク未満から数十スパイクで、ニューロンはほとんどの時間を下限の近くで動作する & 覚醒ラットの聴覚野で、細胞に密着させたピペットによる記録（ニューロンの選択はその活動によらない）：各瞬間に毎秒$20$スパイクを超える発火率を示すニューロンは$5\,\%$未満。一部のニューロンの背景発火率は毎秒$0.01$未満。発火率の分布は対数正規 & \cite{Hromadka2008} & 測定済み \\
2 & スパイク通信路の容量は$R_{\max}\approx r\log_2\frac{e}{r\Delta t}$~\eqref{eq:spikeentropy}で、そのほとんどすべてがスパイクの時刻にある（命題~\ref{prop:capacity}） & 長さ$\Delta t$のセルからなる2進文字列のエントロピー。章の数値：$r=10$スパイク/s、$\Delta t=1\ms$でスパイク1個あたり$8$ビット、$\Delta t=10\ms$で約$5$ビット、スパイクカウンタでは毎秒$2$ビット未満 & \cite{MacKay1952} & 理論 \\
3 & 感覚ニューロンはスパイク1個あたり1ビットから数ビットを伝え、最もよい場合で上限の半分に達する & 刺激がわかっている感覚ニューロンのスパイクから刺激を復元。測定のまとめはこの書籍にある & \cite{Rieke1997} & 測定済み \\
4 & スパイク発生器は精確である。精確な時刻を決めるのは入力の速い変化である & ラット皮質のスライスのニューロン：速く揺らぐ電流を繰り返し与えるとスパイクは$1\ms$未満の精度で再現され、一定の電流では時刻がばらける & \cite{Mainen1995} & 測定済み \\
5 & シナプスは当てにならない。放出のランダム性のため、スパイクの半分以上は受け手に届かない & 海馬スライスで錐体ニューロンの入力を最小刺激：スパイクの半分以上が応答を起こさない。刺激閾値の揺らぎ、軸索分岐での伝導の失敗、実験温度の低さは除外された。残るのは確率的な放出である & \cite{AllenStevens1994} & 測定済み \\
6 & 生きた皮質のスパイク列はランダムに見える。間隔はポアソン列のようにばらつき、スパイク数の分散は平均に近い。「ノイズ」の一部は動物の運動についてのメッセージである & 覚醒サルの視覚野V1とMTでの記録：間隔の変動係数は約$1$、スパイク数の分散と平均の比はランダム過程と同様。独立な小さな入力を足し合わせるモデルでは、ばらつきははるかに小さくなる。マウス視覚野では、ばらつきのかなりの部分が動物自身の運動のビデオ記録から予測できる & \cite{Softky1993,Stringer2019} & 測定済み \\
7 & レート符号は遅い。誤差は$1/\sqrt{rT}$~\eqref{eq:rateerror}だが、脳は画像を$\sim150\ms$で認識する & 式はポアソン統計で、章の導出。実験：$20\ms$だけ見せた見知らぬ写真に動物がいるかどうかを被験者が判断した。「はい」と「いいえ」の脳電位は約$150\ms$後に分かれる。この時間を$10$--$20\ms$ずつ約10段に分けるのは章の推定であり、測定ではない & \cite{Thorpe1996} & 測定済み \\
8 & 群での平均化は相関によって制限される。誤差は$1/\sqrt c$分の1より小さくならない~\eqref{eq:correlated}（命題~\ref{prop:pooling}） & 同時記録したサルMT野のニューロンのペア：スパイク数の相関係数は平均$0.12$で、理論解析によればこの程度の相関でも群による改善は制限される。理論：限界を決めるのは相関のうち信号に似た部分だけである。反対の立場のモデル：$50$--$100$個のニューロンの群の発火率は1スパイク間隔、$10$--$50\ms$で読まれ、個々のスパイクの時刻は皮質では使われない & \cite{Zohary1994,MorenoBote2014,ShadlenNewsome1998} & 測定済み \\
9 & 数は最初のスパイクの時刻~\eqref{eq:latency}、群のスパイクの順序、局所リズムの位相に書くことができる & サンショウウオの網膜：短時間提示した画像の空間構造は、出力ニューロンの最初のスパイクの相対遅延に書かれており、コントラストにほとんどよらない。ラット海馬の場所細胞：「自分の」場所を通過するにつれて、スパイクはシータリズム（$7$--$12$\,Hz）のサイクルのますます早い時点へずれ、ずれは$100^\circ$から$355^\circ$に及ぶ。皮質がスパイクの時刻をどれだけ使うかは未解決の問題（8行目） & \cite{Gollisch2008,OKeefe1993} & 測定済み \\
10 & どの符号が読まれるかは受け手の時定数が決める~\eqref{eq:reader}（命題~\ref{prop:reader}）。活動中の皮質ではニューロンは同時検出器に近い & 式~\eqref{eq:reader}は章の導出。in vivo記録の総説：活動中の皮質では膜のコンダクタンスが静止時の数倍になり、時定数はそれに応じて短い。皮質がまさにこのように符号を切り替えることは直接には示されていない & \cite{Destexhe2003} & 仮説 \\
11 & 枯渇するシナプスは発火率の変化、本質的には$\Delta r/r$を伝える~\eqref{eq:depression}（図~\ref{fig:depression}） & 皮質錐体ニューロンのペア記録：枯渇の速さは放出確率で決まる。枯渇が遅いと応答は発火率を反映し、速いと同時性を反映する。ラット皮質で測定された枯渇に基づくモデル：速い入力と遅い入力で発火率の相対変化が等しければ応答も等しく、つまりシナプスは$\Delta r/r$を伝える。数値$1.7\to2.2$、$6.7$、$90\ms$は章の計算 & \cite{Tsodyks1997,Abbott1997depression} & 測定済み \\
12 & バーストは「ツー」である。失われる確率$(1-p)^k$~\eqref{eq:burst}は小さく、促通がそれをさらに下げる & 総説：当てにならないシナプスでは単発のスパイクはしばしば失われるが、バーストは促通のおかげで確実に伝わる。バーストはシナプスの長期的な変化を起こす。脳がバーストを独立した記号として読むというのは仮説 & \cite{Lisman1997,AllenStevens1994} & 仮説 \\
13 & 速い\nt{GABA}ニューロンはリズムと「句読点」を与える。もう一つのクラスは抑制するものを抑制し、ゲートを開く（命題~\ref{prop:punctuation}） & 皮質の\nt{GABA}ニューロンのクラスの性質の総説。光遺伝学：速い\nt{GABA}ニューロンのリズミックな興奮はガンマリズムを起こし、刺激への応答はサイクルの位相に依存する。マウス視覚野では、PVニューロンは互いを抑制し、SSTニューロンは他のすべてのクラスを、VIPニューロンはおもにSSTを抑制する。覚醒マウスでVIPニューロンを興奮させると\nt{GLUT}ニューロンの抑制が解除される。VIPニューロンをオンにするのは強化の信号である。一般規則としての役割分担は仮説 & \cite{Tremblay2016,Cardin2009,Pfeffer2013,Pi2013} & 仮説 \\
14 & エネルギーは平均発火率を毎秒1スパイク程度に制限する~\eqref{eq:energy} & 解剖学的・生理学的データによるげっ歯類灰白質のエネルギー収支：スパイクとシナプス電流は信号伝達のエネルギーの$47\,\%$と$34\,\%$。同時に活動できるニューロンは$\sim15\,\%$以下。ヒトの皮質では、同時に強く活動できるニューロンは$\sim1\,\%$以下。スパイク1個あたり$\sim10^9$個のATPと信号伝達の電力$\sim10$\,Wは、これらの計算に基づく章の推定 & \cite{Attwell2001,Lennie2003} & 理論 \\
15 & 符号はスパースで、1ジュールあたりのビット数が最大になるよう調整されている。皮質は精度をエネルギーと引き換えにする（命題~\ref{prop:economy}） & 経済的な符号の仮説。根拠：食餌制限したマウスの視覚野では、AMPA受容体のコンダクタンスとシナプスのATP消費が下がり（$-29\,\%$）、スパイクの発火率は保たれ、方位選択性は$32\,\%$広がる & \cite{Barlow1961,Padamsey2022} & 仮説 \\
\bottomrule
\end{longtable}}

\section{第\ref{sec:serocomputer}章　\nt{SERO}ニューロン}

\nt{SERO}ニューロン自体の性質（リズム、自己抑制、受容体による符号、サブシステム）は直接測定されている。本章の計算（調節器、フィルタ、論理）は本章の方程式から導かれる。$\tau_c$、セロトニンの拡散係数、放出部位の数は推定である。脳でループがレベル調節器として働くことは仮説である（縫線核では自己抑制は一対一で、短い休止を生むだけである）。\nt{SERO}が時間幅を担うことも仮説だが、これを支持する行動実験がある。

{\footnotesize
\setlength{\tabcolsep}{3pt}\renewcommand{\arraystretch}{1.15}
\begin{longtable}{>{\raggedright}p{0.5cm}>{\raggedright}p{4.0cm}>{\raggedright}p{5.6cm}>{\raggedright}p{2.3cm}>{\raggedright\arraybackslash}p{1.7cm}}
\toprule
No. & 主張 & 実験と測定内容 & 出典 & 状態 \\
\midrule\endfirsthead
\toprule
No. & 主張 & 実験と測定内容 & 出典 & 状態 \\
\midrule\endhead
1 & セロトニンの作用の符号は受け手の受容体が選ぶ（命題~\ref{prop:receptor}） & セロトニン受容体は薬理と機構によってファミリーに分けられる。5-HT$_{1A}$はGタンパク質を介してカリウムチャネルを開き、細胞を抑制する。5-HT$_{2A}$とイオンチャネル型の5-HT$_3$は興奮させる。一つの伝達物質が細胞ごとに逆の応答を生む。 & \cite{BarnesSharp1999} & 測定済み \\
2 & 覚醒時の\nt{SERO}ニューロンは毎秒数回のスパイクを一定に出す & 自由行動下のネコの背側縫線核ニューロンの記録。静かな覚醒時に$2.8\pm0.2$スパイク/sのゆっくりした規則的発火。徐波睡眠で発火率は$34$--$68\,\%$下がり、レム睡眠では$98\,\%$下がる。麻酔下のラットでは$1.7\pm0.5$スパイク/sで、非常に規則的。 & \cite{TrulsonJacobs1979, GagnonParent2014, JacobsAzmitia1992} & 測定済み \\
3 & \nt{SERO}ニューロンはペースメーカーである。スパイクごとの後押しは要らないが、一定の背景入力は要る（スライスでは$\alpha_1$受容体経由のノルアドレナリン、生体では\nt{NORA}から） & 脳スライスでは細胞はゆっくり規則的に発火し、各スパイクの後に$200$--$800\ms$の後過分極がある。スライスで自発発火する細胞は生体より少ない。規則的なリズムには$\alpha_1$受容体を介したノルアドレナリンの持続的入力が必要で、その遮断でリズムは消える。 & \cite{VandermaelenAghajanian1983} & 測定済み \\
4 & 受容体はほぼすべて代謝型。応答は遅く、およそ1秒のうちに立ち上がって減衰する & 5-HT$_3$を除くすべての受容体ファミリーはGタンパク質共役型である。スライスでは、誘発されたセロトニン放出が5-HT$_{1A}$を介して抑制電流を生み、それは1秒以内に立ち上がって減衰する。 & \cite{BarnesSharp1999, CourtneyFord2016} & 測定済み \\
5 & 投射先領域ではセロトニンは間隙だけでなく体積に出る。縫線核自体では近隣の抑制は一対一 & スライスでの高速サイクリックボルタンメトリー。1スパイクあたりの放出量は、シナプスのある領域と、シナプスがほとんどない縫線核とで同じであり、取り込みや受容体の遮断に依存しない。終末あたりの分子数はトランスポーターや受容体の数よりはるかに少なく、細胞外濃度は主要受容体の親和性と一致する。縫線核では5-HT$_{1A}$を介した抑制は一対一で、トランスポーターがセロトニンを間隙外に溜めない。 & \cite{Bunin1998, CourtneyFord2016} & 測定済み \\
6 & \eqref{eq:seroneuron}の濃度は回収によって減る。$\tau_c$が1秒程度というのは推定 & 生体脳でのボルタンメトリー。細胞外セロトニンを主に制限するのは合成ではなく、再取り込みと分解である。引用した研究に$\tau_c$の直接測定はなく、その大きさは本章の推定である。 & \cite{Hashemi2012serotonin} & 測定済み。$\tau_c$は推定 \\
7 & ペースメーカー1個でNOTとNOR。\nt{SERO}論理の完全性（命題~\ref{prop:serocomplete}、図~\ref{fig:serogates}） & \eqref{eq:seroneuron}から導かれる。抑制されるペースメーカーはNORであり、NORは関数的に完全である。\nt{SERO}ニューロンで論理を組んだ実験はない。体積が分かれているという条件は脳では満たされない。 & 本章で導出 & 理論 \\
8 & セロトニンは\nt{SERO}ニューロン自身をその上の受容体を介して抑制する & 麻酔下ラットで、5-HT$_{1A}$作動薬は静脈内投与でもイオントフォレシスでも、セロトニン自体と同じ用量反応関係で背側縫線核ニューロンの発火を抑制する。5-HT$_{1B}$作動薬はほとんど効かない。スライスでは、近隣細胞の内因性セロトニンが5-HT$_{1A}$を介して電流と発火の短い休止を生むが、平均発火率は変わらない。 & \cite{Sprouse1987, CourtneyFord2016} & 測定済み \\
9 & モデルでは、共通の体積を通るループはレベル調節器であり、ばらつきと時定数を$1+G$分の1にする~\eqref{eq:feedback}。脳でループがそう働くことは仮説 & 負帰還増幅器の古典的な式で、本章で導き直した。\nt{SERO}系のループゲイン$G$は測定されていない。測定されているのは、スライスでの自己抑制が一対一で、短い休止を生み、平均発火率を変えないことである。 & \cite{Black1934feedback, CourtneyFord2016} & 理論。脳では仮説 \\
10 & 濃度は発火率の移動平均、ローパスフィルタ~\eqref{eq:lowpass}、図~\ref{fig:lowpass} & 線形方程式~\eqref{eq:seroneuron}の解。図は$\tau_c=1\,\text{s}$でこの式から計算した。\texttt{models/}に個別のモデルはない。 & 本章で導出 & 理論 \\
11 & 間隙の屈曲は小分子の拡散をおよそ$2.6$倍遅くする & 点源法。スライスと生体脳でテトラメチルアンモニウムをイオントフォレシスで与え、その濃度をイオン選択性電極で記録する。屈曲度$\lambda\approx1.6$、つまり実効$D$は自由拡散の$\lambda^2\approx2.6$分の1。細胞外体積の割合は約$0.2$。セロトニン自体の組織内拡散係数はこれらの研究では測定されておらず、本章では推定である。 & \cite{Sykova2008} & 測定済み \\
12 & 独立な「配線」の長さ$\ell\sim\sqrt{D\tau}\approx20$~\textmu m~\eqref{eq:diffusionlength}。配線の数はそのような領域の数で制限される（命題~\ref{prop:volumewires}） & 拡散の法則に$D$と$\tau$の推定値（6行と11行）を代入したもの。命題~\ref{prop:volumewires}はその帰結である。容積伝達の独立チャネルを数えた直接の実験はない。 & 本章で導出 & 理論 \\
13 & 1個の\nt{SERO}ニューロンの出力は脳の広大な範囲に広がる。放出部位は推定$\sim10^5$ & ラット背側縫線核の単一ニューロンの完全再構成。上行軸索はすべて強く分岐し、1細胞の軸索の総長は最大$18.7$~cm、1細胞の枝が線条体、前頭前皮質、扁桃体に届く。細胞あたりの放出部位数は直接数えられていない。 & \cite{GagnonParent2014} & 測定済み。$10^5$は推定 \\
14 & \nt{SERO}ニューロンは出力と応答の異なるいくつかのサブシステムに分かれる & マウスでのウイルス遺伝学的標識。扁桃体に投射するニューロンと前頭皮質に投射するニューロンは分岐がほとんど重ならず、異なる入力を受け、不快刺激に逆向きに応答する。各群の活性化と抑制は行動を異なる形で変える。 & \cite{Ren2018} & 測定済み \\
15 & 指数割引での忍耐の条件$D<\tau_\gamma\ln(R/r_0)$~\eqref{eq:patience}。測定される双曲割引では$D<\tau_\gamma(R/r_0-1)$ & 割引$e^{-D/\tau_\gamma}$または$1/(1+D/\tau_\gamma)$から導かれる。導出は本章。秒単位の遅延での選択課題のサル。割引は指数より双曲線に近い。遅延報酬の予告信号に対する\nt{DOPA}ニューロンの応答も、遅延とともにやはり双曲線的に下がる。 & 本章で導出、\cite{KobayashiSchultz2008} & 理論 \\
16 & \nt{SERO}が時間幅$\tau_\gamma$を設定する（\ref{sec:horizon}節） & マウスで背側縫線核の\nt{SERO}ニューロンを光遺伝学的に活性化すると、遅延報酬を待つ時間が延び、まれになった報酬の探索の粘り強さが増す。ヒトではセロトニンの低下（トリプトファン欠乏食）で遅延報酬の割引が急になる。濃度がどのように$\tau_\gamma$に変わるかは分かっていない。 & \cite{Doya2002, Miyazaki2014, Lottem2018, Schweighofer2008} & 仮説 \\
\bottomrule
\end{longtable}}

\section{第\ref{sec:dofa}章　\nt{DOPA}による学習}

シナプスの機構（パケット、同時検出器、カルシウム、可塑性の窓）と、予測誤差としてのドーパミンのふるまいは直接測定されている。規則が勾配に導くこと、ブロードキャストの代償は定理である。選択学習の結果は本書のモデルによる計算である。誤差を計算する回路は仮説である。

{\footnotesize
\setlength{\tabcolsep}{3pt}\renewcommand{\arraystretch}{1.15}
\begin{longtable}{>{\raggedright}p{0.5cm}>{\raggedright}p{4.0cm}>{\raggedright}p{5.6cm}>{\raggedright}p{2.3cm}>{\raggedright\arraybackslash}p{1.7cm}}
\toprule
No. & 主張 & 実験と測定内容 & 出典 & 状態 \\
\midrule\endfirsthead
\toprule
No. & 主張 & 実験と測定内容 & 出典 & 状態 \\
\midrule\endhead
1 & 人工ネットワークにおける形の誤差逆伝播は脳では見つかっていない & 直接の実験はない。これは不在についての主張である。レビューは、このアルゴリズムが何を要求するか（順方向をなぞる逆向きの結合、独立したフェーズ）と、どんな生物学的近似が提案されているかを検討している。 & \cite{Lillicrap2020, WhittingtonBogacz2019} & 仮説 \\
2 & 重みは放出部位の数、放出確率、1パケットへの応答に分解される：$w=N_s\,p\,A_1$~\eqref{eq:quantal} & 放出を低下させたカエルの神経筋シナプス。受け手の応答は、1パケットに対する自発的な微小応答の整数倍の段階で変動する。スパイクあたりのパケット数はポアソン則に従う。 & \cite{delCastillo1954} & 測定済み \\
3 & 受け手の受容体は同時検出器であり、カルシウムが重みの変化を起動する & マウスニューロンのパッチクランプ。静止電位ではMg$^{2+}$イオンがグルタミン酸で開いたチャネルを塞ぎ、脱分極で外れる。海馬スライス。受け手内のカルシウムキレート剤は長期増強を阻止し、細胞内で光によりカルシウムを放出すると、それだけで伝達が増強される。 & \cite{Nowak1984, Malenka1988calcium} & 測定済み \\
4 & 決定はどちらの側でも実行される：受け手で$A_1$が増え、送り手で$p$が変わる & 1個のスパインの上で光により放出したグルタミン酸は、スパインとその受容体を通る電流の速く持続的な増大を起こす。増強にはAMPA受容体の挿入が伴う。受け手の脱分極は送り手の放出を一時的に抑える。その効果は間隙を逆向きに進むエンドカンナビノイドが運ぶ（\nt{GABA}入力で示された）。短期的な枯渇は送り手の放出確率に依存する。 & \cite{Matsuzaki2004, Malinow2002, WilsonNicoll2001, Tsodyks1997} & 測定済み \\
5 & 送り手と受け手が同時に活動するとシナプスは強まる~\eqref{eq:hebb} & 麻酔下のウサギ。海馬の入力路への短いスパイク列の後、応答は$30$~分から$10$~時間増強される。海馬スライスでは、受け手のスパイクは同じ瞬間に活動しているシナプスだけを強める。 & \cite{Bliss1973LTP, Kelso1986hebb} & 測定済み \\
6 & 規則~\eqref{eq:oja}は入力の相関の主固有ベクトルを見つける（命題~\ref{prop:oja}） & 定理。証明は本章。ニューロンが入力の主成分を学んだという実験はない。シナプスで重みの増大を制限する機構は確立していない。 & \cite{Oja1982} & 理論 \\
7 & タイミングによるヘッブ則：窓は約$\pm20\ms$（図~\ref{fig:stdp}）。繰り返される系列から連鎖が育つ & 培養した海馬ニューロンのペア。送り手のスパイクの後$20\ms$以内の受け手のスパイクは持続的な増強を、前$20\ms$以内のものは抑圧を生み、どちらもNMDA受容体に依存する。図の比$A_-=0.6A_+$は説明用である。ネットワークでこうして連鎖が育つことは直接測定されていない。 & \cite{BiPoo1998} & 測定済み \\
8 & トレース~\eqref{eq:threefactor}：同時の活動の後$1$--$2\,\text{s}$に届いたドーパミンは結合を強め、それより遅いものは強めない（命題~\ref{prop:threefactor}） & 線条体スライスで、グルタミン酸入力とドーパミン入力を別々に光遺伝学的に刺激。ドーパミンがグルタミン酸入力の$0.3$--$2\,\text{s}$後に来たときだけスパインが大きくなる。窓はcAMPの速い上昇と分解で決まる。皮質では、ペアリングが増強と抑圧の別々のトレースを残し、モノアミン受容体がそれを秒程度の窓で長期的な変化に変える。 & \cite{Yagishita2014, He2015eligibility} & 測定済み \\
9 & ドーパミンなしでも秒単位の窓がある：第三の信号は受け手の強い興奮 & 生きたマウスのCA1野。受け手の1回のプラトー事象が、その前後数秒に来た入力を強め、1回の通過で場所野を作る。 & \cite{Bittner2017} & 測定済み \\
10 & 三要素則の平均ステップは報酬の勾配に沿う~\eqref{eq:perturbation}（命題~\ref{prop:gradient}） & ランダムな摂動による学習の勾配定理。本章では小さな雑音について導いた。 & \cite{Williams1992} & 理論 \\
11 & 雑音は探索である：ネットワークは自分のランダムなずれから学ぶ & 若いキンカチョウ。大脳基底核回路から出る核を不活性化すると、歌の変動性が急激かつ可逆的に消える。成鳥のキンカチョウ。平均より一方の側の高さで歌われた音節に妨害を与えると、鳥はすばやく高さを反対側にずらす。自分のばらつきから学ぶ。 & \cite{Olveczky2005, TumerBrainard2007} & 測定済み \\
12 & 二者択一は$\tau_e=1\,\text{s}$、$\delta=R-\bar R$で学習され、$\tau_e=50\ms$では学習されない。期待値を引かないと重みは際限なく増え、学習率が大きいとネットワークは悪い選択を固定しうる & \texttt{models/dofa\_learning.py}、$\eta=0.05$、$500$試行、シード$6$個。$\tau_e=1\,\text{s}$：$A$を選ぶ割合$0.65$--$0.79\to0.96$--$1.00$、重み$0.85$--$1.33$。$\tau_e=50\ms$：$0.38$--$0.55$、重みは変わらない。$\delta=R$：勝者の重み$860$--$1190$。$\delta=R$、$\eta=0.5$、シード$3$個：1個が$B$に固定（$0.04\to0$）。スクリプトは四つの場合をすべて出力する。再計算は本章と一致した。 & \texttt{models/\allowbreak dofa\_\allowbreak learning.py} & モデル \\
13 & 共通の信号一つによる学習時間は、雑音を出すニューロンの数に比例して伸びる（命題~\ref{prop:broadcastcost}） & 線形ネットワークの学習曲線。ニューロン摂動は正確な勾配より出力ニューロン数倍遅く、重み摂動はさらに入力数倍遅い。 & \cite{Werfel2005} & 理論 \\
14 & \nt{DOPA}の出力は主に行動選択の領域に向かう & ラットの黒質線条体\nt{DOPA}ニューロン個々の軸索の完全再構成。1細胞の枝は線条体体積の$0.45$--$5.7\,\%$（平均$2.7\,\%$）を覆い、線条体の外にはほとんど枝がない。 & \cite{Matsuda2009} & 測定済み \\
15 & 皮質は自分の入力を予測することを学び、誤差はその場で計算される & レビュー。感覚皮質には予期した入力と実際の入力の不一致に対する応答がある。これが皮質の一般的な学習則であることは証明されていない。 & \cite{Keller2018} & 仮説 \\
16 & ドーパミンは誤差~\eqref{eq:ctd}のようにふるまう：バースト、予告信号への移行、落ち込み（図~\ref{fig:ctdcases}） & サルの\nt{DOPA}ニューロン。予期しないジュースにバースト。学習後は予告信号にバーストが出て、約束のジュースには応答しない。ジュースが来ないと落ち込む。連続形~\eqref{eq:ctd}は理論である。 & \cite{Schultz1997, Doya2000} & 測定済み \\
17 & $\dd V/\dd t$を計算し期待を差し引く回路 & マウスの腹側被蓋野での記録と光遺伝学。\nt{DOPA}ニューロンは報酬から期待を差し引く。近くの\nt{GABA}ニューロンは報酬が期待されるときにそれを抑制し、その興奮や抑制は誤差を変える。微分がどう計算されるかは示されていない。 & \cite{Eshel2015arithmetic} & 仮説 \\
18 & \nt{DOPA}ニューロンはペースメーカーで、落ち込みはバーストより浅い & スライスで\nt{DOPA}ニューロンは、ゆっくりしたペースメーカー電位の上で自発的に非常に規則的に発火する。サルではレートは正の誤差を定量的に追うが、負の誤差は弱くしか伝えない。 & \cite{GraceOnn1989, Bayer2005} & 測定済み \\
19 & アクターとクリティック（図~\ref{fig:actorcritic}） & アルゴリズムは強化学習の理論から。ヒト（fMRI）では、腹側線条体は選択があってもなくてもクリティックのようにふるまい、背側線条体は行動を選ぶ必要があるときだけそうふるまう。 & \cite{SuttonBarto2018, ODoherty2004actor} & 仮説 \\
20 & 第二の受容体ファミリーをもつニューロンでは規則の$\delta$の符号が反転する & 線条体スライス。D1受容体をもつニューロンではペアリングがD1を要する増強を、D2をもつニューロンではD2を要する抑圧を生む。 & \cite{Shen2008} & 測定済み \\
\bottomrule
\end{longtable}}

\section{第\ref{sec:dofaalt}章　\nt{DOPA}の他の理論}

拡張された描像はそれぞれ、ドーパミンの独自の直接測定（反転点のばらつき、不快なものと新しいものへの応答、ゆっくりした増加、味の変化への応答）に基づいている。しかしそれらを説明する式は理論であり、そのうち二つの間では実験がまだ互いに矛盾している。

{\footnotesize
\setlength{\tabcolsep}{3pt}\renewcommand{\arraystretch}{1.15}
\begin{longtable}{>{\raggedright}p{0.5cm}>{\raggedright}p{4.0cm}>{\raggedright}p{5.6cm}>{\raggedright}p{2.3cm}>{\raggedright\arraybackslash}p{1.7cm}}
\toprule
No. & 主張 & 実験と測定内容 & 出典 & 状態 \\
\midrule\endfirsthead
\toprule
No. & 主張 & 実験と測定内容 & 出典 & 状態 \\
\midrule\endhead
1 & 非対称な規則~\eqref{eq:asymmetric}は点~\eqref{eq:expectile}に収束する。$\kappa=1/2$ではそれは平均、確率$1/2$の1滴では$V=\kappa$（図~\ref{fig:expectiles}） & 点~\eqref{eq:expectile}はエクスペクタイル、つまり正と負の偏差に異なる重みをつけた最小二乗問題の解である。本章の例の導出は本章内にある。 & \cite{NeweyPowell1987}、本章で導出 & 理論 \\
2 & \nt{DOPA}ニューロンごとに反転点が異なり、非対称性の順に並ぶ（命題~\ref{prop:expectile}） & マウス。光遺伝学的に標識した腹側被蓋野の\nt{DOPA}ニューロン、大きさの異なる報酬。バーストが落ち込みに変わる点はニューロンごとに異なる。応答の非対称性が大きいニューロンほどその点は高い。群の応答から報酬分布の形が復元される。これがばらつきの主な機能であることは仮説。 & \cite{Dabney2020} & 測定済み \\
3 & \nt{DOPA}ニューロンの一部は不快なものにバーストで応答する & サル。ある\nt{DOPA}ニューロンは報酬の予告信号で興奮し、目への空気噴射の予告信号で抑制される。別のものは両方で興奮する。二つの群は中脳の異なる部分にある。 & \cite{Matsumoto2009, BrombergMartin2010} & 測定済み \\
4 & 誤差の絶対値または新奇性が学習率を決める & 引用した研究には、\nt{DOPA}ニューロンの重要性信号が学習率を変えることを示す直接の実験はない。レビューは「注意、重要」という信号の役割を提案している。 & \cite{BrombergMartin2010} & 仮説 \\
5 & 危険の軸のための別の配線 & マウス。線条体の後端に向かう\nt{DOPA}ニューロンは報酬ではなく新しい強い刺激に応答する。その破壊は新しい物体の回避を弱める。 & \cite{Menegas2018} & 測定済み \\
6 & 最適な行動時間$\tau_a^*=\sqrt{C/\bar R}$~\eqref{eq:vigor} & 和$C/\tau_a+\bar R\tau_a$の最小値。導出は本章。元のモデルでは、行動の速さは平均報酬に等しい時間のコストで決まる。 & \cite{Niv2007}、本章で導出 & 理論 \\
7 & ドーパミンのゆっくりしたレベルは$\bar R$を運び、行動の速さを決める（命題~\ref{prop:multiplex}） & ラット、側坐核でのマイクロダイアリシスとボルタンメトリー。ドーパミンのレベルは分単位で報酬の頻度と行動の速さに従う。ヒトではL-DOPAが反応時間の平均報酬頻度への依存を強める。遅いレベルがまさに$\bar R$に等しいことは証明されていない。 & \cite{Hamid2016, Beierholm2013vigor} & 仮説 \\
8 & ゆっくりしたレベルは二つの源からなる。放出はスパイクのレートを変えずに出力で変わるが、ゆっくりした増加はレート自体にもある & ラット、一つの課題。側坐核での放出は変化する報酬の期待に従うが、同定された腹側被蓋野\nt{DOPA}ニューロンのスパイクのレートは従わない。仮想通路のマウス（10行）。報酬への接近に伴うなめらかな増加は、同定された\nt{DOPA}ニューロンのスパイクにもある。 & \cite{Mohebi2019, Kim2020} & 測定済み \\
9 & 動物が遠くの報酬に近づく間、ドーパミンはなめらかに増える & 迷路のラット、線条体でのボルタンメトリー。ドーパミン濃度は目標に近づくにつれ数秒かけて上がり、その傾きは距離と報酬の大きさに依存する。 & \cite{Howe2013ramp, Hamid2016} & 測定済み \\
10 & 増加は推定の精密化に伴う誤差$\delta$である~\eqref{eq:ramp}。通路での跳躍はバーストを生む（図~\ref{fig:ramp}） & 式と毎秒$\delta=0.75\,V$という数値は本章で導出。実験：仮想通路のマウスを目標近くへ瞬時に移す。細胞体のスパイク、細胞体と出力のカルシウム、線条体での濃度は、期待$V$ではなく誤差のように応答する。二つの説明のどちらがすべての出力で正しいかは議論中。 & \cite{Kim2020} & 測定済み \\
11 & 振り返り：ドーパミンは因果的結びつきの尺度~\eqref{eq:retro}を伝える & マウス、この理論と誤差~\eqref{eq:ctd}が異なる予測をする実験での側坐核のドーパミン放出。いくつかの実験で放出は前者に従った。 & \cite{Jeong2022} & 仮説 \\
12 & 学習中のドーパミン応答は報酬から予告信号へ徐々に移る & マウス、学習の進行に伴う腹側線条体での\nt{DOPA}ニューロン出力の信号。バーストは\eqref{eq:ctd}による学習が要求するとおり、報酬の時点から予告信号の時点へ徐々に移る。 & \cite{Amo2022} & 測定済み \\
13 & \nt{DOPA}ニューロンは同じ価値での報酬の味の変化に応答する & ラット、腹側被蓋野ニューロンの記録。ジュースを価値の等しい別のものに替えると、誤差~\eqref{eq:ctd}はゼロなのに応答が生じる。 & \cite{Takahashi2017} & 測定済み \\
14 & 人工的なバーストは二つの中立な刺激の結びつきを教える & ラット、報酬なしに二つの刺激をあらかじめ対にする実験。第一から第二への移行時に\nt{DOPA}ニューロンを光遺伝学的に活性化すると結びつきが生じ、抑制するとその学習が妨げられる。 & \cite{Sharpe2017} & 測定済み \\
15 & 特徴ごとの誤差のベクトル~\eqref{eq:vectortd} & 特徴の予測誤差の理論。ベクトル形式の直接の検証はない。 & \cite{Gardner2018} & 仮説 \\
16 & 一つの課題の中で\nt{DOPA}ニューロンごとに異なる量を運ぶ & 仮想通路のマウス、\nt{DOPA}ニューロンの2光子イメージング。報酬に関係するもの、速度と旋回に関係するもの、予告信号と選択の正確さに関係するものがある。 & \cite{Engelhard2019} & 測定済み \\
17 & 配線ではなく配線の束（命題~\ref{prop:bundle}） & 2--16行のまとめ。各分解は個別に測定されている。それらがすべて一つの信号の部分であるという統一的な描像は実験で示されていない。 & --- & 仮説 \\
\bottomrule
\end{longtable}}

\section{第\ref{sec:nora}章　\nt{NORA}}

\nt{NORA}系の構成、その二つのモード、瞳孔との関係（\nt{NORA}経由だけではない）、信号対背景比の上昇、行動における逆U字は直接測定されている。乗算的なゲイン$g$はモデルである。ゲインが選択のモードを切り替え、逆温度として働くことは本章の導出である。ゲイン調整の利点は本書のモデルによる計算である。「\nt{NORA}は探索のつまみ」と「\nt{NORA}は割り込み」は仮説である。

{\footnotesize
\setlength{\tabcolsep}{3pt}\renewcommand{\arraystretch}{1.15}
\begin{longtable}{>{\raggedright}p{0.5cm}>{\raggedright}p{4.0cm}>{\raggedright}p{5.6cm}>{\raggedright}p{2.3cm}>{\raggedright\arraybackslash}p{1.7cm}}
\toprule
No. & 主張 & 実験と測定内容 & 出典 & 状態 \\
\midrule\endfirsthead
\toprule
No. & 主張 & 実験と測定内容 & 出典 & 状態 \\
\midrule\endhead
1 & ヒトの\nt{NORA}ニューロンは数万個で、ほぼすべてが一つの集団にある & ヒト脳幹の連続切片、チロシン水酸化酵素染色。青斑核に$53\,900$個、隣接する亜核にさらに$6\,260$個。 & \cite{Baker1989LC} & 測定済み \\
2 & 出力は脳のほぼ全体に広がるが、集団の部分ごとにある程度異なる領域を担当する & マウスでのウイルス遺伝学的標識。青斑核の\nt{NORA}ニューロンは脳の広い領域から入力を受け、脳と脊髄の広い領域に出力を送る。ラットでは、扁桃体に向かう群は恐怖の学習を強め、前頭前皮質に向かう別の群はその消去を強める。 & \cite{Schwarz2015LC, Uematsu2017LC, Poe2020} & 測定済み \\
3 & 持続モード：毎秒1回未満から数回の規則的なスパイクで、発火率は覚醒度とともに変わる & ラット、睡眠覚醒サイクルでの青斑核ニューロンの記録。発火率は覚醒時に最大、徐波睡眠で低く、レム睡眠でほぼ0。発火率の変化は状態の切り替わりに先行する。 & \cite{AstonJonesBloom1981} & 測定済み \\
4 & 相動モード：課題にとって重要な出来事に、ほぼ決定の瞬間に短いバースト & まれな標的を見つける課題のサル。青斑核のすべてのニューロンが標的にだけバーストで応答し、その$\sim90\ms$後、手の応答の$\sim200\ms$前に出る。成績が悪いときはバーストが弱い。二者択一の課題では、バーストは刺激より応答に正確に結びつき、正答でも誤答でも出る。 & \cite{AstonJones1994vigilance, Clayton2004LC} & 測定済み \\
5 & 瞳孔径は\nt{NORA}の不正確なテストポイント：\nt{NORA}にも\nt{ACET}にも他の領域にも従う & サル。青斑核のスパイクは、1細胞の個々のスパイクに至るまで瞳孔の散大を予測する。似ているがより遅い関係が丘と帯状皮質にもある。マウス。瞳孔のゆらぎは、皮質のノルアドレナリン作動性出力とコリン作動性出力の両方の活動に従う。 & \cite{Joshi2016, Reimer2016} & 測定済み \\
6 & ノルアドレナリンは背景活動を誘発活動より強く抑える。乗算的ゲイン~\eqref{eq:gain}はこの効果を再現するモデル & 覚醒サルの聴覚皮質、ノルアドレナリンのイオントフォレシス。ニューロンの背景活動は同種の声への応答より強く抑えられ、SN比が上がる。乗算的シグモイド~\eqref{eq:gain}はネットワークモデルからとったもので、傾きの文字どおりの乗算は測定されていない。 & \cite{Foote1975NE, ServanSchreiber1990} & 測定済み。$g$はモデル \\
7 & ゲインが$d'$を上げるのは増幅器の後のノイズに対してだけ~\eqref{eq:gainsnr}（命題~\ref{prop:gainnoise}） & 本章で導出。ネットワークモデルではゲインは信号検出を改善する。増幅器の前と後のノイズを分けてこの違いを検証した実験はない。 & 本章で導出、\cite{ServanSchreiber1990} & 理論 \\
8 & 境界$g\beta=1$は選択ネットワークを「思案中」から「決定済み」に切り替える~\eqref{eq:wtagain}（命題~\ref{prop:gainwta}、図~\ref{fig:pitchfork}） & 互いに抑制し合う二つの群の線形解析。本章で導出。脳でのこの閾値の直接測定はない。 & 本章で導出 & 理論 \\
9 & ゲインは選択の逆温度~\eqref{eq:softmax} & 増幅器の後のロジスティックノイズでは、選択確率は$T=\sigma/g$のソフトマックスになる。本章で導出。 & 本章で導出 & 理論 \\
10 & \nt{NORA}はどれだけ探索するかを決める：高い持続性活動は小さな実効$g$（探索）、決定の瞬間の相動性バーストは大きな実効$g$（決定） & ヒト。当て推量の選択の前には、既知の最良の選択の前よりベースラインの瞳孔が広い。瞳孔が広い人ほど探索の傾向が強い。ラット。ランダムな選択のモードへの移行は、前帯状皮質への\nt{NORA}入力が制御する。バーストが実効$g$を上げることは直接測定されていない。 & \cite{JepmaNieuwenhuis2011, Tervo2014stochastic, AstonJones2005} & 仮説 \\
11 & 報酬は$g$に対して逆U字に依存する。速く変わる世界では最良の$g$は小さい（命題~\ref{prop:explore}、図~\ref{fig:invertedU}） & \texttt{models/nora\_explore.py}、$4000$試行の実行$60$回。$1000$試行に1回の変化：最良$g=16$、割合$0.976$。$20$試行に1回：$g=8$、割合$0.886$。$g=0.5$では$0.77$。再計算は本章と一致した。 & \texttt{models/\allowbreak nora\_\allowbreak explore.py} & モデル \\
12 & 行動における逆U字：明瞭なバーストを伴う中程度の持続性活動で成績が最もよい & 4行と同じ課題のサル。成績のよい時期には持続性発火率が中程度で、標的へのバーストが明瞭である。持続性発火率が高いとバーストがなく、誤反応が多い。持続性活動と誘発活動の変化は成績のゆらぎに従う。 & \cite{Usher1999LC, AstonJones2005} & 測定済み \\
13 & \nt{NORA}は予期しない不確実性を、\nt{ACET}は予期された不確実性を伝える & 二種類の不確実性をもつベイズ推論の理論。引用した研究には、これらの信号を伝達物質で分ける直接の実験はない。 & \cite{YuDayan2005} & 仮説 \\
14 & 急な変化の後、人はより速く学び、瞳孔が広がる & ヒト、急な変化のある予測課題。瞳孔の短い変化は変化確率の推定に従い、ベースラインの瞳孔とともに、新しいデータが推定をどれだけ変えるか（学習率）を予測する。光や課題と無関係な瞳孔の変化もこの率を変える。ここで瞳孔がまさに\nt{NORA}を示すというのは、5行の間接的データからの推論である。 & \cite{Nassar2012} & 測定済み \\
15 & 相動性バーストは割り込みとして働く：ネットワークは状態をリセットする & \nt{NORA}のバーストがネットワークの状態をリセットする直接の実験はない。測定されているのは重要な出来事へのバーストだけである（4行）。 & \cite{BouretSara2005} & 仮説 \\
16 & 驚きによる$g$や学習率の調整は、最良の一定設定よりほとんどよくない & \texttt{models/nora\_explore.py}、各時期$1000$試行の世界（$1000$試行に1回と$20$試行に1回の変化）。最良の一定$g=8$：$0.925$。$g$の調整：最大$0.927$。学習率の調整：最大$0.926$。一定の学習率$0.4$：$0.928$。再計算は本章と一致した。 & \texttt{models/\allowbreak nora\_\allowbreak explore.py} & モデル \\
\bottomrule
\end{longtable}}

\section{第\ref{sec:acet}章　\nt{ACET}を加える}

アセチルコリンの三つの作用はスライスと生体脳で測定されている。書き込みと読み出しの衝突は本書のモデルで示した。\nt{ACET}が書き込み許可線として働き、予期された不確実性を伝えることは、部分的に実験の裏づけがある仮説である。

{\footnotesize
\setlength{\tabcolsep}{3pt}\renewcommand{\arraystretch}{1.15}
\begin{longtable}{>{\raggedright}p{0.5cm}>{\raggedright}p{4.0cm}>{\raggedright}p{5.6cm}>{\raggedright}p{2.3cm}>{\raggedright\arraybackslash}p{1.7cm}}
\toprule
No. & 主張 & 実験と測定内容 & 出典 & 状態 \\
\midrule\endfirsthead
\toprule
No. & 主張 & 実験と測定内容 & 出典 & 状態 \\
\midrule\endhead
1 & 脳の\nt{ACET}ニューロンは少なく、いくつかの群にまとまり、出力は皮質全体に広がる：ブロードキャストチャネル & ラットでのアセチルコリン合成酵素に対する抗体染色と逆行性標識。皮質、海馬、扁桃体、嗅球、視床は、主として前脳基底部と脳幹にある六つのコンパクトな細胞群からアセチルコリンを受け取る & \cite{Mesulam1983} & 測定済み \\
2 & 皮質のアセチルコリンレベルは覚醒時に高く、徐波睡眠で低い & 自由行動下のネコの皮質と海馬でのマイクロダイアリシスとEEG記録。アセチルコリン放出は徐波睡眠に比べ、静かな覚醒で$\sim100\%$、活動的な覚醒で$\sim175\%$高い。レム睡眠の皮質では覚醒時と同程度 & \cite{Marrosu1995,Hasselmo1999} & 測定済み \\
3 & 信号の意味は受容体が決める：アセチルコリンには速いチャネル型受容体と、細胞内の反応を起動する遅い受容体がある（命題~\ref{prop:receptor}） & 測定のレビュー。興奮性、伝達、可塑性に対するアセチルコリンの作用は、放出部位、受容体サブタイプ（ニコチン性はイオンチャネル、ムスカリン性はGタンパク質経由）、受け手の細胞に依存する & \cite{Picciotto2012} & 測定済み \\
4 & 第一の作用：外部入力にほとんど触れずに内部結合を弱める（\eqref{eq:acetnet}の係数$1-a$） & ラット嗅皮質のスライス。$100$\,\textmu Mのアセチルコリン作動薬で、同じ細胞において、内部結合（Ib層）の電位は$60\%$以上、外部入力（Ia層）の電位は$15\%$未満低下する。抑制はシナプス前性。海馬スライス（CA1）では内部経路と入力経路で$89\%$対$40\%$ & \cite{HasselmoBower1992,HasselmoSchnell1994} & 測定済み \\
5 & 第二の作用：入力を強める（係数$\frac{1+a}{2}$） & 新皮質スライス。ニコチン性受容体は視床からの入力のシナプスを強め、皮質内のシナプスには作用しない。ムスカリン性受容体は両方を弱める。アセチルコリンはニューロンの興奮性を高める（レビュー） & \cite{Gil1997,Hasselmo2006} & 測定済み \\
6 & 第三の作用：重みの変化を容易にする（学習率$\eta a$） & ラット。基底核（皮質へのアセチルコリンの供給源）の電気刺激を音と対にすると、成体でも聴覚皮質の地図が提示音に合わせて大きく再編される & \cite{KilgardMerzenich1998,Hasselmo2006} & 測定済み \\
7 & \eqref{eq:acetnet}の第二式のまばらなパターン用ヘッブ則は、一部からパターンを補完する連想記憶を与える & まばらなパターンと共分散則をもつネットワークの容量計算。活動ニューロンの割合$p$が小さいと、ネットワークは$p=1/2$の場合よりはるかに多くのパターンを保存する & \cite{Tsodyks1988} & 理論 \\
8 & 低い$a$での書き込みは記憶を壊す：ネットワークは見たものではなく思い出したものを書き込む。$a=0.1$での破損は$a=0$より弱くない & \texttt{models/\allowbreak acet\_memory.py}：ニューロン$200$個、$20$個ずつのパターン五つ、新パターンはそのうち一つと$10$個のニューロンを共有、$10$回の実行。$a=0$では新パターンは記憶されず、古いパターンは$10$個中$5.9$個の無関係なニューロンを引き連れる。$a=0.1$では新パターンは思い出される（$0.97$）が両者が融合し、新パターンは$7.3$個、古いパターンは$7.9$個の無関係なニューロンを引き連れる。$a=0.2$からは1個未満 & 本章で導出 & モデル \\
9 & 命題~\ref{prop:writeread}：書き込みと読み出しが同じようにうまくいくレベル$a$はない（図~\ref{fig:acetsim}） & 同じスクリプト、学習率$\eta a$。新パターンが1回の提示で丸ごと記憶されるのは$a\ge0.9$、$a=0.75$で$0.8$、$a\le0.5$では記憶されない。パターンの半分からの読み出し：$a\le0.2$で$1.0$、$a=0.3$で$0.96$、$a=0.4$で$0.04$ & 本章で導出 & モデル \\
10 & 脳では1本のブロードキャスト配線、\nt{ACET}が書き込みと読み出しを切り替える & この考えはスライスでの測定に基づくモデルからとった。最も直接的な実験はヒトでのもの。ムスカリン受容体遮断薬スコポラミンは新しい単語対の暗記、特に古いものと重なる対の暗記を悪化させるが、すでに学習した対の想起は悪化させない。ネットワークの二つのモードの直接測定はない & \cite{HasselmoAndersonBower1992,Atri2004} & 仮説 \\
11 & フィルタ~\eqref{eq:kalman}は最適である。入力の重みと学習率は同じ量$P/R$で、定常状態では$P=\sqrt{qR}$、記憶は$\sqrt{R/q}$（命題~\ref{prop:kalman}） & 実験を要しない。カルマン--ビューシーフィルタの最適性は推定理論の古典的結果であり、定常状態は本章で導いた & \cite{KalmanBucy1961} & 理論 \\
12 & \nt{ACET}はネットワークに$P$に似た量、予期された不確実性を伝える & 注意の確率モデルで提案された。アセチルコリンのレベルが推定の不確実性に従うことの直接測定はない & \cite{YuDayan2005} & 仮説 \\
13 & \nt{ACET}ニューロンの一部は報酬と罰に20ミリ秒ほどで応答し、強化が予期しないものであるほど強く応答する & マウス、光遺伝学的に同定した前脳基底部のコリン作動性ニューロン。報酬と罰への応答は$18\pm3$\,msで、大きさは強化の予期しなさとともに増し、二つの核のニューロンで応答は似ている & \cite{Hangya2015} & 測定済み \\
14 & \nt{NORA}は世界の予期しない変化に気づき、\nt{ACET}はネットワークを書き込みモードに保つ & この分業はモデルで提案された。間接的には、ヒトで\nt{NORA}に関係する瞳孔径が変化と学習率を追跡する。\nt{NORA}--\nt{ACET}の組の直接の検証はない & \cite{YuDayan2005,Nassar2012} & 仮説 \\
15 & 徐波睡眠では、読み出しモードのネットワークが日中に書き込んだものを再生し、その再生が長期記憶への書き直しになる & ラット海馬のニューロン集団の記録。日中に一緒に働いた細胞は、その後の徐波睡眠でより頻繁に、同じ順序で一緒に発火する（測定済み）。書き直しについて。学習後に海馬のリップルを抑制すると空間記憶が悪化し、自発リップルを延ばすと改善する。原因が\nt{ACET}の低レベルであることは証明されていない & \cite{WilsonMcNaughton1994,Skaggs1996,Girardeau2009,FernandezRuiz2019,Hasselmo1999} & 仮説 \\
\bottomrule
\end{longtable}}

\section{第\ref{sec:synthesis}章　マシン全体}

まとめの章なので新しい実験は挙げない。各行は、その主張を詳しく扱った章と同じ出典にもとづく。\nt{GLUT}バスと\nt{GABA}による流れの制御は確立しているが、ブロードキャストチャネルの役割とその協調は大部分が仮説である。

{\footnotesize
\setlength{\tabcolsep}{3pt}\renewcommand{\arraystretch}{1.15}
\begin{longtable}{>{\raggedright}p{0.5cm}>{\raggedright}p{4.0cm}>{\raggedright}p{5.6cm}>{\raggedright}p{2.3cm}>{\raggedright\arraybackslash}p{1.7cm}}
\toprule
No. & 主張 & 実験と測定されたもの & 出典 & 状態 \\
\midrule\endfirsthead
\toprule
No. & 主張 & 実験と測定されたもの & 出典 & 状態 \\
\midrule\endhead
1 & $8.6\cdot10^{10}$ 個のニューロンに対し、制御信号はわずか四種類~\eqref{eq:machine} & 脳のホモジネート中の細胞核の計数（等方性分画法）：成人男性でニューロンは $(86.1\pm8.1)\cdot10^9$ 個 & \cite{Azevedo2009} & 測定済み \\
2 & 命題~\ref{prop:architecture}の最初の二層：データは\nt{GLUT}のアドレスバスを流れ、結合の符号は送り手が決め、流れは\nt{GABA}ニューロンが方向づけ正規化する（第\ref{sec:neurontypes}, \ref{sec:gaba}章） & ニューロンあたり主要な伝達物質は一つ（測定の総説）；フィードフォワード抑制は錐体ニューロンが入力を加算する窓を狭める；全体の活動による除算は多くの領域で測定されている & \cite{Strata1999,Pouille2001,Carandini2012} & 測定済み \\
3 & 素子は信頼できない：シナプスはスパイクの大部分を失う（表~\ref{tab:computer}、第\ref{sec:code}章） & 海馬スライス、最小刺激：スパイクの半分以上がCA1で応答を起こさない。閾値や軸索伝導の失敗は除外され、原因は確率的な伝達物質放出 & \cite{AllenStevens1994} & 測定済み \\
4 & 脳は $\sim20$ ワットで動き、ニューロンあたり平均で毎秒約1スパイク（第\ref{sec:code}章）；このようなマシンをエンジニアはまだ作れていない & 測定されたコストからの推定~\eqref{eq:energy}：シナプスの仕事を含め1スパイクあたり約 $10^9$ 個のATP分子（げっ歯類の皮質）；皮質についての詳細な推定ではさらに低い発火率になる。技術の現状は総説による & \cite{Attwell2001,Lennie2003,MehonicKenyon2022} & 理論 \\
5 & 大多数のシナプスの学習は局所的な適格度トレースと一つの共通の数 $\delta$ の積（\eqref{eq:machine}の第2式、第\ref{sec:dofa}章） & サルの\nt{DOPA}ニューロンは報酬予測誤差に応答する；線条体スライスでは、ドーパミンがグルタミン酸入力の $0.3$--$2$\,s 後に来たときだけスパインを大きくする。トレースが測定された & \cite{Schultz1997,Yagishita2014} & 測定済み \\
6 & 出力ごとの専用の誤差配線があるのは運動学習の回路だけ（第\ref{sec:motorlearning}章） & 小脳：登上線維の信号と入力が同時に来るとその入力が弱まる；サルではプルキンエ細胞の複雑スパイクが運動誤差の方向を運び、補正を起こす & \cite{Ito2001,Herzfeld2018} & 測定済み \\
7 & 各ブロードキャストチャネルは数本の線の束（命題~\ref{prop:bundle}） & \nt{DOPA}について：同じ課題で異なるニューロンが報酬、運動、刺激の特徴を運ぶ；速い誤差と遅いドーパミンレベルは別物。\nt{SERO}, \nt{NORA}, \nt{ACET}ではこの分解はあまり調べられていない & \cite{Engelhard2019,Hamid2016,Mohebi2019} & 測定済み \\
8 & \nt{SERO}：レベル調整器であり、仮説では時間幅 $\tau_\gamma$ も決める（第\ref{sec:serocomputer}章） & 自己抑制：自らの5-HT$_{1A}$受容体のアゴニストが縫線核のセロトニンニューロンを抑制する（測定済み）。時間幅は理論で提案された；最も強い実験は、マウスでこれらのニューロンを光遺伝学的に活性化すると遅延報酬を待つ時間が延びること & \cite{Sprouse1987,Doya2002,Miyazaki2014} & 仮説 \\
9 & \nt{NORA}：ゲイン $g$ が探索か決定かを選ぶ（第\ref{sec:nora}章） & SN比の向上：覚醒したサルの聴覚皮質で、ノルアドレナリンは背景活動を同種の鳴き声への応答より強く抑える（測定済み）；乗算的ゲインはモデル；探索と決定の選択での役割は理論 & \cite{Foote1975NE,ServanSchreiber1990,AstonJones2005} & 仮説 \\
10 & \nt{ACET}：書き込みと読み出しを切り替える（第\ref{sec:acet}章） & 三つの作用（内部結合の弱化、入力の強化、可塑性の促進）は測定済み；モードの切り替えはヒトでのスコポラミン実験が間接的に支持 & \cite{HasselmoBower1992,Gil1997,Atri2004} & 仮説 \\
11 & 式~\eqref{eq:machine}はマシン全体を記述する & 第\ref{sec:glutcomputer}--\ref{sec:acet}章のモデルの要約で、各部分はそれぞれの章にもとづく；全体としては実験で検証されていない & 章内の導出 & 仮説 \\
12 & つまみは共通の制御なしに協調して働く：誤差、意外さ、書き込み、復帰（図~\ref{fig:timeline}） & 実験はない：模式図は定性的。個々の環は、\nt{NORA}と\nt{ACET}の分業の理論と、ヒトでの瞳孔と学習速度の関係 & \cite{YuDayan2005,Nassar2012} & 仮説 \\
13 & 一つの共通信号による学習は、結果に影響するニューロンが多いほど遅い（命題~\ref{prop:broadcastcost}） & 線形ネットワークの学習曲線の計算：出力の摂動による学習は勾配降下より出力の数の倍だけ遅い & \cite{Werfel2005} & 理論 \\
14 & 皮質が多層回路をどう調整するかは不明；候補はスパイクのバースト、ニューロンの枝の中での計算、予測による自己教師あり学習 & 誤差の担い手としてのバーストはモデル；皮質ニューロンの詳細モデルの応答を再現できたのは $5$--$8$ 層のネットワークだけ（モデル）；ヒト皮質ニューロンの枝で排他的論理和を生むカルシウムスパイクはスライスで測定済み；自己教師あり学習は証拠の総説 & \cite{Payeur2021,Beniaguev2021,Gidon2020,Keller2018} & 仮説 \\
\bottomrule
\end{longtable}}

\section{第\ref{sec:sensors}章　マシンの入力}

センサの仕組みは、単一細胞の電流記録、蝸牛の振動の記録、網膜の細胞の計数によって直接測定されている。感度の限界とショットノイズの式は物理から導かれる。センサの符号が情報を最大にするよう調整されているというのは、強い裏づけのある仮説である。

{\footnotesize
\setlength{\tabcolsep}{3pt}\renewcommand{\arraystretch}{1.15}
\begin{longtable}{>{\raggedright}p{0.5cm}>{\raggedright}p{4.0cm}>{\raggedright}p{5.6cm}>{\raggedright}p{2.3cm}>{\raggedright\arraybackslash}p{1.7cm}}
\toprule
No. & 主張 & 実験と測定されたもの & 出典 & 状態 \\
\midrule\endfirsthead
\toprule
No. & 主張 & 実験と測定されたもの & 出典 & 状態 \\
\midrule\endhead
1 & センサは変換器である：受容体が電流を生み、目と耳の感覚細胞の出力は\nt{GLUT}バスへのグルタミン酸である；最初の細胞はスパイクを出さない（図~\ref{fig:transducer}） & カメの桿体と錐体は興奮性アミノ酸を放出する：切り取った膜片のグルタミン酸チャネルがそれを捉える。ラット蝸牛の内有毛細胞：神経線維終末の電流はグルタミン酸AMPA受容体を通り、放出頻度は細胞のなめらかな脱分極とともに $150$\,Hz まで増す & \cite{CopenhagenJahr1989,GlowatzkiFuchs2002} & 測定済み \\
2 & 暗所の光センサは光子1個に約1ピコアンペアの電流で応答する & ヒキガエル桿体の外節への吸引電極：弱い閃光への応答は量子化されており、1事象は $\approx1$\,pA（暗電流の $5\%$）、ばらつき $0.2$\,pA、効率 $0.5$。ヒトは光子1個の到来を偶然より高い頻度で報告する & \cite{Baylor1979,Tinsley2016} & 測定済み \\
3 & 音のセンサは1ナノメートル未満の変位を検出する & メスバウアー法、モルモット蝸牛の基底膜の $16$--$19$\,kHz の位置：聴神経の応答しきい値で膜の速度は約 $0.04$\,mm/s、すなわち変位 $\approx0.35$\,nm（換算は本書による）。測ったのは膜であって、センサの毛束ではない & \cite{Sellick1982,Hudspeth1989} & 測定済み \\
4 & 暗所での精度は光子のショットノイズ~\eqref{eq:photonnoise}で制限される；弱い光では蓄積が長くなる & 式はポアソン統計から導かれる。桿体では弱い光のもとでのノイズが独立な量子事象の和と一致する；明るい背景では閃光への応答が小さく短くなる。「十分の数秒」という数値はここでは別に検証されていない & \cite{Baylor1979} & 理論 \\
5 & 入力のレンジは光で10桁、音で12桁だが、スパイクの発火率は3桁未満 & 音について：特徴周波数での基底膜の応答は $40$--$80$\,dB の帯域で傾き $0.2$\,dB/dB まで増し、圧縮は $100$\,dB を超えても見える。桁数そのものは標準的な推定で、章では出典なし & \cite{Ruggero1997} & 測定済み \\
6 & センサの応答は~\eqref{eq:nakarushton}で、$\sigma$ は平均の明るさを追い、曲線を対数軸に沿ってずらす（図~\ref{fig:adaptation}） & 式は最初、魚の網膜のS電位に当てはめられた。カメの錐体：応答は $V=I V_m/(I+\sigma)$ に従い、明るさの $\sim3.5$ 桁をカバーする；背景に $2$--$3$ 分さらすと、曲線は幅を保ったまま水平にずれる。全体の活動による除算との関係は総説 & \cite{NakaRushton1966,NormannPerlman1979,Carandini2012} & 測定済み \\
7 & センサは変化を伝え、その符号は1スパイクあたりの情報が最大になるよう調整されている（命題~\ref{prop:sensors}） & 原理は理論的に提案された。最も強い証拠：ハエの第1層ニューロンの「コントラスト $\to$ 応答」特性は、自然情景のコントラストの累積分布と一致する。こうして情報が最大になる & \cite{Barlow1961,Laughlin1981} & 仮説 \\
8 & オン型とオフ型のセンサは2線式符号である；イベントカメラはそれをシリコンで再現する & 実験の総説：網膜のオンチャネルとオフチャネルは最初のシナプスですでに分かれ、別々に遮断できる。カメラ：フレームなしの $128\times128$ 画素、レンジ $120$\,dB、遅延 $15$\,\textmu s & \cite{Schiller1992,Lichtsteiner2008,Gallego2022} & 測定済み \\
9 & 目には $\sim10^8$ 個の光センサと $\sim10^6$ 個の出力ニューロンがある：$\sim100$ 分の1への圧縮 & ヒト網膜の全載標本での計数：桿体は平均 $92\cdot10^6$ 個、錐体は $4.6\cdot10^6$ 個；出力（神経節）細胞は目によって $0.7\cdot10^6$ から $1.5\cdot10^6$ 個 & \cite{Curcio1990,CurcioAllen1990} & 測定済み \\
10 & マウスでは目の出力チャネルは30種類を超える & マウス：$11\,000$ 個を超える出力細胞の2光子カルシウムイメージングと応答のクラスタリング。機能的な型は $30$ を明らかに超える。ヒトでは数は測定されていない & \cite{Baden2016} & 測定済み \\
11 & モルモットの目は $\sim875\,000$ bit/s を送る；ヒトの目に換算すると $\sim10^7$ bit/s & モルモット、多電極アレイ、自然情景の中の運動：細胞あたり $6$--$13$ bit/s、$\sim10^5$ 個の出力細胞で $\sim875\,000$ bit/s。ヒト（$\sim10^6$ 個の細胞）の $\sim10^7$ bit/s は著者らの換算 & \cite{Koch2006} & 測定済み \\
12 & 耳はほぼ対数的な周波数マップを持つバンドパスフィルタのバンクである（図~\ref{fig:filterbank}） & 基底膜の振動が最大になる場所は周波数によって決まる；「周波数--場所」関数はほぼ指数関数的で、一つの形でヒトと生きた動物の蝸牛を記述する & \cite{Greenwood1990,Robles2001} & 測定済み \\
13 & 共振器は能動的である：一部のセンサが弱い振動を振幅で数千倍（$66$--$76$\,dB）増幅し、音が大きくなると増幅は下がる & チンチラ蝸牛基部のレーザー振動計測：特徴周波数の弱い純音で、アブミ骨に対する増幅は $66$--$76$\,dB；死後は感度が $60$--$81$\,dB（振幅で $10^3$--$10^4$ 倍）下がり、圧縮が消える。力の源は外有毛細胞 & \cite{Ruggero1997,Robles2001} & 測定済み \\
14 & 耳の増幅器は自励発振の境界に置かれている & 計算：ホップ分岐点の近くでは、レンジの圧縮、鋭い同調、結合音が避けられない。これは聴覚の既知の非線形性のすべてである。蝸牛がまさにそう調整されていることは直接には証明されていない & \cite{Eguiluz2000} & 仮説 \\
15 & 低い周波数ではスパイクは音と位相をそろえて出て（モルモットでは同期は $\sim600$\,Hz から弱まり、$\sim3.5$\,kHz まで認められる）、それから方向が求まる；高い周波数では場所符号が残る & モルモットの聴神経：位相同期は約 $600$\,Hz で落ち始め、$3.5$\,kHz を超えると消える；ネコとサルでは境界がほぼ1オクターブ高い。両耳への到達時間差は総説 & \cite{PalmerRussell1986,Grothe2010} & 測定済み \\
16 & その他のセンサ（表~\ref{tab:sensors}）：嗅覚は $\sim400$ 種類の受容体をニューロンあたり1種類持ち、組み合わせ符号を使う；味覚は一部ATPとセロトニンを使う；触覚は速く順応するセンサと遅く順応するセンサ；温と冷は別々 & マウス：一つの受容体が多くの匂いを認識し、一つの匂いが多くの受容体を活性化し、異なる匂いは異なる組を活性化する。ATP受容体がないと味覚神経は応答しない；味蕾はセロトニンを放出する。ヒトの手の皮膚の単一線維記録：順応の速さで四つの型。冷のチャネルは温のチャネルとは別物 & \cite{Mombaerts2004,Malnic1999,Finger2005,Huang2005,JohanssonVallbo1979,McKemy2002} & 測定済み \\
\bottomrule
\end{longtable}}

\section{第\ref{sec:recognition}章　認識}

認識の速さ、最初の段の仕組み、応答の線形な読み出し、時間の連続性による学習は、ヒトとサルで測定されている。連鎖全体を二つの演算に帰着させることと動画からの学習は本書のモデルで検証した。錯視の説明は、十分な根拠のあるものから論争中のものまである。

{\footnotesize
\setlength{\tabcolsep}{3pt}\renewcommand{\arraystretch}{1.15}
\begin{longtable}{>{\raggedright}p{0.5cm}>{\raggedright}p{4.0cm}>{\raggedright}p{5.6cm}>{\raggedright}p{2.3cm}>{\raggedright\arraybackslash}p{1.7cm}}
\toprule
No. & 主張 & 実験と測定されたもの & 出典 & 状態 \\
\midrule\endfirsthead
\toprule
No. & 主張 & 実験と測定されたもの & 出典 & 状態 \\
\midrule\endhead
1 & ずれがあると、画素ごとのテンプレート比較はコイン投げと変わらない；段の対~\eqref{eq:scstage}は図形をどの位置でも認識する & \texttt{models/\allowbreak recognition\_models.py}、$24$ 点の「網膜」、$4000$ 試行：反転した点の割合 $0$, $0.1$, $0.2$ でテンプレートは $0.49$, $0.51$, $0.52$；$S$ と $C$ の対は $1.00$, $0.86$, $0.70$ & 章内の導出 & モデル \\
2 & 絵の中の動物は $\sim150$\,ms で認識される。十数段を1段 $10$--$20$\,ms で1回通過する & ヒト、$20$\,ms 提示の写真で「動物がいるかいないか」の課題：二つの答えの誘発電位は $\sim150$\,ms で分かれる。サル：最初の応答の時刻は段を追って遅くなる（V1が最も早く、次にV2とV4）。段の数と「1段に0個か1個のスパイク」はこれらの数値からの推論 & \cite{Thorpe1996,Schmolesky1998} & 測定済み \\
3 & $S$ 段は部品の AND、$C$ 段は位置についての最大値（命題~\ref{prop:conveyor}） & ネコの一次視覚野：単純細胞は特定の場所の特定の向きのエッジに、複雑細胞は同じ向きにより広い領域で応答する。複雑細胞の細胞内記録：多くの細胞で2本のバーへの応答は二つの応答の大きいほうに近く、平均では max 演算に最も近い。相互抑制による最大値は命題~\ref{prop:wta}、全体の活動による除算は総説 & \cite{HubelWiesel1962,Lampl2004,Carandini2012} & 測定済み \\
4 & 連鎖を上るほど組み合わせは大きく複雑になり、応答は位置にますます依存しなくなる & サル、刺激を「臨界特徴」まで単純化：V4と後部下側頭皮質の細胞は小さな受容野で形、色、テクスチャの複雑な組み合わせに応答する；前部下側頭皮質では受容野が大きい。モデルの $S$ と $C$ の交互の繰り返しはこの描像を再現する & \cite{KobatakeTanaka1994,Riesenhuber1999} & 測定済み \\
5 & 畳み込みネットワークはこの交互の繰り返しを再現し、上位段の応答を最もよく予測する；物体はニューロン集団の発火率から線形に読み出せる & 脳へのあてはめなしに認識を学習したネットワーク：最上層はサル下側頭皮質のニューロン応答を、中間層はV4の応答を予測する。ランダムに選んだ $\sim100$ 個の細胞の $12.5$\,ms 以上の窓での活動から、線形分類器が位置や大きさが違っても物体とカテゴリーを認識する & \cite{Fukushima1980,Yamins2014,Hung2005} & 測定済み \\
6 & トレース付きヘッブ則~\eqref{eq:traceinvariance}は、トレースが $\sim20$\,ms 以上なら、ラベルなしの動画から不変性を学習する & 遅い特徴とトレース付き規則のアイデアは理論。\texttt{models/\allowbreak recognition\_models.py}、二つの $C$ ニューロン、入力 $16$、$10$ 回の実行、物体間の休止 $0.3$\,s。$\tau_{\text{tr}}=5$\,ms で図形との一致は平均 $0.52$、$10$\,ms で $0.56$；$20$, $50$, $200$\,ms ではすべての実行で $1.00$（$30$\,ms では10回中1回誤り） & \cite{Foldiak1991,WiskottSejnowski2002} & モデル \\
7 & 脳の中の不変性は時間の連続性から学習される & サル：視野の特定の場所への眼の跳躍中に物体をこっそり差し替えた；下側頭皮質ニューロンの位置不変性は差し替えに合わせて変わり、1時間後にはすでに有意だった。これが視覚学習の主要な規則だというのは仮説 & \cite{LiDiCarlo2008} & 測定済み \\
8 & 運動：方向検出器、次に広い領域についての同じ AND／OR の対 & ウサギ網膜の方向検出器；サルのMT野では記録した $110$ 個のニューロンすべてが方向選択的で、運動への応答はV1より強い & \cite{Barlow1965,Albright1984} & 測定済み \\
9 & 上位段は下へ予測を送り、上へは誤差が送られる & モデルは一次視覚野の受容野の古典外効果を説明する；個々の観察は整合する（総説）が、パイプラインの一般的な仕組みとしては証明されていない & \cite{RaoBallard1999,Keller2018} & 仮説 \\
10 & 難しい画像にはフィードバックと何周かの処理が要る & サル：「難しい」画像では下側頭皮質での決定が $\sim30$\,ms 遅れて形成される；この遅い応答は順方向ネットワークではうまく予測できず、フィードバックのあるネットワークやより深いネットワークのほうがよく予測する & \cite{Kar2019} & 測定済み \\
11 & 気づかれない画素の改ざんはネットワークをだます；ヒトの視覚ははるかに強いが無敵ではない & ネットワークで：特別に選んだ小さな摂動が答えを変える；「パンダ $\to$ テナガザル」の例は2番目の論文から。ヒトでは短時間提示のとき、ネットワーク間でよく転移する改ざんが選択をずらす & \cite{Szegedy2014,Goodfellow2015,Elsayed2018} & 測定済み \\
12 & 期待の錯視：信頼できない入力は期待に譲る。淡いものは「より遅く」動き、ドレスの見え方は人によって違う（\ref{sec:illusions}節） & コントラストの低い格子はより遅く動いて見える。$1401$ 人の調査：ドレスは三つの安定した見え方に分かれ、照明のヒントで色が切り替わる；$\sim13\,000$ 人の観察者では照明についての想定が決め手になる。遅い運動を期待するベイズモデルが一連の運動錯視を説明する & \cite{Thompson1982,LaferSousa2015,Wallisch2017,Weiss2002,Purves2011} & 仮説 \\
13 & ゲイン調整：滝、傾き、コントラスト；ヘルマン格子は近傍抑制ではない & ウサギ網膜の方向検出器は長い運動の間に発火率を下げ、停止後は背景以下に落ちる。傾きの錯視は近傍の活動による除算のモデルで再現される。網膜の出力ニューロンの応答を変えない格子の変形で斑点が消える。近傍抑制による説明は否定された & \cite{BarlowHill1963,Schwartz2009,SchillerCarvey2005} & 測定済み \\
14 & 遅延の補償：動く物体は閃光より前にあるように見える & 錯視は測定済み。心理物理は運動の前方延長とも遅延の差とも矛盾する；閃光後 $\sim80$\,ms の出来事にもとづく事後的な説明が提案されている。論争は決着していない & \cite{Nijhawan1994,EaglemanSejnowski2000} & 仮説 \\
15 & 静止した「ヘビ」が回る：遅延の差~\eqref{eq:latency}が方向検出器に読まれる & 錯視は色がなくても働く；隣り合う要素の対が単独で錯視を生む；方向選択的なサル皮質のニューロンは、この静止した対に方向性のある応答を示す。ヒトでは眼のマイクロサッカードと瞬きが回転を起動する & \cite{Conway2005,OteroMillan2012} & 測定済み \\
16 & 多義図形：相互抑制、疲労、ノイズ；切り替えは主にノイズが起こす & 競合するニューロン集団のモデル：切り替えを起こすのは\emph{ノイズ}であり、ノイズがなければ切り替えは起きない；刺激の強さとともに切り替え頻度が増すことを説明する。ここでは疲労の役割は小さい & \cite{MorenoBote2007} & 仮説 \\
17 & 錯視の一部はマシンが学習する課題から生じる & 動画の次フレームを予測するよう学習したネットワークは静止した輪の回転を「見て」、対照画像では見ない；画像のノイズ除去と鮮鋭化のネットワークはヒトと同様に色と明るさの錯視にかかる & \cite{Watanabe2018,GomezVilla2020} & 測定済み \\
\bottomrule
\end{longtable}}

\section{第\ref{sec:muscles}章　筋肉への出力}

出力の仕組み、すなわち運動単位、シナプスの安全率、大きさの順の動員、筋肉の対、伸張ループ、リズム発生器は直接測定されている。遅延ループの安定条件は定理である。体の内部モデルは十分な根拠のある仮説である。図の数値は本書のモデルによる計算である。

{\footnotesize
\setlength{\tabcolsep}{3pt}\renewcommand{\arraystretch}{1.15}
\begin{longtable}{>{\raggedright}p{0.5cm}>{\raggedright}p{4.0cm}>{\raggedright}p{5.6cm}>{\raggedright}p{2.3cm}>{\raggedright\arraybackslash}p{1.7cm}}
\toprule
No. & 主張 & 実験と測定されたもの & 出典 & 状態 \\
\midrule\endfirsthead
\toprule
No. & 主張 & 実験と測定されたもの & 出典 & 状態 \\
\midrule\endhead
1 & 運動単位：1個の\nt{ACET}ニューロンが数百から2000本ほどの線維のグループを制御する；細かな調整を要する筋肉では単位が小さい & ヒトの筋肉、運動軸索と筋線維の計数：ニューロンあたりの線維は平均で手の骨間筋で約 $340$、腕橈骨筋で約 $410$、腓腹筋内側頭で約 $1930$。最も小さな単位の正確な数は章では挙げていない。 & \cite{Feinstein1955mu} & 測定済み \\
2 & 筋肉上のシナプスは、線維の発火に必要な量の数倍の伝達物質を放出する & 神経筋シナプスでの測定の総説：安全率（放出された伝達物質と閾値の比）は成体哺乳類で通常 $3$--$5$；ヒトでは安全率は放出量よりも受け手の膜のひだによって保たれている。 & \cite{WoodSlater2001} & 測定済み \\
3 & 単収縮~\eqref{eq:twitch}は $50\ms$ 程度の時間 $T_c$ で立ち上がる & 随意収縮中のヒトの手の単一運動単位、単位のスパイクで加算平均した力：立ち上がり時間 $30$--$100\ms$、$80\,\%$ の単位で $70\ms$ 未満。関数 $(t/T_c)e^{1-t/T_c}$ は運動単位プールのモデルからの近似。 & \cite{MilnerBrown1973rec, Fuglevand1993} & 測定済み \\
4 & 単収縮の和~\eqref{eq:forcesum}：$5$, $15$, $40$ スパイク/s で力は $0.2$--$1.1F_1$、$1.8$--$2.2F_1$、約 $5.4F_1$（図~\ref{fig:tetanus}） & \texttt{models/\allowbreak muscle\_\allowbreak models.py} による計算、$T_c=50\ms$：最後の $0.3\,\text{s}$ で $0.21$--$1.10$、$1.76$--$2.19$、$5.28$--$5.49\,F_1$。線形の加算は単純化：実際の筋肉の飽和はモデルにない。 & \texttt{models/\allowbreak muscle\_\allowbreak models.py} & モデル \\
5 & 単位は大きさの順に動員される；最も弱いものは最も強いものの100分の1 & ネコの前根：反射入力を強めていくと、前根でのスパイクが小さいニューロン、つまり小さいニューロンが最初に発火する。ヒトの手の骨間筋：単位の単収縮の力は $0.1$ から $10$~g で、単位が動員される力にほぼ比例して増える。 & \cite{Henneman1957, MilnerBrown1973rec} & 測定済み \\
6 & 力の刻みは力に比例する、$\Delta F/F\approx q-1$~\eqref{eq:sizeprinciple}：浮動小数点 & 章で力の等比数列から導出。実験と整合する：大きな力では、同じ力の増分に対して新たに動員される単位がずっと少ない。 & 章内の導出；\cite{MilnerBrown1973rec} & 理論 \\
7 & 力のばらつきは力そのものに比例して増える & 母指の筋の随意等尺性収縮：力の標準偏差は平均の力に比例して増える；電気刺激で同じ力を出させるとばらつきは一定。プールのモデル：力の順の動員が線形の増加を生む。 & \cite{Jones2002sdn} & 測定済み \\
8 & 動きのなめらかさは信号依存ノイズの帰結である & モデル：このようなノイズのもとで終点のばらつきを最小にする軌道は、測定された眼と腕の軌道を再現する。この原因を他の原因と区別する直接の実験はない。 & \cite{HarrisWolpert1998} & 仮説 \\
9 & 筋肉の対の力の差がトルクを、和が剛性を決める~\eqref{eq:antagonists} & 同時収縮によるインピーダンス制御の理論。ヒトの腕：小さな摂動で測った各関節の剛性は、その関節のトルクにほぼ比例する；同時収縮の比率を変えると手先の剛性楕円の形と向きが変わる。 & \cite{Hogan1984, GomiOsu1998stiff} & 測定済み \\
10 & 伸張センサは自分の\nt{ACET}ニューロンを興奮させ、抑制性の介在ニューロンを通じて拮抗筋をオフにする（図~\ref{fig:antagonists}） & ネコの脊髄：伸張センサの線維に興奮させられる単一の介在ニューロンが、拮抗筋の\nt{ACET}ニューロンに単シナプス性の抑制性電位を生む。シナプス遅延 $0.28$--$0.42\ms$。介在ニューロンの伝達物質（グリシン）はこの実験では同定されていない。 & \cite{Jankowska1972Ia} & 測定済み \\
11 & ループの遅延：脊髄 $25$--$30\ms$、皮質経由はその1.5〜3倍、目経由 $100$--$200\ms$ & ヒトの腕、関節への摂動：筋肉の短潜時応答は $20$--$45\ms$、長潜時応答は $45$--$105\ms$、随意応答は $120$--$180\ms$。運動中に見える指の位置をずらすと、修正は平均 $160\ms$ 後。 & \cite{Pruszynski2008rapid, SaundersKnill2003vis} & 測定済み \\
12 & $\dd x/\dd t=-Gx(t-d)$ は $Gd<\pi/2$~\eqref{eq:delaystable}で安定；境界での周波数は $1/(4d)$ & 遅延方程式の根についての定理。\texttt{models/\allowbreak muscle\_\allowbreak models.py} による検証、$d=30\ms$：$3\,\text{s}$ の時点で振幅は $G=40$ で $3\cdot10^{-6}$、$G=50$ で $0.13$、$G=55$ で $38$（単位は毎秒、境界は $52$）。 & \cite{Hayes1950} & 理論 \\
13 & 生理的振戦 $8$--$12$~Hz；伸張ループはその源の一つ & 測定の総説：健康な人には約 $10$~Hz 帯の細かな震えがある；その起源は混合的で、中枢の振動、単位の発火の性質、機械的共振、反射ループの共振があり、条件によって優勢なものが異なる。 & \cite{McAuleyMarsden2000} & 仮説 \\
14 & 脳は体の内部モデルによって指令の結果を予測する & ヒトが物体をつまんで持ち、腕を動かす：慣性、粘性、弾性の負荷のいずれでも、把持力はループの遅れなしに負荷と同時に変わる。つまり負荷は予測されている。総説がこうした実験をまとめる；ニューロン内のモデルそのものの直接観察はない。 & \cite{FlanaganWing1997grip, WolpertGhahramani2000} & 仮説 \\
15 & 歩行、呼吸、咀嚼のリズムは、一定の入力のもとで局所的なネットワークが生む & 実験の総説：摘出した神経系（脊髄、神経節）は、リズミカルな入力も筋肉からのフィードバックもなしにリズミカルな運動パターンを出す。 & \cite{MarderBucher2001} & 測定済み \\
16 & 疲労をともなう相互抑制~\eqref{eq:halfcentre}は周期 $0.84\,\text{s}\approx2\tau_a$ のリズムを生む（図~\ref{fig:halfcentre}） & 定理：相互抑制と適応を持ち、一定の入力のもとで安定状態を持たないネットワークは必ず振動する。\texttt{models/\allowbreak muscle\_\allowbreak models.py} による計算（$I=1$, $\beta=2$, $b=2$, $\tau=20\ms$, $\tau_a=0.4\,\text{s}$）：周期 $0.84\,\text{s}$。実際の発生器がまさにこのような構成であることは示されていない。 & \cite{Matsuoka1985osc} & モデル \\
\bottomrule
\end{longtable}}

\section{第\ref{sec:pain}章　痛み}

痛みのセンサ、2本の配線、脊髄のゲート、下行性の音量つまみ、感作、注意の奪取は直接測定されている。ゲートと感作の式は本書の単純化したモデルである。マシンが警報の音量を選ぶ規則は仮説である。

{\footnotesize
\setlength{\tabcolsep}{3pt}\renewcommand{\arraystretch}{1.15}
\begin{longtable}{>{\raggedright}p{0.5cm}>{\raggedright}p{4.0cm}>{\raggedright}p{5.6cm}>{\raggedright}p{2.3cm}>{\raggedright\arraybackslash}p{1.7cm}}
\toprule
No. & 主張 & 実験と測定されたもの & 出典 & 状態 \\
\midrule\endfirsthead
\toprule
No. & 主張 & 実験と測定されたもの & 出典 & 状態 \\
\midrule\endhead
1 & 高い閾値：痛みのセンサの加熱閾値は $41^\circ$ から $46$--$48^\circ$ & 単一線維記録。サルの皮膚：遅い（C）センサの閾値の中央値は $41^\circ$、速いセンサの第2型は $46^\circ$、第1型は $53^\circ$ 超。ヒトの皮膚：圧力にも応答するCセンサは $40.7^\circ$、圧力に応答しないものは $48^\circ$。 & \cite{Treede1995heat, Weidner1999cfib} & 測定済み \\
2 & 痛みのセンサは触覚のセンサよりはるかに順応しにくく、軽い損傷の後にはより敏感になる；強い加熱の後の低下は一つの型で測定 & 皮膚の軽いやけど（$50^\circ$、$100\,\text{s}$）：$5$--$10$~分後、ヒトの痛みの閾値とサルのCセンサの閾値が $2$--$6^\circ$ 下がる；他の皮膚センサにはこのずれはない。速いセンサの第2型では、強い加熱の後に応答が抑えられる。触覚センサとの順応の比較は一般的な推定で、ここでは個別の実験で裏づけられていない。 & \cite{LaMotte1982hyper, Treede1995heat} & 測定済み \\
3 & 2本の配線：速いほうは $5$--$30$~m/s、遅いほうは $0.5$--$2$~m/s；到着時間 $t=L/v$~\eqref{eq:paintime} & 伝導速度：サルの速い痛みのセンサは平均 $15$ と $25$~m/s（二つの型）；ヒトのCセンサは $0.8$--$1.0$~m/s。$L=1$~m ならそれぞれ数十ミリ秒と約1秒。 & \cite{Treede1995heat, Weidner1999cfib} & 測定済み \\
4 & 痛みは2回来る：最初は速く正確に、次に鈍く長く & ヒト前腕の加熱に対する反応時間：温度の上昇とともに $1100$ から $400\ms$；有毛部の皮膚の強い加熱は二重の痛みを生むが、手のひらでは生まない。第一の痛みを運べるのは $6$~m/s より速い線維だけで、潜時からそれに対応するのは手のひらにない速いセンサの第2型である。役割の分担（「どこ」と「どれほど悪いか」）は解釈。 & \cite{CampbellLaMotte1983first, Treede1995heat} & 測定済み \\
5 & ゲート：脊髄の出力ニューロンは痛みと触覚を受け、抑制性の介在ニューロンは触覚で興奮する（図~\ref{fig:gate}） & ゲートの構成は理論として提案された。マウス、遺伝的標識とニューロン群の除去：ソマトスタチンを持つ興奮性ニューロンは機械的な痛みに必要；ダイノルフィンを持つ抑制性ニューロンは触覚センサからの入力がそこへ届くのを防ぐ。これらのニューロンがないと軽い接触が痛みを起こす。 & \cite{MelzackWall1965, Duan2014} & 測定済み \\
6 & 触覚は痛みを弱める & ヒト、腕へのレーザー痛み刺激：同時の触覚は痛みの評定とパルスのエネルギーの識別能力を下げる；効果は、同じ皮膚分節の中で触れる点と痛む点の距離とともに弱まる。 & \cite{Mancini2014touch} & 測定済み \\
7 & ゲート理論の仮定：強い痛みは自ら抑制性介在ニューロンをオフにし、ゲートが開け放たれる & 章では元の理論の仮定と明記。マウスでの研究は、ゲートが触覚を痛みのチャネルに入れない仕組みを記述するが、痛みによる介在ニューロンのオフは示していない。 & \cite{MelzackWall1965} & 仮説 \\
8 & ゲート~\eqref{eq:gate}：閾値は触覚なしで $0.18$、触覚ありで $0.86$、$g=2$ で $0.11$；$\nu=1$ で触覚は出力を $1$ から $0.25$ に下げ、$\nu=2$ では効かない；触覚単独では通らない（図~\ref{fig:gatecurves}） & 本文の重みでの \texttt{models/\allowbreak pain\_\allowbreak gate.py} による計算が、これらの数値をすべて再現する。 & \texttt{models/\allowbreak pain\_\allowbreak gate.py} & モデル \\
9 & 脳幹からの下行路は、ゲートのゲインをどちらの向きにも変える & ラット延髄、尾の加熱中の記録：あるニューロンは引っ込めの前に発火し（「オン」）、別のニューロンは沈黙する（「オフ」）；この領域の弱い刺激は反射を抑える。総説：同じ回路が痛みの伝達を促進も抑制もし、オピオイドはそれを介して働く。 & \cite{Fields1983rvm, Fields2004} & 測定済み \\
10 & 楽になるという期待は、オピオイドのチャネルを通じて痛みを弱める & 急性の痛みを持つ人：楽になるという期待で痛みが減った人では、オピオイド受容体の遮断で痛みが他の人と同じレベルに戻った；それ以外の人では遮断は痛みを変えなかった。 & \cite{Levine1978} & 測定済み \\
11 & 脳は割り込みの損得によって警報の音量を選ぶ & 音量を選ぶ規則を測定した実験はない。 & --- & 仮説 \\
12 & 感作：損傷の後、ゲートの出力は増し、閾値は下がり、その一部は脊髄による；損傷刺激が止んだ後も続く & ラット：皮膚の損傷の後、屈曲反射の閾値が下がり、応答が増す；電気生理学的解析によれば、増大の一部は脊髄そのもので生じる。損傷刺激は刺激そのものより長く続く感度の変化を起こす。正確な回復時間は章では挙げていない。 & \cite{Woolf1983} & 測定済み \\
13 & 感作~\eqref{eq:sensitization}は正帰還である；ループが強いと、警報は治癒後もラッチしうる & モデル~\eqref{eq:sensitization}から導かれる（章末の問題）。治癒後も続く痛みがまさにこの形のラッチしたループであることを示す実験はない。 & 章内の導出 & 仮説 \\
14 & 痛みは負の報酬であり、それによる学習は時間差分誤差に従う & fMRI、ヒト、痛みの前触れの連鎖：腹側線条体と前部島皮質の活動が時間差分学習の信号と一致する。サル：\nt{DOPA}ニューロンの一部は不快なエアパフとその前触れで興奮し、別の一部は抑制される。 & \cite{Seymour2004, Matsumoto2009} & 測定済み \\
15 & 痛みはいまの作業に割り込む & ヒト、電気的な痛み刺激と音のプローブ：痛みの開始 $100\ms$ 後のプローブへの応答は明らかに遅れ、$1.5\,\text{s}$ 後には対照との差はもうない。総説はこうした実験を注意の奪取のモデルにまとめる。 & \cite{Crombez1996disrupt, EcclestonCrombez1999} & 測定済み \\
16 & 引っ込め反射は脊髄で閉じる & ラット：後角深層のニューロンは個々の筋肉の引っ込め反射の空間パターンを再現する；頸部からは逆行性に興奮させられず、脊髄の介在ニューロンである。 & \cite{Schouenborg1995wr} & 測定済み \\
\bottomrule
\end{longtable}}

\section{第\ref{sec:motorlearning}章　運動学習}

最小二乗則と専用の誤差配線の利点は定理である。誤差の配線を持つ回路の構成、一致による弱化、遅延へのトレースの調整、後効果と技能の自発的回復は測定されている。回路全体が教師あり学習として働き、内部モデルを作るというのは有力な仮説である。二過程モデルの数値は本書のモデルによる計算である。

{\footnotesize
\setlength{\tabcolsep}{3pt}\renewcommand{\arraystretch}{1.15}
\begin{longtable}{>{\raggedright}p{0.5cm}>{\raggedright}p{4.0cm}>{\raggedright}p{5.6cm}>{\raggedright}p{2.3cm}>{\raggedright\arraybackslash}p{1.7cm}}
\toprule
No. & 主張 & 実験と測定されたもの & 出典 & 状態 \\
\midrule\endfirsthead
\toprule
No. & 主張 & 実験と測定されたもの & 出典 & 状態 \\
\midrule\endhead
1 & 規則~\eqref{eq:lms}は平均すると二乗誤差の勾配をたどる & 適応フィルタ理論の結果：入力と出力の誤差に比例する補正は、平均二乗誤差の確率的勾配降下である。 & \cite{WidrowHoff1960} & 理論 \\
2 & 出力ごとに誤差の配線があれば速さは出力の数によらない；一つの共通の数では、ノイズを出すニューロンの数とともに落ちる（命題~\ref{prop:teachers}） & 線形ネットワークの学習曲線：ニューロンのランダムな摂動による学習は、摂動を受けるニューロンの数の倍だけ正確な勾配より遅い。 & \cite{Werfel2005} & 理論 \\
3 & 誤差の配線を持つ回路は教師あり学習として働く（命題~\ref{prop:errorwire}） & 回路の構造から導かれた理論。7--10行の実験はそれと整合する；回路全体がまさにそう働くことは証明されていない。 & \cite{Marr1969, Albus1971} & 仮説 \\
4 & 拡張層：約 $7\cdot10^{10}$ 個の小さな\nt{GLUT}ニューロン、脳の残り全部より多い & 溶解した組織の懸濁液での細胞核の計数：ヒトの小脳には脳全体の $86\cdot10^9$ のうち $69\cdot10^9$ のニューロン。立体解析学：高齢男性の小脳に $101\cdot10^9$ 個の小さなニューロン。 & \cite{Azevedo2009, Andersen1992cereb} & 測定済み \\
5 & 入力の次元を上げると、望みの関係が線形になる & 分割の数についての定理：$n$ 個の入力の重み付き和としきい値は、一般の位置にある約 $2n$ 個のベクトルの任意のラベル付けを実現する。小脳がまさにこれを使っているというのは3行目の仮説の一部。 & \cite{Cover1965} & 理論 \\
6 & 約 $3\cdot10^7$ 個の出力\nt{GABA}ニューロン、それぞれ $10^5$ 程度の入力 & ヒト小脳の立体解析学：出力ニューロンは $30.5\cdot10^6$ 個。電子顕微鏡、ラット：出力ニューロン1個あたり拡張層からの入力は約 $1.75\cdot10^5$。出力ニューロンの抑制性伝達物質は総説による。 & \cite{Andersen1992cereb, NapperHarvey1988, Ito2001} & 測定済み \\
7 & 出力ニューロンあたり誤差の配線はちょうど1本、毎秒1回ほど、毎回強力なバースト & 記録の総説：各出力ニューロンは1本の登上\nt{GLUT}入力を受け、それが $1$~Hz 程度の頻度で複雑スパイクを起こす。 & \cite{Ito2001} & 測定済み \\
8 & バーストの確率は運動誤差の方向に依存する & サル、小脳の眼球運動領域、サッカード適応：視覚誤差の方向が複雑スパイクの確率を、大きさがそのタイミングを決める；バーストは次の試行の単純スパイクを変え、細胞の好みの誤差の方向に運動をずらす。 & \cite{Herzfeld2018} & 測定済み \\
9 & 誤差のバーストと活動中の入力が一致すると、その入力が弱まる（$-\eta e_i x_j$） & 実験の総説：拡張層からの入力と誤差の配線を同時に刺激すると、その入力が長期的に弱まる；入力だけの刺激では弱まらない。 & \cite{Ito2001} & 測定済み \\
10 & トレースの長さは誤差の遅延に合わせてある：遅延 $100\ms$ では $100\ms$ 前に活動していた入力が最も強く弱まる & スライスと生きたマウス：眼球運動学習を担う領域では、弱化は入力と誤差の間隔約 $120\ms$（この課題で誤差信号が届くのにかかる時間）に鋭く同調している；別の領域では細胞ごとに異なる間隔に同調している。 & \cite{Suvrathan2016} & 測定済み \\
11 & 回路は体の内部モデルを学習する適応フィルタ~\eqref{eq:adaptivefilter}である & 回路の構造から導かれたモデル。学習したフィルタを体の伝達関数として測定した直接の実験はない。 & \cite{Fujita1982} & 仮説 \\
12 & 予測は自分の動きの結果を差し引くので、自分をくすぐることはできない & ヒト：自分で触れると、他人の同じ接触よりくすぐったくない；fMRIでは体性感覚皮質の自分の接触への応答が弱く、動きが皮膚に触れるとき小脳の活動が少ない。予測の差し引きによる説明は仮説。 & \cite{Blakemore1998} & 仮説 \\
13 & 外力のもとで外れは小さくなり、外力を取り除くと手は反対側に外れる & ヒトが力場の中でロボットのハンドルを持って手を伸ばす：軌道はまっすぐに戻る；力場を取り除いた後の後効果は最初のゆがみのほぼ鏡像で、訓練していない空間の部分にも転移する。 & \cite{ShadmehrMussaIvaldi1994} & 測定済み \\
14 & 二つの過程が学習する~\eqref{eq:tworate}；逆向きの外乱の後、技能はひとりでに戻る & ヒト、力場、次に短い逆向きの力場と誤差を固定した試行：補正はひとりでに古い値へ戻る；二過程モデルはこれを再現し、一過程モデルは再現しない。 & \cite{Smith2006} & 測定済み \\
15 & 図~\ref{fig:tworate}の数値：$200$ 回の動きで $0.84$（遅い過程 $0.63$）；逆向き $6$ 回；速い $-0.53$、遅い $0.46$；$40$ 回で $0.37$ まで回復 & \texttt{models/\allowbreak motor\_\allowbreak learning.py} による計算、$A_f=0.92$, $B_f=0.1$, $A_s=0.996$, $B_s=0.02$：$0.840$（$0.634$ と $0.205$）、$6$ 回、$-0.531$ と $0.457$、$40$ 回目でピーク $0.371$。 & \texttt{models/\allowbreak motor\_\allowbreak learning.py} & モデル \\
16 & 体が異なる速さで変わるので、二つの過程が必要になる & 理論：寿命の異なる複数の外乱のもとでの最適推定は、時定数の異なる複数のフィルタを与え、自発的回復と再学習の加速を説明する。 & \cite{Kording2007} & 仮説 \\
\bottomrule
\end{longtable}}

\section{第\ref{sec:chemoutput}章　化学出力}

心臓への2本の配線とその速さの違い、ホルモンカスケードの負帰還、分量の頻度による符号、概日発振器とその光による同期は直接測定されている。境界$K<8$は本章で導いた制御理論の定理である。カスケードのフィードバックそのものが脈動を生むというのは仮説であり、それを強く支持する実験がある。

{\footnotesize
\setlength{\tabcolsep}{3pt}\renewcommand{\arraystretch}{1.15}
\begin{longtable}{>{\raggedright}p{0.5cm}>{\raggedright}p{4.0cm}>{\raggedright}p{5.6cm}>{\raggedright}p{2.3cm}>{\raggedright\arraybackslash}p{1.7cm}}
\toprule
№ & 主張 & 実験と測定内容 & 出典 & 状態 \\
\midrule\endfirsthead
\toprule
№ & 主張 & 実験と測定内容 & 出典 & 状態 \\
\midrule\endhead
1 & 心臓のペースメーカーは単独で毎分約$100$回拍動する。安静時はブレーキが踏まれており、心拍数はふつう$60$--$70$程度 & ヒト、心臓への2本の配線をともに完全に遮断：固有心拍数は安静時心拍数より高く、加齢とともに下がり、成人で毎分約$100$拍。したがって安静時はブレーキが優位である。安静時心拍数$60$--$70$は典型値で、個別には引用していない。 & \cite{JoseCollison1970ihr} & 測定済み \\
2 & アクセル\nt{NORA}とブレーキ\nt{ACET}はローパスフィルタで、ブレーキの方が速い~\eqref{eq:gasbrake} & イヌ、心臓の迷走神経または交感神経を周波数変調刺激し、心房レートへの伝達関数を測定：どちらの経路もローパスフィルタで、交感神経側は遮断周波数がはるかに低く、約$1.7\,\text{s}$の純粋な遅れが加わる。特性は平均緊張度に依存するので、線形式~\eqref{eq:gasbrake}は近似である。 & \cite{Berger1989} & 測定済み \\
3 & ブレーキが速いのは、\nt{ACET}がセカンドメッセンジャーなしでカリウムチャネルを開き、\nt{NORA}は反応の連鎖を通じて作用するから & 心房細胞の膜パッチ：アセチルコリンで開くカリウムチャネルは、受容体と共役したGタンパク質の$\beta\gamma$サブユニットが、細胞内のセカンドメッセンジャーを介さずに開く。cAMPを介する\nt{NORA}の経路はここでは引用していない。 & \cite{Logothetis1987gbg} & 測定済み \\
4 & 血液はおよそ1分で全身を一周する。血中の\nt{ADRE}は受容体をもつすべての臓器に届く & 一般的な桁の見積もり。本章でもそのように述べており、出典は引用していない。 & --- & 推定 \\
5 & ホルモンカスケードは負帰還で囲まれている：最終段のホルモンが前段を抑える & 実験の総説：副腎皮質ホルモンは、下垂体からのACTH放出と視床下部からの放出ホルモン放出を、速い、中間、遅いの三つの時間域で抑制する。 & \cite{KellerWood1984fb} & 測定済み \\
6 & 同一の3段は$K<8$で安定~\eqref{eq:cascadestable}、境界での周期は$2\pi\tau/\sqrt3\approx3.6\tau$ & 本章で導出：周波数$\sqrt3/\tau$で3段は位相を$180^\circ$ずらし、$8$分の1に減衰させる。 & 本章で導出 & 理論 \\
7 & レベルの誤差は$1/(1+K)$で、この調節器はおよそ$10\,\%$より高い精度を保てない（命題~\ref{prop:cascade}） & 比例フィードバックに対する6行目の帰結。実際の軸は複数の段と複数の時間域にフィードバックをもつ（5行目）ので、この境界はモデル~\eqref{eq:cascade}についてのものであり、測定値ではない。 & 本章で導出 & 理論 \\
8 & $K=4$と$7$ではレベルが落ち着き、$K=12$では周期$3.7$--$3.9\,\tau$の規則正しい脈動（図~\ref{fig:cascade}） & \texttt{models/\allowbreak chem\_\allowbreak models.py}による計算：$K=12$で連続する周期は$3.9$、$3.7$、$3.8$、$3.7\,\tau$。$K=4$では計算終了時の振幅$0.003$。 & \texttt{models/\allowbreak chem\_\allowbreak models.py} & モデル \\
9 & ストレスホルモンは約1時間に1回のパルスで放出される。パルスはフィードバックそのものが生み、独立した発振器はない & ラット、放出ホルモンの持続注入：ACTHとコルチコステロンは生理的なほぼ1時間の頻度で振動する。パルスに視床下部からのリズム入力は必要ない。遅れのあるフィードバックをもつ下垂体--副腎モデルがこの振動を再現する。 & \cite{Walker2012glc, Walker2010} & 仮説 \\
10 & 受け手はレベルではなく分量の頻度を読む & 性腺刺激ホルモン放出ホルモンを自ら分泌できないサル：持続注入では下垂体ホルモンの放出は回復せず、1時間に1回の分量投与で回復する。持続注入に戻すと応答は再び消える。ハイパスフィルタとしての受容体の疲労は解釈である。 & \cite{Belchetz1978} & 測定済み \\
11 & 分量の頻度が違えば、同じ腺の異なる細胞への作用も違う & 同じサル：分量を1時間に$2$--$5$回に増やすと下垂体ホルモンが両方とも下がる。$3$時間に1回に減らすと一方（LH）が下がり他方（FSH）が上がるので、両者の比は頻度で決まる。 & \cite{Wildt1981gnrh} & 測定済み \\
12 & 概日リズムは数時間の遅れをもつ細胞内の負帰還が生む & 総説：時計遺伝子のタンパク質が遅れて自らの遺伝子の転写を抑える。転写と翻訳のループが約1日の周期を生む。 & \cite{TakahashiJS2017} & 測定済み \\
13 & 主発振器は数万個のニューロンからなる群で、体の他の部分を同期させる & 総説：視交叉上核が主概日発振器である。培養した個々のニューロンは独立した発振器で、ネットワークがそれらを同期させる。マウスでは片側に約$10^4$個のニューロン。 & \cite{Welsh2010scn} & 測定済み \\
14 & ヒトの固有周期は$24.18$時間 & 薄暗い光のもと、1日から切り離したスケジュールで暮らすヒト：メラトニン、体温、コルチゾールのリズムの周期は若年者と高齢者で平均$24.18$時間、ばらつきは小さい。 & \cite{Czeisler1999} & 測定済み \\
15 & 光は位相同期ループのように発振器を合わせる~\eqref{eq:pll}。必要な補正は1日約$11$~min & ヒト、1日のさまざまな時刻に$6.7$時間の明るい光を1回照射：振幅$5$時間のタイプ1位相反応曲線。体温最低点より前では遅れ、後では進む。式~\eqref{eq:pll}は標準モデル、$11$~min $=0.18$時間は算術である。 & \cite{Khalsa2003prc} & 測定済み \\
16 & 光がなければ同期はなく、リズムは固有周期で進む & 光をまったく知覚できない全盲者で、通常の社会生活を送る人：約半数でメラトニンとコルチゾールのリズムが1日と一致しない固有周期で進む。 & \cite{Sack1992blind} & 測定済み \\
\bottomrule
\end{longtable}}

\section{第~\ref{sec:debugging}~章　マシンのデバッグ}

電極に何が聞こえるか、活動の低次元性、硬い配列によるインターフェースの性能、脳とデコーダの共同学習、刺激による視覚上の点は、査読付きの実験で測定されている。データ量の見積もりは本章の算術である。ヒトに埋め込んだNeuralinkの装置については、確認できた限り、データを公表しているのは主に同社自身である。

{\footnotesize
\setlength{\tabcolsep}{3pt}\renewcommand{\arraystretch}{1.15}
\begin{longtable}{>{\raggedright}p{0.5cm}>{\raggedright}p{4.0cm}>{\raggedright}p{5.6cm}>{\raggedright}p{2.3cm}>{\raggedright\arraybackslash}p{1.7cm}}
\toprule
№ & 主張 & 実験と測定内容 & 出典 & 状態 \\
\midrule\endfirsthead
\toprule
№ & 主張 & 実験と測定内容 & 出典 & 状態 \\
\midrule\endhead
1 & 電極には半径約100マイクロメートル以内のニューロン、ふつう1個から数個のスパイクが聞こえ、波形で区別できる & ラットCA1野、細胞内と細胞外の同時記録：細胞外スパイクの波形は細胞内スパイクの幅と振幅を反映する。見積もりではテトロードは$60$--$100$個のニューロンを区別できるはずだが、実際に分離できるのは約6個。 & \cite{Henze2000extra} & 測定済み \\
2 & 脳には$8.6\cdot10^{10}$個のニューロンがあり、プローブが聞くのはおよそ1億分の1 & 溶解した組織の懸濁液中の核の計数：ヒトの脳のニューロンは$86\cdot10^9$個。割合は本章の算術。 & \cite{Azevedo2009} & 測定済み \\
3 & 運動皮質の数百個のニューロンの活動は10--20個の共通変数に収まる & サル運動皮質の記録の総説：集団活動は、多くのニューロンに共通する少数の潜在変数で記述できる。 & \cite{Gallego2017} & 測定済み \\
4 & 隣接ニューロンのノイズは相関しており、100個のニューロンは1000個とほぼ同じだけ運ぶ（命題~\ref{prop:pooling}） & サル視覚皮質：隣接ニューロンのノイズ相関係数は約$0.1$で、平均化の利得は100個のニューロンで飽和する。情報制限相関の理論。 & \cite{Zohary1994, MorenoBote2014} & 測定済み \\
5 & 生データ$2\cdot10^8$~bit/sに対しスパイクでは$8\cdot10^4$~bit/s~\eqref{eq:probedata} & スパイク列の容量~\eqref{eq:spikeentropy}による本章の算術。デジタル化のパラメータは実際のもの：Neuralinkのチップは$19.3$~kHz、チャネルあたり$10$~bit。 & 本章で導出；\cite{Rieke1997, Musk2019} & 理論 \\
6 & Neuralinkの最初のシステム：チップは増幅とデジタル化を行い、生データはケーブルで送られ、スパイクは外部のプログラムが検出する。チップ上の検出、密閉筐体、無線、ワイヤレス給電は、同社によればヒト用の装置のもの & 同社の論文：生データ全体をUSB-Cケーブルで送り、チップ外のプログラムがリアルタイムでスパイクを検出する。搭載した圧縮、ワイヤレス給電と伝送は臨床用装置の計画とされている。ヒトに埋め込んだ装置については同社の発表のみ。チップ上でのスパイク検出が必要だというのは、\eqref{eq:probedata}からの本章の結論。 & \cite{Musk2019}；Neuralinkの報告、査読付き論文なし & 測定済み（同社の論文）。ヒトについては同社の報告 \\
7 & $96$本のスレッドに$32$個ずつ最大$3072$個の電極。スレッドはロボットが血管を避けて挿入する & 同社の論文：$96$本のスレッドに$32$個ずつ$3072$個の電極。ロボットは毎分$6$本（$192$電極）を挿入する。配列を長期埋め込みしたラットで最大$70\,\%$の電極がスパイクを聞き取る。 & \cite{Musk2019} & 測定済み（同社の論文） \\
8 & ヒトでは$64$本のスレッドで約1000チャネル。一部のスレッドがずれた。カーソルは約$10$~bit/s & 同社の発表のみ。PubMedの検索ではヒトでの結果を示す査読付き論文は見つからなかった。本文中の但し書きと一致する。 & Neuralinkの報告。査読付き論文なし & 測定済み（同社の報告） \\
9 & 100電極の配列によるインターフェースがカーソルを操作する & 麻痺のある人、一次運動皮質に$96$電極：受傷3年後でも腕を動かす意図がスパイクを変える。デコーダは「ニューラルカーソル」（メール、テレビ）、義手と ロボットアームの操作を可能にする。 & \cite{Hochberg2006} & 測定済み \\
10 & 「想像上の手」による書字：毎分約$90$文字、$8$~bit/s未満 & 手が麻痺した人が手書きを試み、再帰型デコーダを使用：リアルタイムで毎分$90$文字、精度$94.1\,\%$、自動修正後は$99\,\%$超。ビットでの見積もりは本章の算術。 & \cite{Willett2021} & 測定済み \\
11 & 話そうとする試みが毎分約$60$語の速さでテキストに変換される & 明瞭な発話を失った人、皮質に配列：毎分$62$語。語彙$50$語で単語誤り率$9.1\,\%$、語彙$125\,000$語で$23.8\,\%$。 & \cite{Willett2023} & 測定済み \\
12 & インターフェースが取り出すのは容量のごく一部：1ニューロンで毎秒数十ビット、100個で数千ビット（命題~\ref{prop:probe}） & 上限~\eqref{eq:spikeentropy}は理論。8行目、10行目との比較は本章の結論。ボトルネックが配線ではなく低次元性と符号の無知であることは、直接の実験で確かめられていない。 & \cite{Rieke1997} & 理論 \\
13 & 脳とデコーダは互いから学ぶ & サルがカーソルを操作。デコーダは閉ループで適応し、ニューロンも同時に学習し直す。精度は上がり、技能は保持され、自分の腕の動きの影響を受けにくくなる。学習速度を分けるという助言は工学上の経験則であり、本章に個別の実験はない。 & \cite{Orsborn2014coad} & 測定済み \\
14 & 電極を流れる電流は、周囲のニューロンと配線を種類の区別なく興奮させる & マウスとネコの皮質、二光子カルシウムイメージング：弱い電流でも、直径数十マイクロメートルの体積内の軸索を直接興奮させることで、最大数ミリメートルに散らばるまばらなニューロンを興奮させる。つまり興奮は非選択的だが、一様ではない。 & \cite{Histed2009stim} & 測定済み \\
15 & 視覚皮質の刺激は点を灯す。同時に最大$16$点で一部の文字と物体の境界が認識できる & 盲目の人、視覚皮質に$96$電極を$6$か月：単一電極の閾値は$66.8\pm36.5$~\textmu A。電極群の同時刺激（最大$16$電極は刺激装置の上限）で閾値が下がり、区別できるパターンが生じ、一部の文字と物体の境界が認識される。 & \cite{Fernandez2021} & 測定済み \\
16 & Neuralinkの第二の方向は、視覚皮質に信号を入れる装置 & 開発についての同社の発表のみ。その動作に関する査読付きデータは見つからなかった。 & Neuralinkの報告 & 仮説 \\
17 & ブロードキャスト伝達物質の濃度は組織内の化学センサで測れる & ラット脳スライス、炭素繊維による高速サイクリックボルタンメトリー：セロトニンの放出と再取り込みを測定。物質はシナプスの外へ出る。 & \cite{Bunin1998} & 測定済み \\
\bottomrule
\end{longtable}}

\section{第\ref{sec:eetypes}章　素子としてのニューロン}

個々の細胞の動作モードは、電極からの電流注入と電圧記録で直接測定されている。積分器とクラス分けの数学は古典的な理論である。脳がモードの切り替えを計算に使っているというのは、依然として仮説である。

{\footnotesize
\setlength{\tabcolsep}{3pt}\renewcommand{\arraystretch}{1.15}
\begin{longtable}{>{\raggedright}p{0.5cm}>{\raggedright}p{4.0cm}>{\raggedright}p{5.6cm}>{\raggedright}p{2.3cm}>{\raggedright\arraybackslash}p{1.7cm}}
\toprule
№ & 主張 & 実験と測定内容 & 出典 & 状態 \\
\midrule\endfirsthead
\toprule
№ & 主張 & 実験と測定内容 & 出典 & 状態 \\
\midrule\endhead
1 & 共振器：電圧の変化に逆らう遅い電流が、ニューロンを数Hzから数十Hzに最大をもつバンドパスフィルタにする（第\ref{sec:resonator}節） & スライス中の\nt{GLUT}ニューロンに周波数が滑らかに上がる正弦波電流を注入し、インピーダンス$|Z(f)|$を測定：最大は$33^\circ$Cで$2$--$5$~Hz（$38^\circ$Cで約$7$~Hz）。共振は遅いカリウム電流と陽イオン電流が生み、持続性ナトリウム電流が強める。総説は多くの細胞クラスで同じ手法を示す & \cite{HuVervaekeStorm2002,HutcheonYarom2000} & 測定済み \\
2 & 遅い電流はインダクタンスのように振る舞い、膜容量と合わせて共振回路を作る & 電流に対する軸索の閾値下応答を測定し、線形化したコンダクタンス方程式で記述。等価回路は電圧に依存する「現象論的」インダクタンスを含む & \cite{Mauro1970} & 理論 \\
3 & フィードバックのある群の時定数は$\tau_{\text{eff}}=\tau/(1-w)$~\eqref{eq:perfectint}。$\tau=0.1$~s、$w=0.995$で$20$~s & 群の線形方程式から本章で導出。眼の位置を保持する機構として理論研究で提案された & 本章で導出；\cite{Seung1996} & 理論 \\
4 & 眼の位置の積分器は、個々の細胞の性質ではなくネットワークで入力を保持する & 眼の位置を保持する魚の神経核での細胞内記録：急速な回転のあと、ニューロンの発火率は階段状に変わって保たれる。1個の細胞への短い電流は一時的なずれしか起こさず、電圧の階段は閾値下でも保たれる。それを維持するのはシナプス入力である & \cite{Aksay2001} & 測定済み \\
5 & 漏れのない積分器には学習による絶え間ない調整が必要。調整が狂うと忘れるか飽和する & 眼の位置の$\pm$に比例する視覚フィードバックで魚を訓練：保持は不安定または漏れのあるものになり、ニューロンの時定数は1桁以上下がって$<1$~sになった。通常の視覚フィードバックで安定性は徐々に戻った & \cite{Major2004} & 測定済み \\
6 & 双安定ニューロン：短い入力が長いプラトーをオンにし、抑制がオフにする。この性質はセロトニンでオンになる（第\ref{sec:bistable}節） & ネコの筋肉を制御する\nt{ACET}ニューロンの細胞内記録：短いシナプス興奮でニューロンは長く続く発火に移り、短い抑制でリセットされる。脊髄ネコではセロトニン前駆体の投与後にこの状態が現れる & \cite{Hounsgaard1988} & 測定済み \\
7 & 二つのクラス：積分器型（発火率がゼロから上がる）と共振器型（閾値ですぐ有限の発火率） & クラスは分岐理論から導かれる。実験：皮質スライスで\nt{GLUT}ニューロンはいくらでも低い発火率で発火し始め（タイプ1）、速い\nt{GABA}ニューロンは有限の発火率から跳躍的に始まる（タイプ2） & \cite{Izhikevich2007,Tateno2004} & 測定済み \\
8 & 1個のニューロンが積分器にも同時検出器にもなる：抑制が加算の窓を短くする & 海馬スライスで、興奮性入力が閾値まで加算される窓は$2$~ms未満。興奮の直後に届くフィードフォワード抑制がこれを短くする。抑制の弱い樹状突起では窓が広い & \cite{Pouille2001} & 測定済み \\
9 & 軸索による遅延線（表~\ref{tab:eetypes}） & メンフクロウで、同時検出器へ向かう軸索の遅延を測定：軸索の長さの違いが異なる遅延を生み、検出器は音の到達時間差に同調している & \cite{Carr1990} & 測定済み \\
10 & 覚醒中はペースメーカー、睡眠中は黙った細胞 & \nt{SERO}ニューロンは日中ほぼ規則的に発火し、徐波睡眠で遅くなり、レム睡眠ではほぼ沈黙する（詳しくは第\ref{sec:sleep}章） & \cite{JacobsAzmitia1992,McGintyHarper1976} & 測定済み \\
11 & 素子を決めるのはイオンチャネルと、それを調整するブロードキャストチャネル（命題~\ref{prop:eetypes}） & 小さなネットワークで示された：調節物質がニューロン固有の性質とシナプスの強さを変え、同じ回路が異なるリズムを出す & \cite{Marder2012} & 測定済み \\
12 & 脳はモードの切り替えを計算に使っている（命題~\ref{prop:eetypes}） & 細胞のモードの切り替えが課題の解決に必要であることを示す直接の実験はない。調節物質が回路の出力を変える実験があるだけである & \cite{Marder2012} & 仮説 \\
\bottomrule
\end{longtable}}

\section{第\ref{sec:dendrites}章　樹状突起の数}

クラスの構造と入力の数は、解剖学的方法と電気的記録で測定されている。見積もり~\eqref{eq:dendriteload}はコンデンサの単純な物理である。入力数の計算上の説明は、理論とモデルに依拠している。

{\footnotesize
\setlength{\tabcolsep}{3pt}\renewcommand{\arraystretch}{1.15}
\begin{longtable}{>{\raggedright}p{0.5cm}>{\raggedright}p{4.0cm}>{\raggedright}p{5.6cm}>{\raggedright}p{2.3cm}>{\raggedright\arraybackslash}p{1.7cm}}
\toprule
№ & 主張 & 実験と測定内容 & 出典 & 状態 \\
\midrule\endfirsthead
\toprule
№ & 主張 & 実験と測定内容 & 出典 & 状態 \\
\midrule\endhead
1 & 膜の比容量は$c_m\approx1$~\textmu F/cm$^2$ & ニューロンの細胞体から膜をピペットに取り、電圧ステップを与えて容量電流の減衰から全容量を求め、面積で割った：3クラスのニューロンで$0.9$~\textmu F/cm$^2$ & \cite{Gentet2000} & 測定済み \\
2 & 充電時間$t\approx c_mA\Delta V/I$~\eqref{eq:dendriteload}：大きなニューロンには数十個の同時入力が必要で、小さなニューロンには大きなシナプス1個で足りる & コンデンサの充電法則による見積もり。数値（$20$と$300$~pF、$0.1$と数nA）は桁の値 & 本章で導出 & 理論 \\
3 & 1個の巨大シナプスは数nAの電流を流し、小さなニューロンをすぐに閾値に届かせる & スライスで終末と受け手を同時記録：シナプス1個の電流は$2$--$13$~nA、入力スパイクのたびに閾値を超える応答が生じ、遅れは$36^\circ$Cで約$0.4$~ms。受け手は\nt{GLYC}を出す & \cite{Borst1995,BorstSoria2012} & 測定済み \\
4 & 樹状突起ゼロ：触覚と痛みの感覚ニューロンでは、スパイクはセンサから細胞体を経ずに脊髄へ向かう & 構造（細胞体の近くで二つに分かれる軸索）は解剖学的に確立。伝導に細胞体の興奮性が要らないことは、このニューロンの検証済みモデルで示された：興奮性のない細胞体でも、スパイクは分岐点を同じく確実に通過する & \cite{AmirDevor2003} & 理論 \\
5 & 樹状突起1本：嗅覚ニューロンは1種類の受容体をもち、1本の入力線を伝える & 受容体遺伝子の実験の総説：各嗅覚ニューロンは大きな遺伝子ファミリーのうち一つの受容体遺伝子を発現する & \cite{Mombaerts2004} & 測定済み \\
6 & 樹状突起2本：音の方向検出器では、左耳と右耳からの入力が別々の樹状突起に接続する & 総説にまとめられた哺乳類の解剖と記録：両耳からの入力は反対向きの2本の樹状突起に分かれている & \cite{Grothe2010} & 測定済み \\
7 & 入力を2本の樹状突起に分けるとANDが厳密になる & モデルで示された：1本の樹状突起での飽和~\eqref{eq:shunt}により、片側からの入力は閾値に届かず、同時検出が改善する。入力の配置を変えた直接の実験はない & \cite{AgmonSnir1998} & 理論 \\
8 & 運動学習回路の小さな\nt{GLUT}ニューロンは約$7\cdot10^{10}$個で、それぞれ$\sim4$個の入力 & 等方性分画法による細胞計数：ヒトの脳のこの部分に約$690$億個のニューロン。生きたラットでの単一細胞記録：接触に対して少数の入力から離散的な電流のバーストが届き、それらが加算されると細胞が発火する & \cite{Azevedo2009,Chadderton2004} & 測定済み \\
9 & $n$入力の閾値素子が分けられるランダムなパターンは$P_{\max}\approx2n$個まで~\eqref{eq:cover}。運動学習回路の出力ニューロンでは上限$\sim2\cdot10^5$、現実的な見積もりは$\sim4\cdot10^4$個の対応 & 上限は組合せ幾何学の古典的な結果。現実的な見積もりは、正の重みのみとノイズへの余裕を仮定した容量理論を、このニューロンの測定された入力数に適用したもの & \cite{Cover1965,Brunel2004} & 理論 \\
10 & 入力の少ないミキサは入力の拡張に必要で、「4」は最適に近い & 理論：ミキサ層が与える表現の次元は、解剖学的に観察される入力数付近で最大になる。拡張の元の仮説は古典的な研究にある & \cite{Marr1969,Albus1971,LitwinKumar2017} & 仮説 \\
11 & 大きな木：$10^4$--$10^5$個の入力 & 電子顕微鏡像によるシナプスの計数：ラットの運動学習回路の出力\nt{GABA}ニューロン1個に$\approx1.7\cdot10^5$個のシナプス。海馬の\nt{GLUT}ニューロンに約$30\,000$個の興奮性入力と$1\,700$個の抑制性入力 & \cite{NapperHarvey1988,Megias2001} & 測定済み \\
12 & 大きな木の枝は、それぞれ閾値をもつ別々のサブ回路であり、ニューロンは小さな多層ネットワークである & 細い樹状突起の2か所を局所刺激：同じ枝の入力はシグモイド状に、異なる枝の入力は線形に加算される。ヒト皮質ニューロンの樹状突起で段階的なカルシウムスパイクが見つかり、1個の細胞が線形分離不可能な課題を解ける。モデル：1個のニューロンを再現するには深いネットワークが必要 & \cite{Polsky2004,Gidon2020,Beniaguev2021} & 測定済み \\
13 & 出力を兼ねる樹状突起：網膜の\nt{GABA}ニューロンの各枝は独立した方向検出器 & 個々の枝での二光子カルシウム記録：枝の信号は細胞体から先端への動きに選択的だが、細胞体の電圧はそうではない & \cite{Euler2002} & 測定済み \\
14 & 入力数は課題に従う（命題~\ref{prop:dendrites}） & クラスの構造は測定済み（3--13行目）。入力数が速さや分類のために選ばれたことは、個々のクラスについて理論とモデルで示されているだけである & \cite{LitwinKumar2017,AgmonSnir1998} & 仮説 \\
\bottomrule
\end{longtable}}

\section{第~\ref{sec:sleep}~章　睡眠}

睡眠のモード、リプレイ、シナプスの縮小は、ニューロンの記録、マイクロダイアリシス、電子顕微鏡で測定されている。睡眠のスケジュールと遅い振動は本書のモデルで記述され、睡眠の目的は主に仮説である。

{\footnotesize
\setlength{\tabcolsep}{3pt}\renewcommand{\arraystretch}{1.15}
\begin{longtable}{>{\raggedright}p{0.5cm}>{\raggedright}p{4.0cm}>{\raggedright}p{5.6cm}>{\raggedright}p{2.3cm}>{\raggedright\arraybackslash}p{1.7cm}}
\toprule
№ & 主張 & 実験と測定内容 & 出典 & 状態 \\
\midrule\endfirsthead
\toprule
№ & 主張 & 実験と測定内容 & 出典 & 状態 \\
\midrule\endhead
1 & 睡眠中も皮質のニューロンは沈黙しない。変わるのは動作モードである & ラット皮質ニューロンの長時間記録：ノンレム睡眠でも発火率はゼロにならず、活動期と静止期が交互に現れる。長い覚醒のあとはすべての状態で発火率が高く、睡眠のあとは低い & \cite{Vyazovskiy2009} & 測定済み \\
2 & \nt{ACET}は日中とレム睡眠で高く、ノンレム睡眠で低い（表~\ref{tab:sleepmodes}） & 自由に動くネコの皮質でのマイクロダイアリシス：覚醒時のアセチルコリン放出はノンレム睡眠より$100$--$175\%$高く、レム睡眠ではおよそ安静覚醒のレベル & \cite{Marrosu1995} & 測定済み \\
3 & \nt{NORA}と\nt{SERO}はノンレム睡眠で静まり、レム睡眠ではほぼ沈黙する & ラットの\nt{NORA}ニューロンとネコの\nt{SERO}ニューロンを睡眠段階ごとに単一記録：発火率は覚醒時に最大、ノンレム睡眠で低く、レム睡眠ではほぼゼロ & \cite{AstonJonesBloom1981,McGintyHarper1976} & 測定済み \\
4 & レム睡眠では筋肉が\nt{GABA}と\nt{GLYC}による抑制で切り離される & ラットで咀嚼筋の\nt{ACET}ニューロンの活動を測り、そこへ直接遮断薬を与えた：麻痺が解けるのは、両タイプの\nt{GABA}受容体と\nt{GLYC}受容体をすべて遮断した場合だけ & \cite{BrooksPeever2012} & 測定済み \\
5 & 睡眠圧$S$は二つの閾値をもつコンデンサ~\eqref{eq:sleeppressure}、$\tau_r=18.2$~h、$\tau_d=4.2$~h & 2過程モデル。時定数は脳波の徐波パワーが覚醒とともに増え睡眠中に減る様子から、閾値の形は自発的な覚醒の時刻から得られた & \cite{Borbely1982,Daan1984} & 理論 \\
6 & マシンはおのずと$16.1$~hの覚醒と$8.0$~hの睡眠に落ち着く（図~\ref{fig:twoprocess}） & 揺れる閾値を用いた\eqref{eq:sleeppressure}による計算、スクリプト\texttt{models/sleep\_models.py}：覚醒$16.1$~h、睡眠$7.9$--$8.0$~h & --- & モデル \\
7 & $40$~h眠らないと$S=0.90$になるが、睡眠は$8.8$~hしか続かない。借りは長さではなく深さで返される & モデル：\texttt{models/sleep\_models.py}は$S=0.90$と$8.8$~hを与える。実験：$40.5$~hの断眠後、8人の被験者で回復第1夜の脳波の徐波パワーが通常より有意に高い & \cite{Borbely1981} & モデル \\
8 & 物理的には$S$はアデノシンである & ネコでのマイクロダイアリシス：覚醒を制御する領域の一つで細胞外アデノシンが覚醒中に増え、断眠でさらに増え、睡眠中に減る。$S$がアデノシンだけであることは示されていない & \cite{PorkkaHeiskanen1997,Dunwiddie2001} & 仮説 \\
9 & ノンレム睡眠で皮質は振動する：数百msの活動と静止、1秒に1回未満（$<1$~Hz、麻酔下では多くが$0.3$--$0.4$~Hz） & 麻酔下のネコでの細胞内記録：$0.8$--$1.5$~sの脱分極と長い過分極、リズム$<1$~Hz、最多は$0.3$--$0.4$~Hz。同じ活動と静止の交代は自然睡眠でも見られる（1行目） & \cite{Steriade1993,Vyazovskiy2009} & 測定済み \\
10 & 振動はフィードバックと疲労をもつ群~\eqref{eq:updown}。$b=5$で周期$1.1$~s（モデルの値）、活動は時間の約$35\%$。$b=1$では一様な活動（図~\ref{fig:updown}） & 計算、スクリプト\texttt{models/sleep\_models.py}：周期$1.11$~s、活動時間の割合$0.35$（整数個の周期での平均）。$b=1$では活動の割合$1.0$。モデルの周期は測定値（9行目）より短い & --- & モデル \\
11 & \nt{ACET}は振動を止め、皮質を一様な活動に移す & ネコで\nt{ACET}ニューロンの核を刺激すると、皮質ニューロンの$79\%$で遅いリズムが遮断され、主に長い過分極が消えることで一様な活動に移った。アセチルコリンは疲労を生む遅いカリウム電流を抑える & \cite{SteriadeAmzica1993,McCormick1992} & 測定済み \\
12 & $7$--$15$~Hzのリズムのバーストは、皮質と中継核の間の\nt{GLUT}--\nt{GABA}ループが生む & フェレットの視覚中継核のスライスで、バーストは、中継細胞自身が興奮させる隣接\nt{GABA}核からの抑制に対して中継細胞が返すバーストによって生じる & \cite{vonKrosigk1993,SteriadeMcCormickSejnowski1993} & 測定済み \\
13 & 1日のリプレイは早送りで行われる：海馬ではおよそ$20$倍、皮質では$6$--$7$倍速い & ノンレム睡眠で、海馬ニューロンの繰り返す系列が時間的に圧縮されて現れる。トラックを走ったあと、場所ニューロンの系列が同じ順序で$\sim100$~msの短いバーストとして、およそ$20$倍に圧縮されて再生された。皮質での圧縮は$6$--$7$倍 & \cite{Nadasdy1999,LeeWilson2002,Euston2007} & 測定済み \\
14 & リプレイと皮質の協調を強めると記憶が良くなる（因果関係） & ラットで学習後、リプレイに同期した刺激によって皮質の徐波とリズムのバーストとの結びつきを強めた：翌日の記憶は高く、対照群では偶然レベル & \cite{Maingret2016} & 測定済み \\
15 & リプレイの目的は遅い記憶のインターリーブ学習 & 二つの学習系の理論と人工ネットワークでの計算：逐次学習は古いものを壊し、インターリーブ学習は壊さない。リプレイがまさにこのために必要だという脳での直接の実験はない & \cite{McClelland1995} & 仮説 \\
16 & 睡眠後、シナプスは大きさに比例して縮み、大きなものはほとんど変わらない & マウス皮質の$6920$個のシナプスの3次元電子顕微鏡観察：接触面積は睡眠後に覚醒後より$\sim18\%$小さく、縮小は大きさに比例し、大きく安定したシナプスは変わらない & \cite{deVivo2017} & 測定済み \\
17 & 睡眠の主な役割は重みの再正規化 & シナプス恒常性仮説。1行目と16行目はこれと整合するが、それが主な機能であることは証明しない & \cite{TononiCirelli2014} & 仮説 \\
18 & レム睡眠は世界モデルのベンチ試験 & モード（表~\ref{tab:sleepmodes}）は測定済み。ネットワークが世界モデルを試運転しているというのは理論上の推測で、実験はない & \cite{Hobson2009} & 仮説 \\
19 & 睡眠中の脳の洗浄：未解決の問題 & ある実験：睡眠中は細胞間隙が$60\%$広く、除去が速い。測定方法の異なる別の実験：睡眠中と麻酔下では除去が明らかに遅い。結果は互いに矛盾する & \cite{Xie2013,Miao2024} & 仮説 \\
20 & 局所睡眠：覚醒している動物の皮質の部位が短時間静止に落ちる & 長い覚醒のあとのラットで、皮質の一部位のニューロンがノンレム睡眠のように短く沈黙し、局所脳波に徐波が現れる。その瞬間、ラットは課題で誤りやすい & \cite{Vyazovskiy2011} & 測定済み \\
\bottomrule
\end{longtable}}

\section{第\ref{sec:livingmachine}章　生きたニューロンでできたマシン}

電極アレイ上の培養の振る舞いは、記録と刺激を伴う直接の実験で測定されている。一斉発火の機構はシナプス枯渇のモデルと微小培養での実験に基づき、シャーレの中のドーパミンに関する主張は単発の実験に基づく。そのうちいくつかは著者ら自身がデモンストレーションにすぎないとみなしている。

{\footnotesize
\setlength{\tabcolsep}{3pt}\renewcommand{\arraystretch}{1.15}
\begin{longtable}{>{\raggedright}p{0.5cm}>{\raggedright}p{4.0cm}>{\raggedright}p{5.6cm}>{\raggedright}p{2.3cm}>{\raggedright\arraybackslash}p{1.7cm}}
\toprule
No. & 主張 & 実験と測定されたもの & 出典 & 状態 \\
\midrule\endfirsthead
\toprule
No. & 主張 & 実験と測定されたもの & 出典 & 状態 \\
\midrule\endhead
1 & チップ上の培養：大部分は\nt{GLUT}、少数の\nt{GABA}の割合は標本による。海馬の培養では約 $6\%$ & 電極アレイ上の数千ニューロンのネットワークは標準的な実験である（行~3）。海馬の培養での\nt{GABA}ニューロンの割合は GABA 染色で数えられ、ニューロンの約 $6\%$ だった。大脳皮質の培養については検証された数値を挙げない & \cite{Benson1994} & 測定済み \\
2 & オルガノイド：ヒトの幹細胞からできた、大きさ約1ミリの三次元神経組織 & 幹細胞から、複数の脳領域をもつ三次元オルガノイドが育てられた。その中には前駆細胞の層と成熟ニューロンをもつ大脳皮質も含まれる & \cite{Lancaster2013} & 測定済み \\
3 & 入力がないと培養は、培養によって1秒から数分の間隔で一斉発火する & 密度 $3\,000$ から $50\,000$ ニューロンの大脳皮質培養 $58$ 個を5週間記録した（30分の記録 $963$ 本）。2週間後にはほぼすべての培養でネットワーク全体の一斉発火が優勢になる。一斉発火の間隔は $1$ から $300$~s で、標本に強く依存する & \cite{Wagenaar2006} & 測定済み \\
4 & 一斉発火は伝達物質の枯渇で終わり、間隔は $T_{\text{burst}}\approx\tau_{\text{rec}}\ln\frac{1}{1-x^*}$~\eqref{eq:burstinterval}。$2$--$4$~s は $\tau_{\text{rec}}=0.8$~s のモデルによる推定で、測定された回復はもっと遅い & 式は本章で導いた。モデル：枯渇するシナプスをもつネットワークは、自らネットワーク全体の一斉発火を生む。$4$--$30$ ニューロンの微小培養での実験：一斉発火の持続時間は前回からの時間に依存し、伝達物質の小胞を枯渇させると一斉発火は消える。著者らの結論は、一斉発火は伝達物質の蓄えで制限されるというもの。そこでの蓄えの回復は数十秒で、同じ式で数十秒から数分の間隔を与える & 本章の導出、\cite{Tsodyks2000,CohenSegal2011,Tsodyks1997} & 理論 \\
5 & 「黙る教師」：目的の応答が出たら刺激を止め、応答が定着する & 培養を $0.3$--$1$~Hz で刺激し、刺激後 $50\pm10$~ms に選んだ応答が現れたら刺激を $5$~min 止めた。数サイクル後、目的の応答はすぐに誘発されるようになった。学習曲線が描かれた & \cite{Shahaf2001} & 測定済み \\
6 & 培養がテニスをする。セッション内のラリーの有意な伸びは完全なフィードバックのときだけ & 詳細は第\ref{sec:dishbrain}章とその補遺。刺激の代わりに無音にしても、セッションの終わりにはフィードバックなしより上手になるが、効果は小さい。セッション間の学習は不安定 & \cite{Kagan2022} & 測定済み \\
7 & 培養が自分の入力を予測するように学習するというのは仮説。「知性」をめぐる大げさな言葉には反論がある & 自由エネルギー原理は理論的な提案であり、より単純な説明と区別する直接の実験はない。30人の研究者が反論論文で、ループ内での振る舞いの変化は知性も感覚も示さないと指摘した & \cite{Friston2010,Balci2023} & 仮説 \\
8 & 刺激パターンを選べば、培養は仮想の体を目標へ動かす & プログラムが現在の成績に応じて訓練刺激のパターンを自動で選んだ。数十分でネットワークは指定した状態に達し、$2$~h の訓練後、可塑性は $80$~min 続いた。振る舞いと無関係に与えた同じ刺激は効果がなかった & \cite{Bakkum2008} & 測定済み \\
9 & リザバー：学習するのは線形読み出し $y=W_{\text{out}}\mathbf r$~\eqref{eq:reservoir} だけ & リザバー計算の理論：減衰する記憶をもつ任意の非線形システムに、学習させた線形出力を加える & \cite{Maass2002,JaegerHaas2004} & 理論 \\
10 & リザバーとしてのオルガノイドは話者を約 $78\%$ の精度で区別する（著者らのデータ）。誤差で学習するのは読み出しで、オルガノイドの結合は教師なしで変わる & 複数の話者の音声を、電極アレイ上の電流パターンとしてオルガノイドに与えた。学習させた読み出しは約 $78\%$ の精度で話者を区別した。これは著者らの報告である。著者らはまた、訓練中にオルガノイドの機能的結合が変化したと報告し、これをオルガノイド内の教師なし学習と呼んでいる & \cite{Cai2023} & 測定済み \\
11 & 100万個のニューロンがスパイクに使う電力は約 $10^{-4}$~W~\eqref{eq:dishpower} & 推定：大脳皮質のエネルギー収支\eqref{eq:energy}から得たスパイクのコスト $10^{-10}$~J にスパイク数を掛けた。培養の電力の直接測定はない & 本章の導出、\cite{Attwell2001} & 理論 \\
12 & 光のフラッシュで放出するドーパミン：ケージドドーパミン $1$~mM、$240$ 回の繰り返しで誘発スパイク数が増加 & オルガノイドの遠隔プラットフォームで、刺激がより多くのスパイクを誘発したとき、光感受性ケージに入ったドーパミン（$1$~mM）を紫外線（$365$~nm、$0.8$~s）で解放した。$240$ ステップ（約 $5$~h）でスパイク数は全体として増えた。著者らは、実験一つではドーパミンとの関係は確立できないと書いている。対照はない & \cite{Jordan2024} & 仮説 \\
13 & 生きた細胞からのドーパミンが有機トランジスタの重みを長期かつ可逆的に変える & ドーパミンを分泌する細胞を有機ニューロモルフィック素子の上で育てた。電極でのドーパミンの酸化がチャネルのコンダクタンスを変える。重みの長期変化とその復帰が示された & \cite{Keene2020} & 測定済み \\
14 & 働く\nt{DOPA}ニューロンをもつオルガノイドは存在する。そのドーパミンを教師として使った例はない & ヒトの幹細胞からのオルガノイドに、電気的に活動する成熟\nt{DOPA}ニューロンとドーパミン産生が見つかった。このドーパミンが別のネットワークを教える実験は文献に見つからなかった & \cite{Jo2016} & 測定済み \\
15 & オルガノイドが棒を立てる。プログラムによる学習刺激はランダムより良いが定着せず、\nt{GLUT}シナプスを介する & マウス大脳皮質のオルガノイドを「台車の上の棒」課題と閉ループでつないだ。大多数のオルガノイドで、強化学習が選んだ信号はランダムな信号や信号なしより良い成績をもたらした。$45$~min の休止後、改善は残らなかった。AMPA 受容体と NMDA 受容体の遮断で改善は消えた & \cite{Robbins2026} & 測定済み \\
16 & 制御レジスタがなければ、テストベンチ上のネットワークは何を成功とみなすかを自分で決められない（命題\ref{prop:livingmachine}） & 行 5--15 のすべての実験で、学習信号は実験者かプログラムが与える。培養が自ら成功の基準を作り出した実験はない。これは不在からの推論であり、実験ではない & --- & 仮説 \\
\bottomrule
\end{longtable}}

\section{第\ref{sec:dishbrain}章　DishBrain}

テストベンチの構成とゲームの結果は、一つの実験研究とその続編から採った。ノイズと良い設定への漂流の式は本章で導き、本書のモデルで確かめた。シャーレの中の学習の仕組みは確定していない。

{\footnotesize
\setlength{\tabcolsep}{3pt}\renewcommand{\arraystretch}{1.15}
\begin{longtable}{>{\raggedright}p{0.5cm}>{\raggedright}p{4.0cm}>{\raggedright}p{5.6cm}>{\raggedright}p{2.3cm}>{\raggedright\arraybackslash}p{1.7cm}}
\toprule
No. & 主張 & 実験と測定されたもの & 出典 & 状態 \\
\midrule\endfirsthead
\toprule
No. & 主張 & 実験と測定されたもの & 出典 & 状態 \\
\midrule\endhead
1 & テストベンチ：チップあたりマウス大脳皮質の細胞約 $800\,000$ 個またはヒト細胞約 $10^6$ 個。$8$~mm$^2$ に電極 $26\,000$ 本、記録は最大 $1024$、刺激は $8$ 本 & 論文の方法：MaxOne チップ、$8$~mm$^2$ に白金電極 $26\,000$ 本、$1024$ チャネルを毎秒 $20\,000$ サンプルで記録。技術的には最大 $32$ 電極を刺激できるが、独立に刺激できたのは $8$ 本。マウス培養はおよそ $10$ 日目から使用 & \cite{Kagan2022} & 測定済み \\
2 & $10$~ms 周期のループ：運動領域のスパイクがラケットを動かす（図\ref{fig:dishloop}） & スパイクは $10$~ms の窓ごとに数えられる。スパイクから応答刺激までの遅れは約 $5$~ms で、装置のバッファで決まる & \cite{Kagan2022} & 測定済み \\
3 & 入力：場所による符号（$8$ 本のどの電極か）と頻度 $4$--$40$~Hz & 主実験では刺激頻度が、ボールが遠い壁にあるときの $4$~Hz からラケットのそばの $40$~Hz まで線形に増える。ボールのパルスの振幅は $75$~mV、パルスは矩形の二相性 & \cite{Kagan2022} & 測定済み \\
4 & 出力 $\Delta p=c\,(n_\uparrow-n_\downarrow)$~\eqref{eq:dishreadout}、窓あたりのカウントのノイズは $1/\sqrt{RT}$ & 二つの領域の差をとる規則は論文に記述がある（活動に応じた領域の重みつき）が、正確な式はない。ノイズの推定はポアソン計数からの帰結で、本章で導いた & \cite{Kagan2022}、本章の導出 & 理論 \\
5 & 教師：ヒットの後に予測可能な刺激 $75$~mV、$100$~Hz、$0.1$~s。ミスの後に予測不能な刺激 $150$~mV、$5$~Hz、$4$~s をランダムな場所と時刻に、その後 $4$~s 休止 & 方法：ヒットの後、$75$~mV、$100$~Hz、$100$~ms を $8$ 本の電極すべてに同時に。ミスの後、$150$~mV、$5$~Hz をランダムな電極にランダムな時刻で $4$~s、その後 $4$~s 刺激なし。培養にドーパミンはない & \cite{Kagan2022} & 測定済み \\
6 & ネットワークは自分の入力が予測可能になるように行動する & 自由エネルギー原理は理論的な提案で、著者らはそこからプロトコルを導いた。実験はこれと矛盾しないが、より単純な仕組み（行~7--8）と区別できない & \cite{Friston2010,Kagan2022} & 仮説 \\
7 & 失敗がネットワークを揺さぶり成功が揺さぶらないなら、設定で過ごす時間は $\rho\propto1/P_{\text{miss}}$~\eqref{eq:dishdensity}（命題\ref{prop:dishscramble}） & 対称な揺さぶりをもつマルコフ連鎖の流れの釣り合いから従う。本章で導いた。培養での直接の検証はない & 本章の導出 & 理論 \\
8 & 「ミスでかき混ぜる」モデルは学習する：ヒットの割合 $0.21\to0.38$。毎ラリー後に揺さぶると $\approx0.12$、揺さぶりなしで $0.19$（図\ref{fig:dishmodel}） & 計算、スクリプト \texttt{models/dishbrain\_model.py}、モデル培養 $40$ 個に各 $2000$ ラリー：$0.21\to0.38$、$0.13\to0.11$、$0.19\to0.19$ & --- & モデル \\
9 & ラリーが長くなるのは完全なループの生きた培養だけ。対照は学習しない & $20$~min のセッション $399$ 回：マウス培養 $9$ 個（$101$ セッション）、ヒト $11$ 個（$138$）、培地のみのチップ $6$ 枚（$80$）、安静の培養 $20$ 個（$42$）、シミュレーション $3$ 個（$38$）。最後の $15$~min の平均ラリー長が最初の $5$~min より長いのは生きた培養だけ。ニューロンでない細胞（HEK293T）と培地のみのチップでは効果なし & \cite{Kagan2022} & 測定済み \\
10 & セッション内の有意な伸びは完全なフィードバックのときだけ。刺激の代わりに無音でも、セッションの終わりにはフィードバックなしより上手だが、効果は小さい & $3$ 日で $486$ セッション：「刺激」（$n=27$）、「無音」（$17$）、「フィードバックなし」（$15$）。時間による伸びが有意なのは「刺激」条件だけ。セッションの終わりに「無音」条件は「フィードバックなし」を小さな効果で上回った & \cite{Kagan2022} & 測定済み \\
11 & 効果は小さい。ネットワークが長期的に学習するかは未解決 & セッション内の改善は確かだが、著者ら自身によれば、数日にわたるセッション間の学習は安定しては観察されなかった & \cite{Kagan2022} & 測定済み \\
12 & プレイ中、ネットワークは臨界状態に近づき、近いほどプレイが良い & 同じ培養の活動の雪崩の解析：構造化された入力はネットワークを臨界状態に近づけ、プレイの質はそこへの近さと相関する。ただし、行動の結果についての情報がなければ、臨界性だけでは学習に足りない & \cite{Habibollahi2023} & 測定済み \\
13 & 培養の「知性」は示されていない & 30人の研究者による反論論文：閉ループでの振る舞いの変化は知性も感覚も証明しない。それを示す実験はない & \cite{Balci2023} & 仮説 \\
14 & 提案する四つ目の対照：ミスの後に同じ長さとエネルギーの予測可能な刺激（\ref{sec:dishspec}節） & この実験は行われていない。本書の提案である & --- & 仮説 \\
\bottomrule
\end{longtable}}


\begin{thebibliography}{99}
\bibitem{Azevedo2009} F. A. C. Azevedo et al., \emph{Equal numbers of neuronal and nonneuronal cells make the human brain an isometrically scaled-up primate brain}, J. Comp. Neurol. \textbf{513}, 532 (2009).
\bibitem{Schultz1997} W. Schultz, P. Dayan, and P. R. Montague, \emph{A neural substrate of prediction and reward}, Science \textbf{275}, 1593 (1997).
\bibitem{SuttonBarto2018} R. S. Sutton and A. G. Barto, \emph{Reinforcement Learning: An Introduction}, 2nd ed. (MIT Press, Cambridge, MA, 2018).
\bibitem{Doya2002} K. Doya, \emph{Metalearning and neuromodulation}, Neural Networks \textbf{15}, 495 (2002).
\bibitem{Strata1999} P. Strata and R. Harvey, \emph{Dale's principle}, Brain Res. Bull. \textbf{50}, 349 (1999).
\bibitem{vanVreeswijk1996} C. van Vreeswijk and H. Sompolinsky, \emph{Chaos in neuronal networks with balanced excitatory and inhibitory activity}, Science \textbf{274}, 1724 (1996).
\bibitem{AstonJones2005} G. Aston-Jones and J. D. Cohen, \emph{An integrative theory of locus coeruleus-norepinephrine function: adaptive gain and optimal performance}, Annu. Rev. Neurosci. \textbf{28}, 403 (2005).
\bibitem{Beaulieu2011} J.-M. Beaulieu and R. R. Gainetdinov, \emph{The physiology, signaling, and pharmacology of dopamine receptors}, Pharmacol. Rev. \textbf{63}, 182 (2011).
\bibitem{Sparso2001} J. Spars{\o} and S. Furber (eds.), \emph{Principles of Asynchronous Circuit Design: A Systems Perspective} (Kluwer, Boston, 2001).
\bibitem{Carr1990} C. E. Carr and M. Konishi, \emph{A circuit for detection of interaural time differences in the brain stem of the barn owl}, J. Neurosci. \textbf{10}, 3227 (1990).
\bibitem{Wang2001} X.-J. Wang, \emph{Synaptic reverberation underlying mnemonic persistent activity}, Trends Neurosci. \textbf{24}, 455 (2001).
\bibitem{Hahnloser2002} R. H. R. Hahnloser, A. A. Kozhevnikov, and M. S. Fee, \emph{An ultra-sparse code underlies the generation of neural sequences in a songbird}, Nature \textbf{419}, 65 (2002).
\bibitem{Barlow1965} H. B. Barlow and W. R. Levick, \emph{The mechanism of directionally selective units in rabbit's retina}, J. Physiol. \textbf{178}, 477 (1965).
\bibitem{Carandini2012} M. Carandini and D. J. Heeger, \emph{Normalization as a canonical neural computation}, Nat. Rev. Neurosci. \textbf{13}, 51 (2012).
\bibitem{Pouille2001} F. Pouille and M. Scanziani, \emph{Enforcement of temporal fidelity in pyramidal cells by somatic feed-forward inhibition}, Science \textbf{293}, 1159 (2001).
\bibitem{Tsodyks1997isn} M. V. Tsodyks, W. E. Skaggs, T. J. Sejnowski, and B. L. McNaughton, \emph{Paradoxical effects of external modulation of inhibitory interneurons}, J. Neurosci. \textbf{17}, 4382 (1997).
\bibitem{Buzsaki2012} G. Buzs\'aki and X.-J. Wang, \emph{Mechanisms of gamma oscillations}, Annu. Rev. Neurosci. \textbf{35}, 203 (2012).
\bibitem{Rieke1997} F. Rieke, D. Warland, R. de Ruyter van Steveninck, and W. Bialek, \emph{Spikes: Exploring the Neural Code} (MIT Press, Cambridge, MA, 1997).
\bibitem{Mainen1995} Z. F. Mainen and T. J. Sejnowski, \emph{Reliability of spike timing in neocortical neurons}, Science \textbf{268}, 1503 (1995).
\bibitem{Softky1993} W. R. Softky and C. Koch, \emph{The highly irregular firing of cortical cells is inconsistent with temporal integration of random EPSPs}, J. Neurosci. \textbf{13}, 334 (1993).
\bibitem{Thorpe1996} S. Thorpe, D. Fize, and C. Marlot, \emph{Speed of processing in the human visual system}, Nature \textbf{381}, 520 (1996).
\bibitem{Zohary1994} E. Zohary, M. N. Shadlen, and W. T. Newsome, \emph{Correlated neuronal discharge rate and its implications for psychophysical performance}, Nature \textbf{370}, 140 (1994).
\bibitem{MorenoBote2014} R. Moreno-Bote et al., \emph{Information-limiting correlations}, Nat. Neurosci. \textbf{17}, 1410 (2014).
\bibitem{Gollisch2008} T. Gollisch and M. Meister, \emph{Rapid neural coding in the retina with relative spike latencies}, Science \textbf{319}, 1108 (2008).
\bibitem{OKeefe1993} J. O'Keefe and M. L. Recce, \emph{Phase relationship between hippocampal place units and the EEG theta rhythm}, Hippocampus \textbf{3}, 317 (1993).
\bibitem{Destexhe2003} A. Destexhe, M. Rudolph, and D. Par\'e, \emph{The high-conductance state of neocortical neurons in vivo}, Nat. Rev. Neurosci. \textbf{4}, 739 (2003).
\bibitem{Tsodyks1997} M. V. Tsodyks and H. Markram, \emph{The neural code between neocortical pyramidal neurons depends on neurotransmitter release probability}, Proc. Natl. Acad. Sci. USA \textbf{94}, 719 (1997).
\bibitem{Lisman1997} J. E. Lisman, \emph{Bursts as a unit of neural information: making unreliable synapses reliable}, Trends Neurosci. \textbf{20}, 38 (1997).
\bibitem{Attwell2001} D. Attwell and S. B. Laughlin, \emph{An energy budget for signaling in the grey matter of the brain}, J. Cereb. Blood Flow Metab. \textbf{21}, 1133 (2001).
\bibitem{Lennie2003} P. Lennie, \emph{The cost of cortical computation}, Curr. Biol. \textbf{13}, 493 (2003).
\bibitem{Yagishita2014} S. Yagishita et al., \emph{A critical time window for dopamine actions on the structural plasticity of dendritic spines}, Science \textbf{345}, 1616 (2014).
\bibitem{Werfel2005} J. Werfel, X. Xie, and H. S. Seung, \emph{Learning curves for stochastic gradient descent in linear feedforward networks}, Neural Comput. \textbf{17}, 2699 (2005).
\bibitem{Doya2000} K. Doya, \emph{Reinforcement learning in continuous time and space}, Neural Comput. \textbf{12}, 219 (2000).
\bibitem{Oja1982} E. Oja, \emph{Simplified neuron model as a principal component analyzer}, J. Math. Biol. \textbf{15}, 267 (1982).
\bibitem{BiPoo1998} G.-Q. Bi and M.-M. Poo, \emph{Synaptic modifications in cultured hippocampal neurons: dependence on spike timing, synaptic strength, and postsynaptic cell type}, J. Neurosci. \textbf{18}, 10464 (1998).
\bibitem{Williams1992} R. J. Williams, \emph{Simple statistical gradient-following algorithms for connectionist reinforcement learning}, Mach. Learn. \textbf{8}, 229 (1992).
\bibitem{Bayer2005} H. M. Bayer and P. W. Glimcher, \emph{Midbrain dopamine neurons encode a quantitative reward prediction error signal}, Neuron \textbf{47}, 129 (2005).
\bibitem{Shen2008} W. Shen, M. Flajolet, P. Greengard, and D. J. Surmeier, \emph{Dichotomous dopaminergic control of striatal synaptic plasticity}, Science \textbf{321}, 848 (2008).
\bibitem{Dabney2020} W. Dabney, Z. Kurth-Nelson, N. Uchida, C. K. Starkweather, D. Hassabis, R. Munos, and M. Botvinick, \emph{A distributional code for value in dopamine-based reinforcement learning}, Nature \textbf{577}, 671 (2020).
\bibitem{Matsumoto2009} M. Matsumoto and O. Hikosaka, \emph{Two types of dopamine neuron distinctly convey positive and negative motivational signals}, Nature \textbf{459}, 837 (2009).
\bibitem{BrombergMartin2010} E. S. Bromberg-Martin, M. Matsumoto, and O. Hikosaka, \emph{Dopamine in motivational control: rewarding, aversive, and alerting}, Neuron \textbf{68}, 815 (2010).
\bibitem{Menegas2018} W. Menegas, K. Akiti, R. Amo, N. Uchida, and M. Watabe-Uchida, \emph{Dopamine neurons projecting to the posterior striatum reinforce avoidance of threatening stimuli}, Nat. Neurosci. \textbf{21}, 1421 (2018).
\bibitem{Niv2007} Y. Niv, N. D. Daw, D. Joel, and P. Dayan, \emph{Tonic dopamine: opportunity costs and the control of response vigor}, Psychopharmacology \textbf{191}, 507 (2007).
\bibitem{Mohebi2019} A. Mohebi et al., \emph{Dissociable dopamine dynamics for learning and motivation}, Nature \textbf{570}, 65 (2019).
\bibitem{Hamid2016} A. A. Hamid et al., \emph{Mesolimbic dopamine signals the value of work}, Nat. Neurosci. \textbf{19}, 117 (2016).
\bibitem{Kim2020} H. R. Kim, A. N. Malik et al., \emph{A unified framework for dopamine signals across timescales}, Cell \textbf{183}, 1600 (2020).
\bibitem{Jeong2022} H. Jeong et al., \emph{Mesolimbic dopamine release conveys causal associations}, Science \textbf{378}, eabq6740 (2022).
\bibitem{Amo2022} R. Amo et al., \emph{A gradual temporal shift of dopamine responses mirrors the progression of temporal difference error in machine learning}, Nat. Neurosci. \textbf{25}, 1082 (2022).
\bibitem{Takahashi2017} Y. K. Takahashi et al., \emph{Dopamine neurons respond to errors in the prediction of sensory features of expected rewards}, Neuron \textbf{95}, 1395 (2017).
\bibitem{Sharpe2017} M. J. Sharpe et al., \emph{Dopamine transients are sufficient and necessary for acquisition of model-based associations}, Nat. Neurosci. \textbf{20}, 735 (2017).
\bibitem{Gardner2018} M. P. H. Gardner, G. Schoenbaum, and S. J. Gershman, \emph{Rethinking dopamine as generalized prediction error}, Proc. R. Soc. B \textbf{285}, 20181645 (2018).
\bibitem{Engelhard2019} B. Engelhard et al., \emph{Specialized coding of sensory, motor and cognitive variables in VTA dopamine neurons}, Nature \textbf{570}, 509 (2019).
\bibitem{Clements1992} J. D. Clements, R. A. J. Lester, G. Tong, C. E. Jahr, and G. L. Westbrook, \emph{The time course of glutamate in the synaptic cleft}, Science \textbf{258}, 1498 (1992).
\bibitem{Hasselmo2006} M. E. Hasselmo, \emph{The role of acetylcholine in learning and memory}, Curr. Opin. Neurobiol. \textbf{16}, 710 (2006).
\bibitem{JacobsAzmitia1992} B. L. Jacobs and E. C. Azmitia, \emph{Structure and function of the brain serotonin system}, Physiol. Rev. \textbf{72}, 165 (1992).
\bibitem{Schiller1992} P. H. Schiller, \emph{The ON and OFF channels of the visual system}, Trends Neurosci. \textbf{15}, 86 (1992).
\bibitem{Grothe2010} B. Grothe, M. Pecka, and D. McAlpine, \emph{Mechanisms of sound localization in mammals}, Physiol. Rev. \textbf{90}, 983 (2010).
\bibitem{Markram2004} H. Markram et al., \emph{Interneurons of the neocortical inhibitory system}, Nat. Rev. Neurosci. \textbf{5}, 793 (2004).
\bibitem{Joshi2016} S. Joshi, Y. Li, R. M. Kalwani, and J. I. Gold, \emph{Relationships between pupil diameter and neuronal activity in the locus coeruleus, colliculi, and cingulate cortex}, Neuron \textbf{89}, 221 (2016).
\bibitem{ServanSchreiber1990} D. Servan-Schreiber, H. Printz, and J. D. Cohen, \emph{A network model of catecholamine effects: gain, signal-to-noise ratio, and behavior}, Science \textbf{249}, 892 (1990).
\bibitem{YuDayan2005} A. J. Yu and P. Dayan, \emph{Uncertainty, neuromodulation, and attention}, Neuron \textbf{46}, 681 (2005).
\bibitem{Nassar2012} M. R. Nassar et al., \emph{Rational regulation of learning dynamics by pupil-linked arousal systems}, Nat. Neurosci. \textbf{15}, 1040 (2012).
\bibitem{BouretSara2005} S. Bouret and S. J. Sara, \emph{Network reset: a simplified overarching theory of locus coeruleus noradrenaline function}, Trends Neurosci. \textbf{28}, 574 (2005).
\bibitem{Hebb1949} D. O. Hebb, \emph{The Organization of Behavior: A Neuropsychological Theory} (Wiley, New York, 1949).
\bibitem{MacKay1952} D. M. MacKay and W. S. McCulloch, \emph{The limiting information capacity of a neuronal link}, Bull. Math. Biophys. \textbf{14}, 127 (1952).
\bibitem{WilsonCowan1972} H. R. Wilson and J. D. Cowan, \emph{Excitatory and inhibitory interactions in localized populations of model neurons}, Biophys. J. \textbf{12}, 1 (1972).
\bibitem{AllenStevens1994} C. Allen and C. F. Stevens, \emph{An evaluation of causes for unreliability of synaptic transmission}, Proc. Natl. Acad. Sci. USA \textbf{91}, 10380 (1994).
\bibitem{Barlow1961} H. B. Barlow, \emph{Possible principles underlying the transformations of sensory messages}, in \emph{Sensory Communication}, ed. W. A. Rosenblith (MIT Press, Cambridge, MA, 1961), pp.~217--234.
\bibitem{Cardin2009} J. A. Cardin et al., \emph{Driving fast-spiking cells induces gamma rhythm and controls sensory responses}, Nature \textbf{459}, 663 (2009).
\bibitem{Lillicrap2020} T. P. Lillicrap, A. Santoro, L. Marris, C. J. Akerman, and G. Hinton, \emph{Backpropagation and the brain}, Nat. Rev. Neurosci. \textbf{21}, 335 (2020).
\bibitem{Fremaux2016} N. Fr\'emaux and W. Gerstner, \emph{Neuromodulated spike-timing-dependent plasticity, and theory of three-factor learning rules}, Front. Neural Circuits \textbf{9}, 85 (2016).
\bibitem{Haider2006} B. Haider, A. Duque, A. R. Hasenstaub, and D. A. McCormick, \emph{Neocortical network activity in vivo is generated through a dynamic balance of excitation and inhibition}, J. Neurosci. \textbf{26}, 4535 (2006).
\bibitem{HoltKoch1997} G. R. Holt and C. Koch, \emph{Shunting inhibition does not have a divisive effect on firing rates}, Neural Comput. \textbf{9}, 1001 (1997).
\bibitem{Chance2002} F. S. Chance, L. F. Abbott, and A. D. Reyes, \emph{Gain modulation from background synaptic input}, Neuron \textbf{35}, 773 (2002).
\bibitem{Ozeki2009} H. Ozeki, I. M. Finn, E. S. Schaffer, K. D. Miller, and D. Ferster, \emph{Inhibitory stabilization of the cortical network underlies visual surround suppression}, Neuron \textbf{62}, 578 (2009).
\bibitem{HasselmoBower1992} M. E. Hasselmo and J. M. Bower, \emph{Cholinergic suppression specific to intrinsic not afferent fiber synapses in rat piriform (olfactory) cortex}, J. Neurophysiol. \textbf{67}, 1222 (1992).
\bibitem{HasselmoAndersonBower1992} M. E. Hasselmo, B. P. Anderson, and J. M. Bower, \emph{Cholinergic modulation of cortical associative memory function}, J. Neurophysiol. \textbf{67}, 1230 (1992).
\bibitem{Hasselmo1999} M. E. Hasselmo, \emph{Neuromodulation: acetylcholine and memory consolidation}, Trends Cogn. Sci. \textbf{3}, 351 (1999).
\bibitem{Tsodyks1988} M. V. Tsodyks and M. V. Feigel'man, \emph{The enhanced storage capacity in neural networks with low activity level}, Europhys. Lett. \textbf{6}, 101 (1988).
\bibitem{KalmanBucy1961} R. E. Kalman and R. S. Bucy, \emph{New results in linear filtering and prediction theory}, J. Basic Eng. \textbf{83}, 95 (1961).
\bibitem{WilsonMcNaughton1994} M. A. Wilson and B. L. McNaughton, \emph{Reactivation of hippocampal ensemble memories during sleep}, Science \textbf{265}, 676 (1994).
\bibitem{BarnesSharp1999} N. M. Barnes and T. Sharp, \emph{A review of central 5-HT receptors and their function}, Neuropharmacology \textbf{38}, 1083 (1999).
\bibitem{Bunin1998} M. A. Bunin and R. M. Wightman, \emph{Quantitative evaluation of 5-hydroxytryptamine (serotonin) neuronal release and uptake: an investigation of extrasynaptic transmission}, J. Neurosci. \textbf{18}, 4854 (1998).
\bibitem{Sprouse1987} J. S. Sprouse and G. K. Aghajanian, \emph{Electrophysiological responses of serotoninergic dorsal raphe neurons to 5-HT$_{1A}$ and 5-HT$_{1B}$ agonists}, Synapse \textbf{1}, 3 (1987).
\bibitem{Sykova2008} E. Syková and C. Nicholson, \emph{Diffusion in brain extracellular space}, Physiol. Rev. \textbf{88}, 1277 (2008).
\bibitem{WilsonNicoll2001} R. I. Wilson and R. A. Nicoll, \emph{Endogenous cannabinoids mediate retrograde signalling at hippocampal synapses}, Nature \textbf{410}, 588 (2001).
\bibitem{Garthwaite2008} J. Garthwaite, \emph{Concepts of neural nitric oxide-mediated transmission}, Eur. J. Neurosci. \textbf{27}, 2783 (2008).
\bibitem{vandenPol2012} A. N. van den Pol, \emph{Neuropeptide transmission in brain circuits}, Neuron \textbf{76}, 98 (2012).
\bibitem{Dunwiddie2001} T. V. Dunwiddie and S. A. Masino, \emph{The role and regulation of adenosine in the central nervous system}, Annu. Rev. Neurosci. \textbf{24}, 31 (2001).
\bibitem{Skaggs1996} W. E. Skaggs and B. L. McNaughton, \emph{Replay of neuronal firing sequences in rat hippocampus during sleep following spatial experience}, Science \textbf{271}, 1870 (1996).
\bibitem{Baylor1979} D. A. Baylor, T. D. Lamb, and K.-W. Yau, \emph{Responses of retinal rods to single photons}, J. Physiol. \textbf{288}, 613 (1979).
\bibitem{Hudspeth1989} A. J. Hudspeth, \emph{How the ear's works work}, Nature \textbf{341}, 397 (1989).
\bibitem{NakaRushton1966} K. I. Naka and W. A. H. Rushton, \emph{S-potentials from colour units in the retina of fish (Cyprinidae)}, J. Physiol. \textbf{185}, 536 (1966).
\bibitem{Lichtsteiner2008} P. Lichtsteiner, C. Posch, and T. Delbruck, \emph{A $128\times128$ 120 dB 15 $\mu$s latency asynchronous temporal contrast vision sensor}, IEEE J. Solid-State Circuits \textbf{43}, 566 (2008).
\bibitem{Koch2006} K. Koch et al., \emph{How much the eye tells the brain}, Curr. Biol. \textbf{16}, 1428 (2006).
\bibitem{Robles2001} L. Robles and M. A. Ruggero, \emph{Mechanics of the mammalian cochlea}, Physiol. Rev. \textbf{81}, 1305 (2001).
\bibitem{Mombaerts2004} P. Mombaerts, \emph{Genes and ligands for odorant, vomeronasal and taste receptors}, Nat. Rev. Neurosci. \textbf{5}, 263 (2004).
\bibitem{WoodSlater2001} S. J. Wood and C. R. Slater, \emph{Safety factor at the neuromuscular junction}, Prog. Neurobiol. \textbf{64}, 393 (2001).
\bibitem{Henneman1957} E. Henneman, \emph{Relation between size of neurons and their susceptibility to discharge}, Science \textbf{126}, 1345 (1957).
\bibitem{HarrisWolpert1998} C. M. Harris and D. M. Wolpert, \emph{Signal-dependent noise determines motor planning}, Nature \textbf{394}, 780 (1998).
\bibitem{Hogan1984} N. Hogan, \emph{Adaptive control of mechanical impedance by coactivation of antagonist muscles}, IEEE Trans. Autom. Control \textbf{29}, 681 (1984).
\bibitem{McAuleyMarsden2000} J. H. McAuley and C. D. Marsden, \emph{Physiological and pathological tremors and rhythmic central motor control}, Brain \textbf{123}, 1545 (2000).
\bibitem{WolpertGhahramani2000} D. M. Wolpert and Z. Ghahramani, \emph{Computational principles of movement neuroscience}, Nat. Neurosci. \textbf{3}, 1212 (2000).
\bibitem{MarderBucher2001} E. Marder and D. Bucher, \emph{Central pattern generators and the control of rhythmic movements}, Curr. Biol. \textbf{11}, R986 (2001).
\bibitem{Miyazaki2014} K. W. Miyazaki et al., \emph{Optogenetic activation of dorsal raphe serotonin neurons enhances patience for future rewards}, Curr. Biol. \textbf{24}, 2033 (2014).
\bibitem{Fuglevand1993} A. J. Fuglevand, D. A. Winter, and A. E. Patla, \emph{Models of recruitment and rate coding organization in motor-unit pools}, J. Neurophysiol. \textbf{70}, 2470 (1993).
\bibitem{Curcio1990} C. A. Curcio, K. R. Sloan, R. E. Kalina, and A. E. Hendrickson, \emph{Human photoreceptor topography}, J. Comp. Neurol. \textbf{292}, 497 (1990).
\bibitem{Hayes1950} N. D. Hayes, \emph{Roots of the transcendental equation associated with a certain difference-differential equation}, J. London Math. Soc. \textbf{25}, 226 (1950).
\bibitem{MelzackWall1965} R. Melzack and P. D. Wall, \emph{Pain mechanisms: a new theory}, Science \textbf{150}, 971 (1965).
\bibitem{Fields2004} H. Fields, \emph{State-dependent opioid control of pain}, Nat. Rev. Neurosci. \textbf{5}, 565 (2004).
\bibitem{Levine1978} J. D. Levine, N. C. Gordon, and H. L. Fields, \emph{The mechanism of placebo analgesia}, Lancet \textbf{312}, 654 (1978).
\bibitem{Woolf1983} C. J. Woolf, \emph{Evidence for a central component of post-injury pain hypersensitivity}, Nature \textbf{306}, 686 (1983).
\bibitem{Seymour2004} B. Seymour et al., \emph{Temporal difference models describe higher-order learning in humans}, Nature \textbf{429}, 664 (2004).
\bibitem{EcclestonCrombez1999} C. Eccleston and G. Crombez, \emph{Pain demands attention: a cognitive-affective model of the interruptive function of pain}, Psychol. Bull. \textbf{125}, 356 (1999).
\bibitem{WidrowHoff1960} B. Widrow and M. E. Hoff, \emph{Adaptive switching circuits}, IRE WESCON Convention Record, part 4, 96 (1960).
\bibitem{Marr1969} D. Marr, \emph{A theory of cerebellar cortex}, J. Physiol. \textbf{202}, 437 (1969).
\bibitem{Albus1971} J. S. Albus, \emph{A theory of cerebellar function}, Math. Biosci. \textbf{10}, 25 (1971).
\bibitem{Cover1965} T. M. Cover, \emph{Geometrical and statistical properties of systems of linear inequalities with applications in pattern recognition}, IEEE Trans. Electron. Comput. \textbf{EC-14}, 326 (1965).
\bibitem{Ito2001} M. Ito, \emph{Cerebellar long-term depression: characterization, signal transduction, and functional roles}, Physiol. Rev. \textbf{81}, 1143 (2001).
\bibitem{Suvrathan2016} A. Suvrathan, H. L. Payne, and J. L. Raymond, \emph{Timing rules for synaptic plasticity matched to behavioral function}, Neuron \textbf{92}, 959 (2016).
\bibitem{Fujita1982} M. Fujita, \emph{Adaptive filter model of the cerebellum}, Biol. Cybern. \textbf{45}, 195 (1982).
\bibitem{Blakemore1998} S.-J. Blakemore, D. M. Wolpert, and C. D. Frith, \emph{Central cancellation of self-produced tickle sensation}, Nat. Neurosci. \textbf{1}, 635 (1998).
\bibitem{Smith2006} M. A. Smith, A. Ghazizadeh, and R. Shadmehr, \emph{Interacting adaptive processes with different timescales underlie short-term motor learning}, PLoS Biol. \textbf{4}, e179 (2006).
\bibitem{Kording2007} K. P. K\"ording, J. B. Tenenbaum, and R. Shadmehr, \emph{The dynamics of memory as a consequence of optimal adaptation to a changing body}, Nat. Neurosci. \textbf{10}, 779 (2007).
\bibitem{Yao2023} Z. Yao et al., \emph{A high-resolution transcriptomic and spatial atlas of cell types in the whole mouse brain}, Nature \textbf{624}, 317 (2023).
\bibitem{Sanzeni2020} A. Sanzeni, B. Akitake, H. C. Goldbach, C. E. Leedy, N. Brunel, and M. H. Histed, \emph{Inhibition stabilization is a widespread property of cortical networks}, eLife \textbf{9}, e54875 (2020).
\bibitem{Padamsey2022} Z. Padamsey, D. Katsanevaki, N. Dupuy, and N. L. Rochefort, \emph{Neocortex saves energy by reducing coding precision during food scarcity}, Neuron \textbf{110}, 280 (2022).
\bibitem{Stringer2019} C. Stringer, M. Pachitariu, N. Steinmetz, C. B. Reddy, M. Carandini, and K. D. Harris, \emph{Spontaneous behaviors drive multidimensional, brainwide activity}, Science \textbf{364}, eaav7893 (2019).
\bibitem{Bittner2017} K. C. Bittner, A. D. Milstein, C. Grienberger, S. Romani, and J. C. Magee, \emph{Behavioral time scale synaptic plasticity underlies CA1 place fields}, Science \textbf{357}, 1033 (2017).
\bibitem{WhittingtonBogacz2019} J. C. R. Whittington and R. Bogacz, \emph{Theories of error back-propagation in the brain}, Trends Cogn. Sci. \textbf{23}, 235 (2019).
\bibitem{Reimer2016} J. Reimer et al., \emph{Pupil fluctuations track rapid changes in adrenergic and cholinergic activity in cortex}, Nat. Commun. \textbf{7}, 13289 (2016).
\bibitem{Beniaguev2021} D. Beniaguev, I. Segev, and M. London, \emph{Single cortical neurons as deep artificial neural networks}, Neuron \textbf{109}, 2727 (2021).
\bibitem{Gidon2020} A. Gidon et al., \emph{Dendritic action potentials and computation in human layer 2/3 cortical neurons}, Science \textbf{367}, 83 (2020).
\bibitem{Gallego2022} G. Gallego et al., \emph{Event-based vision: a survey}, IEEE Trans. Pattern Anal. Mach. Intell. \textbf{44}, 154 (2022).
\bibitem{Berger1989} R. D. Berger, J. P. Saul, and R. J. Cohen, \emph{Transfer function analysis of autonomic regulation. I. Canine atrial rate response}, Am. J. Physiol. \textbf{256}, H142 (1989).
\bibitem{Walker2010} J. J. Walker, J. R. Terry, and S. L. Lightman, \emph{Origin of ultradian pulsatility in the hypothalamic--pituitary--adrenal axis}, Proc. R. Soc. B \textbf{277}, 1627 (2010).
\bibitem{Belchetz1978} P. E. Belchetz, T. M. Plant, Y. Nakai, E. J. Keogh, and E. Knobil, \emph{Hypophysial responses to continuous and intermittent delivery of hypothalamic gonadotropin-releasing hormone}, Science \textbf{202}, 631 (1978).
\bibitem{Czeisler1999} C. A. Czeisler et al., \emph{Stability, precision, and near-24-hour period of the human circadian pacemaker}, Science \textbf{284}, 2177 (1999).
\bibitem{TakahashiJS2017} J. S. Takahashi, \emph{Transcriptional architecture of the mammalian circadian clock}, Nat. Rev. Genet. \textbf{18}, 164 (2017).
\bibitem{Musk2019} E. Musk and Neuralink, \emph{An integrated brain-machine interface platform with thousands of channels}, J. Med. Internet Res. \textbf{21}, e16194 (2019).
\bibitem{Hochberg2006} L. R. Hochberg et al., \emph{Neuronal ensemble control of prosthetic devices by a human with tetraplegia}, Nature \textbf{442}, 164 (2006).
\bibitem{Willett2021} F. R. Willett et al., \emph{High-performance brain-to-text communication via handwriting}, Nature \textbf{593}, 249 (2021).
\bibitem{Willett2023} F. R. Willett et al., \emph{A high-performance speech neuroprosthesis}, Nature \textbf{620}, 1031 (2023).
\bibitem{Gallego2017} J. A. Gallego, M. G. Perich, L. E. Miller, and S. A. Solla, \emph{Neural manifolds for the control of movement}, Neuron \textbf{94}, 978 (2017).
\bibitem{Fernandez2021} E. Fern\'andez et al., \emph{Visual percepts evoked with an intracortical 96-channel microelectrode array inserted in human occipital cortex}, J. Clin. Invest. \textbf{131}, e151331 (2021).
\bibitem{Tritsch2012} N. X. Tritsch, J. B. Ding, and B. L. Sabatini, \emph{Dopaminergic neurons inhibit striatal output through non-canonical release of GABA}, Nature \textbf{490}, 262 (2012).
\bibitem{Zhang2015} S. Zhang et al., \emph{Dopaminergic and glutamatergic microdomains in a subset of rodent mesoaccumbens axons}, Nat. Neurosci. \textbf{18}, 386 (2015).
\bibitem{Mongillo2008} G. Mongillo, O. Barak, and M. Tsodyks, \emph{Synaptic theory of working memory}, Science \textbf{319}, 1543 (2008).
\bibitem{Stokes2015} M. G. Stokes, \emph{`Activity-silent' working memory in prefrontal cortex: a dynamic coding framework}, Trends Cogn. Sci. \textbf{19}, 394 (2015).
\bibitem{Smith2013} S. L. Smith, I. T. Smith, T. Branco, and M. H\"ausser, \emph{Dendritic spikes enhance stimulus selectivity in cortical neurons in vivo}, Nature \textbf{503}, 115 (2013).
\bibitem{Baden2016} T. Baden et al., \emph{The functional diversity of retinal ganglion cells in the mouse}, Nature \textbf{529}, 345 (2016).
\bibitem{Lottem2018} E. Lottem et al., \emph{Activation of serotonin neurons promotes active persistence in a probabilistic foraging task}, Nat. Commun. \textbf{9}, 1000 (2018).
\bibitem{Payeur2021} A. Payeur, J. Guerguiev, F. Zenke, B. A. Richards, and R. Naud, \emph{Burst-dependent synaptic plasticity can coordinate learning in hierarchical circuits}, Nat. Neurosci. \textbf{24}, 1010 (2021).
\bibitem{MehonicKenyon2022} A. Mehonic and A. J. Kenyon, \emph{Brain-inspired computing needs a master plan}, Nature \textbf{604}, 255 (2022).
\bibitem{Poe2020} G. R. Poe et al., \emph{Locus coeruleus: a new look at the blue spot}, Nat. Rev. Neurosci. \textbf{21}, 644 (2020).
\bibitem{Hangya2015} B. Hangya, S. P. Ranade, M. Lorenc, and A. Kepecs, \emph{Central cholinergic neurons are rapidly recruited by reinforcement feedback}, Cell \textbf{162}, 1155 (2015).
\bibitem{FernandezRuiz2019} A. Fern\'andez-Ruiz et al., \emph{Long-duration hippocampal sharp wave ripples improve memory}, Science \textbf{364}, 1082 (2019).
\bibitem{Herzfeld2018} D. J. Herzfeld, Y. Kojima, R. Soetedjo, and R. Shadmehr, \emph{Encoding of error and learning to correct that error by the Purkinje cells of the cerebellum}, Nat. Neurosci. \textbf{21}, 736 (2018).
\bibitem{Duan2014} B. Duan et al., \emph{Identification of spinal circuits transmitting and gating mechanical pain}, Cell \textbf{159}, 1417 (2014).
\bibitem{Tremblay2016} R. Tremblay, S. Lee, and B. Rudy, \emph{GABAergic interneurons in the neocortex: from cellular properties to circuits}, Neuron \textbf{91}, 260 (2016).
\bibitem{Ren2018} J. Ren et al., \emph{Anatomically defined and functionally distinct dorsal raphe serotonin sub-systems}, Cell \textbf{175}, 472 (2018).
\bibitem{Pfeffer2013} C. K. Pfeffer, M. Xue, M. He, Z. J. Huang, and M. Scanziani, \emph{Inhibition of inhibition in visual cortex: the logic of connections between molecularly distinct interneurons}, Nat. Neurosci. \textbf{16}, 1068 (2013).
\bibitem{Pi2013} H.-J. Pi et al., \emph{Cortical interneurons that specialize in disinhibitory control}, Nature \textbf{503}, 521 (2013).
\bibitem{Keller2018} G. B. Keller and T. D. Mrsic-Flogel, \emph{Predictive processing: a canonical cortical computation}, Neuron \textbf{100}, 424 (2018).
\bibitem{HubelWiesel1962} D. H. Hubel and T. N. Wiesel, \emph{Receptive fields, binocular interaction and functional architecture in the cat's visual cortex}, J. Physiol. \textbf{160}, 106 (1962).
\bibitem{Riesenhuber1999} M. Riesenhuber and T. Poggio, \emph{Hierarchical models of object recognition in cortex}, Nat. Neurosci. \textbf{2}, 1019 (1999).
\bibitem{Fukushima1980} K. Fukushima, \emph{Neocognitron: a self-organizing neural network model for a mechanism of pattern recognition unaffected by shift in position}, Biol. Cybern. \textbf{36}, 193 (1980).
\bibitem{Yamins2014} D. L. K. Yamins et al., \emph{Performance-optimized hierarchical models predict neural responses in higher visual cortex}, Proc. Natl. Acad. Sci. USA \textbf{111}, 8619 (2014).
\bibitem{Hung2005} C. P. Hung, G. Kreiman, T. Poggio, and J. J. DiCarlo, \emph{Fast readout of object identity from macaque inferior temporal cortex}, Science \textbf{310}, 863 (2005).
\bibitem{WiskottSejnowski2002} L. Wiskott and T. J. Sejnowski, \emph{Slow feature analysis: unsupervised learning of invariances}, Neural Comput. \textbf{14}, 715 (2002).
\bibitem{Foldiak1991} P. F\"oldi\'ak, \emph{Learning invariance from transformation sequences}, Neural Comput. \textbf{3}, 194 (1991).
\bibitem{LiDiCarlo2008} N. Li and J. J. DiCarlo, \emph{Unsupervised natural experience rapidly alters invariant object representation in visual cortex}, Science \textbf{321}, 1502 (2008).
\bibitem{RaoBallard1999} R. P. N. Rao and D. H. Ballard, \emph{Predictive coding in the visual cortex: a functional interpretation of some extra-classical receptive-field effects}, Nat. Neurosci. \textbf{2}, 79 (1999).
\bibitem{Kar2019} K. Kar, J. Kubilius, K. Schmidt, E. B. Issa, and J. J. DiCarlo, \emph{Evidence that recurrent circuits are critical to the ventral stream's execution of core object recognition behavior}, Nat. Neurosci. \textbf{22}, 974 (2019).
\bibitem{Szegedy2014} C. Szegedy et al., \emph{Intriguing properties of neural networks}, in \emph{Proc. ICLR} (2014), arXiv:1312.6199.
\bibitem{Weiss2002} Y. Weiss, E. P. Simoncelli, and E. H. Adelson, \emph{Motion illusions as optimal percepts}, Nat. Neurosci. \textbf{5}, 598 (2002).
\bibitem{Purves2011} D. Purves, W. T. Wojtach, and R. B. Lotto, \emph{Understanding vision in wholly empirical terms}, Proc. Natl. Acad. Sci. USA \textbf{108}, Suppl.~3, 15588 (2011).
\bibitem{LaferSousa2015} R. Lafer-Sousa, K. L. Hermann, and B. R. Conway, \emph{Striking individual differences in color perception uncovered by `the dress' photograph}, Curr. Biol. \textbf{25}, R545 (2015).
\bibitem{Wallisch2017} P. Wallisch, \emph{Illumination assumptions account for individual differences in the perceptual interpretation of a profoundly ambiguous stimulus in the color domain: ``The dress''}, J. Vis. \textbf{17}(4), 5 (2017).
\bibitem{Schwartz2009} O. Schwartz, T. J. Sejnowski, and P. Dayan, \emph{Perceptual organization in the tilt illusion}, J. Vis. \textbf{9}(4), 19 (2009).
\bibitem{SchillerCarvey2005} P. H. Schiller and C. E. Carvey, \emph{The Hermann grid illusion revisited}, Perception \textbf{34}, 1375 (2005).
\bibitem{Nijhawan1994} R. Nijhawan, \emph{Motion extrapolation in catching}, Nature \textbf{370}, 256 (1994).
\bibitem{Conway2005} B. R. Conway, A. Kitaoka, A. Yazdanbakhsh, C. C. Pack, and M. S. Livingstone, \emph{Neural basis for a powerful static motion illusion}, J. Neurosci. \textbf{25}, 5651 (2005).
\bibitem{OteroMillan2012} J. Otero-Millan, S. L. Macknik, and S. Martinez-Conde, \emph{Microsaccades and blinks trigger illusory rotation in the ``rotating snakes'' illusion}, J. Neurosci. \textbf{32}, 6043 (2012).
\bibitem{MorenoBote2007} R. Moreno-Bote, J. Rinzel, and N. Rubin, \emph{Noise-induced alternations in an attractor network model of perceptual bistability}, J. Neurophysiol. \textbf{98}, 1125 (2007).
\bibitem{Watanabe2018} E. Watanabe, A. Kitaoka, K. Sakamoto, M. Yasugi, and K. Tanaka, \emph{Illusory motion reproduced by deep neural networks trained for prediction}, Front. Psychol. \textbf{9}, 345 (2018).
\bibitem{GomezVilla2020} A. G\'omez-Villa, A. Mart\'in, J. Vazquez-Corral, M. Bertalm\'io, and J. Malo, \emph{Color illusions also deceive CNNs for low-level vision tasks: analysis and implications}, Vision Res. \textbf{176}, 156 (2020).
\bibitem{delCastillo1954} J. del Castillo and B. Katz, \emph{Quantal components of the end-plate potential}, J. Physiol. \textbf{124}, 560 (1954).
\bibitem{Nowak1984} L. Nowak, P. Bregestovski, P. Ascher, A. Herbet, and A. Prochiantz, \emph{Magnesium gates glutamate-activated channels in mouse central neurones}, Nature \textbf{307}, 462 (1984).
\bibitem{Malinow2002} R. Malinow and R. C. Malenka, \emph{AMPA receptor trafficking and synaptic plasticity}, Annu. Rev. Neurosci. \textbf{25}, 103 (2002).
\bibitem{Matsuzaki2004} M. Matsuzaki, N. Honkura, G. C. R. Ellis-Davies, and H. Kasai, \emph{Structural basis of long-term potentiation in single dendritic spines}, Nature \textbf{429}, 761 (2004).
\bibitem{Rudomin1999} P. Rudomin and R. F. Schmidt, \emph{Presynaptic inhibition in the vertebrate spinal cord revisited}, Exp. Brain Res. \textbf{129}, 1 (1999).
\bibitem{HutcheonYarom2000} B. Hutcheon and Y. Yarom, \emph{Resonance, oscillation and the intrinsic frequency preferences of neurons}, Trends Neurosci. \textbf{23}, 216 (2000).
\bibitem{Seung1996} H. S. Seung, \emph{How the brain keeps the eyes still}, Proc. Natl. Acad. Sci. USA \textbf{93}, 13339 (1996).
\bibitem{Hounsgaard1988} J. Hounsgaard, H. Hultborn, B. Jespersen, and O. Kiehn, \emph{Bistability of alpha-motoneurones in the decerebrate cat and in the acute spinal cat after intravenous 5-hydroxytryptophan}, J. Physiol. \textbf{405}, 345 (1988).
\bibitem{Izhikevich2007} E. M. Izhikevich, \emph{Dynamical Systems in Neuroscience: The Geometry of Excitability and Bursting} (MIT Press, Cambridge, MA, 2007).
\bibitem{AgmonSnir1998} H. Agmon-Snir, C. E. Carr, and J. Rinzel, \emph{The role of dendrites in auditory coincidence detection}, Nature \textbf{393}, 268 (1998).
\bibitem{BorstSoria2012} J. G. G. Borst and J. Soria van Hoeve, \emph{The calyx of Held synapse: from model synapse to auditory relay}, Annu. Rev. Physiol. \textbf{74}, 199 (2012).
\bibitem{Euler2002} T. Euler, P. B. Detwiler, and W. Denk, \emph{Directionally selective calcium signals in dendrites of starburst amacrine cells}, Nature \textbf{418}, 845 (2002).
\bibitem{AstonJonesBloom1981} G. Aston-Jones and F. E. Bloom, \emph{Activity of norepinephrine-containing locus coeruleus neurons in behaving rats anticipates fluctuations in the sleep-waking cycle}, J. Neurosci. \textbf{1}, 876 (1981).
\bibitem{BrooksPeever2012} P. L. Brooks and J. H. Peever, \emph{Identification of the transmitter and receptor mechanisms responsible for REM sleep paralysis}, J. Neurosci. \textbf{32}, 9785 (2012).
\bibitem{Borbely1982} A. A. Borb\'ely, \emph{A two process model of sleep regulation}, Hum. Neurobiol. \textbf{1}, 195 (1982).
\bibitem{PorkkaHeiskanen1997} T. Porkka-Heiskanen et al., \emph{Adenosine: a mediator of the sleep-inducing effects of prolonged wakefulness}, Science \textbf{276}, 1265 (1997).
\bibitem{Steriade1993} M. Steriade, A. N\'u\~nez, and F. Amzica, \emph{A novel slow ($<1$ Hz) oscillation of neocortical neurons in vivo: depolarizing and hyperpolarizing components}, J. Neurosci. \textbf{13}, 3252 (1993).
\bibitem{McCormick1992} D. A. McCormick, \emph{Neurotransmitter actions in the thalamus and cerebral cortex and their role in neuromodulation of thalamocortical activity}, Prog. Neurobiol. \textbf{39}, 337 (1992).
\bibitem{SteriadeMcCormickSejnowski1993} M. Steriade, D. A. McCormick, and T. J. Sejnowski, \emph{Thalamocortical oscillations in the sleeping and aroused brain}, Science \textbf{262}, 679 (1993).
\bibitem{Nadasdy1999} Z. N\'adasdy et al., \emph{Replay and time compression of recurring spike sequences in the hippocampus}, J. Neurosci. \textbf{19}, 9497 (1999).
\bibitem{Maingret2016} N. Maingret et al., \emph{Hippocampo-cortical coupling mediates memory consolidation during sleep}, Nat. Neurosci. \textbf{19}, 959 (2016).
\bibitem{McClelland1995} J. L. McClelland, B. L. McNaughton, and R. C. O'Reilly, \emph{Why there are complementary learning systems in the hippocampus and neocortex}, Psychol. Rev. \textbf{102}, 419 (1995).
\bibitem{TononiCirelli2014} G. Tononi and C. Cirelli, \emph{Sleep and the price of plasticity: from synaptic and cellular homeostasis to memory consolidation and integration}, Neuron \textbf{81}, 12 (2014).
\bibitem{deVivo2017} L. de Vivo et al., \emph{Ultrastructural evidence for synaptic scaling across the wake/sleep cycle}, Science \textbf{355}, 507 (2017).
\bibitem{Xie2013} L. Xie et al., \emph{Sleep drives metabolite clearance from the adult brain}, Science \textbf{342}, 373 (2013).
\bibitem{Miao2024} A. Miao et al., \emph{Brain clearance is reduced during sleep and anesthesia}, Nat. Neurosci. \textbf{27}, 1046 (2024).
\bibitem{Vyazovskiy2011} V. V. Vyazovskiy et al., \emph{Local sleep in awake rats}, Nature \textbf{472}, 443 (2011).
\bibitem{Lancaster2013} M. A. Lancaster et al., \emph{Cerebral organoids model human brain development and microcephaly}, Nature \textbf{501}, 373 (2013).
\bibitem{Jordan2024} F. D. Jordan et al., \emph{Open and remotely accessible Neuroplatform for research in wetware computing}, Front. Artif. Intell. \textbf{7}, 1376042 (2024).
\bibitem{Smirnova2023} L. Smirnova et al., \emph{Organoid intelligence (OI): the new frontier in biocomputing and intelligence-in-a-dish}, Front. Sci. \textbf{1}, 1017235 (2023).
\bibitem{Wagenaar2006} D. A. Wagenaar, J. Pine, and S. M. Potter, \emph{An extremely rich repertoire of bursting patterns during the development of cortical cultures}, BMC Neurosci. \textbf{7}, 11 (2006).
\bibitem{Shahaf2001} G. Shahaf and S. Marom, \emph{Learning in networks of cortical neurons}, J. Neurosci. \textbf{21}, 8782 (2001).
\bibitem{Kagan2022} B. J. Kagan et al., \emph{In vitro neurons learn and exhibit sentience when embodied in a simulated game-world}, Neuron \textbf{110}, 3952 (2022).
\bibitem{Friston2010} K. Friston, \emph{The free-energy principle: a unified brain theory?}, Nat. Rev. Neurosci. \textbf{11}, 127 (2010).
\bibitem{Balci2023} F. Balci et al., \emph{A response to claims of emergent intelligence and sentience in a dish}, Neuron \textbf{111}, 604 (2023).
\bibitem{Bakkum2008} D. J. Bakkum, Z. C. Chao, and S. M. Potter, \emph{Spatio-temporal electrical stimuli shape behavior of an embodied cortical network in a goal-directed learning task}, J. Neural Eng. \textbf{5}, 310 (2008).
\bibitem{Maass2002} W. Maass, T. Natschl\"ager, and H. Markram, \emph{Real-time computing without stable states: a new framework for neural computation based on perturbations}, Neural Comput. \textbf{14}, 2531 (2002).
\bibitem{JaegerHaas2004} H. Jaeger and H. Haas, \emph{Harnessing nonlinearity: predicting chaotic systems and saving energy in wireless communication}, Science \textbf{304}, 78 (2004).
\bibitem{Cai2023} H. Cai et al., \emph{Brain organoid reservoir computing for artificial intelligence}, Nat. Electron. \textbf{6}, 1032 (2023).
\bibitem{Habibollahi2023} F. Habibollahi, B. J. Kagan, A. N. Burkitt, and C. French, \emph{Critical dynamics arise during structured information presentation within embodied in vitro neuronal networks}, Nat. Commun. \textbf{14}, 5287 (2023).
\bibitem{Robbins2026} A. Robbins et al., \emph{Goal-directed learning in cortical organoids}, Cell Rep. \textbf{45}, 116984 (2026).
\bibitem{Keene2020} S. T. Keene et al., \emph{A biohybrid synapse with neurotransmitter-mediated plasticity}, Nat. Mater. \textbf{19}, 969 (2020).
\bibitem{Jo2016} J. Jo et al., \emph{Midbrain-like organoids from human pluripotent stem cells contain functional dopaminergic and neuromelanin-producing neurons}, Cell Stem Cell \textbf{19}, 248 (2016).
\bibitem{Hendry1987} S. H. Hendry, H. D. Schwark, E. G. Jones, and J. Yan, \emph{Numbers and proportions of GABA-immunoreactive neurons in different areas of monkey cerebral cortex}, J. Neurosci. \textbf{7}, 1503 (1987).
\bibitem{Tang2001synapses} Y. Tang, J. R. Nyengaard, D. M. G. De Groot, and H. J. G. Gundersen, \emph{Total regional and global number of synapses in the human brain neocortex}, Synapse \textbf{41}, 258 (2001).
\bibitem{Pakkenberg1997} B. Pakkenberg and H. J. G. Gundersen, \emph{Neocortical neuron number in humans: effect of sex and age}, J. Comp. Neurol. \textbf{384}, 312 (1997).
\bibitem{Matsuda2009} W. Matsuda, T. Furuta, K. C. Nakamura, H. Hioki, F. Fujiyama, R. Arai, and T. Kaneko, \emph{Single nigrostriatal dopaminergic neurons form widely spread and highly dense axonal arborizations in the neostriatum}, J. Neurosci. \textbf{29}, 444 (2009).
\bibitem{Pakkenberg1991} B. Pakkenberg, A. M{\o}ller, H. J. G. Gundersen, A. Mouritzen Dam, and H. Pakkenberg, \emph{The absolute number of nerve cells in substantia nigra in normal subjects and in patients with Parkinson's disease estimated with an unbiased stereological method}, J. Neurol. Neurosurg. Psychiatry \textbf{54}, 30 (1991).
\bibitem{Baker1990} K. G. Baker, G. M. Halliday, and I. T\"ork, \emph{Cytoarchitecture of the human dorsal raphe nucleus}, J. Comp. Neurol. \textbf{301}, 147 (1990).
\bibitem{Baker1991} K. G. Baker, G. M. Halliday, P. Halasz, J.-P. Hornung, L. B. Geffen, R. G. H. Cotton, and I. T\"ork, \emph{Cytoarchitecture of serotonin-synthesizing neurons in the pontine tegmentum of the human brain}, Synapse \textbf{7}, 301 (1991).
\bibitem{Airaksinen1991} M. S. Airaksinen, A. Paetau, L. Palj\"arvi, K. Reinikainen, P. Riekkinen, R. Suomalainen, and P. Panula, \emph{Histamine neurons in human hypothalamus: anatomy in normal and Alzheimer diseased brains}, Neuroscience \textbf{44}, 465 (1991).
\bibitem{Colquhoun1992} D. Colquhoun, P. Jonas, and B. Sakmann, \emph{Action of brief pulses of glutamate on AMPA/kainate receptors in patches from different neurones of rat hippocampal slices}, J. Physiol. \textbf{458}, 261 (1992).
\bibitem{DutarNicoll1988} P. Dutar and R. A. Nicoll, \emph{A physiological role for GABA$_B$ receptors in the central nervous system}, Nature \textbf{332}, 156 (1988).

\bibitem{Rauch2003} A. Rauch, G. La Camera, H.-R. L\"uscher, W. Senn, and S. Fusi, \emph{Neocortical pyramidal cells respond as integrate-and-fire neurons to in vivo-like input currents}, J. Neurophysiol. \textbf{90}, 1598 (2003).
\bibitem{KlumppEady1956} R. G. Klumpp and H. R. Eady, \emph{Some measurements of interaural time difference thresholds}, J. Acoust. Soc. Am. \textbf{28}, 859 (1956).
\bibitem{Brand2002} A. Brand, O. Behrend, T. Marquardt, D. McAlpine, and B. Grothe, \emph{Precise inhibition is essential for microsecond interaural time difference coding}, Nature \textbf{417}, 543 (2002).
\bibitem{Funahashi1989} S. Funahashi, C. J. Bruce, and P. S. Goldman-Rakic, \emph{Mnemonic coding of visual space in the monkey's dorsolateral prefrontal cortex}, J. Neurophysiol. \textbf{61}, 331 (1989).
\bibitem{WangMarkram2006} Y. Wang, H. Markram, P. H. Goodman, T. K. Berger, J. Ma, and P. S. Goldman-Rakic, \emph{Heterogeneity in the pyramidal network of the medial prefrontal cortex}, Nat. Neurosci. \textbf{9}, 534 (2006).

\bibitem{McCullochPitts1943} W. S. McCulloch and W. Pitts, \emph{A logical calculus of the ideas immanent in nervous activity}, Bull. Math. Biophys. \textbf{5}, 115 (1943).
\bibitem{Fried2002} S. I. Fried, T. A. M\"unch, and F. S. Werblin, \emph{Mechanisms and circuitry underlying directional selectivity in the retina}, Nature \textbf{420}, 411 (2002).

\bibitem{Hromadka2008} T. Hrom\'adka, M. R. DeWeese, and A. M. Zador, \emph{Sparse representation of sounds in the unanesthetized auditory cortex}, PLoS Biol. \textbf{6}, e16 (2008).
\bibitem{Abbott1997depression} L. F. Abbott, J. A. Varela, K. Sen, and S. B. Nelson, \emph{Synaptic depression and cortical gain control}, Science \textbf{275}, 221 (1997).
\bibitem{ShadlenNewsome1998} M. N. Shadlen and W. T. Newsome, \emph{The variable discharge of cortical neurons: implications for connectivity, computation, and information coding}, J. Neurosci. \textbf{18}, 3870 (1998).

\bibitem{TrulsonJacobs1979} M. E. Trulson and B. L. Jacobs, \emph{Raphe unit activity in freely moving cats: correlation with level of behavioral arousal}, Brain Res. \textbf{163}, 135 (1979).
\bibitem{GagnonParent2014} D. Gagnon and M. Parent, \emph{Distribution of VGLUT3 in highly collateralized axons from the rat dorsal raphe nucleus as revealed by single-neuron reconstructions}, PLoS One \textbf{9}, e87709 (2014).
\bibitem{VandermaelenAghajanian1983} C. P. Vandermaelen and G. K. Aghajanian, \emph{Electrophysiological and pharmacological characterization of serotonergic dorsal raphe neurons recorded extracellularly and intracellularly in rat brain slices}, Brain Res. \textbf{289}, 109 (1983).
\bibitem{CourtneyFord2016} N. A. Courtney and C. P. Ford, \emph{Mechanisms of 5-HT$_{1A}$ receptor-mediated transmission in dorsal raphe serotonin neurons}, J. Physiol. \textbf{594}, 953 (2016).
\bibitem{Hashemi2012serotonin} P. Hashemi et al., \emph{Brain dopamine and serotonin differ in regulation and its consequences}, Proc. Natl. Acad. Sci. USA \textbf{109}, 11510 (2012).
\bibitem{Black1934feedback} H. S. Black, \emph{Stabilized feedback amplifiers}, Bell Syst. Tech. J. \textbf{13}, 1 (1934).
\bibitem{KobayashiSchultz2008} S. Kobayashi and W. Schultz, \emph{Influence of reward delays on responses of dopamine neurons}, J. Neurosci. \textbf{28}, 7837 (2008).
\bibitem{Schweighofer2008} N. Schweighofer et al., \emph{Low-serotonin levels increase delayed reward discounting in humans}, J. Neurosci. \textbf{28}, 4528 (2008).

\bibitem{Malenka1988calcium} R. C. Malenka, J. A. Kauer, R. S. Zucker, and R. A. Nicoll, \emph{Postsynaptic calcium is sufficient for potentiation of hippocampal synaptic transmission}, Science \textbf{242}, 81 (1988).
\bibitem{Bliss1973LTP} T. V. P. Bliss and T. L{\o}mo, \emph{Long-lasting potentiation of synaptic transmission in the dentate area of the anaesthetized rabbit following stimulation of the perforant path}, J. Physiol. \textbf{232}, 331 (1973).
\bibitem{Kelso1986hebb} S. R. Kelso, A. H. Ganong, and T. H. Brown, \emph{Hebbian synapses in hippocampus}, Proc. Natl. Acad. Sci. USA \textbf{83}, 5326 (1986).
\bibitem{He2015eligibility} K. He et al., \emph{Distinct eligibility traces for LTP and LTD in cortical synapses}, Neuron \textbf{88}, 528 (2015).
\bibitem{Olveczky2005} B. P. \"Olveczky, A. S. Andalman, and M. S. Fee, \emph{Vocal experimentation in the juvenile songbird requires a basal ganglia circuit}, PLoS Biol. \textbf{3}, e153 (2005).
\bibitem{TumerBrainard2007} E. C. Tumer and M. S. Brainard, \emph{Performance variability enables adaptive plasticity of `crystallized' adult birdsong}, Nature \textbf{450}, 1240 (2007).
\bibitem{Eshel2015arithmetic} N. Eshel et al., \emph{Arithmetic and local circuitry underlying dopamine prediction errors}, Nature \textbf{525}, 243 (2015).
\bibitem{GraceOnn1989} A. A. Grace and S.-P. Onn, \emph{Morphology and electrophysiological properties of immunocytochemically identified rat dopamine neurons recorded in vitro}, J. Neurosci. \textbf{9}, 3463 (1989).
\bibitem{ODoherty2004actor} J. O'Doherty et al., \emph{Dissociable roles of ventral and dorsal striatum in instrumental conditioning}, Science \textbf{304}, 452 (2004).

\bibitem{NeweyPowell1987} W. K. Newey and J. L. Powell, \emph{Asymmetric least squares estimation and testing}, Econometrica \textbf{55}, 819 (1987).
\bibitem{Beierholm2013vigor} U. Beierholm et al., \emph{Dopamine modulates reward-related vigor}, Neuropsychopharmacology \textbf{38}, 1495 (2013).
\bibitem{Howe2013ramp} M. W. Howe, P. L. Tierney, S. G. Sandberg, P. E. M. Phillips, and A. M. Graybiel, \emph{Prolonged dopamine signalling in striatum signals proximity and value of distant rewards}, Nature \textbf{500}, 575 (2013).

\bibitem{Baker1989LC} K. G. Baker, I. T\"ork, J.-P. Hornung, and P. Halasz, \emph{The human locus coeruleus complex: an immunohistochemical and three dimensional reconstruction study}, Exp. Brain Res. \textbf{77}, 257 (1989).
\bibitem{Schwarz2015LC} L. A. Schwarz et al., \emph{Viral-genetic tracing of the input--output organization of a central noradrenaline circuit}, Nature \textbf{524}, 88 (2015).
\bibitem{Uematsu2017LC} A. Uematsu et al., \emph{Modular organization of the brainstem noradrenaline system coordinates opposing learning states}, Nat. Neurosci. \textbf{20}, 1602 (2017).
\bibitem{AstonJones1994vigilance} G. Aston-Jones, J. Rajkowski, P. Kubiak, and T. Alexinsky, \emph{Locus coeruleus neurons in monkey are selectively activated by attended cues in a vigilance task}, J. Neurosci. \textbf{14}, 4467 (1994).
\bibitem{Clayton2004LC} E. C. Clayton, J. Rajkowski, J. D. Cohen, and G. Aston-Jones, \emph{Phasic activation of monkey locus ceruleus neurons by simple decisions in a forced-choice task}, J. Neurosci. \textbf{24}, 9914 (2004).
\bibitem{Foote1975NE} S. L. Foote, R. Freedman, and A. P. Oliver, \emph{Effects of putative neurotransmitters on neuronal activity in monkey auditory cortex}, Brain Res. \textbf{86}, 229 (1975).
\bibitem{JepmaNieuwenhuis2011} M. Jepma and S. Nieuwenhuis, \emph{Pupil diameter predicts changes in the exploration--exploitation trade-off: evidence for the adaptive gain theory}, J. Cogn. Neurosci. \textbf{23}, 1587 (2011).
\bibitem{Tervo2014stochastic} D. G. R. Tervo et al., \emph{Behavioral variability through stochastic choice and its gating by anterior cingulate cortex}, Cell \textbf{159}, 21 (2014).
\bibitem{Usher1999LC} M. Usher, J. D. Cohen, D. Servan-Schreiber, J. Rajkowski, and G. Aston-Jones, \emph{The role of locus coeruleus in the regulation of cognitive performance}, Science \textbf{283}, 549 (1999).

\bibitem{Mesulam1983} M.-M. Mesulam, E. J. Mufson, B. H. Wainer, and A. I. Levey, \emph{Central cholinergic pathways in the rat: an overview based on an alternative nomenclature (Ch1--Ch6)}, Neuroscience \textbf{10}, 1185 (1983).
\bibitem{Marrosu1995} F. Marrosu et al., \emph{Microdialysis measurement of cortical and hippocampal acetylcholine release during sleep-wake cycle in freely moving cats}, Brain Res. \textbf{671}, 329 (1995).
\bibitem{Picciotto2012} M. R. Picciotto, M. J. Higley, and Y. S. Mineur, \emph{Acetylcholine as a neuromodulator: cholinergic signaling shapes nervous system function and behavior}, Neuron \textbf{76}, 116 (2012).
\bibitem{HasselmoSchnell1994} M. E. Hasselmo and E. Schnell, \emph{Laminar selectivity of the cholinergic suppression of synaptic transmission in rat hippocampal region CA1: computational modeling and brain slice physiology}, J. Neurosci. \textbf{14}, 3898 (1994).
\bibitem{Gil1997} Z. Gil, B. W. Connors, and Y. Amitai, \emph{Differential regulation of neocortical synapses by neuromodulators and activity}, Neuron \textbf{19}, 679 (1997).
\bibitem{KilgardMerzenich1998} M. P. Kilgard and M. M. Merzenich, \emph{Cortical map reorganization enabled by nucleus basalis activity}, Science \textbf{279}, 1714 (1998).
\bibitem{Atri2004} A. Atri et al., \emph{Blockade of central cholinergic receptors impairs new learning and increases proactive interference in a word paired-associate memory task}, Behav. Neurosci. \textbf{118}, 223 (2004).
\bibitem{Girardeau2009} G. Girardeau, K. Benchenane, S. I. Wiener, G. Buzs\'aki, and M. B. Zugaro, \emph{Selective suppression of hippocampal ripples impairs spatial memory}, Nat. Neurosci. \textbf{12}, 1222 (2009).

\input{supplement/bib_10}
\bibitem{CopenhagenJahr1989} D. R. Copenhagen and C. E. Jahr, \emph{Release of endogenous excitatory amino acids from turtle photoreceptors}, Nature \textbf{341}, 536 (1989).
\bibitem{GlowatzkiFuchs2002} E. Glowatzki and P. A. Fuchs, \emph{Transmitter release at the hair cell ribbon synapse}, Nat. Neurosci. \textbf{5}, 147 (2002).
\bibitem{Tinsley2016} J. N. Tinsley et al., \emph{Direct detection of a single photon by humans}, Nat. Commun. \textbf{7}, 12172 (2016).
\bibitem{Sellick1982} P. M. Sellick, R. Patuzzi, and B. M. Johnstone, \emph{Measurement of basilar membrane motion in the guinea pig using the M\"ossbauer technique}, J. Acoust. Soc. Am. \textbf{72}, 131 (1982).
\bibitem{Ruggero1997} M. A. Ruggero, N. C. Rich, A. Recio, S. S. Narayan, and L. Robles, \emph{Basilar-membrane responses to tones at the base of the chinchilla cochlea}, J. Acoust. Soc. Am. \textbf{101}, 2151 (1997).
\bibitem{NormannPerlman1979} R. A. Normann and I. Perlman, \emph{The effects of background illumination on the photoresponses of red and green cones}, J. Physiol. \textbf{286}, 491 (1979).
\bibitem{Laughlin1981} S. Laughlin, \emph{A simple coding procedure enhances a neuron's information capacity}, Z. Naturforsch. C \textbf{36}, 910 (1981).
\bibitem{CurcioAllen1990} C. A. Curcio and K. A. Allen, \emph{Topography of ganglion cells in human retina}, J. Comp. Neurol. \textbf{300}, 5 (1990).
\bibitem{Greenwood1990} D. D. Greenwood, \emph{A cochlear frequency-position function for several species --- 29 years later}, J. Acoust. Soc. Am. \textbf{87}, 2592 (1990).
\bibitem{Eguiluz2000} V. M. Egu\'iluz, M. Ospeck, Y. Choe, A. J. Hudspeth, and M. O. Magnasco, \emph{Essential nonlinearities in hearing}, Phys. Rev. Lett. \textbf{84}, 5232 (2000).
\bibitem{PalmerRussell1986} A. R. Palmer and I. J. Russell, \emph{Phase-locking in the cochlear nerve of the guinea-pig and its relation to the receptor potential of inner hair-cells}, Hear. Res. \textbf{24}, 1 (1986).
\bibitem{Malnic1999} B. Malnic, J. Hirono, T. Sato, and L. B. Buck, \emph{Combinatorial receptor codes for odors}, Cell \textbf{96}, 713 (1999).
\bibitem{Finger2005} T. E. Finger et al., \emph{ATP signaling is crucial for communication from taste buds to gustatory nerves}, Science \textbf{310}, 1495 (2005).
\bibitem{Huang2005} Y.-J. Huang et al., \emph{Mouse taste buds use serotonin as a neurotransmitter}, J. Neurosci. \textbf{25}, 843 (2005).
\bibitem{JohanssonVallbo1979} R. S. Johansson and A. B. Vallbo, \emph{Tactile sensibility in the human hand: relative and absolute densities of four types of mechanoreceptive units in glabrous skin}, J. Physiol. \textbf{286}, 283 (1979).
\bibitem{McKemy2002} D. D. McKemy, W. M. Neuhausser, and D. Julius, \emph{Identification of a cold receptor reveals a general role for TRP channels in thermosensation}, Nature \textbf{416}, 52 (2002).

\bibitem{Schmolesky1998} M. T. Schmolesky et al., \emph{Signal timing across the macaque visual system}, J. Neurophysiol. \textbf{79}, 3272 (1998).
\bibitem{Lampl2004} I. Lampl, D. Ferster, T. Poggio, and M. Riesenhuber, \emph{Intracellular measurements of spatial integration and the MAX operation in complex cells of the cat primary visual cortex}, J. Neurophysiol. \textbf{92}, 2704 (2004).
\bibitem{KobatakeTanaka1994} E. Kobatake and K. Tanaka, \emph{Neuronal selectivities to complex object features in the ventral visual pathway of the macaque cerebral cortex}, J. Neurophysiol. \textbf{71}, 856 (1994).
\bibitem{Albright1984} T. D. Albright, \emph{Direction and orientation selectivity of neurons in visual area MT of the macaque}, J. Neurophysiol. \textbf{52}, 1106 (1984).
\bibitem{Goodfellow2015} I. J. Goodfellow, J. Shlens, and C. Szegedy, \emph{Explaining and harnessing adversarial examples}, in \emph{Proc. ICLR} (2015), arXiv:1412.6572.
\bibitem{Elsayed2018} G. F. Elsayed et al., \emph{Adversarial examples that fool both computer vision and time-limited humans}, in \emph{Advances in Neural Information Processing Systems 31} (2018), arXiv:1802.08195.
\bibitem{Thompson1982} P. Thompson, \emph{Perceived rate of movement depends on contrast}, Vision Res. \textbf{22}, 377 (1982).
\bibitem{BarlowHill1963} H. B. Barlow and R. M. Hill, \emph{Evidence for a physiological explanation of the waterfall phenomenon and figural after-effects}, Nature \textbf{200}, 1345 (1963).
\bibitem{EaglemanSejnowski2000} D. M. Eagleman and T. J. Sejnowski, \emph{Motion integration and postdiction in visual awareness}, Science \textbf{287}, 2036 (2000).

\bibitem{Feinstein1955mu} B. Feinstein, B. Lindeg{\aa}rd, E. Nyman, and G. Wohlfart, \emph{Morphologic studies of motor units in normal human muscles}, Acta Anat. \textbf{23}, 127 (1955).
\bibitem{MilnerBrown1973rec} H. S. Milner-Brown, R. B. Stein, and R. Yemm, \emph{The orderly recruitment of human motor units during voluntary isometric contractions}, J. Physiol. \textbf{230}, 359 (1973).
\bibitem{Jones2002sdn} K. E. Jones, A. F. de C. Hamilton, and D. M. Wolpert, \emph{Sources of signal-dependent noise during isometric force production}, J. Neurophysiol. \textbf{88}, 1533 (2002).
\bibitem{GomiOsu1998stiff} H. Gomi and R. Osu, \emph{Task-dependent viscoelasticity of human multijoint arm and its spatial characteristics for interaction with environments}, J. Neurosci. \textbf{18}, 8965 (1998).
\bibitem{Jankowska1972Ia} E. Jankowska and W. J. Roberts, \emph{Synaptic actions of single interneurones mediating reciprocal Ia inhibition of motoneurones}, J. Physiol. \textbf{222}, 623 (1972).
\bibitem{Pruszynski2008rapid} J. A. Pruszynski, I. Kurtzer, and S. H. Scott, \emph{Rapid motor responses are appropriately tuned to the metrics of a visuospatial task}, J. Neurophysiol. \textbf{100}, 224 (2008).
\bibitem{SaundersKnill2003vis} J. A. Saunders and D. C. Knill, \emph{Humans use continuous visual feedback from the hand to control fast reaching movements}, Exp. Brain Res. \textbf{152}, 341 (2003).
\bibitem{FlanaganWing1997grip} J. R. Flanagan and A. M. Wing, \emph{The role of internal models in motion planning and control: evidence from grip force adjustments during movements of hand-held loads}, J. Neurosci. \textbf{17}, 1519 (1997).
\bibitem{Matsuoka1985osc} K. Matsuoka, \emph{Sustained oscillations generated by mutually inhibiting neurons with adaptation}, Biol. Cybern. \textbf{52}, 367 (1985).

\bibitem{Treede1995heat} R.-D. Treede, R. A. Meyer, S. N. Raja, and J. N. Campbell, \emph{Evidence for two different heat transduction mechanisms in nociceptive primary afferents innervating monkey skin}, J. Physiol. \textbf{483}, 747 (1995).
\bibitem{Weidner1999cfib} C. Weidner et al., \emph{Functional attributes discriminating mechano-insensitive and mechano-responsive C nociceptors in human skin}, J. Neurosci. \textbf{19}, 10184 (1999).
\bibitem{CampbellLaMotte1983first} J. N. Campbell and R. H. LaMotte, \emph{Latency to detection of first pain}, Brain Res. \textbf{266}, 203 (1983).
\bibitem{LaMotte1982hyper} R. H. LaMotte, J. G. Thalhammer, H. E. Torebj{\"o}rk, and C. J. Robinson, \emph{Peripheral neural mechanisms of cutaneous hyperalgesia following mild injury by heat}, J. Neurosci. \textbf{2}, 765 (1982).
\bibitem{Mancini2014touch} F. Mancini, T. Nash, G. D. Iannetti, and P. Haggard, \emph{Pain relief by touch: a quantitative approach}, Pain \textbf{155}, 635 (2014).
\bibitem{Fields1983rvm} H. L. Fields, J. Bry, I. Hentall, and G. Zorman, \emph{The activity of neurons in the rostral medulla of the rat during withdrawal from noxious heat}, J. Neurosci. \textbf{3}, 2545 (1983).
\bibitem{Crombez1996disrupt} G. Crombez, C. Eccleston, F. Baeyens, and P. Eelen, \emph{The disruptive nature of pain: an experimental investigation}, Behav. Res. Ther. \textbf{34}, 911 (1996).
\bibitem{Schouenborg1995wr} J. Schouenborg, H.-R. Weng, J. Kalliom{\"a}ki, and H. Holmberg, \emph{A survey of spinal dorsal horn neurones encoding the spatial organization of withdrawal reflexes in the rat}, Exp. Brain Res. \textbf{106}, 19 (1995).

\bibitem{Andersen1992cereb} B. B. Andersen, L. Korbo, and B. Pakkenberg, \emph{A quantitative study of the human cerebellum with unbiased stereological techniques}, J. Comp. Neurol. \textbf{326}, 549 (1992).
\bibitem{ShadmehrMussaIvaldi1994} R. Shadmehr and F. A. Mussa-Ivaldi, \emph{Adaptive representation of dynamics during learning of a motor task}, J. Neurosci. \textbf{14}, 3208 (1994).

\bibitem{JoseCollison1970ihr} A. D. Jose and D. Collison, \emph{The normal range and determinants of the intrinsic heart rate in man}, Cardiovasc. Res. \textbf{4}, 160 (1970).
\bibitem{Logothetis1987gbg} D. E. Logothetis, Y. Kurachi, J. Galper, E. J. Neer, and D. E. Clapham, \emph{The $\beta\gamma$ subunits of GTP-binding proteins activate the muscarinic K$^+$ channel in heart}, Nature \textbf{325}, 321 (1987).
\bibitem{KellerWood1984fb} M. E. Keller-Wood and M. F. Dallman, \emph{Corticosteroid inhibition of ACTH secretion}, Endocr. Rev. \textbf{5}, 1 (1984).
\bibitem{Walker2012glc} J. J. Walker et al., \emph{The origin of glucocorticoid hormone oscillations}, PLoS Biol. \textbf{10}, e1001341 (2012).
\bibitem{Wildt1981gnrh} L. Wildt et al., \emph{Frequency and amplitude of gonadotropin-releasing hormone stimulation and gonadotropin secretion in the rhesus monkey}, Endocrinology \textbf{109}, 376 (1981).
\bibitem{Welsh2010scn} D. K. Welsh, J. S. Takahashi, and S. A. Kay, \emph{Suprachiasmatic nucleus: cell autonomy and network properties}, Annu. Rev. Physiol. \textbf{72}, 551 (2010).
\bibitem{Khalsa2003prc} S. B. S. Khalsa, M. E. Jewett, C. Cajochen, and C. A. Czeisler, \emph{A phase response curve to single bright light pulses in human subjects}, J. Physiol. \textbf{549}, 945 (2003).
\bibitem{Sack1992blind} R. L. Sack, A. J. Lewy, M. L. Blood, L. D. Keith, and H. Nakagawa, \emph{Circadian rhythm abnormalities in totally blind people: incidence and clinical significance}, J. Clin. Endocrinol. Metab. \textbf{75}, 127 (1992).

\bibitem{Henze2000extra} D. A. Henze et al., \emph{Intracellular features predicted by extracellular recordings in the hippocampus in vivo}, J. Neurophysiol. \textbf{84}, 390 (2000).
\bibitem{Histed2009stim} M. H. Histed, V. Bonin, and R. C. Reid, \emph{Direct activation of sparse, distributed populations of cortical neurons by electrical microstimulation}, Neuron \textbf{63}, 508 (2009).
\bibitem{Orsborn2014coad} A. L. Orsborn et al., \emph{Closed-loop decoder adaptation shapes neural plasticity for skillful neuroprosthetic control}, Neuron \textbf{82}, 1380 (2014).

\bibitem{HuVervaekeStorm2002} H. Hu, K. Vervaeke, and J. F. Storm, \emph{Two forms of electrical resonance at theta frequencies, generated by M-current, h-current and persistent Na$^+$ current in rat hippocampal pyramidal cells}, J. Physiol. \textbf{545}, 783 (2002).
\bibitem{Mauro1970} A. Mauro, F. Conti, F. Dodge, and R. Schor, \emph{Subthreshold behavior and phenomenological impedance of the squid giant axon}, J. Gen. Physiol. \textbf{55}, 497 (1970).
\bibitem{Aksay2001} E. Aksay et al., \emph{In vivo intracellular recording and perturbation of persistent activity in a neural integrator}, Nat. Neurosci. \textbf{4}, 184 (2001).
\bibitem{Major2004} G. Major et al., \emph{Plasticity and tuning of the time course of analog persistent firing in a neural integrator}, Proc. Natl. Acad. Sci. USA \textbf{101}, 7745 (2004).
\bibitem{Tateno2004} T. Tateno, A. Harsch, and H. P. C. Robinson, \emph{Threshold firing frequency-current relationships of neurons in rat somatosensory cortex: type 1 and type 2 dynamics}, J. Neurophysiol. \textbf{92}, 2283 (2004).
\bibitem{Marder2012} E. Marder, \emph{Neuromodulation of neuronal circuits: back to the future}, Neuron \textbf{76}, 1 (2012).
\bibitem{McGintyHarper1976} D. J. McGinty and R. M. Harper, \emph{Dorsal raphe neurons: depression of firing during sleep in cats}, Brain Res. \textbf{101}, 569 (1976).

\bibitem{Gentet2000} L. J. Gentet, G. J. Stuart, and J. D. Clements, \emph{Direct measurement of specific membrane capacitance in neurons}, Biophys. J. \textbf{79}, 314 (2000).
\bibitem{Borst1995} J. G. G. Borst, F. Helmchen, and B. Sakmann, \emph{Pre- and postsynaptic whole-cell recordings in the medial nucleus of the trapezoid body of the rat}, J. Physiol. \textbf{489}, 825 (1995).
\bibitem{AmirDevor2003} R. Amir and M. Devor, \emph{Electrical excitability of the soma of sensory neurons is required for spike invasion of the soma, but not for through-conduction}, Biophys. J. \textbf{84}, 2181 (2003).
\bibitem{Chadderton2004} P. Chadderton, T. W. Margrie, and M. H\"ausser, \emph{Integration of quanta in cerebellar granule cells during sensory processing}, Nature \textbf{428}, 856 (2004).
\bibitem{LitwinKumar2017} A. Litwin-Kumar et al., \emph{Optimal degrees of synaptic connectivity}, Neuron \textbf{93}, 1153 (2017).
\bibitem{NapperHarvey1988} R. M. A. Napper and R. J. Harvey, \emph{Number of parallel fiber synapses on an individual Purkinje cell in the cerebellum of the rat}, J. Comp. Neurol. \textbf{274}, 168 (1988).
\bibitem{Megias2001} M. Meg\'{\i}as, Z. Emri, T. F. Freund, and A. I. Guly\'as, \emph{Total number and distribution of inhibitory and excitatory synapses on hippocampal CA1 pyramidal cells}, Neuroscience \textbf{102}, 527 (2001).
\bibitem{Polsky2004} A. Polsky, B. W. Mel, and J. Schiller, \emph{Computational subunits in thin dendrites of pyramidal cells}, Nat. Neurosci. \textbf{7}, 621 (2004).
\bibitem{Brunel2004} N. Brunel et al., \emph{Optimal information storage and the distribution of synaptic weights: perceptron versus Purkinje cell}, Neuron \textbf{43}, 745 (2004).

\bibitem{Vyazovskiy2009} V. V. Vyazovskiy et al., \emph{Cortical firing and sleep homeostasis}, Neuron \textbf{63}, 865 (2009).
\bibitem{Daan1984} S. Daan, D. G. Beersma, and A. A. Borb\'ely, \emph{Timing of human sleep: recovery process gated by a circadian pacemaker}, Am. J. Physiol. \textbf{246}, R161 (1984).
\bibitem{Borbely1981} A. A. Borb\'ely et al., \emph{Sleep deprivation: effect on sleep stages and EEG power density in man}, Electroencephalogr. Clin. Neurophysiol. \textbf{51}, 483 (1981).
\bibitem{SteriadeAmzica1993} M. Steriade, F. Amzica, and A. Nu\~nez, \emph{Cholinergic and noradrenergic modulation of the slow ($\approx0.3$ Hz) oscillation in neocortical cells}, J. Neurophysiol. \textbf{70}, 1385 (1993).
\bibitem{vonKrosigk1993} M. von Krosigk, T. Bal, and D. A. McCormick, \emph{Cellular mechanisms of a synchronized oscillation in the thalamus}, Science \textbf{261}, 361 (1993).
\bibitem{LeeWilson2002} A. K. Lee and M. A. Wilson, \emph{Memory of sequential experience in the hippocampus during slow wave sleep}, Neuron \textbf{36}, 1183 (2002).
\bibitem{Euston2007} D. R. Euston, M. Tatsuno, and B. L. McNaughton, \emph{Fast-forward playback of recent memory sequences in prefrontal cortex during sleep}, Science \textbf{318}, 1147 (2007).
\bibitem{Hobson2009} J. A. Hobson, \emph{REM sleep and dreaming: towards a theory of protoconsciousness}, Nat. Rev. Neurosci. \textbf{10}, 803 (2009).

\bibitem{CohenSegal2011} D. Cohen and M. Segal, \emph{Network bursts in hippocampal microcultures are terminated by exhaustion of vesicle pools}, J. Neurophysiol. \textbf{106}, 2314 (2011).
\bibitem{Tsodyks2000} M. Tsodyks, A. Uziel, and H. Markram, \emph{Synchrony generation in recurrent networks with frequency-dependent synapses}, J. Neurosci. \textbf{20}, RC50 (2000).
\bibitem{Benson1994} D. L. Benson, F. H. Watkins, O. Steward, and G. Banker, \emph{Characterization of GABAergic neurons in hippocampal cell cultures}, J. Neurocytol. \textbf{23}, 279 (1994).

\input{supplement/bib_22}
\end{thebibliography}


\end{document}
